% Italian Learning Lab — book edition. Generated by book/generate.mjs; edit the
% generator or this preamble, not italian-learning-lab.tex.
\documentclass[11pt,a4paper,twoside,openright]{book}

\usepackage{fontspec}
\setmainfont{TeX Gyre Pagella}
\setsansfont{TeX Gyre Heros}[Scale=MatchLowercase]
\setmonofont{Latin Modern Mono}
\newfontfamily\fallbackfont{DejaVu Serif}[Scale=MatchLowercase]
\usepackage[english]{babel}
\babelprovide[import]{italian}
\usepackage{newunicodechar}
% Characters missing from TeX Gyre Pagella (IPA symbols), printed with DejaVu Serif.
\newunicodechar{ɛ}{{\fallbackfont ɛ}}
\newunicodechar{ɔ}{{\fallbackfont ɔ}}
\newunicodechar{ʎ}{{\fallbackfont ʎ}}
\newunicodechar{ɲ}{{\fallbackfont ɲ}}
\newunicodechar{ʃ}{{\fallbackfont ʃ}}
\newunicodechar{ʒ}{{\fallbackfont ʒ}}
\newunicodechar{ʧ}{{\fallbackfont ʧ}}
\newunicodechar{ʤ}{{\fallbackfont ʤ}}
\newunicodechar{ˈ}{{\fallbackfont ˈ}}
\newunicodechar{ː}{{\fallbackfont ː}}


\usepackage[a4paper,inner=26mm,outer=22mm,top=24mm,bottom=26mm,headheight=14pt]{geometry}
\usepackage{microtype}
\usepackage[dvipsnames,table]{xcolor}
\usepackage{array,booktabs,longtable,tabularx,xltabular}
\usepackage[inline]{enumitem}
\usepackage{multicol}
\usepackage[most]{tcolorbox}
\usepackage{titlesec}
\usepackage{fancyhdr}
\usepackage{imakeidx}
\makeindex[title=Index of Italian words,columns=3,intoc]
\usepackage[hidelinks,bookmarksnumbered,pdfusetitle]{hyperref}
\usepackage{bookmark}

% Colours (the website's palette)
\definecolor{accent}{HTML}{1D6A50}
\definecolor{rulecol}{HTML}{1D6A50}
\definecolor{tendcol}{HTML}{8A6A12}
\definecolor{exccol}{HTML}{A64B32}
\definecolor{tipcol}{HTML}{285A8F}
\definecolor{muted}{HTML}{5F646C}
\definecolor{paper}{HTML}{F8F6F1}

% Headings
\titleformat{\part}[display]{\centering\sffamily\bfseries\Huge\color{accent}}{\partname~\thepart}{18pt}{\Huge}
\titleformat{\chapter}[display]{\sffamily\bfseries\huge}{\color{accent}\large\MakeUppercase{\chaptertitlename}~\thechapter}{6pt}{\Huge}
\titlespacing*{\chapter}{0pt}{0pt}{24pt}
\titleformat{\section}{\sffamily\bfseries\Large}{\color{accent}\thesection}{0.8em}{}
\titleformat{\subsection}{\sffamily\bfseries\large}{\color{accent}\thesubsection}{0.7em}{}
\titleformat{\subsubsection}{\sffamily\bfseries\normalsize}{}{0em}{}
\setcounter{secnumdepth}{2}
\setcounter{tocdepth}{1}

\pagestyle{fancy}
\fancyhf{}
\fancyhead[LE]{\sffamily\small\leftmark}
\fancyhead[RO]{\sffamily\small\rightmark}
\fancyfoot[C]{\sffamily\small\thepage}
\renewcommand{\headrulewidth}{0.4pt}
\fancypagestyle{plain}{\fancyhf{}\fancyfoot[C]{\sffamily\small\thepage}\renewcommand{\headrulewidth}{0pt}}

\setlength{\parindent}{0pt}
\setlength{\parskip}{4pt plus 2pt}
\setlength{\emergencystretch}{3em}
\raggedbottom

% Tables
\newcolumntype{L}{>{\raggedright\arraybackslash}X}
\renewcommand{\arraystretch}{1.15}
\newcommand{\tablecaption}[1]{\par\smallskip\noindent{\sffamily\small\bfseries #1}\par\nopagebreak\vspace{1pt}}

% Rule boxes: rule / tendency / exception / tip
\newtcolorbox{notebox}[1]{breakable,enhanced,colback=#1!6!white,colframe=#1,boxrule=0pt,leftrule=2.5pt,arc=0pt,outer arc=0pt,left=7pt,right=7pt,top=4pt,bottom=4pt,before skip=8pt,after skip=8pt}
\newcommand{\notelabel}[3]{{\sffamily\small\bfseries\color{#1}#2}\if\relax\detokenize{#3}\relax\else{\sffamily\small\bfseries\ ---\ #3}\fi\par\nopagebreak}
\newtcolorbox{summarybox}{breakable,enhanced,colback=paper,colframe=accent!50,boxrule=0.5pt,arc=2pt,left=6pt,right=6pt,top=4pt,bottom=4pt,before upper={{\sffamily\bfseries\color{accent}Summary}\par\setlist[itemize]{topsep=2pt,itemsep=1pt}}}

% Verb tables
\newcommand{\irr}[1]{{\color{exccol}\bfseries #1}}
\newcommand{\hlending}[1]{{\color{accent}\bfseries #1}}
\newcommand{\paradigmbox}[3]{{\small\sffamily\bfseries #1}\par\nopagebreak\vspace{2pt}{\small #2}\par\if\relax\detokenize{#3}\relax\else{\footnotesize\color{muted}#3\par}\fi}
\newcommand{\moodheading}[1]{\par\medskip\noindent{\sffamily\bfseries\small\color{muted}\MakeUppercase{#1}}\par\nopagebreak}

% Lessons
\newcommand{\lessonref}[2]{#2 (\S\ref{les:#1})}
\newcommand{\lessonmeta}[2]{\noindent{\sffamily\small\color{accent}Approximate level: #1}\par{\itshape #2}\par\medskip}
\newcommand{\moduleintro}[3]{\noindent{\sffamily\color{accent}Approximate level: #1}\par\smallskip{\large #2}\par\medskip\noindent{\small\textbf{Lessons:} #3}\par\bigskip}
\newenvironment{objectives}{\par\noindent{\sffamily\bfseries Objectives}\begin{itemize}[topsep=2pt,itemsep=1pt]}{\end{itemize}}
\newlist{examples}{enumerate}{1}
\setlist[examples]{label=\arabic*.,leftmargin=2em,itemsep=4pt}
\newlist{exercises}{enumerate}{1}
\setlist[exercises]{label=\textbf{\arabic*.},leftmargin=2em,itemsep=10pt,topsep=4pt}
\newlist{options}{enumerate}{1}
\setlist[options]{label=(\alph*),leftmargin=2.2em,itemsep=0pt,topsep=2pt,parsep=0pt}
\newlist{answers}{enumerate}{1}
\setlist[answers]{label=\textbf{\arabic*.},leftmargin=2em,itemsep=3pt,topsep=2pt}
\newlist{sentences}{enumerate}{1}
\setlist[sentences]{label=\arabic*.,leftmargin=2em,itemsep=6pt}
\newlist{glosslist}{itemize}{1}
\setlist[glosslist]{label={},leftmargin=1.2em,itemindent=-1.2em,itemsep=3pt,topsep=2pt}
\newcommand{\relatedwords}[1]{\par\medskip\noindent{\sffamily\small\bfseries Related vocabulary}\par\nopagebreak{\small #1}\par}
\newcommand{\applyit}[1]{\par\smallskip\noindent{\sffamily\small\bfseries Apply it}\par\nopagebreak{\small #1}\par}
\newcommand{\answerheading}[1]{\par\medskip\noindent{\sffamily\bfseries\small #1}\par\nopagebreak}

% Exercises on paper
\newcommand{\blank}{\underline{\hspace{2.6cm}}}
\newcommand{\clozeblank}[1]{\underline{\hspace{2.4cm}}\textsuperscript{\,(#1)}}
\newcommand{\blankcell}{\underline{\hspace{4.5cm}}}
\newcommand{\answerline}{\par\vspace{6pt}\noindent\rule{0.75\linewidth}{0.4pt}\par}
\newcommand{\token}[1]{\fbox{\strut\itshape #1}}

% Reader
\newcommand{\storymeta}[3]{\noindent{\large\itshape #1}\quad{\sffamily\small\color{accent}Level #2}\par\smallskip{\small #3}\par}
\newenvironment{storytext}{\par\begingroup\large\setlength{\parskip}{8pt}\linespread{1.15}\selectfont}{\par\endgroup}
\newcommand{\sentnum}[1]{\textsuperscript{\color{muted}\tiny #1}}

% Glossary
\newcommand{\glossentry}[5]{\item #1 {\small\color{muted}(#2)} --- #3.\newline{\small\itshape #4}{\small\ --- #5}}
\newcommand{\glossnote}[1]{\newline{\footnotesize\color{muted}#1}}

\hypersetup{pdftitle={Italian Learning Lab: Grammar, Reader and Workbook},pdfauthor={Italian Learning Lab contributors}}

\begin{document}
\frontmatter
\begin{titlepage}
\centering
\vspace*{30mm}
{\sffamily\bfseries\color{accent}\fontsize{40}{44}\selectfont Italian Learning Lab\par}
\vspace{10mm}
{\sffamily\LARGE Grammar, Reader and Workbook\par}
\vspace{6mm}
{\large From the first sounds to the complete verb system\par}
\vspace{4mm}
{\large\color{muted} Levels A0 to about B1--B2\par}
\vfill
{\large Companion to the Italian Learning Lab website\par}
\vspace{4mm}
{\color{muted}Edition generated from course data version 1.0.0\par}
\vspace{20mm}
\end{titlepage}

\thispagestyle{empty}
\vspace*{\fill}
{\small
\textbf{Italian Learning Lab: Grammar, Reader and Workbook}

Written by the Italian Learning Lab contributors. The educational content (lessons, explanations, examples, exercises, vocabulary, stories and rhymes) is released under the Creative Commons Attribution 4.0 International licence (CC BY 4.0): you may copy, adapt and share it, including commercially, if you give appropriate credit. The software that generates this book is released under the MIT licence.

This book is generated automatically from the same data and the same grammar engine as the website, so the two always contain the same lessons, rules, exercises, answers, vocabulary and texts. Verb tables were produced by the conjugation engine and checked by its automated tests.

\textbf{Verification.} No native speaker, Italian teacher or linguist has reviewed this course. Statements about the passato remoto and its regional use, past-participle agreement, the auxiliary of modal verbs, the periodo ipotetico and the choice between congiuntivo and indicativo were compared with the reference works listed in Appendix~\ref{sources}; everything else is the authors' own work based on standard grammar. Appendix~\ref{sources} records, topic by topic, what has and has not been checked. Corrections are welcome.

Levels (A0, A1, A2, B1, B2) are approximate and describe intended difficulty only. This book is not affiliated with, or certified by, any examination body.
}
\clearpage

\chapter{How to use this book}

This book contains everything in the Italian Learning Lab website except listening and speaking: there is no audio, no read-aloud and no pronunciation practice. It can be used on its own, as a printed textbook and workbook, or next to the website. Where the website is interactive, the book prints the result instead: every story sentence comes with its translation and a verb analysis, every verb in the course is conjugated in full or listed with its principal parts, and every exercise has an answer and an explanation in the answer key.

\section*{The teaching principle}

The course teaches accurate, standard Italian from the first lesson. Difficulty is controlled by the choice of vocabulary, the length of sentences and the reading material --- not by leaving out grammatical distinctions or by stating rules that later have to be unlearned. Exceptions, alternative forms, differences of register and regional differences are stated where they occur. Every rule box says what kind of statement it is:

\begin{notebox}{rulecol}\notelabel{rulecol}{Rule}{}A grammatical rule of standard Italian. Departing from it is an error in the standard language.\end{notebox}
\begin{notebox}{tendcol}\notelabel{tendcol}{Tendency}{}A strong regularity, a stylistic preference or a difference of register or region. Alternatives exist and are not errors.\end{notebox}
\begin{notebox}{exccol}\notelabel{exccol}{Exception}{}A form or word that does not follow the rule just stated.\end{notebox}
\begin{notebox}{tipcol}\notelabel{tipcol}{Tip}{}A practical way to remember or apply something.\end{notebox}

\section*{What is in the book}

\begin{description}[leftmargin=0pt,itemsep=3pt]
\item[Part I, Course.] 20 chapters (one per module) with 49 lessons, from spelling and pronunciation rules to the congiuntivo, the sequence of tenses, indirect speech and the passive. Each lesson has objectives, explanations, verb tables, worked examples, exercises, a summary and its related vocabulary.
\item[Part II, Reader.] 10 graded stories and 4 rhymes. Each story is printed first as continuous text, then sentence by sentence with a translation, notes and an analysis of every verb.
\item[Part III, Practice.] A diagnostic test to find your starting point and mixed revision sets.
\item[Part IV, Grammar reference.] An overview of the whole verb system, the sequence of tenses, the types of hypothetical sentence, the articulated prepositions, full conjugations of the auxiliaries, the regular models and every irregular verb in the course, and the principal parts of all 144 verbs.
\item[Appendices.] The answer key for all 437 exercises; the 503 words of the course grouped by topic; sources and the verification register; an alphabetical index of Italian words.
\end{description}

\section*{Suggested ways of working}

\begin{itemize}
\item \textbf{In order.} Work through Part I chapter by chapter. After each lesson, do its exercises in writing, check them in the answer key, and read the story or rhyme suggested under \emph{Apply it}. Every lesson can be revisited at any time; nothing depends on finishing the previous one, but later lessons assume the earlier ones.
\item \textbf{From a placement.} Take the diagnostic test in Chapter~\ref{diag}. Start with the first chapter in which you made a mistake.
\item \textbf{As a reference.} Use Part IV and the index. Every tense in the overview table points to the lesson that explains it.
\item \textbf{As a reader.} The stories in Part II can be read before the grammar they use: the verb analysis explains every form and gives the section where it is taught.
\end{itemize}

\section*{Conventions}

Italian words in English text are in \textit{italics}. In verb tables the persons are \textit{io, tu, lui/lei, noi, voi, loro}; the formal \textit{Lei} and \textit{Loro} use the third-person forms. Endings are printed in \hlending{green} where the lesson asks you to notice them; forms that differ from the regular pattern of their family are printed in \irr{red-brown}. Where two forms are separated by a slash (\textit{credei/credetti}), both are standard. A cross-reference such as \S\ref{les:passato-remoto} points to a section of the book.

\section*{Exercises and answers}

Exercises are designed to be written out. Blanks are shown as lines; in matching exercises the right-hand column has been put in a different order; in word-order exercises the words are given in alphabetical order. The answer key gives every answer that the website accepts, the explanation, and, where the website anticipates a typical mistake, what is wrong with it. When checking your answers, accents, articles, endings and double consonants count; capital letters and final punctuation do not.

\section*{Reviewing vocabulary on paper}

The website schedules vocabulary with a Leitner system, which works equally well with paper cards. Write each word on a card (Italian with its article on one side, English on the other) and keep six boxes, reviewed after 1, 2, 4, 7, 14 and 30 days. A new card starts in a \emph{learning} pile that is reviewed every day. After each review: if you did not remember the word, the card goes back to the learning pile; if you remembered it with difficulty, it stays in its box and comes back after the interval of the previous box; if you remembered it, it moves up one box; if it was effortless, it moves up two boxes. A card in the last box stays there and comes back after 30 days. A word in box 4 or higher (an interval of a week or more) can be counted as known.

\tableofcontents

\mainmatter

\part{Course}\label{part:course}

\chapter{Pronunciation and orthography}\label{mod:m01}
\moduleintro{A0}{The alphabet, vowels, c and g, sc, gn, gli, stress, double consonants, accents and the apostrophe.}{The alphabet and the vowels (\S\ref{les:alphabet-vowels}); Hard and soft c and g; sc, gn and gli (\S\ref{les:c-and-g}); Syllables, stress, double consonants and accents (\S\ref{les:stress-doubles-accents})}

\section{The alphabet and the vowels}\label{les:alphabet-vowels}
\lessonmeta{A0}{Letter names, the five vowels, and why Italian spelling is so regular.}
\begin{objectives}
\item\relax Name the 21 letters of the Italian alphabet and recognise the 5 foreign letters.
\item\relax Pronounce the five vowels as pure, unreduced sounds.
\item\relax Understand that final vowels carry grammatical information (gender and number).
\item\relax Know that h is always silent.
\end{objectives}

\subsection{The 21 letters}\label{les:alphabet-vowels:alphabet}

The native Italian alphabet has 21 letters. Five more — \textit{j, k, w, x, y} — occur only in loanwords and foreign names (\textit{il jazz, il kiwi, il weekend, lo yogurt}).\par

{\small\setlength{\LTpre}{2pt}\setlength{\LTpost}{6pt}
\tablecaption{Letter names}
\begin{longtable}{llllll}
\toprule
\textbf{Letter} & \textbf{Name} & \textbf{Letter} & \textbf{Name} & \textbf{Letter} & \textbf{Name} \\
\midrule
a & a & h & acca & q & cu \\
b & bi & i & i & r & erre \\
c & ci & l & elle & s & esse \\
d & di & m & emme & t & ti \\
e & e & n & enne & u & u \\
f & effe & o & o & v & vu (or vi) \\
g & gi & p & pi & z & zeta \\
\bottomrule
\end{longtable}}

Foreign letters: \textit{j} (i lunga), \textit{k} (cappa), \textit{w} (doppia vu), \textit{x} (ics), \textit{y} (ipsilon or i greca).\par

\begin{notebox}{rulecol}\notelabel{rulecol}{Rule}{The letter h is always silent}
\textit{h} is never pronounced. It has two jobs: inside \textit{ch} and \textit{gh} it keeps \textit{c} and \textit{g} hard (see the next lesson), and at the start of four forms of \textit{avere} — \textit{ho, hai, ha, hanno} — it distinguishes them in writing from \textit{o} (or), \textit{ai} (to the), \textit{a} (to) and \textit{anno} (year).
\end{notebox}

\subsection{Five pure vowels}\label{les:alphabet-vowels:vowels}

Italian has five vowel letters: \textit{a, e, i, o, u}. Each is pronounced fully and clearly, whether stressed or not. English often weakens unstressed vowels to a neutral 'uh' sound; Italian never does.\par

{\small\setlength{\LTpre}{2pt}\setlength{\LTpost}{6pt}
\tablecaption{Approximate vowel sounds}
\begin{xltabular}{\linewidth}{lLl}
\toprule
\textbf{Vowel} & \textbf{Like English…} & \textbf{Examples} \\
\midrule
a & the a in father & casa, mamma, gatto \\
e & the e in bed (open) or a closed 'ay' without the glide & mela, bene, sera \\
i & the ee in see, but short & vino, libri, idea \\
o & the o in for (open) or a pure closed 'oh' & notte, sole, cosa \\
u & the oo in moon, but short & luna, uva, tu \\
\bottomrule
\end{xltabular}}

\begin{notebox}{tendcol}\notelabel{tendcol}{Tendency}{Open and closed e and o}
Stressed \textit{e} and \textit{o} can be open or closed (\textit{bène} versus \textit{séra}; \textit{còsa} versus \textit{sóle}). Standard Italian therefore has seven stressed vowel sounds: a, è [ɛ], é [e], i, ò [ɔ], ó [o], u. In unstressed syllables only five occur. The distribution of open and closed vowels follows Florentine usage in the standard, but it varies a great deal by region (many speakers in the North and South distribute them differently), and only a few word pairs depend on it (\textit{pèsca} 'peach' / \textit{pésca} 'fishing'; \textit{vènti} 'winds' / \textit{vénti} 'twenty'). A regional distribution is not an error in everyday speech; dictionaries mark the standard one with grave (è, ò) and acute (é, ó) accents, and you can check any word there.
\end{notebox}

\begin{notebox}{rulecol}\notelabel{rulecol}{Rule}{Final vowels matter}
Many words differ only in their final vowel: \textit{gatto} (cat) and \textit{gatti} (cats); \textit{bambino} (boy) and \textit{bambina} (girl). The final vowel often marks gender and number, so it must be pronounced clearly.
\end{notebox}

\subsection{Spelling and sound}\label{les:alphabet-vowels:spelling}

Italian spelling is highly regular: once you know a handful of rules (mostly about \textit{c}, \textit{g}, \textit{sc}, \textit{gn} and \textit{gli}), you can pronounce almost any written word correctly, and spell almost any word you hear.\par

\begin{itemize}
\item\relax Every written vowel is pronounced: \textit{idea} has three syllables (i-de-a), \textit{aereo} has four (a-e-re-o).
\item\relax \textit{u} after \textit{q} is always pronounced: \textit{quando} (kwando), \textit{quattro}.
\item\relax \textit{z} sounds like ts or dz: \textit{zio} (tsio), \textit{zaino} (dzaino) — the choice varies by word and region.
\end{itemize}

\subsection*{Worked examples}
\begin{examples}
\item\relax \textit{la casa}\newline the house\newline {\color{muted}\small Two clear a sounds: CA-sa.}
\item\relax \textit{la luna}\newline the moon\newline {\color{muted}\small u as in 'moon', but short.}
\item\relax \textit{Ho un gatto.}\newline I have a cat.\newline {\color{muted}\small ho: the h is silent — pronounced 'o'.}
\item\relax \textit{il gatto, i gatti}\newline the cat, the cats\newline {\color{muted}\small Only the final vowel (and the article) changes.}
\item\relax \textit{il bambino, la bambina}\newline the boy, the girl\newline {\color{muted}\small The final vowel marks gender.}
\item\relax \textit{l'idea}\newline the idea\newline {\color{muted}\small Three vowels, all pronounced: i-de-a.}
\end{examples}

\subsection*{Exercises}\label{les:alphabet-vowels:exercises}
{\color{muted}\small Answers: page~\pageref{ans:alphabet-vowels}.}

\begin{exercises}
\item\label{ex:alphabet-vowels-1}\relax How is the h in \textit{ho} (I have) pronounced?\par\nopagebreak\begin{options}
\item\relax Like English h
\item\relax Like the ch in Scottish loch
\item\relax It is silent
\item\relax Like k
\end{options}
\item\label{ex:alphabet-vowels-2}\relax Match each letter to its Italian name.\par\smallskip\begin{tabularx}{\linewidth}{@{}L L@{}}
1.\ h & (a)\ acca \\
2.\ z & (b)\ erre \\
3.\ q & (c)\ gi \\
4.\ r & (d)\ zeta \\
5.\ g & (e)\ cu \\
\end{tabularx}
\item\label{ex:alphabet-vowels-3}\relax Which of these letters appears only in loanwords?\par\nopagebreak\begin{options}
\item\relax z
\item\relax q
\item\relax k
\item\relax v
\end{options}
\item\label{ex:alphabet-vowels-4}\relax \textit{gatto} and \textit{gatti} differ only in the final vowel. What does that vowel show here?\par\nopagebreak\begin{options}
\item\relax Formality
\item\relax Number (singular or plural)
\item\relax Tense
\item\relax Nothing — it is optional
\end{options}
\item\label{ex:alphabet-vowels-5}\relax How many syllables does \textit{idea} have?\par\nopagebreak\begin{options}
\item\relax 1
\item\relax 4
\item\relax 3
\item\relax 2
\end{options}
\item\label{ex:alphabet-vowels-6}\relax Translate: \textit{la bambina}\par\answerline
\end{exercises}

\begin{summarybox}
\begin{itemize}
\item\relax 21 native letters; j, k, w, x, y only in loanwords.
\item\relax Five pure vowels, always fully pronounced — including at the end of words.
\item\relax h is always silent.
\item\relax Spelling is regular: learn the few rules for c, g, sc, gn, gli and you can read aloud almost anything.
\end{itemize}
\end{summarybox}

\relatedwords{\textit{la casa} house; home \textperiodcentered{} \textit{la luna} moon \textperiodcentered{} \textit{il gatto} cat \textperiodcentered{} \textit{il bambino} child; little boy \textperiodcentered{} \textit{la bambina} little girl \textperiodcentered{} \textit{avere} to have}
\applyit{Read \textit{Il gatto e la luna} (A0–A1), \S\ref{story:il-gatto-e-la-luna}.}

\section{Hard and soft c and g; sc, gn and gli}\label{les:c-and-g}
\lessonmeta{A0}{When c and g are hard or soft, the role of h and i, and the sounds sc, gn and gli.}
\begin{objectives}
\item\relax Predict whether c and g are hard (k, g) or soft (ch, j) from the following letter.
\item\relax Use ch/gh and ci/gi spellings correctly.
\item\relax Pronounce sc, gn, gli and qu.
\end{objectives}

\subsection{c and g before different vowels}\label{les:c-and-g:c-g}

\begin{notebox}{rulecol}\notelabel{rulecol}{Rule}{The basic rule}
Before \textit{a, o, u} (and consonants), \textit{c} and \textit{g} are \textbf{hard}: c = k, g = g as in 'go'. Before \textit{e, i} they are \textbf{soft}: c = ch as in 'church', g = j as in 'jam'.
\end{notebox}

{\small\setlength{\LTpre}{2pt}\setlength{\LTpost}{6pt}
\tablecaption{c and g}
\begin{xltabular}{\linewidth}{llL}
\toprule
\textbf{Spelling} & \textbf{Sound} & \textbf{Examples} \\
\midrule
ca, co, cu & k & casa, cosa, cuore \\
ce, ci & ch (church) & cena, cielo \\
che, chi & k & che, chiave, perché \\
cia, cio, ciu & ch + a/o/u & ciao, arancia, cioccolato \\
ga, go, gu & g (go) & gatto, gola, gufo \\
ge, gi & j (jam) & gelato, giro \\
ghe, ghi & g (go) & laghi, spaghetti \\
gia, gio, giu & j + a/o/u & giallo, giorno, giugno \\
\bottomrule
\end{xltabular}}

\begin{notebox}{rulecol}\notelabel{rulecol}{Rule}{h makes c and g hard before e and i}
To write the hard sound before \textit{e} or \textit{i}, insert \textit{h}: \textit{che} (ke), \textit{chi} (ki), \textit{ghe}, \textit{ghi}. This is why the plural of \textit{lago} is \textit{laghi} and \textit{tu cerchi} is spelt with h.
\end{notebox}

\begin{notebox}{rulecol}\notelabel{rulecol}{Rule}{i makes c and g soft before a, o, u}
To write the soft sound before \textit{a, o, u}, insert an \textit{i} that is not pronounced as a separate vowel: \textit{ciao} (chao), \textit{giallo} (jallo), \textit{mangio} (manjo).
\end{notebox}

\subsection{sc}\label{les:c-and-g:sc}

{\small\setlength{\LTpre}{2pt}\setlength{\LTpost}{6pt}
\tablecaption{sc}
\begin{longtable}{lll}
\toprule
\textbf{Spelling} & \textbf{Sound} & \textbf{Examples} \\
\midrule
sca, sco, scu & sk & scuola, scala, tasca \\
sce, sci & sh (ship) & pesce, scimmia, uscire \\
sche, schi & sk & scherzo, schiena \\
scia, scio, sciu & sh + a/o/u & sciarpa, lascio \\
\bottomrule
\end{longtable}}

The same logic as c: h keeps it hard, i keeps it soft. That is why \textit{uscire} has the forms \textit{esco} (sk) and \textit{esci} (sh).\par

\subsection{gn, gli and qu}\label{les:c-and-g:gn-gli}

\begin{itemize}
\item\relax \textit{gn} = the ny sound of Spanish ñ or the ni in 'onion': \textit{montagna, sogno, bagno}.
\item\relax \textit{gli} = a palatal l, close to the lli in 'million': \textit{famiglia, figlio, foglia, coniglio}. Before \textit{a, o, u} the i is not pronounced separately: \textit{foglia} = fo-lya.
\item\relax \textit{qu} = kw: \textit{quando, quattro, questo}. The u is always pronounced.
\end{itemize}

\begin{notebox}{exccol}\notelabel{exccol}{Exception}{A few hard gl words}
In a small number of learned words \textit{gl} + i is pronounced g + l, as in English: \textit{glicine} (wisteria), \textit{glicerina, glifo, anglicano, negligente} and some proper names. (\textit{globo} is hard simply because gl is followed by o, not i.) They are few, and a dictionary shows the pronunciation.
\end{notebox}

\subsection*{Worked examples}
\begin{examples}
\item\relax \textit{Ciao, come stai?}\newline Hi, how are you?\newline {\color{muted}\small ciao: soft c, the i is silent.}
\item\relax \textit{La chiave è qui.}\newline The key is here.\newline {\color{muted}\small chiave: hard c (ki-a-ve).}
\item\relax \textit{Il gelato è giallo.}\newline The ice cream is yellow.\newline {\color{muted}\small Soft g in both words.}
\item\relax \textit{Il pesce nuota.}\newline The fish swims.\newline {\color{muted}\small pesce: sce = she.}
\item\relax \textit{La mia famiglia vive in montagna.}\newline My family lives in the mountains.\newline {\color{muted}\small famiglia (gli) and montagna (gn).}
\item\relax \textit{Vado a scuola.}\newline I go to school.\newline {\color{muted}\small scuola: sc before u = sk.}
\end{examples}

\subsection*{Exercises}\label{les:c-and-g:exercises}
{\color{muted}\small Answers: page~\pageref{ans:c-and-g}.}

\begin{exercises}
\item\label{ex:c-and-g-1}\relax In which word is c pronounced like English k?\par\nopagebreak\begin{options}
\item\relax chiave
\item\relax cena
\item\relax ciao
\item\relax cielo
\end{options}
\item\label{ex:c-and-g-2}\relax In which word is g pronounced like the j in 'jam'?\par\nopagebreak\begin{options}
\item\relax gufo
\item\relax laghi
\item\relax gatto
\item\relax giorno
\end{options}
\item\label{ex:c-and-g-3}\relax \textit{il lago} (the lake) becomes plural. Which spelling keeps the hard g sound?\par\nopagebreak\begin{options}
\item\relax i laghe
\item\relax i lagi
\item\relax i laghi
\item\relax i lagii
\end{options}
\item\label{ex:c-and-g-4}\relax Match each word to its sound.\par\smallskip\begin{tabularx}{\linewidth}{@{}L L@{}}
1.\ scuola & (a)\ sk \\
2.\ pesce & (b)\ ny \\
3.\ montagna & (c)\ ly (palatal l) \\
4.\ famiglia & (d)\ sh \\
\end{tabularx}
\item\label{ex:c-and-g-5}\relax Complete the spellings to match the sounds given in brackets.\par\smallskip\begingroup\setlength{\parindent}{0pt}\clozeblank{1}iave (ki-a-ve) · \clozeblank{2}allo (jal-lo) · \clozeblank{3}uola (skwo-la) · \clozeblank{4}erzo (sker-tso)\par\endgroup
\item\label{ex:c-and-g-6}\relax How do you pronounce the \textit{u} in \textit{quando}?\par\nopagebreak\begin{options}
\item\relax It is silent: kando
\item\relax As oo in a separate syllable: ku-an-do
\item\relax As a w: kwando
\item\relax As v: kvando
\end{options}
\end{exercises}

\begin{summarybox}
\begin{itemize}
\item\relax c, g, sc are hard before a, o, u and soft before e, i.
\item\relax Insert h to keep them hard before e/i (che, laghi, scherzo).
\item\relax Insert a silent i to make them soft before a/o/u (ciao, giallo, sciarpa).
\item\relax gn = ny, gli = palatal l, qu = kw.
\end{itemize}
\end{summarybox}

\relatedwords{\textit{la chiave} key \textperiodcentered{} \textit{il gelato} ice cream \textperiodcentered{} \textit{giallo} yellow \textperiodcentered{} \textit{il pesce} fish \textperiodcentered{} \textit{la famiglia} family \textperiodcentered{} \textit{la montagna} mountain \textperiodcentered{} \textit{la scuola} school \textperiodcentered{} \textit{il lago} lake \textperiodcentered{} \textit{la foglia} leaf \textperiodcentered{} \textit{il coniglio} rabbit}
\applyit{Read \textit{La colazione} (A0–A1), \S\ref{story:la-colazione}.}

\section{Syllables, stress, double consonants and accents}\label{les:stress-doubles-accents}
\lessonmeta{A0}{Where the stress falls, why double consonants matter, when accents are written, and elision with the apostrophe.}
\begin{objectives}
\item\relax Divide words into syllables and locate the stress.
\item\relax Pronounce and spell double consonants and recognise minimal pairs.
\item\relax Know when an accent must be written, and the difference between è and e.
\item\relax Use the apostrophe for elision (l'albero, un'amica, c'è).
\end{objectives}

\subsection{Syllables and stress}\label{les:stress-doubles-accents:syllables-stress}

Italian syllables tend to end in a vowel: \textit{ca-sa, ma-ti-ta, fi-ne-stra}. A double consonant is split between two syllables (\textit{gat-to, mam-ma}), and \textit{s} + consonant goes with the following syllable (\textit{pa-sta, fi-ne-stra}).\par

\begin{notebox}{tendcol}\notelabel{tendcol}{Tendency}{Stress usually falls on the second-to-last syllable}
Most words are stressed on the penultimate syllable: \textit{CA-sa, bam-BI-no, fi-NE-stra}. Many common words, however, are stressed on the third-to-last: \textit{AL-be-ro, TA-vo-lo, MAC-chi-na, PEC-co-ra}. This is not shown in writing, so learn the stress with each new word.
\end{notebox}

\begin{notebox}{rulecol}\notelabel{rulecol}{Rule}{Final stress is always written}
If the stress falls on the last vowel of a word of two or more syllables, it is always marked with an accent: \textit{città, caffè, perché, papà, più, felicità}.
\end{notebox}

\subsection{Written accents}\label{les:stress-doubles-accents:accents}

Italian writes a grave accent (à, è, ì, ò, ù) on most stressed final vowels. A small set of words take an acute accent on closed \textit{é}: \textit{perché, poiché, né, sé}, and compounds like \textit{ventitré}.\par

{\small\setlength{\LTpre}{2pt}\setlength{\LTpost}{6pt}
\tablecaption{One-syllable words distinguished only by the accent}
\begin{xltabular}{\linewidth}{lllL}
\toprule
\textbf{With accent} & \textbf{Meaning} & \textbf{Without accent} & \textbf{Meaning} \\
\midrule
è & is & e & and \\
là & there & la & the; her \\
sì & yes & si & himself, herself (reflexive) \\
dà & gives & da & from \\
tè & tea & te & you (after a preposition) \\
\bottomrule
\end{xltabular}}

\begin{notebox}{rulecol}\notelabel{rulecol}{Rule}{Never write e' for è}
The apostrophe is not an accent. Writing \textit{e'} for \textit{è} is a typing habit, not standard spelling. In this course an answer without the correct accent is marked wrong.
\end{notebox}

\subsection{Double consonants}\label{les:stress-doubles-accents:doubles}

\begin{notebox}{rulecol}\notelabel{rulecol}{Rule}{Double consonants are held longer}
A written double consonant is pronounced longer (or with more tension) than a single one. The difference can change meaning entirely.
\end{notebox}

{\small\setlength{\LTpre}{2pt}\setlength{\LTpost}{6pt}
\tablecaption{Minimal pairs}
\begin{longtable}{llll}
\toprule
\textbf{Single} & \textbf{Meaning} & \textbf{Double} & \textbf{Meaning} \\
\midrule
nono & ninth & nonno & grandfather \\
sete & thirst & sette & seven \\
caro & dear & carro & cart \\
pala & shovel & palla & ball \\
capello & (a) hair & cappello & hat \\
\bottomrule
\end{longtable}}

\subsection{The apostrophe: elision}\label{les:stress-doubles-accents:elision}

When certain short words end in a vowel and the next word begins with a vowel, the first vowel is dropped and replaced by an apostrophe, with no space after it.\par

\begin{itemize}
\item\relax \textit{lo / la} → \textit{l'}: \textit{l'albero, l'acqua}
\item\relax \textit{una} → \textit{un'}: \textit{un'amica, un'ape}
\item\relax \textit{ci} + \textit{è} → \textit{c'è} (there is); \textit{c'era} (there was)
\item\relax \textit{di} → \textit{d'}: \textit{d'estate, pieno d'acqua}
\item\relax \textit{dove, come, cosa} + \textit{è} → \textit{dov'è, com'è, cos'è}
\end{itemize}

\begin{notebox}{exccol}\notelabel{exccol}{Exception}{un amico — no apostrophe}
The masculine \textit{un} is not an elided form (it is also used before consonants: \textit{un gatto}), so it never takes an apostrophe: \textit{un amico}, but \textit{un'amica}. Likewise \textit{qual è} is written without an apostrophe.
\end{notebox}

\subsection*{Worked examples}
\begin{examples}
\item\relax \textit{L'albero è alto.}\newline The tree is tall.\newline {\color{muted}\small l' (from lo) before a vowel; AL-be-ro is stressed on the third-to-last syllable.}
\item\relax \textit{Il nonno beve il caffè.}\newline Grandfather drinks coffee.\newline {\color{muted}\small nonno: double n. caffè: written final stress.}
\item\relax \textit{Sara è la mia amica.}\newline Sara is my friend.\newline {\color{muted}\small è (is) carries an accent; e (and) does not.}
\item\relax \textit{C'è un'ape sul fiore.}\newline There is a bee on the flower.\newline {\color{muted}\small c'è = ci + è; un'ape = una + ape.}
\item\relax \textit{Ho sette anni.}\newline I am seven years old.\newline {\color{muted}\small sette (seven) versus sete (thirst).}
\item\relax \textit{Perché la città è bella.}\newline Because the city is beautiful.\newline {\color{muted}\small perché (acute accent), città (grave accent).}
\end{examples}

\subsection*{Exercises}\label{les:stress-doubles-accents:exercises}
{\color{muted}\small Answers: page~\pageref{ans:stress-doubles-accents}.}

\begin{exercises}
\item\label{ex:stress-doubles-accents-1}\relax Write the missing word meaning 'is': Il gatto \blank{} nero.
\item\label{ex:stress-doubles-accents-2}\relax What does \textit{nonno} mean?\par\nopagebreak\begin{options}
\item\relax not
\item\relax grandfather
\item\relax nobody
\item\relax ninth
\end{options}
\item\label{ex:stress-doubles-accents-3}\relax Choose the correct spelling: 'a (female) friend'\par\nopagebreak\begin{options}
\item\relax un' amica
\item\relax un amica
\item\relax un'amica
\item\relax una amica
\end{options}
\item\label{ex:stress-doubles-accents-4}\relax Choose the correct spelling: 'a (male) friend'\par\nopagebreak\begin{options}
\item\relax un amico
\item\relax uno amico
\item\relax un'amico
\item\relax una amico
\end{options}
\item\label{ex:stress-doubles-accents-5}\relax Which word is stressed on the third-to-last syllable?\par\nopagebreak\begin{options}
\item\relax finestra
\item\relax albero
\item\relax bambino
\item\relax città
\end{options}
\item\label{ex:stress-doubles-accents-6}\relax Add the elided forms.\par\smallskip\begingroup\setlength{\parindent}{0pt}\clozeblank{1}acqua (the water) · \clozeblank{2} un orso (there is a bear) · \clozeblank{3} la mamma? (where is Mum?)\par\endgroup
\end{exercises}

\begin{summarybox}
\begin{itemize}
\item\relax Stress usually falls on the second-to-last syllable, but many words are stressed earlier — learn it with the word.
\item\relax Final stress is always written: città, caffè, perché.
\item\relax Accents distinguish è/e, là/la, sì/si, dà/da, tè/te.
\item\relax Double consonants are pronounced longer and can change meaning.
\item\relax Elision: l'albero, un'amica, c'è — but un amico.
\end{itemize}
\end{summarybox}

\relatedwords{\textit{l'albero} tree \textperiodcentered{} \textit{il nonno} grandfather \textperiodcentered{} \textit{il caffè} coffee \textperiodcentered{} \textit{la città} city; town \textperiodcentered{} \textit{perché} because; why \textperiodcentered{} \textit{l'amica} friend (female) \textperiodcentered{} \textit{l'amico} friend (male) \textperiodcentered{} \textit{sette} seven \textperiodcentered{} \textit{l'ape} bee}
\applyit{Read \textit{Il gatto e la luna} (A0–A1), \S\ref{story:il-gatto-e-la-luna}. Study the rhyme \textit{La pioggia sul tetto}, \S\ref{rhyme:la-pioggia-sul-tetto}.}

\chapter{Nouns and articles}\label{mod:m02}
\moduleintro{A0}{Gender, plurals, definite and indefinite articles, irregular nouns.}{Noun gender and number (\S\ref{les:noun-gender-number}); The definite article (\S\ref{les:definite-articles}); The indefinite article; c'è and ci sono (\S\ref{les:indefinite-articles}); Spelling changes and irregular plurals (\S\ref{les:irregular-nouns})}

\section{Noun gender and number}\label{les:noun-gender-number}
\lessonmeta{A0}{Masculine and feminine nouns, the three main endings, and regular plurals.}
\begin{objectives}
\item\relax Identify the gender of most nouns from their ending.
\item\relax Form regular plurals of nouns in -o, -a and -e.
\item\relax Recognise nouns that do not change in the plural.
\end{objectives}

\subsection{Every noun has a gender}\label{les:noun-gender-number:gender}

Every Italian noun is either masculine or feminine — even nouns for objects and ideas. The gender decides which article and which adjective forms go with the noun, so always learn a noun together with its article: \textit{il gatto}, \textit{la casa}.\par

{\small\setlength{\LTpre}{2pt}\setlength{\LTpost}{6pt}
\tablecaption{Gender by ending}
\begin{xltabular}{\linewidth}{llLl}
\toprule
\textbf{Ending} & \textbf{Usually} & \textbf{Examples} & \textbf{Status} \\
\midrule
-o & masculine & il gatto, il libro, il bosco & strong tendency \\
-a & feminine & la casa, la mela, la luna & strong tendency \\
-e & either & il fiore, il cane (m); la notte, la chiave (f) & must be learned \\
-tà, -tù & feminine & la città, la felicità, la virtù & rule \\
consonant & masculine & il film, lo sport, il bar & strong tendency \\
\bottomrule
\end{xltabular}}

\begin{notebox}{exccol}\notelabel{exccol}{Exception}{Common exceptions}
\textit{la mano} (hand) is feminine despite -o. \textit{il problema}, \textit{il papà} and \textit{il poeta} are masculine despite -a. Shortened words keep the gender of the full word: \textit{la foto} (fotografia), \textit{la radio}, \textit{la moto}.
\end{notebox}

\begin{notebox}{rulecol}\notelabel{rulecol}{Rule}{Natural gender}
For people and many animals, gender follows sex where a pair exists: \textit{il ragazzo / la ragazza}, \textit{il nonno / la nonna}, \textit{il gatto / la gatta}.
\end{notebox}

\subsection{Regular plurals}\label{les:noun-gender-number:plural}

{\small\setlength{\LTpre}{2pt}\setlength{\LTpost}{6pt}
\tablecaption{Plural endings}
\begin{xltabular}{\linewidth}{llL}
\toprule
\textbf{Singular} & \textbf{Plural} & \textbf{Examples} \\
\midrule
-o (m) & -i & il libro → i libri \\
-a (f) & -e & la casa → le case \\
-e (m or f) & -i & il fiore → i fiori; la notte → le notti \\
-a (m) & -i & il problema → i problemi \\
\bottomrule
\end{xltabular}}

\begin{notebox}{rulecol}\notelabel{rulecol}{Rule}{Invariable nouns}
These do not change in the plural: nouns with a stressed final vowel (\textit{la città → le città, il caffè → i caffè}); nouns ending in a consonant (\textit{il film → i film}); one-syllable nouns (\textit{il re → i re}); shortened words (\textit{la foto → le foto}). Only the article shows the number.
\end{notebox}

\subsection*{Worked examples}
\begin{examples}
\item\relax \textit{il libro, i libri}\newline the book, the books
\item\relax \textit{la stella, le stelle}\newline the star, the stars
\item\relax \textit{il fiore, i fiori}\newline the flower, the flowers\newline {\color{muted}\small -e → -i, masculine.}
\item\relax \textit{la notte, le notti}\newline the night, the nights\newline {\color{muted}\small -e → -i, feminine.}
\item\relax \textit{la città, le città}\newline the city, the cities\newline {\color{muted}\small Invariable: stressed final vowel.}
\item\relax \textit{La mano è piccola.}\newline The hand is small.\newline {\color{muted}\small mano is feminine: la mano, piccola.}
\end{examples}

\subsection*{Exercises}\label{les:noun-gender-number:exercises}
{\color{muted}\small Answers: page~\pageref{ans:noun-gender-number}.}

\begin{exercises}
\item\label{ex:noun-gender-number-1}\relax Write the plural: \textit{il libro} → i \blank{}
\item\label{ex:noun-gender-number-2}\relax Write the plural: \textit{la stella} → le \blank{}
\item\label{ex:noun-gender-number-3}\relax Write the plural: \textit{la notte} → le \blank{}
\item\label{ex:noun-gender-number-4}\relax Write the plural: \textit{la città} → le \blank{}
\item\label{ex:noun-gender-number-5}\relax What is the gender of \textit{mano} (hand)?\par\nopagebreak\begin{options}
\item\relax masculine
\item\relax feminine
\end{options}
\item\label{ex:noun-gender-number-6}\relax \textit{fiore} ends in -e. What does that tell you about its gender?\par\nopagebreak\begin{options}
\item\relax Nothing certain: -e nouns can be masculine or feminine
\item\relax It has no gender
\item\relax It is feminine
\item\relax It is masculine
\end{options}
\item\label{ex:noun-gender-number-7}\relax Match each singular to its plural.\par\smallskip\begin{tabularx}{\linewidth}{@{}L L@{}}
1.\ il gatto & (a)\ i re \\
2.\ la mela & (b)\ i gatti \\
3.\ il cane & (c)\ i cani \\
4.\ il re & (d)\ le mele \\
\end{tabularx}
\end{exercises}

\begin{summarybox}
\begin{itemize}
\item\relax -o: usually masculine; -a: usually feminine; -e: either — learn the article.
\item\relax Plurals: -o → -i, -a (f) → -e, -e → -i, -a (m) → -i.
\item\relax Invariable: final stressed vowel, final consonant, one syllable, shortened words.
\item\relax Exceptions to remember: la mano, il problema, il papà.
\end{itemize}
\end{summarybox}

\relatedwords{\textit{il libro} book \textperiodcentered{} \textit{la stella} star \textperiodcentered{} \textit{il fiore} flower \textperiodcentered{} \textit{la notte} night \textperiodcentered{} \textit{la città} city; town \textperiodcentered{} \textit{la mano} hand \textperiodcentered{} \textit{il problema} problem \textperiodcentered{} \textit{il re} king \textperiodcentered{} \textit{il cane} dog \textperiodcentered{} \textit{la mela} apple}
\applyit{Read \textit{La famiglia di Marco} (A0–A1), \S\ref{story:la-famiglia-di-marco}. Study the rhyme \textit{Conta delle stelle}, \S\ref{rhyme:conta-delle-stelle}.}

\section{The definite article}\label{les:definite-articles}
\lessonmeta{A0}{il, lo, l', la, i, gli, le — choosing the article from the noun's gender, number and first letters.}
\begin{objectives}
\item\relax Choose the correct definite article for any noun whose gender you know.
\item\relax Apply the lo/gli rule for s + consonant, z, gn, ps, x, y.
\item\relax Use l' before vowels and know its plural (gli or le).
\end{objectives}

\subsection{Seven forms}\label{les:definite-articles:forms}

The definite article (the) agrees in gender and number with the noun, and its form also depends on the sound at the start of the \textbf{next word}.\par

{\small\setlength{\LTpre}{2pt}\setlength{\LTpost}{6pt}
\tablecaption{Definite articles}
\begin{xltabular}{\linewidth}{lllL}
\toprule
\textbf{} & \textbf{Singular} & \textbf{Plural} & \textbf{Used before} \\
\midrule
masculine & il & i & most consonants: il gatto → i gatti \\
masculine & lo & gli & s + consonant, z, gn, ps, pn, x, y: lo zaino → gli zaini; lo specchio → gli specchi \\
masculine & l' & gli & a vowel: l'albero → gli alberi \\
feminine & la & le & a consonant: la casa → le case \\
feminine & l' & le & a vowel: l'ape → le api \\
\bottomrule
\end{xltabular}}

\begin{notebox}{rulecol}\notelabel{rulecol}{Rule}{The article depends on the next word}
The choice is made by the word immediately after the article, which may be an adjective: \textit{lo zio} but \textit{il mio zio}; \textit{l'amico} but \textit{il vecchio amico}.
\end{notebox}

\begin{notebox}{rulecol}\notelabel{rulecol}{Rule}{Feminine plural never elides}
In the plural, use \textit{le} before every feminine noun, including those beginning with a vowel: \textit{le api, le arance} (not \textit{l'api}).
\end{notebox}

\subsection{Using the definite article}\label{les:definite-articles:use}

Italian uses the definite article more than English does:\par

\begin{itemize}
\item\relax With general statements: \textit{Il miele è dolce} = Honey is sweet.
\item\relax With categories you like or dislike: \textit{Amo i gatti} = I love cats.
\item\relax With possessives: \textit{il mio libro} (see the possessives lesson).
\item\relax Not with names of people: \textit{Marco}, not \textit{il Marco} (except in some regional speech).
\end{itemize}

\subsection*{Worked examples}
\begin{examples}
\item\relax \textit{Il gatto dorme.}\newline The cat is sleeping.
\item\relax \textit{Lo zaino è pesante.}\newline The backpack is heavy.\newline {\color{muted}\small lo before z.}
\item\relax \textit{L'albero è alto.}\newline The tree is tall.\newline {\color{muted}\small l' before a vowel.}
\item\relax \textit{Gli scoiattoli salgono sugli alberi.}\newline The squirrels climb the trees.\newline {\color{muted}\small gli: plural of lo and of masculine l'.}
\item\relax \textit{Le api fanno il miele.}\newline Bees make honey.\newline {\color{muted}\small General statement: Italian uses the article.}
\item\relax \textit{Il miele è dolce.}\newline Honey is sweet.
\end{examples}

\subsection*{Exercises}\label{les:definite-articles:exercises}
{\color{muted}\small Answers: page~\pageref{ans:definite-articles}.}

\begin{exercises}
\item\label{ex:definite-articles-1}\relax \blank{} zaino (backpack, masculine)\par\nopagebreak\begin{options}
\item\relax l'
\item\relax il
\item\relax lo
\item\relax la
\end{options}
\item\label{ex:definite-articles-2}\relax \blank{} albero (tree, masculine)\par\nopagebreak\begin{options}
\item\relax la
\item\relax l'
\item\relax lo
\item\relax il
\end{options}
\item\label{ex:definite-articles-3}\relax \blank{} api (bees, feminine plural)\par\nopagebreak\begin{options}
\item\relax le
\item\relax i
\item\relax gli
\item\relax l'
\end{options}
\item\label{ex:definite-articles-4}\relax Give the plural with its article.\par\smallskip{\begin{tabular}{@{}l l@{}}
il cane & \blankcell \\
lo specchio & \blankcell \\
l'orso & \blankcell \\
la foglia & \blankcell \\
l'isola & \blankcell \\
\end{tabular}}
\item\label{ex:definite-articles-5}\relax \textit{zio} takes \textit{lo} (lo zio). Which article goes in: \blank{} vecchio zio?\par\nopagebreak\begin{options}
\item\relax lo
\item\relax il
\item\relax la
\item\relax l'
\end{options}
\item\label{ex:definite-articles-6}\relax Translate: The bees make honey.\par\answerline
\end{exercises}

\begin{summarybox}
\begin{itemize}
\item\relax Masculine: il/i (most consonants), lo/gli (s + consonant, z, gn, ps…), l'/gli (vowel).
\item\relax Feminine: la/le (consonant), l'/le (vowel).
\item\relax The article is chosen by the word right after it.
\item\relax General statements and categories take the article: Il miele è dolce.
\end{itemize}
\end{summarybox}

\relatedwords{\textit{lo zaino} backpack; rucksack \textperiodcentered{} \textit{l'albero} tree \textperiodcentered{} \textit{l'ape} bee \textperiodcentered{} \textit{lo specchio} mirror \textperiodcentered{} \textit{lo scoiattolo} squirrel \textperiodcentered{} \textit{lo zio} uncle \textperiodcentered{} \textit{l'isola} island \textperiodcentered{} \textit{il miele} honey \textperiodcentered{} \textit{l'orso} bear}
\applyit{Read \textit{Il gatto e la luna} (A0–A1), \S\ref{story:il-gatto-e-la-luna}. Read \textit{Nel bosco} (A1), \S\ref{story:nel-bosco}.}

\section{The indefinite article; c'è and ci sono}\label{les:indefinite-articles}
\lessonmeta{A0}{un, uno, una, un'; saying 'there is' and 'there are'.}
\begin{objectives}
\item\relax Choose un, uno, una or un' correctly.
\item\relax Say what there is somewhere with c'è and ci sono.
\item\relax Recognise dei, degli, delle as 'some'.
\end{objectives}

\subsection{un, uno, una, un'}\label{les:indefinite-articles:forms}

{\small\setlength{\LTpre}{2pt}\setlength{\LTpost}{6pt}
\tablecaption{Indefinite articles (a / an)}
\begin{xltabular}{\linewidth}{llLL}
\toprule
\textbf{} & \textbf{Form} & \textbf{Used before} & \textbf{Examples} \\
\midrule
masculine & un & most consonants and all vowels & un gatto, un albero \\
masculine & uno & s + consonant, z, gn, ps, x, y & uno zaino, uno scoiattolo \\
feminine & una & consonants & una mela, una stella \\
feminine & un' & vowels & un'ape, un'isola \\
\bottomrule
\end{xltabular}}

\begin{notebox}{rulecol}\notelabel{rulecol}{Rule}{The parallel with the definite article}
Wherever the definite article is \textit{lo}, the indefinite is \textit{uno} (lo zaino / uno zaino). Before a masculine vowel, \textit{l'} corresponds to \textit{un} without an apostrophe (l'albero / un albero).
\end{notebox}

\subsection{c'è and ci sono}\label{les:indefinite-articles:ce}

\textit{c'è} (= \textit{ci è}) means 'there is'; \textit{ci sono} means 'there are'. The verb agrees with what follows.\par

\begin{itemize}
\item\relax \textit{C'è un gatto in giardino.} = There is a cat in the garden.
\item\relax \textit{Ci sono tre gatti in giardino.} = There are three cats in the garden.
\item\relax Past (imperfetto): \textit{c'era / c'erano} = there was / there were — as in \textit{C'era una volta…} (Once upon a time there was…).
\end{itemize}

\subsection{'Some': dei, degli, delle}\label{les:indefinite-articles:partitive}

The indefinite article has no plural, but \textit{di} + the definite article gives the meaning 'some': \textit{dei gatti} (some cats), \textit{degli alberi} (some trees), \textit{delle mele} (some apples). These are the partitive article: \textit{di} contracted with the definite article (\textit{di + i = dei, di + gli = degli, di + le = delle}; in the singular \textit{del pane, della farina} = some bread, some flour, with uncountable nouns). The contractions themselves are covered in the lesson on \lessonref{articulated-prepositions}{articulated prepositions}. The partitive is optional: \textit{Ho comprato mele} and \textit{Ho comprato delle mele} are both correct, and in negative sentences it is usually omitted: \textit{Non ho comprato mele}. \textit{alcuni/alcune} + plural and \textit{un po' di} + singular are common alternatives.\par

\subsection*{Worked examples}
\begin{examples}
\item\relax \textit{Ho un cane e una gatta.}\newline I have a (male) dog and a (female) cat.
\item\relax \textit{C'è uno scoiattolo sull'albero.}\newline There is a squirrel in the tree.
\item\relax \textit{C'è un'ape sul fiore.}\newline There is a bee on the flower.
\item\relax \textit{Ci sono molte stelle in cielo.}\newline There are many stars in the sky.
\item\relax \textit{Un albero è caduto.}\newline A tree has fallen.\newline {\color{muted}\small un before a masculine vowel, no apostrophe.}
\item\relax \textit{Mangio delle fragole.}\newline I eat some strawberries.
\end{examples}

\subsection*{Exercises}\label{les:indefinite-articles:exercises}
{\color{muted}\small Answers: page~\pageref{ans:indefinite-articles}.}

\begin{exercises}
\item\label{ex:indefinite-articles-1}\relax \blank{} zaino\par\nopagebreak\begin{options}
\item\relax una
\item\relax uno
\item\relax un
\item\relax un'
\end{options}
\item\label{ex:indefinite-articles-2}\relax \blank{} isola (feminine)\par\nopagebreak\begin{options}
\item\relax un'
\item\relax una
\item\relax uno
\item\relax un
\end{options}
\item\label{ex:indefinite-articles-3}\relax \blank{} orso (masculine)\par\nopagebreak\begin{options}
\item\relax un
\item\relax una
\item\relax uno
\item\relax un'
\end{options}
\item\label{ex:indefinite-articles-4}\relax Complete with c'è or ci sono.\par\smallskip\begingroup\setlength{\parindent}{0pt}In giardino \clozeblank{1} un albero. Sull'albero \clozeblank{2} due uccelli.\par\endgroup
\item\label{ex:indefinite-articles-5}\relax Translate: \textit{C'è un'ape sul fiore.}\par\answerline
\item\label{ex:indefinite-articles-6}\relax Translate: There are three cats.\par\answerline
\end{exercises}

\begin{summarybox}
\begin{itemize}
\item\relax un (most masculine nouns, including vowels), uno (s + consonant, z…), una (feminine consonant), un' (feminine vowel).
\item\relax c'è + singular = there is; ci sono + plural = there are; c'era = there was.
\item\relax dei / degli / delle = some.
\end{itemize}
\end{summarybox}

\relatedwords{\textit{lo scoiattolo} squirrel \textperiodcentered{} \textit{l'ape} bee \textperiodcentered{} \textit{l'isola} island \textperiodcentered{} \textit{la stella} star \textperiodcentered{} \textit{la fragola} strawberry \textperiodcentered{} \textit{essere} to be \textperiodcentered{} \textit{c'era una volta} once upon a time (there was)}
\applyit{Read \textit{Il cane al mare} (A0–A1), \S\ref{story:il-cane-al-mare}. Read \textit{Nel bosco} (A1), \S\ref{story:nel-bosco}. Study the rhyme \textit{Conta delle stelle}, \S\ref{rhyme:conta-delle-stelle}.}

\section{Spelling changes and irregular plurals}\label{les:irregular-nouns}
\lessonmeta{A1}{Plurals of -ca/-ga, -co/-go, -cia/-gia and -io, and the common irregular nouns.}
\begin{objectives}
\item\relax Form plurals with the spelling adjustments for hard c and g.
\item\relax Know which plural patterns are rules and which are only tendencies.
\item\relax Learn the high-frequency irregular nouns: uomo, uovo, braccio, mano.
\end{objectives}

\subsection{-ca and -ga: always -che, -ghe}\label{les:irregular-nouns:ca-ga}

\begin{notebox}{rulecol}\notelabel{rulecol}{Rule}{Feminine -ca/-ga → -che/-ghe}
The h keeps the hard sound in the plural: \textit{la mucca → le mucche}, \textit{la barca → le barche}, \textit{la tartaruga → le tartarughe}.
\end{notebox}

\subsection{-co and -go: a tendency, not a rule}\label{les:irregular-nouns:co-go}

\begin{notebox}{tendcol}\notelabel{tendcol}{Tendency}{It depends on the stress}
If the stress falls on the second-to-last syllable, -co/-go usually become -chi/-ghi: \textit{il BO-sco → i boschi}, \textit{il LA-go → i laghi}, \textit{il GIO-co → i giochi}. If it falls earlier, they usually become -ci/-gi: \textit{il ME-di-co → i medici}, \textit{sim-PA-ti-co → simpatici}.
\end{notebox}

\begin{notebox}{exccol}\notelabel{exccol}{Exception}{Exceptions}
\textit{l'amico → gli amici}, \textit{il nemico → i nemici}, \textit{il greco → i greci} are stressed on the second-to-last syllable but still take -ci. Learn these individually.
\end{notebox}

\subsection{-cia, -gia and -io}\label{les:irregular-nouns:cia-gia-io}

{\small\setlength{\LTpre}{2pt}\setlength{\LTpost}{6pt}
\tablecaption{Plurals with i}
\begin{xltabular}{\linewidth}{LlLl}
\toprule
\textbf{Pattern} & \textbf{Plural} & \textbf{Examples} & \textbf{Status} \\
\midrule
-cia/-gia after a consonant & -ce/-ge & l'arancia → le arance; la spiaggia → le spiagge & tendency \\
-cia/-gia after a vowel & -cie/-gie & la camicia → le camicie; la valigia → le valigie & tendency \\
-io (i unstressed) & -i & il figlio → i figli; l'occhio → gli occhi & rule \\
-ìo (i stressed) & -ii & lo zio → gli zii & rule \\
\bottomrule
\end{xltabular}}

\subsection{Irregular nouns}\label{les:irregular-nouns:irregular}

{\small\setlength{\LTpre}{2pt}\setlength{\LTpost}{6pt}
\tablecaption{Common irregular plurals}
\begin{xltabular}{\linewidth}{llL}
\toprule
\textbf{Singular} & \textbf{Plural} & \textbf{Note} \\
\midrule
l'uomo & gli uomini & completely irregular \\
la mano & le mani & feminine noun in -o \\
l'uovo & le uova & masculine singular, feminine plural in -a \\
il braccio & le braccia & body parts: also il dito → le dita, il labbro → le labbra \\
la gente & — & collective noun, always singular: la gente è… \\
\bottomrule
\end{xltabular}}

\subsection*{Worked examples}
\begin{examples}
\item\relax \textit{Le mucche sono nel prato.}\newline The cows are in the meadow.
\item\relax \textit{Nei boschi ci sono molti funghi.}\newline There are many mushrooms in the woods.\newline {\color{muted}\small bosco → boschi; fungo → funghi.}
\item\relax \textit{Gli amici giocano insieme.}\newline The friends play together.\newline {\color{muted}\small amico → amici (exception).}
\item\relax \textit{Le arance sono dolci.}\newline The oranges are sweet.\newline {\color{muted}\small arancia → arance.}
\item\relax \textit{La gallina fa le uova.}\newline The hen lays eggs.\newline {\color{muted}\small uovo → le uova.}
\item\relax \textit{Il gatto ha gli occhi verdi.}\newline The cat has green eyes.\newline {\color{muted}\small occhio → occhi.}
\end{examples}

\subsection*{Exercises}\label{les:irregular-nouns:exercises}
{\color{muted}\small Answers: page~\pageref{ans:irregular-nouns}.}

\begin{exercises}
\item\label{ex:irregular-nouns-1}\relax Plural of \textit{la barca}: le \blank{}
\item\label{ex:irregular-nouns-2}\relax Plural of \textit{l'amico}: gli \blank{}
\item\label{ex:irregular-nouns-3}\relax Plural of \textit{il gioco}: i \blank{}
\item\label{ex:irregular-nouns-4}\relax Write the plural with its article.\par\smallskip{\begin{tabular}{@{}l l@{}}
l'uomo & \blankcell \\
l'uovo & \blankcell \\
la mano & \blankcell \\
lo zio & \blankcell \\
la spiaggia & \blankcell \\
\end{tabular}}
\item\label{ex:irregular-nouns-5}\relax Which statement is accurate?\par\nopagebreak\begin{options}
\item\relax -ca and -co both always keep the hard sound with h.
\item\relax -ca always becomes -che, but -co becomes -chi or -ci depending largely on stress.
\item\relax -ca becomes -ce.
\item\relax -co always becomes -chi.
\end{options}
\item\label{ex:irregular-nouns-6}\relax Translate: The cat has green eyes.\par\answerline
\end{exercises}

\begin{summarybox}
\begin{itemize}
\item\relax -ca/-ga → -che/-ghe (rule).
\item\relax -co/-go → -chi/-ghi or -ci/-gi (tendency by stress; amico → amici).
\item\relax -cia/-gia: -ce/-ge after a consonant, -cie/-gie after a vowel (tendency).
\item\relax -io → -i; stressed -ìo → -ii.
\item\relax Irregular: uomini, uova, braccia, mani.
\end{itemize}
\end{summarybox}

\relatedwords{\textit{la barca} boat \textperiodcentered{} \textit{la mucca} cow \textperiodcentered{} \textit{il bosco} wood; forest \textperiodcentered{} \textit{il lago} lake \textperiodcentered{} \textit{il gioco} game; toy \textperiodcentered{} \textit{l'amico} friend (male) \textperiodcentered{} \textit{l'arancia} orange (fruit) \textperiodcentered{} \textit{la spiaggia} beach \textperiodcentered{} \textit{l'uomo} man \textperiodcentered{} \textit{l'uovo} egg \textperiodcentered{} \textit{il braccio} arm \textperiodcentered{} \textit{la mano} hand \textperiodcentered{} \textit{l'occhio} eye \textperiodcentered{} \textit{lo zio} uncle}
\applyit{Read \textit{Nel bosco} (A1), \S\ref{story:nel-bosco}. Read \textit{Il cane al mare} (A0–A1), \S\ref{story:il-cane-al-mare}.}

\chapter{Adjectives}\label{mod:m03}
\moduleintro{A0–A1}{Agreement and position, demonstratives, possessives, bello/buono, comparison.}{Adjectives: agreement and position (\S\ref{les:adjective-agreement}); Demonstratives and possessives (\S\ref{les:demonstratives-possessives}); bello and buono; comparing things (\S\ref{les:bello-buono-comparison})}

\section{Adjectives: agreement and position}\label{les:adjective-agreement}
\lessonmeta{A0}{Four-form and two-form adjectives, invariable colours, and where adjectives go.}
\begin{objectives}
\item\relax Make adjectives agree in gender and number with their noun.
\item\relax Distinguish -o adjectives (four forms) from -e adjectives (two forms).
\item\relax Place adjectives correctly, usually after the noun.
\end{objectives}

\subsection{Two classes of adjective}\label{les:adjective-agreement:classes}

\begin{notebox}{rulecol}\notelabel{rulecol}{Rule}{Agreement}
An adjective takes the gender and number of the noun it describes: \textit{il gatto nero, la gatta nera, i gatti neri, le gatte nere}.
\end{notebox}

{\small\setlength{\LTpre}{2pt}\setlength{\LTpost}{6pt}
\tablecaption{Adjective endings}
\begin{longtable}{lllll}
\toprule
\textbf{Class} & \textbf{masc. sg.} & \textbf{fem. sg.} & \textbf{masc. pl.} & \textbf{fem. pl.} \\
\midrule
-o (four forms) & rosso & rossa & rossi & rosse \\
-e (two forms) & verde & verde & verdi & verdi \\
invariable & blu & blu & blu & blu \\
\bottomrule
\end{longtable}}

Invariable colour adjectives include \textit{blu, rosa, viola}: \textit{le farfalle rosa} (the pink butterflies).\par

\begin{notebox}{rulecol}\notelabel{rulecol}{Rule}{Mixed groups take the masculine plural}
One adjective describing a masculine and a feminine noun together is masculine plural: \textit{Il gatto e la volpe sono stanchi.}
\end{notebox}

\begin{notebox}{rulecol}\notelabel{rulecol}{Rule}{-co/-go adjectives follow the noun pattern}
\textit{bianco → bianchi, bianche}; \textit{lungo → lunghi, lunghe}. But stress-early adjectives take -ci: \textit{simpatico → simpatici} (feminine always -che: \textit{simpatiche}).
\end{notebox}

\subsection{Position}\label{les:adjective-agreement:position}

\begin{notebox}{tendcol}\notelabel{tendcol}{Tendency}{Most adjectives follow the noun}
Colours, shapes, nationalities and most descriptive adjectives come after the noun: \textit{un gatto nero, una canzone italiana, un fiume lungo}.
\end{notebox}

\begin{notebox}{tendcol}\notelabel{tendcol}{Tendency}{Some common adjectives often precede}
\textit{bello, buono, brutto, grande, piccolo, nuovo, vecchio, giovane, lungo} often come before the noun, especially in a subjective or emotional sense: \textit{una bella giornata, un piccolo orso}.
\end{notebox}

{\small\setlength{\LTpre}{2pt}\setlength{\LTpost}{6pt}
\tablecaption{Position can change the meaning}
\begin{xltabular}{\linewidth}{lLll}
\toprule
\textbf{Before the noun} & \textbf{Meaning} & \textbf{After the noun} & \textbf{Meaning} \\
\midrule
un grande uomo & a great man & un uomo grande & a big (tall, large) man \\
un vecchio amico & a friend of many years & un amico vecchio & an elderly friend \\
un nuovo libro & another / a different book & un libro nuovo & a brand-new book \\
\bottomrule
\end{xltabular}}

\subsection*{Worked examples}
\begin{examples}
\item\relax \textit{La casa è piccola.}\newline The house is small.
\item\relax \textit{I fiori sono gialli.}\newline The flowers are yellow.
\item\relax \textit{Le notti sono fredde.}\newline The nights are cold.
\item\relax \textit{Le ragazze sono felici.}\newline The girls are happy.\newline {\color{muted}\small felice: two-form adjective.}
\item\relax \textit{Ho uno zaino blu.}\newline I have a blue backpack.\newline {\color{muted}\small blu is invariable.}
\item\relax \textit{C'era un piccolo orso.}\newline There was a little bear.\newline {\color{muted}\small piccolo before the noun: affectionate, storybook tone.}
\end{examples}

\subsection*{Exercises}\label{les:adjective-agreement:exercises}
{\color{muted}\small Answers: page~\pageref{ans:adjective-agreement}.}

\begin{exercises}
\item\label{ex:adjective-agreement-1}\relax la mela (rosso) → la mela \blank{}
\item\label{ex:adjective-agreement-2}\relax le notti (freddo) → le notti \blank{}
\item\label{ex:adjective-agreement-3}\relax le ragazze (felice) → le ragazze \blank{}
\item\label{ex:adjective-agreement-4}\relax le farfalle (rosa)\par\nopagebreak\begin{options}
\item\relax le farfalle rose
\item\relax le farfalle rosae
\item\relax le farfalle rosi
\item\relax le farfalle rosa
\end{options}
\item\label{ex:adjective-agreement-5}\relax Write all four forms of \textit{bianco}.\par\smallskip{\begin{tabular}{@{}l l@{}}
masculine singular & \blankcell \\
feminine singular & \blankcell \\
masculine plural & \blankcell \\
feminine plural & \blankcell \\
\end{tabular}}
\item\label{ex:adjective-agreement-6}\relax What does \textit{un grande uomo} usually mean?\par\nopagebreak\begin{options}
\item\relax a tall boy
\item\relax a great man
\item\relax an old man
\item\relax a big man
\end{options}
\item\label{ex:adjective-agreement-7}\relax Put the words in order: 'The black cat sleeps.'\par\smallskip \token{dorme} \token{gatto} \token{Il} \token{nero}\par\answerline
\end{exercises}

\begin{summarybox}
\begin{itemize}
\item\relax Adjectives agree with the noun in gender and number.
\item\relax -o adjectives: -o, -a, -i, -e. -e adjectives: -e, -i. Some colours are invariable.
\item\relax Mixed masculine + feminine → masculine plural.
\item\relax Most adjectives follow the noun; a few common ones often precede it, sometimes with a different meaning.
\end{itemize}
\end{summarybox}

\relatedwords{\textit{rosso} red \textperiodcentered{} \textit{freddo} cold \textperiodcentered{} \textit{felice} happy \textperiodcentered{} \textit{rosa} pink \textperiodcentered{} \textit{bianco} white \textperiodcentered{} \textit{grande} big; large; great \textperiodcentered{} \textit{piccolo} small; little \textperiodcentered{} \textit{nero} black \textperiodcentered{} \textit{blu} blue (dark) \textperiodcentered{} \textit{simpatico} nice; likeable \textperiodcentered{} \textit{giallo} yellow}
\applyit{Read \textit{Il gatto e la luna} (A0–A1), \S\ref{story:il-gatto-e-la-luna}. Read \textit{La tartaruga lenta} (A1), \S\ref{story:la-tartaruga-lenta}. Study the rhyme \textit{Il valzer degli animali}, \S\ref{rhyme:il-valzer-degli-animali}.}

\section{Demonstratives and possessives}\label{les:demonstratives-possessives}
\lessonmeta{A1}{questo and quello; mio, tuo, suo, nostro, vostro, loro — with and without the article.}
\begin{objectives}
\item\relax Use questo and quello in all their forms.
\item\relax Use possessives with the definite article.
\item\relax Know when family nouns drop the article.
\end{objectives}

\subsection{questo (this) and quello (that)}\label{les:demonstratives-possessives:questo-quello}

\textit{questo} behaves like an ordinary -o adjective: \textit{questo libro, questa casa, questi libri, queste case}. Before a vowel in the singular it may elide: \textit{quest'albero}.\par

\begin{notebox}{rulecol}\notelabel{rulecol}{Rule}{quello before a noun copies the definite article}
Take the article the noun would have and attach it to \textit{que-}: il → \textit{quel}, lo → \textit{quello}, l' → \textit{quell'}, la → \textit{quella}, i → \textit{quei}, gli → \textit{quegli}, le → \textit{quelle}.
\end{notebox}

{\small\setlength{\LTpre}{2pt}\setlength{\LTpost}{6pt}
\tablecaption{quello + noun}
\begin{longtable}{lll}
\toprule
\textbf{Article} & \textbf{quello form} & \textbf{Example} \\
\midrule
il gatto & quel & quel gatto \\
lo zaino & quello & quello zaino \\
l'albero & quell' & quell'albero \\
la casa & quella & quella casa \\
i gatti & quei & quei gatti \\
gli alberi & quegli & quegli alberi \\
le case & quelle & quelle case \\
\bottomrule
\end{longtable}}

\subsection{Possessives}\label{les:demonstratives-possessives:possessives}

{\small\setlength{\LTpre}{2pt}\setlength{\LTpost}{6pt}
\tablecaption{Possessive adjectives}
\begin{longtable}{lllll}
\toprule
\textbf{Owner} & \textbf{masc. sg.} & \textbf{fem. sg.} & \textbf{masc. pl.} & \textbf{fem. pl.} \\
\midrule
io & il mio & la mia & i miei & le mie \\
tu & il tuo & la tua & i tuoi & le tue \\
lui / lei & il suo & la sua & i suoi & le sue \\
noi & il nostro & la nostra & i nostri & le nostre \\
voi & il vostro & la vostra & i vostri & le vostre \\
loro & il loro & la loro & i loro & le loro \\
\bottomrule
\end{longtable}}

\begin{notebox}{rulecol}\notelabel{rulecol}{Rule}{Agreement with the thing possessed}
The possessive agrees with the noun it accompanies, not with the owner: \textit{il suo libro} = his book or her book; \textit{la sua casa} = his house or her house.
\end{notebox}

\begin{notebox}{rulecol}\notelabel{rulecol}{Rule}{The article is normally used}
Possessives are normally preceded by the definite article: \textit{il mio cane, la nostra casa}.
\end{notebox}

\begin{notebox}{exccol}\notelabel{exccol}{Exception}{Singular family nouns}
With a \textbf{singular}, unmodified family noun the article is dropped: \textit{mio padre, mia sorella, tuo fratello}. It returns in the plural (\textit{i miei fratelli}), with \textit{loro} (\textit{il loro padre}), with an adjective (\textit{la mia sorella piccola}), and usually with affectionate forms (\textit{la mia mamma, il mio papà}).
\end{notebox}

\subsection*{Worked examples}
\begin{examples}
\item\relax \textit{Questo libro è nuovo.}\newline This book is new.
\item\relax \textit{Quell'albero è molto alto.}\newline That tree is very tall.
\item\relax \textit{Il mio cane è piccolo.}\newline My dog is small.
\item\relax \textit{Mia sorella ama i gatti.}\newline My sister loves cats.\newline {\color{muted}\small Singular family noun: no article.}
\item\relax \textit{I miei nonni abitano al mare.}\newline My grandparents live by the sea.\newline {\color{muted}\small Plural: the article returns.}
\item\relax \textit{Marco cerca la sua palla.}\newline Marco is looking for his ball.\newline {\color{muted}\small sua agrees with palla (feminine), not with Marco.}
\end{examples}

\subsection*{Exercises}\label{les:demonstratives-possessives:exercises}
{\color{muted}\small Answers: page~\pageref{ans:demonstratives-possessives}.}

\begin{exercises}
\item\label{ex:demonstratives-possessives-1}\relax \blank{} alberi (those trees)\par\nopagebreak\begin{options}
\item\relax quegli
\item\relax quelle
\item\relax quelli
\item\relax quei
\end{options}
\item\label{ex:demonstratives-possessives-2}\relax \blank{} gatto (that cat)\par\nopagebreak\begin{options}
\item\relax quei
\item\relax quell'
\item\relax quello
\item\relax quel
\end{options}
\item\label{ex:demonstratives-possessives-3}\relax Sara and her book: Sara legge \blank{} libro.\par\nopagebreak\begin{options}
\item\relax suo
\item\relax la sua
\item\relax il sua
\item\relax il suo
\end{options}
\item\label{ex:demonstratives-possessives-4}\relax 'my father'\par\nopagebreak\begin{options}
\item\relax mia padre
\item\relax il mio padre
\item\relax la mia padre
\item\relax mio padre
\end{options}
\item\label{ex:demonstratives-possessives-5}\relax 'their house'\par\nopagebreak\begin{options}
\item\relax loro casa
\item\relax la sua casa
\item\relax la lora casa
\item\relax la loro casa
\end{options}
\item\label{ex:demonstratives-possessives-6}\relax Translate: My grandparents live by the sea.\par\answerline
\end{exercises}

\begin{summarybox}
\begin{itemize}
\item\relax questo: -o, -a, -i, -e. quello before a noun copies the article: quel, quello, quell', quella, quei, quegli, quelle.
\item\relax Possessives agree with the thing possessed and normally take the article.
\item\relax No article with a singular unmodified family noun (mia madre); loro always takes the article.
\end{itemize}
\end{summarybox}

\relatedwords{\textit{il padre} father \textperiodcentered{} \textit{la madre} mother \textperiodcentered{} \textit{la sorella} sister \textperiodcentered{} \textit{il fratello} brother \textperiodcentered{} \textit{il nonno} grandfather \textperiodcentered{} \textit{il libro} book \textperiodcentered{} \textit{l'albero} tree}
\applyit{Read \textit{La famiglia di Marco} (A0–A1), \S\ref{story:la-famiglia-di-marco}.}

\section{bello and buono; comparing things}\label{les:bello-buono-comparison}
\lessonmeta{A1}{The special forms of bello and buono before nouns, and elementary comparatives.}
\begin{objectives}
\item\relax Use bel, bello, bell', bei, begli and buon before nouns.
\item\relax Make comparisons with più … di, meno … di and come.
\item\relax Know the irregular comparatives migliore and peggiore.
\end{objectives}

\subsection{bello and buono before a noun}\label{les:bello-buono-comparison:bello-buono}

\begin{notebox}{rulecol}\notelabel{rulecol}{Rule}{bello copies the definite article}
Like \textit{quello}, \textit{bello} before a noun takes the form of the article: \textit{bel} (il), \textit{bello} (lo), \textit{bell'} (l'), \textit{bella} (la), \textit{bei} (i), \textit{begli} (gli), \textit{belle} (le). After the noun or after \textit{essere} it is a regular -o adjective: \textit{Il fiore è bello. I fiori sono belli.}
\end{notebox}

\begin{notebox}{rulecol}\notelabel{rulecol}{Rule}{buono copies the indefinite article}
In the singular before a noun: \textit{buon} (un), \textit{buono} (uno), \textit{buona} (una), \textit{buon'} (un'): \textit{buon giorno, buono zucchero, buona notte, buon'idea}. Plurals are regular: \textit{buoni, buone}.
\end{notebox}

\subsection{Comparisons}\label{les:bello-buono-comparison:comparison}

{\small\setlength{\LTpre}{2pt}\setlength{\LTpost}{6pt}
\tablecaption{Comparing two things}
\begin{xltabular}{\linewidth}{lLL}
\toprule
\textbf{Meaning} & \textbf{Structure} & \textbf{Example} \\
\midrule
more … than & più + adjective + di & La lepre è più veloce della tartaruga. \\
less … than & meno + adjective + di & La tartaruga è meno veloce della lepre. \\
as … as & (così) + adjective + come & Luca è alto come suo padre. \\
\bottomrule
\end{xltabular}}

\begin{notebox}{rulecol}\notelabel{rulecol}{Rule}{di combines with the article}
\textit{di} + article contracts: \textit{della tartaruga, del cane, dei gatti}. Before a name or pronoun there is no article: \textit{più alto di Marco, più alta di me}.
\end{notebox}

\begin{notebox}{rulecol}\notelabel{rulecol}{Rule}{di or che?}
Use \textit{di} when two different people or things are compared for the same quality, and the second term is a noun, name or pronoun: \textit{Marco è più alto di Luca. Sei più brava di me.} Use \textit{che} in all the other cases: two qualities of the same thing (\textit{È più bello che utile}), two verbs (\textit{Leggere è più facile che scrivere}), two quantities (\textit{Ho più libri che quaderni}), two terms introduced by a preposition (\textit{Fa più caldo a Roma che a Milano}), and two adverbs (\textit{Meglio tardi che mai}).
\end{notebox}

\begin{notebox}{rulecol}\notelabel{rulecol}{Rule}{Comparing with a clause}
When the second term is a whole clause, use \textit{di quanto} or \textit{di quello che}: \textit{È più difficile di quanto pensassi / di quello che pensavo.} After \textit{di quanto} the congiuntivo is common in careful style; the indicativo is also used.
\end{notebox}

\begin{notebox}{exccol}\notelabel{exccol}{Exception}{migliore and peggiore}
\textit{buono} has the comparative \textit{migliore} (better), \textit{cattivo} has \textit{peggiore} (worse). \textit{più buono} and \textit{più cattivo} are also correct, especially about taste and character. The adverbs are \textit{meglio} (better) and \textit{peggio} (worse).
\end{notebox}

\subsection*{Worked examples}
\begin{examples}
\item\relax \textit{Che bel fiore!}\newline What a beautiful flower!
\item\relax \textit{Che begli alberi!}\newline What beautiful trees!
\item\relax \textit{Buon giorno a tutti!}\newline Good morning, everyone!\newline {\color{muted}\small Also written buongiorno.}
\item\relax \textit{La lepre è più veloce della tartaruga.}\newline The hare is faster than the tortoise.
\item\relax \textit{Il gatto è meno grande del cane.}\newline The cat is smaller (less big) than the dog.
\item\relax \textit{Questa torta è migliore di quella.}\newline This cake is better than that one.
\end{examples}

\subsection*{Exercises}\label{les:bello-buono-comparison:exercises}
{\color{muted}\small Answers: page~\pageref{ans:bello-buono-comparison}.}

\begin{exercises}
\item\label{ex:bello-buono-comparison-1}\relax Che \blank{} albero! (beautiful)\par\nopagebreak\begin{options}
\item\relax bello
\item\relax bella
\item\relax bel
\item\relax bell'
\end{options}
\item\label{ex:bello-buono-comparison-2}\relax un \blank{} giorno (a lovely day)\par\nopagebreak\begin{options}
\item\relax bell'
\item\relax bello
\item\relax bei
\item\relax bel
\end{options}
\item\label{ex:bello-buono-comparison-3}\relax I fiori sono \blank{}.\par\nopagebreak\begin{options}
\item\relax bei
\item\relax belli
\item\relax bel
\item\relax begli
\end{options}
\item\label{ex:bello-buono-comparison-4}\relax Complete the comparison.\par\smallskip\begingroup\setlength{\parindent}{0pt}La lepre è \clozeblank{1} veloce \clozeblank{2} tartaruga.\par\endgroup
\item\label{ex:bello-buono-comparison-5}\relax Translate: Luca is as tall as his father.\par\answerline
\item\label{ex:bello-buono-comparison-6}\relax Which word means 'better' as an adjective (a better cake)?\par\nopagebreak\begin{options}
\item\relax meglio
\item\relax buonissimo
\item\relax più bene
\item\relax migliore
\end{options}
\end{exercises}

\begin{summarybox}
\begin{itemize}
\item\relax bello before a noun: bel, bello, bell', bella, bei, begli, belle (like the definite article).
\item\relax buono before a singular noun: buon, buono, buona, buon' (like the indefinite article).
\item\relax più / meno + adjective + di (+ article); (così) … come for equality.
\item\relax migliore / peggiore (adjectives), meglio / peggio (adverbs).
\end{itemize}
\end{summarybox}

\relatedwords{\textit{bello} beautiful; lovely; handsome \textperiodcentered{} \textit{buono} good; tasty; kind \textperiodcentered{} \textit{veloce} fast; quick \textperiodcentered{} \textit{lento} slow \textperiodcentered{} \textit{la lepre} hare \textperiodcentered{} \textit{la tartaruga} tortoise; turtle \textperiodcentered{} \textit{alto} tall; high}
\applyit{Read \textit{La tartaruga lenta} (A1), \S\ref{story:la-tartaruga-lenta}. Study the rhyme \textit{Il valzer degli animali}, \S\ref{rhyme:il-valzer-degli-animali}.}

\chapter{Regular verbs in the present}\label{mod:m04}
\moduleintro{A0–A1}{A map of the whole verb system; infinitives, stems, subject pronouns, and the three regular conjugations.}{Map of the verb system (\S\ref{les:verb-system-overview}); Infinitives, stems and subject pronouns (\S\ref{les:infinitives-pronouns}); Present tense of -are verbs (\S\ref{les:present-are}); Present tense of -ere and -ire verbs (\S\ref{les:present-ere-ire})}

\section{Map of the verb system}\label{les:verb-system-overview}
\lessonmeta{A0–A1}{What a finite verb form tells you (person, number, tense, mood), how simple and compound tenses are built, and where each tense is taught. A map to come back to, not something to memorise now.}
\begin{objectives}
\item\relax Know the four finite moods (indicativo, congiuntivo, condizionale, imperativo) and the three non-finite forms (infinito, gerundio, participio).
\item\relax Understand that every compound tense = a form of avere or essere + the past participle.
\item\relax Read a verb label such as 'congiuntivo imperfetto, 3rd person singular' and know what each part means.
\item\relax Know where in the course each tense is studied in detail.
\end{objectives}

\subsection{What a verb form tells you}\label{les:verb-system-overview:what-a-form-encodes}

An Italian finite verb form carries four pieces of information in its ending: \textbf{person} (1st, 2nd, 3rd), \textbf{number} (singular, plural), \textbf{tense} (when, and how the event is viewed) and \textbf{mood} (how the speaker presents it: as fact, as possibility or wish, as a condition, as a command). \textit{parlavamo} = parl- (lemma \textit{parlare}) + -avamo (1st person plural, imperfetto, indicativo).\par

\begin{notebox}{rulecol}\notelabel{rulecol}{Rule}{Subject pronouns are usually omitted}
Because the ending shows the person, \textit{io, tu, lui …} are normally left out: \textit{Parliamo italiano.} They are used for emphasis, contrast or to avoid ambiguity (see the next lesson).
\end{notebox}

\begin{notebox}{tipcol}\notelabel{tipcol}{Tip}{One form, several analyses}
Some endings are shared between persons or tenses. \textit{parli} is tu in the presente indicativo, but io/tu/lui/lei in the congiuntivo presente; \textit{cantò} and \textit{canto} differ only by the accent. Context decides. The conjugation tables in Part IV let you check every possible analysis of an ambiguous form.
\end{notebox}

\subsection{Moods and non-finite forms}\label{les:verb-system-overview:moods}

{\small\setlength{\LTpre}{2pt}\setlength{\LTpost}{6pt}
\tablecaption{The moods of Italian}
\begin{xltabular}{\linewidth}{lLL}
\toprule
\textbf{Mood} & \textbf{Presents the event as …} & \textbf{Example} \\
\midrule
indicativo & a fact, a question about facts, a certain prediction & \textit{Luca \textbf{parla} italiano.} \\
congiuntivo & something thought, wished, feared, doubted or not yet real; mostly in subordinate clauses & \textit{Penso che Luca \textbf{parli} italiano.} \\
condizionale & dependent on a condition; also polite, attenuated or unconfirmed & \textit{Luca \textbf{parlerebbe} se potesse.} \\
imperativo & a command, request or instruction & \textit{\textbf{Parla} più piano!} \\
infinito & the action itself, without person (non-finite) & \textit{Voglio \textbf{parlare}.} \\
gerundio & an action accompanying or explaining another (non-finite) & \textit{Sto \textbf{parlando}.} \\
participio & the result of an action; builds compound tenses and the passive (non-finite) & \textit{Ho \textbf{parlato}.} \\
\bottomrule
\end{xltabular}}

\begin{notebox}{tendcol}\notelabel{tendcol}{Tendency}{Mood is partly a matter of register}
In some subordinate clauses the congiuntivo is required by the norm; in others informal speech often uses the indicativo instead (\textit{penso che è vero} for \textit{penso che sia vero}). The course always teaches the standard form first and labels the informal alternative as such.
\end{notebox}

\subsection{Simple and compound tenses}\label{les:verb-system-overview:simple-compound}

Each mood has simple tenses (one word) and compound tenses (two words). Every compound tense is built in the same way: the auxiliary \textit{avere} or \textit{essere} in the matching simple tense + the past participle.\par

{\small\setlength{\LTpre}{2pt}\setlength{\LTpost}{6pt}
\tablecaption{Each simple tense has a compound partner}
\begin{xltabular}{\linewidth}{LlL}
\toprule
\textbf{Simple tense of the auxiliary} & \textbf{Compound tense} & \textbf{Example (io)} \\
\midrule
presente: \textit{ho, sono} & passato prossimo & \textit{ho parlato, sono andato/a} \\
imperfetto: \textit{avevo, ero} & trapassato prossimo & \textit{avevo parlato, ero andato/a} \\
passato remoto: \textit{ebbi, fui} & trapassato remoto & \textit{ebbi parlato, fui andato/a} \\
futuro: \textit{avrò, sarò} & futuro anteriore & \textit{avrò parlato, sarò andato/a} \\
congiuntivo presente: \textit{abbia, sia} & congiuntivo passato & \textit{abbia parlato, sia andato/a} \\
congiuntivo imperfetto: \textit{avessi, fossi} & congiuntivo trapassato & \textit{avessi parlato, fossi andato/a} \\
condizionale presente: \textit{avrei, sarei} & condizionale passato & \textit{avrei parlato, sarei andato/a} \\
infinito: \textit{avere, essere} & infinito passato & \textit{aver(e) parlato, esser(e) andato/a} \\
gerundio: \textit{avendo, essendo} & gerundio passato & \textit{avendo parlato, essendo andato/a} \\
\bottomrule
\end{xltabular}}

\begin{notebox}{rulecol}\notelabel{rulecol}{Rule}{Compound tense = anteriority}
A compound tense normally presents the action as completed before the moment expressed by its simple partner: \textit{quando arrivo, lui è già uscito} (present → passato prossimo), \textit{quando arrivai, lui era già uscito} (past → trapassato), \textit{quando arriverò, lui sarà già uscito} (future → futuro anteriore).
\end{notebox}

\begin{notebox}{rulecol}\notelabel{rulecol}{Rule}{Two auxiliaries}
Transitive verbs (with a direct object) take \textit{avere}; most verbs of movement or change of state, reflexive verbs and the passive take \textit{essere}, and then the participle agrees with the subject: \textit{Sara è andata}. Several verbs can take either, with a difference in meaning. This is covered fully in the lesson on auxiliaries and agreement.
\end{notebox}

\subsection{The full map}\label{les:verb-system-overview:map}

The table lists every tense with the io form of \textit{parlare} (auxiliary avere) and \textit{andare} (auxiliary essere), its main use, the approximate level at which it is studied, and the lesson where it is taught. The levels refer to when you are likely to \textit{produce} a form; you will meet every tense in reading earlier than that, and the conjugation tables in Part IV let you identify any form you meet.\par

{\footnotesize\setlength{\LTpre}{2pt}
\tablecaption{The Italian verb system}
\begin{xltabular}{\linewidth}{@{}>{\raggedright\arraybackslash}p{1.6cm} >{\raggedright\arraybackslash}p{2.1cm} l l L l l@{}}
\toprule
\textbf{Mood} & \textbf{Tense} & \textbf{\textit{parlare}} (io) & \textbf{\textit{andare}} (io) & \textbf{Main use} & \textbf{Level} & \textbf{Lesson} \\
\midrule
\endhead
indicativo & \textbf{presente} & \textit{parlo} & \textit{vado} & Facts, habits and actions in progress; also the near future (domani parto) and, in narration, the "historic present". & A0 & \S\ref{les:present-are} \\
indicativo & \textbf{passato prossimo} {\color{muted}(compound)} & \textit{ho parlato} & \textit{sono andato} & Completed past events, especially with a link to the present. In the North it is the ordinary spoken past; see also passato remoto. & A1 & \S\ref{les:passato-prossimo-avere} \\
indicativo & \textbf{imperfetto} & \textit{parlavo} & \textit{andavo} & Background, description, habits and actions in progress in the past; also polite requests (volevo chiedere) and, colloquially, counterfactuals. & A2 & \S\ref{les:imperfetto} \\
indicativo & \textbf{trapassato prossimo} {\color{muted}(compound)} & \textit{avevo parlato} & \textit{ero andato} & An action completed before another past moment: quando sono arrivato, il treno era già partito. & A2 & \S\ref{les:trapassato-prossimo} \\
indicativo & \textbf{passato remoto} & \textit{parlai} & \textit{andai} & Completed past events viewed as cut off from the present: standard in written narration, history and fairy tales; in speech its use varies strongly by region. & B1 & \S\ref{les:passato-remoto} \\
indicativo & \textbf{trapassato remoto} {\color{muted}(compound)} & \textit{ebbi parlato} & \textit{fui andato} & Literary: immediate anteriority in a temporal clause (dopo che, appena, quando) whose main verb is in the passato remoto. & B2 & \S\ref{les:trapassato-remoto} \\
indicativo & \textbf{futuro semplice} & \textit{parlerò} & \textit{andrò} & Future events; also supposition about the present (sarà stanco = he is probably tired). & A2 & \S\ref{les:futuro-semplice} \\
indicativo & \textbf{futuro anteriore} {\color{muted}(compound)} & \textit{avrò parlato} & \textit{sarò andato} & An action completed before a future moment; also supposition about the past (avrà perso il treno). & B1 & \S\ref{les:futuro-anteriore} \\
congiuntivo & \textbf{congiuntivo presente} & \textit{parli} & \textit{vada} & In subordinate clauses after verbs of opinion, wish, emotion and doubt, impersonal expressions and certain conjunctions, when the main verb is present or future. & B1 & \S\ref{les:congiuntivo-presente} \\
congiuntivo & \textbf{congiuntivo passato} {\color{muted}(compound)} & \textit{abbia parlato} & \textit{sia andato} & As the congiuntivo presente, for an action prior to the main verb: penso che sia partito. & B1 & \S\ref{les:congiuntivo-passato} \\
congiuntivo & \textbf{congiuntivo imperfetto} & \textit{parlassi} & \textit{andassi} & Subjunctive contexts with a past or conditional main verb; the "se" clause of possible hypotheses; wishes (magari fosse vero). & B1 & \S\ref{les:congiuntivo-imperfetto-trapassato} \\
congiuntivo & \textbf{congiuntivo trapassato} {\color{muted}(compound)} & \textit{avessi parlato} & \textit{fossi andato} & Anteriority with a past main verb; the "se" clause of counterfactual hypotheses (se avessi saputo…). & B2 & \S\ref{les:congiuntivo-imperfetto-trapassato} \\
condizionale & \textbf{condizionale presente} & \textit{parlerei} & \textit{andrei} & Politeness, wishes, advice, hypothetical consequences, and reported unconfirmed information. & A2 & \S\ref{les:condizionale-presente} \\
condizionale & \textbf{condizionale passato} {\color{muted}(compound)} & \textit{avrei parlato} & \textit{sarei andato} & Unrealised past possibilities; and the "future in the past": disse che sarebbe venuto. & B1 & \S\ref{les:condizionale-passato} \\
imperativo & \textbf{imperativo} & \textit{parla (tu)} & \textit{va' / vai (tu)} & Commands, requests and instructions; formal Lei/Loro forms are taken from the congiuntivo presente; negative tu = non + infinitive. & A2 & \S\ref{les:imperativo} \\
\bottomrule
\end{xltabular}}

\subsection*{Worked examples}
\begin{examples}
\item\relax \textit{Parliamo italiano.}\newline We speak Italian.\newline {\color{muted}\small presente indicativo, 1st person plural; the subject noi is understood.}
\item\relax \textit{Ieri abbiamo parlato con Sara.}\newline Yesterday we spoke with Sara.\newline {\color{muted}\small passato prossimo = presente of avere + participle.}
\item\relax \textit{Spero che Sara parli con te.}\newline I hope Sara talks to you.\newline {\color{muted}\small congiuntivo presente after sperare.}
\item\relax \textit{Parlerei con lei, se potessi.}\newline I would speak to her if I could.\newline {\color{muted}\small condizionale presente + congiuntivo imperfetto.}
\item\relax \textit{Parla più piano, per favore!}\newline Speak more slowly, please!\newline {\color{muted}\small imperativo, tu.}
\item\relax \textit{Sara è andata a scuola.}\newline Sara went to school.\newline {\color{muted}\small andare takes essere; andata agrees with Sara.}
\item\relax \textit{Sto parlando al telefono.}\newline I'm talking on the phone.\newline {\color{muted}\small stare + gerundio: an action in progress.}
\end{examples}

\subsection*{Exercises}\label{les:verb-system-overview:exercises}
{\color{muted}\small Answers: page~\pageref{ans:verb-system-overview}.}

\begin{exercises}
\item\label{ex:verb-system-overview-1}\relax Match each mood with how it presents the event.\par\smallskip\begin{tabularx}{\linewidth}{@{}L L@{}}
1.\ indicativo & (a)\ as dependent on a condition \\
2.\ congiuntivo & (b)\ as a fact \\
3.\ condizionale & (c)\ as something thought, wished or doubted \\
4.\ imperativo & (d)\ as a command \\
\end{tabularx}
\item\label{ex:verb-system-overview-2}\relax Every compound tense is formed with …\par\nopagebreak\begin{options}
\item\relax the infinitive + an ending
\item\relax avere or essere + past participle
\item\relax the present + a time word
\item\relax stare + gerund
\end{options}
\item\label{ex:verb-system-overview-3}\relax Match the simple tense of the auxiliary with the compound tense it forms.\par\smallskip\begin{tabularx}{\linewidth}{@{}L L@{}}
1.\ ho / sono & (a)\ passato prossimo \\
2.\ avevo / ero & (b)\ condizionale passato \\
3.\ avrò / sarò & (c)\ trapassato prossimo \\
4.\ avrei / sarei & (d)\ congiuntivo passato \\
5.\ abbia / sia & (e)\ futuro anteriore \\
\end{tabularx}
\item\label{ex:verb-system-overview-4}\relax \textit{parlavamo}: which information does the ending -avamo carry?\par\nopagebreak\begin{options}
\item\relax 1st person plural, futuro
\item\relax 1st person plural, imperfetto, indicativo
\item\relax 3rd person plural, congiuntivo
\item\relax 2nd person plural, presente, indicativo
\end{options}
\item\label{ex:verb-system-overview-5}\relax Which of these is a non-finite form (it shows no person)?\par\nopagebreak\begin{options}
\item\relax parlino
\item\relax parlano
\item\relax parlando
\item\relax parlerò
\end{options}
\item\label{ex:verb-system-overview-6}\relax \textit{parli} can be …\par\nopagebreak\begin{options}
\item\relax only tu (indicativo)
\item\relax only io (congiuntivo)
\item\relax tu (indicativo) or io/tu/lui/lei (congiuntivo presente)
\item\relax the imperativo of noi
\end{options}
\end{exercises}

\begin{summarybox}
\begin{itemize}
\item\relax A finite form encodes person, number, tense and mood.
\item\relax Four finite moods: indicativo, congiuntivo, condizionale, imperativo; three non-finite forms: infinito, gerundio, participio.
\item\relax Compound tense = auxiliary (avere/essere) in the matching simple tense + participle; it usually expresses anteriority.
\item\relax Use this map to locate any form you meet in reading.
\end{itemize}
\end{summarybox}

\relatedwords{\textit{parlare} to speak; to talk \textperiodcentered{} \textit{andare} to go \textperiodcentered{} \textit{essere} to be \textperiodcentered{} \textit{avere} to have}

\section{Infinitives, stems and subject pronouns}\label{les:infinitives-pronouns}
\lessonmeta{A0}{The three verb families, how to find the stem, the subject pronouns, and why Italian usually omits them.}
\begin{objectives}
\item\relax Recognise the three conjugation families: -are, -ere, -ire.
\item\relax Split an infinitive into stem and ending.
\item\relax Know the subject pronouns, including formal Lei.
\item\relax Understand subject omission.
\end{objectives}

\subsection{Three families}\label{les:infinitives-pronouns:families}

The dictionary form of a verb is the infinitive. Almost every infinitive ends in \textit{-are}, \textit{-ere} or \textit{-ire}, and that ending tells you which set of endings the verb uses. A small group ends in \textit{-rre}: \textit{porre} (to put), \textit{trarre} (to draw), \textit{tradurre} (to translate), \textit{condurre} (to lead) and their compounds (\textit{proporre, comporre, produrre, ridurre} …). They are contracted -ere verbs and take -ere endings on a longer stem: \textit{tradurre → traduco, traduceva} (from Latin \textit{traducere}). \textit{fare, dire} and \textit{bere} also conjugate on longer stems (\textit{facevo, dicevo, bevevo}); see \lessonref{other-irregulars}{other irregular verbs}.\par

{\small\setlength{\LTpre}{2pt}\setlength{\LTpost}{6pt}
\tablecaption{Infinitive = stem + ending}
\begin{longtable}{lllll}
\toprule
\textbf{Infinitive} & \textbf{Stem} & \textbf{Ending} & \textbf{Family} & \textbf{Meaning} \\
\midrule
parlare & parl- & -are & first & to speak \\
vedere & ved- & -ere & second & to see \\
dormire & dorm- & -ire & third & to sleep \\
\bottomrule
\end{longtable}}

\begin{notebox}{rulecol}\notelabel{rulecol}{Rule}{Conjugating = stem + personal ending}
To conjugate a regular verb, remove the infinitive ending and add the ending for the person: \textit{parl-} + \textit{-o} = \textit{parlo} (I speak).
\end{notebox}

\begin{notebox}{tendcol}\notelabel{tendcol}{Tendency}{Which family is largest?}
Most verbs — and almost all new verbs — belong to -are. -ere verbs include many of the most irregular verbs. -ire verbs split into two groups (dormire type and capire type), studied two lessons from now.
\end{notebox}

\subsection{Subject pronouns}\label{les:infinitives-pronouns:pronouns}

{\small\setlength{\LTpre}{2pt}\setlength{\LTpost}{6pt}
\tablecaption{Subject pronouns}
\begin{xltabular}{\linewidth}{llllL}
\toprule
\textbf{Person} & \textbf{Singular} & \textbf{} & \textbf{Plural} & \textbf{} \\
\midrule
1st & io & I & noi & we \\
2nd & tu & you (informal) & voi & you (plural) \\
3rd & lui / lei & he / she & loro & they \\
formal & Lei & you (formal, one person) & (voi) & you (plural, normally voi) \\
\bottomrule
\end{xltabular}}

\textit{io} is written with a small letter except at the start of a sentence. Formal \textit{Lei} uses the same verb form as \textit{lui/lei}: \textit{Lei parla italiano?} = Do you speak Italian? (formal).\par

\subsection{Subject omission}\label{les:infinitives-pronouns:omission}

\begin{notebox}{rulecol}\notelabel{rulecol}{Rule}{The ending shows the subject}
Because each person has a distinct verb ending, subject pronouns are normally left out: \textit{Parlo italiano} = I speak Italian. Adding \textit{io} is not wrong, but it adds emphasis or contrast: \textit{Io parlo italiano, lui no} = I speak Italian; he doesn't.
\end{notebox}

\begin{notebox}{tendcol}\notelabel{tendcol}{Tendency}{When pronouns are used}
Use the pronoun for contrast, emphasis, after \textit{anche} (\textit{anche io}), when the subject would otherwise be unclear, or alone without a verb (\textit{Chi è? — Io!}).
\end{notebox}

\subsection*{Worked examples}
\begin{examples}
\item\relax \textit{Parlo italiano.}\newline I speak Italian.\newline {\color{muted}\small No pronoun needed: -o shows io.}
\item\relax \textit{Io parlo italiano, lui parla inglese.}\newline I speak Italian; he speaks English.\newline {\color{muted}\small Pronouns for contrast.}
\item\relax \textit{Anche io dormo.}\newline I'm sleeping too.
\item\relax \textit{Lei abita a Roma?}\newline Do you live in Rome? (formal)\newline {\color{muted}\small Formal Lei with third-person verb.}
\item\relax \textit{Noi camminiamo nel bosco.}\newline We walk in the wood.
\item\relax \textit{Chi canta? — Io!}\newline Who is singing? — Me!
\end{examples}

\subsection*{Exercises}\label{les:infinitives-pronouns:exercises}
{\color{muted}\small Answers: page~\pageref{ans:infinitives-pronouns}.}

\begin{exercises}
\item\label{ex:infinitives-pronouns-1}\relax What is the stem of \textit{cantare}?\par\nopagebreak\begin{options}
\item\relax canta-
\item\relax cant-
\item\relax cantar-
\item\relax can-
\end{options}
\item\label{ex:infinitives-pronouns-2}\relax Match each infinitive to its family.\par\smallskip\begin{tabularx}{\linewidth}{@{}L L@{}}
1.\ leggere & (a)\ -are \\
2.\ aprire & (b)\ -ere \\
3.\ guardare & (c)\ -ere \\
4.\ credere & (d)\ -ire \\
5.\ sentire & (e)\ -ire \\
\end{tabularx}
\item\label{ex:infinitives-pronouns-3}\relax Match the pronouns.\par\smallskip\begin{tabularx}{\linewidth}{@{}L L@{}}
1.\ io & (a)\ they \\
2.\ noi & (b)\ you (plural) \\
3.\ voi & (c)\ I \\
4.\ loro & (d)\ we \\
5.\ Lei & (e)\ you (formal) \\
\end{tabularx}
\item\label{ex:infinitives-pronouns-4}\relax Why does \textit{Parlo italiano} not need \textit{io}?\par\nopagebreak\begin{options}
\item\relax It is a question.
\item\relax The ending -o already shows the subject is io.
\item\relax Because italiano is the subject.
\item\relax io is never used in Italian.
\end{options}
\item\label{ex:infinitives-pronouns-5}\relax In which sentence is the pronoun really needed?\par\nopagebreak\begin{options}
\item\relax Noi camminiamo.
\item\relax Io parlo italiano.
\item\relax Tu dormi.
\item\relax Io parlo italiano, lui parla inglese.
\end{options}
\item\label{ex:infinitives-pronouns-6}\relax Write the infinitive whose stem is \textit{dorm-} and which belongs to the third family.\par\answerline
\end{exercises}

\begin{summarybox}
\begin{itemize}
\item\relax Infinitives end in -are, -ere or -ire; stem + personal ending = conjugated form.
\item\relax Pronouns: io, tu, lui/lei, noi, voi, loro; formal Lei uses the third-person singular.
\item\relax Subject pronouns are usually omitted; use them for contrast or emphasis.
\end{itemize}
\end{summarybox}

\relatedwords{\textit{parlare} to speak; to talk \textperiodcentered{} \textit{vedere} to see \textperiodcentered{} \textit{dormire} to sleep \textperiodcentered{} \textit{io} I \textperiodcentered{} \textit{tu} you (singular, informal) \textperiodcentered{} \textit{lui} he; him \textperiodcentered{} \textit{lei} she; her; you (formal) \textperiodcentered{} \textit{noi} we; us \textperiodcentered{} \textit{voi} you (plural) \textperiodcentered{} \textit{loro} they; them; their}
\applyit{Read \textit{La famiglia di Marco} (A0–A1), \S\ref{story:la-famiglia-di-marco}.}

\section{Present tense of -are verbs}\label{les:present-are}
\lessonmeta{A0}{The endings -o, -i, -a, -iamo, -ate, -ano, with the spelling adjustments of -care/-gare and -iare verbs.}
\begin{objectives}
\item\relax Conjugate any regular -are verb in the present.
\item\relax Apply the h of cercare/giocare and the single i of mangiare/studiare.
\item\relax Understand the meanings of the Italian present tense.
\end{objectives}

\subsection{The endings}\label{les:present-are:endings}

\par\noindent\begin{minipage}[t]{0.48\linewidth}\raggedright
\paradigmbox{parlare --- to speak; to talk \textperiodcentered{} indicativo presente}{%
\begin{tabular}{@{}l l l@{}}
\textit{io} & parl\hlending{o} & {\color{muted}I speak} \\
\textit{tu} & parl\hlending{i} & {\color{muted}you speak} \\
\textit{lui/lei} & parl\hlending{a} & {\color{muted}he/she speaks} \\
\textit{noi} & parl\hlending{iamo} & {\color{muted}we speak} \\
\textit{voi} & parl\hlending{ate} & {\color{muted}you (pl.) speak} \\
\textit{loro} & parl\hlending{ano} & {\color{muted}they speak} \\
\end{tabular}}{Stem \textit{parl-} + \hlending{ending}.}
\end{minipage}\par\medskip

\begin{notebox}{rulecol}\notelabel{rulecol}{Rule}{Stress}
In \textit{io, tu, lui/lei} and \textit{loro} the stress falls on the stem: \textit{PAR-lo, PAR-li, PAR-la, PAR-la-no}. In \textit{noi} and \textit{voi} it moves to the ending: \textit{par-LIA-mo, par-LA-te}. The \textit{loro} form keeps the stress on the same syllable as the \textit{io} form: \textit{àbito → àbitano} (accents shown here only to mark the stress; they are not written).
\end{notebox}

\subsection{Spelling adjustments}\label{les:present-are:spelling}

\begin{notebox}{rulecol}\notelabel{rulecol}{Rule}{-care and -gare add h}
To keep c and g hard before the endings -i and -iamo, add h: \textit{cercare → cerchi, cerchiamo}; \textit{giocare → giochi, giochiamo}.
\end{notebox}

\par\noindent\begin{minipage}[t]{0.48\linewidth}\raggedright
\paradigmbox{giocare --- to play \textperiodcentered{} indicativo presente}{%
\begin{tabular}{@{}l l l@{}}
\textit{io} & gioco & {\color{muted}I play} \\
\textit{tu} & giochi & {\color{muted}you play} \\
\textit{lui/lei} & gioca & {\color{muted}he/she plays} \\
\textit{noi} & giochiamo & {\color{muted}we play} \\
\textit{voi} & giocate & {\color{muted}you (pl.) play} \\
\textit{loro} & giocano & {\color{muted}they play} \\
\end{tabular}}{}
\end{minipage}\par\medskip

\begin{notebox}{rulecol}\notelabel{rulecol}{Rule}{-iare verbs do not double the i}
If the stem already ends in -i, do not add another one in \textit{tu} and \textit{noi}: \textit{mangiare → mangi, mangiamo}; \textit{studiare → studi, studiamo}. (A few verbs with a stressed i, like \textit{inviare → invii}, keep both.)
\end{notebox}

\par\noindent\begin{minipage}[t]{0.48\linewidth}\raggedright
\paradigmbox{mangiare --- to eat \textperiodcentered{} indicativo presente}{%
\begin{tabular}{@{}l l l@{}}
\textit{io} & mangio & {\color{muted}I eat} \\
\textit{tu} & mangi & {\color{muted}you eat} \\
\textit{lui/lei} & mangia & {\color{muted}he/she eats} \\
\textit{noi} & mangiamo & {\color{muted}we eat} \\
\textit{voi} & mangiate & {\color{muted}you (pl.) eat} \\
\textit{loro} & mangiano & {\color{muted}they eat} \\
\end{tabular}}{}
\end{minipage}\par\medskip

\subsection{What the present tense means}\label{les:present-are:meaning}

\begin{itemize}
\item\relax Habitual actions: \textit{Lavoro in città} = I work in town.
\item\relax Actions happening now: \textit{Che cosa guardi?} = What are you looking at?
\item\relax The near future, especially with a time expression: \textit{Domani compro il pane} = Tomorrow I'll buy bread.
\end{itemize}

\subsection*{Worked examples}
\begin{examples}
\item\relax \textit{Parlo con la nonna.}\newline I'm talking with Grandma.
\item\relax \textit{Che cosa cerchi?}\newline What are you looking for?\newline {\color{muted}\small cercare: tu cerchi.}
\item\relax \textit{Giochiamo in giardino.}\newline We play in the garden.\newline {\color{muted}\small giocare: noi giochiamo.}
\item\relax \textit{I bambini mangiano la pizza.}\newline The children are eating pizza.
\item\relax \textit{Studi l'italiano?}\newline Do you study Italian?\newline {\color{muted}\small studiare: tu studi (one i).}
\item\relax \textit{Domani compriamo il pane.}\newline Tomorrow we'll buy bread.\newline {\color{muted}\small Present for the near future.}
\end{examples}

\subsection*{Exercises}\label{les:present-are:exercises}
{\color{muted}\small Answers: page~\pageref{ans:present-are}.}

\begin{exercises}
\item\label{ex:present-are-1}\relax Conjugate \textit{cantare} in the present.\par\smallskip{\begin{tabular}{@{}l l@{}}
io & \blankcell \\
tu & \blankcell \\
lui/lei & \blankcell \\
noi & \blankcell \\
voi & \blankcell \\
loro & \blankcell \\
\end{tabular}}
\item\label{ex:present-are-2}\relax cercare, tu: \blank{}
\item\label{ex:present-are-3}\relax mangiare, noi: \blank{}
\item\label{ex:present-are-4}\relax giocare, noi: \blank{}
\item\label{ex:present-are-5}\relax Complete with the present tense.\par\smallskip\begingroup\setlength{\parindent}{0pt}I bambini \clozeblank{1} (giocare) in giardino e la mamma \clozeblank{2} (guardare) dalla finestra. Noi \clozeblank{3} (ascoltare) la musica.\par\endgroup
\item\label{ex:present-are-6}\relax Translate: What are you looking for? (tu)\par\answerline
\item\label{ex:present-are-7}\relax What is the infinitive of \textit{studiamo}?\par\answerline
\end{exercises}

\begin{summarybox}
\begin{itemize}
\item\relax -are endings: -o, -i, -a, -iamo, -ate, -ano.
\item\relax -care/-gare: add h before i (cerchi, giochiamo).
\item\relax -iare: no double i (mangi, mangiamo).
\item\relax The present covers 'I eat', 'I am eating' and often 'I'll eat (soon)'.
\end{itemize}
\end{summarybox}

\relatedwords{\textit{parlare} to speak; to talk \textperiodcentered{} \textit{cantare} to sing \textperiodcentered{} \textit{cercare} to look for; to try \textperiodcentered{} \textit{giocare} to play \textperiodcentered{} \textit{mangiare} to eat \textperiodcentered{} \textit{studiare} to study \textperiodcentered{} \textit{guardare} to look at; to watch \textperiodcentered{} \textit{ascoltare} to listen to \textperiodcentered{} \textit{comprare} to buy \textperiodcentered{} \textit{lavorare} to work \textperiodcentered{} \textit{abitare} to live (reside)}
\applyit{Read \textit{La colazione} (A0–A1), \S\ref{story:la-colazione}. Read \textit{La famiglia di Marco} (A0–A1), \S\ref{story:la-famiglia-di-marco}. Study the rhyme \textit{Il valzer degli animali}, \S\ref{rhyme:il-valzer-degli-animali}.}

\section{Present tense of -ere and -ire verbs}\label{les:present-ere-ire}
\lessonmeta{A1}{vedere and dormire, and the -isc- verbs like capire.}
\begin{objectives}
\item\relax Conjugate regular -ere and -ire verbs.
\item\relax Recognise -isc- verbs and conjugate them.
\item\relax Compare the three families side by side.
\end{objectives}

\subsection{-ere verbs}\label{les:present-ere-ire:ere}

\par\noindent\begin{minipage}[t]{0.48\linewidth}\raggedright
\paradigmbox{vedere --- to see \textperiodcentered{} indicativo presente}{%
\begin{tabular}{@{}l l l@{}}
\textit{io} & ved\hlending{o} & {\color{muted}I see} \\
\textit{tu} & ved\hlending{i} & {\color{muted}you see} \\
\textit{lui/lei} & ved\hlending{e} & {\color{muted}he/she sees} \\
\textit{noi} & ved\hlending{iamo} & {\color{muted}we see} \\
\textit{voi} & ved\hlending{ete} & {\color{muted}you (pl.) see} \\
\textit{loro} & ved\hlending{ono} & {\color{muted}they see} \\
\end{tabular}}{Stem \textit{ved-} + \hlending{ending}.}
\end{minipage}\par\medskip

\begin{notebox}{rulecol}\notelabel{rulecol}{Rule}{Hard and soft sounds may alternate}
Because endings change, the pronunciation of c/g in the stem can change too: \textit{leggere → leggo} (hard g), \textit{leggi} (soft g); \textit{vincere → vinco, vinci}. The spelling stays regular — no h is added in -ere verbs.
\end{notebox}

\subsection{-ire verbs: two groups}\label{les:present-ere-ire:ire}

\par\noindent\begin{minipage}[t]{0.48\linewidth}\raggedright
\paradigmbox{dormire --- to sleep \textperiodcentered{} indicativo presente}{%
\begin{tabular}{@{}l l l@{}}
\textit{io} & dorm\hlending{o} & {\color{muted}I sleep} \\
\textit{tu} & dorm\hlending{i} & {\color{muted}you sleep} \\
\textit{lui/lei} & dorm\hlending{e} & {\color{muted}he/she sleeps} \\
\textit{noi} & dorm\hlending{iamo} & {\color{muted}we sleep} \\
\textit{voi} & dorm\hlending{ite} & {\color{muted}you (pl.) sleep} \\
\textit{loro} & dorm\hlending{ono} & {\color{muted}they sleep} \\
\end{tabular}}{Stem \textit{dorm-} + \hlending{ending}.}
\end{minipage}\par\medskip

\begin{notebox}{rulecol}\notelabel{rulecol}{Rule}{The -isc- group}
Many -ire verbs insert \textit{-isc-} between stem and ending in \textit{io, tu, lui/lei} and \textit{loro}. Noi and voi are normal.
\end{notebox}

\par\noindent\begin{minipage}[t]{0.48\linewidth}\raggedright
\paradigmbox{capire --- to understand \textperiodcentered{} indicativo presente}{%
\begin{tabular}{@{}l l l@{}}
\textit{io} & cap\hlending{isco} & {\color{muted}I understand} \\
\textit{tu} & cap\hlending{isci} & {\color{muted}you understand} \\
\textit{lui/lei} & cap\hlending{isce} & {\color{muted}he/she understands} \\
\textit{noi} & cap\hlending{iamo} & {\color{muted}we understand} \\
\textit{voi} & cap\hlending{ite} & {\color{muted}you (pl.) understand} \\
\textit{loro} & cap\hlending{iscono} & {\color{muted}they understand} \\
\end{tabular}}{Stem \textit{cap-} + \hlending{ending}.}
\end{minipage}\par\medskip

\begin{notebox}{tendcol}\notelabel{tendcol}{Tendency}{You cannot see which group a verb belongs to}
The infinitive does not show whether a verb takes -isc-. Common -isc- verbs: \textit{capire, finire, preferire, pulire, costruire, spedire}. Common verbs without it: \textit{dormire, partire, sentire, aprire, seguire, offrire}. The dictionary entry tells you.
\end{notebox}

\subsection{Three families compared}\label{les:present-ere-ire:comparison}

{\small\setlength{\LTpre}{2pt}\setlength{\LTpost}{6pt}
\tablecaption{Present-tense endings}
\begin{longtable}{lllll}
\toprule
\textbf{} & \textbf{-are} & \textbf{-ere} & \textbf{-ire} & \textbf{-ire (-isc-)} \\
\midrule
io & -o & -o & -o & -isco \\
tu & -i & -i & -i & -isci \\
lui/lei & -a & -e & -e & -isce \\
noi & -iamo & -iamo & -iamo & -iamo \\
voi & -ate & -ete & -ite & -ite \\
loro & -ano & -ono & -ono & -iscono \\
\bottomrule
\end{longtable}}

Note what is shared: \textit{io -o}, \textit{tu -i} and \textit{noi -iamo} are the same in all families. The \textit{voi} ending simply uses the infinitive's vowel: \textit{-ate, -ete, -ite}.\par

\subsection*{Worked examples}
\begin{examples}
\item\relax \textit{Vedo il mare.}\newline I see the sea.
\item\relax \textit{Leggi un libro?}\newline Are you reading a book?
\item\relax \textit{Il gatto dorme sul letto.}\newline The cat sleeps on the bed.
\item\relax \textit{Capisco l'italiano.}\newline I understand Italian.\newline {\color{muted}\small -isc-: capisco.}
\item\relax \textit{Gli uccelli costruiscono il nido.}\newline The birds build their nest.\newline {\color{muted}\small -isc-: costruiscono.}
\item\relax \textit{Partiamo domani.}\newline We leave tomorrow.
\end{examples}

\subsection*{Exercises}\label{les:present-ere-ire:exercises}
{\color{muted}\small Answers: page~\pageref{ans:present-ere-ire}.}

\begin{exercises}
\item\label{ex:present-ere-ire-1}\relax Conjugate \textit{prendere} in the present.\par\smallskip{\begin{tabular}{@{}l l@{}}
io & \blankcell \\
tu & \blankcell \\
lui/lei & \blankcell \\
noi & \blankcell \\
voi & \blankcell \\
loro & \blankcell \\
\end{tabular}}
\item\label{ex:present-ere-ire-2}\relax Conjugate \textit{finire} (an -isc- verb).\par\smallskip{\begin{tabular}{@{}l l@{}}
io & \blankcell \\
tu & \blankcell \\
lui/lei & \blankcell \\
noi & \blankcell \\
voi & \blankcell \\
loro & \blankcell \\
\end{tabular}}
\item\label{ex:present-ere-ire-3}\relax dormire, voi: \blank{}
\item\label{ex:present-ere-ire-4}\relax Gli uccelli \blank{} il nido. (costruire)\par\nopagebreak\begin{options}
\item\relax costruiscano
\item\relax costruisciono
\item\relax costruono
\item\relax costruiscono
\end{options}
\item\label{ex:present-ere-ire-5}\relax Infinitive of \textit{capisce}?\par\answerline
\item\label{ex:present-ere-ire-6}\relax Complete.\par\smallskip\begingroup\setlength{\parindent}{0pt}Io \clozeblank{1} (leggere) e tu \clozeblank{2} (scrivere). Loro \clozeblank{3} (partire) domani.\par\endgroup
\item\label{ex:present-ere-ire-7}\relax Translate: \textit{Preferisco il tè.}\par\answerline
\end{exercises}

\begin{summarybox}
\begin{itemize}
\item\relax -ere: -o, -i, -e, -iamo, -ete, -ono.
\item\relax -ire: -o, -i, -e, -iamo, -ite, -ono.
\item\relax -isc- verbs: -isco, -isci, -isce, -iamo, -ite, -iscono (capire, finire, preferire, pulire, costruire).
\item\relax io -o, tu -i, noi -iamo are shared by all three families.
\end{itemize}
\end{summarybox}

\relatedwords{\textit{vedere} to see \textperiodcentered{} \textit{leggere} to read \textperiodcentered{} \textit{prendere} to take; to catch; to have (food) \textperiodcentered{} \textit{scrivere} to write \textperiodcentered{} \textit{dormire} to sleep \textperiodcentered{} \textit{partire} to leave; to set off \textperiodcentered{} \textit{sentire} to hear; to feel \textperiodcentered{} \textit{capire} to understand \textperiodcentered{} \textit{finire} to finish; to end \textperiodcentered{} \textit{preferire} to prefer \textperiodcentered{} \textit{costruire} to build \textperiodcentered{} \textit{pulire} to clean}
\applyit{Read \textit{Il gatto e la luna} (A0–A1), \S\ref{story:il-gatto-e-la-luna}. Read \textit{Nel bosco} (A1), \S\ref{story:nel-bosco}. Study the rhyme \textit{Ninna nanna della luna}, \S\ref{rhyme:ninna-nanna-della-luna}.}

\chapter{Essential irregular verbs}\label{mod:m05}
\moduleintro{A0–A1}{essere, avere, andare, fare, stare, dire, and other high-frequency irregular verbs.}{essere and avere (\S\ref{les:essere-avere}); andare, fare, stare, dire (\S\ref{les:andare-fare-stare-dire}); venire, uscire, bere, dare and piacere (\S\ref{les:other-irregulars})}

\section{essere and avere}\label{les:essere-avere}
\lessonmeta{A0}{The two most important verbs: their present tense, their main uses, and expressions with avere.}
\begin{objectives}
\item\relax Conjugate essere and avere in the present.
\item\relax Use essere for identity, description, origin and location.
\item\relax Use avere for possession, age and sensations (fame, sete, sonno, paura).
\end{objectives}

\subsection{essere — to be}\label{les:essere-avere:essere}

\par\noindent\begin{minipage}[t]{0.48\linewidth}\raggedright
\paradigmbox{essere --- to be \textperiodcentered{} indicativo presente}{%
\begin{tabular}{@{}l l l@{}}
\textit{io} & \irr{sono} & {\color{muted}I am} \\
\textit{tu} & \irr{sei} & {\color{muted}you are} \\
\textit{lui/lei} & \irr{è} & {\color{muted}he/she is} \\
\textit{noi} & \irr{siamo} & {\color{muted}we are} \\
\textit{voi} & \irr{siete} & {\color{muted}you (pl.) are} \\
\textit{loro} & \irr{sono} & {\color{muted}they are} \\
\end{tabular}}{\irr{Highlighted} forms differ from the regular pattern of the verb's family.}
\end{minipage}\par\medskip

\begin{notebox}{rulecol}\notelabel{rulecol}{Rule}{sono has two readings}
\textit{sono} is both 'I am' and 'they are'. Context or the subject decides: \textit{Sono stanco} (I am tired) — \textit{I bambini sono stanchi} (The children are tired).
\end{notebox}

\begin{itemize}
\item\relax Identity: \textit{Sono Marco.}
\item\relax Description: \textit{Il mare è calmo.}
\item\relax Origin: \textit{Sono di Roma} (I'm from Rome).
\item\relax Location: \textit{Il gatto è in cucina.}
\item\relax There is / there are: \textit{c'è, ci sono}.
\end{itemize}

\subsection{avere — to have}\label{les:essere-avere:avere}

\par\noindent\begin{minipage}[t]{0.48\linewidth}\raggedright
\paradigmbox{avere --- to have \textperiodcentered{} indicativo presente}{%
\begin{tabular}{@{}l l l@{}}
\textit{io} & \irr{ho} & {\color{muted}I have} \\
\textit{tu} & \irr{hai} & {\color{muted}you have} \\
\textit{lui/lei} & \irr{ha} & {\color{muted}he/she has} \\
\textit{noi} & \irr{abbiamo} & {\color{muted}we have} \\
\textit{voi} & avete & {\color{muted}you (pl.) have} \\
\textit{loro} & \irr{hanno} & {\color{muted}they have} \\
\end{tabular}}{\irr{Highlighted} forms differ from the regular pattern of the verb's family.}
\end{minipage}\par\medskip

\begin{notebox}{rulecol}\notelabel{rulecol}{Rule}{The silent h}
\textit{ho, hai, ha, hanno} are written with h to distinguish them from \textit{o} (or), \textit{ai} (to the), \textit{a} (to), \textit{anno} (year). The h is never pronounced.
\end{notebox}

{\small\setlength{\LTpre}{2pt}\setlength{\LTpost}{6pt}
\tablecaption{avere where English uses 'to be'}
\begin{longtable}{ll}
\toprule
\textbf{Italian} & \textbf{English} \\
\midrule
avere fame & to be hungry \\
avere sete & to be thirsty \\
avere sonno & to be sleepy \\
avere paura (di) & to be afraid (of) \\
avere caldo / freddo & to feel hot / cold \\
avere … anni & to be … years old \\
avere ragione / torto & to be right / wrong \\
\bottomrule
\end{longtable}}

\subsection{Side by side}\label{les:essere-avere:comparison}

\tablecaption{essere \textperiodcentered{} avere --- indicativo presente}
{\small\begin{tabular}{@{}l l l@{}}
\toprule
\textbf{Person} & \textbf{\textit{essere}} & \textbf{\textit{avere}} \\
\midrule
\textit{io} & \irr{sono} & \irr{ho} \\
\textit{tu} & \irr{sei} & \irr{hai} \\
\textit{lui/lei} & \irr{è} & \irr{ha} \\
\textit{noi} & \irr{siamo} & \irr{abbiamo} \\
\textit{voi} & \irr{siete} & avete \\
\textit{loro} & \irr{sono} & \irr{hanno} \\
\bottomrule
\end{tabular}}\par\medskip

\subsection*{Worked examples}
\begin{examples}
\item\relax \textit{Sono stanco.}\newline I am tired.
\item\relax \textit{I bambini sono in giardino.}\newline The children are in the garden.
\item\relax \textit{Ho sette anni.}\newline I am seven years old.
\item\relax \textit{Il topo ha paura del gatto.}\newline The mouse is afraid of the cat.
\item\relax \textit{Hai fame? — No, ho sete.}\newline Are you hungry? — No, I'm thirsty.
\item\relax \textit{Siamo di Milano.}\newline We are from Milan.
\end{examples}

\subsection*{Exercises}\label{les:essere-avere:exercises}
{\color{muted}\small Answers: page~\pageref{ans:essere-avere}.}

\begin{exercises}
\item\label{ex:essere-avere-1}\relax Conjugate \textit{essere} in the present.\par\smallskip{\begin{tabular}{@{}l l@{}}
io & \blankcell \\
tu & \blankcell \\
lui/lei & \blankcell \\
noi & \blankcell \\
voi & \blankcell \\
loro & \blankcell \\
\end{tabular}}
\item\label{ex:essere-avere-2}\relax Conjugate \textit{avere} in the present.\par\smallskip{\begin{tabular}{@{}l l@{}}
io & \blankcell \\
tu & \blankcell \\
lui/lei & \blankcell \\
noi & \blankcell \\
voi & \blankcell \\
loro & \blankcell \\
\end{tabular}}
\item\label{ex:essere-avere-3}\relax Luca \blank{} dieci anni.\par\nopagebreak\begin{options}
\item\relax è
\item\relax ha
\item\relax hanno
\item\relax a
\end{options}
\item\label{ex:essere-avere-4}\relax Il topo \blank{} paura del gatto.\par\nopagebreak\begin{options}
\item\relax è
\item\relax ha
\item\relax sta
\item\relax fa
\end{options}
\item\label{ex:essere-avere-5}\relax Complete with essere or avere.\par\smallskip\begingroup\setlength{\parindent}{0pt}Noi \clozeblank{1} fame. Voi \clozeblank{2} stanchi? Loro \clozeblank{3} una casa al mare.\par\endgroup
\item\label{ex:essere-avere-6}\relax Translate: Are you hungry? (tu)\par\answerline
\item\label{ex:essere-avere-7}\relax Translate: \textit{Siamo di Milano.}\par\answerline
\end{exercises}

\begin{summarybox}
\begin{itemize}
\item\relax essere: sono, sei, è, siamo, siete, sono.
\item\relax avere: ho, hai, ha, abbiamo, avete, hanno (silent h).
\item\relax avere fame / sete / sonno / paura / … anni where English uses 'to be'.
\end{itemize}
\end{summarybox}

\relatedwords{\textit{essere} to be \textperiodcentered{} \textit{avere} to have \textperiodcentered{} \textit{la paura} fear \textperiodcentered{} \textit{il sonno} sleep; sleepiness \textperiodcentered{} \textit{l'anno} year \textperiodcentered{} \textit{stanco} tired}
\applyit{Read \textit{La famiglia di Marco} (A0–A1), \S\ref{story:la-famiglia-di-marco}. Read \textit{Il gatto e la luna} (A0–A1), \S\ref{story:il-gatto-e-la-luna}. Study the rhyme \textit{Conta delle stelle}, \S\ref{rhyme:conta-delle-stelle}.}

\section{andare, fare, stare, dire}\label{les:andare-fare-stare-dire}
\lessonmeta{A1}{Four short, very common irregular verbs and the expressions built on them.}
\begin{objectives}
\item\relax Conjugate andare, fare, stare and dire in the present.
\item\relax Use andare a / andare in correctly with places.
\item\relax Use common expressions with fare and stare.
\end{objectives}

\subsection{andare — to go}\label{les:andare-fare-stare-dire:andare}

\par\noindent\begin{minipage}[t]{0.48\linewidth}\raggedright
\paradigmbox{andare --- to go \textperiodcentered{} indicativo presente}{%
\begin{tabular}{@{}l l l@{}}
\textit{io} & \irr{vado} & {\color{muted}I go} \\
\textit{tu} & \irr{vai} & {\color{muted}you go} \\
\textit{lui/lei} & \irr{va} & {\color{muted}he/she goes} \\
\textit{noi} & andiamo & {\color{muted}we go} \\
\textit{voi} & andate & {\color{muted}you (pl.) go} \\
\textit{loro} & \irr{vanno} & {\color{muted}they go} \\
\end{tabular}}{\irr{Highlighted} forms differ from the regular pattern of the verb's family.}
\end{minipage}\par\medskip

\begin{notebox}{rulecol}\notelabel{rulecol}{Rule}{Two stems}
\textit{vad-/va-} in the stressed forms (vado, vai, va, vanno); regular \textit{and-} in \textit{andiamo, andate}.
\end{notebox}

\begin{notebox}{tendcol}\notelabel{tendcol}{Tendency}{andare a or in?}
\textit{a} with towns and cities and with an infinitive (\textit{vado a Roma, vado a dormire}), and in fixed phrases (\textit{a scuola, a casa, a letto, al mare}). \textit{in} with countries, regions and many places (\textit{in Italia, in giardino, in cucina, in montagna, in città}). These are tendencies; learn the fixed phrases.
\end{notebox}

\subsection{fare — to do, to make; stare — to stay, to be}\label{les:andare-fare-stare-dire:fare-stare}

\par\noindent\begin{minipage}[t]{0.48\linewidth}\raggedright
\paradigmbox{fare --- to do; to make \textperiodcentered{} indicativo presente}{%
\begin{tabular}{@{}l l l@{}}
\textit{io} & \irr{faccio} & {\color{muted}I do} \\
\textit{tu} & \irr{fai} & {\color{muted}you do} \\
\textit{lui/lei} & fa & {\color{muted}he/she does} \\
\textit{noi} & \irr{facciamo} & {\color{muted}we do} \\
\textit{voi} & fate & {\color{muted}you (pl.) do} \\
\textit{loro} & \irr{fanno} & {\color{muted}they do} \\
\end{tabular}}{\irr{Highlighted} forms differ from the regular pattern of the verb's family.}
\end{minipage}\hfill\begin{minipage}[t]{0.48\linewidth}\raggedright
\paradigmbox{stare --- to stay; to be (of health) \textperiodcentered{} indicativo presente}{%
\begin{tabular}{@{}l l l@{}}
\textit{io} & sto & {\color{muted}I stay} \\
\textit{tu} & \irr{stai} & {\color{muted}you stay} \\
\textit{lui/lei} & sta & {\color{muted}he/she stays} \\
\textit{noi} & stiamo & {\color{muted}we stay} \\
\textit{voi} & state & {\color{muted}you (pl.) stay} \\
\textit{loro} & \irr{stanno} & {\color{muted}they stay} \\
\end{tabular}}{\irr{Highlighted} forms differ from the regular pattern of the verb's family.}
\end{minipage}\par\medskip

{\small\setlength{\LTpre}{2pt}\setlength{\LTpost}{6pt}
\tablecaption{Expressions}
\begin{xltabular}{\linewidth}{lL}
\toprule
\textbf{Italian} & \textbf{English} \\
\midrule
fare colazione & to have breakfast \\
fare una passeggiata & to go for a walk \\
Fa caldo / freddo. & It's hot / cold (weather). \\
Che tempo fa? & What's the weather like? \\
Come stai? Sto bene. & How are you? I'm well. \\
stare a casa & to stay at home \\
\bottomrule
\end{xltabular}}

\subsection{dire — to say, to tell}\label{les:andare-fare-stare-dire:dire}

\par\noindent\begin{minipage}[t]{0.48\linewidth}\raggedright
\paradigmbox{dire --- to say; to tell \textperiodcentered{} indicativo presente}{%
\begin{tabular}{@{}l l l@{}}
\textit{io} & \irr{dico} & {\color{muted}I say} \\
\textit{tu} & dici & {\color{muted}you say} \\
\textit{lui/lei} & \irr{dice} & {\color{muted}he/she says} \\
\textit{noi} & \irr{diciamo} & {\color{muted}we say} \\
\textit{voi} & dite & {\color{muted}you (pl.) say} \\
\textit{loro} & \irr{dicono} & {\color{muted}they say} \\
\end{tabular}}{\irr{Highlighted} forms differ from the regular pattern of the verb's family.}
\end{minipage}\par\medskip

\begin{notebox}{tendcol}\notelabel{tendcol}{Tendency}{dire versus parlare}
\textit{dire} takes what is said as its object (\textit{dire la verità, dire ciao}); \textit{parlare} means to speak or talk (\textit{parlare con qualcuno, parlare italiano}).
\end{notebox}

\subsection*{Worked examples}
\begin{examples}
\item\relax \textit{Vado a scuola a piedi.}\newline I go to school on foot.
\item\relax \textit{Andiamo in giardino!}\newline Let's go into the garden!
\item\relax \textit{Facciamo colazione insieme.}\newline We have breakfast together.
\item\relax \textit{Oggi fa freddo.}\newline It's cold today.
\item\relax \textit{Come stai? — Sto bene, grazie.}\newline How are you? — I'm fine, thanks.
\item\relax \textit{Il nonno dice che domani piove.}\newline Grandfather says it's going to rain tomorrow.
\end{examples}

\subsection*{Exercises}\label{les:andare-fare-stare-dire:exercises}
{\color{muted}\small Answers: page~\pageref{ans:andare-fare-stare-dire}.}

\begin{exercises}
\item\label{ex:andare-fare-stare-dire-1}\relax Conjugate \textit{andare}.\par\smallskip{\begin{tabular}{@{}l l@{}}
io & \blankcell \\
tu & \blankcell \\
lui/lei & \blankcell \\
noi & \blankcell \\
voi & \blankcell \\
loro & \blankcell \\
\end{tabular}}
\item\label{ex:andare-fare-stare-dire-2}\relax Conjugate \textit{fare}.\par\smallskip{\begin{tabular}{@{}l l@{}}
io & \blankcell \\
tu & \blankcell \\
lui/lei & \blankcell \\
noi & \blankcell \\
voi & \blankcell \\
loro & \blankcell \\
\end{tabular}}
\item\label{ex:andare-fare-stare-dire-3}\relax Vado \blank{} Roma.\par\nopagebreak\begin{options}
\item\relax in
\item\relax a
\item\relax da
\item\relax di
\end{options}
\item\label{ex:andare-fare-stare-dire-4}\relax Andiamo \blank{} montagna.\par\nopagebreak\begin{options}
\item\relax su
\item\relax in
\item\relax a
\item\relax al
\end{options}
\item\label{ex:andare-fare-stare-dire-5}\relax Complete.\par\smallskip\begingroup\setlength{\parindent}{0pt}Come \clozeblank{1} (stare, tu)? — \clozeblank{2} (stare, io) bene. Oggi \clozeblank{3} (fare) caldo.\par\endgroup
\item\label{ex:andare-fare-stare-dire-6}\relax dire, loro: \blank{}
\item\label{ex:andare-fare-stare-dire-7}\relax Translate: We have breakfast together.\par\answerline
\end{exercises}

\begin{summarybox}
\begin{itemize}
\item\relax andare: vado, vai, va, andiamo, andate, vanno.
\item\relax fare: faccio, fai, fa, facciamo, fate, fanno.
\item\relax stare: sto, stai, sta, stiamo, state, stanno.
\item\relax dire: dico, dici, dice, diciamo, dite, dicono.
\item\relax andare a + city / infinitive; in + country, region and many places.
\end{itemize}
\end{summarybox}

\relatedwords{\textit{andare} to go \textperiodcentered{} \textit{fare} to do; to make \textperiodcentered{} \textit{stare} to stay; to be (of health) \textperiodcentered{} \textit{dire} to say; to tell \textperiodcentered{} \textit{la colazione} breakfast \textperiodcentered{} \textit{la scuola} school \textperiodcentered{} \textit{la montagna} mountain \textperiodcentered{} \textit{il tempo} weather; time}
\applyit{Read \textit{La colazione} (A0–A1), \S\ref{story:la-colazione}. Read \textit{Il cane al mare} (A0–A1), \S\ref{story:il-cane-al-mare}.}

\section{venire, uscire, bere, dare and piacere}\label{les:other-irregulars}
\lessonmeta{A1}{More high-frequency irregular verbs and the pattern behind them; saying what you like.}
\begin{objectives}
\item\relax Conjugate venire, uscire, bere, dare, tenere and rimanere.
\item\relax Notice the recurring 'g in io and loro' pattern.
\item\relax Say what you like with mi piace / mi piacciono.
\end{objectives}

\subsection{A recurring pattern: g in io and loro}\label{les:other-irregulars:pattern}

Several irregular verbs add a \textit{g} in the \textit{io} and \textit{loro} forms: \textit{venire → vengo, vengono}; \textit{tenere → tengo, tengono}; \textit{rimanere → rimango, rimangono}; \textit{salire → salgo, salgono}. Venire and tenere also change the stressed \textit{e} to \textit{ie} in \textit{tu} and \textit{lui/lei}.\par

\par\noindent\begin{minipage}[t]{0.48\linewidth}\raggedright
\paradigmbox{venire --- to come \textperiodcentered{} indicativo presente}{%
\begin{tabular}{@{}l l l@{}}
\textit{io} & \irr{vengo} & {\color{muted}I come} \\
\textit{tu} & \irr{vieni} & {\color{muted}you come} \\
\textit{lui/lei} & \irr{viene} & {\color{muted}he/she comes} \\
\textit{noi} & veniamo & {\color{muted}we come} \\
\textit{voi} & venite & {\color{muted}you (pl.) come} \\
\textit{loro} & \irr{vengono} & {\color{muted}they come} \\
\end{tabular}}{\irr{Highlighted} forms differ from the regular pattern of the verb's family.}
\end{minipage}\hfill\begin{minipage}[t]{0.48\linewidth}\raggedright
\paradigmbox{rimanere --- to remain; to stay \textperiodcentered{} indicativo presente}{%
\begin{tabular}{@{}l l l@{}}
\textit{io} & \irr{rimango} & {\color{muted}I remain} \\
\textit{tu} & rimani & {\color{muted}you remain} \\
\textit{lui/lei} & rimane & {\color{muted}he/she remains} \\
\textit{noi} & rimaniamo & {\color{muted}we remain} \\
\textit{voi} & rimanete & {\color{muted}you (pl.) remain} \\
\textit{loro} & \irr{rimangono} & {\color{muted}they remain} \\
\end{tabular}}{\irr{Highlighted} forms differ from the regular pattern of the verb's family.}
\end{minipage}\par\medskip

\subsection{uscire, bere, dare}\label{les:other-irregulars:uscire-bere-dare}

\par\noindent\begin{minipage}[t]{0.48\linewidth}\raggedright
\paradigmbox{uscire --- to go out; to leave \textperiodcentered{} indicativo presente}{%
\begin{tabular}{@{}l l l@{}}
\textit{io} & \irr{esco} & {\color{muted}I go out} \\
\textit{tu} & \irr{esci} & {\color{muted}you go out} \\
\textit{lui/lei} & \irr{esce} & {\color{muted}he/she goes out} \\
\textit{noi} & usciamo & {\color{muted}we go out} \\
\textit{voi} & uscite & {\color{muted}you (pl.) go out} \\
\textit{loro} & \irr{escono} & {\color{muted}they go out} \\
\end{tabular}}{\irr{Highlighted} forms differ from the regular pattern of the verb's family.}
\end{minipage}\hfill\begin{minipage}[t]{0.48\linewidth}\raggedright
\paradigmbox{bere --- to drink \textperiodcentered{} indicativo presente}{%
\begin{tabular}{@{}l l l@{}}
\textit{io} & \irr{bevo} & {\color{muted}I drink} \\
\textit{tu} & \irr{bevi} & {\color{muted}you drink} \\
\textit{lui/lei} & \irr{beve} & {\color{muted}he/she drinks} \\
\textit{noi} & \irr{beviamo} & {\color{muted}we drink} \\
\textit{voi} & \irr{bevete} & {\color{muted}you (pl.) drink} \\
\textit{loro} & \irr{bevono} & {\color{muted}they drink} \\
\end{tabular}}{\irr{Highlighted} forms differ from the regular pattern of the verb's family.}
\end{minipage}\par\medskip\noindent\begin{minipage}[t]{0.48\linewidth}\raggedright
\paradigmbox{dare --- to give \textperiodcentered{} indicativo presente}{%
\begin{tabular}{@{}l l l@{}}
\textit{io} & do & {\color{muted}I give} \\
\textit{tu} & \irr{dai} & {\color{muted}you give} \\
\textit{lui/lei} & \irr{dà} & {\color{muted}he/she gives} \\
\textit{noi} & diamo & {\color{muted}we give} \\
\textit{voi} & date & {\color{muted}you (pl.) give} \\
\textit{loro} & \irr{danno} & {\color{muted}they give} \\
\end{tabular}}{\irr{Highlighted} forms differ from the regular pattern of the verb's family.}
\end{minipage}\par\medskip

\begin{notebox}{rulecol}\notelabel{rulecol}{Rule}{dà with an accent}
\textit{dà} (gives) is written with an accent to distinguish it from the preposition \textit{da} (from).
\end{notebox}

\subsection{piacere — to like}\label{les:other-irregulars:piacere}

\begin{notebox}{rulecol}\notelabel{rulecol}{Rule}{The thing liked is the subject}
\textit{piacere} means 'to please'. \textit{Mi piace il gelato} is literally 'Ice cream pleases me'. So the verb agrees with the thing liked: \textit{piace} + singular or infinitive, \textit{piacciono} + plural.
\end{notebox}

{\small\setlength{\LTpre}{2pt}\setlength{\LTpost}{6pt}
\tablecaption{mi piace / mi piacciono}
\begin{longtable}{ll}
\toprule
\textbf{Italian} & \textbf{English} \\
\midrule
Mi piace il gelato. & I like ice cream. \\
Mi piacciono le fragole. & I like strawberries. \\
Ti piace cantare? & Do you like singing? \\
Al gatto piace il latte. & The cat likes milk. \\
Non mi piace la pioggia. & I don't like the rain. \\
\bottomrule
\end{longtable}}

\textit{mi, ti, gli, le, ci, vi} are indirect object pronouns ('to me, to you…'); with \textit{piacere} they show the person who likes something: \textit{mi piace} = 'it pleases me', \textit{ti piace} = 'it pleases you', \textit{gli piace} = 'it pleases him', \textit{le piace} = 'it pleases her'. The thing liked is the grammatical subject, so the verb agrees with it: \textit{mi piace il gelato}, \textit{mi piacciono le mele}. With a noun the person is introduced by \textit{a}: \textit{A Marco piace il mare}. The full set of indirect pronouns is in \lessonref{indirect-object-pronouns}{indirect object pronouns}.\par

\subsection*{Worked examples}
\begin{examples}
\item\relax \textit{Vieni al mare con noi?}\newline Are you coming to the seaside with us?
\item\relax \textit{Esco con gli amici.}\newline I'm going out with my friends.
\item\relax \textit{Il gatto beve il latte.}\newline The cat drinks the milk.
\item\relax \textit{La mucca dà il latte.}\newline The cow gives milk.
\item\relax \textit{Rimaniamo a casa oggi.}\newline We're staying at home today.
\item\relax \textit{Mi piacciono le storie.}\newline I like stories.
\end{examples}

\subsection*{Exercises}\label{les:other-irregulars:exercises}
{\color{muted}\small Answers: page~\pageref{ans:other-irregulars}.}

\begin{exercises}
\item\label{ex:other-irregulars-1}\relax Conjugate \textit{venire}.\par\smallskip{\begin{tabular}{@{}l l@{}}
io & \blankcell \\
tu & \blankcell \\
lui/lei & \blankcell \\
noi & \blankcell \\
voi & \blankcell \\
loro & \blankcell \\
\end{tabular}}
\item\label{ex:other-irregulars-2}\relax uscire, io: \blank{}
\item\label{ex:other-irregulars-3}\relax bere, noi: \blank{}
\item\label{ex:other-irregulars-4}\relax Mi \blank{} le fragole.\par\nopagebreak\begin{options}
\item\relax piaccio
\item\relax piacciamo
\item\relax piace
\item\relax piacciono
\end{options}
\item\label{ex:other-irregulars-5}\relax Ti \blank{} cantare?\par\nopagebreak\begin{options}
\item\relax piacciono
\item\relax piaci
\item\relax piaccio
\item\relax piace
\end{options}
\item\label{ex:other-irregulars-6}\relax Infinitive of \textit{tengono}?\par\answerline
\item\label{ex:other-irregulars-7}\relax Translate: I don't like the rain.\par\answerline
\end{exercises}

\begin{summarybox}
\begin{itemize}
\item\relax g in io and loro: vengo/vengono, tengo/tengono, rimango/rimangono, salgo/salgono.
\item\relax uscire: esco, esci, esce, usciamo, uscite, escono. bere: bev- + regular endings. dare: do, dai, dà…
\item\relax mi piace + singular / infinitive; mi piacciono + plural.
\end{itemize}
\end{summarybox}

\relatedwords{\textit{venire} to come \textperiodcentered{} \textit{uscire} to go out; to leave \textperiodcentered{} \textit{bere} to drink \textperiodcentered{} \textit{dare} to give \textperiodcentered{} \textit{tenere} to hold; to keep \textperiodcentered{} \textit{rimanere} to remain; to stay \textperiodcentered{} \textit{piacere} to please; to be liked \textperiodcentered{} \textit{salire} to go up; to climb; to get on}
\applyit{Read \textit{Il cane al mare} (A0–A1), \S\ref{story:il-cane-al-mare}. Read \textit{La colazione} (A0–A1), \S\ref{story:la-colazione}. Study the rhyme \textit{Il valzer degli animali}, \S\ref{rhyme:il-valzer-degli-animali}.}

\chapter{Building sentences}\label{mod:m06}
\moduleintro{A1}{Word order, negation, questions, question words and conjunctions.}{Word order, negation and yes/no questions (\S\ref{les:word-order-negation-questions}); Question words and conjunctions (\S\ref{les:question-words-conjunctions})}

\section{Word order, negation and yes/no questions}\label{les:word-order-negation-questions}
\lessonmeta{A1}{Subject–verb–object, non before the verb, double negatives, and questions by intonation.}
\begin{objectives}
\item\relax Build basic subject–verb–object sentences.
\item\relax Negate with non and use non … mai / niente / nessuno.
\item\relax Ask yes/no questions without changing word order.
\end{objectives}

\subsection{Basic word order}\label{les:word-order-negation-questions:order}

The neutral order is subject – verb – object, as in English: \textit{La bambina mangia una mela}. Because verb endings show the subject, the subject is often omitted (\textit{Mangia una mela}), and it can move after the verb for emphasis or when it is new information: \textit{Arriva il treno} (Here comes the train).\par

\begin{notebox}{rulecol}\notelabel{rulecol}{Rule}{Adjectives and noun phrases}
Inside a noun phrase the order is article – (some adjectives) – noun – (most adjectives): \textit{la piccola casa bianca}.
\end{notebox}

\subsection{Negation}\label{les:word-order-negation-questions:negation}

\begin{notebox}{rulecol}\notelabel{rulecol}{Rule}{non goes directly before the verb}
Place \textit{non} immediately before the conjugated verb — and before any object or reflexive pronoun attached to it: \textit{Non mangio. Non lo so. Non mi sveglio presto.}
\end{notebox}

\begin{notebox}{rulecol}\notelabel{rulecol}{Rule}{Double negatives are required}
With \textit{mai} (never), \textit{niente} (nothing), \textit{nessuno} (nobody) and \textit{più} (no longer), Italian keeps \textit{non} before the verb: \textit{Non piango mai. Non vedo niente. Non c'è nessuno. Non piove più.} If the negative word comes first, \textit{non} is dropped: \textit{Nessuno risponde.}
\end{notebox}

\subsection{Yes/no questions}\label{les:word-order-negation-questions:questions}

\begin{notebox}{rulecol}\notelabel{rulecol}{Rule}{Intonation, not word order}
A statement becomes a question by rising intonation (and a question mark in writing). There is no equivalent of English 'do': \textit{Parli italiano.} (You speak Italian.) → \textit{Parli italiano?} (Do you speak Italian?)
\end{notebox}

An explicit subject often moves to the end: \textit{È arrivato Marco?} (Has Marco arrived?). Answers: \textit{sì} (yes) — with an accent — and \textit{no}.\par

\subsection*{Worked examples}
\begin{examples}
\item\relax \textit{La bambina mangia una mela.}\newline The little girl eats an apple.
\item\relax \textit{La bambina non mangia la mela.}\newline The little girl doesn't eat the apple.
\item\relax \textit{Non piango mai.}\newline I never cry.
\item\relax \textit{Nel bosco non c'è nessuno.}\newline There is nobody in the wood.
\item\relax \textit{Parli italiano?}\newline Do you speak Italian?
\item\relax \textit{Non lo so.}\newline I don't know (it).\newline {\color{muted}\small non before the pronoun lo.}
\end{examples}

\subsection*{Exercises}\label{les:word-order-negation-questions:exercises}
{\color{muted}\small Answers: page~\pageref{ans:word-order-negation-questions}.}

\begin{exercises}
\item\label{ex:word-order-negation-questions-1}\relax Make the sentence negative and put the words in order: 'The cat does not sleep.'\par\smallskip \token{dorme} \token{gatto} \token{Il} \token{non}\par\answerline
\item\label{ex:word-order-negation-questions-2}\relax Order the words: 'I don't wake up early.'\par\smallskip \token{mi} \token{Non} \token{presto} \token{sveglio}\par\answerline
\item\label{ex:word-order-negation-questions-3}\relax Translate: I never cry.\par\answerline
\item\label{ex:word-order-negation-questions-4}\relax Translate: Do you speak Italian? (tu)\par\answerline
\item\label{ex:word-order-negation-questions-5}\relax Which sentence is correct?\par\nopagebreak\begin{options}
\item\relax Non vedo non niente.
\item\relax Niente non vedo.
\item\relax Non vedo niente.
\item\relax Vedo niente.
\end{options}
\item\label{ex:word-order-negation-questions-6}\relax In \textit{Arriva il treno}, what is \textit{il treno}?\par\nopagebreak\begin{options}
\item\relax a prepositional phrase
\item\relax the direct object
\item\relax the subject
\item\relax an adverb
\end{options}
\end{exercises}

\begin{summarybox}
\begin{itemize}
\item\relax Neutral order: subject – verb – object; subjects are often omitted or placed after the verb.
\item\relax non goes immediately before the verb and its pronouns.
\item\relax non … mai / niente / nessuno / più: double negation is normal.
\item\relax Yes/no questions keep statement order; intonation does the work.
\end{itemize}
\end{summarybox}

\relatedwords{\textit{non} not \textperiodcentered{} \textit{mai} never; ever \textperiodcentered{} \textit{sì} yes \textperiodcentered{} \textit{no} no \textperiodcentered{} \textit{parlare} to speak; to talk \textperiodcentered{} \textit{piangere} to cry}
\applyit{Read \textit{La tartaruga lenta} (A1), \S\ref{story:la-tartaruga-lenta}. Read \textit{Nel bosco} (A1), \S\ref{story:nel-bosco}.}

\section{Question words and conjunctions}\label{les:question-words-conjunctions}
\lessonmeta{A1}{chi, che cosa, dove, quando, come, perché, quanto, quale — and joining clauses with e, ma, o, perché, quando, se, mentre.}
\begin{objectives}
\item\relax Ask open questions with the main question words.
\item\relax Join two clauses with common conjunctions.
\item\relax Use perché both as 'why' and 'because'.
\end{objectives}

\subsection{Question words}\label{les:question-words-conjunctions:question-words}

{\small\setlength{\LTpre}{2pt}\setlength{\LTpost}{6pt}
\tablecaption{Interrogatives}
\begin{xltabular}{\linewidth}{llL}
\toprule
\textbf{Italian} & \textbf{English} & \textbf{Example} \\
\midrule
chi? & who? & Chi è? — Who is it? \\
che cosa? / cosa? / che? & what? & Che cosa mangi? \\
dove? & where? & Dove abiti? Dov'è il gatto? \\
quando? & when? & Quando parti? \\
come? & how? & Come stai? Come ti chiami? \\
perché? & why? & Perché piangi? \\
quanto/a/i/e? & how much / how many? & Quanti anni hai? \\
quale/i? & which? & Quale libro preferisci? Qual è…? \\
\bottomrule
\end{xltabular}}

\begin{notebox}{rulecol}\notelabel{rulecol}{Rule}{The subject usually follows the verb}
In open questions an explicit subject normally comes after the verb: \textit{Dove abita Marco?} \textit{Che cosa fa la mamma?}
\end{notebox}

\begin{notebox}{rulecol}\notelabel{rulecol}{Rule}{quanto agrees, qual è has no apostrophe}
\textit{quanto} agrees with its noun: \textit{quanta acqua, quanti anni, quante stelle}. \textit{quale} is truncated before \textit{è} without an apostrophe: \textit{Qual è il tuo nome?}
\end{notebox}

\subsection{Conjunctions}\label{les:question-words-conjunctions:conjunctions}

{\small\setlength{\LTpre}{2pt}\setlength{\LTpost}{6pt}
\tablecaption{Joining clauses}
\begin{xltabular}{\linewidth}{llL}
\toprule
\textbf{Italian} & \textbf{English} & \textbf{Example} \\
\midrule
e & and & Mangio e poi dormo. \\
ma & but & È piccolo ma forte. \\
o & or & Vuoi il tè o il caffè? \\
perché & because & Resto a casa perché piove. \\
quando & when & Quando piove, leggo. \\
se & if & Se piove, prendo l'ombrello. \\
mentre & while & Canto mentre cucino. \\
poi & then & Prima mangio, poi gioco. \\
\bottomrule
\end{xltabular}}

\begin{notebox}{rulecol}\notelabel{rulecol}{Rule}{perché = why and because}
The same word asks and answers: \textit{Perché piangi? — Perché sono stanco.} Always written with an acute accent: \textit{perché}.
\end{notebox}

\subsection*{Worked examples}
\begin{examples}
\item\relax \textit{Dove abita la volpe?}\newline Where does the fox live?
\item\relax \textit{Quanti anni hai?}\newline How old are you?
\item\relax \textit{Come ti chiami?}\newline What's your name?
\item\relax \textit{Resto a casa perché piove.}\newline I'm staying at home because it is raining.
\item\relax \textit{Se piove, prendo l'ombrello.}\newline If it rains, I'll take the umbrella.
\item\relax \textit{L'orso è grande ma gentile.}\newline The bear is big but gentle.
\end{examples}

\subsection*{Exercises}\label{les:question-words-conjunctions:exercises}
{\color{muted}\small Answers: page~\pageref{ans:question-words-conjunctions}.}

\begin{exercises}
\item\label{ex:question-words-conjunctions-1}\relax Match question and answer.\par\smallskip\begin{tabularx}{\linewidth}{@{}L L@{}}
1.\ Dove abiti? & (a)\ A Roma. \\
2.\ Quanti anni hai? & (b)\ Perché sono stanco. \\
3.\ Chi è? & (c)\ È la nonna. \\
4.\ Quando parti? & (d)\ Sette. \\
5.\ Perché piangi? & (e)\ Domani. \\
\end{tabularx}
\item\label{ex:question-words-conjunctions-2}\relax \blank{} stelle ci sono in cielo?\par\nopagebreak\begin{options}
\item\relax Quanto
\item\relax Quale
\item\relax Quanti
\item\relax Quante
\end{options}
\item\label{ex:question-words-conjunctions-3}\relax \blank{} il tuo nome?\par\nopagebreak\begin{options}
\item\relax Quale è'
\item\relax Qual è
\item\relax Quali è
\item\relax Qual'è
\end{options}
\item\label{ex:question-words-conjunctions-4}\relax Complete with ma, perché, se or mentre.\par\smallskip\begingroup\setlength{\parindent}{0pt}Resto a casa \clozeblank{1} piove. \clozeblank{2} fa bel tempo, andiamo al mare. Il cane è vecchio \clozeblank{3} corre veloce. Il gatto dorme \clozeblank{4} io leggo.\par\endgroup
\item\label{ex:question-words-conjunctions-5}\relax Translate: Where does the fox live?\par\answerline
\item\label{ex:question-words-conjunctions-6}\relax Translate: \textit{Come ti chiami?}\par\answerline
\end{exercises}

\begin{summarybox}
\begin{itemize}
\item\relax chi, che cosa, dove, quando, come, perché, quanto, quale.
\item\relax Explicit subjects usually follow the verb in questions.
\item\relax e, ma, o, perché, quando, se, mentre, poi join clauses.
\item\relax perché = why / because; qual è without apostrophe.
\end{itemize}
\end{summarybox}

\relatedwords{\textit{chi} who; whom \textperiodcentered{} \textit{che} that; which; who; what \textperiodcentered{} \textit{cosa} what (question) \textperiodcentered{} \textit{dove} where \textperiodcentered{} \textit{come} how; like; as \textperiodcentered{} \textit{quando} when \textperiodcentered{} \textit{perché} because; why \textperiodcentered{} \textit{quanto} how much; how many \textperiodcentered{} \textit{quale} which; what \textperiodcentered{} \textit{e} and \textperiodcentered{} \textit{ma} but \textperiodcentered{} \textit{o} or \textperiodcentered{} \textit{se} if; whether \textperiodcentered{} \textit{mentre} while \textperiodcentered{} \textit{poi} then; afterwards}
\applyit{Read \textit{Una giornata di pioggia} (A1), \S\ref{story:una-giornata-di-pioggia}. Study the rhyme \textit{La pioggia sul tetto}, \S\ref{rhyme:la-pioggia-sul-tetto}.}

\chapter{Prepositions}\label{mod:m07}
\moduleintro{A1}{The simple prepositions and their contractions with the article.}{The simple prepositions (\S\ref{les:simple-prepositions}); Articulated prepositions (\S\ref{les:articulated-prepositions})}

\section{The simple prepositions}\label{les:simple-prepositions}
\lessonmeta{A1}{a, di, da, in, con, su, per, tra/fra — for place, movement, time and possession.}
\begin{objectives}
\item\relax Know the core meanings of the nine simple prepositions.
\item\relax Express place and movement (a, in, da).
\item\relax Express time (in, a, per, tra, da) and possession (di).
\end{objectives}

\subsection{Nine prepositions}\label{les:simple-prepositions:overview}

{\small\setlength{\LTpre}{2pt}\setlength{\LTpost}{6pt}
\tablecaption{Core meanings}
\begin{xltabular}{\linewidth}{lLL}
\toprule
\textbf{Preposition} & \textbf{Core meanings} & \textbf{Examples} \\
\midrule
di & of; possession; origin with essere; material & il libro di Marco; sono di Roma \\
a & to, at; cities; time of day; + infinitive after some verbs & vado a Roma; a mezzanotte; comincio a leggere \\
da & from; at/to someone's place; since (time) & vengo da Milano; vado dal nonno; studio da tre mesi \\
in & in, into; countries, regions; means of transport; seasons & in Italia; in giardino; in treno; in estate \\
con & with & gioco con il cane \\
su & on; about & sul tavolo; un libro sugli animali \\
per & for; through; in order to; duration & per te; per due ore; studio per imparare \\
tra / fra & between, among; in (time from now) & tra i fiori; fra un'ora \\
\bottomrule
\end{xltabular}}

\begin{notebox}{tendcol}\notelabel{tendcol}{Tendency}{Prepositions do not map one-to-one}
No Italian preposition has a single English equivalent, and fixed expressions are common (\textit{a casa, a scuola, a letto, in montagna, in vacanza, al mare}). Learn prepositions in phrases rather than as single words.
\end{notebox}

\subsection{Place and movement}\label{les:simple-prepositions:place}

\begin{notebox}{rulecol}\notelabel{rulecol}{Rule}{No difference between 'in' and 'to'}
Italian uses the same preposition for being somewhere and going there: \textit{Sono a Roma / Vado a Roma}; \textit{Sono in giardino / Vado in giardino}.
\end{notebox}

\begin{notebox}{rulecol}\notelabel{rulecol}{Rule}{da + person}
\textit{da} + a person or a shop means 'at/to someone's place': \textit{Sono dalla nonna} (I'm at Grandma's), \textit{Vado dal medico} (I'm going to the doctor's).
\end{notebox}

\subsection{Time}\label{les:simple-prepositions:time}

\begin{itemize}
\item\relax \textit{in} + season or month: \textit{in estate, in inverno, in dicembre}.
\item\relax \textit{a} + hour, festival: \textit{a mezzogiorno, alle tre, a Natale}.
\item\relax \textit{per} + duration: \textit{per un'ora} (for an hour).
\item\relax \textit{tra / fra} + amount of time = in … from now: \textit{tra un'ora} (in an hour).
\item\relax \textit{da} + amount of time with the present = for/since (still continuing): \textit{Abito qui da due anni} (I have lived here for two years).
\end{itemize}

\subsection*{Worked examples}
\begin{examples}
\item\relax \textit{Il cane di Sara è nero.}\newline Sara's dog is black.\newline {\color{muted}\small di for possession.}
\item\relax \textit{Vado a scuola con mio fratello.}\newline I go to school with my brother.
\item\relax \textit{Vengo da Milano.}\newline I come from Milan.
\item\relax \textit{In estate andiamo al mare.}\newline In summer we go to the seaside.
\item\relax \textit{Arrivo fra un'ora.}\newline I'll arrive in an hour.
\item\relax \textit{Il regalo è per te.}\newline The present is for you.
\end{examples}

\subsection*{Exercises}\label{les:simple-prepositions:exercises}
{\color{muted}\small Answers: page~\pageref{ans:simple-prepositions}.}

\begin{exercises}
\item\label{ex:simple-prepositions-1}\relax Abito \blank{} Napoli.\par\nopagebreak\begin{options}
\item\relax a
\item\relax in
\item\relax su
\item\relax da
\end{options}
\item\label{ex:simple-prepositions-2}\relax Abito \blank{} Italia.\par\nopagebreak\begin{options}
\item\relax di
\item\relax in
\item\relax per
\item\relax a
\end{options}
\item\label{ex:simple-prepositions-3}\relax Complete with di, da, con or per.\par\smallskip\begingroup\setlength{\parindent}{0pt}Il libro \clozeblank{1} Marco è sul tavolo. Vengo \clozeblank{2} Torino. Gioco \clozeblank{3} il gatto. Questo regalo è \clozeblank{4} te.\par\endgroup
\item\label{ex:simple-prepositions-4}\relax 'In an hour (from now)': \blank{} un'ora\par\nopagebreak\begin{options}
\item\relax fra
\item\relax in
\item\relax da
\item\relax per
\end{options}
\item\label{ex:simple-prepositions-5}\relax Translate: Sara's dog is black.\par\answerline
\item\label{ex:simple-prepositions-6}\relax Translate: \textit{Abito qui da due anni.}\par\answerline
\end{exercises}

\begin{summarybox}
\begin{itemize}
\item\relax di (of, possession), a (to/at, cities), da (from, someone's place, since), in (in/into, countries), con (with), su (on), per (for), tra/fra (between, in … time).
\item\relax The same preposition expresses position and movement: sono a Roma / vado a Roma.
\item\relax da + present for ongoing duration: abito qui da due anni.
\end{itemize}
\end{summarybox}

\relatedwords{\textit{a} to; at; in (cities) \textperiodcentered{} \textit{di} of; from (origin with essere) \textperiodcentered{} \textit{da} from; at (someone's place); since \textperiodcentered{} \textit{in} in; into; to (countries, rooms) \textperiodcentered{} \textit{con} with \textperiodcentered{} \textit{su} on; onto; about \textperiodcentered{} \textit{per} for; through; in order to \textperiodcentered{} \textit{tra} between; among; in (time from now) \textperiodcentered{} \textit{fra} between; among; in (time from now)}
\applyit{Read \textit{Il cane al mare} (A0–A1), \S\ref{story:il-cane-al-mare}. Read \textit{Nel bosco} (A1), \S\ref{story:nel-bosco}.}

\section{Articulated prepositions}\label{les:articulated-prepositions}
\lessonmeta{A1}{di, a, da, in, su combined with the definite article: del, al, dal, nel, sul…}
\begin{objectives}
\item\relax Form all articulated prepositions from the table.
\item\relax Recognise them in reading: nel bosco, sulla spiaggia, dalla nonna.
\item\relax Use di + article for 'some'.
\end{objectives}

\subsection{The table}\label{les:articulated-prepositions:table}

When \textit{di, a, da, in, su} are followed by a definite article, they merge with it into one word. \textit{in} becomes \textit{ne-} and \textit{di} becomes \textit{de-}.\par

\tablecaption{Preposition + definite article}
{\small\begin{tabular}{@{}l l l l l l l l@{}}
\toprule
 & \textbf{\textit{il}} & \textbf{\textit{lo}} & \textbf{\textit{l'}} & \textbf{\textit{la}} & \textbf{\textit{i}} & \textbf{\textit{gli}} & \textbf{\textit{le}} \\
\midrule
\textbf{\textit{di}} & \textit{del} & \textit{dello} & \textit{dell'} & \textit{della} & \textit{dei} & \textit{degli} & \textit{delle} \\
\textbf{\textit{a}} & \textit{al} & \textit{allo} & \textit{all'} & \textit{alla} & \textit{ai} & \textit{agli} & \textit{alle} \\
\textbf{\textit{da}} & \textit{dal} & \textit{dallo} & \textit{dall'} & \textit{dalla} & \textit{dai} & \textit{dagli} & \textit{dalle} \\
\textbf{\textit{in}} & \textit{nel} & \textit{nello} & \textit{nell'} & \textit{nella} & \textit{nei} & \textit{negli} & \textit{nelle} \\
\textbf{\textit{su}} & \textit{sul} & \textit{sullo} & \textit{sull'} & \textit{sulla} & \textit{sui} & \textit{sugli} & \textit{sulle} \\
\bottomrule
\end{tabular}}\par\medskip

\begin{notebox}{rulecol}\notelabel{rulecol}{Rule}{Building the form}
Take the article and attach it: \textit{a + il = al}; with an article beginning with l, double the l: \textit{a + la = alla}, \textit{a + lo = allo}, \textit{a + l' = all'}.
\end{notebox}

\begin{notebox}{rulecol}\notelabel{rulecol}{Rule}{con, per, tra, fra do not contract}
Modern Italian writes \textit{con il, con la, per il, tra i}. (The forms \textit{col} and \textit{coi} exist but are optional and informal.)
\end{notebox}

\subsection{Using them}\label{les:articulated-prepositions:use}

\begin{itemize}
\item\relax Place: \textit{nel bosco, sulla spiaggia, nella tana, sull'albero}.
\item\relax Movement: \textit{al mare, alla finestra, dal nonno}.
\item\relax Possession with a noun that takes the article: \textit{la casa della nonna, il nido dell'uccello}.
\item\relax Time: \textit{alle tre} (at three o'clock), \textit{dalle otto} (from eight o'clock).
\item\relax 'Some': \textit{del pane, delle mele, degli amici}.
\end{itemize}

\begin{notebox}{rulecol}\notelabel{rulecol}{Rule}{No article, no contraction}
Where the noun normally has no article, the simple preposition stays: \textit{a Roma, di Marco, in Italia, a casa, a scuola}. Compare \textit{in giardino} (fixed phrase) and \textit{nel giardino della nonna} (a specific garden).
\end{notebox}

\subsection*{Worked examples}
\begin{examples}
\item\relax \textit{La volpe torna nella tana.}\newline The fox goes back into its den.\newline {\color{muted}\small in + la = nella.}
\item\relax \textit{Lo scoiattolo è sull'albero.}\newline The squirrel is in the tree.\newline {\color{muted}\small su + l' = sull'.}
\item\relax \textit{Andiamo dal nonno.}\newline We're going to Grandpa's.\newline {\color{muted}\small da + il = dal.}
\item\relax \textit{Il nido dell'uccello è sul ramo.}\newline The bird's nest is on the branch.
\item\relax \textit{Ai bambini piacciono i giochi.}\newline Children like games.\newline {\color{muted}\small a + i = ai.}
\item\relax \textit{Compro del pane e delle mele.}\newline I'm buying some bread and some apples.
\end{examples}

\subsection*{Exercises}\label{les:articulated-prepositions:exercises}
{\color{muted}\small Answers: page~\pageref{ans:articulated-prepositions}.}

\begin{exercises}
\item\label{ex:articulated-prepositions-1}\relax in + il bosco → \blank{} bosco
\item\label{ex:articulated-prepositions-2}\relax su + la spiaggia → \blank{} spiaggia
\item\label{ex:articulated-prepositions-3}\relax Write the contraction of \textit{di} with each article.\par\smallskip{\begin{tabular}{@{}l l@{}}
di + il & \blankcell \\
di + lo & \blankcell \\
di + l' & \blankcell \\
di + la & \blankcell \\
di + i & \blankcell \\
di + gli & \blankcell \\
di + le & \blankcell \\
\end{tabular}}
\item\label{ex:articulated-prepositions-4}\relax Gli uccelli sono \blank{} alberi.\par\nopagebreak\begin{options}
\item\relax sui
\item\relax sulle
\item\relax sugli
\item\relax su gli
\end{options}
\item\label{ex:articulated-prepositions-5}\relax Vado \blank{} Roma.\par\nopagebreak\begin{options}
\item\relax nella
\item\relax alla
\item\relax al
\item\relax a
\end{options}
\item\label{ex:articulated-prepositions-6}\relax Complete with articulated prepositions.\par\smallskip\begingroup\setlength{\parindent}{0pt}La casa \clozeblank{1} (di + la) nonna è vicino \clozeblank{2} (a + il) fiume. Andiamo \clozeblank{3} (da + i) cugini.\par\endgroup
\item\label{ex:articulated-prepositions-7}\relax Translate: The fox goes back into its den.\par\answerline
\end{exercises}

\begin{summarybox}
\begin{itemize}
\item\relax di, a, da, in, su + article merge: del, al, dal, nel, sul; della, alla, dalla, nella, sulla…
\item\relax Double l before articles beginning with l (alla, dello, nell').
\item\relax No article (a Roma, di Marco, a casa) → no contraction.
\item\relax con, per, tra, fra are written separately.
\end{itemize}
\end{summarybox}

\relatedwords{\textit{il bosco} wood; forest \textperiodcentered{} \textit{la tana} den; burrow \textperiodcentered{} \textit{la spiaggia} beach \textperiodcentered{} \textit{il nido} nest \textperiodcentered{} \textit{il ramo} branch \textperiodcentered{} \textit{il fiume} river \textperiodcentered{} \textit{di} of; from (origin with essere) \textperiodcentered{} \textit{a} to; at; in (cities) \textperiodcentered{} \textit{da} from; at (someone's place); since \textperiodcentered{} \textit{in} in; into; to (countries, rooms) \textperiodcentered{} \textit{su} on; onto; about}
\applyit{Read \textit{Nel bosco} (A1), \S\ref{story:nel-bosco}. Read \textit{Il cane al mare} (A0–A1), \S\ref{story:il-cane-al-mare}. Study the rhyme \textit{La pioggia sul tetto}, \S\ref{rhyme:la-pioggia-sul-tetto}.}

\chapter{Modal verbs}\label{mod:m08}
\moduleintro{A1}{potere, volere, dovere, sapere + infinitive.}{Modal verbs: potere, volere, dovere, sapere (\S\ref{les:modal-verbs})}

\section{Modal verbs: potere, volere, dovere, sapere}\label{les:modal-verbs}
\lessonmeta{A1}{Can, want, must and know how to — with an infinitive.}
\begin{objectives}
\item\relax Conjugate potere, volere, dovere and sapere in the present.
\item\relax Build modal + infinitive sentences, including negatives and pronouns.
\item\relax Distinguish potere from sapere, and use vorrei for polite requests.
\end{objectives}

\subsection{Forms}\label{les:modal-verbs:forms}

\par\noindent\begin{minipage}[t]{0.48\linewidth}\raggedright
\paradigmbox{potere --- to be able to; can; may \textperiodcentered{} indicativo presente}{%
\begin{tabular}{@{}l l l@{}}
\textit{io} & \irr{posso} & {\color{muted}I can} \\
\textit{tu} & \irr{puoi} & {\color{muted}you can} \\
\textit{lui/lei} & \irr{può} & {\color{muted}he/she can} \\
\textit{noi} & \irr{possiamo} & {\color{muted}we can} \\
\textit{voi} & potete & {\color{muted}you (pl.) can} \\
\textit{loro} & \irr{possono} & {\color{muted}they can} \\
\end{tabular}}{\irr{Highlighted} forms differ from the regular pattern of the verb's family.}
\end{minipage}\hfill\begin{minipage}[t]{0.48\linewidth}\raggedright
\paradigmbox{volere --- to want \textperiodcentered{} indicativo presente}{%
\begin{tabular}{@{}l l l@{}}
\textit{io} & \irr{voglio} & {\color{muted}I want} \\
\textit{tu} & \irr{vuoi} & {\color{muted}you want} \\
\textit{lui/lei} & \irr{vuole} & {\color{muted}he/she wants} \\
\textit{noi} & \irr{vogliamo} & {\color{muted}we want} \\
\textit{voi} & volete & {\color{muted}you (pl.) want} \\
\textit{loro} & \irr{vogliono} & {\color{muted}they want} \\
\end{tabular}}{\irr{Highlighted} forms differ from the regular pattern of the verb's family.}
\end{minipage}\par\medskip\noindent\begin{minipage}[t]{0.48\linewidth}\raggedright
\paradigmbox{dovere --- to have to; must \textperiodcentered{} indicativo presente}{%
\begin{tabular}{@{}l l l@{}}
\textit{io} & \irr{devo} & {\color{muted}I must} \\
\textit{tu} & \irr{devi} & {\color{muted}you must} \\
\textit{lui/lei} & \irr{deve} & {\color{muted}he/she must} \\
\textit{noi} & \irr{dobbiamo} & {\color{muted}we must} \\
\textit{voi} & dovete & {\color{muted}you (pl.) must} \\
\textit{loro} & \irr{devono} & {\color{muted}they must} \\
\end{tabular}}{\irr{Highlighted} forms differ from the regular pattern of the verb's family.}
\end{minipage}\hfill\begin{minipage}[t]{0.48\linewidth}\raggedright
\paradigmbox{sapere --- to know (a fact); to know how to \textperiodcentered{} indicativo presente}{%
\begin{tabular}{@{}l l l@{}}
\textit{io} & \irr{so} & {\color{muted}I know} \\
\textit{tu} & \irr{sai} & {\color{muted}you know} \\
\textit{lui/lei} & \irr{sa} & {\color{muted}he/she knows} \\
\textit{noi} & \irr{sappiamo} & {\color{muted}we know} \\
\textit{voi} & sapete & {\color{muted}you (pl.) know} \\
\textit{loro} & \irr{sanno} & {\color{muted}they know} \\
\end{tabular}}{\irr{Highlighted} forms differ from the regular pattern of the verb's family.}
\end{minipage}\par\medskip

\subsection{Modal + infinitive}\label{les:modal-verbs:construction}

\begin{notebox}{rulecol}\notelabel{rulecol}{Rule}{The second verb stays in the infinitive}
A modal is followed directly by an infinitive, with no preposition: \textit{Posso entrare? Voglio dormire. Devo studiare. So nuotare.}
\end{notebox}

\begin{notebox}{rulecol}\notelabel{rulecol}{Rule}{Negation and pronouns}
\textit{non} goes before the modal: \textit{Non posso venire.} An object or reflexive pronoun can go before the modal or be attached to the end of the infinitive (which drops its final -e): \textit{Lo voglio vedere} = \textit{Voglio vederlo}; \textit{Mi devo alzare} = \textit{Devo alzarmi}.
\end{notebox}

\begin{notebox}{rulecol}\notelabel{rulecol}{Rule}{volere with a noun; vorrei}
\textit{volere} can also take a noun: \textit{Voglio un gelato}. For polite requests use the conditional \textit{vorrei} (I would like): \textit{Vorrei un bicchiere d'acqua, per favore.} Saying \textit{voglio} to a waiter sounds blunt.
\end{notebox}

\subsection{potere or sapere?}\label{les:modal-verbs:potere-sapere}

{\small\setlength{\LTpre}{2pt}\setlength{\LTpost}{6pt}
\tablecaption{Two kinds of 'can'}
\begin{xltabular}{\linewidth}{lLL}
\toprule
\textbf{Italian} & \textbf{Meaning} & \textbf{Example} \\
\midrule
sapere + infinitive & to know how to (a learned skill) & So nuotare. = I can (know how to) swim. \\
potere + infinitive & to be able / allowed to (circumstances, permission) & Oggi posso nuotare: il mare è calmo. = Today I can swim: the sea is calm. \\
\bottomrule
\end{xltabular}}

Alone, \textit{sapere} means to know a fact: \textit{So che piove. Non lo so.} (To know a person or place is \textit{conoscere}.)\par

\subsection*{Worked examples}
\begin{examples}
\item\relax \textit{Posso entrare?}\newline May I come in?
\item\relax \textit{Voglio un gelato.}\newline I want an ice cream.
\item\relax \textit{Dobbiamo tornare a casa.}\newline We have to go back home.
\item\relax \textit{Sai nuotare?}\newline Can you swim? (Do you know how to swim?)
\item\relax \textit{Non posso venire oggi.}\newline I can't come today.
\item\relax \textit{Vorrei un bicchiere d'acqua, per favore.}\newline I'd like a glass of water, please.
\end{examples}

\subsection*{Exercises}\label{les:modal-verbs:exercises}
{\color{muted}\small Answers: page~\pageref{ans:modal-verbs}.}

\begin{exercises}
\item\label{ex:modal-verbs-1}\relax Conjugate \textit{dovere}.\par\smallskip{\begin{tabular}{@{}l l@{}}
io & \blankcell \\
tu & \blankcell \\
lui/lei & \blankcell \\
noi & \blankcell \\
voi & \blankcell \\
loro & \blankcell \\
\end{tabular}}
\item\label{ex:modal-verbs-2}\relax potere, lui/lei: \blank{}
\item\label{ex:modal-verbs-3}\relax volere, noi: \blank{}
\item\label{ex:modal-verbs-4}\relax I know how to swim: \blank{} nuotare.\par\nopagebreak\begin{options}
\item\relax So
\item\relax Posso
\item\relax Conosco
\item\relax Devo
\end{options}
\item\label{ex:modal-verbs-5}\relax Order: 'I have to get up early.' (pronoun attached to the infinitive)\par\smallskip \token{alzarmi} \token{Devo} \token{presto}\par\answerline
\item\label{ex:modal-verbs-6}\relax Translate: I can't come today.\par\answerline
\item\label{ex:modal-verbs-7}\relax Translate politely: I'd like a glass of water, please.\par\answerline
\end{exercises}

\begin{summarybox}
\begin{itemize}
\item\relax potere: posso, puoi, può, possiamo, potete, possono.
\item\relax volere: voglio, vuoi, vuole, vogliamo, volete, vogliono.
\item\relax dovere: devo, devi, deve, dobbiamo, dovete, devono.
\item\relax sapere: so, sai, sa, sappiamo, sapete, sanno.
\item\relax Modal + bare infinitive; pronouns before the modal or attached to the infinitive.
\end{itemize}
\end{summarybox}

\relatedwords{\textit{potere} to be able to; can; may \textperiodcentered{} \textit{volere} to want \textperiodcentered{} \textit{dovere} to have to; must \textperiodcentered{} \textit{sapere} to know (a fact); to know how to \textperiodcentered{} \textit{nuotare} to swim \textperiodcentered{} \textit{conoscere} to know (a person or place)}
\applyit{Read \textit{L'orso che non poteva dormire} (A1–A2), \S\ref{story:lorso-che-non-poteva-dormire}. Read \textit{Una giornata di pioggia} (A1), \S\ref{story:una-giornata-di-pioggia}.}

\chapter{Reflexive verbs}\label{mod:m09}
\moduleintro{A1}{Reflexive pronouns and daily routines.}{Reflexive verbs and daily routine (\S\ref{les:reflexive-verbs})}

\section{Reflexive verbs and daily routine}\label{les:reflexive-verbs}
\lessonmeta{A1}{mi, ti, si, ci, vi, si with the verb; describing a day.}
\begin{objectives}
\item\relax Recognise reflexive infinitives in -arsi, -ersi, -irsi.
\item\relax Conjugate reflexive verbs with the correct pronoun.
\item\relax Describe a daily routine.
\end{objectives}

\subsection{Form}\label{les:reflexive-verbs:form}

A reflexive infinitive ends in \textit{-si}: \textit{svegliarsi} (to wake up), \textit{sedersi} (to sit down), \textit{vestirsi} (to get dressed). Remove \textit{-si}, put back the \textit{-e}, and you have an ordinary verb that conjugates normally (\textit{svegliare}, \textit{sedere}, \textit{vestire}).\par

\begin{notebox}{rulecol}\notelabel{rulecol}{Rule}{Pronoun + conjugated verb}
Each person takes its reflexive pronoun before the verb: \textit{mi, ti, si, ci, vi, si}.
\end{notebox}

\par\noindent\begin{minipage}[t]{0.48\linewidth}\raggedright
\paradigmbox{svegliarsi --- to wake up \textperiodcentered{} indicativo presente}{%
\begin{tabular}{@{}l l l@{}}
\textit{io} & mi sveglio & {\color{muted}I wake up} \\
\textit{tu} & ti svegli & {\color{muted}you wake up} \\
\textit{lui/lei} & si sveglia & {\color{muted}he/she wakes up} \\
\textit{noi} & ci svegliamo & {\color{muted}we wake up} \\
\textit{voi} & vi svegliate & {\color{muted}you (pl.) wake up} \\
\textit{loro} & si svegliano & {\color{muted}they wake up} \\
\end{tabular}}{}
\end{minipage}\hfill\begin{minipage}[t]{0.48\linewidth}\raggedright
\paradigmbox{vestirsi --- to get dressed \textperiodcentered{} indicativo presente}{%
\begin{tabular}{@{}l l l@{}}
\textit{io} & mi vesto & {\color{muted}I get dressed} \\
\textit{tu} & ti vesti & {\color{muted}you get dressed} \\
\textit{lui/lei} & si veste & {\color{muted}he/she gets dressed} \\
\textit{noi} & ci vestiamo & {\color{muted}we get dressed} \\
\textit{voi} & vi vestite & {\color{muted}you (pl.) get dressed} \\
\textit{loro} & si vestono & {\color{muted}they get dressed} \\
\end{tabular}}{}
\end{minipage}\par\medskip

\begin{notebox}{rulecol}\notelabel{rulecol}{Rule}{Position of the pronoun}
Before the conjugated verb (\textit{mi alzo}), after \textit{non} (\textit{non mi alzo}). With an infinitive it attaches to the end: \textit{alzarmi, alzarti, alzarsi…} — \textit{Devo alzarmi} or \textit{Mi devo alzare}.
\end{notebox}

\subsection{Meaning}\label{les:reflexive-verbs:meaning}

\begin{notebox}{tendcol}\notelabel{tendcol}{Tendency}{Not always 'oneself'}
Some reflexives really mean 'oneself' (\textit{lavarsi} = to wash oneself). Many simply correspond to an English intransitive verb: \textit{svegliarsi} = to wake up, \textit{addormentarsi} = to fall asleep, \textit{chiamarsi} = to be called, \textit{divertirsi} = to have fun.
\end{notebox}

\begin{notebox}{rulecol}\notelabel{rulecol}{Rule}{Body parts take the article, not a possessive}
\textit{Mi lavo le mani} = I wash my hands. The reflexive pronoun already shows whose hands they are.
\end{notebox}

\subsection{A day}\label{les:reflexive-verbs:routine}

\begin{itemize}
\item\relax \textit{Mi sveglio alle sette.} — I wake up at seven.
\item\relax \textit{Mi alzo e mi lavo.} — I get up and wash.
\item\relax \textit{Mi vesto e faccio colazione.} — I get dressed and have breakfast.
\item\relax \textit{La sera mi siedo sul divano e leggo.} — In the evening I sit on the sofa and read.
\item\relax \textit{Alle dieci mi addormento.} — At ten I fall asleep.
\end{itemize}

\subsection*{Worked examples}
\begin{examples}
\item\relax \textit{Come ti chiami? — Mi chiamo Sara.}\newline What's your name? — My name is Sara.
\item\relax \textit{Il gatto si lava.}\newline The cat washes itself.
\item\relax \textit{Ci divertiamo al mare.}\newline We have fun at the seaside.
\item\relax \textit{I bambini si addormentano presto.}\newline The children fall asleep early.
\item\relax \textit{Mi siedo sulla sedia.}\newline I sit down on the chair.\newline {\color{muted}\small sedersi: irregular stem sied-.}
\item\relax \textit{Non mi alzo presto la domenica.}\newline I don't get up early on Sundays.
\end{examples}

\subsection*{Exercises}\label{les:reflexive-verbs:exercises}
{\color{muted}\small Answers: page~\pageref{ans:reflexive-verbs}.}

\begin{exercises}
\item\label{ex:reflexive-verbs-1}\relax Conjugate \textit{alzarsi} in the present (include the pronoun).\par\smallskip{\begin{tabular}{@{}l l@{}}
io & \blankcell \\
tu & \blankcell \\
lui/lei & \blankcell \\
noi & \blankcell \\
voi & \blankcell \\
loro & \blankcell \\
\end{tabular}}
\item\label{ex:reflexive-verbs-2}\relax Noi \blank{} divertiamo.\par\nopagebreak\begin{options}
\item\relax ti
\item\relax ci
\item\relax si
\item\relax vi
\end{options}
\item\label{ex:reflexive-verbs-3}\relax sedersi, io: \blank{}
\item\label{ex:reflexive-verbs-4}\relax Order: 'I don't get up early.'\par\smallskip \token{alzo} \token{mi} \token{Non} \token{presto}\par\answerline
\item\label{ex:reflexive-verbs-5}\relax Translate: I wash my hands.\par\answerline
\item\label{ex:reflexive-verbs-6}\relax Infinitive of \textit{si addormentano}?\par\answerline
\item\label{ex:reflexive-verbs-7}\relax Complete the routine.\par\smallskip\begingroup\setlength{\parindent}{0pt}La mattina io \clozeblank{1} (svegliarsi) alle sette, \clozeblank{2} (vestirsi) e faccio colazione. La sera i bambini \clozeblank{3} (addormentarsi) presto.\par\endgroup
\end{exercises}

\begin{summarybox}
\begin{itemize}
\item\relax Reflexive infinitives end in -si; conjugate the base verb and add mi, ti, si, ci, vi, si before it.
\item\relax With an infinitive the pronoun attaches: alzarmi.
\item\relax Many reflexives translate as simple English verbs (wake up, fall asleep, be called).
\item\relax Mi lavo le mani: article, not possessive.
\end{itemize}
\end{summarybox}

\relatedwords{\textit{svegliarsi} to wake up \textperiodcentered{} \textit{alzarsi} to get up \textperiodcentered{} \textit{lavarsi} to wash (oneself) \textperiodcentered{} \textit{chiamarsi} to be called \textperiodcentered{} \textit{divertirsi} to have fun; to enjoy oneself \textperiodcentered{} \textit{addormentarsi} to fall asleep \textperiodcentered{} \textit{vestirsi} to get dressed \textperiodcentered{} \textit{sedersi} to sit down}
\applyit{Read \textit{Una giornata di pioggia} (A1), \S\ref{story:una-giornata-di-pioggia}. Study the rhyme \textit{Ninna nanna della luna}, \S\ref{rhyme:ninna-nanna-della-luna}.}

\chapter{Imperativo}\label{mod:m15}
\moduleintro{A2}{Commands for tu, noi, voi and the formal Lei/Loro; the negative imperative; pronouns with the imperative; short forms (va’, fa’, di’).}{The imperative (\S\ref{les:imperativo})}

\section{The imperative}\label{les:imperativo}
\lessonmeta{A2}{Affirmative and negative commands for tu, noi, voi and formal Lei/Loro; irregular short forms; position of pronouns.}
\begin{objectives}
\item\relax Form the imperative of regular and irregular verbs.
\item\relax Form the negative tu imperative with non + infinitive.
\item\relax Use the formal Lei imperative (from the congiuntivo presente).
\item\relax Place object and reflexive pronouns correctly, including the doubled consonant after va', fa', da', di', sta'.
\end{objectives}

\subsection{Affirmative forms}\label{les:imperativo:forms}

\begin{notebox}{rulecol}\notelabel{rulecol}{Rule}{tu: -a for -are verbs, -i for the others}
For tu, -are verbs use the lui/lei present form (\textit{parla!}); -ere and -ire verbs use the tu present form (\textit{prendi! dormi! finisci!}). noi and voi are identical to the present (\textit{parliamo! parlate!}).
\end{notebox}

\par\noindent\begin{minipage}[t]{0.48\linewidth}\raggedright
\paradigmbox{parlare --- to speak; to talk \textperiodcentered{} imperativo}{%
\begin{tabular}{@{}l l@{}}
\textit{tu} & parla \\
\textit{Lei} & parli \\
\textit{noi} & parliamo \\
\textit{voi} & parlate \\
\textit{Loro} & parlino \\
\end{tabular}}{Negative \textit{tu}: \textit{non parlare}. \textit{Lei} and \textit{Loro} are the congiuntivo forms.}
\end{minipage}\hfill\begin{minipage}[t]{0.48\linewidth}\raggedright
\paradigmbox{prendere --- to take; to catch; to have (food) \textperiodcentered{} imperativo}{%
\begin{tabular}{@{}l l@{}}
\textit{tu} & prendi \\
\textit{Lei} & prenda \\
\textit{noi} & prendiamo \\
\textit{voi} & prendete \\
\textit{Loro} & prendano \\
\end{tabular}}{Negative \textit{tu}: \textit{non prendere}. \textit{Lei} and \textit{Loro} are the congiuntivo forms.}
\end{minipage}\par\medskip

\begin{notebox}{rulecol}\notelabel{rulecol}{Rule}{Formal Lei and Loro}
The formal forms are borrowed from the congiuntivo presente: \textit{(Lei) parli! prenda! venga!} Formal plural \textit{Loro} (\textit{parlino!}) is now limited to very formal service language; ordinarily \textit{voi} is used for several people.
\end{notebox}

\subsection{Irregular forms}\label{les:imperativo:irregular}

{\small\setlength{\LTpre}{2pt}\setlength{\LTpost}{6pt}
\tablecaption{Irregular tu forms}
\begin{longtable}{lll}
\toprule
\textbf{Verb} & \textbf{tu} & \textbf{Note} \\
\midrule
andare & va' / vai & both standard \\
dare & da' / dai & both standard \\
fare & fa' / fai & both standard \\
stare & sta' / stai & both standard \\
dire & di' & only form \\
essere & sii & voi: siate \\
avere & abbi & voi: abbiate \\
sapere & sappi & voi: sappiate \\
\bottomrule
\end{longtable}}

\par\noindent\begin{minipage}[t]{0.48\linewidth}\raggedright
\paradigmbox{essere --- to be \textperiodcentered{} imperativo}{%
\begin{tabular}{@{}l l@{}}
\textit{tu} & \irr{sii} \\
\textit{Lei} & \irr{sia} \\
\textit{noi} & \irr{siamo} \\
\textit{voi} & \irr{siate} \\
\textit{Loro} & \irr{siano} \\
\end{tabular}}{\irr{Highlighted} forms differ from the regular pattern of the verb's family. Negative \textit{tu}: \textit{non essere}. \textit{Lei} and \textit{Loro} are the congiuntivo forms.}
\end{minipage}\par\medskip

\begin{notebox}{exccol}\notelabel{exccol}{Exception}{No imperative}
\textit{potere} and \textit{dovere} have no imperative. The imperative of \textit{volere} (\textit{vogli, vogliate}) is rare and literary.
\end{notebox}

\subsection{Negative imperative}\label{les:imperativo:negative}

\begin{notebox}{rulecol}\notelabel{rulecol}{Rule}{tu: non + infinitive}
\textit{Non parlare! Non prendere! Non andare!} For all other persons, put non before the affirmative form: \textit{Non parliamo! Non parlate! Non parli!} (Lei).
\end{notebox}

\subsection{Pronouns with the imperative}\label{les:imperativo:pronouns}

\begin{notebox}{rulecol}\notelabel{rulecol}{Rule}{Attached to tu, noi, voi}
Object and reflexive pronouns follow and attach to the tu, noi and voi forms: \textit{guardalo! alzati! alziamoci! sedetevi! dimmelo!} With Lei and Loro they precede: \textit{lo guardi! si alzi!}
\end{notebox}

\par\noindent\begin{minipage}[t]{0.48\linewidth}\raggedright
\paradigmbox{alzarsi --- to get up \textperiodcentered{} imperativo}{%
\begin{tabular}{@{}l l@{}}
\textit{tu} & alzati \\
\textit{Lei} & si alzi \\
\textit{noi} & alziamoci \\
\textit{voi} & alzatevi \\
\textit{Loro} & si alzino \\
\end{tabular}}{Negative \textit{tu}: \textit{non alzarti / non ti alzare}. \textit{Lei} and \textit{Loro} are the congiuntivo forms.}
\end{minipage}\par\medskip

\begin{notebox}{rulecol}\notelabel{rulecol}{Rule}{Doubling after the short forms}
After \textit{va', da', fa', di', sta'} the consonant of the pronoun doubles: \textit{dammi} (give me), \textit{fallo} (do it), \textit{dimmi} (tell me), \textit{vattene} (go away), \textit{stammi bene}. Exception: \textit{gli} does not double (\textit{dagli}).
\end{notebox}

\begin{notebox}{rulecol}\notelabel{rulecol}{Rule}{Negative tu: two positions}
With the negative tu imperative the pronoun may attach to the infinitive or precede it: \textit{Non alzarti! = Non ti alzare!} \textit{Non dirlo! = Non lo dire!} Both are standard.
\end{notebox}

\subsection*{Worked examples}
\begin{examples}
\item\relax \textit{Chiudi la porta, per favore!}\newline Close the door, please!
\item\relax \textit{Non toccare il forno: è caldo!}\newline Don't touch the oven: it's hot!\newline {\color{muted}\small Negative tu = non + infinitive.}
\item\relax \textit{Signora, si accomodi!}\newline Madam, please take a seat!\newline {\color{muted}\small Formal Lei: pronoun before the verb.}
\item\relax \textit{Dimmi la verità.}\newline Tell me the truth.\newline {\color{muted}\small di' + mi → dimmi (doubled m).}
\item\relax \textit{Alziamoci presto domani.}\newline Let's get up early tomorrow.\newline {\color{muted}\small noi: pronoun attached.}
\item\relax \textit{Ragazzi, sedetevi e ascoltate.}\newline Children, sit down and listen.
\item\relax \textit{Sii gentile con tua sorella.}\newline Be kind to your sister.
\end{examples}

\subsection*{Exercises}\label{les:imperativo:exercises}
{\color{muted}\small Answers: page~\pageref{ans:imperativo}.}

\begin{exercises}
\item\label{ex:imperativo-1}\relax Give the tu imperative.\par\smallskip{\begin{tabular}{@{}l l@{}}
parlare & \blankcell \\
scrivere & \blankcell \\
finire & \blankcell \\
essere & \blankcell \\
avere & \blankcell \\
dire & \blankcell \\
\end{tabular}}
\item\label{ex:imperativo-2}\relax 'Don't go!' (tu)\par\nopagebreak\begin{options}
\item\relax Non vada!
\item\relax Non vai!
\item\relax Non andare!
\item\relax Non va'!
\end{options}
\item\label{ex:imperativo-3}\relax venire, Lei (formal) imperative: \blank{}
\item\label{ex:imperativo-4}\relax fa' + lo (Do it!): \blank{}
\item\label{ex:imperativo-5}\relax 'Sit down!' (voi)\par\nopagebreak\begin{options}
\item\relax Sedeteci!
\item\relax Sedetevi!
\item\relax Siedetevi!
\item\relax Vi sedete!
\end{options}
\item\label{ex:imperativo-6}\relax Translate (tu): Don't say it!\par\answerline
\item\label{ex:imperativo-7}\relax Give the imperative of \textit{alzarsi}.\par\smallskip{\begin{tabular}{@{}l l@{}}
tu & \blankcell \\
Lei & \blankcell \\
noi & \blankcell \\
voi & \blankcell \\
\end{tabular}}
\end{exercises}

\begin{summarybox}
\begin{itemize}
\item\relax tu: -a (-are), -i (-ere/-ire); noi and voi = present; Lei/Loro = congiuntivo presente.
\item\relax Irregular: va'/vai, da'/dai, fa'/fai, sta'/stai, di', sii, abbi, sappi.
\item\relax Negative tu: non + infinitive. Other persons: non + affirmative form.
\item\relax Pronouns attach to tu/noi/voi (dimmi, alziamoci), precede Lei/Loro (si alzi); consonants double after va', da', fa', di', sta'.
\end{itemize}
\end{summarybox}

\relatedwords{\textit{parlare} to speak; to talk \textperiodcentered{} \textit{prendere} to take; to catch; to have (food) \textperiodcentered{} \textit{chiudere} to close \textperiodcentered{} \textit{dire} to say; to tell \textperiodcentered{} \textit{essere} to be \textperiodcentered{} \textit{avere} to have \textperiodcentered{} \textit{alzarsi} to get up \textperiodcentered{} \textit{sedersi} to sit down \textperiodcentered{} \textit{andare} to go \textperiodcentered{} \textit{fare} to do; to make}
\applyit{Read \textit{La colazione} (A0–A1), \S\ref{story:la-colazione}. Study the rhyme \textit{Ninna nanna della luna}, \S\ref{rhyme:ninna-nanna-della-luna}.}

\chapter{Passato prossimo and auxiliaries}\label{mod:m10}
\moduleintro{A2}{The compound past; choosing avere or essere (including verbs that take both and modal verbs); full past-participle agreement.}{Passato prossimo with avere (\S\ref{les:passato-prossimo-avere}); Passato prossimo with essere; agreement (\S\ref{les:passato-prossimo-essere}); Choosing the auxiliary and agreeing the participle (\S\ref{les:ausiliari-accordo})}

\section{Passato prossimo with avere}\label{les:passato-prossimo-avere}
\lessonmeta{A2}{The compound past: avere + past participle; regular and irregular participles.}
\begin{objectives}
\item\relax Form regular past participles (-ato, -uto, -ito).
\item\relax Learn the common irregular participles.
\item\relax Form the passato prossimo with avere and use it for completed events.
\end{objectives}

\subsection{Formation}\label{les:passato-prossimo-avere:formation}

\begin{notebox}{rulecol}\notelabel{rulecol}{Rule}{auxiliary + past participle}
The passato prossimo is formed with the present of \textit{avere} (or \textit{essere}, next lesson) + the past participle: \textit{ho mangiato} = I ate / I have eaten.
\end{notebox}

{\small\setlength{\LTpre}{2pt}\setlength{\LTpost}{6pt}
\tablecaption{Regular participles}
\begin{longtable}{lll}
\toprule
\textbf{Infinitive} & \textbf{Ending} & \textbf{Participle} \\
\midrule
parlare & -are → -ato & parlato \\
credere & -ere → -uto & creduto \\
dormire & -ire → -ito & dormito \\
\bottomrule
\end{longtable}}

\par\noindent\begin{minipage}[t]{0.48\linewidth}\raggedright
\paradigmbox{mangiare --- to eat \textperiodcentered{} indicativo passato prossimo (with avere)}{%
\begin{tabular}{@{}l l@{}}
\textit{io} & ho mangiato \\
\textit{tu} & hai mangiato \\
\textit{lui/lei} & ha mangiato \\
\textit{noi} & abbiamo mangiato \\
\textit{voi} & avete mangiato \\
\textit{loro} & hanno mangiato \\
\end{tabular}}{}
\end{minipage}\par\medskip

\begin{notebox}{rulecol}\notelabel{rulecol}{Rule}{With avere the participle does not change}
With \textit{avere} the participle normally stays in -o, whoever the subject is: \textit{Sara ha mangiato, i bambini hanno mangiato}. The main exception is a 3rd-person direct object pronoun placed before the verb: \textit{Le mele? Le ho mangiate.} All the agreement rules are in \lessonref{ausiliari-accordo}{auxiliaries and agreement}.
\end{notebox}

\subsection{Irregular participles}\label{les:passato-prossimo-avere:irregular}

Many -ere verbs, and a few others, have irregular participles. They must be learned, but they fall into families.\par

{\small\setlength{\LTpre}{2pt}\setlength{\LTpost}{6pt}
\tablecaption{Common irregular participles}
\begin{longtable}{llll}
\toprule
\textbf{Infinitive} & \textbf{Participle} & \textbf{Infinitive} & \textbf{Participle} \\
\midrule
fare & fatto & dire & detto \\
vedere & visto & prendere & preso \\
leggere & letto & scrivere & scritto \\
mettere & messo & chiudere & chiuso \\
aprire & aperto & rispondere & risposto \\
perdere & perso & vincere & vinto \\
bere & bevuto & piangere & pianto \\
ridere & riso & correre & corso \\
\bottomrule
\end{longtable}}

\subsection{Use}\label{les:passato-prossimo-avere:use}

\begin{notebox}{tendcol}\notelabel{tendcol}{Tendency}{Completed events}
The passato prossimo expresses a completed action in the past — both 'I ate' and 'I have eaten'. In much of Italy it is the normal past tense of conversation. In Tuscany and much of the South the \lessonref{passato-remoto}{passato remoto} is still used in speech for events felt as concluded; in written narration it is the standard narrative tense everywhere.
\end{notebox}

\begin{itemize}
\item\relax Typical time expressions: \textit{ieri} (yesterday), \textit{stamattina} (this morning), \textit{la settimana scorsa} (last week), \textit{un'ora fa} (an hour ago), \textit{già} (already), \textit{non … ancora} (not yet), \textit{mai} (ever / never).
\item\relax \textit{già}, \textit{mai} and \textit{ancora} usually go between the auxiliary and the participle: \textit{Ho già mangiato. Non ho mai visto un orso.}
\end{itemize}

\subsection*{Worked examples}
\begin{examples}
\item\relax \textit{Ieri ho mangiato una pizza.}\newline Yesterday I ate a pizza.
\item\relax \textit{Hai visto il gufo?}\newline Did you see the owl? / Have you seen the owl?
\item\relax \textit{Il nonno ha letto una storia.}\newline Grandfather read a story.
\item\relax \textit{Abbiamo fatto un castello di sabbia.}\newline We made a sandcastle.
\item\relax \textit{Non ho mai visto un lupo.}\newline I have never seen a wolf.
\item\relax \textit{Hanno già finito.}\newline They have already finished.
\end{examples}

\subsection*{Exercises}\label{les:passato-prossimo-avere:exercises}
{\color{muted}\small Answers: page~\pageref{ans:passato-prossimo-avere}.}

\begin{exercises}
\item\label{ex:passato-prossimo-avere-1}\relax Give the past participle.\par\smallskip{\begin{tabular}{@{}l l@{}}
giocare & \blankcell \\
credere & \blankcell \\
sentire & \blankcell \\
fare & \blankcell \\
vedere & \blankcell \\
scrivere & \blankcell \\
\end{tabular}}
\item\label{ex:passato-prossimo-avere-2}\relax Passato prossimo of \textit{mangiare}, noi: \blank{}
\item\label{ex:passato-prossimo-avere-3}\relax Passato prossimo of \textit{prendere}, lui/lei: \blank{}
\item\label{ex:passato-prossimo-avere-4}\relax Which infinitive does the participle \textit{detto} come from?\par\answerline
\item\label{ex:passato-prossimo-avere-5}\relax Order: 'I have never seen a wolf.'\par\smallskip \token{ho} \token{lupo} \token{mai} \token{Non} \token{un} \token{visto}\par\answerline
\item\label{ex:passato-prossimo-avere-6}\relax Translate: Yesterday Grandfather read a story.\par\answerline
\item\label{ex:passato-prossimo-avere-7}\relax Complete with the passato prossimo.\par\smallskip\begingroup\setlength{\parindent}{0pt}Ieri io \clozeblank{1} (fare) una torta e i bambini \clozeblank{2} (bere) il latte. Tu \clozeblank{3} (vedere) il film?\par\endgroup
\end{exercises}

\begin{summarybox}
\begin{itemize}
\item\relax passato prossimo = present of avere/essere + past participle.
\item\relax Regular participles: -ato, -uto, -ito.
\item\relax Common irregulars: fatto, detto, visto, preso, letto, scritto, messo, chiuso, aperto, risposto, perso, vinto, bevuto.
\item\relax With avere the participle stays in -o.
\end{itemize}
\end{summarybox}

\relatedwords{\textit{avere} to have \textperiodcentered{} \textit{mangiare} to eat \textperiodcentered{} \textit{fare} to do; to make \textperiodcentered{} \textit{vedere} to see \textperiodcentered{} \textit{prendere} to take; to catch; to have (food) \textperiodcentered{} \textit{leggere} to read \textperiodcentered{} \textit{scrivere} to write \textperiodcentered{} \textit{bere} to drink \textperiodcentered{} \textit{ieri} yesterday \textperiodcentered{} \textit{già} already \textperiodcentered{} \textit{mai} never; ever}
\applyit{Read \textit{L'orso che non poteva dormire} (A1–A2), \S\ref{story:lorso-che-non-poteva-dormire}. Read \textit{Il cane al mare} (A0–A1), \S\ref{story:il-cane-al-mare}.}

\section{Passato prossimo with essere; agreement}\label{les:passato-prossimo-essere}
\lessonmeta{A2}{Which verbs take essere, agreement of the participle, and reflexive verbs in the past.}
\begin{objectives}
\item\relax Identify verbs that take essere.
\item\relax Make the participle agree with the subject.
\item\relax Form the passato prossimo of reflexive verbs.
\end{objectives}

\subsection{Which verbs take essere?}\label{les:passato-prossimo-essere:which}

\begin{notebox}{tendcol}\notelabel{tendcol}{Tendency}{Movement, change of state, staying}
Most verbs that take \textit{essere} express movement to or from a place, a change of state, or remaining: \textit{andare, venire, arrivare, partire, entrare, uscire, tornare, cadere, salire, scendere, nascere, morire, diventare, restare, rimanere, stare} — and \textit{essere} itself. Not every verb of movement takes essere (\textit{camminare, nuotare, ballare} take \textit{avere}), so check the dictionary.
\end{notebox}

\begin{notebox}{rulecol}\notelabel{rulecol}{Rule}{All reflexive verbs take essere}
\textit{mi sono svegliato, ti sei alzata, si è lavato…}
\end{notebox}

\begin{notebox}{rulecol}\notelabel{rulecol}{Rule}{piacere takes essere}
\textit{Il film mi è piaciuto. Le storie mi sono piaciute.}
\end{notebox}

\subsection{Agreement with essere}\label{les:passato-prossimo-essere:agreement}

\begin{notebox}{rulecol}\notelabel{rulecol}{Rule}{The participle agrees with the subject}
With \textit{essere} the participle behaves like an adjective: \textit{-o, -a, -i, -e}.
\end{notebox}

\par\noindent\begin{minipage}[t]{0.48\linewidth}\raggedright
\paradigmbox{andare --- to go \textperiodcentered{} indicativo passato prossimo (with essere, masculine subject)}{%
\begin{tabular}{@{}l l@{}}
\textit{io} & sono andato \\
\textit{tu} & sei andato \\
\textit{lui/lei} & è andato \\
\textit{noi} & siamo andati \\
\textit{voi} & siete andati \\
\textit{loro} & sono andati \\
\end{tabular}}{}
\end{minipage}\par\medskip

{\small\setlength{\LTpre}{2pt}\setlength{\LTpost}{6pt}
\tablecaption{Agreement}
\begin{xltabular}{\linewidth}{lL}
\toprule
\textbf{Subject} & \textbf{Form} \\
\midrule
Marco & è andato \\
Sara & è andata \\
Marco e Luca & sono andati \\
Sara e Giulia & sono andate \\
Marco e Sara & sono andati (mixed → masculine) \\
\bottomrule
\end{xltabular}}

\subsection{Reflexive verbs}\label{les:passato-prossimo-essere:reflexive}

\par\noindent\begin{minipage}[t]{0.48\linewidth}\raggedright
\paradigmbox{svegliarsi --- to wake up \textperiodcentered{} indicativo passato prossimo (with essere, masculine subject)}{%
\begin{tabular}{@{}l l@{}}
\textit{io} & mi sono svegliato \\
\textit{tu} & ti sei svegliato \\
\textit{lui/lei} & si è svegliato \\
\textit{noi} & ci siamo svegliati \\
\textit{voi} & vi siete svegliati \\
\textit{loro} & si sono svegliati \\
\end{tabular}}{}
\end{minipage}\par\medskip

For a female speaker: \textit{mi sono svegliata}; for a group of women: \textit{ci siamo svegliate}.\par

\subsection*{Worked examples}
\begin{examples}
\item\relax \textit{Sara è andata al mare.}\newline Sara went to the seaside.
\item\relax \textit{I bambini sono tornati a casa.}\newline The children came back home.
\item\relax \textit{La foglia è caduta dall'albero.}\newline The leaf fell from the tree.
\item\relax \textit{Mi sono svegliata presto.}\newline I woke up early. (female speaker)
\item\relax \textit{Ci siamo divertiti molto.}\newline We had a lot of fun.
\item\relax \textit{La storia ti è piaciuta?}\newline Did you like the story?
\end{examples}

\subsection*{Exercises}\label{les:passato-prossimo-essere:exercises}
{\color{muted}\small Answers: page~\pageref{ans:passato-prossimo-essere}.}

\begin{exercises}
\item\label{ex:passato-prossimo-essere-1}\relax Ieri Sara \blank{} al mare.\par\nopagebreak\begin{options}
\item\relax è andato
\item\relax è andata
\item\relax ha andato
\item\relax ha andata
\end{options}
\item\label{ex:passato-prossimo-essere-2}\relax Which auxiliary does each verb take?\par\smallskip\begin{tabularx}{\linewidth}{@{}L L@{}}
1.\ arrivare & (a)\ essere \\
2.\ mangiare & (b)\ avere \\
3.\ partire & (c)\ essere \\
4.\ camminare & (d)\ essere \\
5.\ alzarsi & (e)\ avere \\
\end{tabularx}
\item\label{ex:passato-prossimo-essere-3}\relax Le ragazze (tornare) → Le ragazze \blank{} a casa.
\item\label{ex:passato-prossimo-essere-4}\relax svegliarsi, io (male speaker): \blank{}
\item\label{ex:passato-prossimo-essere-5}\relax Complete with the passato prossimo.\par\smallskip\begingroup\setlength{\parindent}{0pt}La foglia \clozeblank{1} (cadere) dall'albero. Noi \clozeblank{2} (divertirsi) molto. Il gatto \clozeblank{3} (salire) sul tetto.\par\endgroup
\item\label{ex:passato-prossimo-essere-6}\relax Translate: The children came back home.\par\answerline
\end{exercises}

\begin{summarybox}
\begin{itemize}
\item\relax essere: most verbs of movement to/from a place, change of state, staying; all reflexives; piacere.
\item\relax With essere the participle agrees with the subject (-o, -a, -i, -e).
\item\relax Reflexive: mi sono svegliato/a, ci siamo divertiti/e.
\end{itemize}
\end{summarybox}

\relatedwords{\textit{essere} to be \textperiodcentered{} \textit{andare} to go \textperiodcentered{} \textit{tornare} to return; to come back \textperiodcentered{} \textit{arrivare} to arrive \textperiodcentered{} \textit{partire} to leave; to set off \textperiodcentered{} \textit{cadere} to fall \textperiodcentered{} \textit{salire} to go up; to climb; to get on \textperiodcentered{} \textit{svegliarsi} to wake up \textperiodcentered{} \textit{divertirsi} to have fun; to enjoy oneself \textperiodcentered{} \textit{piacere} to please; to be liked}
\applyit{Read \textit{L'orso che non poteva dormire} (A1–A2), \S\ref{story:lorso-che-non-poteva-dormire}. Read \textit{Una giornata di pioggia} (A1), \S\ref{story:una-giornata-di-pioggia}.}

\section{Choosing the auxiliary and agreeing the participle}\label{les:ausiliari-accordo}
\lessonmeta{A2}{The complete rules for avere vs essere — including verbs that take both, weather verbs and modal verbs — and for past-participle agreement with the subject, with object pronouns, with ne and in relative clauses.}
\begin{objectives}
\item\relax Choose avere or essere for any verb, and recognise the verbs that change auxiliary with their meaning.
\item\relax Choose the auxiliary of a modal verb followed by an infinitive.
\item\relax Make the participle agree with the subject (essere) and with a preceding lo, la, li, le or ne (avere).
\item\relax Know where agreement is optional (mi, ti, ci, vi) and where it is no longer standard (relative che).
\end{objectives}

\subsection{avere or essere?}\label{les:ausiliari-accordo:avere-essere}

The choice of auxiliary is the same in every compound tense (passato prossimo, trapassato, futuro anteriore, condizionale passato, congiuntivo passato …), so it is learned once.\par

\begin{notebox}{rulecol}\notelabel{rulecol}{Rule}{avere}
All transitive verbs used transitively, whether or not the object is expressed: \textit{ho mangiato (la pizza), abbiamo letto}. Also many intransitive verbs that describe an activity of the subject: \textit{ho dormito, hai camminato, abbiamo riso, hanno parlato, ho lavorato, ho viaggiato, ho telefonato}.
\end{notebox}

\begin{notebox}{rulecol}\notelabel{rulecol}{Rule}{essere}
1) Most intransitive verbs of movement with a direction or of change of state: \textit{andare, venire, arrivare, partire, entrare, uscire, tornare, cadere, nascere, morire, diventare}. 2) Verbs of state: \textit{essere, stare, rimanere, restare}. 3) Verbs whose subject is the thing experienced: \textit{piacere, sembrare, bastare, mancare, costare, succedere, servire} (to be needed): \textit{Il film mi è piaciuto.} 4) All reflexive and pronominal verbs: \textit{mi sono lavato, ci siamo divertiti}. 5) The passive: \textit{La casa è stata venduta}. 6) Impersonal \textit{si} in compound tenses: \textit{Si è mangiato bene}.
\end{notebox}

\begin{notebox}{tipcol}\notelabel{tipcol}{Tip}{A useful test}
If the participle can be used as an adjective describing the subject (\textit{una ragazza partita, un bambino nato ieri}), the verb usually takes essere. It is a guide, not a rule: learn the auxiliary with each new verb — the vocabulary list in Appendix B always shows it.
\end{notebox}

\subsection{Verbs that take both}\label{les:ausiliari-accordo:both}

Many verbs can be used transitively (avere) or intransitively (essere). The auxiliary follows the use, not the verb.\par

{\small\setlength{\LTpre}{2pt}\setlength{\LTpost}{6pt}
\tablecaption{One verb, two auxiliaries}
\begin{xltabular}{\linewidth}{lLL}
\toprule
\textbf{Verb} & \textbf{with avere (transitive)} & \textbf{with essere (intransitive)} \\
\midrule
finire & \textit{Ho finito il libro.} & \textit{Il film è finito.} \\
cominciare & \textit{Abbiamo cominciato la lezione.} & \textit{La lezione è cominciata.} \\
passare & \textit{Ho passato l'estate al mare.} & \textit{Sono passato da te.} \\
salire / scendere & \textit{Ho sceso le scale.} & \textit{Sono sceso dal treno.} \\
crescere & \textit{Ha cresciuto tre figli.} & \textit{Sono cresciuto in campagna.} \\
cambiare & \textit{Ho cambiato casa.} & \textit{La città è cambiata.} \\
correre & \textit{Ho corso per un'ora.} (activity) & \textit{Sono corso a casa.} (movement to a goal) \\
\bottomrule
\end{xltabular}}

\begin{notebox}{tendcol}\notelabel{tendcol}{Tendency}{Weather verbs}
\textit{piovere, nevicare} and similar can take either auxiliary: \textit{è piovuto} and \textit{ha piovuto} are both standard. Traditional grammars prefer essere; avere is very common and not an error.
\end{notebox}

\begin{notebox}{tendcol}\notelabel{tendcol}{Tendency}{vivere}
\textit{ho vissuto} and \textit{sono vissuto} are both accepted when there is no object; avere is more common. With an object only avere: \textit{ho vissuto un'avventura}.
\end{notebox}

\subsection{Modal verbs + infinitive}\label{les:ausiliari-accordo:modals}

\begin{notebox}{rulecol}\notelabel{rulecol}{Rule}{The modal can take the auxiliary of the infinitive}
\textit{dovere, potere, volere} followed by an infinitive can take the auxiliary the infinitive would take: \textit{Sono dovuto partire} (partire → essere), \textit{Ho dovuto lavorare} (lavorare → avere). With an essere-infinitive, \textit{avere} is also accepted: \textit{Ho dovuto partire}. Both are correct; essere is the more traditional choice.
\end{notebox}

\begin{notebox}{rulecol}\notelabel{rulecol}{Rule}{With a pronoun before the modal}
If a pronoun linked to an essere-verb is placed before the modal, essere is required: \textit{Non ci sono voluti andare.} If the pronoun is attached to the infinitive, avere is normal: \textit{Non hanno voluto andarci.}
\end{notebox}

\begin{notebox}{rulecol}\notelabel{rulecol}{Rule}{When the infinitive is essere, or passive}
With \textit{essere} itself or a passive infinitive the modal takes avere: \textit{Ho dovuto essere paziente. Avrebbe voluto essere invitato.}
\end{notebox}

\subsection{Past participle agreement}\label{les:ausiliari-accordo:agreement}

\begin{notebox}{rulecol}\notelabel{rulecol}{Rule}{With essere: agree with the subject}
\textit{Marco è partito, Sara è partita, i ragazzi sono partiti, le ragazze sono partite.} A mixed group takes the masculine plural. With formal \textit{Lei} the participle agrees with the real sex of the person: \textit{Signora, è arrivata?}
\end{notebox}

\begin{notebox}{rulecol}\notelabel{rulecol}{Rule}{With avere: normally invariable}
\textit{Sara ha mangiato, le ragazze hanno mangiato.} The participle does not agree with the subject or with a following object: \textit{Ho comprato le mele.}
\end{notebox}

\begin{notebox}{rulecol}\notelabel{rulecol}{Rule}{Obligatory: preceding lo, la, li, le}
When a 3rd-person direct object pronoun comes before the verb, the participle agrees with it: \textit{Le mele? Le ho comprate. La lettera? L'ho scritta. I libri? Li ho letti.} \textit{lo} and \textit{la} elide before \textit{ho, hai, ha, hanno}: \textit{l'ho visto, l'ho vista}.
\end{notebox}

\begin{notebox}{tendcol}\notelabel{tendcol}{Tendency}{Usual: partitive ne}
With \textit{ne} ('of it/of them') the participle normally agrees with the noun ne refers to: \textit{Di mele, ne ho mangiate tre.} The invariable form (\textit{ne ho mangiato tre}) is also found, especially in speech.
\end{notebox}

\begin{notebox}{tendcol}\notelabel{tendcol}{Tendency}{Optional: mi, ti, ci, vi}
With 1st and 2nd person direct object pronouns agreement is optional: \textit{Ci hanno visti} and \textit{Ci hanno visto} are both correct; so are \textit{Ti ho vista} and \textit{Ti ho visto} (said to a woman).
\end{notebox}

\begin{notebox}{rulecol}\notelabel{rulecol}{Rule}{Never: indirect object pronouns}
\textit{gli, le} (to him, to her), \textit{mi, ti, ci, vi} used as indirect objects do not trigger agreement: \textit{Le ho parlato} (I spoke to her) — compare \textit{L'ho vista} (I saw her).
\end{notebox}

\begin{notebox}{exccol}\notelabel{exccol}{Exception}{Relative che: no longer standard}
Agreement with the relative pronoun \textit{che} (\textit{la ragazza che ho conosciuta}) is found in older and literary texts; in modern standard Italian the participle stays invariable: \textit{la ragazza che ho conosciuto}.
\end{notebox}

\begin{notebox}{tendcol}\notelabel{tendcol}{Tendency}{Reflexive verbs with an object}
With a reflexive verb and a direct object, the participle normally agrees with the subject: \textit{Sara si è lavata le mani.} Agreement with the object (\textit{si è lavate le mani}) also occurs but is less common.
\end{notebox}

\subsection*{Worked examples}
\begin{examples}
\item\relax \textit{Le ragazze sono uscite alle otto.}\newline The girls went out at eight.\newline {\color{muted}\small essere: uscite agrees with le ragazze.}
\item\relax \textit{Il film è finito tardi.}\newline The film finished late.\newline {\color{muted}\small Intransitive finire → essere.}
\item\relax \textit{Abbiamo finito i compiti.}\newline We have finished our homework.\newline {\color{muted}\small Transitive finire → avere.}
\item\relax \textit{Le chiavi? Le ho perse.}\newline The keys? I've lost them.\newline {\color{muted}\small le before the verb: perse agrees.}
\item\relax \textit{Le ho telefonato ieri.}\newline I phoned her yesterday.\newline {\color{muted}\small le = to her (indirect): no agreement.}
\item\relax \textit{Sono dovuta partire presto.}\newline I had to leave early.\newline {\color{muted}\small Female speaker; essere because partire takes essere. Ho dovuto partire is also correct.}
\item\relax \textit{Di biscotti, ne abbiamo mangiati tanti.}\newline We ate lots of biscuits.\newline {\color{muted}\small ne: agreement with biscotti (mangiato also found).}
\item\relax \textit{Ieri è piovuto tutto il giorno.}\newline Yesterday it rained all day.\newline {\color{muted}\small Also: ha piovuto.}
\end{examples}

\subsection*{Exercises}\label{les:ausiliari-accordo:exercises}
{\color{muted}\small Answers: page~\pageref{ans:ausiliari-accordo}.}

\begin{exercises}
\item\label{ex:ausiliari-accordo-1}\relax Which auxiliary? Match each sentence with its auxiliary.\par\smallskip\begin{tabularx}{\linewidth}{@{}L L@{}}
1.\ Il concerto \blank{} cominciato. & (a)\ è \\
2.\ Noi \blank{} cominciato il lavoro. & (b)\ è \\
3.\ Il prezzo \blank{} cambiato. & (c)\ hanno \\
4.\ Loro \blank{} cambiato idea. & (d)\ abbiamo \\
\end{tabularx}
\item\label{ex:ausiliari-accordo-2}\relax Complete with the passato prossimo.\par\smallskip\begingroup\setlength{\parindent}{0pt}Sara \clozeblank{1} (nascere) a Napoli, ma \clozeblank{2} (crescere) a Torino. Il film ci \clozeblank{3} (piacere) molto.\par\endgroup
\item\label{ex:ausiliari-accordo-3}\relax Replace the object with a pronoun: \textit{Ho comprato le mele.} → \blank{}
\item\label{ex:ausiliari-accordo-4}\relax \textit{Signora, Le ho \blank{} una lettera.} (scrivere: I wrote a letter to you)\par\nopagebreak\begin{options}
\item\relax scritta
\item\relax scritti
\item\relax scritto
\item\relax scritte
\end{options}
\item\label{ex:ausiliari-accordo-5}\relax Passato prossimo: \textit{Ci (vedere, loro) al mercato.} Give the verb (we are a group of men and women).\par\answerline
\item\label{ex:ausiliari-accordo-6}\relax Which sentence is NOT standard modern Italian?\par\nopagebreak\begin{options}
\item\relax Non ci hanno voluti andare.
\item\relax Sono dovuti partire.
\item\relax Non hanno voluto andarci.
\item\relax Non ci sono voluti andare.
\end{options}
\item\label{ex:ausiliari-accordo-7}\relax Translate (speaker is a woman): I had to leave.\par\answerline
\item\label{ex:ausiliari-accordo-8}\relax Complete with the participle.\par\smallskip\begingroup\setlength{\parindent}{0pt}Di torte, ne abbiamo \clozeblank{1} (fare) due. La ragazza che ho \clozeblank{2} (conoscere) ieri è simpatica.\par\endgroup
\end{exercises}

\begin{summarybox}
\begin{itemize}
\item\relax avere: transitive verbs and activity intransitives; essere: movement/change of state, state, piacere-type verbs, reflexives, passive, impersonal si.
\item\relax Verbs like finire, cominciare, passare, cambiare, crescere take avere with an object and essere without.
\item\relax Modal + essere-verb: sono dovuto partire or ho dovuto partire; pronoun before the modal → essere.
\item\relax Agreement: with essere always (subject); with avere only with preceding lo/la/li/le (obligatory), ne (usual), mi/ti/ci/vi (optional); never with indirect objects or (in modern standard) relative che.
\end{itemize}
\end{summarybox}

\relatedwords{\textit{finire} to finish; to end \textperiodcentered{} \textit{cominciare} to begin; to start \textperiodcentered{} \textit{passare} to pass; to spend (time); to drop by \textperiodcentered{} \textit{cambiare} to change \textperiodcentered{} \textit{crescere} to grow \textperiodcentered{} \textit{scendere} to go down; to get off \textperiodcentered{} \textit{correre} to run \textperiodcentered{} \textit{piovere} to rain \textperiodcentered{} \textit{dovere} to have to; must \textperiodcentered{} \textit{piacere} to please; to be liked}

\chapter{Imperfetto and trapassato prossimo}\label{mod:m11}
\moduleintro{A2}{The imperfect and its uses, choosing between imperfetto and passato prossimo, and the pluperfect.}{The imperfetto (\S\ref{les:imperfetto}); Imperfetto or passato prossimo? (\S\ref{les:imperfetto-passato}); Trapassato prossimo (\S\ref{les:trapassato-prossimo})}

\section{The imperfetto}\label{les:imperfetto}
\lessonmeta{A2}{Formation of the imperfect tense and its uses: habits, descriptions and background.}
\begin{objectives}
\item\relax Form the imperfetto of regular verbs and of essere, fare, dire, bere.
\item\relax Use it for habits in the past and for descriptions.
\item\relax Recognise it in stories: C'era una volta…
\item\relax Recognise the polite, intentional, counterfactual (colloquial) and reported uses of the imperfetto.
\end{objectives}

\subsection{Formation}\label{les:imperfetto:formation}

\begin{notebox}{rulecol}\notelabel{rulecol}{Rule}{Remove -re, add -vo, -vi, -va, -vamo, -vate, -vano}
The imperfetto is the most regular tense in Italian: \textit{parla-re → parlavo}, \textit{vede-re → vedevo}, \textit{dormi-re → dormivo}. Even most irregular verbs are regular here.
\end{notebox}

\par\noindent\begin{minipage}[t]{0.48\linewidth}\raggedright
\paradigmbox{parlare --- to speak; to talk \textperiodcentered{} indicativo imperfetto}{%
\begin{tabular}{@{}l l@{}}
\textit{io} & parl\hlending{avo} \\
\textit{tu} & parl\hlending{avi} \\
\textit{lui/lei} & parl\hlending{ava} \\
\textit{noi} & parl\hlending{avamo} \\
\textit{voi} & parl\hlending{avate} \\
\textit{loro} & parl\hlending{avano} \\
\end{tabular}}{Stem \textit{parl-} + \hlending{ending}.}
\end{minipage}\hfill\begin{minipage}[t]{0.48\linewidth}\raggedright
\paradigmbox{dormire --- to sleep \textperiodcentered{} indicativo imperfetto}{%
\begin{tabular}{@{}l l@{}}
\textit{io} & dormivo \\
\textit{tu} & dormivi \\
\textit{lui/lei} & dormiva \\
\textit{noi} & dormivamo \\
\textit{voi} & dormivate \\
\textit{loro} & dormivano \\
\end{tabular}}{}
\end{minipage}\par\medskip

\begin{notebox}{exccol}\notelabel{exccol}{Exception}{essere is irregular}
\textit{ero, eri, era, eravamo, eravate, erano}.
\end{notebox}

\par\noindent\begin{minipage}[t]{0.48\linewidth}\raggedright
\paradigmbox{essere --- to be \textperiodcentered{} indicativo imperfetto}{%
\begin{tabular}{@{}l l@{}}
\textit{io} & \irr{ero} \\
\textit{tu} & \irr{eri} \\
\textit{lui/lei} & \irr{era} \\
\textit{noi} & \irr{eravamo} \\
\textit{voi} & \irr{eravate} \\
\textit{loro} & \irr{erano} \\
\end{tabular}}{\irr{Highlighted} forms differ from the regular pattern of the verb's family.}
\end{minipage}\par\medskip

\begin{notebox}{exccol}\notelabel{exccol}{Exception}{fare, dire, bere use their old stems}
\textit{fare → facevo}, \textit{dire → dicevo}, \textit{bere → bevevo} — from the older forms \textit{facere, dicere, bevere}.
\end{notebox}

\subsection{Uses}\label{les:imperfetto:uses}

\begin{itemize}
\item\relax Habits and repeated actions in the past: \textit{Da bambino giocavo sempre in giardino.} (As a child I always played in the garden.)
\item\relax Descriptions of people, places, weather, time and age: \textit{Era notte. Il bosco era buio. L'orso aveva sonno.}
\item\relax States of mind and feelings: \textit{Il topo aveva paura.}
\item\relax Actions in progress at a moment in the past: \textit{Alle otto dormivo.} (At eight I was sleeping.)
\item\relax The fairy-tale opening: \textit{C'era una volta…} (Once upon a time there was…).
\end{itemize}

\begin{notebox}{tendcol}\notelabel{tendcol}{Tendency}{English equivalents}
The imperfetto corresponds to English 'used to + verb', 'was/were + -ing', or a simple past of description ('the wood was dark').
\end{notebox}

\begin{notebox}{rulecol}\notelabel{rulecol}{Rule}{Polite or attenuated requests}
With verbs such as \textit{volere, venire, chiamare} the imperfetto softens a present request: \textit{Volevo un caffè, per favore. Telefonavo per chiedere un’informazione.} The meaning is present; the past tense makes it less direct, like the condizionale (\textit{vorrei}).
\end{notebox}

\begin{notebox}{rulecol}\notelabel{rulecol}{Rule}{Intentions that did not happen}
\textit{Stavo per uscire} (I was about to go out), \textit{Dovevo partire ieri} (I was supposed to leave yesterday — and perhaps did not).
\end{notebox}

\begin{notebox}{tendcol}\notelabel{tendcol}{Tendency}{Colloquial counterfactuals}
In informal speech the imperfetto can replace both verbs of an unreal past hypothesis: \textit{Se lo sapevo, venivo} for standard \textit{Se l’avessi saputo, sarei venuto.} It is very common in conversation and is attested in literature, but it is not standard in formal writing. See the \lessonref{periodo-ipotetico}{periodo ipotetico}.
\end{notebox}

\begin{notebox}{rulecol}\notelabel{rulecol}{Rule}{Reported speech}
After a past verb of saying, a present of direct speech becomes imperfetto: \textit{«Ho fame»} → \textit{Ha detto che aveva fame.} See \lessonref{discorso-indiretto}{indirect speech}.
\end{notebox}

\subsection*{Worked examples}
\begin{examples}
\item\relax \textit{C'era una volta un piccolo orso.}\newline Once upon a time there was a little bear.
\item\relax \textit{Era inverno e faceva freddo.}\newline It was winter and it was cold.
\item\relax \textit{Da bambina abitavo al mare.}\newline As a child I lived by the sea.
\item\relax \textit{Il nonno ci raccontava sempre delle storie.}\newline Grandfather always used to tell us stories.
\item\relax \textit{L'orso aveva sonno ma non dormiva.}\newline The bear was sleepy but wasn't sleeping.
\item\relax \textit{Alle otto i bambini dormivano.}\newline At eight the children were sleeping.
\item\relax \textit{Volevo un caffè, per favore.}\newline I’d like a coffee, please.\newline {\color{muted}\small Polite imperfetto: the request is in the present.}
\item\relax \textit{Se lo sapevo, venivo.}\newline If I had known, I would have come.\newline {\color{muted}\small Colloquial; standard: Se l’avessi saputo, sarei venuto.}
\end{examples}

\subsection*{Exercises}\label{les:imperfetto:exercises}
{\color{muted}\small Answers: page~\pageref{ans:imperfetto}.}

\begin{exercises}
\item\label{ex:imperfetto-1}\relax Conjugate \textit{cantare} in the imperfetto.\par\smallskip{\begin{tabular}{@{}l l@{}}
io & \blankcell \\
tu & \blankcell \\
lui/lei & \blankcell \\
noi & \blankcell \\
voi & \blankcell \\
loro & \blankcell \\
\end{tabular}}
\item\label{ex:imperfetto-2}\relax Conjugate \textit{essere} in the imperfetto.\par\smallskip{\begin{tabular}{@{}l l@{}}
io & \blankcell \\
tu & \blankcell \\
lui/lei & \blankcell \\
noi & \blankcell \\
voi & \blankcell \\
loro & \blankcell \\
\end{tabular}}
\item\label{ex:imperfetto-3}\relax fare, lui/lei, imperfetto: \blank{}
\item\label{ex:imperfetto-4}\relax Infinitive of \textit{dicevano}?\par\answerline
\item\label{ex:imperfetto-5}\relax Complete the description with the imperfetto.\par\smallskip\begingroup\setlength{\parindent}{0pt}\clozeblank{1} (essere) notte. Il bosco \clozeblank{2} (essere) buio e \clozeblank{3} (fare) freddo. L'orso \clozeblank{4} (avere) sonno.\par\endgroup
\item\label{ex:imperfetto-6}\relax Translate: As a child I lived by the sea. (male speaker)\par\answerline
\item\label{ex:imperfetto-7}\relax Translate: \textit{C'era una volta un piccolo orso.}\par\answerline
\end{exercises}

\begin{summarybox}
\begin{itemize}
\item\relax Imperfetto: infinitive minus -re + -vo, -vi, -va, -vamo, -vate, -vano.
\item\relax Irregular: essere (ero…), fare (facevo), dire (dicevo), bere (bevevo).
\item\relax Used for habits, descriptions, states, actions in progress, and C'era una volta.
\end{itemize}
\end{summarybox}

\relatedwords{\textit{essere} to be \textperiodcentered{} \textit{avere} to have \textperiodcentered{} \textit{fare} to do; to make \textperiodcentered{} \textit{dire} to say; to tell \textperiodcentered{} \textit{bere} to drink \textperiodcentered{} \textit{c'era una volta} once upon a time (there was) \textperiodcentered{} \textit{la volta} time (occasion)}
\applyit{Read \textit{L'orso che non poteva dormire} (A1–A2), \S\ref{story:lorso-che-non-poteva-dormire}. Study the rhyme \textit{Ninna nanna della luna}, \S\ref{rhyme:ninna-nanna-della-luna}.}

\section{Imperfetto or passato prossimo?}\label{les:imperfetto-passato}
\lessonmeta{A2}{Background versus events: choosing the past tense in a story.}
\begin{objectives}
\item\relax Distinguish background (imperfetto) from events (passato prossimo).
\item\relax Use mentre with the imperfetto and quando with an interrupting event.
\item\relax Narrate a short sequence in the past.
\end{objectives}

\subsection{Background and events}\label{les:imperfetto-passato:contrast}

{\small\setlength{\LTpre}{2pt}\setlength{\LTpost}{6pt}
\tablecaption{Which tense?}
\begin{xltabular}{\linewidth}{LL}
\toprule
\textbf{Imperfetto} & \textbf{Passato prossimo} \\
\midrule
Background, scene-setting: \textit{Era sera. Pioveva.} & Events that move the story forward: \textit{Il telefono ha squillato.} \\
What was going on: \textit{Dormivo.} & What happened: \textit{Mi sono svegliato.} \\
Habits: \textit{Ogni estate andavamo al mare.} & A single occasion: \textit{Quell'estate siamo andati in montagna.} \\
States without a clear end: \textit{Avevo paura.} & A completed, bounded action: \textit{Ho avuto paura per un momento.} \\
\bottomrule
\end{xltabular}}

\begin{notebox}{tendcol}\notelabel{tendcol}{Tendency}{A useful test}
Ask: is this the stage (imperfetto) or something that happens on the stage (passato prossimo)? A completed action with a clear beginning and end takes the passato prossimo; an ongoing situation or repeated habit takes the imperfetto.
\end{notebox}

\subsection{mentre and quando}\label{les:imperfetto-passato:mentre-quando}

\begin{notebox}{rulecol}\notelabel{rulecol}{Rule}{Interrupted action}
An action in progress (imperfetto) is interrupted by an event (passato prossimo): \textit{Mentre dormivo, è arrivato il nonno.} / \textit{Dormivo quando è arrivato il nonno.}
\end{notebox}

\begin{notebox}{rulecol}\notelabel{rulecol}{Rule}{Simultaneous background actions}
Two things going on at the same time both take the imperfetto: \textit{Mentre la mamma cucinava, io leggevo.}
\end{notebox}

\subsection*{Worked examples}
\begin{examples}
\item\relax \textit{Era notte e pioveva.}\newline It was night and it was raining.
\item\relax \textit{Mentre dormivo, è arrivato il nonno.}\newline While I was sleeping, Grandfather arrived.
\item\relax \textit{Ogni estate andavamo al mare.}\newline Every summer we used to go to the seaside.
\item\relax \textit{Quell'estate siamo andati in montagna.}\newline That summer we went to the mountains.
\item\relax \textit{Il topo mangiava il formaggio quando è entrato il gatto.}\newline The mouse was eating the cheese when the cat came in.
\item\relax \textit{Mentre la mamma cucinava, io leggevo.}\newline While Mum was cooking, I was reading.
\end{examples}

\subsection*{Exercises}\label{les:imperfetto-passato:exercises}
{\color{muted}\small Answers: page~\pageref{ans:imperfetto-passato}.}

\begin{exercises}
\item\label{ex:imperfetto-passato-1}\relax Mentre \blank{} (dormire, io), è arrivato il nonno.\par\nopagebreak\begin{options}
\item\relax dormivo
\item\relax dormo
\item\relax dormii
\item\relax ho dormito
\end{options}
\item\label{ex:imperfetto-passato-2}\relax Il topo mangiava il formaggio quando il gatto \blank{} (entrare).\par\nopagebreak\begin{options}
\item\relax entrava
\item\relax è entrato
\item\relax entra
\item\relax ha entrato
\end{options}
\item\label{ex:imperfetto-passato-3}\relax Ogni estate \blank{} (andare, noi) al mare.\par\nopagebreak\begin{options}
\item\relax siamo andati
\item\relax andiamo
\item\relax andremo
\item\relax andavamo
\end{options}
\item\label{ex:imperfetto-passato-4}\relax Quell'estate \blank{} (andare, noi) in montagna.\par\nopagebreak\begin{options}
\item\relax andavamo
\item\relax andiamo
\item\relax siamo andati
\item\relax abbiamo andato
\end{options}
\item\label{ex:imperfetto-passato-5}\relax Choose the tense.\par\smallskip\begingroup\setlength{\parindent}{0pt}\clozeblank{1} (essere) notte e \clozeblank{2} (piovere). All'improvviso l'orso \clozeblank{3} (sentire) un rumore.\par\endgroup
\item\label{ex:imperfetto-passato-6}\relax Translate: While Mum was cooking, I was reading.\par\answerline
\end{exercises}

\begin{summarybox}
\begin{itemize}
\item\relax Imperfetto = the stage: background, description, habits, actions in progress.
\item\relax Passato prossimo = events on the stage: completed actions.
\item\relax mentre + imperfetto; an interrupting event in the passato prossimo.
\end{itemize}
\end{summarybox}

\relatedwords{\textit{mentre} while \textperiodcentered{} \textit{quando} when \textperiodcentered{} \textit{l'estate} summer \textperiodcentered{} \textit{piovere} to rain \textperiodcentered{} \textit{sentire} to hear; to feel \textperiodcentered{} \textit{il rumore} noise}
\applyit{Read \textit{L'orso che non poteva dormire} (A1–A2), \S\ref{story:lorso-che-non-poteva-dormire}.}

\section{Trapassato prossimo}\label{les:trapassato-prossimo}
\lessonmeta{A2}{The pluperfect: an action completed before another past moment. Formation, uses in narration and reported speech, and when the passato prossimo is used instead.}
\begin{objectives}
\item\relax Form the trapassato prossimo with the imperfetto of avere or essere + past participle.
\item\relax Use it for an action that precedes another past action or moment.
\item\relax Use it in reported speech after a past reporting verb.
\item\relax Distinguish it from the trapassato remoto, which is literary.
\end{objectives}

\subsection{Formation}\label{les:trapassato-prossimo:formation}

\begin{notebox}{rulecol}\notelabel{rulecol}{Rule}{Imperfetto of the auxiliary + participle}
\textit{avevo mangiato} (I had eaten), \textit{ero partito/a} (I had left). The auxiliary and the agreement follow exactly the same rules as in the passato prossimo.
\end{notebox}

\par\noindent\begin{minipage}[t]{0.48\linewidth}\raggedright
\paradigmbox{mangiare --- to eat \textperiodcentered{} indicativo trapassato prossimo (with avere)}{%
\begin{tabular}{@{}l l@{}}
\textit{io} & avevo mangiato \\
\textit{tu} & avevi mangiato \\
\textit{lui/lei} & aveva mangiato \\
\textit{noi} & avevamo mangiato \\
\textit{voi} & avevate mangiato \\
\textit{loro} & avevano mangiato \\
\end{tabular}}{}
\end{minipage}\hfill\begin{minipage}[t]{0.48\linewidth}\raggedright
\paradigmbox{partire --- to leave; to set off \textperiodcentered{} indicativo trapassato prossimo (with essere, feminine subject)}{%
\begin{tabular}{@{}l l@{}}
\textit{io} & ero partita \\
\textit{tu} & eri partita \\
\textit{lui/lei} & era partita \\
\textit{noi} & eravamo partite \\
\textit{voi} & eravate partite \\
\textit{loro} & erano partite \\
\end{tabular}}{}
\end{minipage}\par\medskip

\subsection{Uses}\label{les:trapassato-prossimo:uses}

\begin{notebox}{rulecol}\notelabel{rulecol}{Rule}{Before another past moment}
The trapassato prossimo places an action before a reference point that is already in the past: \textit{Quando sono arrivato alla stazione, il treno era già partito.} The order of events is: the train left → I arrived.
\end{notebox}

\begin{notebox}{rulecol}\notelabel{rulecol}{Rule}{Background in narration}
In a story told in the passato prossimo or the passato remoto, it fills in what had happened earlier: \textit{Sara era stanca: aveva lavorato tutto il giorno.}
\end{notebox}

\begin{notebox}{rulecol}\notelabel{rulecol}{Rule}{Reported speech}
After a past verb of saying or thinking, a passato prossimo of the direct speech becomes a trapassato prossimo: \textit{«Ho finito»} → \textit{Ha detto che aveva finito.} See the lesson on indirect speech.
\end{notebox}

\begin{notebox}{tendcol}\notelabel{tendcol}{Tendency}{When the order is obvious}
If the sequence is already clear from the context or from words like \textit{prima} or \textit{dopo}, the passato prossimo is often used instead, especially in conversation: \textit{Prima ho mangiato, poi sono uscito.} The trapassato is needed when the earlier event is mentioned after the later one.
\end{notebox}

\begin{notebox}{rulecol}\notelabel{rulecol}{Rule}{Trapassato remoto}
In literary narration in the passato remoto, an action that immediately precedes another, after \textit{dopo che, appena, quando, (non) appena}, can be expressed with the trapassato remoto: \textit{Dopo che ebbe mangiato, uscì.} It is limited to this construction and to writing; the trapassato prossimo is also correct there.
\end{notebox}

\subsection*{Worked examples}
\begin{examples}
\item\relax \textit{Quando siamo arrivati, il film era già cominciato.}\newline When we arrived, the film had already started.\newline {\color{muted}\small cominciare without an object → essere.}
\item\relax \textit{Non avevo mai visto il mare prima di quel giorno.}\newline I had never seen the sea before that day.
\item\relax \textit{Sara era stanca perché aveva lavorato tutto il giorno.}\newline Sara was tired because she had worked all day.
\item\relax \textit{Marco ha detto che aveva perso le chiavi.}\newline Marco said that he had lost his keys.\newline {\color{muted}\small Reported speech: «Ho perso le chiavi».}
\item\relax \textit{Avevo appena finito di mangiare quando è arrivata la nonna.}\newline I had just finished eating when Grandma arrived.
\item\relax \textit{Dopo che ebbe finito la lettera, la chiuse.}\newline After she had finished the letter, she closed it.\newline {\color{muted}\small Trapassato remoto: literary.}
\end{examples}

\subsection*{Exercises}\label{les:trapassato-prossimo:exercises}
{\color{muted}\small Answers: page~\pageref{ans:trapassato-prossimo}.}

\begin{exercises}
\item\label{ex:trapassato-prossimo-1}\relax Conjugate \textit{uscire} in the trapassato prossimo (masculine).\par\smallskip{\begin{tabular}{@{}l l@{}}
io & \blankcell \\
tu & \blankcell \\
lui & \blankcell \\
noi & \blankcell \\
voi & \blankcell \\
loro & \blankcell \\
\end{tabular}}
\item\label{ex:trapassato-prossimo-2}\relax Complete with the trapassato prossimo.\par\smallskip\begingroup\setlength{\parindent}{0pt}Quando sono arrivato, i miei amici \clozeblank{1} (mangiare) già e Sara \clozeblank{2} (andare) a casa.\par\endgroup
\item\label{ex:trapassato-prossimo-3}\relax \textit{Quando ho chiamato, Luca era già uscito.} What happened first?\par\nopagebreak\begin{options}
\item\relax it is not clear
\item\relax I called
\item\relax Luca went out
\item\relax they happened at the same time
\end{options}
\item\label{ex:trapassato-prossimo-4}\relax Report it: \textit{«Ho finito i compiti»}, ha detto Anna. → Anna ha detto che \blank{} i compiti.
\item\label{ex:trapassato-prossimo-5}\relax Translate: I had never seen a wolf.\par\answerline
\item\label{ex:trapassato-prossimo-6}\relax \textit{Dopo che \blank{} la lettera, la chiuse.} Which option is the traditional literary choice?\par\nopagebreak\begin{options}
\item\relax ha scritto
\item\relax ebbe scritto
\item\relax avrà scritto
\item\relax scriveva
\end{options}
\end{exercises}

\begin{summarybox}
\begin{itemize}
\item\relax Trapassato prossimo = imperfetto of avere/essere + participle (avevo fatto, ero partito/a).
\item\relax It marks an action completed before another past moment.
\item\relax In reported speech, passato prossimo → trapassato prossimo after a past reporting verb.
\item\relax With prima/poi or clear order, conversation often uses the passato prossimo; the trapassato remoto is literary.
\end{itemize}
\end{summarybox}

\relatedwords{\textit{mangiare} to eat \textperiodcentered{} \textit{partire} to leave; to set off \textperiodcentered{} \textit{uscire} to go out; to leave \textperiodcentered{} \textit{finire} to finish; to end \textperiodcentered{} \textit{perdere} to lose; to miss \textperiodcentered{} \textit{arrivare} to arrive}
\applyit{Read \textit{Il segreto del faro} (A2–B1), \S\ref{story:il-segreto-del-faro}.}

\chapter{Pronouns and complex sentences}\label{mod:m12}
\moduleintro{A2}{Direct and indirect object pronouns, ne and combined pronouns, relative clauses (che, cui, il quale), adverbs and story expressions.}{Direct object pronouns (\S\ref{les:direct-object-pronouns}); Indirect object pronouns (\S\ref{les:indirect-object-pronouns}); Relative clauses, adverbs and expressions in stories (\S\ref{les:relative-clauses-adverbs})}

\section{Direct object pronouns}\label{les:direct-object-pronouns}
\lessonmeta{A2}{mi, ti, lo, la, ci, vi, li, le — position, elision and agreement in the passato prossimo.}
\begin{objectives}
\item\relax Replace a direct object noun with the right pronoun.
\item\relax Place the pronoun before the verb or attach it to an infinitive.
\item\relax Make the participle agree after lo, la, li, le.
\item\relax Use ne for quantities and for di + something.
\end{objectives}

\subsection{Forms}\label{les:direct-object-pronouns:forms}

{\small\setlength{\LTpre}{2pt}\setlength{\LTpost}{6pt}
\tablecaption{Direct object pronouns}
\begin{longtable}{lllll}
\toprule
\textbf{} & \textbf{Singular} & \textbf{} & \textbf{Plural} & \textbf{} \\
\midrule
1st & mi & me & ci & us \\
2nd & ti & you & vi & you (pl.) \\
3rd masc. & lo & him, it & li & them (m.) \\
3rd fem. & la & her, it & le & them (f.) \\
formal & La & you (formal) &  &  \\
\bottomrule
\end{longtable}}

\begin{notebox}{rulecol}\notelabel{rulecol}{Rule}{The pronoun replaces the whole noun phrase}
\textit{Vedo il gatto → Lo vedo. Mangio la mela → La mangio. Chiamo i bambini → Li chiamo. Guardo le stelle → Le guardo.} Gender and number come from the noun replaced.
\end{notebox}

\subsection{Position}\label{les:direct-object-pronouns:position}

\begin{notebox}{rulecol}\notelabel{rulecol}{Rule}{Before the conjugated verb}
\textit{Lo vedo. Non lo vedo. Ti chiamo domani.}
\end{notebox}

\begin{notebox}{rulecol}\notelabel{rulecol}{Rule}{Attached to an infinitive}
The infinitive loses its final -e: \textit{vedere + lo → vederlo}. With modals either position works: \textit{Lo voglio vedere = Voglio vederlo.}
\end{notebox}

\begin{notebox}{rulecol}\notelabel{rulecol}{Rule}{Elision before avere}
\textit{lo} and \textit{la} usually elide before a vowel or h: \textit{l'ho visto, l'ha presa}. \textit{li} and \textit{le} do not: \textit{li ho visti}.
\end{notebox}

\subsection{Agreement in the passato prossimo}\label{les:direct-object-pronouns:agreement}

\begin{notebox}{rulecol}\notelabel{rulecol}{Rule}{Agreement with lo, la, li, le}
When \textit{lo, la, li, le} come before \textit{avere}, the participle agrees with the pronoun: \textit{Hai visto la volpe? — Sì, l'ho vista. Hai visto i gufi? — Sì, li ho visti. Le stelle? — Le ho guardate.}
\end{notebox}

\begin{notebox}{tendcol}\notelabel{tendcol}{Tendency}{Optional with mi, ti, ci, vi}
With \textit{mi, ti, ci, vi} agreement is optional: \textit{Sara, ti ho vista / ti ho visto.}
\end{notebox}

\begin{notebox}{rulecol}\notelabel{rulecol}{Rule}{Not with indirect objects}
There is no agreement with indirect object pronouns: \textit{Hai scritto a Sara? — Sì, le ho scritto.} For the complete rules (ne, reflexive verbs, relative che, modal verbs) see \lessonref{ausiliari-accordo}{auxiliaries and agreement}.
\end{notebox}

\subsection{The pronoun ne}\label{les:direct-object-pronouns:ne}

\begin{notebox}{rulecol}\notelabel{rulecol}{Rule}{Partitive ne: a quantity of something}
\textit{ne} replaces a noun when you mention a quantity of it, or none: \textit{Quante mele hai? — Ne ho tre. Vuoi del pane? — No, grazie, non ne voglio.} The number or quantity stays in the sentence; \textit{ne} replaces the noun. English often has nothing: 'I have three (of them)'.
\end{notebox}

\begin{notebox}{rulecol}\notelabel{rulecol}{Rule}{ne = di + something}
\textit{ne} also replaces \textit{di} + a thing or idea: \textit{Ne parliamo domani} (we will talk about it tomorrow), \textit{Che ne pensi?} (what do you think of it?). Less often it replaces \textit{da} + a place: \textit{Ne è uscito} (he came out of it).
\end{notebox}

\begin{notebox}{rulecol}\notelabel{rulecol}{Rule}{Position}
Like the other unstressed pronouns, \textit{ne} goes before a finite verb and is attached to an infinitive, a gerund or an affirmative imperative (tu, noi, voi): \textit{Voglio comprarne due. Prendine uno!} With \textit{c'è / ci sono}: \textit{ce n'è uno, ce ne sono tre}.
\end{notebox}

\begin{notebox}{tendcol}\notelabel{tendcol}{Tendency}{Agreement}
In compound tenses the participle normally agrees with the noun ne replaces: \textit{Di mele, ne ho mangiate tre.} The invariable \textit{mangiato} is also found.
\end{notebox}

\subsection*{Worked examples}
\begin{examples}
\item\relax \textit{Dov'è il gatto? — Non lo vedo.}\newline Where's the cat? — I can't see it.
\item\relax \textit{La mela? La mangio io.}\newline The apple? I'll eat it.
\item\relax \textit{La mamma chiama i bambini e li porta a casa.}\newline Mum calls the children and takes them home.
\item\relax \textit{Ti aspetto qui.}\newline I'll wait for you here.
\item\relax \textit{Voglio vederlo.}\newline I want to see him / it.
\item\relax \textit{Hai visto la volpe? — Sì, l'ho vista.}\newline Did you see the fox? — Yes, I saw it.
\item\relax \textit{Quanti fratelli hai? — Ne ho due.}\newline How many brothers do you have? — I have two.\newline {\color{muted}\small Partitive ne.}
\end{examples}

\subsection*{Exercises}\label{les:direct-object-pronouns:exercises}
{\color{muted}\small Answers: page~\pageref{ans:direct-object-pronouns}.}

\begin{exercises}
\item\label{ex:direct-object-pronouns-1}\relax Vedo le stelle. → \blank{} vedo.\par\nopagebreak\begin{options}
\item\relax Gli
\item\relax Li
\item\relax La
\item\relax Le
\end{options}
\item\label{ex:direct-object-pronouns-2}\relax Chiamo i bambini. → \blank{} chiamo.\par\nopagebreak\begin{options}
\item\relax Lo
\item\relax Le
\item\relax Li
\item\relax Gli
\end{options}
\item\label{ex:direct-object-pronouns-3}\relax Order: 'I don't see it.' (il gatto)\par\smallskip \token{lo} \token{Non} \token{vedo}\par\answerline
\item\label{ex:direct-object-pronouns-4}\relax Hai visto la volpe? — Sì, \blank{} (I saw it).
\item\label{ex:direct-object-pronouns-5}\relax Rewrite with the pronoun attached: Voglio vedere il mare. → Voglio \blank{}.
\item\label{ex:direct-object-pronouns-6}\relax In \textit{La mangio}, what is \textit{la}?\par\nopagebreak\begin{options}
\item\relax direct object pronoun (it, feminine)
\item\relax indirect object pronoun
\item\relax subject pronoun
\item\relax definite article
\end{options}
\item\label{ex:direct-object-pronouns-ne1}\relax Complete with ne or a direct object pronoun.\par\smallskip\begingroup\setlength{\parindent}{0pt}Quante torte fai? — \clozeblank{1} faccio due. Il libro? \clozeblank{2} leggo stasera.\par\endgroup
\item\label{ex:direct-object-pronouns-ne2}\relax Answer with ne and the passato prossimo: \textit{Quante mele hai mangiato?} (tre) → \blank{} tre.
\end{exercises}

\begin{summarybox}
\begin{itemize}
\item\relax mi, ti, lo/la, ci, vi, li/le replace the direct object.
\item\relax Before the conjugated verb (after non); attached to an infinitive (vederlo).
\item\relax lo/la elide before avere (l'ho); the participle agrees with lo, la, li, le.
\item\relax ne replaces a noun after a quantity (ne ho tre) or di + something (ne parliamo); its position is like the other pronouns.
\end{itemize}
\end{summarybox}

\relatedwords{\textit{vedere} to see \textperiodcentered{} \textit{chiamare} to call \textperiodcentered{} \textit{aspettare} to wait (for) \textperiodcentered{} \textit{guardare} to look at; to watch \textperiodcentered{} \textit{la volpe} fox}
\applyit{Read \textit{Nel bosco} (A1), \S\ref{story:nel-bosco}. Read \textit{L'orso che non poteva dormire} (A1–A2), \S\ref{story:lorso-che-non-poteva-dormire}.}

\section{Indirect object pronouns}\label{les:indirect-object-pronouns}
\lessonmeta{A2}{mi, ti, gli, le, ci, vi, gli — 'to me, to him…', with dare, dire, piacere.}
\begin{objectives}
\item\relax Use indirect object pronouns with verbs of giving and saying.
\item\relax Distinguish gli/le (indirect) from lo/la/li (direct).
\item\relax Use piacere with all persons.
\item\relax Combine indirect and direct pronouns (me lo, glielo).
\end{objectives}

\subsection{Forms}\label{les:indirect-object-pronouns:forms}

{\small\setlength{\LTpre}{2pt}\setlength{\LTpost}{6pt}
\tablecaption{Indirect object pronouns}
\begin{longtable}{lllll}
\toprule
\textbf{} & \textbf{Singular} & \textbf{} & \textbf{Plural} & \textbf{} \\
\midrule
1st & mi & to me & ci & to us \\
2nd & ti & to you & vi & to you (pl.) \\
3rd & gli / le & to him / to her & gli (loro) & to them \\
formal & Le & to you (formal) &  &  \\
\bottomrule
\end{longtable}}

\begin{notebox}{rulecol}\notelabel{rulecol}{Rule}{Replacing a + person}
An indirect object pronoun replaces \textit{a} + a person: \textit{Do un regalo a Marco → Gli do un regalo. Scrivo a Sara → Le scrivo.}
\end{notebox}

\begin{notebox}{tendcol}\notelabel{tendcol}{Tendency}{'To them'}
In everyday Italian \textit{gli} also means 'to them': \textit{Gli parlo domani}. The formal alternative is \textit{loro} after the verb: \textit{Parlo loro domani}.
\end{notebox}

\subsection{Verbs that take an indirect object}\label{les:indirect-object-pronouns:verbs}

\begin{itemize}
\item\relax \textit{dare, portare, regalare} (to give something to someone)
\item\relax \textit{dire, raccontare, scrivere, rispondere, telefonare} (to say / write / reply to someone)
\item\relax \textit{piacere, sembrare} (to please, to seem to someone)
\end{itemize}

\begin{notebox}{exccol}\notelabel{exccol}{Exception}{Different from English}
Some verbs take a direct object in Italian where English has 'to' or 'for': \textit{ascoltare la musica} (listen to music), \textit{aspettare il treno} (wait for the train), \textit{cercare la chiave} (look for the key). Others take an indirect object where English has none: \textit{telefonare a Sara → le telefono} (I phone her).
\end{notebox}

\subsection{piacere with all persons}\label{les:indirect-object-pronouns:piacere}

{\small\setlength{\LTpre}{2pt}\setlength{\LTpost}{6pt}
\tablecaption{Who likes it?}
\begin{longtable}{ll}
\toprule
\textbf{Italian} & \textbf{English} \\
\midrule
Mi piace il mare. & I like the sea. \\
Ti piace il mare? & Do you like the sea? \\
Gli piace il mare. & He likes the sea. \\
Le piace il mare. & She likes the sea. \\
Ci piacciono i gatti. & We like cats. \\
Vi piacciono i gatti? & Do you (pl.) like cats? \\
Gli piacciono i gatti. & They like cats. \\
\bottomrule
\end{longtable}}

\subsection{Combined pronouns}\label{les:indirect-object-pronouns:combined}

\begin{notebox}{rulecol}\notelabel{rulecol}{Rule}{Indirect + direct}
When an indirect pronoun and \textit{lo, la, li, le} or \textit{ne} occur together, the indirect one comes first and its -i becomes -e: \textit{mi → me, ti → te, ci → ce, vi → ve, si → se}. \textit{Mi dai il libro? → Me lo dai? Ti porto le mele → Te le porto.}
\end{notebox}

\begin{notebox}{rulecol}\notelabel{rulecol}{Rule}{glielo}
\textit{gli} and \textit{le} (to him, to her, to you formal) both become \textit{glie-} and are written as one word with the following pronoun: \textit{glielo, gliela, glieli, gliele, gliene}. \textit{Do il libro a Marco / a Sara → Glielo do.} Context shows who receives it.
\end{notebox}

{\small\setlength{\LTpre}{2pt}\setlength{\LTpost}{6pt}
\tablecaption{Combined forms}
\begin{longtable}{llllll}
\toprule
\textbf{} & \textbf{+ lo} & \textbf{+ la} & \textbf{+ li} & \textbf{+ le} & \textbf{+ ne} \\
\midrule
mi & me lo & me la & me li & me le & me ne \\
ti & te lo & te la & te li & te le & te ne \\
gli / le & glielo & gliela & glieli & gliele & gliene \\
ci & ce lo & ce la & ce li & ce le & ce ne \\
vi & ve lo & ve la & ve li & ve le & ve ne \\
si & se lo & se la & se li & se le & se ne \\
\bottomrule
\end{longtable}}

\begin{notebox}{rulecol}\notelabel{rulecol}{Rule}{Position and agreement}
The pair behaves like a single pronoun: before a finite verb, attached to an infinitive, gerund or affirmative imperative (\textit{Puoi darmelo? Dammelo!}). In compound tenses the participle agrees with the direct object pronoun: \textit{Le chiavi? Gliele ho date.}
\end{notebox}

\subsection*{Worked examples}
\begin{examples}
\item\relax \textit{Il nonno ci racconta una storia.}\newline Grandfather tells us a story.
\item\relax \textit{Gli do la palla.}\newline I give him the ball.
\item\relax \textit{Le scrivo una lettera.}\newline I'm writing her a letter.
\item\relax \textit{Ti piace questa canzone?}\newline Do you like this song?
\item\relax \textit{Al cane piace correre. Gli piace molto.}\newline The dog likes running. He likes it a lot.
\item\relax \textit{Mi dici la verità?}\newline Will you tell me the truth?
\item\relax \textit{Il libro? Te lo porto domani.}\newline The book? I'll bring it to you tomorrow.
\item\relax \textit{Le chiavi? Gliele ho già date.}\newline The keys? I've already given them to him / her.\newline {\color{muted}\small glie- + le; date agrees with le.}
\end{examples}

\subsection*{Exercises}\label{les:indirect-object-pronouns:exercises}
{\color{muted}\small Answers: page~\pageref{ans:indirect-object-pronouns}.}

\begin{exercises}
\item\label{ex:indirect-object-pronouns-1}\relax Do il libro a Sara. → \blank{} do il libro.\par\nopagebreak\begin{options}
\item\relax Lo
\item\relax Le
\item\relax Gli
\item\relax La
\end{options}
\item\label{ex:indirect-object-pronouns-2}\relax Scrivo a Marco. → \blank{} scrivo.\par\nopagebreak\begin{options}
\item\relax Le
\item\relax Lo
\item\relax Li
\item\relax Gli
\end{options}
\item\label{ex:indirect-object-pronouns-3}\relax Aspetto il treno. → \blank{} aspetto.\par\nopagebreak\begin{options}
\item\relax Ci
\item\relax Gli
\item\relax Le
\item\relax Lo
\end{options}
\item\label{ex:indirect-object-pronouns-4}\relax Complete with mi, ti, gli, le, ci.\par\smallskip\begingroup\setlength{\parindent}{0pt}Sara ama i gatti: \clozeblank{1} piacciono molto. Marco ama il mare: \clozeblank{2} piace nuotare. A noi \clozeblank{3} piace leggere.\par\endgroup
\item\label{ex:indirect-object-pronouns-5}\relax Translate: Grandfather tells us a story.\par\answerline
\item\label{ex:indirect-object-pronouns-6}\relax Translate: \textit{Gli piacciono i gatti.} (gli = to them)\par\answerline
\item\label{ex:indirect-object-pronouns-c1}\relax Replace both objects: \textit{Do il libro a Marco.} → \blank{}
\item\label{ex:indirect-object-pronouns-c2}\relax \textit{Mi presti la penna?} — Sì, \blank{} presto.\par\nopagebreak\begin{options}
\item\relax ti la
\item\relax la ti
\item\relax te lo
\item\relax te la
\end{options}
\end{exercises}

\begin{summarybox}
\begin{itemize}
\item\relax mi, ti, gli/le, ci, vi, gli replace a + person.
\item\relax gli = to him (and colloquially to them); le = to her; Le = to you (formal).
\item\relax ascoltare, aspettare, cercare take direct objects; telefonare, rispondere, piacere take indirect ones.
\item\relax Combined pronouns: me lo, te la, glielo, ce ne …; indirect first, -i → -e, gli/le → glie-.
\end{itemize}
\end{summarybox}

\relatedwords{\textit{dare} to give \textperiodcentered{} \textit{dire} to say; to tell \textperiodcentered{} \textit{raccontare} to tell (a story) \textperiodcentered{} \textit{scrivere} to write \textperiodcentered{} \textit{rispondere} to answer; to reply \textperiodcentered{} \textit{piacere} to please; to be liked \textperiodcentered{} \textit{aspettare} to wait (for) \textperiodcentered{} \textit{ascoltare} to listen to}
\applyit{Read \textit{La famiglia di Marco} (A0–A1), \S\ref{story:la-famiglia-di-marco}. Read \textit{L'orso che non poteva dormire} (A1–A2), \S\ref{story:lorso-che-non-poteva-dormire}.}

\section{Relative clauses, adverbs and expressions in stories}\label{les:relative-clauses-adverbs}
\lessonmeta{A2}{che and cui; adverbs in -mente; infinitive constructions; idioms you meet in children's stories.}
\begin{objectives}
\item\relax Join sentences with the relative pronouns che and cui.
\item\relax Form adverbs with -mente.
\item\relax Recognise infinitive constructions and common narrative expressions.
\end{objectives}

\subsection{Relative pronouns: che and cui}\label{les:relative-clauses-adverbs:relatives}

\begin{notebox}{rulecol}\notelabel{rulecol}{Rule}{che = who, which, that (subject or direct object)}
\textit{Il ragazzo che canta è mio fratello.} (subject) — \textit{Il libro che leggo è bello.} (object). \textit{che} is invariable and, unlike English 'that', can never be omitted.
\end{notebox}

\begin{notebox}{rulecol}\notelabel{rulecol}{Rule}{cui after a preposition}
After a preposition use \textit{cui}: \textit{la casa in cui abito} (the house I live in), \textit{l'amica con cui gioco} (the friend I play with), \textit{il ragazzo a cui scrivo} (the boy I'm writing to).
\end{notebox}

\begin{notebox}{tendcol}\notelabel{tendcol}{Tendency}{il quale}
\textit{il quale, la quale, i quali, le quali} can replace \textit{che} or \textit{cui}, and agree in gender and number with the antecedent: \textit{la ragazza con la quale parlavo} = \textit{la ragazza con cui parlavo}. They are typical of formal or written style, and useful to avoid ambiguity, because the article shows which noun is meant: \textit{la sorella di Marco, la quale vive a Roma} (the sister lives in Rome). After a preposition the article contracts: \textit{del quale, alla quale, nei quali}. As the subject or object of a restrictive relative clause, \textit{che} is normal and \textit{il quale} sounds heavy or bureaucratic.
\end{notebox}

\subsection{Adverbs in -mente}\label{les:relative-clauses-adverbs:adverbs}

\begin{notebox}{rulecol}\notelabel{rulecol}{Rule}{Feminine singular + -mente}
\textit{lento → lenta → lentamente}; \textit{veloce → velocemente}. Adjectives in \textit{-le} and \textit{-re} drop the final e: \textit{gentile → gentilmente, regolare → regolarmente}.
\end{notebox}

\begin{notebox}{exccol}\notelabel{exccol}{Exception}{Irregular adverbs}
\textit{buono → bene}, \textit{cattivo → male}. Many common adverbs do not end in -mente: \textit{presto, tardi, sempre, piano, insieme, subito}.
\end{notebox}

\subsection{Infinitive constructions and story expressions}\label{les:relative-clauses-adverbs:infinitives-idioms}

\begin{itemize}
\item\relax Verb + \textit{a} + infinitive: \textit{cominciare a, imparare a, aiutare a} — \textit{Comincia a piovere.}
\item\relax Verb + \textit{di} + infinitive: \textit{cercare di} (to try to), \textit{decidere di, finire di} — \textit{Cerca di dormire.}
\item\relax \textit{per} + infinitive = in order to: \textit{Esce per giocare.}
\item\relax \textit{prima di} + infinitive = before …-ing: \textit{prima di dormire}.
\end{itemize}

{\small\setlength{\LTpre}{2pt}\setlength{\LTpost}{6pt}
\tablecaption{Expressions common in children's stories}
\begin{xltabular}{\linewidth}{LL}
\toprule
\textbf{Italian} & \textbf{Meaning} \\
\midrule
C'era una volta… & Once upon a time there was… \\
all'improvviso & suddenly \\
a un certo punto & at a certain point \\
piano piano & little by little; very slowly \\
in fretta & in a hurry \\
e vissero felici e contenti & and they lived happily ever after (passato remoto of vivere, the standard tense of fairy-tale narration) \\
\bottomrule
\end{xltabular}}

\subsection*{Worked examples}
\begin{examples}
\item\relax \textit{Il ragazzo che canta è mio fratello.}\newline The boy who is singing is my brother.
\item\relax \textit{Questa è la casa in cui abito.}\newline This is the house I live in.
\item\relax \textit{La tartaruga cammina lentamente.}\newline The tortoise walks slowly.
\item\relax \textit{Comincia a piovere.}\newline It's starting to rain.
\item\relax \textit{L'orso cerca di dormire.}\newline The bear tries to sleep.
\item\relax \textit{All'improvviso è arrivato il vento.}\newline Suddenly the wind came.
\end{examples}

\subsection*{Exercises}\label{les:relative-clauses-adverbs:exercises}
{\color{muted}\small Answers: page~\pageref{ans:relative-clauses-adverbs}.}

\begin{exercises}
\item\label{ex:relative-clauses-adverbs-1}\relax Il libro \blank{} leggo è bello.\par\nopagebreak\begin{options}
\item\relax quale
\item\relax chi
\item\relax cui
\item\relax che
\end{options}
\item\label{ex:relative-clauses-adverbs-2}\relax L'amica con \blank{} gioco si chiama Sara.\par\nopagebreak\begin{options}
\item\relax che
\item\relax chi
\item\relax cui
\item\relax la quale che
\end{options}
\item\label{ex:relative-clauses-adverbs-3}\relax Make an adverb from \textit{veloce}: \blank{}
\item\label{ex:relative-clauses-adverbs-4}\relax Make an adverb from \textit{gentile}: \blank{}
\item\label{ex:relative-clauses-adverbs-5}\relax L'orso cerca \blank{} dormire.\par\nopagebreak\begin{options}
\item\relax —
\item\relax a
\item\relax per
\item\relax di
\end{options}
\item\label{ex:relative-clauses-adverbs-6}\relax Comincia \blank{} piovere.\par\nopagebreak\begin{options}
\item\relax di
\item\relax —
\item\relax da
\item\relax a
\end{options}
\item\label{ex:relative-clauses-adverbs-7}\relax Translate: The boy who is singing is my brother.\par\answerline
\end{exercises}

\begin{summarybox}
\begin{itemize}
\item\relax che (subject or object) and cui (after a preposition); che is never omitted.
\item\relax Adverbs: feminine adjective + -mente; -le/-re drop the e; bene, male irregular.
\item\relax cominciare a, imparare a; cercare di, finire di; per + infinitive.
\item\relax Story expressions: C'era una volta, all'improvviso, piano piano.
\end{itemize}
\end{summarybox}

\relatedwords{\textit{che} that; which; who; what \textperiodcentered{} \textit{lentamente} slowly \textperiodcentered{} \textit{veloce} fast; quick \textperiodcentered{} \textit{gentile} kind; polite \textperiodcentered{} \textit{cominciare} to begin; to start \textperiodcentered{} \textit{cercare} to look for; to try \textperiodcentered{} \textit{finalmente} at last; finally \textperiodcentered{} \textit{c'era una volta} once upon a time (there was)}
\applyit{Read \textit{L'orso che non poteva dormire} (A1–A2), \S\ref{story:lorso-che-non-poteva-dormire}. Read \textit{La tartaruga lenta} (A1), \S\ref{story:la-tartaruga-lenta}. Study the rhyme \textit{La pioggia sul tetto}, \S\ref{rhyme:la-pioggia-sul-tetto}.}

\chapter{Futuro}\label{mod:m13}
\moduleintro{A2–B1}{Futuro semplice and futuro anteriore: future time, anteriority in the future, and supposition about the present and the past.}{Futuro semplice (\S\ref{les:futuro-semplice}); Futuro anteriore (\S\ref{les:futuro-anteriore})}

\section{Futuro semplice}\label{les:futuro-semplice}
\lessonmeta{A2}{Formation of the simple future, its irregular stems, and its two main uses: future time and supposition about the present.}
\begin{objectives}
\item\relax Form the futuro semplice of regular verbs, including the spelling rules for -care/-gare and -ciare/-giare.
\item\relax Learn the irregular future stems (sarò, avrò, andrò, verrò, vorrò …).
\item\relax Use the futuro for future events and know when the present is used instead.
\item\relax Use the futuro to express a supposition about the present (sarà stanco = he must be tired).
\end{objectives}

\subsection{Formation}\label{les:futuro-semplice:formation}

\begin{notebox}{rulecol}\notelabel{rulecol}{Rule}{Stem + -ò, -ai, -à, -emo, -ete, -anno}
The future stem is the infinitive without its final -e, and -are verbs change -ar- to -er-: \textit{parlare → parler-}, \textit{credere → creder-}, \textit{dormire → dormir-}. The endings are the same for every verb.
\end{notebox}

\par\noindent\begin{minipage}[t]{0.48\linewidth}\raggedright
\paradigmbox{parlare --- to speak; to talk \textperiodcentered{} indicativo futuro semplice}{%
\begin{tabular}{@{}l l@{}}
\textit{io} & parler\hlending{ò} \\
\textit{tu} & parler\hlending{ai} \\
\textit{lui/lei} & parler\hlending{à} \\
\textit{noi} & parler\hlending{emo} \\
\textit{voi} & parler\hlending{ete} \\
\textit{loro} & parler\hlending{anno} \\
\end{tabular}}{Stem \textit{parler-} + \hlending{ending}.}
\end{minipage}\par\medskip

\begin{notebox}{rulecol}\notelabel{rulecol}{Rule}{Spelling adjustments}
\textit{-care/-gare} add h to keep the hard sound: \textit{cercherò, pagherò}. \textit{-ciare/-giare} drop the i, which is no longer needed to soften c/g before e: \textit{comincerò, mangerò}. Other -iare verbs keep it: \textit{studierò}.
\end{notebox}

\begin{notebox}{rulecol}\notelabel{rulecol}{Rule}{Final stress is written}
The io and lui/lei forms are stressed on the last syllable and carry a written accent: \textit{parlerò, parlerà}. \textit{parlero} without the accent is a spelling error.
\end{notebox}

\subsection{Irregular stems}\label{les:futuro-semplice:irregular}

Some verbs have a shortened or changed stem; the endings stay regular. The same stem is used for the \lessonref{condizionale-presente}{condizionale presente}.\par

{\small\setlength{\LTpre}{2pt}\setlength{\LTpost}{6pt}
\tablecaption{Irregular future stems}
\begin{xltabular}{\linewidth}{LLL}
\toprule
\textbf{Pattern} & \textbf{Verbs} & \textbf{Futuro (io)} \\
\midrule
loss of the vowel & andare, avere, cadere, dovere, potere, sapere, vedere, vivere & andrò, avrò, cadrò, dovrò, potrò, saprò, vedrò, vivrò \\
-rr- & bere, rimanere, tenere, venire, volere & berrò, rimarrò, terrò, verrò, vorrò \\
short -are verbs keep -a- & dare, fare, stare & darò, farò, starò \\
unique & essere & sarò \\
\bottomrule
\end{xltabular}}

\par\noindent\begin{minipage}[t]{0.48\linewidth}\raggedright
\paradigmbox{essere --- to be \textperiodcentered{} indicativo futuro semplice}{%
\begin{tabular}{@{}l l@{}}
\textit{io} & \irr{sarò} \\
\textit{tu} & \irr{sarai} \\
\textit{lui/lei} & \irr{sarà} \\
\textit{noi} & \irr{saremo} \\
\textit{voi} & \irr{sarete} \\
\textit{loro} & \irr{saranno} \\
\end{tabular}}{\irr{Highlighted} forms differ from the regular pattern of the verb's family.}
\end{minipage}\par\medskip

\begin{notebox}{tendcol}\notelabel{tendcol}{Tendency}{Alternative forms}
A few verbs have two accepted futures: \textit{morirò / morrò} (the first is far more common today), \textit{sederò / siederò}.
\end{notebox}

\subsection{Uses}\label{les:futuro-semplice:uses}

\begin{itemize}
\item\relax Future events: \textit{Domani pioverà.} (Tomorrow it will rain.)
\item\relax Promises and predictions: \textit{Ti scriverò ogni settimana.}
\item\relax Supposition about the present (futuro epistemico): \textit{Dov'è Marco? — Sarà a casa.} (He must be / is probably at home.) \textit{Che ore sono? — Saranno le otto.} (It must be about eight.)
\item\relax Concession: \textit{Sarà anche bravo, ma è antipatico.} (He may well be good, but he's unpleasant.)
\item\relax Instructions and orders, especially in formal or firm style: \textit{Farai i compiti prima di uscire.}
\end{itemize}

\begin{notebox}{tendcol}\notelabel{tendcol}{Tendency}{The present for the future}
For planned or near events Italian very often uses the present: \textit{Domani parto per Roma} is more natural in conversation than \textit{Domani partirò per Roma}. The future stresses prediction, uncertainty or distance in time. Both are correct; the choice is stylistic and pragmatic, not grammatical.
\end{notebox}

\begin{notebox}{rulecol}\notelabel{rulecol}{Rule}{After quando, appena, se with future reference}
When the main verb is future, a temporal clause about the future uses the futuro, unlike English: \textit{Quando arriverai, ti chiamerò} (When you arrive, I'll call you). In conversation the present is also common: \textit{Quando arrivi, ti chiamo}. After se both are found: \textit{Se verrai / Se vieni, sarò contento}.
\end{notebox}

\subsection*{Worked examples}
\begin{examples}
\item\relax \textit{Domani pioverà tutto il giorno.}\newline Tomorrow it will rain all day.
\item\relax \textit{Quando arriverai, ti chiamerò.}\newline When you arrive, I'll call you.\newline {\color{muted}\small Futuro in both clauses; English uses the present after 'when'.}
\item\relax \textit{Dov'è la nonna? — Sarà in giardino.}\newline Where's Grandma? — She must be in the garden.\newline {\color{muted}\small Futuro of supposition: a guess about the present.}
\item\relax \textit{Che ore sono? — Saranno le otto.}\newline What time is it? — It must be about eight.
\item\relax \textit{L'anno prossimo andremo in Sicilia.}\newline Next year we'll go to Sicily.\newline {\color{muted}\small andare → andr-.}
\item\relax \textit{Domani parto per Roma.}\newline I'm leaving for Rome tomorrow.\newline {\color{muted}\small The present for a planned near-future event.}
\item\relax \textit{Sarà anche bravo, ma non mi piace.}\newline He may well be good, but I don't like him.\newline {\color{muted}\small Concessive futuro.}
\end{examples}

\subsection*{Exercises}\label{les:futuro-semplice:exercises}
{\color{muted}\small Answers: page~\pageref{ans:futuro-semplice}.}

\begin{exercises}
\item\label{ex:futuro-semplice-1}\relax Conjugate \textit{cercare} in the futuro semplice.\par\smallskip{\begin{tabular}{@{}l l@{}}
io & \blankcell \\
tu & \blankcell \\
lui/lei & \blankcell \\
noi & \blankcell \\
voi & \blankcell \\
loro & \blankcell \\
\end{tabular}}
\item\label{ex:futuro-semplice-2}\relax mangiare, io, futuro: \blank{}
\item\label{ex:futuro-semplice-3}\relax Give the io form of the futuro.\par\smallskip{\begin{tabular}{@{}l l@{}}
essere & \blankcell \\
venire & \blankcell \\
andare & \blankcell \\
volere & \blankcell \\
fare & \blankcell \\
vedere & \blankcell \\
\end{tabular}}
\item\label{ex:futuro-semplice-4}\relax \textit{Non risponde al telefono. Sarà in palestra.} What does \textit{sarà} express here?\par\nopagebreak\begin{options}
\item\relax a past event
\item\relax a future event
\item\relax an order
\item\relax a supposition about the present
\end{options}
\item\label{ex:futuro-semplice-5}\relax Which sentence uses the futuro to express a supposition about the present rather than a future event?\par\nopagebreak\begin{options}
\item\relax L'anno prossimo saremo in Italia.
\item\relax Alle otto sarò a casa.
\item\relax A quest'ora sarà in ufficio.
\item\relax Domani sarò in ufficio.
\end{options}
\item\label{ex:futuro-semplice-6}\relax Complete with the futuro.\par\smallskip\begingroup\setlength{\parindent}{0pt}L'anno prossimo noi \clozeblank{1} (andare) al mare e voi \clozeblank{2} (rimanere) in città. I bambini \clozeblank{3} (essere) contenti.\par\endgroup
\item\label{ex:futuro-semplice-7}\relax Translate (supposition): It must be about eight o'clock.\par\answerline
\end{exercises}

\begin{summarybox}
\begin{itemize}
\item\relax Futuro = future stem + -ò, -ai, -à, -emo, -ete, -anno; -are → -er- (parlerò).
\item\relax cercherò, pagherò (h); mangerò, comincerò (no i).
\item\relax Irregular stems: sar-, avr-, andr-, verr-, vorr-, berr-, rimarr-, terr-, potr-, dovr-, sapr-, vedr-, vivr-, far-, dar-, star-.
\item\relax Uses: future, supposition about the present, concession; the present is very common for near or planned events.
\end{itemize}
\end{summarybox}

\relatedwords{\textit{essere} to be \textperiodcentered{} \textit{avere} to have \textperiodcentered{} \textit{andare} to go \textperiodcentered{} \textit{venire} to come \textperiodcentered{} \textit{volere} to want \textperiodcentered{} \textit{fare} to do; to make \textperiodcentered{} \textit{vedere} to see \textperiodcentered{} \textit{cercare} to look for; to try \textperiodcentered{} \textit{mangiare} to eat \textperiodcentered{} \textit{piovere} to rain}
\applyit{Read \textit{Una lettera dalla nonna} (A2–B1), \S\ref{story:una-lettera-dalla-nonna}.}

\section{Futuro anteriore}\label{les:futuro-anteriore}
\lessonmeta{B1}{The future perfect: an action completed before a future moment, and supposition about the past.}
\begin{objectives}
\item\relax Form the futuro anteriore with the future of avere or essere + past participle.
\item\relax Express anteriority in the future: dopo che avrai finito, …
\item\relax Express a supposition about a past event: avrà perso il treno.
\end{objectives}

\subsection{Formation}\label{les:futuro-anteriore:formation}

\begin{notebox}{rulecol}\notelabel{rulecol}{Rule}{Futuro of the auxiliary + participle}
\textit{avrò mangiato, sarò partito/a}. The choice of auxiliary and the agreement of the participle are exactly as in the passato prossimo.
\end{notebox}

\par\noindent\begin{minipage}[t]{0.48\linewidth}\raggedright
\paradigmbox{mangiare --- to eat \textperiodcentered{} indicativo futuro anteriore (with avere)}{%
\begin{tabular}{@{}l l@{}}
\textit{io} & avrò mangiato \\
\textit{tu} & avrai mangiato \\
\textit{lui/lei} & avrà mangiato \\
\textit{noi} & avremo mangiato \\
\textit{voi} & avrete mangiato \\
\textit{loro} & avranno mangiato \\
\end{tabular}}{}
\end{minipage}\hfill\begin{minipage}[t]{0.48\linewidth}\raggedright
\paradigmbox{partire --- to leave; to set off \textperiodcentered{} indicativo futuro anteriore (with essere, feminine subject)}{%
\begin{tabular}{@{}l l@{}}
\textit{io} & sarò partita \\
\textit{tu} & sarai partita \\
\textit{lui/lei} & sarà partita \\
\textit{noi} & saremo partite \\
\textit{voi} & sarete partite \\
\textit{loro} & saranno partite \\
\end{tabular}}{}
\end{minipage}\par\medskip

\subsection{Uses}\label{les:futuro-anteriore:uses}

\begin{notebox}{rulecol}\notelabel{rulecol}{Rule}{Anteriority in the future}
It presents an action as completed before another future action or moment, typically after \textit{quando, dopo che, appena, se}: \textit{Quando avrai finito, usciremo.} In everyday speech the passato prossimo is also used in this context: \textit{Quando hai finito, usciamo.}
\end{notebox}

\begin{notebox}{rulecol}\notelabel{rulecol}{Rule}{Supposition about the past}
Just as the futuro semplice expresses a guess about the present, the futuro anteriore expresses a guess about the past: \textit{Non è ancora arrivato: avrà perso il treno} (he must have missed the train).
\end{notebox}

\begin{notebox}{tendcol}\notelabel{tendcol}{Tendency}{By a future moment}
\textit{Entro domani avrò finito il libro} (By tomorrow I will have finished the book). English 'will have done' corresponds closely here.
\end{notebox}

\subsection*{Worked examples}
\begin{examples}
\item\relax \textit{Quando avrai finito i compiti, potrai uscire.}\newline When you have finished your homework, you can go out.\newline {\color{muted}\small English uses the present perfect; Italian uses the futuro anteriore (or, informally, the passato prossimo).}
\item\relax \textit{Appena sarò arrivata, ti scriverò.}\newline As soon as I have arrived, I'll write to you.\newline {\color{muted}\small Female speaker: arrivata.}
\item\relax \textit{Marco non risponde: sarà già uscito.}\newline Marco isn't answering: he must have gone out already.\newline {\color{muted}\small Supposition about the past.}
\item\relax \textit{Entro giugno avremo letto tutti i libri.}\newline By June we will have read all the books.
\item\relax \textit{Chi avrà rotto il vaso?}\newline Who can have broken the vase? / I wonder who broke the vase.\newline {\color{muted}\small A question expressing wonder or doubt.}
\item\relax \textit{Dopo che avranno mangiato, i bambini andranno a letto.}\newline After they have eaten, the children will go to bed.
\end{examples}

\subsection*{Exercises}\label{les:futuro-anteriore:exercises}
{\color{muted}\small Answers: page~\pageref{ans:futuro-anteriore}.}

\begin{exercises}
\item\label{ex:futuro-anteriore-1}\relax finire, tu, futuro anteriore: \blank{}
\item\label{ex:futuro-anteriore-2}\relax arrivare, noi (masculine), futuro anteriore: \blank{}
\item\label{ex:futuro-anteriore-3}\relax \textit{Il treno è in ritardo: ci sarà stato un guasto.} What does \textit{ci sarà stato} express?\par\nopagebreak\begin{options}
\item\relax a wish
\item\relax a future event
\item\relax a certainty about the past
\item\relax a guess about the past
\end{options}
\item\label{ex:futuro-anteriore-4}\relax Quando \blank{} (finire, voi), andremo al parco.\par\nopagebreak\begin{options}
\item\relax avevate finito
\item\relax avrete finito
\item\relax finivate
\item\relax aveste finito
\end{options}
\item\label{ex:futuro-anteriore-5}\relax Translate: He must have missed the train.\par\answerline
\item\label{ex:futuro-anteriore-6}\relax Complete with the futuro anteriore.\par\smallskip\begingroup\setlength{\parindent}{0pt}Entro stasera io \clozeblank{1} (scrivere) la lettera. Sara non è qui: \clozeblank{2} (andare) a casa.\par\endgroup
\end{exercises}

\begin{summarybox}
\begin{itemize}
\item\relax Futuro anteriore = futuro of avere/essere + participle (avrò fatto, sarò partito/a).
\item\relax Anteriority in the future: quando avrai finito, …
\item\relax Supposition about the past: avrà perso il treno.
\item\relax In conversation the passato prossimo often replaces it after quando/appena.
\end{itemize}
\end{summarybox}

\relatedwords{\textit{finire} to finish; to end \textperiodcentered{} \textit{arrivare} to arrive \textperiodcentered{} \textit{perdere} to lose; to miss \textperiodcentered{} \textit{scrivere} to write \textperiodcentered{} \textit{rompere} to break \textperiodcentered{} \textit{uscire} to go out; to leave}
\applyit{Read \textit{Una lettera dalla nonna} (A2–B1), \S\ref{story:una-lettera-dalla-nonna}.}

\chapter{Condizionale}\label{mod:m14}
\moduleintro{A2–B1}{Condizionale presente and passato: politeness, wishes, advice, hypothetical consequences, reported information and the future in the past.}{Condizionale presente (\S\ref{les:condizionale-presente}); Condizionale passato and the future in the past (\S\ref{les:condizionale-passato})}

\section{Condizionale presente}\label{les:condizionale-presente}
\lessonmeta{A2}{Formation of the present conditional and its uses: politeness, wishes, advice, hypothetical consequences and reported information.}
\begin{objectives}
\item\relax Form the condizionale presente from the future stem.
\item\relax Make polite requests (vorrei, potrebbe…) and give advice (dovresti).
\item\relax Use the conditional for the consequence of a hypothetical condition.
\item\relax Recognise the 'journalistic' conditional for unconfirmed information.
\end{objectives}

\subsection{Formation}\label{les:condizionale-presente:formation}

\begin{notebox}{rulecol}\notelabel{rulecol}{Rule}{Future stem + -ei, -esti, -ebbe, -emmo, -este, -ebbero}
The conditional uses exactly the same stem as the futuro, with its own endings: \textit{parler-ei, sar-ei, vorr-ei, andr-ei}. Every irregularity of the future therefore reappears here.
\end{notebox}

\par\noindent\begin{minipage}[t]{0.48\linewidth}\raggedright
\paradigmbox{parlare --- to speak; to talk \textperiodcentered{} condizionale presente}{%
\begin{tabular}{@{}l l@{}}
\textit{io} & parler\hlending{ei} \\
\textit{tu} & parler\hlending{esti} \\
\textit{lui/lei} & parler\hlending{ebbe} \\
\textit{noi} & parler\hlending{emmo} \\
\textit{voi} & parler\hlending{este} \\
\textit{loro} & parler\hlending{ebbero} \\
\end{tabular}}{Stem \textit{parler-} + \hlending{ending}.}
\end{minipage}\hfill\begin{minipage}[t]{0.48\linewidth}\raggedright
\paradigmbox{volere --- to want \textperiodcentered{} condizionale presente}{%
\begin{tabular}{@{}l l@{}}
\textit{io} & \irr{vorrei} \\
\textit{tu} & \irr{vorresti} \\
\textit{lui/lei} & \irr{vorrebbe} \\
\textit{noi} & \irr{vorremmo} \\
\textit{voi} & \irr{vorreste} \\
\textit{loro} & \irr{vorrebbero} \\
\end{tabular}}{\irr{Highlighted} forms differ from the regular pattern of the verb's family.}
\end{minipage}\par\medskip

\begin{notebox}{exccol}\notelabel{exccol}{Exception}{noi: -emmo, not -emo}
The noi forms of futuro and condizionale differ only by a double m: \textit{parleremo} (we will speak) / \textit{parleremmo} (we would speak). The difference is audible: the mm is pronounced long.
\end{notebox}

\subsection{Uses}\label{les:condizionale-presente:uses}

\begin{itemize}
\item\relax Politeness and attenuation: \textit{Vorrei un caffè. Potrebbe aiutarmi? Mi daresti una mano?}
\item\relax Wishes: \textit{Mi piacerebbe vivere al mare.}
\item\relax Advice and obligation softened: \textit{Dovresti dormire di più.} (You ought to sleep more.)
\item\relax Consequence of a hypothesis: \textit{Se avessi tempo, leggerei di più.} (see the lesson on hypothetical sentences)
\item\relax Opinion stated cautiously: \textit{Direi che hai ragione.}
\item\relax Unconfirmed or reported information, typical of journalism: \textit{Secondo alcune fonti, il ministro sarebbe a Roma.} (The minister is reportedly in Rome.)
\end{itemize}

\begin{notebox}{rulecol}\notelabel{rulecol}{Rule}{Never after se in a hypothesis}
In the se-clause of a hypothetical sentence the standard language uses the congiuntivo, not the conditional: \textit{Se potessi, verrei} — not \textit{se potrei}. The conditional in the se-clause is heard in nonstandard speech.
\end{notebox}

\subsection*{Worked examples}
\begin{examples}
\item\relax \textit{Vorrei un bicchiere d'acqua, per favore.}\newline I'd like a glass of water, please.\newline {\color{muted}\small Polite request.}
\item\relax \textit{Potresti chiudere la finestra?}\newline Could you close the window?
\item\relax \textit{Dovresti dormire di più.}\newline You ought to sleep more.\newline {\color{muted}\small Softened obligation = advice.}
\item\relax \textit{Mi piacerebbe vivere in campagna.}\newline I'd like to live in the country.
\item\relax \textit{Se avessi una barca, andrei in Sardegna.}\newline If I had a boat, I'd go to Sardinia.\newline {\color{muted}\small Congiuntivo imperfetto in the se-clause, conditional in the main clause.}
\item\relax \textit{Secondo il giornale, il ladro sarebbe scappato in Francia.}\newline According to the newspaper, the thief reportedly escaped to France.\newline {\color{muted}\small Conditional of reported information (here condizionale passato).}
\item\relax \textit{Noi parleremmo volentieri con lui.}\newline We would gladly speak with him.\newline {\color{muted}\small parleremmo (conditional) ≠ parleremo (future).}
\end{examples}

\subsection*{Exercises}\label{les:condizionale-presente:exercises}
{\color{muted}\small Answers: page~\pageref{ans:condizionale-presente}.}

\begin{exercises}
\item\label{ex:condizionale-presente-1}\relax Conjugate \textit{potere} in the condizionale presente.\par\smallskip{\begin{tabular}{@{}l l@{}}
io & \blankcell \\
tu & \blankcell \\
lui/lei & \blankcell \\
noi & \blankcell \\
voi & \blankcell \\
loro & \blankcell \\
\end{tabular}}
\item\label{ex:condizionale-presente-2}\relax 'We would go' = noi \blank{}\par\nopagebreak\begin{options}
\item\relax andassimo
\item\relax andremmo
\item\relax andremo
\item\relax andavamo
\end{options}
\item\label{ex:condizionale-presente-3}\relax What does \textit{Il presidente sarebbe malato} mean in a newspaper?\par\nopagebreak\begin{options}
\item\relax The president would be ill.
\item\relax The president is reportedly ill (not confirmed).
\item\relax The president was ill.
\item\relax The president will be ill.
\end{options}
\item\label{ex:condizionale-presente-4}\relax Se \blank{} (potere, io), verrei con te.\par\nopagebreak\begin{options}
\item\relax potrei
\item\relax posso
\item\relax potessi
\item\relax potevo
\end{options}
\item\label{ex:condizionale-presente-5}\relax Translate: You ought to sleep more. (tu)\par\answerline
\item\label{ex:condizionale-presente-6}\relax Complete with the condizionale presente.\par\smallskip\begingroup\setlength{\parindent}{0pt}Mi \clozeblank{1} (piacere) un gelato. \clozeblank{2} (Volere, io) anche un caffè. Voi \clozeblank{3} (venire) con noi?\par\endgroup
\end{exercises}

\begin{summarybox}
\begin{itemize}
\item\relax Condizionale presente = future stem + -ei, -esti, -ebbe, -emmo, -este, -ebbero.
\item\relax Same irregular stems as the futuro: sarei, avrei, vorrei, andrei, verrei, potrei, dovrei.
\item\relax Uses: politeness, wishes, advice, hypothetical consequences, unconfirmed information.
\item\relax Not after se in standard Italian: se potessi, verrei.
\end{itemize}
\end{summarybox}

\relatedwords{\textit{volere} to want \textperiodcentered{} \textit{potere} to be able to; can; may \textperiodcentered{} \textit{dovere} to have to; must \textperiodcentered{} \textit{piacere} to please; to be liked \textperiodcentered{} \textit{andare} to go \textperiodcentered{} \textit{venire} to come}
\applyit{Read \textit{Una lettera dalla nonna} (A2–B1), \S\ref{story:una-lettera-dalla-nonna}.}

\section{Condizionale passato and the future in the past}\label{les:condizionale-passato}
\lessonmeta{B1}{Unrealised past possibilities, and the condizionale passato as the standard 'future in the past'.}
\begin{objectives}
\item\relax Form the condizionale passato.
\item\relax Express what would have happened, or should have happened, but did not.
\item\relax Use the condizionale passato for a future seen from the past (disse che sarebbe venuto).
\item\relax Avoid transferring the English/French/Spanish 'would + verb' pattern into reported speech.
\end{objectives}

\subsection{Formation}\label{les:condizionale-passato:formation}

\begin{notebox}{rulecol}\notelabel{rulecol}{Rule}{Conditional of the auxiliary + participle}
\textit{avrei fatto, sarei andato/a}. Auxiliary choice and agreement as in the passato prossimo.
\end{notebox}

\par\noindent\begin{minipage}[t]{0.48\linewidth}\raggedright
\paradigmbox{venire --- to come \textperiodcentered{} condizionale passato (with essere, masculine subject)}{%
\begin{tabular}{@{}l l@{}}
\textit{io} & sarei venuto \\
\textit{tu} & saresti venuto \\
\textit{lui/lei} & sarebbe venuto \\
\textit{noi} & saremmo venuti \\
\textit{voi} & sareste venuti \\
\textit{loro} & sarebbero venuti \\
\end{tabular}}{}
\end{minipage}\par\medskip

\subsection{Unrealised possibilities}\label{les:condizionale-passato:unrealised}

\begin{itemize}
\item\relax What would have happened: \textit{Sarei venuto, ma ero malato.} (I would have come, but I was ill.)
\item\relax Regret and reproach: \textit{Avresti dovuto dirmelo.} (You should have told me.)
\item\relax Polite past wishes: \textit{Avrei voluto un tè.}
\item\relax Consequence of a counterfactual condition: \textit{Se fossi partito prima, sarei arrivato in tempo.}
\end{itemize}

\begin{notebox}{tendcol}\notelabel{tendcol}{Tendency}{Modal verbs}
\textit{avrei dovuto, avrei potuto, avrei voluto} + infinitive are extremely common: 'I should have / could have / would have liked to'. With an essere-verb infinitive, both \textit{avrei dovuto partire} and \textit{sarei dovuto partire} are standard (Accademia della Crusca); avere is increasingly frequent.
\end{notebox}

\subsection{The future in the past}\label{les:condizionale-passato:future-in-past}

\begin{notebox}{rulecol}\notelabel{rulecol}{Rule}{condizionale passato, not presente}
To report, from a past point of view, an event that was then in the future, standard Italian uses the \textbf{condizionale passato}: \textit{Marco ha detto che sarebbe venuto} (Marco said he would come) — even if the event may still happen. English 'would come', French \textit{viendrait} and Spanish \textit{vendría} use a simple conditional; Italian \textit{verrebbe} in this sense is a common learner error.
\end{notebox}

{\small\setlength{\LTpre}{2pt}\setlength{\LTpost}{6pt}
\tablecaption{From direct to reported speech}
\begin{xltabular}{\linewidth}{lL}
\toprule
\textbf{Direct speech} & \textbf{Reported after a past verb} \\
\midrule
«Verrò domani.» & Ha detto che sarebbe venuto il giorno dopo. \\
«Ti scriverò.» & Mi ha promesso che mi avrebbe scritto. \\
«Pioverà.» & Pensavo che sarebbe piovuto. \\
\bottomrule
\end{xltabular}}

\subsection*{Worked examples}
\begin{examples}
\item\relax \textit{Sarei venuta volentieri, ma dovevo lavorare.}\newline I would gladly have come, but I had to work.\newline {\color{muted}\small Female speaker: venuta.}
\item\relax \textit{Avresti dovuto dirmelo prima.}\newline You should have told me earlier.
\item\relax \textit{Marco ha detto che sarebbe arrivato alle otto.}\newline Marco said he would arrive at eight.\newline {\color{muted}\small Future in the past = condizionale passato.}
\item\relax \textit{Ero sicuro che avrebbe vinto la tartaruga.}\newline I was sure the tortoise would win.\newline {\color{muted}\small Subject after the verb: la tartaruga.}
\item\relax \textit{Se avessi saputo, ti avrei aiutato.}\newline If I had known, I would have helped you.
\item\relax \textit{Avrei voluto vedere il mare.}\newline I would have liked to see the sea.
\end{examples}

\subsection*{Exercises}\label{les:condizionale-passato:exercises}
{\color{muted}\small Answers: page~\pageref{ans:condizionale-passato}.}

\begin{exercises}
\item\label{ex:condizionale-passato-1}\relax andare, io (male speaker), condizionale passato: \blank{}
\item\label{ex:condizionale-passato-2}\relax Sara ha detto che \blank{} (venire) alla festa.\par\nopagebreak\begin{options}
\item\relax verrebbe
\item\relax sarebbe venuta
\item\relax verrà
\item\relax veniva
\end{options}
\item\label{ex:condizionale-passato-3}\relax Translate: You should have told me. (tu)\par\answerline
\item\label{ex:condizionale-passato-4}\relax Report in the past.\par\smallskip\begingroup\setlength{\parindent}{0pt}«Ti scriverò.» → Mi ha promesso che mi \clozeblank{1}. «Pioverà.» → Pensavo che \clozeblank{2}.\par\endgroup
\item\label{ex:condizionale-passato-5}\relax \textit{Sarei venuto, ma ero malato.} Did the speaker come?\par\nopagebreak\begin{options}
\item\relax Yes
\item\relax He will come later
\item\relax No
\item\relax We cannot know
\end{options}
\item\label{ex:condizionale-passato-6}\relax Match each sentence with its function.\par\smallskip\begin{tabularx}{\linewidth}{@{}L L@{}}
1.\ Avrei voluto un tè. & (a)\ unconfirmed report \\
2.\ Avresti dovuto studiare. & (b)\ reproach \\
3.\ Disse che sarebbe tornato. & (c)\ future in the past \\
4.\ Il ladro sarebbe fuggito. & (d)\ polite past wish \\
\end{tabularx}
\end{exercises}

\begin{summarybox}
\begin{itemize}
\item\relax Condizionale passato = conditional of avere/essere + participle.
\item\relax Unrealised past possibilities, regrets, reproaches (avresti dovuto…).
\item\relax The future in the past is the condizionale passato: disse che sarebbe venuto.
\item\relax dovere/potere/volere + essere-verb infinitive: both avrei dovuto partire and sarei dovuto partire are standard.
\end{itemize}
\end{summarybox}

\relatedwords{\textit{venire} to come \textperiodcentered{} \textit{dovere} to have to; must \textperiodcentered{} \textit{volere} to want \textperiodcentered{} \textit{dire} to say; to tell \textperiodcentered{} \textit{promettere} to promise \textperiodcentered{} \textit{vincere} to win}
\applyit{Read \textit{Una lettera dalla nonna} (A2–B1), \S\ref{story:una-lettera-dalla-nonna}. Read \textit{Il segreto del faro} (A2–B1), \S\ref{story:il-segreto-del-faro}.}

\chapter{Infinitive, gerund and participle}\label{mod:m16}
\moduleintro{A2–B1}{The non-finite forms: infinito presente and passato, gerundio presente and passato, stare + gerundio, participio passato and the participio assoluto.}{Infinitive and gerund (\S\ref{les:infinito-gerundio}); The participle (\S\ref{les:participio})}

\section{Infinitive and gerund}\label{les:infinito-gerundio}
\lessonmeta{B1}{Present and past infinitive; present and past gerund; stare + gerundio; what the gerund can and cannot do in Italian.}
\begin{objectives}
\item\relax Use the infinito presente and infinito passato (aver mangiato, essere partito).
\item\relax Know which prepositions introduce an infinitive, and that dopo requires the past infinitive.
\item\relax Form the gerundio presente and passato and use stare + gerundio.
\item\relax Avoid using the gerund where Italian requires an infinitive (after prepositions, as a subject).
\end{objectives}

\subsection{The infinitive}\label{les:infinito-gerundio:infinitive}

\begin{notebox}{rulecol}\notelabel{rulecol}{Rule}{Present and past infinitive}
The past infinitive is \textit{avere / essere} + participle: \textit{aver mangiato, essere partito/a/i/e}. It expresses an action completed before the main verb. The final -e of \textit{avere/essere} is often dropped before a consonant (\textit{aver detto, esser stato}); both forms are correct.
\end{notebox}

\begin{itemize}
\item\relax After modal verbs and verbs like \textit{sapere, preferire, lasciare, fare}: \textit{Voglio partire. Lascia parlare la nonna.}
\item\relax After prepositions: \textit{prima di partire, invece di studiare, senza parlare, per imparare}.
\item\relax \textit{dopo} requires the \textbf{past} infinitive: \textit{dopo aver mangiato} (not \textit{dopo mangiare} in standard Italian).
\item\relax As a noun (subject or object): \textit{Leggere è bello. Il mangiare.}
\item\relax In instructions and signs: \textit{Non fumare. Mescolare bene.}
\item\relax Negative tu imperative: \textit{Non toccare!}
\end{itemize}

\begin{notebox}{rulecol}\notelabel{rulecol}{Rule}{Pronouns attach to the infinitive}
The infinitive loses its -e and the pronoun attaches: \textit{vederlo, dirtelo, alzarsi}. With modals the pronoun may also precede the modal (\textit{lo voglio vedere}).
\end{notebox}

\subsection{The gerund}\label{les:infinito-gerundio:gerund}

\begin{notebox}{rulecol}\notelabel{rulecol}{Rule}{Formation}
-are → -ando, -ere/-ire → -endo: \textit{parlando, vedendo, dormendo}. Verbs with an extended imperfetto stem use it: \textit{facendo, dicendo, bevendo}. Gerundio passato: \textit{avendo mangiato, essendo partito/a}.
\end{notebox}

\begin{itemize}
\item\relax Manner or means: \textit{Ho imparato l'italiano leggendo.} (by reading)
\item\relax Time: \textit{Tornando a casa, ho visto un gufo.} (while coming home)
\item\relax Cause: \textit{Essendo stanco, è andato a letto.} (since he was tired)
\item\relax Condition: \textit{Volendo, potresti finire oggi.} (if you wanted to)
\item\relax Progressive with \textit{stare}: \textit{Sto leggendo. Stavamo cenando quando è arrivato.}
\end{itemize}

\begin{notebox}{rulecol}\notelabel{rulecol}{Rule}{The gerund's subject}
The gerund normally has the same subject as the main verb. \textit{Ho visto Marco uscendo dal negozio} means that \textit{I} was coming out of the shop; to say Marco was, use a relative clause: \textit{Ho visto Marco che usciva dal negozio.}
\end{notebox}

\begin{notebox}{exccol}\notelabel{exccol}{Exception}{Not like the English -ing}
Italian never uses the gerund after a preposition or as a noun: 'before leaving' = \textit{prima di partire}; 'I like swimming' = \textit{mi piace nuotare}; 'Swimming is healthy' = \textit{Nuotare fa bene}.
\end{notebox}

\begin{notebox}{tendcol}\notelabel{tendcol}{Tendency}{stare + gerundio is optional}
The simple present already means 'I am doing': \textit{Che cosa fai?} = What are you doing? \textit{stare + gerundio} stresses that the action is in progress right now. It is not used with \textit{essere, avere} or \textit{stare} itself, nor in compound tenses.
\end{notebox}

\subsection*{Worked examples}
\begin{examples}
\item\relax \textit{Dopo aver mangiato, il nonno dormì un'ora.}\newline After eating, Grandfather slept for an hour.\newline {\color{muted}\small dopo + past infinitive.}
\item\relax \textit{Prima di uscire, chiudi la finestra.}\newline Before going out, close the window.
\item\relax \textit{Ho imparato l'italiano leggendo storie.}\newline I learned Italian by reading stories.\newline {\color{muted}\small Gerund of means.}
\item\relax \textit{Stavamo cenando quando è squillato il telefono.}\newline We were having dinner when the phone rang.\newline {\color{muted}\small stare (imperfetto) + gerundio.}
\item\relax \textit{Essendo arrivati tardi, non abbiamo trovato posto.}\newline Having arrived late, we didn't find a seat.\newline {\color{muted}\small Gerundio passato with essere: agreement with the subject.}
\item\relax \textit{Nuotare fa bene alla salute.}\newline Swimming is good for your health.\newline {\color{muted}\small Infinitive as subject (not the gerund).}
\end{examples}

\subsection*{Exercises}\label{les:infinito-gerundio:exercises}
{\color{muted}\small Answers: page~\pageref{ans:infinito-gerundio}.}

\begin{exercises}
\item\label{ex:infinito-gerundio-1}\relax \blank{} i compiti, Luca è uscito.\par\nopagebreak\begin{options}
\item\relax Dopo finendo
\item\relax Dopo finito
\item\relax Dopo finire
\item\relax Dopo aver finito
\end{options}
\item\label{ex:infinito-gerundio-2}\relax 'I like swimming.'\par\nopagebreak\begin{options}
\item\relax Mi piace nuotare.
\item\relax Mi piace nuotando.
\item\relax Mi piacciono nuotare.
\item\relax Mi piace il nuotando.
\end{options}
\item\label{ex:infinito-gerundio-3}\relax Give the gerundio presente.\par\smallskip{\begin{tabular}{@{}l l@{}}
cantare & \blankcell \\
leggere & \blankcell \\
partire & \blankcell \\
fare & \blankcell \\
dire & \blankcell \\
bere & \blankcell \\
\end{tabular}}
\item\label{ex:infinito-gerundio-4}\relax Translate: We were having dinner. (use stare + gerundio)\par\answerline
\item\label{ex:infinito-gerundio-5}\relax \textit{Ho visto Sara uscendo dal negozio.} Who was coming out of the shop?\par\nopagebreak\begin{options}
\item\relax nobody
\item\relax both
\item\relax the speaker
\item\relax Sara
\end{options}
\item\label{ex:infinito-gerundio-6}\relax Complete with the past infinitive.\par\smallskip\begingroup\setlength{\parindent}{0pt}Dopo \clozeblank{1} (arrivare, lei), Sara ha telefonato. Dopo \clozeblank{2} (leggere, noi) il libro, abbiamo visto il film.\par\endgroup
\end{exercises}

\begin{summarybox}
\begin{itemize}
\item\relax Infinito passato: avere/essere + participle; dopo + infinito passato.
\item\relax Infinitive after modals and prepositions, as a noun, in instructions.
\item\relax Gerundio: -ando/-endo (facendo, dicendo, bevendo); passato: avendo/essendo + participle.
\item\relax The gerund shares the main subject; never after prepositions; stare + gerundio for actions in progress.
\end{itemize}
\end{summarybox}

\relatedwords{\textit{mangiare} to eat \textperiodcentered{} \textit{uscire} to go out; to leave \textperiodcentered{} \textit{leggere} to read \textperiodcentered{} \textit{nuotare} to swim \textperiodcentered{} \textit{finire} to finish; to end \textperiodcentered{} \textit{arrivare} to arrive \textperiodcentered{} \textit{stare} to stay; to be (of health)}
\applyit{Read \textit{Il segreto del faro} (A2–B1), \S\ref{story:il-segreto-del-faro}.}

\section{The participle}\label{les:participio}
\lessonmeta{B1}{The past participle as verb, adjective and participial clause; the present participle.}
\begin{objectives}
\item\relax Recognise the past participle in compound tenses, the passive and as an adjective.
\item\relax Understand participial clauses (finita la cena, …) and their agreement.
\item\relax Recognise the present participle (-ante, -ente) and know it is mostly used as a noun or adjective.
\end{objectives}

\subsection{The past participle}\label{les:participio:past}

\begin{itemize}
\item\relax Compound tenses: \textit{ho mangiato, sono partita}.
\item\relax Passive: \textit{la lettera è stata scritta}.
\item\relax Adjective: \textit{una porta chiusa, i bambini stanchi e addormentati}.
\item\relax Noun: \textit{il passato, gli invitati}.
\end{itemize}

\begin{notebox}{rulecol}\notelabel{rulecol}{Rule}{Participial clauses}
A past participle can form a short clause expressing an action completed before the main one: \textit{Finita la cena, uscimmo} (When dinner was over, we went out). With transitive verbs the participle agrees with the noun it refers to (\textit{Finita la cena}, \textit{Letti i giornali}); with intransitive verbs it agrees with the subject (\textit{Arrivata a casa, Sara si riposò}). This construction is mainly written.
\end{notebox}

\begin{notebox}{tendcol}\notelabel{tendcol}{Tendency}{Pronouns attach}
In participial clauses pronouns attach to the participle: \textit{Salutatolo, se ne andò} (Having said goodbye to him, she left). This is literary.
\end{notebox}

\subsection{The present participle}\label{les:participio:present}

\begin{notebox}{rulecol}\notelabel{rulecol}{Rule}{-ante, -ente}
\textit{cantare → cantante, dormire → dormiente, bollire → bollente}. In modern Italian it survives mainly as an adjective or noun: \textit{un cantante} (a singer), \textit{acqua corrente} (running water), \textit{il presidente}. As a real verb form (\textit{le persone abitanti in città}) it belongs to bureaucratic or literary language; use a relative clause instead (\textit{le persone che abitano in città}).
\end{notebox}

\subsection*{Worked examples}
\begin{examples}
\item\relax \textit{Finita la cena, uscimmo in giardino.}\newline When dinner was over, we went out into the garden.\newline {\color{muted}\small Participial clause; finita agrees with la cena.}
\item\relax \textit{Arrivati alla stazione, i ragazzi comprarono i biglietti.}\newline Having arrived at the station, the boys bought the tickets.\newline {\color{muted}\small Intransitive verb: agreement with the subject.}
\item\relax \textit{La porta è chiusa.}\newline The door is closed.\newline {\color{muted}\small Participle as an adjective describing a state.}
\item\relax \textit{Il cantante ha una bella voce.}\newline The singer has a beautiful voice.\newline {\color{muted}\small Present participle used as a noun.}
\item\relax \textit{Attenzione: l'acqua è bollente!}\newline Careful: the water is boiling hot!
\item\relax \textit{Letta la lettera, la nonna sorrise.}\newline Having read the letter, Grandmother smiled.
\end{examples}

\subsection*{Exercises}\label{les:participio:exercises}
{\color{muted}\small Answers: page~\pageref{ans:participio}.}

\begin{exercises}
\item\label{ex:participio-1}\relax \blank{} le lettere, la nonna andò alla posta.\par\nopagebreak\begin{options}
\item\relax Scrivendo le
\item\relax Scritti
\item\relax Scritte
\item\relax Scritto
\end{options}
\item\label{ex:participio-2}\relax \blank{} a casa, Sara si riposò.\par\nopagebreak\begin{options}
\item\relax Arrivando
\item\relax Arrivata
\item\relax Arrivati
\item\relax Arrivato
\end{options}
\item\label{ex:participio-3}\relax In \textit{una finestra aperta}, what is \textit{aperta}?\par\nopagebreak\begin{options}
\item\relax a present participle
\item\relax a gerund
\item\relax a passato prossimo
\item\relax a past participle used as an adjective
\end{options}
\item\label{ex:participio-4}\relax Match the present participle to its meaning.\par\smallskip\begin{tabularx}{\linewidth}{@{}L L@{}}
1.\ cantante & (a)\ boiling hot \\
2.\ insegnante & (b)\ running, current \\
3.\ bollente & (c)\ teacher \\
4.\ corrente & (d)\ singer \\
\end{tabularx}
\item\label{ex:participio-5}\relax Translate: \textit{Finita la lezione, i bambini uscirono.}\par\answerline
\end{exercises}

\begin{summarybox}
\begin{itemize}
\item\relax Past participle: compound tenses, passive, adjective, noun, participial clause.
\item\relax Participial clauses: agreement with the object (transitive) or subject (intransitive).
\item\relax Present participle (-ante, -ente): mostly nouns and adjectives today.
\end{itemize}
\end{summarybox}

\relatedwords{\textit{chiudere} to close \textperiodcentered{} \textit{aprire} to open \textperiodcentered{} \textit{scrivere} to write \textperiodcentered{} \textit{arrivare} to arrive \textperiodcentered{} \textit{finire} to finish; to end \textperiodcentered{} \textit{cantare} to sing}
\applyit{Read \textit{Il segreto del faro} (A2–B1), \S\ref{story:il-segreto-del-faro}.}

\chapter{Passato remoto}\label{mod:m17}
\moduleintro{B1}{The narrative past: regular and irregular forms (including alternatives such as credei/credetti), its use in writing and its very different use in speech across Italy; the trapassato remoto.}{Passato remoto: forms, meaning and regional use (\S\ref{les:passato-remoto}); Trapassato remoto (\S\ref{les:trapassato-remoto})}

\section{Passato remoto: forms, meaning and regional use}\label{les:passato-remoto}
\lessonmeta{B1}{The narrative past: regular and irregular forms, its difference from the passato prossimo, and how its use varies by register and region.}
\begin{objectives}
\item\relax Form the regular passato remoto of all three conjugations, including the two -ere series (-ei/-etti).
\item\relax Recognise the '1-3-3' pattern of strong verbs (ebbi, ebbe, ebbero) and learn the most frequent ones.
\item\relax Explain the aspectual difference between passato remoto and passato prossimo.
\item\relax Describe how usage differs between written narration and speech, and across Italian regions.
\end{objectives}

\subsection{Regular forms}\label{les:passato-remoto:regular}

\par\noindent\begin{minipage}[t]{0.48\linewidth}\raggedright
\paradigmbox{parlare --- to speak; to talk \textperiodcentered{} indicativo passato remoto}{%
\begin{tabular}{@{}l l@{}}
\textit{io} & parl\hlending{ai} \\
\textit{tu} & parl\hlending{asti} \\
\textit{lui/lei} & parl\hlending{ò} \\
\textit{noi} & parl\hlending{ammo} \\
\textit{voi} & parl\hlending{aste} \\
\textit{loro} & parl\hlending{arono} \\
\end{tabular}}{Stem \textit{parl-} + \hlending{ending}.}
\end{minipage}\hfill\begin{minipage}[t]{0.48\linewidth}\raggedright
\paradigmbox{credere --- to believe \textperiodcentered{} indicativo passato remoto}{%
\begin{tabular}{@{}l l@{}}
\textit{io} & cred\hlending{ei} / cred\hlending{etti} \\
\textit{tu} & cred\hlending{esti} \\
\textit{lui/lei} & cred\hlending{é} / cred\hlending{ette} \\
\textit{noi} & cred\hlending{emmo} \\
\textit{voi} & cred\hlending{este} \\
\textit{loro} & cred\hlending{erono} / cred\hlending{ettero} \\
\end{tabular}}{Stem \textit{cred-} + \hlending{ending}. Where two forms are given, both are standard.}
\end{minipage}\par\medskip\noindent\begin{minipage}[t]{0.48\linewidth}\raggedright
\paradigmbox{dormire --- to sleep \textperiodcentered{} indicativo passato remoto}{%
\begin{tabular}{@{}l l@{}}
\textit{io} & dorm\hlending{ii} \\
\textit{tu} & dorm\hlending{isti} \\
\textit{lui/lei} & dorm\hlending{ì} \\
\textit{noi} & dorm\hlending{immo} \\
\textit{voi} & dorm\hlending{iste} \\
\textit{loro} & dorm\hlending{irono} \\
\end{tabular}}{Stem \textit{dorm-} + \hlending{ending}.}
\end{minipage}\par\medskip

\begin{notebox}{rulecol}\notelabel{rulecol}{Rule}{Two series in -ere}
Regular -ere verbs have two equally standard series in io, lui/lei and loro: \textit{credei / credetti, credé / credette, crederono / credettero}. Which series is preferred varies from verb to verb and writer to writer; both are correct (Treccani gives \textit{vendei o vendetti}). Final stress is written: \textit{parlò, credé, dormì}.
\end{notebox}

\begin{notebox}{rulecol}\notelabel{rulecol}{Rule}{No -isc-}
-isc- verbs do not insert -isc- in the passato remoto: \textit{capii, capisti, capì}.
\end{notebox}

\subsection{Irregular (strong) verbs}\label{les:passato-remoto:irregular}

\begin{notebox}{rulecol}\notelabel{rulecol}{Rule}{The 1-3-3 pattern}
Most irregular verbs are irregular only in the 1st person singular and the 3rd persons singular and plural, which share a special stem ending in -i, -e, -ero: \textit{ebbi, ebbe, ebbero}. The other three persons are regular, built on the normal stem: \textit{avesti, avemmo, aveste}.
\end{notebox}

\par\noindent\begin{minipage}[t]{0.48\linewidth}\raggedright
\paradigmbox{avere --- to have \textperiodcentered{} indicativo passato remoto}{%
\begin{tabular}{@{}l l@{}}
\textit{io} & \irr{ebbi} \\
\textit{tu} & avesti \\
\textit{lui/lei} & \irr{ebbe} \\
\textit{noi} & avemmo \\
\textit{voi} & aveste \\
\textit{loro} & \irr{ebbero} \\
\end{tabular}}{\irr{Highlighted} forms differ from the regular pattern of the verb's family.}
\end{minipage}\hfill\begin{minipage}[t]{0.48\linewidth}\raggedright
\paradigmbox{fare --- to do; to make \textperiodcentered{} indicativo passato remoto}{%
\begin{tabular}{@{}l l@{}}
\textit{io} & \irr{feci} \\
\textit{tu} & \irr{facesti} \\
\textit{lui/lei} & \irr{fece} \\
\textit{noi} & \irr{facemmo} \\
\textit{voi} & \irr{faceste} \\
\textit{loro} & \irr{fecero} \\
\end{tabular}}{\irr{Highlighted} forms differ from the regular pattern of the verb's family.}
\end{minipage}\par\medskip

{\small\setlength{\LTpre}{2pt}\setlength{\LTpost}{6pt}
\tablecaption{Frequent strong verbs (io · lui/lei · loro)}
\begin{xltabular}{\linewidth}{lLlL}
\toprule
\textbf{Infinitive} & \textbf{Forms} & \textbf{Infinitive} & \textbf{Forms} \\
\midrule
avere & ebbi · ebbe · ebbero & venire & venni · venne · vennero \\
fare & feci · fece · fecero & volere & volli · volle · vollero \\
dire & dissi · disse · dissero & sapere & seppi · seppe · seppero \\
vedere & vidi · vide · videro & conoscere & conobbi · conobbe · conobbero \\
prendere & presi · prese · presero & nascere & nacqui · nacque · nacquero \\
mettere & misi · mise · misero & rimanere & rimasi · rimase · rimasero \\
scrivere & scrissi · scrisse · scrissero & bere & bevvi · bevve · bevvero \\
leggere & lessi · lesse · lessero & cadere & caddi · cadde · caddero \\
chiedere & chiesi · chiese · chiesero & piacere & piacqui · piacque · piacquero \\
rispondere & risposi · rispose · risposero & vivere & vissi · visse · vissero \\
\bottomrule
\end{xltabular}}

\par\noindent\begin{minipage}[t]{0.48\linewidth}\raggedright
\paradigmbox{essere --- to be \textperiodcentered{} indicativo passato remoto}{%
\begin{tabular}{@{}l l@{}}
\textit{io} & \irr{fui} \\
\textit{tu} & \irr{fosti} \\
\textit{lui/lei} & \irr{fu} \\
\textit{noi} & \irr{fummo} \\
\textit{voi} & \irr{foste} \\
\textit{loro} & \irr{furono} \\
\end{tabular}}{\irr{Highlighted} forms differ from the regular pattern of the verb's family.}
\end{minipage}\par\medskip

\begin{notebox}{exccol}\notelabel{exccol}{Exception}{essere, dare, stare}
\textit{essere} is irregular throughout (\textit{fui, fosti, fu, fummo, foste, furono}). \textit{dare}: \textit{diedi/detti, desti, diede/dette, demmo, deste, diedero/dettero}. \textit{stare}: \textit{stetti, stesti, stette, stemmo, steste, stettero}.
\end{notebox}

\subsection{Passato remoto or passato prossimo?}\label{les:passato-remoto:meaning}

\begin{notebox}{rulecol}\notelabel{rulecol}{Rule}{The aspectual distinction}
Both tenses present an action as completed. The \textbf{passato prossimo} links the past event to the present (its result or relevance still holds, or the time period includes the present: \textit{quest'anno, oggi}). The \textbf{passato remoto} presents the event as closed and detached from the present. Hence the norm: \textit{Oggi ho visto Marco} (not \textit{vidi}), but \textit{Dante nacque nel 1265}.
\end{notebox}

\begin{notebox}{rulecol}\notelabel{rulecol}{Rule}{Written narration}
The passato remoto is the standard tense of written narration — novels, fairy tales, history books — where it carries the main sequence of events while the imperfetto provides the background: \textit{Era notte. Il lupo entrò nella casa.}
\end{notebox}

\begin{notebox}{tendcol}\notelabel{tendcol}{Tendency}{Speech varies by region}
In spoken Italian the choice depends strongly on region. In the North the passato prossimo has largely replaced the passato remoto in conversation, even for distant events. In Tuscany the traditional distinction is kept. In much of the South the passato remoto is used in speech even where standard Italian would prefer the passato prossimo (\textit{Stamattina andai al mercato}); regional Italian generally mirrors the local dialects. These are tendencies of regional Italian, not rules: a speaker from Milan and one from Palermo may describe the same event differently and both be understood.
\end{notebox}

\begin{notebox}{tipcol}\notelabel{tipcol}{Tip}{What to do as a learner}
Learn to recognise every form: you will meet them in every book. For your own speech, the passato prossimo is always acceptable for recent events; use the passato remoto for historical facts and written narration.
\end{notebox}

\subsection*{Worked examples}
\begin{examples}
\item\relax \textit{C'era una volta un re che aveva tre figlie.}\newline Once upon a time there was a king who had three daughters.\newline {\color{muted}\small Fairy-tale background: imperfetto.}
\item\relax \textit{Un giorno il re chiamò le figlie e disse: «Ascoltatemi.»}\newline One day the king called his daughters and said: “Listen to me.”\newline {\color{muted}\small Events in the passato remoto: chiamò, disse.}
\item\relax \textit{Dante nacque a Firenze nel 1265.}\newline Dante was born in Florence in 1265.\newline {\color{muted}\small Historical fact: passato remoto.}
\item\relax \textit{Oggi ho visto Marco al mercato.}\newline Today I saw Marco at the market.\newline {\color{muted}\small oggi includes the present: passato prossimo in standard Italian.}
\item\relax \textit{I ragazzi videro una luce e si fermarono.}\newline The boys saw a light and stopped.\newline {\color{muted}\small vedere: videro (strong); fermarsi: si fermarono (regular).}
\item\relax \textit{Fu una notte lunga e fredda.}\newline It was a long, cold night.\newline {\color{muted}\small essere: fu presents the whole night as a closed event.}
\item\relax \textit{Quando arrivammo al porto, il sole tramontava.}\newline When we reached the harbour, the sun was setting.\newline {\color{muted}\small Event (passato remoto) against background (imperfetto).}
\end{examples}

\subsection*{Exercises}\label{les:passato-remoto:exercises}
{\color{muted}\small Answers: page~\pageref{ans:passato-remoto}.}

\begin{exercises}
\item\label{ex:passato-remoto-1}\relax Conjugate \textit{cantare} in the passato remoto.\par\smallskip{\begin{tabular}{@{}l l@{}}
io & \blankcell \\
tu & \blankcell \\
lui/lei & \blankcell \\
noi & \blankcell \\
voi & \blankcell \\
loro & \blankcell \\
\end{tabular}}
\item\label{ex:passato-remoto-2}\relax vendere, io, passato remoto: \blank{}
\item\label{ex:passato-remoto-3}\relax Give the infinitive.\par\smallskip{\begin{tabular}{@{}l l@{}}
ebbe & \blankcell \\
fece & \blankcell \\
disse & \blankcell \\
vide & \blankcell \\
nacque & \blankcell \\
fu & \blankcell \\
\end{tabular}}
\item\label{ex:passato-remoto-4}\relax Conjugate \textit{prendere} in the passato remoto.\par\smallskip{\begin{tabular}{@{}l l@{}}
io & \blankcell \\
tu & \blankcell \\
lui/lei & \blankcell \\
noi & \blankcell \\
voi & \blankcell \\
loro & \blankcell \\
\end{tabular}}
\item\label{ex:passato-remoto-5}\relax In standard Italian, which is the expected form? \textit{Stamattina \blank{} Marco al bar.}\par\nopagebreak\begin{options}
\item\relax avevo visto
\item\relax ho visto
\item\relax vedevo
\item\relax vidi
\end{options}
\item\label{ex:passato-remoto-6}\relax Which statement about the passato remoto is accurate?\par\nopagebreak\begin{options}
\item\relax It is used only in the South of Italy.
\item\relax It is used for actions still happening now.
\item\relax It is standard in written narration; in speech its use varies strongly by region.
\item\relax It is obsolete and never used in modern Italian.
\end{options}
\item\label{ex:passato-remoto-7}\relax Complete the story with the passato remoto.\par\smallskip\begingroup\setlength{\parindent}{0pt}Il lupo \clozeblank{1} (entrare) nella casa e \clozeblank{2} (vedere) la nonna. La nonna \clozeblank{3} (avere) paura e \clozeblank{4} (scappare).\par\endgroup
\end{exercises}

\begin{summarybox}
\begin{itemize}
\item\relax Regular: -ai/-asti/-ò…; -ei(-etti)/-esti/-é(-ette)…; -ii/-isti/-ì…
\item\relax Strong verbs: irregular stem in io, lui/lei, loro (ebbi, ebbe, ebbero); regular elsewhere.
\item\relax essere: fui…; dare: diedi/detti…; stare: stetti…
\item\relax Passato remoto = closed, detached past; standard in written narration and history.
\item\relax Spoken use varies by region: rare in the North, traditional in Tuscany, very common in the South.
\end{itemize}
\end{summarybox}

\relatedwords{\textit{avere} to have \textperiodcentered{} \textit{fare} to do; to make \textperiodcentered{} \textit{dire} to say; to tell \textperiodcentered{} \textit{vedere} to see \textperiodcentered{} \textit{prendere} to take; to catch; to have (food) \textperiodcentered{} \textit{nascere} to be born \textperiodcentered{} \textit{essere} to be \textperiodcentered{} \textit{venire} to come \textperiodcentered{} \textit{vendere} to sell}
\applyit{Read \textit{Il segreto del faro} (A2–B1), \S\ref{story:il-segreto-del-faro}.}

\section{Trapassato remoto}\label{les:trapassato-remoto}
\lessonmeta{B2}{The literary past anterior, used in temporal clauses with a main verb in the passato remoto.}
\begin{objectives}
\item\relax Form the trapassato remoto (ebbi mangiato, fui partito).
\item\relax Use it after dopo che, (non) appena, quando, una volta che with a main verb in the passato remoto.
\item\relax Know that outside this construction the trapassato prossimo is used.
\end{objectives}

\subsection{Formation}\label{les:trapassato-remoto:formation}

\begin{notebox}{rulecol}\notelabel{rulecol}{Rule}{Passato remoto of the auxiliary + participle}
\textit{ebbi finito, avesti finito, ebbe finito…; fui partito/a, fosti partito/a, fu partito/a…}
\end{notebox}

\par\noindent\begin{minipage}[t]{0.48\linewidth}\raggedright
\paradigmbox{finire --- to finish; to end \textperiodcentered{} indicativo trapassato remoto (with avere)}{%
\begin{tabular}{@{}l l@{}}
\textit{io} & ebbi finito \\
\textit{tu} & avesti finito \\
\textit{lui/lei} & ebbe finito \\
\textit{noi} & avemmo finito \\
\textit{voi} & aveste finito \\
\textit{loro} & ebbero finito \\
\end{tabular}}{}
\end{minipage}\hfill\begin{minipage}[t]{0.48\linewidth}\raggedright
\paradigmbox{partire --- to leave; to set off \textperiodcentered{} indicativo trapassato remoto (with essere, masculine subject)}{%
\begin{tabular}{@{}l l@{}}
\textit{io} & fui partito \\
\textit{tu} & fosti partito \\
\textit{lui/lei} & fu partito \\
\textit{noi} & fummo partiti \\
\textit{voi} & foste partiti \\
\textit{loro} & furono partiti \\
\end{tabular}}{}
\end{minipage}\par\medskip

\subsection{Use}\label{les:trapassato-remoto:use}

\begin{notebox}{rulecol}\notelabel{rulecol}{Rule}{Only in one construction}
The trapassato remoto expresses an action immediately preceding another action in the passato remoto, in a temporal clause introduced by \textit{dopo che, (non) appena, quando, come, una volta che}: \textit{Dopo che ebbe cenato, uscì.} Elsewhere, anteriority in the past is expressed with the trapassato prossimo.
\end{notebox}

\begin{notebox}{tendcol}\notelabel{tendcol}{Tendency}{Register}
It belongs to written, literary narration and is rare even there; many writers prefer alternatives with the same meaning: \textit{Dopo aver cenato, uscì} or \textit{Cenato, uscì}. In speech it is practically never used.
\end{notebox}

\subsection*{Worked examples}
\begin{examples}
\item\relax \textit{Dopo che ebbe cenato, il vecchio uscì sul molo.}\newline When he had had dinner, the old man went out onto the pier.
\item\relax \textit{Appena fu partito il treno, cominciò a piovere.}\newline As soon as the train had left, it began to rain.
\item\relax \textit{Quando ebbero finito il lavoro, i pescatori tornarono a casa.}\newline When they had finished their work, the fishermen went home.
\item\relax \textit{Dopo aver cenato, il vecchio uscì sul molo.}\newline After having dinner, the old man went out onto the pier.\newline {\color{muted}\small Equivalent, more common: dopo + infinito passato.}
\item\relax \textit{Non appena l'ebbe vista, la riconobbe.}\newline As soon as he had seen her, he recognised her.\newline {\color{muted}\small Agreement with la (l'ebbe vista).}
\end{examples}

\subsection*{Exercises}\label{les:trapassato-remoto:exercises}
{\color{muted}\small Answers: page~\pageref{ans:trapassato-remoto}.}

\begin{exercises}
\item\label{ex:trapassato-remoto-1}\relax mangiare, lui, trapassato remoto: \blank{}
\item\label{ex:trapassato-remoto-2}\relax arrivare, loro (masculine), trapassato remoto: \blank{}
\item\label{ex:trapassato-remoto-3}\relax Dopo che \blank{} (leggere) la lettera, la principessa pianse.\par\nopagebreak\begin{options}
\item\relax lesse
\item\relax ebbe letto
\item\relax ha letto
\item\relax aveva letto
\end{options}
\item\label{ex:trapassato-remoto-4}\relax Which sentence uses the trapassato remoto correctly?\par\nopagebreak\begin{options}
\item\relax Appena ebbe finito, se ne andò.
\item\relax Se ebbe finito, uscirebbe.
\item\relax Quando ebbi finito, sono uscito.
\item\relax Ebbe finito il lavoro ieri.
\end{options}
\item\label{ex:trapassato-remoto-5}\relax Translate: \textit{Quando ebbero finito, tornarono a casa.}\par\answerline
\end{exercises}

\begin{summarybox}
\begin{itemize}
\item\relax Trapassato remoto = passato remoto of avere/essere + participle.
\item\relax Only in temporal clauses (dopo che, appena, quando) before a main verb in the passato remoto.
\item\relax Literary and rare; dopo + infinito passato is the usual alternative.
\end{itemize}
\end{summarybox}

\relatedwords{\textit{finire} to finish; to end \textperiodcentered{} \textit{partire} to leave; to set off \textperiodcentered{} \textit{mangiare} to eat \textperiodcentered{} \textit{leggere} to read \textperiodcentered{} \textit{piangere} to cry}
\applyit{Read \textit{Il segreto del faro} (A2–B1), \S\ref{story:il-segreto-del-faro}.}

\chapter{Congiuntivo}\label{mod:m18}
\moduleintro{B1}{All four tenses of the subjunctive: when it is required, when it alternates with the indicative, and what changes with register.}{Congiuntivo presente: forms and main uses (\S\ref{les:congiuntivo-presente}); Congiuntivo passato (\S\ref{les:congiuntivo-passato}); Congiuntivo imperfetto and trapassato (\S\ref{les:congiuntivo-imperfetto-trapassato})}

\section{Congiuntivo presente: forms and main uses}\label{les:congiuntivo-presente}
\lessonmeta{B1}{How to form the present subjunctive, which contexts require it, and how register affects its use.}
\begin{objectives}
\item\relax Form the congiuntivo presente of regular and irregular verbs.
\item\relax Use it after verbs of opinion, wish, emotion and doubt, after impersonal expressions and after certain conjunctions.
\item\relax Use di + infinitive instead when the two clauses have the same subject.
\item\relax Understand that the indicative after verbs of opinion is informal, and that some contexts allow a real choice of meaning.
\end{objectives}

\subsection{Formation}\label{les:congiuntivo-presente:formation}

\begin{notebox}{rulecol}\notelabel{rulecol}{Rule}{Endings}
-are verbs: \textit{-i, -i, -i, -iamo, -iate, -ino}. -ere and -ire verbs: \textit{-a, -a, -a, -iamo, -iate, -ano}. The three singular forms are identical, so a subject pronoun is often added for clarity: \textit{che io parli, che tu parli, che lui parli}.
\end{notebox}

\par\noindent\begin{minipage}[t]{0.48\linewidth}\raggedright
\paradigmbox{parlare --- to speak; to talk \textperiodcentered{} congiuntivo presente}{%
\begin{tabular}{@{}l l@{}}
\textit{io} & parl\hlending{i} \\
\textit{tu} & parl\hlending{i} \\
\textit{lui/lei} & parl\hlending{i} \\
\textit{noi} & parl\hlending{iamo} \\
\textit{voi} & parl\hlending{iate} \\
\textit{loro} & parl\hlending{ino} \\
\end{tabular}}{Stem \textit{parl-} + \hlending{ending}.}
\end{minipage}\hfill\begin{minipage}[t]{0.48\linewidth}\raggedright
\paradigmbox{vedere --- to see \textperiodcentered{} congiuntivo presente}{%
\begin{tabular}{@{}l l@{}}
\textit{io} & ved\hlending{a} \\
\textit{tu} & ved\hlending{a} \\
\textit{lui/lei} & ved\hlending{a} \\
\textit{noi} & ved\hlending{iamo} \\
\textit{voi} & ved\hlending{iate} \\
\textit{loro} & ved\hlending{ano} \\
\end{tabular}}{Stem \textit{ved-} + \hlending{ending}.}
\end{minipage}\par\medskip\noindent\begin{minipage}[t]{0.48\linewidth}\raggedright
\paradigmbox{capire --- to understand \textperiodcentered{} congiuntivo presente}{%
\begin{tabular}{@{}l l@{}}
\textit{io} & capisca \\
\textit{tu} & capisca \\
\textit{lui/lei} & capisca \\
\textit{noi} & capiamo \\
\textit{voi} & capiate \\
\textit{loro} & capiscano \\
\end{tabular}}{}
\end{minipage}\par\medskip

\begin{notebox}{rulecol}\notelabel{rulecol}{Rule}{Irregular verbs: start from the io form}
Most irregular verbs build the singular and loro on the io form of the present, and noi/voi on the noi form: \textit{vado → vada, vadano; andiamo → andiamo, andiate}. So: \textit{faccia, dica, venga, esca, possa, voglia, beva, rimanga, piaccia, scelga}. Fully irregular: \textit{essere (sia), avere (abbia), sapere (sappia), dare (dia), stare (stia)}. \textit{dovere}: \textit{debba} (the alternative \textit{deva} exists but is much less common).
\end{notebox}

\par\noindent\begin{minipage}[t]{0.48\linewidth}\raggedright
\paradigmbox{essere --- to be \textperiodcentered{} congiuntivo presente}{%
\begin{tabular}{@{}l l@{}}
\textit{io} & \irr{sia} \\
\textit{tu} & \irr{sia} \\
\textit{lui/lei} & \irr{sia} \\
\textit{noi} & \irr{siamo} \\
\textit{voi} & \irr{siate} \\
\textit{loro} & \irr{siano} \\
\end{tabular}}{\irr{Highlighted} forms differ from the regular pattern of the verb's family.}
\end{minipage}\hfill\begin{minipage}[t]{0.48\linewidth}\raggedright
\paradigmbox{andare --- to go \textperiodcentered{} congiuntivo presente}{%
\begin{tabular}{@{}l l@{}}
\textit{io} & \irr{vada} \\
\textit{tu} & \irr{vada} \\
\textit{lui/lei} & \irr{vada} \\
\textit{noi} & andiamo \\
\textit{voi} & andiate \\
\textit{loro} & \irr{vadano} \\
\end{tabular}}{\irr{Highlighted} forms differ from the regular pattern of the verb's family.}
\end{minipage}\par\medskip

\subsection{When the congiuntivo is required}\label{les:congiuntivo-presente:uses}

{\small\setlength{\LTpre}{2pt}\setlength{\LTpost}{6pt}
\tablecaption{Main contexts}
\begin{xltabular}{\linewidth}{LL}
\toprule
\textbf{Context} & \textbf{Examples} \\
\midrule
opinion, belief, supposition & Penso che sia vero. Credo che Marco abiti a Roma. Immagino che tu sia stanco. \\
wish, will, request & Voglio che tu venga. Preferisco che restiate a casa. \\
emotion & Sono contento che tu sia qui. Mi dispiace che piova. \\
doubt, uncertainty & Dubito che arrivi in tempo. Non sono sicuro che sia aperto. \\
impersonal expressions & Bisogna che tu parta. È importante che studiate. È possibile che nevichi. \\
conjunctions & benché/sebbene (although), affinché/perché (so that), purché (provided that), prima che (before), senza che (without), a meno che (unless), nonostante (although) \\
indefinite or hypothetical antecedent & Cerco qualcuno che sappia il cinese. \\
after a superlative & È il libro più bello che io abbia mai letto. \\
\bottomrule
\end{xltabular}}

\begin{notebox}{rulecol}\notelabel{rulecol}{Rule}{Same subject: di + infinitive}
If the subject of both clauses is the same, use \textit{di} + infinitive, not \textit{che} + congiuntivo: \textit{Penso di partire} (I think I'll leave), not \textit{Penso che io parta}. \textit{Spero di vederti}, \textit{Sono contento di essere qui}.
\end{notebox}

\begin{notebox}{rulecol}\notelabel{rulecol}{Rule}{Indicative after statements of certainty}
Verbs and expressions of certainty and perception take the indicative: \textit{So che è vero. È certo che verrà. Vedo che sei stanco.} \textit{perché} meaning 'because' takes the indicative; \textit{perché} meaning 'so that' takes the congiuntivo.
\end{notebox}

\subsection{Register and variation}\label{les:congiuntivo-presente:register}

\begin{notebox}{tendcol}\notelabel{tendcol}{Tendency}{Indicative in informal speech}
In everyday speech the indicative often replaces the congiuntivo after verbs of opinion (\textit{credo che è vero}). This is understood everywhere but belongs to a lower, informal register; in writing and careful speech the congiuntivo is expected (Treccani, \textit{Enciclopedia dell'italiano}, s.v. congiuntivo).
\end{notebox}

\begin{notebox}{tendcol}\notelabel{tendcol}{Tendency}{sperare che + futuro}
After \textit{sperare} a future event can also be expressed with the futuro indicativo: \textit{Spero che verrai} alongside \textit{Spero che tu venga}. Both are widely used; the congiuntivo is the more formal choice.
\end{notebox}

\begin{notebox}{tendcol}\notelabel{tendcol}{Tendency}{A real choice in relative clauses}
In restrictive relative clauses the mood changes the meaning: \textit{Cerco lo studente che ha perso il libretto} (a specific student who exists) / \textit{Cerco uno studente che legga il cinese} (anyone who might). (Treccani, s.v. congiuntivo)
\end{notebox}

\subsection*{Worked examples}
\begin{examples}
\item\relax \textit{Penso che Marco sia stanco.}\newline I think Marco is tired.
\item\relax \textit{Voglio che tu venga con me.}\newline I want you to come with me.\newline {\color{muted}\small English uses an infinitive; Italian needs che + congiuntivo with a different subject.}
\item\relax \textit{Penso di partire domani.}\newline I think I'll leave tomorrow.\newline {\color{muted}\small Same subject: di + infinitive.}
\item\relax \textit{Benché piova, usciamo.}\newline Although it's raining, we're going out.
\item\relax \textit{Ti presto il libro purché tu me lo restituisca.}\newline I'll lend you the book provided you give it back to me.
\item\relax \textit{So che è vero.}\newline I know it's true.\newline {\color{muted}\small Certainty: indicative.}
\item\relax \textit{È il libro più bello che io abbia mai letto.}\newline It's the most beautiful book I have ever read.\newline {\color{muted}\small After a superlative: congiuntivo passato.}
\end{examples}

\subsection*{Exercises}\label{les:congiuntivo-presente:exercises}
{\color{muted}\small Answers: page~\pageref{ans:congiuntivo-presente}.}

\begin{exercises}
\item\label{ex:congiuntivo-presente-1}\relax Give the congiuntivo presente (che lui/lei …).\par\smallskip{\begin{tabular}{@{}l l@{}}
parlare & \blankcell \\
prendere & \blankcell \\
finire & \blankcell \\
fare & \blankcell \\
andare & \blankcell \\
essere & \blankcell \\
\end{tabular}}
\item\label{ex:congiuntivo-presente-2}\relax Credo che Sara \blank{} ragione.\par\nopagebreak\begin{options}
\item\relax avrebbe
\item\relax abbia
\item\relax ha
\item\relax avesse
\end{options}
\item\label{ex:congiuntivo-presente-3}\relax 'I hope to see you soon.'\par\nopagebreak\begin{options}
\item\relax Spero vederti presto.
\item\relax Spero di vederti presto.
\item\relax Spero che io ti veda presto.
\item\relax Spero che ti vedo presto.
\end{options}
\item\label{ex:congiuntivo-presente-4}\relax Indicative or congiuntivo?\par\smallskip\begin{tabularx}{\linewidth}{@{}L L@{}}
1.\ So che… & (a)\ congiuntivo (impersonal) \\
2.\ Penso che… & (b)\ congiuntivo (conjunction) \\
3.\ Benché… & (c)\ indicative \\
4.\ È certo che… & (d)\ indicative (certainty) \\
5.\ Bisogna che… & (e)\ congiuntivo \\
\end{tabularx}
\item\label{ex:congiuntivo-presente-5}\relax Complete with the congiuntivo presente.\par\smallskip\begingroup\setlength{\parindent}{0pt}Voglio che voi \clozeblank{1} (venire) alla festa. È importante che tu \clozeblank{2} (dormire). Dubito che loro \clozeblank{3} (sapere) la risposta.\par\endgroup
\item\label{ex:congiuntivo-presente-6}\relax Translate: I want you to come with me. (tu)\par\answerline
\item\label{ex:congiuntivo-presente-7}\relax Ti spiego la regola \blank{} tu la capisca.\par\nopagebreak\begin{options}
\item\relax affinché
\item\relax perché (= because)
\item\relax quando
\item\relax siccome
\end{options}
\end{exercises}

\begin{summarybox}
\begin{itemize}
\item\relax -are: -i, -i, -i, -iamo, -iate, -ino; -ere/-ire: -a, -a, -a, -iamo, -iate, -ano.
\item\relax Irregulars from the io form (vada, faccia, venga); sia, abbia, sappia, dia, stia.
\item\relax After opinion, wish, emotion, doubt, impersonal expressions, benché/affinché/purché/prima che/senza che, indefinite antecedents, superlatives.
\item\relax Same subject → di + infinitive. Certainty → indicative.
\item\relax Informal speech often uses the indicative after verbs of opinion; standard writing uses the congiuntivo.
\end{itemize}
\end{summarybox}

\relatedwords{\textit{pensare} to think \textperiodcentered{} \textit{credere} to believe \textperiodcentered{} \textit{volere} to want \textperiodcentered{} \textit{sperare} to hope \textperiodcentered{} \textit{dubitare} to doubt \textperiodcentered{} \textit{bisognare} to be necessary (impersonal) \textperiodcentered{} \textit{essere} to be \textperiodcentered{} \textit{avere} to have \textperiodcentered{} \textit{temere} to fear; to be afraid \textperiodcentered{} \textit{contento} glad; pleased}
\applyit{Read \textit{Una lettera dalla nonna} (A2–B1), \S\ref{story:una-lettera-dalla-nonna}.}

\section{Congiuntivo passato}\label{les:congiuntivo-passato}
\lessonmeta{B1}{The perfect subjunctive: an action before the time of a present main verb.}
\begin{objectives}
\item\relax Form the congiuntivo passato (abbia fatto, sia partito).
\item\relax Choose between congiuntivo presente and passato according to time relations.
\item\relax Use it after superlatives and in subjunctive contexts referring to the past.
\end{objectives}

\subsection{Formation}\label{les:congiuntivo-passato:formation}

\begin{notebox}{rulecol}\notelabel{rulecol}{Rule}{Congiuntivo presente of the auxiliary + participle}
\textit{(che) io abbia mangiato, (che) lei sia partita}. Auxiliary and agreement as in the passato prossimo.
\end{notebox}

\par\noindent\begin{minipage}[t]{0.48\linewidth}\raggedright
\paradigmbox{partire --- to leave; to set off \textperiodcentered{} congiuntivo passato (with essere, feminine subject)}{%
\begin{tabular}{@{}l l@{}}
\textit{io} & sia partita \\
\textit{tu} & sia partita \\
\textit{lui/lei} & sia partita \\
\textit{noi} & siamo partite \\
\textit{voi} & siate partite \\
\textit{loro} & siano partite \\
\end{tabular}}{}
\end{minipage}\par\medskip

\subsection{Presente or passato?}\label{les:congiuntivo-passato:use}

\begin{notebox}{rulecol}\notelabel{rulecol}{Rule}{Time relative to the main verb}
With a present (or future) main verb: the congiuntivo \textbf{presente} for an action at the same time or later; the congiuntivo \textbf{passato} for an earlier action. \textit{Penso che Marco sia a casa} (now) / \textit{Penso che Marco sia uscito} (earlier).
\end{notebox}

\begin{notebox}{tendcol}\notelabel{tendcol}{Tendency}{Past situations seen from the present}
For a past state or ongoing past situation, the congiuntivo imperfetto is also used after a present main verb: \textit{Penso che da bambino fosse molto timido} (I think he was very shy as a child).
\end{notebox}

\subsection*{Worked examples}
\begin{examples}
\item\relax \textit{Penso che Marco sia già partito.}\newline I think Marco has already left.
\item\relax \textit{Mi dispiace che tu non abbia vinto.}\newline I'm sorry you didn't win.
\item\relax \textit{È il film più bello che io abbia mai visto.}\newline It's the best film I have ever seen.
\item\relax \textit{Non credo che abbiano capito.}\newline I don't think they understood.
\item\relax \textit{Sono contenta che siate venuti.}\newline I'm glad you came.\newline {\color{muted}\small Female speaker; venuti agrees with voi (mixed or masculine group).}
\item\relax \textit{Benché abbia studiato, non ha superato l'esame.}\newline Although he studied, he didn't pass the exam.
\end{examples}

\subsection*{Exercises}\label{les:congiuntivo-passato:exercises}
{\color{muted}\small Answers: page~\pageref{ans:congiuntivo-passato}.}

\begin{exercises}
\item\label{ex:congiuntivo-passato-1}\relax che io (fare), congiuntivo passato: \blank{}
\item\label{ex:congiuntivo-passato-2}\relax Penso che ieri Luca \blank{} (andare) al cinema.\par\nopagebreak\begin{options}
\item\relax andasse
\item\relax è andato
\item\relax vada
\item\relax sia andato
\end{options}
\item\label{ex:congiuntivo-passato-3}\relax Penso che adesso Luca \blank{} (essere) al cinema.\par\nopagebreak\begin{options}
\item\relax sia stato
\item\relax sia
\item\relax è stato
\item\relax fosse
\end{options}
\item\label{ex:congiuntivo-passato-4}\relax Complete with the congiuntivo passato.\par\smallskip\begingroup\setlength{\parindent}{0pt}Sono contento che voi \clozeblank{1} (arrivare) in tempo. È strano che Sara non \clozeblank{2} (telefonare).\par\endgroup
\item\label{ex:congiuntivo-passato-5}\relax Translate: It's the most beautiful book I have ever read.\par\answerline
\end{exercises}

\begin{summarybox}
\begin{itemize}
\item\relax Congiuntivo passato = congiuntivo presente of avere/essere + participle.
\item\relax Present main verb: presente for simultaneous/later, passato for earlier actions.
\item\relax Common after superlatives: il più bello che abbia mai visto.
\end{itemize}
\end{summarybox}

\relatedwords{\textit{partire} to leave; to set off \textperiodcentered{} \textit{andare} to go \textperiodcentered{} \textit{vincere} to win \textperiodcentered{} \textit{capire} to understand \textperiodcentered{} \textit{telefonare} to phone}
\applyit{Read \textit{Una lettera dalla nonna} (A2–B1), \S\ref{story:una-lettera-dalla-nonna}.}

\section{Congiuntivo imperfetto and trapassato}\label{les:congiuntivo-imperfetto-trapassato}
\lessonmeta{B1}{The past subjunctives: after past and conditional main verbs, in hypotheses and in wishes.}
\begin{objectives}
\item\relax Form the congiuntivo imperfetto (parlassi, fossi, facessi) and trapassato (avessi parlato, fossi partito).
\item\relax Use them after a past or conditional main verb.
\item\relax Use them in hypothetical se-clauses and in wishes with magari / se solo / vorrei che.
\end{objectives}

\subsection{Formation}\label{les:congiuntivo-imperfetto-trapassato:formation}

\begin{notebox}{rulecol}\notelabel{rulecol}{Rule}{Imperfetto: stem + -ssi, -ssi, -sse, -ssimo, -ste, -ssero}
Take the imperfetto indicativo, remove -vo, and add the endings: \textit{parlavo → parlassi, vedevo → vedessi, dormivo → dormissi, facevo → facessi, dicevo → dicessi, bevevo → bevessi}. Irregular: \textit{essere (fossi), dare (dessi), stare (stessi)}.
\end{notebox}

\par\noindent\begin{minipage}[t]{0.48\linewidth}\raggedright
\paradigmbox{parlare --- to speak; to talk \textperiodcentered{} congiuntivo imperfetto}{%
\begin{tabular}{@{}l l@{}}
\textit{io} & parl\hlending{assi} \\
\textit{tu} & parl\hlending{assi} \\
\textit{lui/lei} & parl\hlending{asse} \\
\textit{noi} & parl\hlending{assimo} \\
\textit{voi} & parl\hlending{aste} \\
\textit{loro} & parl\hlending{assero} \\
\end{tabular}}{Stem \textit{parl-} + \hlending{ending}.}
\end{minipage}\hfill\begin{minipage}[t]{0.48\linewidth}\raggedright
\paradigmbox{essere --- to be \textperiodcentered{} congiuntivo imperfetto}{%
\begin{tabular}{@{}l l@{}}
\textit{io} & \irr{fossi} \\
\textit{tu} & \irr{fossi} \\
\textit{lui/lei} & \irr{fosse} \\
\textit{noi} & \irr{fossimo} \\
\textit{voi} & \irr{foste} \\
\textit{loro} & \irr{fossero} \\
\end{tabular}}{\irr{Highlighted} forms differ from the regular pattern of the verb's family.}
\end{minipage}\par\medskip

\begin{notebox}{rulecol}\notelabel{rulecol}{Rule}{Trapassato: avessi / fossi + participle}
\textit{avessi mangiato, fossi partito/a}.
\end{notebox}

\par\noindent\begin{minipage}[t]{0.48\linewidth}\raggedright
\paradigmbox{andare --- to go \textperiodcentered{} congiuntivo trapassato (with essere, masculine subject)}{%
\begin{tabular}{@{}l l@{}}
\textit{io} & fossi andato \\
\textit{tu} & fossi andato \\
\textit{lui/lei} & fosse andato \\
\textit{noi} & fossimo andati \\
\textit{voi} & foste andati \\
\textit{loro} & fossero andati \\
\end{tabular}}{}
\end{minipage}\par\medskip

\begin{notebox}{exccol}\notelabel{exccol}{Exception}{voi: -ste, the same as the passato remoto}
The voi forms of congiuntivo imperfetto and passato remoto coincide: \textit{parlaste, credeste, dormiste}. Context decides.
\end{notebox}

\subsection{Use}\label{les:congiuntivo-imperfetto-trapassato:use}

\begin{notebox}{rulecol}\notelabel{rulecol}{Rule}{After a past or conditional main verb}
When the main verb is in a past tense or in the conditional, a congiuntivo clause uses the \textbf{imperfetto} for a simultaneous action and the \textbf{trapassato} for an earlier one: \textit{Pensavo che fosse a casa / che fosse uscito. Vorrei che tu venissi. Avrei voluto che tu fossi venuto.}
\end{notebox}

\begin{notebox}{rulecol}\notelabel{rulecol}{Rule}{Hypotheses}
In se-clauses: \textit{Se avessi tempo, leggerei} (possible); \textit{Se avessi avuto tempo, avrei letto} (unreal past). See the lesson on hypothetical sentences.
\end{notebox}

\begin{notebox}{rulecol}\notelabel{rulecol}{Rule}{Wishes}
\textit{Magari fosse vero!} (If only it were true!) \textit{Se solo avessi ascoltato!} (If only I had listened!) \textit{Fosse vero!}
\end{notebox}

\begin{notebox}{tendcol}\notelabel{tendcol}{Tendency}{come se}
\textit{come se} (as if) always takes the congiuntivo imperfetto or trapassato: \textit{Parla come se fosse il capo.}
\end{notebox}

\subsection*{Worked examples}
\begin{examples}
\item\relax \textit{Pensavo che fossi in vacanza.}\newline I thought you were on holiday.\newline {\color{muted}\small Past main verb, simultaneous: imperfetto.}
\item\relax \textit{Pensavo che fossi già partito.}\newline I thought you had already left.\newline {\color{muted}\small Earlier action: trapassato.}
\item\relax \textit{Vorrei che tu restassi.}\newline I'd like you to stay.\newline {\color{muted}\small Conditional main verb: imperfetto.}
\item\relax \textit{Magari avessi più tempo!}\newline If only I had more time!
\item\relax \textit{Parla come se sapesse tutto.}\newline He talks as if he knew everything.
\item\relax \textit{Se avessi saputo, sarei venuta.}\newline If I had known, I would have come.\newline {\color{muted}\small Counterfactual past: trapassato + condizionale passato.}
\end{examples}

\subsection*{Exercises}\label{les:congiuntivo-imperfetto-trapassato:exercises}
{\color{muted}\small Answers: page~\pageref{ans:congiuntivo-imperfetto-trapassato}.}

\begin{exercises}
\item\label{ex:congiuntivo-imperfetto-trapassato-1}\relax Give the congiuntivo imperfetto (che io …).\par\smallskip{\begin{tabular}{@{}l l@{}}
parlare & \blankcell \\
avere & \blankcell \\
essere & \blankcell \\
fare & \blankcell \\
dire & \blankcell \\
stare & \blankcell \\
\end{tabular}}
\item\label{ex:congiuntivo-imperfetto-trapassato-2}\relax Credevo che Marco \blank{} (lavorare) a Roma.\par\nopagebreak\begin{options}
\item\relax lavorerebbe
\item\relax abbia lavorato
\item\relax lavorasse
\item\relax lavori
\end{options}
\item\label{ex:congiuntivo-imperfetto-trapassato-3}\relax Non sapevo che voi \blank{} (partire).\par\nopagebreak\begin{options}
\item\relax siate partiti
\item\relax foste partiti
\item\relax partireste
\item\relax partiste
\end{options}
\item\label{ex:congiuntivo-imperfetto-trapassato-4}\relax Complete.\par\smallskip\begingroup\setlength{\parindent}{0pt}Magari \clozeblank{1} (essere) vero! Vorrei che tu \clozeblank{2} (venire) con me. Parla come se \clozeblank{3} (sapere) tutto.\par\endgroup
\item\label{ex:congiuntivo-imperfetto-trapassato-5}\relax Translate: I thought you had already left. (tu, male)\par\answerline
\item\label{ex:congiuntivo-imperfetto-trapassato-6}\relax \textit{parlaste} can be…\par\nopagebreak\begin{options}
\item\relax only congiuntivo imperfetto
\item\relax only passato remoto
\item\relax condizionale presente
\item\relax passato remoto or congiuntivo imperfetto (voi)
\end{options}
\end{exercises}

\begin{summarybox}
\begin{itemize}
\item\relax Congiuntivo imperfetto = imperfetto stem + -ssi, -ssi, -sse, -ssimo, -ste, -ssero (fossi, dessi, stessi irregular).
\item\relax Congiuntivo trapassato = avessi/fossi + participle.
\item\relax After past or conditional main verbs: imperfetto (simultaneous), trapassato (earlier).
\item\relax Also in se-clauses of hypotheses, in wishes (magari…) and after come se.
\end{itemize}
\end{summarybox}

\relatedwords{\textit{pensare} to think \textperiodcentered{} \textit{credere} to believe \textperiodcentered{} \textit{volere} to want \textperiodcentered{} \textit{sapere} to know (a fact); to know how to \textperiodcentered{} \textit{essere} to be \textperiodcentered{} \textit{partire} to leave; to set off}
\applyit{Read \textit{Una lettera dalla nonna} (A2–B1), \S\ref{story:una-lettera-dalla-nonna}.}

\chapter{Complex sentences}\label{mod:m19}
\moduleintro{B1–B2}{The periodo ipotetico, the sequence of tenses (concordanza dei tempi) and indirect speech.}{Hypothetical sentences (periodo ipotetico) (\S\ref{les:periodo-ipotetico}); Sequence of tenses (concordanza dei tempi) (\S\ref{les:concordanza-dei-tempi}); Reported speech (discorso indiretto) (\S\ref{les:discorso-indiretto})}

\section{Hypothetical sentences (periodo ipotetico)}\label{les:periodo-ipotetico}
\lessonmeta{B1}{Real, possible and unreal conditions; mixed types; the colloquial double imperfetto; what is not standard.}
\begin{objectives}
\item\relax Build the three standard types of hypothetical sentence.
\item\relax Combine a past condition with a present consequence (mixed type).
\item\relax Recognise the colloquial 'se lo sapevo, venivo' and know its register.
\item\relax Avoid the nonstandard conditional in the se-clause.
\end{objectives}

\subsection{The three standard types}\label{les:periodo-ipotetico:types}

{\small\setlength{\LTpre}{2pt}\setlength{\LTpost}{6pt}
\tablecaption{Periodo ipotetico}
\begin{xltabular}{\linewidth}{lllL}
\toprule
\textbf{Type} & \textbf{se-clause (protasi)} & \textbf{Main clause (apodosi)} & \textbf{Example} \\
\midrule
realtà (real) & indicative & indicative or imperative & \textit{Se piove, resto a casa. Se hai fame, mangia!} \\
possibilità (possible) & congiuntivo imperfetto & condizionale presente & \textit{Se avessi tempo, leggerei di più.} \\
irrealtà (unreal, past) & congiuntivo trapassato & condizionale passato & \textit{Se avessi studiato, avrei superato l'esame.} \\
\bottomrule
\end{xltabular}}

\begin{notebox}{rulecol}\notelabel{rulecol}{Rule}{Future reference in the real type}
With future reference both clauses may be in the futuro: \textit{Se verrai, sarò felice}, or in the present: \textit{Se vieni, sono felice}.
\end{notebox}

\begin{notebox}{tendcol}\notelabel{tendcol}{Tendency}{Possible or unreal in the present?}
The congiuntivo imperfetto + condizionale presente covers both what is possible (\textit{se vincessi alla lotteria…}) and what is contrary to present fact (\textit{se fossi in te…} — I am not you). The context makes the difference.
\end{notebox}

\subsection{Mixed types}\label{les:periodo-ipotetico:mixed}

\begin{itemize}
\item\relax Past condition, present consequence: \textit{Se avessi studiato, ora sarei più tranquillo.} (trapassato + condizionale presente)
\item\relax Permanent condition, past consequence: \textit{Se fossi più paziente, non ti saresti arrabbiato.} (imperfetto + condizionale passato)
\end{itemize}

\subsection{Register: what is colloquial and what is not standard}\label{les:periodo-ipotetico:register}

\begin{notebox}{tendcol}\notelabel{tendcol}{Tendency}{The double imperfetto}
In conversation an unreal past condition is very often expressed with the imperfetto indicativo in both clauses: \textit{Se lo sapevo, venivo} (= \textit{se l'avessi saputo, sarei venuto}). It has a long history (it appears in Dante, Machiavelli and Manzoni) and is understood everywhere, but it belongs to an informal register; in writing and careful speech use the congiuntivo and condizionale (Treccani, \textit{Enciclopedia dell'italiano}, s.v. periodo ipotetico).
\end{notebox}

\begin{notebox}{rulecol}\notelabel{rulecol}{Rule}{Not standard: conditional after se}
\textit{Se avrei tempo, leggerei} is not standard Italian. It is common in the speech of less educated speakers and is widely stigmatised. Use \textit{se avessi tempo}.
\end{notebox}

\begin{notebox}{rulecol}\notelabel{rulecol}{Rule}{se meaning 'whether'}
In indirect questions \textit{se} means 'whether' and is not a hypothesis: \textit{Non so se verrà} (I don't know whether he'll come). Here a future or conditional after se is normal: \textit{Mi chiedo se sarebbe d'accordo.}
\end{notebox}

\subsection*{Worked examples}
\begin{examples}
\item\relax \textit{Se piove, restiamo a casa.}\newline If it rains, we'll stay at home.\newline {\color{muted}\small Real: indicative + indicative.}
\item\relax \textit{Se avessi una casa al mare, ci andrei ogni estate.}\newline If I had a house by the sea, I'd go there every summer.
\item\relax \textit{Se fossi in te, non lo farei.}\newline If I were you, I wouldn't do it.
\item\relax \textit{Se avessimo preso il treno delle nove, saremmo già arrivati.}\newline If we had taken the nine o'clock train, we would already have arrived.
\item\relax \textit{Se avessi dormito di più, oggi non sarei così stanca.}\newline If I had slept more, I wouldn't be so tired today.\newline {\color{muted}\small Mixed: past condition, present consequence.}
\item\relax \textit{Se lo sapevo, venivo.}\newline If I'd known, I would have come.\newline {\color{muted}\small Colloquial double imperfetto.}
\item\relax \textit{Non so se verrà.}\newline I don't know whether he'll come.\newline {\color{muted}\small se = whether (indirect question): futuro is normal.}
\end{examples}

\subsection*{Exercises}\label{les:periodo-ipotetico:exercises}
{\color{muted}\small Answers: page~\pageref{ans:periodo-ipotetico}.}

\begin{exercises}
\item\label{ex:periodo-ipotetico-1}\relax Se \blank{} (avere) più soldi, comprerei una barca.\par\nopagebreak\begin{options}
\item\relax avevo
\item\relax avessi
\item\relax avrei
\item\relax ho
\end{options}
\item\label{ex:periodo-ipotetico-2}\relax Se avessi studiato, \blank{} (superare) l'esame.\par\nopagebreak\begin{options}
\item\relax avrei superato
\item\relax avessi superato
\item\relax supererei
\item\relax superavo
\end{options}
\item\label{ex:periodo-ipotetico-3}\relax Complete.\par\smallskip\begingroup\setlength{\parindent}{0pt}Se \clozeblank{1} (essere, io) in te, non \clozeblank{2} (partire). Se \clozeblank{3} (piovere) domani, restiamo a casa.\par\endgroup
\item\label{ex:periodo-ipotetico-4}\relax Classify each sentence.\par\smallskip\begin{tabularx}{\linewidth}{@{}L L@{}}
1.\ Se hai fame, mangia. & (a)\ colloquial \\
2.\ Se potessi, verrei. & (b)\ irrealtà (past) \\
3.\ Se fossi venuto, ti saresti divertito. & (c)\ realtà \\
4.\ Se l'avessi saputo, ora sarei lì. & (d)\ possibilità \\
5.\ Se lo sapevo, venivo. & (e)\ mixed \\
\end{tabularx}
\item\label{ex:periodo-ipotetico-5}\relax Translate (standard written Italian): If I had known, I would have come. (male speaker)\par\answerline
\item\label{ex:periodo-ipotetico-6}\relax Which sentence is correct standard Italian?\par\nopagebreak\begin{options}
\item\relax Se saresti qui, lo vedresti.
\item\relax Mi chiedo se Marco verrebbe.
\item\relax Se Marco verrebbe, sarei contento.
\item\relax Se avrei tempo, verrei.
\end{options}
\end{exercises}

\begin{summarybox}
\begin{itemize}
\item\relax realtà: se + indicativo, indicativo/imperativo.
\item\relax possibilità: se + congiuntivo imperfetto, condizionale presente.
\item\relax irrealtà: se + congiuntivo trapassato, condizionale passato; mixed types combine them.
\item\relax Colloquial: se + imperfetto, imperfetto. Not standard: se + condizionale.
\item\relax se = whether (indirect question) can take futuro or condizionale.
\end{itemize}
\end{summarybox}

\relatedwords{\textit{se} if; whether \textperiodcentered{} \textit{avere} to have \textperiodcentered{} \textit{essere} to be \textperiodcentered{} \textit{sapere} to know (a fact); to know how to \textperiodcentered{} \textit{venire} to come \textperiodcentered{} \textit{piovere} to rain}
\applyit{Read \textit{Una lettera dalla nonna} (A2–B1), \S\ref{story:una-lettera-dalla-nonna}.}

\section{Sequence of tenses (concordanza dei tempi)}\label{les:concordanza-dei-tempi}
\lessonmeta{B2}{How the tense of a subordinate clause depends on the main verb and on the time relation between the two actions.}
\begin{objectives}
\item\relax Choose the subordinate tense for simultaneous, earlier and later actions after a present main verb.
\item\relax Do the same after a past main verb, in both indicative and congiuntivo clauses.
\item\relax Use the condizionale passato for posteriority in the past.
\end{objectives}

\subsection{Clauses with the congiuntivo}\label{les:concordanza-dei-tempi:subjunctive}

{\small\setlength{\LTpre}{2pt}\setlength{\LTpost}{6pt}
\tablecaption{After a verb that requires the congiuntivo (pensare, credere, sperare, temere…)}
\begin{xltabular}{\linewidth}{Llll}
\toprule
\textbf{Main verb} & \textbf{Earlier} & \textbf{Simultaneous} & \textbf{Later} \\
\midrule
present / future: \textit{Penso che…} & \textit{sia partito} & \textit{parta} & \textit{parta / partirà} \\
past: \textit{Pensavo che…} & \textit{fosse partito} & \textit{partisse} & \textit{sarebbe partito} \\
conditional: \textit{Vorrei che…} & — & \textit{partisse} & \textit{partisse} \\
conditional past: \textit{Avrei voluto che…} & \textit{fosse partito} & \textit{partisse} & — \\
\bottomrule
\end{xltabular}}

\begin{notebox}{rulecol}\notelabel{rulecol}{Rule}{Posteriority in the past}
After a past main verb, a later action is expressed with the \textbf{condizionale passato} in both congiuntivo and indicative contexts: \textit{Pensavo che sarebbe partito. Sapevo che sarebbe partito.}
\end{notebox}

\begin{notebox}{rulecol}\notelabel{rulecol}{Rule}{Volition after the conditional}
Verbs of wish in the conditional (\textit{vorrei, preferirei, mi piacerebbe}) take the congiuntivo imperfetto: \textit{Vorrei che tu venissi}, not \textit{che tu venga}.
\end{notebox}

\subsection{Clauses with the indicative}\label{les:concordanza-dei-tempi:indicative}

{\small\setlength{\LTpre}{2pt}\setlength{\LTpost}{6pt}
\tablecaption{After verbs of saying and knowing (dire, sapere, scrivere…)}
\begin{xltabular}{\linewidth}{Llll}
\toprule
\textbf{Main verb} & \textbf{Earlier} & \textbf{Simultaneous} & \textbf{Later} \\
\midrule
present: \textit{Dice che…} & \textit{è partito / partì} & \textit{parte} & \textit{partirà} \\
past: \textit{Ha detto che… / Disse che…} & \textit{era partito} & \textit{partiva} & \textit{sarebbe partito} \\
\bottomrule
\end{xltabular}}

\begin{notebox}{tendcol}\notelabel{tendcol}{Tendency}{When the content is still true}
After a past main verb, if the reported fact is still true or still future at the moment of speaking, the present or futuro can be kept: \textit{Mi ha detto che è sposato. Ha detto che verrà domani.} The shifted forms (\textit{era sposato, sarebbe venuto}) are always correct.
\end{notebox}

\subsection*{Worked examples}
\begin{examples}
\item\relax \textit{Penso che Luca sia già arrivato.}\newline I think Luca has already arrived.\newline {\color{muted}\small Present main verb, earlier action.}
\item\relax \textit{Pensavo che Luca fosse già arrivato.}\newline I thought Luca had already arrived.\newline {\color{muted}\small Past main verb, earlier action.}
\item\relax \textit{Credevo che piovesse.}\newline I thought it was raining.\newline {\color{muted}\small Past main verb, simultaneous.}
\item\relax \textit{Ero sicuro che avrebbe piovuto.}\newline I was sure it would rain.\newline {\color{muted}\small Past main verb, later: condizionale passato.}
\item\relax \textit{Vorrei che tu mi aiutassi.}\newline I'd like you to help me.\newline {\color{muted}\small Conditional main verb: congiuntivo imperfetto.}
\item\relax \textit{Mi ha detto che sarebbe tornato presto.}\newline He told me he would come back soon.
\item\relax \textit{Mi ha detto che è sposato.}\newline He told me he is married.\newline {\color{muted}\small Still true now: the present is kept.}
\end{examples}

\subsection*{Exercises}\label{les:concordanza-dei-tempi:exercises}
{\color{muted}\small Answers: page~\pageref{ans:concordanza-dei-tempi}.}

\begin{exercises}
\item\label{ex:concordanza-dei-tempi-1}\relax Complete \textit{Pensavo che Marco … (partire)} for each time relation.\par\smallskip{\begin{tabular}{@{}l l@{}}
earlier & \blankcell \\
simultaneous & \blankcell \\
later & \blankcell \\
\end{tabular}}
\item\label{ex:concordanza-dei-tempi-2}\relax Complete \textit{Penso che Marco … (partire)} for each time relation.\par\smallskip{\begin{tabular}{@{}l l@{}}
earlier & \blankcell \\
simultaneous & \blankcell \\
later & \blankcell \\
\end{tabular}}
\item\label{ex:concordanza-dei-tempi-3}\relax Vorrei che voi \blank{} (restare) ancora un po'.\par\nopagebreak\begin{options}
\item\relax siete restati
\item\relax restereste
\item\relax restaste
\item\relax restiate
\end{options}
\item\label{ex:concordanza-dei-tempi-4}\relax Sapevo che Sara \blank{} (vincere): era la più veloce. (later than 'sapevo')\par\nopagebreak\begin{options}
\item\relax vincerebbe
\item\relax vincesse
\item\relax avrebbe vinto
\item\relax vincerà
\end{options}
\item\label{ex:concordanza-dei-tempi-5}\relax Complete after a past verb of saying.\par\smallskip\begingroup\setlength{\parindent}{0pt}Ha detto che \clozeblank{1} (essere) stanco (simultaneous) e che \clozeblank{2} (dormire) poco la notte prima (earlier).\par\endgroup
\item\label{ex:concordanza-dei-tempi-6}\relax \textit{Mi ha detto che è sposato.} Is this correct?\par\nopagebreak\begin{options}
\item\relax Yes: the fact is still true now, so the present can be kept.
\item\relax No: it must be 'fosse sposato'.
\item\relax No: it must be 'era sposato'.
\item\relax No: it must be 'sarebbe sposato'.
\end{options}
\end{exercises}

\begin{summarybox}
\begin{itemize}
\item\relax Present main verb + congiuntivo: passato (earlier), presente (simultaneous), presente/futuro (later).
\item\relax Past main verb + congiuntivo: trapassato (earlier), imperfetto (simultaneous), condizionale passato (later).
\item\relax Vorrei che + congiuntivo imperfetto.
\item\relax Indicative after a past verb of saying: trapassato prossimo / imperfetto / condizionale passato; the present can be kept for facts still true.
\end{itemize}
\end{summarybox}

\relatedwords{\textit{pensare} to think \textperiodcentered{} \textit{credere} to believe \textperiodcentered{} \textit{dire} to say; to tell \textperiodcentered{} \textit{sapere} to know (a fact); to know how to \textperiodcentered{} \textit{volere} to want \textperiodcentered{} \textit{partire} to leave; to set off}
\applyit{Read \textit{Una lettera dalla nonna} (A2–B1), \S\ref{story:una-lettera-dalla-nonna}. Read \textit{Il segreto del faro} (A2–B1), \S\ref{story:il-segreto-del-faro}.}

\section{Reported speech (discorso indiretto)}\label{les:discorso-indiretto}
\lessonmeta{B1}{Turning direct speech into reported speech: tenses, persons, time and place words, questions and commands.}
\begin{objectives}
\item\relax Shift tenses after a past verb of saying.
\item\relax Change persons, possessives and time/place expressions consistently.
\item\relax Report questions (se, che cosa, dove…) and commands (di + infinitive or che + congiuntivo).
\end{objectives}

\subsection{Tense changes after a past verb}\label{les:discorso-indiretto:tenses}

{\small\setlength{\LTpre}{2pt}\setlength{\LTpost}{6pt}
\tablecaption{Direct → reported (after disse / ha detto)}
\begin{xltabular}{\linewidth}{LL}
\toprule
\textbf{Direct speech} & \textbf{Reported speech} \\
\midrule
presente: «Sono stanco.» & imperfetto: disse che era stanco \\
passato prossimo / remoto: «Ho perso le chiavi.» & trapassato prossimo: disse che aveva perso le chiavi \\
imperfetto: «Ero felice.» & imperfetto (unchanged): disse che era felice \\
futuro: «Partirò.» & condizionale passato: disse che sarebbe partito \\
condizionale presente: «Vorrei un tè.» & condizionale passato: disse che avrebbe voluto un tè \\
imperativo: «Chiudi la porta!» & di + infinitive: gli disse di chiudere la porta \\
congiuntivo presente: «Spero che venga.» & congiuntivo imperfetto: disse che sperava che venisse \\
\bottomrule
\end{xltabular}}

\begin{notebox}{rulecol}\notelabel{rulecol}{Rule}{After a present verb nothing shifts}
\textit{Dice che è stanco. Dice che partirà.} Only persons and deictic words change.
\end{notebox}

\subsection{Persons, time and place}\label{les:discorso-indiretto:deixis}

{\small\setlength{\LTpre}{2pt}\setlength{\LTpost}{6pt}
\tablecaption{Typical changes}
\begin{xltabular}{\linewidth}{lL}
\toprule
\textbf{Direct} & \textbf{Reported} \\
\midrule
io, mio & lui/lei, suo \\
oggi & quel giorno \\
ieri & il giorno prima \\
domani & il giorno dopo / l'indomani \\
adesso, ora & allora, in quel momento \\
qui, qua & lì, là \\
questo & quello \\
fra due giorni & dopo due giorni \\
venire & andare (often) \\
\bottomrule
\end{xltabular}}

\subsection{Questions and commands}\label{les:discorso-indiretto:questions}

\begin{notebox}{rulecol}\notelabel{rulecol}{Rule}{Yes/no questions: se}
«Hai fame?» → \textit{Mi chiese se avevo fame.} In indirect questions the congiuntivo is also possible and more formal: \textit{mi chiese se avessi fame}.
\end{notebox}

\begin{notebox}{rulecol}\notelabel{rulecol}{Rule}{Open questions keep the question word}
«Dove abiti?» → \textit{Mi chiese dove abitavo / abitassi.}
\end{notebox}

\begin{notebox}{rulecol}\notelabel{rulecol}{Rule}{Commands}
di + infinitive (same reference) or che + congiuntivo: «Vieni qui!» → \textit{Mi disse di andare lì} / \textit{Disse che andassi lì}.
\end{notebox}

\subsection*{Worked examples}
\begin{examples}
\item\relax \textit{«Sono stanca», disse Sara. → Sara disse che era stanca.}\newline “I'm tired,” said Sara. → Sara said she was tired.
\item\relax \textit{«Domani partirò.» → Disse che il giorno dopo sarebbe partito.}\newline “I'll leave tomorrow.” → He said he would leave the next day.
\item\relax \textit{«Ho perso la chiave.» → Disse che aveva perso la chiave.}\newline “I've lost the key.” → He said he had lost the key.
\item\relax \textit{«Hai fame?» → Mi chiese se avevo fame.}\newline “Are you hungry?” → He asked me if I was hungry.
\item\relax \textit{«Chiudi la finestra!» → Mi disse di chiudere la finestra.}\newline “Close the window!” → She told me to close the window.
\item\relax \textit{«Verrò qui domani.» → Disse che sarebbe andato lì il giorno dopo.}\newline “I'll come here tomorrow.” → He said he would go there the next day.\newline {\color{muted}\small venire → andare, qui → lì, domani → il giorno dopo.}
\end{examples}

\subsection*{Exercises}\label{les:discorso-indiretto:exercises}
{\color{muted}\small Answers: page~\pageref{ans:discorso-indiretto}.}

\begin{exercises}
\item\label{ex:discorso-indiretto-1}\relax «Ho fame.» → Luca disse che \blank{} fame.
\item\label{ex:discorso-indiretto-2}\relax «Partirò domani.» → Disse che \blank{} il giorno dopo. (male speaker)
\item\label{ex:discorso-indiretto-3}\relax «Ho visto un lupo.» → Disse che \blank{} un lupo.
\item\label{ex:discorso-indiretto-4}\relax Match direct and reported expressions.\par\smallskip\begin{tabularx}{\linewidth}{@{}L L@{}}
1.\ oggi & (a)\ lì \\
2.\ domani & (b)\ quel giorno \\
3.\ ieri & (c)\ quello \\
4.\ qui & (d)\ il giorno prima \\
5.\ questo & (e)\ il giorno dopo \\
\end{tabularx}
\item\label{ex:discorso-indiretto-5}\relax Translate: \textit{Mi chiese se avevo fame.}\par\answerline
\item\label{ex:discorso-indiretto-6}\relax «Chiudi la porta!» → La mamma mi disse \blank{} la porta.\par\nopagebreak\begin{options}
\item\relax di chiudere
\item\relax che chiuderò
\item\relax chiudere
\item\relax che chiudo
\end{options}
\end{exercises}

\begin{summarybox}
\begin{itemize}
\item\relax After a past verb: presente → imperfetto, passato → trapassato, futuro/condizionale → condizionale passato, imperativo → di + infinitive.
\item\relax Persons and deictics change: oggi → quel giorno, domani → il giorno dopo, qui → lì, venire → andare.
\item\relax Questions: se / question word, indicative or (more formal) congiuntivo.
\end{itemize}
\end{summarybox}

\relatedwords{\textit{dire} to say; to tell \textperiodcentered{} \textit{chiedere} to ask \textperiodcentered{} \textit{rispondere} to answer; to reply \textperiodcentered{} \textit{raccontare} to tell (a story) \textperiodcentered{} \textit{promettere} to promise}
\applyit{Read \textit{Il segreto del faro} (A2–B1), \S\ref{story:il-segreto-del-faro}. Read \textit{Una lettera dalla nonna} (A2–B1), \S\ref{story:una-lettera-dalla-nonna}.}

\chapter{Passive and impersonal constructions}\label{mod:m20}
\moduleintro{B1}{The passive with essere, venire and andare; si passivante and si impersonale.}{The passive voice (\S\ref{les:forma-passiva}); Impersonal si (\S\ref{les:si-impersonale})}

\section{The passive voice}\label{les:forma-passiva}
\lessonmeta{B1}{Passive with essere, venire and andare; the agent with da; the si passivante.}
\begin{objectives}
\item\relax Form the passive with essere + participle in any tense, with agreement.
\item\relax Use venire + participle in simple tenses and andare + participle for obligation.
\item\relax Express the agent with da.
\item\relax Use the si passivante (si vendono case).
\end{objectives}

\subsection{essere + participle}\label{les:forma-passiva:essere}

\begin{notebox}{rulecol}\notelabel{rulecol}{Rule}{Only transitive verbs}
Only verbs that take a direct object can be made passive. The object becomes the subject; the participle agrees with it; the tense is carried by essere: \textit{Il gatto mangia la torta → La torta è mangiata dal gatto}. \textit{La torta è stata mangiata} (passato prossimo), \textit{sarà mangiata} (futuro), \textit{fu mangiata} (passato remoto).
\end{notebox}

\begin{notebox}{rulecol}\notelabel{rulecol}{Rule}{The agent: da}
Who does the action is introduced by \textit{da} (with the article where needed): \textit{dal gatto, dalla maestra, da Marco}.
\end{notebox}

\begin{notebox}{tendcol}\notelabel{tendcol}{Tendency}{State or action?}
\textit{La porta è chiusa} can describe a state (the door is closed) or an action (the door is being closed). Context — or the choice of venire — decides.
\end{notebox}

\subsection{venire and andare}\label{les:forma-passiva:venire-andare}

\begin{notebox}{rulecol}\notelabel{rulecol}{Rule}{venire + participle}
\textit{venire} can replace \textit{essere} in \textbf{simple} tenses only and makes the action, not the state, clear: \textit{La porta viene chiusa alle otto} (is closed every day at eight). Not in compound tenses: \textit{è stata chiusa}, never \textit{è venuta chiusa}.
\end{notebox}

\begin{notebox}{rulecol}\notelabel{rulecol}{Rule}{andare + participle}
In simple tenses \textit{andare} + participle expresses obligation: \textit{Il modulo va compilato in stampatello} (must be filled in). With verbs of losing it is simply passive: \textit{La lettera è andata persa.}
\end{notebox}

\subsection{The si passivante}\label{les:forma-passiva:si-passivante}

\begin{notebox}{rulecol}\notelabel{rulecol}{Rule}{si + third person, agreeing with the subject}
\textit{Si vende la casa} (the house is for sale), \textit{Si vendono case}. The noun is the grammatical subject, so the verb is plural with a plural noun. In compound tenses it takes essere and agreement: \textit{Si sono vendute molte case.}
\end{notebox}

\begin{notebox}{tendcol}\notelabel{tendcol}{Tendency}{Very common in everyday language}
Signs, recipes and general statements use it constantly: \textit{Qui si parla italiano. Si cuociono le patate per venti minuti.}
\end{notebox}

\subsection*{Worked examples}
\begin{examples}
\item\relax \textit{La torta è stata mangiata dal gatto.}\newline The cake was eaten by the cat.\newline {\color{muted}\small Passato prossimo passive.}
\item\relax \textit{Il castello fu costruito nel Trecento.}\newline The castle was built in the fourteenth century.\newline {\color{muted}\small Passato remoto passive.}
\item\relax \textit{La posta viene consegnata ogni mattina.}\newline The post is delivered every morning.\newline {\color{muted}\small venire stresses the action.}
\item\relax \textit{Questo libro va letto con calma.}\newline This book must be read slowly.\newline {\color{muted}\small andare = obligation.}
\item\relax \textit{In questo negozio si vendono libri usati.}\newline Second-hand books are sold in this shop.\newline {\color{muted}\small si passivante: plural verb.}
\item\relax \textit{La notizia è stata annunciata ieri.}\newline The news was announced yesterday.
\end{examples}

\subsection*{Exercises}\label{les:forma-passiva:exercises}
{\color{muted}\small Answers: page~\pageref{ans:forma-passiva}.}

\begin{exercises}
\item\label{ex:forma-passiva-1}\relax Il nonno racconta la storia. → La storia \blank{} dal nonno.
\item\label{ex:forma-passiva-2}\relax I bambini hanno mangiato le mele. → Le mele \blank{} dai bambini.
\item\label{ex:forma-passiva-3}\relax In questo negozio \blank{} libri usati.\par\nopagebreak\begin{options}
\item\relax si venduti
\item\relax si vende
\item\relax si vendono
\item\relax vendono si
\end{options}
\item\label{ex:forma-passiva-4}\relax \textit{Il modulo va firmato.} What does it mean?\par\nopagebreak\begin{options}
\item\relax The form was signed.
\item\relax The form has gone.
\item\relax The form is going to be signed.
\item\relax The form must be signed.
\end{options}
\item\label{ex:forma-passiva-5}\relax Which verb cannot be made passive?\par\nopagebreak\begin{options}
\item\relax mangiare
\item\relax leggere
\item\relax arrivare
\item\relax costruire
\end{options}
\item\label{ex:forma-passiva-6}\relax Translate: The castle was built in 1300. (passato remoto)\par\answerline
\end{exercises}

\begin{summarybox}
\begin{itemize}
\item\relax Passive = essere (any tense) + participle agreeing with the subject; agent with da.
\item\relax venire + participle in simple tenses (action); andare + participle = obligation.
\item\relax si passivante: si vende la casa / si vendono le case; compound tenses with essere and agreement.
\end{itemize}
\end{summarybox}

\relatedwords{\textit{costruire} to build \textperiodcentered{} \textit{mangiare} to eat \textperiodcentered{} \textit{vendere} to sell \textperiodcentered{} \textit{annunciare} to announce \textperiodcentered{} \textit{raccontare} to tell (a story) \textperiodcentered{} \textit{scrivere} to write}
\applyit{Read \textit{Il segreto del faro} (A2–B1), \S\ref{story:il-segreto-del-faro}.}

\section{Impersonal si}\label{les:si-impersonale}
\lessonmeta{B1}{si + verb meaning 'one, people, you, we': agreement, reflexive verbs (ci si) and compound tenses.}
\begin{objectives}
\item\relax Use impersonal si with intransitive verbs (si vive bene).
\item\relax Make predicate adjectives plural after impersonal si (si è stanchi).
\item\relax Use ci si with reflexive verbs (ci si alza presto).
\item\relax Distinguish impersonal si from the si passivante.
\end{objectives}

\subsection{Basics}\label{les:si-impersonale:basics}

\begin{notebox}{rulecol}\notelabel{rulecol}{Rule}{si + third person singular}
Impersonal si refers to people in general: \textit{In Italia si mangia bene. Qui si sta tranquilli. Come si dice 'cat' in italiano?} English uses 'one', 'you', 'we' or 'people'.
\end{notebox}

\begin{notebox}{rulecol}\notelabel{rulecol}{Rule}{Plural agreement of adjectives}
Although the verb is singular, a predicate adjective or noun is plural (masculine by default): \textit{Quando si è stanchi, si dorme bene. Da giovani si è più coraggiosi.}
\end{notebox}

\begin{notebox}{rulecol}\notelabel{rulecol}{Rule}{Reflexive verbs: ci si}
To avoid \textit{si si}, reflexive verbs use \textit{ci si}: \textit{La mattina ci si alza presto. Ci si diverte molto.}
\end{notebox}

\subsection{Compound tenses}\label{les:si-impersonale:compound}

\begin{notebox}{rulecol}\notelabel{rulecol}{Rule}{Always essere}
Impersonal si always takes \textit{essere}. With verbs that normally take essere, the participle is plural: \textit{Si è arrivati tardi.} With verbs that normally take avere, the participle stays singular: \textit{Si è mangiato bene.}
\end{notebox}

\begin{notebox}{tendcol}\notelabel{tendcol}{Tendency}{Impersonal or passive?}
With a transitive verb and a noun, si is passivante and the verb agrees: \textit{si mangiano le mele}. Without a direct object, si is impersonal and the verb is singular: \textit{si mangia bene}. In practice the two are close in meaning.
\end{notebox}

\begin{notebox}{tendcol}\notelabel{tendcol}{Tendency}{noi in Tuscan speech}
In Tuscany \textit{si} + singular verb is also used in speech to mean 'we': \textit{noi si va al mare} = \textit{noi andiamo al mare}. This is regional, not standard.
\end{notebox}

\subsection*{Worked examples}
\begin{examples}
\item\relax \textit{In questa città si vive bene.}\newline People live well in this city. / Life is good in this city.
\item\relax \textit{Come si dice «cat» in italiano?}\newline How do you say 'cat' in Italian?
\item\relax \textit{Quando si è stanchi, si dorme bene.}\newline When you're tired, you sleep well.\newline {\color{muted}\small stanchi is plural.}
\item\relax \textit{In vacanza ci si alza tardi.}\newline On holiday people get up late.\newline {\color{muted}\small ci si with a reflexive verb.}
\item\relax \textit{Ieri si è mangiato benissimo.}\newline We ate extremely well yesterday.\newline {\color{muted}\small avere-verb: participle singular.}
\item\relax \textit{Si è arrivati tardi alla festa.}\newline We got to the party late.\newline {\color{muted}\small essere-verb: participle plural.}
\end{examples}

\subsection*{Exercises}\label{les:si-impersonale:exercises}
{\color{muted}\small Answers: page~\pageref{ans:si-impersonale}.}

\begin{exercises}
\item\label{ex:si-impersonale-1}\relax Quando \blank{} giovani, si ha molta energia.\par\nopagebreak\begin{options}
\item\relax si sono
\item\relax si siamo
\item\relax ci si è
\item\relax si è
\end{options}
\item\label{ex:si-impersonale-2}\relax In montagna \blank{} presto. (alzarsi)\par\nopagebreak\begin{options}
\item\relax si alzano
\item\relax si alza
\item\relax si si alza
\item\relax ci si alza
\end{options}
\item\label{ex:si-impersonale-3}\relax Quando si è \blank{} (stanco), si dorme bene.
\item\label{ex:si-impersonale-4}\relax Ieri \blank{} tardi. (arrivare, impersonal)\par\nopagebreak\begin{options}
\item\relax si sono arrivato
\item\relax si ha arrivato
\item\relax si è arrivati
\item\relax si è arrivato
\end{options}
\item\label{ex:si-impersonale-5}\relax Impersonal or passivante?\par\smallskip\begin{tabularx}{\linewidth}{@{}L L@{}}
1.\ Si mangia bene qui. & (a)\ impersonal \\
2.\ Si mangiano molte mele. & (b)\ passivante \\
3.\ Si vive bene. & (c)\ passivante \\
4.\ Si vende la casa. & (d)\ impersonal \\
\end{tabularx}
\item\label{ex:si-impersonale-6}\relax Translate: How do you say 'cat' in Italian?\par\answerline
\end{exercises}

\begin{summarybox}
\begin{itemize}
\item\relax Impersonal si + 3rd person singular: si vive bene.
\item\relax Predicate adjectives plural: si è stanchi.
\item\relax Reflexive verbs: ci si alza.
\item\relax Compound tenses with essere: si è mangiato (avere-verb) / si è arrivati (essere-verb).
\item\relax Tuscan 'noi si va' is regional.
\end{itemize}
\end{summarybox}

\relatedwords{\textit{vivere} to live \textperiodcentered{} \textit{dire} to say; to tell \textperiodcentered{} \textit{alzarsi} to get up \textperiodcentered{} \textit{divertirsi} to have fun; to enjoy oneself \textperiodcentered{} \textit{arrivare} to arrive \textperiodcentered{} \textit{stanco} tired}

\part{Reader}\label{part:reader}
\chapter{Stories}\label{stories}
The stories are graded. Read the text first without help; then use the sentence-by-sentence translation, notes and verb analysis. The verb analysis gives, for each verb group, the tense, mood, person and the lesson that explains it --- including tenses you have not studied yet. Where a sentence cannot be analysed clause by clause, every possible reading of each verb form is listed instead.\par

\section{Il gatto e la luna}\label{story:il-gatto-e-la-luna}
\storymeta{The cat and the moon}{A0–A1}{A little black cat tries to talk to the moon.}
\subsection*{Text}
\begin{storytext}
\sentnum{1}Il gatto si chiama Tobia. \sentnum{2}Tobia è un gatto nero e piccolo. \sentnum{3}Di giorno Tobia dorme sul letto.\par
\sentnum{4}Di notte Tobia guarda dalla finestra. \sentnum{5}Nel cielo c'è la luna. \sentnum{6}La luna è grande e bianca. \sentnum{7}«Ciao, luna!» dice Tobia. \sentnum{8}La luna non risponde.\par
\sentnum{9}Tobia salta sul tavolo, ma la luna è sempre lontana. \sentnum{10}Allora Tobia beve il suo latte e torna a letto. \sentnum{11}Buonanotte, Tobia. Buonanotte, luna.\par
\end{storytext}
\subsection*{Translation, notes and verb analysis}
\begin{sentences}
\item\relax \textit{Il gatto si chiama Tobia.}\newline The cat is called Tobia.\newline {\small\color{muted}si chiama: from chiamarsi, to be called.}\newline {\footnotesize Verbs: \textit{chiama} = indicativo presente of chiamare, lui/lei; main clause. (\S\ref{les:present-are})}
\item\relax \textit{Tobia è un gatto nero e piccolo.}\newline Tobia is a small black cat.\newline {\small\color{muted}Two adjectives after the noun, joined by e.}\newline {\footnotesize Verbs: \textit{è} = indicativo presente of essere, lui/lei; main clause. (\S\ref{les:present-are})}
\item\relax \textit{Di giorno Tobia dorme sul letto.}\newline During the day Tobia sleeps on the bed.\newline {\small\color{muted}sul = su + il.}\newline {\footnotesize Verbs: \textit{dorme} = indicativo presente of dormire, lui/lei; main clause. (\S\ref{les:present-are})}
\item\relax \textit{Di notte Tobia guarda dalla finestra.}\newline At night Tobia looks out of the window.\newline {\small\color{muted}dalla = da + la: looking 'from' the window.}\newline {\footnotesize Verbs: \textit{guarda} = indicativo presente of guardare, lui/lei; main clause. (\S\ref{les:present-are})}
\item\relax \textit{Nel cielo c'è la luna.}\newline In the sky there is the moon.\newline {\small\color{muted}nel = in + il; c'è = there is.}\newline {\footnotesize Verbs: \textit{è} = indicativo presente of essere, lui/lei; main clause. (\S\ref{les:present-are})}
\item\relax \textit{La luna è grande e bianca.}\newline The moon is big and white.\newline {\small\color{muted}bianca agrees with luna (feminine).}\newline {\footnotesize Verbs: \textit{è} = indicativo presente of essere, lui/lei; main clause. (\S\ref{les:present-are})}
\item\relax \textit{«Ciao, luna!» dice Tobia.}\newline “Hello, moon!” says Tobia.\newline {\footnotesize Verbs: \textit{dice} = indicativo presente of dire, lui/lei; main clause. (\S\ref{les:present-are})}
\item\relax \textit{La luna non risponde.}\newline The moon doesn't answer.\newline {\footnotesize Verbs: \textit{risponde} = indicativo presente of rispondere, lui/lei; main clause. (\S\ref{les:present-are})}
\item\relax \textit{Tobia salta sul tavolo, ma la luna è sempre lontana.}\newline Tobia jumps onto the table, but the moon is still far away.\newline {\footnotesize Verbs: \textit{salta} = indicativo presente of saltare, lui/lei; main clause. (\S\ref{les:present-are}); \textit{è} = indicativo presente of essere, lui/lei; main clause. (\S\ref{les:present-are})}
\item\relax \textit{Allora Tobia beve il suo latte e torna a letto.}\newline So Tobia drinks his milk and goes back to bed.\newline {\small\color{muted}il suo latte: the possessive agrees with latte.}\newline {\footnotesize Verbs: \textit{beve} = indicativo presente of bere, lui/lei; main clause. (\S\ref{les:present-are}); \textit{torna} = indicativo presente of tornare, lui/lei; main clause. (\S\ref{les:present-are})}
\item\relax \textit{Buonanotte, Tobia. Buonanotte, luna.}\newline Good night, Tobia. Good night, moon.
\end{sentences}
\subsection*{Grammar in this story}
\begin{itemize}
\item\relax \textbf{Adjectives after the noun.} un gatto nero, la luna è grande e bianca: colour and size adjectives agree with the noun. (\S\ref{les:adjective-agreement})
\item\relax \textbf{Articulated prepositions.} sul letto (su + il), dalla finestra (da + la), nel cielo (in + il). (\S\ref{les:articulated-prepositions})
\item\relax \textbf{c'è.} c'è + singular noun = there is. (\S\ref{les:indefinite-articles})
\end{itemize}
\subsection*{Verb forms in this story}
\begin{itemize}
\item\relax \textbf{indicativo presente}: \textit{chiama}, \textit{è}, \textit{dorme}, \textit{guarda}, \textit{dice}, \textit{risponde} (\S\ref{les:present-are})
\item\relax \textbf{participio passato}: \textit{letto} (\S\ref{les:participio})
\end{itemize}
\relatedwords{\textit{il gatto} cat \textperiodcentered{} \textit{la luna} moon \textperiodcentered{} \textit{il letto} bed \textperiodcentered{} \textit{la finestra} window \textperiodcentered{} \textit{il cielo} sky \textperiodcentered{} \textit{il tavolo} table \textperiodcentered{} \textit{il latte} milk \textperiodcentered{} \textit{il giorno} day \textperiodcentered{} \textit{la notte} night \textperiodcentered{} \textit{nero} black \textperiodcentered{} \textit{piccolo} small; little \textperiodcentered{} \textit{grande} big; large; great \textperiodcentered{} \textit{bianco} white \textperiodcentered{} \textit{lontano} far; distant \textperiodcentered{} \textit{dormire} to sleep \textperiodcentered{} \textit{guardare} to look at; to watch \textperiodcentered{} \textit{saltare} to jump \textperiodcentered{} \textit{bere} to drink \textperiodcentered{} \textit{tornare} to return; to come back \textperiodcentered{} \textit{rispondere} to answer; to reply \textperiodcentered{} \textit{chiamarsi} to be called \textperiodcentered{} \textit{buonanotte} good night}
\subsection*{Questions}\label{story:il-gatto-e-la-luna:questions}
{\color{muted}\small Answers: page~\pageref{ans:story:il-gatto-e-la-luna}.}

\begin{exercises}
\item\label{ex:story-gatto-luna-1}\relax Di che colore è Tobia?\par\nopagebreak\begin{options}
\item\relax bianco
\item\relax rosso
\item\relax nero
\item\relax giallo
\end{options}
\item\label{ex:story-gatto-luna-2}\relax Dove dorme Tobia di giorno?\par\nopagebreak\begin{options}
\item\relax dalla finestra
\item\relax nel cielo
\item\relax sul letto
\item\relax sul tavolo
\end{options}
\item\label{ex:story-gatto-luna-3}\relax Che cosa fa la luna quando Tobia dice «Ciao»?\par\nopagebreak\begin{options}
\item\relax Dice «Ciao, Tobia!»
\item\relax Canta.
\item\relax Non risponde.
\item\relax Va via.
\end{options}
\item\label{ex:story-gatto-luna-4}\relax What does Tobia drink?\par\nopagebreak\begin{options}
\item\relax milk
\item\relax coffee
\item\relax water
\item\relax juice
\end{options}
\item\label{ex:story-gatto-luna-5}\relax Translate: \textit{La luna è grande e bianca.}\par\answerline
\end{exercises}

\section{La famiglia di Marco}\label{story:la-famiglia-di-marco}
\storymeta{Marco's family}{A0–A1}{Marco introduces his family, his pets and his grandparents.}
\subsection*{Text}
\begin{storytext}
\sentnum{1}Ciao! Mi chiamo Marco e ho sette anni. \sentnum{2}Abito in una piccola casa con la mia famiglia.\par
\sentnum{3}Mio padre si chiama Paolo e lavora in città. \sentnum{4}Mia madre si chiama Anna ed è maestra. \sentnum{5}Ho una sorella, Giulia. \sentnum{6}Giulia ha dieci anni e ama i gatti.\par
\sentnum{7}Abbiamo un cane, Bruno, e un gatto, Tobia. \sentnum{8}Il nonno e la nonna abitano al mare. \sentnum{9}In estate andiamo dai nonni. \sentnum{10}Il nonno racconta molte storie e la nonna fa una torta molto buona. \sentnum{11}Io amo la mia famiglia!\par
\end{storytext}
\subsection*{Translation, notes and verb analysis}
\begin{sentences}
\item\relax \textit{Ciao! Mi chiamo Marco e ho sette anni.}\newline Hi! My name is Marco and I'm seven years old.\newline {\small\color{muted}avere … anni = to be … years old.}\newline {\footnotesize Verbs: \textit{chiamo} = indicativo presente of chiamare, io; main clause. (\S\ref{les:present-are}); \textit{ho} = indicativo presente of avere, io; main clause. (\S\ref{les:present-are})}
\item\relax \textit{Abito in una piccola casa con la mia famiglia.}\newline I live in a small house with my family.\newline {\footnotesize Verbs: \textit{Abito} = indicativo presente of abitare, io; main clause. (\S\ref{les:present-are})}
\item\relax \textit{Mio padre si chiama Paolo e lavora in città.}\newline My father is called Paolo and works in town.\newline {\small\color{muted}mio padre: no article with a singular family noun.}\newline {\footnotesize Verbs: \textit{chiama} = indicativo presente of chiamare, lui/lei; main clause. (\S\ref{les:present-are}); \textit{lavora} = indicativo presente of lavorare, lui/lei; main clause. (\S\ref{les:present-are})}
\item\relax \textit{Mia madre si chiama Anna ed è maestra.}\newline My mother is called Anna and she is a teacher.\newline {\small\color{muted}No article before a profession after essere: è maestra.}\newline {\footnotesize Verbs: \textit{chiama} = indicativo presente of chiamare, lui/lei; main clause. (\S\ref{les:present-are}); \textit{è} = indicativo presente of essere, lui/lei; main clause. (\S\ref{les:present-are})}
\item\relax \textit{Ho una sorella, Giulia.}\newline I have a sister, Giulia.\newline {\footnotesize Verbs: \textit{Ho} = indicativo presente of avere, io; main clause. (\S\ref{les:present-are})}
\item\relax \textit{Giulia ha dieci anni e ama i gatti.}\newline Giulia is ten years old and loves cats.\newline {\footnotesize Verbs: \textit{ha} = indicativo presente of avere, lui/lei; main clause. (\S\ref{les:present-are}); \textit{ama} = indicativo presente of amare, lui/lei; main clause. (\S\ref{les:present-are})}
\item\relax \textit{Abbiamo un cane, Bruno, e un gatto, Tobia.}\newline We have a dog, Bruno, and a cat, Tobia.\newline {\footnotesize Verbs: \textit{Abbiamo} = indicativo presente of avere, noi; main clause. (\S\ref{les:present-are})}
\item\relax \textit{Il nonno e la nonna abitano al mare.}\newline Grandpa and Grandma live by the sea.\newline {\footnotesize Verbs: \textit{abitano} = indicativo presente of abitare, loro; main clause. (\S\ref{les:present-are})}
\item\relax \textit{In estate andiamo dai nonni.}\newline In summer we go to our grandparents' house.\newline {\small\color{muted}dai = da + i: at/to someone's place.}\newline {\footnotesize Verbs: \textit{andiamo} = indicativo presente of andare, noi; main clause. (\S\ref{les:present-are})}
\item\relax \textit{Il nonno racconta molte storie e la nonna fa una torta molto buona.}\newline Grandpa tells lots of stories and Grandma makes a very good cake.\newline {\footnotesize Verbs: \textit{racconta}: verb, indicativo presente, 3rd person singular (lui/lei) of raccontare | verb, imperativo, tu of raccontare; \textit{fa}: verb, indicativo presente, 3rd person singular (lui/lei) of fare}
\item\relax \textit{Io amo la mia famiglia!}\newline I love my family!\newline {\small\color{muted}io for emphasis.}\newline {\footnotesize Verbs: \textit{amo} = indicativo presente of amare, io; main clause. (\S\ref{les:present-are})}
\end{sentences}
\subsection*{Grammar in this story}
\begin{itemize}
\item\relax \textbf{Possessives with family nouns.} mio padre, mia madre (no article), but la mia famiglia, i nonni. (\S\ref{les:demonstratives-possessives})
\item\relax \textbf{avere … anni.} Age is expressed with avere: ho sette anni. (\S\ref{les:essere-avere})
\item\relax \textbf{da + person.} andiamo dai nonni = we go to the grandparents' place. (\S\ref{les:simple-prepositions})
\end{itemize}
\subsection*{Verb forms in this story}
\begin{itemize}
\item\relax \textbf{indicativo presente}: \textit{chiamo}, \textit{ho}, \textit{Abito}, \textit{chiama}, \textit{lavora}, \textit{è} (\S\ref{les:present-are})
\end{itemize}
\relatedwords{\textit{la famiglia} family \textperiodcentered{} \textit{il padre} father \textperiodcentered{} \textit{la madre} mother \textperiodcentered{} \textit{la sorella} sister \textperiodcentered{} \textit{il nonno} grandfather \textperiodcentered{} \textit{la nonna} grandmother \textperiodcentered{} \textit{il cane} dog \textperiodcentered{} \textit{il gatto} cat \textperiodcentered{} \textit{la casa} house; home \textperiodcentered{} \textit{la città} city; town \textperiodcentered{} \textit{il mare} sea \textperiodcentered{} \textit{l'estate} summer \textperiodcentered{} \textit{la storia} story; history \textperiodcentered{} \textit{la torta} cake \textperiodcentered{} \textit{la maestra} (primary-school) teacher (female) \textperiodcentered{} \textit{l'anno} year \textperiodcentered{} \textit{abitare} to live (reside) \textperiodcentered{} \textit{lavorare} to work \textperiodcentered{} \textit{amare} to love \textperiodcentered{} \textit{raccontare} to tell (a story) \textperiodcentered{} \textit{chiamarsi} to be called \textperiodcentered{} \textit{sette} seven \textperiodcentered{} \textit{dieci} ten}
\subsection*{Questions}\label{story:la-famiglia-di-marco:questions}
{\color{muted}\small Answers: page~\pageref{ans:story:la-famiglia-di-marco}.}

\begin{exercises}
\item\label{ex:story-famiglia-1}\relax Quanti anni ha Marco?\par\nopagebreak\begin{options}
\item\relax dieci
\item\relax sette
\item\relax sei
\item\relax otto
\end{options}
\item\label{ex:story-famiglia-2}\relax Che lavoro fa la madre di Marco?\par\nopagebreak\begin{options}
\item\relax È maestra.
\item\relax Racconta storie.
\item\relax Lavora in città.
\item\relax Fa le torte.
\end{options}
\item\label{ex:story-famiglia-3}\relax Chi è Bruno?\par\nopagebreak\begin{options}
\item\relax il nonno
\item\relax il padre
\item\relax il gatto
\item\relax il cane
\end{options}
\item\label{ex:story-famiglia-4}\relax Dove abitano i nonni?\par\nopagebreak\begin{options}
\item\relax in città
\item\relax in montagna
\item\relax al mare
\item\relax con Marco
\end{options}
\item\label{ex:story-famiglia-5}\relax Translate: My sister loves cats.\par\answerline
\end{exercises}

\section{La colazione}\label{story:la-colazione}
\storymeta{Breakfast}{A0–A1}{Sara has breakfast — and so does the cat.}
\subsection*{Text}
\begin{storytext}
\sentnum{1}È mattina. Il sole entra dalla finestra. \sentnum{2}Sara si sveglia e va in cucina. \sentnum{3}La mamma prepara la colazione.\par
\sentnum{4}Sul tavolo ci sono il pane, la marmellata, il latte e le mele. \sentnum{5}«Che cosa vuoi, Sara?» chiede la mamma. \sentnum{6}«Vorrei pane e marmellata, per favore.» \sentnum{7}Sara beve il latte e mangia il pane.\par
\sentnum{8}Anche il gatto vuole la colazione! \sentnum{9}La mamma dà al gatto un po' di latte. \sentnum{10}Poi Sara prende lo zaino e va a scuola.\par
\end{storytext}
\subsection*{Translation, notes and verb analysis}
\begin{sentences}
\item\relax \textit{È mattina. Il sole entra dalla finestra.}\newline It is morning. The sun comes in through the window.\newline {\small\color{muted}È mattina: impersonal 'it is'.}\newline {\footnotesize Verbs: \textit{È}: verb, indicativo presente, 3rd person singular (lui/lei) of essere; \textit{entra}: verb, indicativo presente, 3rd person singular (lui/lei) of entrare | verb, imperativo, tu of entrare}
\item\relax \textit{Sara si sveglia e va in cucina.}\newline Sara wakes up and goes into the kitchen.\newline {\footnotesize Verbs: \textit{sveglia}: verb, indicativo presente, 3rd person singular (lui/lei) (reflexive: used with mi, ti, si…) of svegliarsi; \textit{va}: verb, indicativo presente, 3rd person singular (lui/lei) of andare; \textit{cucina}: verb, indicativo presente, 3rd person singular (lui/lei) of cucinare | verb, imperativo, tu of cucinare}
\item\relax \textit{La mamma prepara la colazione.}\newline Mum is making breakfast.\newline {\footnotesize Verbs: \textit{prepara} = indicativo presente of preparare, lui/lei; main clause. (\S\ref{les:present-are})}
\item\relax \textit{Sul tavolo ci sono il pane, la marmellata, il latte e le mele.}\newline On the table there are bread, jam, milk and apples.\newline {\small\color{muted}ci sono + plural (several things).}\newline {\footnotesize Verbs: \textit{sono} = indicativo presente of essere, io; main clause. (\S\ref{les:present-are})}
\item\relax \textit{«Che cosa vuoi, Sara?» chiede la mamma.}\newline “What would you like, Sara?” Mum asks.\newline {\footnotesize Verbs: \textit{vuoi} = indicativo presente of volere, tu; subordinate clause (che). (\S\ref{les:present-are}); \textit{chiede} = indicativo presente of chiedere, lui/lei; main clause. (\S\ref{les:present-are})}
\item\relax \textit{«Vorrei pane e marmellata, per favore.»}\newline “I'd like bread and jam, please.”\newline {\small\color{muted}vorrei: polite form of volere.}\newline {\footnotesize Verbs: \textit{Vorrei} = condizionale presente of volere, io; main clause. (\S\ref{les:condizionale-presente})}
\item\relax \textit{Sara beve il latte e mangia il pane.}\newline Sara drinks the milk and eats the bread.\newline {\footnotesize Verbs: \textit{beve} = indicativo presente of bere, lui/lei; main clause. (\S\ref{les:present-are}); \textit{mangia} = indicativo presente of mangiare, lui/lei; main clause. (\S\ref{les:present-are})}
\item\relax \textit{Anche il gatto vuole la colazione!}\newline The cat wants breakfast too!\newline {\footnotesize Verbs: \textit{vuole} = indicativo presente of volere, lui/lei; main clause. (\S\ref{les:present-are})}
\item\relax \textit{La mamma dà al gatto un po' di latte.}\newline Mum gives the cat a little milk.\newline {\small\color{muted}dà (gives) has an accent; al = a + il; un po' di = a little.}\newline {\footnotesize Verbs: \textit{dà} = indicativo presente of dare, lui/lei; main clause. (\S\ref{les:present-are})}
\item\relax \textit{Poi Sara prende lo zaino e va a scuola.}\newline Then Sara takes her backpack and goes to school.\newline {\small\color{muted}lo zaino: lo before z.}\newline {\footnotesize Verbs: \textit{prende} = indicativo presente of prendere, lui/lei; main clause. (\S\ref{les:present-are}); \textit{va} = indicativo presente of andare, lui/lei; main clause. (\S\ref{les:present-are})}
\end{sentences}
\subsection*{Grammar in this story}
\begin{itemize}
\item\relax \textbf{c'è / ci sono.} Sul tavolo ci sono … (plural). (\S\ref{les:indefinite-articles})
\item\relax \textbf{Irregular verbs.} va (andare), vuoi / vuole (volere), dà (dare), beve (bere). (\S\ref{les:other-irregulars})
\item\relax \textbf{vorrei.} A polite request: vorrei = I would like. (\S\ref{les:modal-verbs})
\end{itemize}
\subsection*{Verb forms in this story}
\begin{itemize}
\item\relax \textbf{indicativo presente}: \textit{È}, \textit{sveglia}, \textit{va}, \textit{prepara}, \textit{sono}, \textit{vuoi} (\S\ref{les:present-are})
\item\relax \textbf{condizionale presente}: \textit{Vorrei} (\S\ref{les:condizionale-presente})
\end{itemize}
\relatedwords{\textit{la colazione} breakfast \textperiodcentered{} \textit{la mattina} morning \textperiodcentered{} \textit{il sole} sun \textperiodcentered{} \textit{la finestra} window \textperiodcentered{} \textit{la cucina} kitchen \textperiodcentered{} \textit{la mamma} mum; mom \textperiodcentered{} \textit{il tavolo} table \textperiodcentered{} \textit{il pane} bread \textperiodcentered{} \textit{la marmellata} jam \textperiodcentered{} \textit{il latte} milk \textperiodcentered{} \textit{la mela} apple \textperiodcentered{} \textit{lo zaino} backpack; rucksack \textperiodcentered{} \textit{la scuola} school \textperiodcentered{} \textit{svegliarsi} to wake up \textperiodcentered{} \textit{andare} to go \textperiodcentered{} \textit{preparare} to prepare; to make \textperiodcentered{} \textit{volere} to want \textperiodcentered{} \textit{chiedere} to ask \textperiodcentered{} \textit{bere} to drink \textperiodcentered{} \textit{mangiare} to eat \textperiodcentered{} \textit{dare} to give \textperiodcentered{} \textit{prendere} to take; to catch; to have (food) \textperiodcentered{} \textit{per favore} please}
\subsection*{Questions}\label{story:la-colazione:questions}
{\color{muted}\small Answers: page~\pageref{ans:story:la-colazione}.}

\begin{exercises}
\item\label{ex:story-colazione-1}\relax Chi prepara la colazione?\par\nopagebreak\begin{options}
\item\relax Sara
\item\relax il papà
\item\relax la mamma
\item\relax il gatto
\end{options}
\item\label{ex:story-colazione-2}\relax Che cosa mangia Sara?\par\nopagebreak\begin{options}
\item\relax niente
\item\relax una torta
\item\relax una mela
\item\relax pane e marmellata
\end{options}
\item\label{ex:story-colazione-3}\relax Che cosa dà la mamma al gatto?\par\nopagebreak\begin{options}
\item\relax un po' di latte
\item\relax una mela
\item\relax il pane
\item\relax la marmellata
\end{options}
\item\label{ex:story-colazione-4}\relax Where does Sara go at the end?\par\nopagebreak\begin{options}
\item\relax to school
\item\relax to bed
\item\relax to the kitchen
\item\relax to the seaside
\end{options}
\item\label{ex:story-colazione-5}\relax What is the infinitive of \textit{beve}?\par\answerline
\end{exercises}

\section{Il cane al mare}\label{story:il-cane-al-mare}
\storymeta{The dog at the seaside}{A0–A1}{Luca and his dog Pippo spend a hot day on the beach.}
\subsection*{Text}
\begin{storytext}
\sentnum{1}Oggi fa caldo. \sentnum{2}Luca e il suo cane, Pippo, vanno al mare. \sentnum{3}Il mare è blu e calmo.\par
\sentnum{4}Pippo corre sulla spiaggia e salta nelle onde. \sentnum{5}Luca fa un castello di sabbia. \sentnum{6}Sulla sabbia c'è un granchio piccolo. \sentnum{7}Pippo guarda il granchio e il granchio guarda Pippo. \sentnum{8}Poi il granchio scappa sotto un sasso!\par
\sentnum{9}Luca trova una conchiglia bianca e la porta a casa. \sentnum{10}La sera Pippo è stanco e bagnato, ma felice.\par
\end{storytext}
\subsection*{Translation, notes and verb analysis}
\begin{sentences}
\item\relax \textit{Oggi fa caldo.}\newline It is hot today.\newline {\small\color{muted}fa caldo: weather expression with fare.}\newline {\footnotesize Verbs: \textit{fa} = indicativo presente of fare, lui/lei; main clause. (\S\ref{les:present-are})}
\item\relax \textit{Luca e il suo cane, Pippo, vanno al mare.}\newline Luca and his dog, Pippo, go to the seaside.\newline {\small\color{muted}Two subjects → plural verb vanno.}\newline {\footnotesize Verbs: \textit{vanno} = indicativo presente of andare, loro; main clause. (\S\ref{les:present-are})}
\item\relax \textit{Il mare è blu e calmo.}\newline The sea is blue and calm.\newline {\footnotesize Verbs: \textit{è} = indicativo presente of essere, lui/lei; main clause. (\S\ref{les:present-are})}
\item\relax \textit{Pippo corre sulla spiaggia e salta nelle onde.}\newline Pippo runs on the beach and jumps into the waves.\newline {\footnotesize Verbs: \textit{corre} = indicativo presente of correre, lui/lei; main clause. (\S\ref{les:present-are}); \textit{salta} = indicativo presente of saltare, lui/lei; main clause. (\S\ref{les:present-are})}
\item\relax \textit{Luca fa un castello di sabbia.}\newline Luca makes a sandcastle.\newline {\small\color{muted}castello di sabbia: di for material.}\newline {\footnotesize Verbs: \textit{fa} = indicativo presente of fare, lui/lei; main clause. (\S\ref{les:present-are})}
\item\relax \textit{Sulla sabbia c'è un granchio piccolo.}\newline On the sand there is a little crab.\newline {\footnotesize Verbs: \textit{è} = indicativo presente of essere, lui/lei; main clause. (\S\ref{les:present-are})}
\item\relax \textit{Pippo guarda il granchio e il granchio guarda Pippo.}\newline Pippo looks at the crab and the crab looks at Pippo.\newline {\footnotesize Verbs: \textit{guarda}: verb, indicativo presente, 3rd person singular (lui/lei) of guardare | verb, imperativo, tu of guardare; \textit{guarda}: verb, indicativo presente, 3rd person singular (lui/lei) of guardare | verb, imperativo, tu of guardare}
\item\relax \textit{Poi il granchio scappa sotto un sasso!}\newline Then the crab runs off under a stone!\newline {\footnotesize Verbs: \textit{scappa} = indicativo presente of scappare, lui/lei; main clause. (\S\ref{les:present-are})}
\item\relax \textit{Luca trova una conchiglia bianca e la porta a casa.}\newline Luca finds a white shell and takes it home.\newline {\small\color{muted}la = it (the shell), a direct object pronoun.}\newline {\footnotesize Verbs: \textit{trova} = indicativo presente of trovare, lui/lei; main clause. (\S\ref{les:present-are})}
\item\relax \textit{La sera Pippo è stanco e bagnato, ma felice.}\newline In the evening Pippo is tired and wet, but happy.\newline {\footnotesize Verbs: \textit{è} = indicativo presente of essere, lui/lei; main clause. (\S\ref{les:present-are})}
\end{sentences}
\subsection*{Grammar in this story}
\begin{itemize}
\item\relax \textbf{fare + weather.} fa caldo = it is hot; fa freddo = it is cold. (\S\ref{les:andare-fare-stare-dire})
\item\relax \textbf{Articulated prepositions.} al mare (a + il), sulla spiaggia (su + la), nelle onde (in + le). (\S\ref{les:articulated-prepositions})
\item\relax \textbf{A first object pronoun.} la porta a casa: la replaces una conchiglia. (\S\ref{les:direct-object-pronouns})
\end{itemize}
\subsection*{Verb forms in this story}
\begin{itemize}
\item\relax \textbf{indicativo presente}: \textit{fa}, \textit{vanno}, \textit{è}, \textit{corre}, \textit{salta}, \textit{scappa} (\S\ref{les:present-are})
\end{itemize}
\relatedwords{\textit{il mare} sea \textperiodcentered{} \textit{il cane} dog \textperiodcentered{} \textit{la spiaggia} beach \textperiodcentered{} \textit{l'onda} wave \textperiodcentered{} \textit{il castello} castle \textperiodcentered{} \textit{la sabbia} sand \textperiodcentered{} \textit{il granchio} crab \textperiodcentered{} \textit{il sasso} stone; pebble \textperiodcentered{} \textit{la conchiglia} shell; seashell \textperiodcentered{} \textit{la sera} evening \textperiodcentered{} \textit{caldo} hot; warm \textperiodcentered{} \textit{blu} blue (dark) \textperiodcentered{} \textit{calmo} calm \textperiodcentered{} \textit{stanco} tired \textperiodcentered{} \textit{bagnato} wet \textperiodcentered{} \textit{felice} happy \textperiodcentered{} \textit{bianco} white \textperiodcentered{} \textit{correre} to run \textperiodcentered{} \textit{saltare} to jump \textperiodcentered{} \textit{trovare} to find \textperiodcentered{} \textit{portare} to bring; to carry; to wear \textperiodcentered{} \textit{scappare} to run away; to escape \textperiodcentered{} \textit{guardare} to look at; to watch}
\subsection*{Questions}\label{story:il-cane-al-mare:questions}
{\color{muted}\small Answers: page~\pageref{ans:story:il-cane-al-mare}.}

\begin{exercises}
\item\label{ex:story-cane-mare-1}\relax Che tempo fa?\par\nopagebreak\begin{options}
\item\relax Fa caldo.
\item\relax Nevica.
\item\relax Fa freddo.
\item\relax Piove.
\end{options}
\item\label{ex:story-cane-mare-2}\relax Che cosa fa Luca sulla spiaggia?\par\nopagebreak\begin{options}
\item\relax un castello di sabbia
\item\relax il bagno
\item\relax un nido
\item\relax una torta
\end{options}
\item\label{ex:story-cane-mare-3}\relax Dove scappa il granchio?\par\nopagebreak\begin{options}
\item\relax nelle onde
\item\relax sotto un sasso
\item\relax nel mare
\item\relax a casa
\end{options}
\item\label{ex:story-cane-mare-4}\relax In \textit{la porta a casa}, what does \textit{la} refer to?\par\nopagebreak\begin{options}
\item\relax la sera
\item\relax la spiaggia
\item\relax la conchiglia
\item\relax la sabbia
\end{options}
\item\label{ex:story-cane-mare-5}\relax How is Pippo in the evening?\par\nopagebreak\begin{options}
\item\relax tired, wet and happy
\item\relax hungry
\item\relax sad and cold
\item\relax angry
\end{options}
\end{exercises}

\section{Nel bosco}\label{story:nel-bosco}
\storymeta{In the wood}{A1}{Giulia, Luca and their grandfather walk quietly through the wood and meet its animals.}
\subsection*{Text}
\begin{storytext}
\sentnum{1}Giulia e suo fratello Luca vanno nel bosco con il nonno. \sentnum{2}Nel bosco c'è un sentiero lungo. \sentnum{3}Gli alberi sono alti e le foglie sono verdi.\par
\sentnum{4}Il nonno dice: «Che cosa sentite?» \sentnum{5}I bambini sentono un uccello che canta. \sentnum{6}Su un ramo c'è un nido con tre piccoli uccelli. \sentnum{7}Poi vedono uno scoiattolo. \sentnum{8}Lo scoiattolo sale veloce sull'albero.\par
\sentnum{9}Vicino al fiume c'è una volpe rossa. \sentnum{10}La volpe guarda i bambini e poi torna nella sua tana. \sentnum{11}«Il bosco è la casa di molti animali» dice il nonno. \sentnum{12}I bambini camminano piano per non fare rumore.\par
\end{storytext}
\subsection*{Translation, notes and verb analysis}
\begin{sentences}
\item\relax \textit{Giulia e suo fratello Luca vanno nel bosco con il nonno.}\newline Giulia and her brother Luca go into the wood with their grandfather.\newline {\footnotesize Verbs: \textit{vanno} = indicativo presente of andare, loro; main clause. (\S\ref{les:present-are})}
\item\relax \textit{Nel bosco c'è un sentiero lungo.}\newline In the wood there is a long path.\newline {\footnotesize Verbs: \textit{è} = indicativo presente of essere, lui/lei; main clause. (\S\ref{les:present-are})}
\item\relax \textit{Gli alberi sono alti e le foglie sono verdi.}\newline The trees are tall and the leaves are green.\newline {\footnotesize Verbs: \textit{sono} = indicativo presente of essere, loro; main clause. (\S\ref{les:present-are})}
\item\relax \textit{Il nonno dice: «Che cosa sentite?»}\newline Grandfather says: “What can you hear?”\newline {\small\color{muted}sentite: voi form of sentire.}\newline {\footnotesize Verbs: \textit{dice} = indicativo presente of dire, lui/lei; main clause. (\S\ref{les:present-are}); \textit{sentite} = indicativo presente of sentire, voi; subordinate clause (che). (\S\ref{les:present-are})}
\item\relax \textit{I bambini sentono un uccello che canta.}\newline The children hear a bird singing.\newline {\small\color{muted}che canta: relative clause, 'that is singing'.}\newline {\footnotesize Verbs: \textit{sentono} = indicativo presente of sentire, loro; main clause. (\S\ref{les:present-are}); \textit{canta} = indicativo presente of cantare, lui/lei; relative clause. (\S\ref{les:present-are})}
\item\relax \textit{Su un ramo c'è un nido con tre piccoli uccelli.}\newline On a branch there is a nest with three little birds.\newline {\footnotesize Verbs: \textit{è} = indicativo presente of essere, lui/lei; main clause. (\S\ref{les:present-are})}
\item\relax \textit{Poi vedono uno scoiattolo.}\newline Then they see a squirrel.\newline {\small\color{muted}uno before s + consonant.}\newline {\footnotesize Verbs: \textit{vedono} = indicativo presente of vedere, loro; main clause. (\S\ref{les:present-are})}
\item\relax \textit{Lo scoiattolo sale veloce sull'albero.}\newline The squirrel climbs quickly up the tree.\newline {\small\color{muted}sale: salire (sal-, salgo, sali, sale…).}\newline {\footnotesize Verbs: \textit{sale}: verb, indicativo presente, 3rd person singular (lui/lei) of salire}
\item\relax \textit{Vicino al fiume c'è una volpe rossa.}\newline Near the river there is a red fox.\newline {\footnotesize Verbs: \textit{è} = indicativo presente of essere, lui/lei; main clause. (\S\ref{les:present-are})}
\item\relax \textit{La volpe guarda i bambini e poi torna nella sua tana.}\newline The fox looks at the children and then goes back into its den.\newline {\footnotesize Verbs: \textit{guarda} = indicativo presente of guardare, lui/lei; main clause. (\S\ref{les:present-are}); \textit{torna} = indicativo presente of tornare, lui/lei; main clause. (\S\ref{les:present-are})}
\item\relax \textit{«Il bosco è la casa di molti animali» dice il nonno.}\newline “The wood is home to many animals,” says Grandfather.\newline {\footnotesize Verbs: \textit{è} = indicativo presente of essere, lui/lei; main clause. (\S\ref{les:present-are}); \textit{dice} = indicativo presente of dire, lui/lei; main clause. (\S\ref{les:present-are})}
\item\relax \textit{I bambini camminano piano per non fare rumore.}\newline The children walk quietly so as not to make any noise.\newline {\small\color{muted}per non + infinitive = so as not to.}\newline {\footnotesize Verbs: \textit{camminano} = indicativo presente of camminare, loro; main clause. (\S\ref{les:present-are})}
\end{sentences}
\subsection*{Grammar in this story}
\begin{itemize}
\item\relax \textbf{Articulated prepositions.} nel bosco, sull'albero, al fiume, nella sua tana. (\S\ref{les:articulated-prepositions})
\item\relax \textbf{Relative che.} un uccello che canta = a bird that is singing. (\S\ref{les:relative-clauses-adverbs})
\item\relax \textbf{uno scoiattolo / lo scoiattolo.} s + consonant takes uno and lo. (\S\ref{les:definite-articles})
\end{itemize}
\subsection*{Verb forms in this story}
\begin{itemize}
\item\relax \textbf{indicativo presente}: \textit{vanno}, \textit{è}, \textit{sono}, \textit{dice}, \textit{sentite}, \textit{sentono} (\S\ref{les:present-are})
\item\relax \textbf{infinito}: \textit{fare} (\S\ref{les:infinito-gerundio})
\end{itemize}
\relatedwords{\textit{il bosco} wood; forest \textperiodcentered{} \textit{il sentiero} path; trail \textperiodcentered{} \textit{l'albero} tree \textperiodcentered{} \textit{la foglia} leaf \textperiodcentered{} \textit{l'uccello} bird \textperiodcentered{} \textit{il ramo} branch \textperiodcentered{} \textit{il nido} nest \textperiodcentered{} \textit{lo scoiattolo} squirrel \textperiodcentered{} \textit{il fiume} river \textperiodcentered{} \textit{la volpe} fox \textperiodcentered{} \textit{la tana} den; burrow \textperiodcentered{} \textit{l'animale} animal \textperiodcentered{} \textit{il rumore} noise \textperiodcentered{} \textit{il fratello} brother \textperiodcentered{} \textit{lungo} long \textperiodcentered{} \textit{alto} tall; high \textperiodcentered{} \textit{verde} green \textperiodcentered{} \textit{rosso} red \textperiodcentered{} \textit{sentire} to hear; to feel \textperiodcentered{} \textit{vedere} to see \textperiodcentered{} \textit{salire} to go up; to climb; to get on \textperiodcentered{} \textit{camminare} to walk \textperiodcentered{} \textit{piano} slowly; softly; quietly \textperiodcentered{} \textit{vicino} near; nearby}
\subsection*{Questions}\label{story:nel-bosco:questions}
{\color{muted}\small Answers: page~\pageref{ans:story:nel-bosco}.}

\begin{exercises}
\item\label{ex:story-bosco-1}\relax Con chi vanno nel bosco Giulia e Luca?\par\nopagebreak\begin{options}
\item\relax con la mamma
\item\relax da soli
\item\relax con il cane
\item\relax con il nonno
\end{options}
\item\label{ex:story-bosco-2}\relax Quanti uccelli ci sono nel nido?\par\nopagebreak\begin{options}
\item\relax uno
\item\relax dieci
\item\relax due
\item\relax tre
\end{options}
\item\label{ex:story-bosco-3}\relax Dove torna la volpe?\par\nopagebreak\begin{options}
\item\relax al fiume
\item\relax a casa del nonno
\item\relax sull'albero
\item\relax nella sua tana
\end{options}
\item\label{ex:story-bosco-4}\relax Why do the children walk quietly?\par\nopagebreak\begin{options}
\item\relax so as not to make noise
\item\relax because the path is long
\item\relax because it is night
\item\relax because they are tired
\end{options}
\item\label{ex:story-bosco-5}\relax What is the infinitive of \textit{sale}?\par\answerline
\item\label{ex:story-bosco-6}\relax Complete: Il bosco è la casa \blank{} molti animali.
\end{exercises}

\section{La tartaruga lenta}\label{story:la-tartaruga-lenta}
\storymeta{The slow tortoise}{A1}{A fast hare laughs at a slow tortoise — until they have a race.}
\subsection*{Text}
\begin{storytext}
\sentnum{1}In un prato vivono una lepre e una tartaruga. \sentnum{2}La lepre è molto veloce. La tartaruga è molto lenta. \sentnum{3}Ogni giorno la lepre ride della tartaruga. \sentnum{4}«Sei così lenta!» dice la lepre.\par
\sentnum{5}Un giorno la tartaruga risponde: «Facciamo una gara!» \sentnum{6}Tutti gli animali vengono a guardare. \sentnum{7}«La gara finisce al grande albero» dice il gufo.\par
\sentnum{8}La lepre corre veloce, poi si siede sotto un albero. \sentnum{9}«La tartaruga è lontana. Ho tempo per dormire» pensa la lepre. \sentnum{10}E la lepre si addormenta. \sentnum{11}La tartaruga cammina lentamente, ma non si ferma mai.\par
\sentnum{12}Quando la lepre si sveglia, la tartaruga è già vicino al grande albero! \sentnum{13}La lepre corre, ma è tardi. \sentnum{14}La tartaruga vince la gara e tutti gli animali cantano e ballano. \sentnum{15}Da quel giorno la lepre non ride più della tartaruga.\par
\end{storytext}
\subsection*{Translation, notes and verb analysis}
\begin{sentences}
\item\relax \textit{In un prato vivono una lepre e una tartaruga.}\newline A hare and a tortoise live in a meadow.\newline {\small\color{muted}The subject follows the verb: vivono una lepre e una tartaruga.}\newline {\footnotesize Verbs: \textit{vivono} = indicativo presente of vivere, loro; main clause. (\S\ref{les:present-are})}
\item\relax \textit{La lepre è molto veloce. La tartaruga è molto lenta.}\newline The hare is very fast. The tortoise is very slow.\newline {\footnotesize Verbs: \textit{è}: verb, indicativo presente, 3rd person singular (lui/lei) of essere; \textit{è}: verb, indicativo presente, 3rd person singular (lui/lei) of essere}
\item\relax \textit{Ogni giorno la lepre ride della tartaruga.}\newline Every day the hare laughs at the tortoise.\newline {\small\color{muted}ridere di = to laugh at; della = di + la.}\newline {\footnotesize Verbs: \textit{ride} = indicativo presente of ridere, lui/lei; main clause. (\S\ref{les:present-are})}
\item\relax \textit{«Sei così lenta!» dice la lepre.}\newline “You're so slow!” says the hare.\newline {\footnotesize Verbs: \textit{Sei} = indicativo presente of essere, tu; main clause. (\S\ref{les:present-are}); \textit{dice} = indicativo presente of dire, lui/lei; main clause. (\S\ref{les:present-are})}
\item\relax \textit{Un giorno la tartaruga risponde: «Facciamo una gara!»}\newline One day the tortoise replies: “Let's have a race!”\newline {\small\color{muted}facciamo: the noi form is also used for suggestions (let's…).}\newline {\footnotesize Verbs: \textit{risponde} = indicativo presente of rispondere, lui/lei; main clause. (\S\ref{les:present-are}); \textit{Facciamo} = imperativo of fare, noi; main clause. (\S\ref{les:imperativo})}
\item\relax \textit{Tutti gli animali vengono a guardare.}\newline All the animals come to watch.\newline {\small\color{muted}venire a + infinitive.}\newline {\footnotesize Verbs: \textit{vengono} = indicativo presente of venire, loro; main clause. (\S\ref{les:present-are})}
\item\relax \textit{«La gara finisce al grande albero» dice il gufo.}\newline “The race finishes at the big tree,” says the owl.\newline {\footnotesize Verbs: \textit{finisce} = indicativo presente of finire, lui/lei; main clause. (\S\ref{les:present-are}); \textit{dice} = indicativo presente of dire, lui/lei; main clause. (\S\ref{les:present-are})}
\item\relax \textit{La lepre corre veloce, poi si siede sotto un albero.}\newline The hare runs fast, then sits down under a tree.\newline {\small\color{muted}si siede: sedersi.}\newline {\footnotesize Verbs: \textit{corre} = indicativo presente of correre, lui/lei; main clause. (\S\ref{les:present-are}); \textit{siede} = indicativo presente of sedersi, lui/lei; main clause. (\S\ref{les:present-are})}
\item\relax \textit{«La tartaruga è lontana. Ho tempo per dormire» pensa la lepre.}\newline “The tortoise is far away. I have time to sleep,” thinks the hare.\newline {\footnotesize Verbs: \textit{è}: verb, indicativo presente, 3rd person singular (lui/lei) of essere; \textit{Ho}: verb, indicativo presente, 1st person singular (io) of avere; \textit{dormire}: verb, infinito presente of dormire; \textit{pensa}: verb, indicativo presente, 3rd person singular (lui/lei) of pensare | verb, imperativo, tu of pensare}
\item\relax \textit{E la lepre si addormenta.}\newline And the hare falls asleep.\newline {\footnotesize Verbs: \textit{addormenta} = indicativo presente of addormentarsi, lui/lei; main clause. (\S\ref{les:present-are})}
\item\relax \textit{La tartaruga cammina lentamente, ma non si ferma mai.}\newline The tortoise walks slowly, but never stops.\newline {\small\color{muted}non … mai = never; lentamente = lenta + -mente.}\newline {\footnotesize Verbs: \textit{cammina} = indicativo presente of camminare, lui/lei; main clause. (\S\ref{les:present-are}); \textit{ferma} = indicativo presente of fermarsi, lui/lei; main clause. (\S\ref{les:present-are})}
\item\relax \textit{Quando la lepre si sveglia, la tartaruga è già vicino al grande albero!}\newline When the hare wakes up, the tortoise is already near the big tree!\newline {\footnotesize Verbs: \textit{sveglia} = indicativo presente of svegliarsi, lui/lei; subordinate clause (quando). (\S\ref{les:present-are}); \textit{è} = indicativo presente of essere, lui/lei; main clause. (\S\ref{les:present-are})}
\item\relax \textit{La lepre corre, ma è tardi.}\newline The hare runs, but it is too late.\newline {\footnotesize Verbs: \textit{corre} = indicativo presente of correre, lui/lei; main clause. (\S\ref{les:present-are}); \textit{è} = indicativo presente of essere, lui/lei; main clause. (\S\ref{les:present-are})}
\item\relax \textit{La tartaruga vince la gara e tutti gli animali cantano e ballano.}\newline The tortoise wins the race and all the animals sing and dance.\newline {\footnotesize Verbs: \textit{vince} = indicativo presente of vincere, lui/lei; main clause. (\S\ref{les:present-are}); \textit{cantano} = indicativo presente of cantare, loro; main clause. (\S\ref{les:present-are}); \textit{ballano} = indicativo presente of ballare, loro; main clause. (\S\ref{les:present-are})}
\item\relax \textit{Da quel giorno la lepre non ride più della tartaruga.}\newline From that day on, the hare no longer laughs at the tortoise.\newline {\small\color{muted}non … più = no longer.}\newline {\footnotesize Verbs: \textit{ride} = indicativo presente of ridere, lui/lei; main clause. (\S\ref{les:present-are})}
\end{sentences}
\subsection*{Grammar in this story}
\begin{itemize}
\item\relax \textbf{Reflexive verbs.} si siede, si addormenta, si ferma, si sveglia. (\S\ref{les:reflexive-verbs})
\item\relax \textbf{Double negation.} non si ferma mai; non ride più. (\S\ref{les:word-order-negation-questions})
\item\relax \textbf{Adverbs in -mente.} lenta → lentamente. (\S\ref{les:relative-clauses-adverbs})
\end{itemize}
\subsection*{Verb forms in this story}
\begin{itemize}
\item\relax \textbf{indicativo presente}: \textit{vivono}, \textit{è}, \textit{ride}, \textit{Sei}, \textit{dice}, \textit{risponde} (\S\ref{les:present-are})
\item\relax \textbf{imperativo}: \textit{Facciamo} (\S\ref{les:imperativo})
\item\relax \textbf{infinito}: \textit{guardare}, \textit{dormire} (\S\ref{les:infinito-gerundio})
\end{itemize}
\relatedwords{\textit{la lepre} hare \textperiodcentered{} \textit{la tartaruga} tortoise; turtle \textperiodcentered{} \textit{il prato} meadow; lawn \textperiodcentered{} \textit{la gara} race; competition \textperiodcentered{} \textit{il gufo} owl \textperiodcentered{} \textit{l'albero} tree \textperiodcentered{} \textit{l'animale} animal \textperiodcentered{} \textit{il tempo} weather; time \textperiodcentered{} \textit{veloce} fast; quick \textperiodcentered{} \textit{lento} slow \textperiodcentered{} \textit{lontano} far; distant \textperiodcentered{} \textit{lentamente} slowly \textperiodcentered{} \textit{vivere} to live \textperiodcentered{} \textit{ridere} to laugh \textperiodcentered{} \textit{rispondere} to answer; to reply \textperiodcentered{} \textit{venire} to come \textperiodcentered{} \textit{finire} to finish; to end \textperiodcentered{} \textit{correre} to run \textperiodcentered{} \textit{sedersi} to sit down \textperiodcentered{} \textit{pensare} to think \textperiodcentered{} \textit{addormentarsi} to fall asleep \textperiodcentered{} \textit{svegliarsi} to wake up \textperiodcentered{} \textit{vincere} to win \textperiodcentered{} \textit{ballare} to dance \textperiodcentered{} \textit{tardi} late \textperiodcentered{} \textit{già} already}
\subsection*{Questions}\label{story:la-tartaruga-lenta:questions}
{\color{muted}\small Answers: page~\pageref{ans:story:la-tartaruga-lenta}.}

\begin{exercises}
\item\label{ex:story-tartaruga-1}\relax Chi ride della tartaruga?\par\nopagebreak\begin{options}
\item\relax nessuno
\item\relax la lepre
\item\relax il gufo
\item\relax tutti gli animali
\end{options}
\item\label{ex:story-tartaruga-2}\relax Dove finisce la gara?\par\nopagebreak\begin{options}
\item\relax al grande albero
\item\relax nel prato
\item\relax nella tana
\item\relax al fiume
\end{options}
\item\label{ex:story-tartaruga-3}\relax Perché la lepre perde la gara?\par\nopagebreak\begin{options}
\item\relax Perché è stanca della tartaruga.
\item\relax Perché si addormenta.
\item\relax Perché ha paura.
\item\relax Perché è lenta.
\end{options}
\item\label{ex:story-tartaruga-4}\relax What is the infinitive of \textit{si siede}?\par\answerline
\item\label{ex:story-tartaruga-5}\relax Translate: \textit{La tartaruga non si ferma mai.}\par\answerline
\end{exercises}

\section{Una giornata di pioggia}\label{story:una-giornata-di-pioggia}
\storymeta{A rainy day}{A1}{It's raining, but Sara and Luca find a way to enjoy it.}
\subsection*{Text}
\begin{storytext}
\sentnum{1}Oggi piove e il cielo è grigio. \sentnum{2}Sara si sveglia e guarda dalla finestra. \sentnum{3}«Non possiamo giocare in giardino» dice suo fratello Luca. \sentnum{4}Sara è un po' triste.\par
\sentnum{5}Poi ha un'idea. \sentnum{6}«Mettiamo gli stivali e prendiamo l'ombrello!» \sentnum{7}I bambini escono di casa. \sentnum{8}Sulla strada ci sono molte pozzanghere.\par
\sentnum{9}Sara salta in una pozzanghera e ride. \sentnum{10}Luca salta in una pozzanghera più grande. \sentnum{11}Anche un'anatra cammina sotto la pioggia. \sentnum{12}Quando tornano a casa, i bambini sono bagnati ma contenti. \sentnum{13}La mamma prepara due tazze di latte caldo. \sentnum{14}«La pioggia non è così brutta» dice Sara.\par
\end{storytext}
\subsection*{Translation, notes and verb analysis}
\begin{sentences}
\item\relax \textit{Oggi piove e il cielo è grigio.}\newline Today it is raining and the sky is grey.\newline {\small\color{muted}piove: weather verb with no subject.}\newline {\footnotesize Verbs: \textit{piove} = indicativo presente of piovere, lui/lei; main clause. (\S\ref{les:present-are}); \textit{è} = indicativo presente of essere, lui/lei; main clause. (\S\ref{les:present-are})}
\item\relax \textit{Sara si sveglia e guarda dalla finestra.}\newline Sara wakes up and looks out of the window.\newline {\footnotesize Verbs: \textit{sveglia} = indicativo presente of svegliarsi, lui/lei; main clause. (\S\ref{les:present-are}); \textit{guarda} = indicativo presente of guardare, lui/lei; main clause. (\S\ref{les:present-are})}
\item\relax \textit{«Non possiamo giocare in giardino» dice suo fratello Luca.}\newline “We can't play in the garden,” says her brother Luca.\newline {\small\color{muted}potere + infinitive.}\newline {\footnotesize Verbs: \textit{possiamo giocare} = indicativo presente of potere + giocare, modal + infinitive, noi; main clause. (\S\ref{les:present-are}); \textit{dice} = indicativo presente of dire, lui/lei; main clause. (\S\ref{les:present-are})}
\item\relax \textit{Sara è un po' triste.}\newline Sara is a little sad.\newline {\footnotesize Verbs: \textit{è} = indicativo presente of essere, lui/lei; main clause. (\S\ref{les:present-are})}
\item\relax \textit{Poi ha un'idea.}\newline Then she has an idea.\newline {\small\color{muted}un'idea: una elided before a vowel.}\newline {\footnotesize Verbs: \textit{ha} = indicativo presente of avere, lui/lei; main clause. (\S\ref{les:present-are})}
\item\relax \textit{«Mettiamo gli stivali e prendiamo l'ombrello!»}\newline “Let's put on our boots and take the umbrella!”\newline {\small\color{muted}noi form for a suggestion: let's.}\newline {\footnotesize Verbs: \textit{Mettiamo} = imperativo of mettere, noi; main clause. (\S\ref{les:imperativo}); \textit{prendiamo} = indicativo presente of prendere, noi; main clause. (\S\ref{les:present-are})}
\item\relax \textit{I bambini escono di casa.}\newline The children go out of the house.\newline {\small\color{muted}uscire: escono.}\newline {\footnotesize Verbs: \textit{escono} = indicativo presente of uscire, loro; main clause. (\S\ref{les:present-are})}
\item\relax \textit{Sulla strada ci sono molte pozzanghere.}\newline There are lots of puddles in the street.\newline {\footnotesize Verbs: \textit{sono} = indicativo presente of essere, io; main clause. (\S\ref{les:present-are})}
\item\relax \textit{Sara salta in una pozzanghera e ride.}\newline Sara jumps into a puddle and laughs.\newline {\footnotesize Verbs: \textit{salta} = indicativo presente of saltare, lui/lei; main clause. (\S\ref{les:present-are}); \textit{ride} = indicativo presente of ridere, lui/lei; main clause. (\S\ref{les:present-are})}
\item\relax \textit{Luca salta in una pozzanghera più grande.}\newline Luca jumps into a bigger puddle.\newline {\small\color{muted}più grande = bigger.}\newline {\footnotesize Verbs: \textit{salta} = indicativo presente of saltare, lui/lei; main clause. (\S\ref{les:present-are})}
\item\relax \textit{Anche un'anatra cammina sotto la pioggia.}\newline A duck is walking in the rain too.\newline {\footnotesize Verbs: \textit{cammina} = indicativo presente of camminare, lui/lei; main clause. (\S\ref{les:present-are})}
\item\relax \textit{Quando tornano a casa, i bambini sono bagnati ma contenti.}\newline When they get home, the children are wet but happy.\newline {\footnotesize Verbs: \textit{tornano} = indicativo presente of tornare, loro; subordinate clause (quando). (\S\ref{les:present-are}); \textit{sono} = indicativo presente of essere, loro; main clause. (\S\ref{les:present-are})}
\item\relax \textit{La mamma prepara due tazze di latte caldo.}\newline Mum makes two cups of hot milk.\newline {\footnotesize Verbs: \textit{prepara} = indicativo presente of preparare, lui/lei; main clause. (\S\ref{les:present-are})}
\item\relax \textit{«La pioggia non è così brutta» dice Sara.}\newline “The rain isn't so bad,” says Sara.\newline {\footnotesize Verbs: \textit{è} = indicativo presente of essere, lui/lei; main clause. (\S\ref{les:present-are}); \textit{dice} = indicativo presente of dire, lui/lei; main clause. (\S\ref{les:present-are})}
\end{sentences}
\subsection*{Grammar in this story}
\begin{itemize}
\item\relax \textbf{Modal + infinitive.} Non possiamo giocare = We can't play. (\S\ref{les:modal-verbs})
\item\relax \textbf{Irregular uscire.} esco, esci, esce, usciamo, uscite, escono. (\S\ref{les:other-irregulars})
\item\relax \textbf{Comparatives.} una pozzanghera più grande = a bigger puddle. (\S\ref{les:bello-buono-comparison})
\end{itemize}
\subsection*{Verb forms in this story}
\begin{itemize}
\item\relax \textbf{indicativo presente}: \textit{piove}, \textit{è}, \textit{sveglia}, \textit{guarda}, \textit{possiamo giocare}, \textit{dice} (\S\ref{les:present-are})
\item\relax \textbf{imperativo}: \textit{Mettiamo} (\S\ref{les:imperativo})
\end{itemize}
\relatedwords{\textit{la pioggia} rain \textperiodcentered{} \textit{il cielo} sky \textperiodcentered{} \textit{la finestra} window \textperiodcentered{} \textit{il giardino} garden \textperiodcentered{} \textit{l'idea} idea \textperiodcentered{} \textit{lo stivale} boot \textperiodcentered{} \textit{l'ombrello} umbrella \textperiodcentered{} \textit{la strada} street; road \textperiodcentered{} \textit{la pozzanghera} puddle \textperiodcentered{} \textit{l'anatra} duck \textperiodcentered{} \textit{la tazza} cup; mug \textperiodcentered{} \textit{grigio} grey \textperiodcentered{} \textit{triste} sad \textperiodcentered{} \textit{bagnato} wet \textperiodcentered{} \textit{contento} glad; pleased \textperiodcentered{} \textit{brutto} ugly; bad (weather) \textperiodcentered{} \textit{piovere} to rain \textperiodcentered{} \textit{svegliarsi} to wake up \textperiodcentered{} \textit{potere} to be able to; can; may \textperiodcentered{} \textit{mettere} to put \textperiodcentered{} \textit{prendere} to take; to catch; to have (food) \textperiodcentered{} \textit{uscire} to go out; to leave \textperiodcentered{} \textit{saltare} to jump \textperiodcentered{} \textit{ridere} to laugh \textperiodcentered{} \textit{la giornata} day (its duration, the daytime)}
\subsection*{Questions}\label{story:una-giornata-di-pioggia:questions}
{\color{muted}\small Answers: page~\pageref{ans:story:una-giornata-di-pioggia}.}

\begin{exercises}
\item\label{ex:story-pioggia-1}\relax Che tempo fa?\par\nopagebreak\begin{options}
\item\relax Piove.
\item\relax Nevica.
\item\relax Fa caldo.
\item\relax C'è il sole.
\end{options}
\item\label{ex:story-pioggia-2}\relax Che cosa mettono i bambini?\par\nopagebreak\begin{options}
\item\relax il cappello
\item\relax le scarpe
\item\relax lo zaino
\item\relax gli stivali
\end{options}
\item\label{ex:story-pioggia-3}\relax Come sono i bambini quando tornano a casa?\par\nopagebreak\begin{options}
\item\relax tristi e stanchi
\item\relax bagnati ma contenti
\item\relax asciutti
\item\relax arrabbiati
\end{options}
\item\label{ex:story-pioggia-4}\relax What is the infinitive of \textit{escono}?\par\answerline
\item\label{ex:story-pioggia-5}\relax Translate: We can't play in the garden.\par\answerline
\end{exercises}

\section{L'orso che non poteva dormire}\label{story:lorso-che-non-poteva-dormire}
\storymeta{The bear who couldn't sleep}{A1–A2}{In the middle of winter a little bear cannot sleep. An owl helps him.}
\subsection*{Text}
\begin{storytext}
\sentnum{1}C'era una volta un piccolo orso che si chiamava Bruno. \sentnum{2}Bruno abitava in una grotta nel bosco con la sua mamma. \sentnum{3}Era inverno e faceva molto freddo. \sentnum{4}Tutti gli orsi dormivano, ma Bruno non poteva dormire. \sentnum{5}«Mamma, non ho sonno» diceva Bruno ogni sera.\par
\sentnum{6}Una notte Bruno è uscito dalla grotta. \sentnum{7}Fuori c'era la neve e nel bosco c'era silenzio. \sentnum{8}Sopra un ramo c'era un gufo. \sentnum{9}«Perché non dormi?» ha chiesto il gufo. \sentnum{10}«Non lo so» ha risposto Bruno. «Ho la pancia vuota.»\par
\sentnum{11}Il gufo ha riso. \sentnum{12}«Gli orsi mangiano molto in autunno, prima dell'inverno.» \sentnum{13}Allora Bruno ha pensato al miele. \sentnum{14}«Le api hanno il miele!» ha detto Bruno. \sentnum{15}Ma il gufo ha detto: «In inverno anche le api dormono.»\par
\sentnum{16}Bruno era triste. \sentnum{17}Il gufo ha guardato la luna e ha cominciato a cantare piano piano. \sentnum{18}Era una canzone dolce e lenta. \sentnum{19}Bruno ha chiuso gli occhi. \sentnum{20}Poi è tornato nella grotta, si è messo vicino alla mamma e si è addormentato.\par
\sentnum{21}Ha dormito fino alla primavera. \sentnum{22}Quando si è svegliato, il sole era caldo e nel prato c'erano molti fiori. \sentnum{23}E la mamma lo aspettava con il miele.\par
\end{storytext}
\subsection*{Translation, notes and verb analysis}
\begin{sentences}
\item\relax \textit{C'era una volta un piccolo orso che si chiamava Bruno.}\newline Once upon a time there was a little bear called Bruno.\newline {\small\color{muted}Fairy-tale opening with the imperfetto.}\newline {\footnotesize Verbs: \textit{era} = indicativo imperfetto of essere, lui/lei; main clause. (\S\ref{les:imperfetto}); \textit{chiamava} = indicativo imperfetto of chiamare, lui/lei; relative clause. (\S\ref{les:imperfetto})}
\item\relax \textit{Bruno abitava in una grotta nel bosco con la sua mamma.}\newline Bruno lived in a cave in the wood with his mother.\newline {\footnotesize Verbs: \textit{abitava} = indicativo imperfetto of abitare, lui/lei; main clause. (\S\ref{les:imperfetto})}
\item\relax \textit{Era inverno e faceva molto freddo.}\newline It was winter and it was very cold.\newline {\small\color{muted}Background description → imperfetto.}\newline {\footnotesize Verbs: \textit{Era} = indicativo imperfetto of essere, lui/lei; main clause. (\S\ref{les:imperfetto}); \textit{faceva} = indicativo imperfetto of fare, lui/lei; main clause. (\S\ref{les:imperfetto})}
\item\relax \textit{Tutti gli orsi dormivano, ma Bruno non poteva dormire.}\newline All the bears were sleeping, but Bruno couldn't sleep.\newline {\footnotesize Verbs: \textit{dormivano} = indicativo imperfetto of dormire, loro; main clause. (\S\ref{les:imperfetto}); \textit{poteva dormire} = indicativo imperfetto of potere + dormire, modal + infinitive, lui/lei; main clause. (\S\ref{les:imperfetto})}
\item\relax \textit{«Mamma, non ho sonno» diceva Bruno ogni sera.}\newline “Mum, I'm not sleepy,” Bruno said every evening.\newline {\small\color{muted}diceva: a repeated action → imperfetto.}\newline {\footnotesize Verbs: \textit{ho} = indicativo presente of avere, io; main clause. (\S\ref{les:present-are}); \textit{diceva} = indicativo imperfetto of dire, lui/lei; main clause. (\S\ref{les:imperfetto})}
\item\relax \textit{Una notte Bruno è uscito dalla grotta.}\newline One night Bruno went out of the cave.\newline {\small\color{muted}A single event → passato prossimo; uscire takes essere.}\newline {\footnotesize Verbs: \textit{è uscito} = indicativo passato prossimo of uscire, lui/lei; main clause. (\S\ref{les:passato-prossimo-avere})}
\item\relax \textit{Fuori c'era la neve e nel bosco c'era silenzio.}\newline Outside there was snow and the wood was silent.\newline {\footnotesize Verbs: \textit{era} = indicativo imperfetto of essere, lui/lei; main clause. (\S\ref{les:imperfetto})}
\item\relax \textit{Sopra un ramo c'era un gufo.}\newline On a branch there was an owl.\newline {\footnotesize Verbs: \textit{era} = indicativo imperfetto of essere, lui/lei; main clause. (\S\ref{les:imperfetto})}
\item\relax \textit{«Perché non dormi?» ha chiesto il gufo.}\newline “Why aren't you asleep?” asked the owl.\newline {\small\color{muted}chiedere → chiesto (irregular participle).}\newline {\footnotesize Verbs: \textit{dormi} = indicativo presente of dormire, tu; subordinate clause (perché). (\S\ref{les:present-are}); \textit{ha chiesto} = indicativo passato prossimo of chiedere, lui/lei; main clause. (\S\ref{les:passato-prossimo-avere})}
\item\relax \textit{«Non lo so» ha risposto Bruno. «Ho la pancia vuota.»}\newline “I don't know,” answered Bruno. “My tummy is empty.”\newline {\small\color{muted}rispondere → risposto.}\newline {\footnotesize Verbs: \textit{so} = indicativo presente of sapere, io; main clause. (\S\ref{les:present-are}); \textit{ha risposto} = indicativo passato prossimo of rispondere, lui/lei; main clause. (\S\ref{les:passato-prossimo-avere}); \textit{Ho} = indicativo presente of avere, io; main clause. (\S\ref{les:present-are})}
\item\relax \textit{Il gufo ha riso.}\newline The owl laughed.\newline {\small\color{muted}ridere → riso.}\newline {\footnotesize Verbs: \textit{ha riso} = indicativo passato prossimo of ridere, lui/lei; main clause. (\S\ref{les:passato-prossimo-avere})}
\item\relax \textit{«Gli orsi mangiano molto in autunno, prima dell'inverno.»}\newline “Bears eat a lot in autumn, before the winter.”\newline {\footnotesize Verbs: \textit{mangiano} = indicativo presente of mangiare, loro; main clause. (\S\ref{les:present-are})}
\item\relax \textit{Allora Bruno ha pensato al miele.}\newline So Bruno thought about honey.\newline {\small\color{muted}pensare a = to think about.}\newline {\footnotesize Verbs: \textit{ha pensato} = indicativo passato prossimo of pensare, lui/lei; main clause. (\S\ref{les:passato-prossimo-avere})}
\item\relax \textit{«Le api hanno il miele!» ha detto Bruno.}\newline “The bees have honey!” said Bruno.\newline {\footnotesize Verbs: \textit{hanno} = indicativo presente of avere, loro; main clause. (\S\ref{les:present-are}); \textit{ha detto} = indicativo passato prossimo of dire, lui/lei; main clause. (\S\ref{les:passato-prossimo-avere})}
\item\relax \textit{Ma il gufo ha detto: «In inverno anche le api dormono.»}\newline But the owl said: “In winter the bees sleep too.”\newline {\footnotesize Verbs: \textit{ha detto} = indicativo passato prossimo of dire, lui/lei; main clause. (\S\ref{les:passato-prossimo-avere}); \textit{dormono} = indicativo presente of dormire, loro; main clause. (\S\ref{les:present-are})}
\item\relax \textit{Bruno era triste.}\newline Bruno was sad.\newline {\footnotesize Verbs: \textit{era} = indicativo imperfetto of essere, lui/lei; main clause. (\S\ref{les:imperfetto})}
\item\relax \textit{Il gufo ha guardato la luna e ha cominciato a cantare piano piano.}\newline The owl looked at the moon and began to sing very softly.\newline {\small\color{muted}cominciare a + infinitive.}\newline {\footnotesize Verbs: \textit{ha guardato} = indicativo passato prossimo of guardare, lui/lei; main clause. (\S\ref{les:passato-prossimo-avere}); \textit{ha cominciato} = indicativo passato prossimo of cominciare, lui/lei; main clause. (\S\ref{les:passato-prossimo-avere})}
\item\relax \textit{Era una canzone dolce e lenta.}\newline It was a sweet, slow song.\newline {\footnotesize Verbs: \textit{Era} = indicativo imperfetto of essere, lui/lei; main clause. (\S\ref{les:imperfetto})}
\item\relax \textit{Bruno ha chiuso gli occhi.}\newline Bruno closed his eyes.\newline {\small\color{muted}chiudere → chiuso.}\newline {\footnotesize Verbs: \textit{ha chiuso} = indicativo passato prossimo of chiudere, lui/lei; main clause. (\S\ref{les:passato-prossimo-avere})}
\item\relax \textit{Poi è tornato nella grotta, si è messo vicino alla mamma e si è addormentato.}\newline Then he went back into the cave, settled down next to his mother and fell asleep.\newline {\small\color{muted}Reflexive verbs in the passato prossimo take essere.}\newline {\footnotesize Verbs: \textit{è tornato} = indicativo passato prossimo of tornare, lui/lei; main clause. (\S\ref{les:passato-prossimo-avere}); \textit{è messo} = indicativo passato prossimo of mettere, lui/lei; main clause. (\S\ref{les:passato-prossimo-avere}); \textit{è addormentato} = indicativo passato prossimo of addormentarsi, lui/lei; main clause. (\S\ref{les:passato-prossimo-avere})}
\item\relax \textit{Ha dormito fino alla primavera.}\newline He slept until spring.\newline {\footnotesize Verbs: \textit{Ha dormito} = indicativo passato prossimo of dormire, lui/lei; main clause. (\S\ref{les:passato-prossimo-avere})}
\item\relax \textit{Quando si è svegliato, il sole era caldo e nel prato c'erano molti fiori.}\newline When he woke up, the sun was warm and there were many flowers in the meadow.\newline {\small\color{muted}Event (si è svegliato) + description (era, c'erano).}\newline {\footnotesize Verbs: \textit{è svegliato} = indicativo passato prossimo of svegliarsi, lui/lei; subordinate clause (quando). (\S\ref{les:passato-prossimo-avere}); \textit{era} = indicativo imperfetto of essere, lui/lei; main clause. (\S\ref{les:imperfetto}); \textit{erano} = indicativo imperfetto of essere, loro; main clause. (\S\ref{les:imperfetto})}
\item\relax \textit{E la mamma lo aspettava con il miele.}\newline And his mother was waiting for him with the honey.\newline {\small\color{muted}lo = him; aspettare takes a direct object.}\newline {\footnotesize Verbs: \textit{aspettava} = indicativo imperfetto of aspettare, lui/lei; main clause. (\S\ref{les:imperfetto})}
\end{sentences}
\subsection*{Grammar in this story}
\begin{itemize}
\item\relax \textbf{Imperfetto for background.} Era inverno, faceva freddo, tutti gli orsi dormivano. (\S\ref{les:imperfetto})
\item\relax \textbf{Passato prossimo for events.} è uscito, ha chiesto, ha risposto, ha riso, si è addormentato. (\S\ref{les:imperfetto-passato})
\item\relax \textbf{Irregular participles.} chiesto, risposto, riso, detto, chiuso, messo. (\S\ref{les:passato-prossimo-avere})
\item\relax \textbf{Essere and agreement.} è uscito, è tornato, si è addormentato (Bruno is masculine). (\S\ref{les:passato-prossimo-essere})
\end{itemize}
\subsection*{Verb forms in this story}
\begin{itemize}
\item\relax \textbf{indicativo presente}: \textit{ho}, \textit{dormi}, \textit{so}, \textit{Ho}, \textit{mangiano}, \textit{hanno} (\S\ref{les:present-are})
\item\relax \textbf{indicativo passato prossimo}: \textit{è uscito}, \textit{ha chiesto}, \textit{ha risposto}, \textit{ha riso}, \textit{ha pensato}, \textit{ha detto} (\S\ref{les:passato-prossimo-avere})
\item\relax \textbf{indicativo imperfetto}: \textit{era}, \textit{chiamava}, \textit{abitava}, \textit{Era}, \textit{faceva}, \textit{dormivano} (\S\ref{les:imperfetto})
\item\relax \textbf{infinito}: \textit{cantare} (\S\ref{les:infinito-gerundio})
\end{itemize}
\relatedwords{\textit{l'orso} bear \textperiodcentered{} \textit{la grotta} cave \textperiodcentered{} \textit{il bosco} wood; forest \textperiodcentered{} \textit{l'inverno} winter \textperiodcentered{} \textit{il sonno} sleep; sleepiness \textperiodcentered{} \textit{la neve} snow \textperiodcentered{} \textit{il silenzio} silence \textperiodcentered{} \textit{il ramo} branch \textperiodcentered{} \textit{il gufo} owl \textperiodcentered{} \textit{la pancia} belly; tummy \textperiodcentered{} \textit{l'autunno} autumn \textperiodcentered{} \textit{il miele} honey \textperiodcentered{} \textit{l'ape} bee \textperiodcentered{} \textit{la luna} moon \textperiodcentered{} \textit{la canzone} song \textperiodcentered{} \textit{l'occhio} eye \textperiodcentered{} \textit{la primavera} spring \textperiodcentered{} \textit{il sole} sun \textperiodcentered{} \textit{il fiore} flower \textperiodcentered{} \textit{il prato} meadow; lawn \textperiodcentered{} \textit{vuoto} empty \textperiodcentered{} \textit{dolce} sweet; gentle \textperiodcentered{} \textit{lento} slow \textperiodcentered{} \textit{triste} sad \textperiodcentered{} \textit{abitare} to live (reside) \textperiodcentered{} \textit{potere} to be able to; can; may \textperiodcentered{} \textit{uscire} to go out; to leave \textperiodcentered{} \textit{chiedere} to ask \textperiodcentered{} \textit{rispondere} to answer; to reply \textperiodcentered{} \textit{ridere} to laugh \textperiodcentered{} \textit{pensare} to think \textperiodcentered{} \textit{cominciare} to begin; to start \textperiodcentered{} \textit{chiudere} to close \textperiodcentered{} \textit{addormentarsi} to fall asleep \textperiodcentered{} \textit{svegliarsi} to wake up \textperiodcentered{} \textit{aspettare} to wait (for) \textperiodcentered{} \textit{c'era una volta} once upon a time (there was)}
\subsection*{Questions}\label{story:lorso-che-non-poteva-dormire:questions}
{\color{muted}\small Answers: page~\pageref{ans:story:lorso-che-non-poteva-dormire}.}

\begin{exercises}
\item\label{ex:story-orso-1}\relax Perché Bruno non poteva dormire?\par\nopagebreak\begin{options}
\item\relax Aveva la pancia vuota.
\item\relax C'era rumore.
\item\relax Aveva paura del buio.
\item\relax Faceva troppo caldo.
\end{options}
\item\label{ex:story-orso-2}\relax Chi ha incontrato Bruno fuori dalla grotta?\par\nopagebreak\begin{options}
\item\relax la mamma
\item\relax una volpe
\item\relax un'ape
\item\relax un gufo
\end{options}
\item\label{ex:story-orso-3}\relax Che cosa ha fatto il gufo per aiutare Bruno?\par\nopagebreak\begin{options}
\item\relax Ha cantato una canzone.
\item\relax Gli ha dato il miele.
\item\relax Lo ha portato a casa.
\item\relax Ha svegliato le api.
\end{options}
\item\label{ex:story-orso-4}\relax Why is \textit{Era inverno} in the imperfetto?\par\nopagebreak\begin{options}
\item\relax It is in the future.
\item\relax Because of the word inverno.
\item\relax It is a single completed event.
\item\relax It describes the background.
\end{options}
\item\label{ex:story-orso-5}\relax What is the infinitive of \textit{riso} (in \textit{ha riso})?\par\answerline
\item\label{ex:story-orso-6}\relax Until when did Bruno sleep?\par\nopagebreak\begin{options}
\item\relax until spring
\item\relax until morning
\item\relax until autumn
\item\relax for one hour
\end{options}
\end{exercises}

\section{Il segreto del faro}\label{story:il-segreto-del-faro}
\storymeta{The secret of the lighthouse}{A2–B1}{An old fisherman lights an abandoned lighthouse every night. During a storm his granddaughter discovers why.}
\subsection*{Text}
\begin{storytext}
\sentnum{1}Molti anni fa, in un piccolo paese sul mare, viveva un vecchio pescatore che si chiamava Tommaso. \sentnum{2}Ogni sera Tommaso saliva sul faro e accendeva la luce. \sentnum{3}Nessuno capiva perché lo facesse: il faro era vecchio e le navi non passavano più da lì.\par
\sentnum{4}Una notte d'inverno ci fu una grande tempesta. \sentnum{5}Il vento soffiava forte e il mare era nero. \sentnum{6}Elena, la nipote di Tommaso, vide una piccola luce in mezzo alle onde. \sentnum{7}«Nonno, c'è una barca là fuori!» gridò. \sentnum{8}Tommaso prese la lanterna e corse verso il faro. \sentnum{9}Quando arrivarono in cima, la barca era già vicina agli scogli. \sentnum{10}Tommaso accese la grande luce e la barca cambiò direzione.\par
\sentnum{11}La mattina dopo un giovane pescatore bussò alla porta. \sentnum{12}Disse che si era perso nella tempesta e che senza la luce del faro sarebbe morto. \sentnum{13}Tommaso sorrise, ma non disse niente. \sentnum{14}Dopo che il giovane fu partito, Elena chiese al nonno perché avesse acceso la luce per tanti anni.\par
\sentnum{15}Tommaso le raccontò che, quando era giovane, suo padre era uscito in mare una sera e non era più tornato. \sentnum{16}«Quella notte il faro era spento», disse. \sentnum{17}«Quel giorno ho promesso che avrei acceso la luce ogni sera, per tutte le barche del mare.» \sentnum{18}Così Elena capì il segreto del faro. \sentnum{19}Oggi Tommaso non c'è più, ma ogni sera Elena sale sul faro e accende la luce.\par
\end{storytext}
\subsection*{Translation, notes and verb analysis}
\begin{sentences}
\item\relax \textit{Molti anni fa, in un piccolo paese sul mare, viveva un vecchio pescatore che si chiamava Tommaso.}\newline Many years ago, in a small village by the sea, there lived an old fisherman called Tommaso.\newline {\small\color{muted}Imperfetto for the setting (viveva, si chiamava).}\newline {\footnotesize Verbs: \textit{viveva} = indicativo imperfetto of vivere, lui/lei; main clause. (\S\ref{les:imperfetto}); \textit{chiamava} = indicativo imperfetto of chiamare, lui/lei; relative clause. (\S\ref{les:imperfetto})}
\item\relax \textit{Ogni sera Tommaso saliva sul faro e accendeva la luce.}\newline Every evening Tommaso climbed up the lighthouse and switched on the light.\newline {\small\color{muted}Imperfetto for a habit.}\newline {\footnotesize Verbs: \textit{saliva} = indicativo imperfetto of salire, lui/lei; main clause. (\S\ref{les:imperfetto}); \textit{accendeva} = indicativo imperfetto of accendere, lui/lei; main clause. (\S\ref{les:imperfetto})}
\item\relax \textit{Nessuno capiva perché lo facesse: il faro era vecchio e le navi non passavano più da lì.}\newline Nobody understood why he did it: the lighthouse was old and ships no longer passed that way.\newline {\small\color{muted}facesse: congiuntivo imperfetto in an indirect question after a past verb. The indicative (perché lo faceva) is also used, in a less formal register.}\newline {\footnotesize Verbs: \textit{capiva} = indicativo imperfetto of capire, lui/lei; main clause. (\S\ref{les:imperfetto}); \textit{facesse} = congiuntivo imperfetto of fare, lui/lei; subordinate clause (perché). (\S\ref{les:congiuntivo-imperfetto-trapassato}); \textit{era} = indicativo imperfetto of essere, lui/lei; main clause. (\S\ref{les:imperfetto}); \textit{passavano} = indicativo imperfetto of passare, loro; main clause. (\S\ref{les:imperfetto})}
\item\relax \textit{Una notte d'inverno ci fu una grande tempesta.}\newline One winter night there was a great storm.\newline {\small\color{muted}Passato remoto (ci fu): the story's first event. In speech many Italians would say c'è stata.}\newline {\footnotesize Verbs: \textit{fu} = indicativo passato remoto of essere, lui/lei; main clause. (\S\ref{les:passato-remoto})}
\item\relax \textit{Il vento soffiava forte e il mare era nero.}\newline The wind was blowing hard and the sea was black.\newline {\small\color{muted}Imperfetto: the background to the events.}\newline {\footnotesize Verbs: \textit{soffiava} = indicativo imperfetto of soffiare, lui/lei; main clause. (\S\ref{les:imperfetto}); \textit{era} = indicativo imperfetto of essere, lui/lei; main clause. (\S\ref{les:imperfetto})}
\item\relax \textit{Elena, la nipote di Tommaso, vide una piccola luce in mezzo alle onde.}\newline Elena, Tommaso's granddaughter, saw a small light among the waves.\newline {\small\color{muted}vide: passato remoto of vedere (vid-, 1-3-3 pattern).}\newline {\footnotesize Verbs: \textit{vide} = indicativo passato remoto of vedere, lui/lei; main clause. (\S\ref{les:passato-remoto})}
\item\relax \textit{«Nonno, c'è una barca là fuori!» gridò.}\newline "Grandpa, there's a boat out there!" she shouted.\newline {\small\color{muted}Direct speech in the present; the narration around it is in the passato remoto.}\newline {\footnotesize Verbs: \textit{è} = indicativo presente of essere, lui/lei; main clause. (\S\ref{les:present-are}); \textit{gridò} = indicativo passato remoto of gridare, lui/lei; main clause. (\S\ref{les:passato-remoto})}
\item\relax \textit{Tommaso prese la lanterna e corse verso il faro.}\newline Tommaso took the lantern and ran towards the lighthouse.\newline {\small\color{muted}prese, corse: strong passato remoto forms.}\newline {\footnotesize Verbs: \textit{prese} = indicativo passato remoto of prendere, lui/lei; main clause. (\S\ref{les:passato-remoto}); \textit{corse} = indicativo passato remoto of correre, lui/lei; main clause. (\S\ref{les:passato-remoto})}
\item\relax \textit{Quando arrivarono in cima, la barca era già vicina agli scogli.}\newline When they reached the top, the boat was already close to the rocks.\newline {\footnotesize Verbs: \textit{arrivarono} = indicativo passato remoto of arrivare, loro; subordinate clause (quando). (\S\ref{les:passato-remoto}); \textit{era} = indicativo imperfetto of essere, lui/lei; main clause. (\S\ref{les:imperfetto})}
\item\relax \textit{Tommaso accese la grande luce e la barca cambiò direzione.}\newline Tommaso switched on the great light and the boat changed direction.\newline {\small\color{muted}accese (accendere), cambiò (regular -are, final stress).}\newline {\footnotesize Verbs: \textit{accese}: verb, indicativo passato remoto, 3rd person singular (lui/lei) of accendere | verb, participio passato (feminine plural) — compound tenses take avere; also passive and adjectival use of accendere; \textit{cambiò}: verb, indicativo passato remoto, 3rd person singular (lui/lei) of cambiare}
\item\relax \textit{La mattina dopo un giovane pescatore bussò alla porta.}\newline The next morning a young fisherman knocked at the door.\newline {\footnotesize Verbs: \textit{bussò} = indicativo passato remoto of bussare, lui/lei; main clause. (\S\ref{les:passato-remoto})}
\item\relax \textit{Disse che si era perso nella tempesta e che senza la luce del faro sarebbe morto.}\newline He said that he had got lost in the storm and that without the light of the lighthouse he would have died.\newline {\small\color{muted}Reported speech: «Mi sono perso» → si era perso (trapassato); «sarei morto» → sarebbe morto (condizionale passato).}\newline {\footnotesize Verbs: \textit{Disse} = indicativo passato remoto of dire, lui/lei; main clause. (\S\ref{les:passato-remoto}); \textit{era perso} = indicativo trapassato prossimo of perdere, lui/lei; subordinate clause (che). (\S\ref{les:trapassato-prossimo}); \textit{sarebbe morto} = condizionale passato of morire, lui/lei; subordinate clause (che). (\S\ref{les:condizionale-passato})}
\item\relax \textit{Tommaso sorrise, ma non disse niente.}\newline Tommaso smiled, but said nothing.\newline {\small\color{muted}sorrise, disse: strong passato remoto forms.}\newline {\footnotesize Verbs: \textit{sorrise} = indicativo passato remoto of sorridere, lui/lei; main clause. (\S\ref{les:passato-remoto}); \textit{disse} = indicativo passato remoto of dire, lui/lei; main clause. (\S\ref{les:passato-remoto})}
\item\relax \textit{Dopo che il giovane fu partito, Elena chiese al nonno perché avesse acceso la luce per tanti anni.}\newline After the young man had left, Elena asked her grandfather why he had switched on the light for so many years.\newline {\small\color{muted}fu partito: trapassato remoto (literary anteriority after dopo che). avesse acceso: congiuntivo trapassato in an indirect question.}\newline {\footnotesize Verbs: \textit{fu partito} = indicativo trapassato remoto of partire, lui/lei; main clause. (\S\ref{les:trapassato-remoto}); \textit{chiese} = indicativo passato remoto of chiedere, lui/lei; main clause. (\S\ref{les:passato-remoto}); \textit{avesse acceso} = congiuntivo trapassato of accendere, lui/lei; subordinate clause (perché). (\S\ref{les:congiuntivo-imperfetto-trapassato})}
\item\relax \textit{Tommaso le raccontò che, quando era giovane, suo padre era uscito in mare una sera e non era più tornato.}\newline Tommaso told her that, when he was young, his father had gone out to sea one evening and had never come back.\newline {\small\color{muted}Trapassato prossimo for events before the time of the story.}\newline {\footnotesize Verbs: \textit{raccontò} = indicativo passato remoto of raccontare, lui/lei; main clause. (\S\ref{les:passato-remoto}); \textit{era} = indicativo imperfetto of essere, lui/lei; subordinate clause (quando). (\S\ref{les:imperfetto}); \textit{era uscito} = indicativo trapassato prossimo of uscire, lui/lei; main clause. (\S\ref{les:trapassato-prossimo}); \textit{era tornato} = indicativo trapassato prossimo of tornare, lui/lei; main clause. (\S\ref{les:trapassato-prossimo})}
\item\relax \textit{«Quella notte il faro era spento», disse.}\newline "That night the lighthouse was dark," he said.\newline {\small\color{muted}spento: participle of spegnere used as an adjective.}\newline {\footnotesize Verbs: \textit{era spento} = indicativo imperfetto of spegnere, passive, lui/lei; main clause. (\S\ref{les:imperfetto}); \textit{disse} = indicativo passato remoto of dire, lui/lei; main clause. (\S\ref{les:passato-remoto})}
\item\relax \textit{«Quel giorno ho promesso che avrei acceso la luce ogni sera, per tutte le barche del mare.»}\newline "That day I promised that I would switch on the light every evening, for all the boats on the sea."\newline {\small\color{muted}In his own speech Tommaso uses the passato prossimo; avrei acceso is the future in the past (condizionale passato).}\newline {\footnotesize Verbs: \textit{ho promesso} = indicativo passato prossimo of promettere, io; main clause. (\S\ref{les:passato-prossimo-avere}); \textit{avrei acceso} = condizionale passato of accendere, io; subordinate clause (che). (\S\ref{les:condizionale-passato})}
\item\relax \textit{Così Elena capì il segreto del faro.}\newline That is how Elena understood the secret of the lighthouse.\newline {\footnotesize Verbs: \textit{capì} = indicativo passato remoto of capire, lui/lei; main clause. (\S\ref{les:passato-remoto})}
\item\relax \textit{Oggi Tommaso non c'è più, ma ogni sera Elena sale sul faro e accende la luce.}\newline Today Tommaso is no longer here, but every evening Elena climbs the lighthouse and switches on the light.\newline {\small\color{muted}The present brings the story up to today.}\newline {\footnotesize Verbs: \textit{è} = indicativo presente of essere, lui/lei; main clause. (\S\ref{les:present-are}); \textit{sale} = indicativo presente of salire, lui/lei; main clause. (\S\ref{les:present-are}); \textit{accende} = indicativo presente of accendere, lui/lei; main clause. (\S\ref{les:present-are})}
\end{sentences}
\subsection*{Grammar in this story}
\begin{itemize}
\item\relax \textbf{Passato remoto for the events.} ci fu, vide, gridò, prese, corse, arrivarono, accese, cambiò, bussò, disse, sorrise, chiese, raccontò, capì. (\S\ref{les:passato-remoto})
\item\relax \textbf{Imperfetto for background and habits.} viveva, saliva, accendeva, capiva, soffiava, era. (\S\ref{les:imperfetto})
\item\relax \textbf{Trapassato for earlier events.} si era perso, era uscito, non era più tornato; trapassato remoto: fu partito. (\S\ref{les:trapassato-prossimo})
\item\relax \textbf{Reported speech and the future in the past.} Disse che si era perso … e che sarebbe morto; ho promesso che avrei acceso. (\S\ref{les:discorso-indiretto})
\item\relax \textbf{Congiuntivo in indirect questions.} perché lo facesse; perché avesse acceso. (\S\ref{les:congiuntivo-imperfetto-trapassato})
\end{itemize}
\subsection*{Verb forms in this story}
\begin{itemize}
\item\relax \textbf{indicativo presente}: \textit{fa}, \textit{è}, \textit{sale}, \textit{accende} (\S\ref{les:present-are})
\item\relax \textbf{indicativo passato prossimo}: \textit{ho promesso} (\S\ref{les:passato-prossimo-avere})
\item\relax \textbf{indicativo imperfetto}: \textit{viveva}, \textit{chiamava}, \textit{saliva}, \textit{accendeva}, \textit{capiva}, \textit{era} (\S\ref{les:imperfetto})
\item\relax \textbf{indicativo imperfetto (passive)}: \textit{era spento} (\S\ref{les:forma-passiva})
\item\relax \textbf{indicativo trapassato prossimo}: \textit{era perso}, \textit{era uscito}, \textit{era tornato} (\S\ref{les:trapassato-prossimo})
\item\relax \textbf{indicativo passato remoto}: \textit{fu}, \textit{vide}, \textit{gridò}, \textit{prese}, \textit{corse}, \textit{arrivarono} (\S\ref{les:passato-remoto})
\item\relax \textbf{indicativo trapassato remoto}: \textit{fu partito} (\S\ref{les:trapassato-remoto})
\item\relax \textbf{congiuntivo imperfetto}: \textit{facesse} (\S\ref{les:congiuntivo-imperfetto-trapassato})
\item\relax \textbf{congiuntivo trapassato}: \textit{avesse acceso} (\S\ref{les:congiuntivo-imperfetto-trapassato})
\item\relax \textbf{condizionale passato}: \textit{sarebbe morto}, \textit{avrei acceso} (\S\ref{les:condizionale-passato})
\end{itemize}
\relatedwords{\textit{il pescatore} fisherman \textperiodcentered{} \textit{il faro} lighthouse \textperiodcentered{} \textit{la luce} light \textperiodcentered{} \textit{la nave} ship \textperiodcentered{} \textit{la tempesta} storm \textperiodcentered{} \textit{il nipote} grandchild; nephew/niece \textperiodcentered{} \textit{la lanterna} lantern \textperiodcentered{} \textit{la cima} top; summit \textperiodcentered{} \textit{lo scoglio} rock (in the sea) \textperiodcentered{} \textit{la direzione} direction \textperiodcentered{} \textit{il segreto} secret \textperiodcentered{} \textit{il mare} sea \textperiodcentered{} \textit{la barca} boat \textperiodcentered{} \textit{il vento} wind \textperiodcentered{} \textit{l'onda} wave \textperiodcentered{} \textit{la porta} door \textperiodcentered{} \textit{vecchio} old \textperiodcentered{} \textit{vicino} near; close \textperiodcentered{} \textit{vivere} to live \textperiodcentered{} \textit{salire} to go up; to climb; to get on \textperiodcentered{} \textit{accendere} to switch on; to light \textperiodcentered{} \textit{spegnere} to switch off; to put out \textperiodcentered{} \textit{soffiare} to blow \textperiodcentered{} \textit{gridare} to shout \textperiodcentered{} \textit{bussare} to knock \textperiodcentered{} \textit{correre} to run \textperiodcentered{} \textit{sorridere} to smile \textperiodcentered{} \textit{morire} to die \textperiodcentered{} \textit{raccontare} to tell (a story) \textperiodcentered{} \textit{promettere} to promise \textperiodcentered{} \textit{capire} to understand \textperiodcentered{} \textit{lì} there \textperiodcentered{} \textit{verso} towards; around (time)}
\subsection*{Questions}\label{story:il-segreto-del-faro:questions}
{\color{muted}\small Answers: page~\pageref{ans:story:il-segreto-del-faro}.}

\begin{exercises}
\item\label{ex:story-faro-1}\relax Che cosa faceva Tommaso ogni sera?\par\nopagebreak\begin{options}
\item\relax Accendeva la luce del faro.
\item\relax Andava a pescare.
\item\relax Scriveva una lettera.
\item\relax Guardava le navi nel porto.
\end{options}
\item\label{ex:story-faro-2}\relax Che cosa vide Elena durante la tempesta?\par\nopagebreak\begin{options}
\item\relax il nonno sugli scogli
\item\relax una nave grande nel porto
\item\relax una piccola luce in mezzo alle onde
\item\relax un pescatore alla porta
\end{options}
\item\label{ex:story-faro-3}\relax Perché Tommaso accendeva la luce ogni sera?\par\nopagebreak\begin{options}
\item\relax Perché aveva paura del buio.
\item\relax Perché suo padre non era tornato una notte in cui il faro era spento.
\item\relax Perché lo voleva Elena.
\item\relax Perché lo pagavano.
\end{options}
\item\label{ex:story-faro-4}\relax \textit{Disse che si era perso nella tempesta.} Which tense is \textit{si era perso}, and why?\par\nopagebreak\begin{options}
\item\relax passato prossimo: it is linked to the present
\item\relax imperfetto: it is background
\item\relax trapassato prossimo: it happened before disse
\item\relax passato remoto: it is a main event
\end{options}
\item\label{ex:story-faro-5}\relax Which infinitive does the passato remoto \textit{sorrise} come from?\par\answerline
\item\label{ex:story-faro-6}\relax \textit{… sarebbe morto.} What does the condizionale passato express here?\par\nopagebreak\begin{options}
\item\relax an unreal past consequence (he did not die)
\item\relax a future event
\item\relax a supposition about the present
\item\relax a polite request
\end{options}
\end{exercises}

\section{Una lettera dalla nonna}\label{story:una-lettera-dalla-nonna}
\storymeta{A letter from Grandma}{A2–B1}{Grandma Rosa writes from Sicily to her granddaughter, who is coming to stay for the summer.}
\subsection*{Text}
\begin{storytext}
\sentnum{1}Cara Giulia, ti scrivo da Palermo, dove fa già caldo anche se siamo solo ad aprile. \sentnum{2}La mamma mi ha detto che a giugno finirai la scuola e che vorresti venire in Sicilia. \sentnum{3}Sono molto contenta e spero che tu possa restare almeno un mese.\par
\sentnum{4}Quando arriverai, ti porterò al mare e ti insegnerò a nuotare. \sentnum{5}Se farà bel tempo, andremo anche in montagna. \sentnum{6}Il nonno ti farebbe vedere il suo giardino, ma pensa che tu preferisca il mare. \sentnum{7}Se avessi la tua età, passerei tutte le giornate in spiaggia!\par
\sentnum{8}Quando ero giovane, non sapevo nuotare. \sentnum{9}Se qualcuno mi avesse insegnato, avrei nuotato tutta la vita. \sentnum{10}Per questo voglio che tu impari bene.\par
\sentnum{11}Ti chiedo solo una cosa: porta un cappello e non dimenticare la crema per il sole. \sentnum{12}È importante che tu non stia troppo al sole nelle ore più calde. \sentnum{13}Prima che tu parta, ti manderò una mappa dell'isola. \sentnum{14}Il nonno vuole sapere se ti piacerebbe imparare a pescare. \sentnum{15}Rispondi presto, per favore, e scrivi quando arriverà il tuo treno. \sentnum{16}Un abbraccio forte, la tua nonna Rosa.\par
\end{storytext}
\subsection*{Translation, notes and verb analysis}
\begin{sentences}
\item\relax \textit{Cara Giulia, ti scrivo da Palermo, dove fa già caldo anche se siamo solo ad aprile.}\newline Dear Giulia, I'm writing to you from Palermo, where it is already hot even though it's only April.\newline {\footnotesize Verbs: \textit{scrivo} = indicativo presente of scrivere, io; main clause. (\S\ref{les:present-are}); \textit{fa} = indicativo presente of fare, lui/lei; main clause. (\S\ref{les:present-are}); \textit{siamo} = congiuntivo presente of essere, noi; subordinate clause (se). (\S\ref{les:congiuntivo-presente})}
\item\relax \textit{La mamma mi ha detto che a giugno finirai la scuola e che vorresti venire in Sicilia.}\newline Mum told me that you will finish school in June and that you would like to come to Sicily.\newline {\small\color{muted}Reported speech: futuro (finirai) and condizionale of politeness/wish (vorresti) keep their tense because the information is still valid.}\newline {\footnotesize Verbs: \textit{ha detto} = indicativo passato prossimo of dire, lui/lei; main clause. (\S\ref{les:passato-prossimo-avere}); \textit{finirai} = indicativo futuro semplice of finire, tu; subordinate clause (che). (\S\ref{les:futuro-semplice}); \textit{vorresti venire} = condizionale presente of volere + venire, modal + infinitive, tu; subordinate clause (che). (\S\ref{les:condizionale-presente})}
\item\relax \textit{Sono molto contenta e spero che tu possa restare almeno un mese.}\newline I am very happy and I hope you can stay at least a month.\newline {\small\color{muted}sperare che + congiuntivo presente (possa).}\newline {\footnotesize Verbs: \textit{Sono} = indicativo presente of essere, io; main clause. (\S\ref{les:present-are}); \textit{spero} = indicativo presente of sperare, io; main clause. (\S\ref{les:present-are}); \textit{possa restare} = congiuntivo presente of potere + restare, modal + infinitive, tu; subordinate clause (che). (\S\ref{les:congiuntivo-presente})}
\item\relax \textit{Quando arriverai, ti porterò al mare e ti insegnerò a nuotare.}\newline When you arrive, I'll take you to the sea and teach you to swim.\newline {\small\color{muted}Futuro in both clauses: Italian uses the future after quando when the time is future.}\newline {\footnotesize Verbs: \textit{arriverai} = indicativo futuro semplice of arrivare, tu; subordinate clause (quando). (\S\ref{les:futuro-semplice}); \textit{porterò} = indicativo futuro semplice of portare, io; main clause. (\S\ref{les:futuro-semplice}); \textit{insegnerò} = indicativo futuro semplice of insegnare, io; main clause. (\S\ref{les:futuro-semplice})}
\item\relax \textit{Se farà bel tempo, andremo anche in montagna.}\newline If the weather is good, we'll also go to the mountains.\newline {\small\color{muted}Periodo ipotetico of reality; se fa bel tempo is equally correct.}\newline {\footnotesize Verbs: \textit{farà} = indicativo futuro semplice of fare, lui/lei; subordinate clause (se). (\S\ref{les:futuro-semplice}); \textit{andremo} = indicativo futuro semplice of andare, noi; main clause. (\S\ref{les:futuro-semplice})}
\item\relax \textit{Il nonno ti farebbe vedere il suo giardino, ma pensa che tu preferisca il mare.}\newline Grandpa would show you his garden, but he thinks you prefer the sea.\newline {\small\color{muted}farebbe: condizionale; pensare che + congiuntivo (preferisca).}\newline {\footnotesize Verbs: \textit{farebbe vedere} = condizionale presente of fare + vedere, verb + infinitive, lui/lei; main clause. (\S\ref{les:condizionale-presente}); \textit{pensa} = indicativo presente of pensare, lui/lei; main clause. (\S\ref{les:present-are}); \textit{preferisca} = congiuntivo presente of preferire, tu; subordinate clause (che). (\S\ref{les:congiuntivo-presente})}
\item\relax \textit{Se avessi la tua età, passerei tutte le giornate in spiaggia!}\newline If I were your age, I would spend every day on the beach!\newline {\small\color{muted}Periodo ipotetico of possibility/unreality: se + congiuntivo imperfetto, condizionale presente.}\newline {\footnotesize Verbs: \textit{avessi} = congiuntivo imperfetto of avere, io; subordinate clause (se). (\S\ref{les:congiuntivo-imperfetto-trapassato}); \textit{passerei} = condizionale presente of passare, io; main clause. (\S\ref{les:condizionale-presente})}
\item\relax \textit{Quando ero giovane, non sapevo nuotare.}\newline When I was young, I couldn't swim.\newline {\footnotesize Verbs: \textit{ero} = indicativo imperfetto of essere, io; subordinate clause (quando). (\S\ref{les:imperfetto}); \textit{sapevo nuotare} = indicativo imperfetto of sapere + nuotare, modal + infinitive, io; main clause. (\S\ref{les:imperfetto})}
\item\relax \textit{Se qualcuno mi avesse insegnato, avrei nuotato tutta la vita.}\newline If someone had taught me, I would have swum all my life.\newline {\small\color{muted}Counterfactual past: se + congiuntivo trapassato, condizionale passato. Informally also: se qualcuno mi insegnava, nuotavo (substandard double imperfetto).}\newline {\footnotesize Verbs: \textit{avesse insegnato} = congiuntivo trapassato of insegnare, lui/lei; subordinate clause (se). (\S\ref{les:congiuntivo-imperfetto-trapassato}); \textit{avrei nuotato} = condizionale passato of nuotare, io; main clause. (\S\ref{les:condizionale-passato})}
\item\relax \textit{Per questo voglio che tu impari bene.}\newline That's why I want you to learn well.\newline {\small\color{muted}volere che + congiuntivo (impari).}\newline {\footnotesize Verbs: \textit{voglio} = indicativo presente of volere, io; main clause. (\S\ref{les:present-are}); \textit{impari} = congiuntivo presente of imparare, tu; subordinate clause (che). (\S\ref{les:congiuntivo-presente})}
\item\relax \textit{Ti chiedo solo una cosa: porta un cappello e non dimenticare la crema per il sole.}\newline I ask you only one thing: bring a hat and don't forget the sun cream.\newline {\small\color{muted}Imperativo (porta) and negative imperativo for tu: non + infinitive (non dimenticare).}\newline {\footnotesize Verbs: \textit{chiedo} = indicativo presente of chiedere, io; main clause. (\S\ref{les:present-are})}
\item\relax \textit{È importante che tu non stia troppo al sole nelle ore più calde.}\newline It's important that you don't stay in the sun too long in the hottest hours.\newline {\small\color{muted}È importante che + congiuntivo (stia, from stare).}\newline {\footnotesize Verbs: \textit{È} = indicativo presente of essere, lui/lei; main clause. (\S\ref{les:present-are}); \textit{stia} = congiuntivo presente of stare, tu; subordinate clause (che). (\S\ref{les:congiuntivo-presente})}
\item\relax \textit{Prima che tu parta, ti manderò una mappa dell'isola.}\newline Before you leave, I'll send you a map of the island.\newline {\small\color{muted}prima che always takes the congiuntivo.}\newline {\footnotesize Verbs: \textit{parta} = congiuntivo presente of partire, tu; main clause. (\S\ref{les:congiuntivo-presente}); \textit{manderò} = indicativo futuro semplice of mandare, io; main clause. (\S\ref{les:futuro-semplice})}
\item\relax \textit{Il nonno vuole sapere se ti piacerebbe imparare a pescare.}\newline Grandpa wants to know if you would like to learn to fish.\newline {\small\color{muted}Indirect question with se; piacerebbe: condizionale of piacere.}\newline {\footnotesize Verbs: \textit{vuole sapere} = indicativo presente of volere + sapere, modal + infinitive, lui/lei; main clause. (\S\ref{les:present-are}); \textit{piacerebbe imparare} = condizionale presente of piacere + imparare, verb + infinitive, lui/lei; subordinate clause (se). (\S\ref{les:condizionale-presente})}
\item\relax \textit{Rispondi presto, per favore, e scrivi quando arriverà il tuo treno.}\newline Reply soon, please, and write when your train will arrive.\newline {\footnotesize Verbs: \textit{Rispondi} = indicativo presente of rispondere, tu; main clause. (\S\ref{les:present-are}); \textit{scrivi} = indicativo presente of scrivere, tu; main clause. (\S\ref{les:present-are}); \textit{arriverà} = indicativo futuro semplice of arrivare, lui/lei; subordinate clause (quando). (\S\ref{les:futuro-semplice})}
\item\relax \textit{Un abbraccio forte, la tua nonna Rosa.}\newline A big hug, your Grandma Rosa.
\end{sentences}
\subsection*{Grammar in this story}
\begin{itemize}
\item\relax \textbf{Futuro.} finirai, arriverai, porterò, insegnerò, farà, andremo, manderò, arriverà. (\S\ref{les:futuro-semplice})
\item\relax \textbf{Condizionale.} vorresti, farebbe, passerei, avrei nuotato, piacerebbe. (\S\ref{les:condizionale-presente})
\item\relax \textbf{Congiuntivo.} spero che tu possa; pensa che tu preferisca; voglio che tu impari; è importante che tu non stia; prima che tu parta. (\S\ref{les:congiuntivo-presente})
\item\relax \textbf{Periodo ipotetico.} Se farà bel tempo, andremo …; Se avessi la tua età, passerei …; Se qualcuno mi avesse insegnato, avrei nuotato … (\S\ref{les:periodo-ipotetico})
\item\relax \textbf{Imperativo.} porta, non dimenticare, rispondi, scrivi. (\S\ref{les:imperativo})
\end{itemize}
\subsection*{Verb forms in this story}
\begin{itemize}
\item\relax \textbf{indicativo presente}: \textit{scrivo}, \textit{fa}, \textit{Sono}, \textit{spero}, \textit{pensa}, \textit{voglio} (\S\ref{les:present-are})
\item\relax \textbf{indicativo passato prossimo}: \textit{ha detto} (\S\ref{les:passato-prossimo-avere})
\item\relax \textbf{indicativo imperfetto}: \textit{ero}, \textit{sapevo nuotare} (\S\ref{les:imperfetto})
\item\relax \textbf{indicativo futuro semplice}: \textit{finirai}, \textit{arriverai}, \textit{porterò}, \textit{insegnerò}, \textit{farà}, \textit{andremo} (\S\ref{les:futuro-semplice})
\item\relax \textbf{congiuntivo presente}: \textit{siamo}, \textit{possa restare}, \textit{preferisca}, \textit{impari}, \textit{stia}, \textit{parta} (\S\ref{les:congiuntivo-presente})
\item\relax \textbf{congiuntivo imperfetto}: \textit{avessi} (\S\ref{les:congiuntivo-imperfetto-trapassato})
\item\relax \textbf{congiuntivo trapassato}: \textit{avesse insegnato} (\S\ref{les:congiuntivo-imperfetto-trapassato})
\item\relax \textbf{condizionale presente}: \textit{vorresti venire}, \textit{farebbe vedere}, \textit{passerei}, \textit{piacerebbe imparare} (\S\ref{les:condizionale-presente})
\item\relax \textbf{condizionale passato}: \textit{avrei nuotato} (\S\ref{les:condizionale-passato})
\item\relax \textbf{infinito}: \textit{nuotare}, \textit{dimenticare}, \textit{pescare} (\S\ref{les:infinito-gerundio})
\end{itemize}
\relatedwords{\textit{la lettera} letter \textperiodcentered{} \textit{la scuola} school \textperiodcentered{} \textit{il mare} sea \textperiodcentered{} \textit{il giardino} garden \textperiodcentered{} \textit{l'età} age \textperiodcentered{} \textit{la spiaggia} beach \textperiodcentered{} \textit{la vita} life \textperiodcentered{} \textit{il cappello} hat \textperiodcentered{} \textit{la crema} cream \textperiodcentered{} \textit{il sole} sun \textperiodcentered{} \textit{la mappa} map \textperiodcentered{} \textit{l'isola} island \textperiodcentered{} \textit{il treno} train \textperiodcentered{} \textit{l'abbraccio} hug \textperiodcentered{} \textit{l'aprile} April \textperiodcentered{} \textit{il giugno} June \textperiodcentered{} \textit{importante} important \textperiodcentered{} \textit{sperare} to hope \textperiodcentered{} \textit{restare} to stay; to remain \textperiodcentered{} \textit{insegnare} to teach \textperiodcentered{} \textit{nuotare} to swim \textperiodcentered{} \textit{preferire} to prefer \textperiodcentered{} \textit{dimenticare} to forget \textperiodcentered{} \textit{mandare} to send \textperiodcentered{} \textit{pescare} to fish; to catch (fish) \textperiodcentered{} \textit{imparare} to learn \textperiodcentered{} \textit{rispondere} to answer; to reply \textperiodcentered{} \textit{almeno} at least \textperiodcentered{} \textit{solo} only}
\subsection*{Questions}\label{story:una-lettera-dalla-nonna:questions}
{\color{muted}\small Answers: page~\pageref{ans:story:una-lettera-dalla-nonna}.}

\begin{exercises}
\item\label{ex:story-lettera-1}\relax Che cosa vuole fare la nonna quando arriverà Giulia?\par\nopagebreak\begin{options}
\item\relax Comprarle un cappello.
\item\relax Mandarla in montagna da sola.
\item\relax Insegnarle a nuotare.
\item\relax Portarla a scuola.
\end{options}
\item\label{ex:story-lettera-2}\relax La nonna sapeva nuotare da giovane?\par\nopagebreak\begin{options}
\item\relax Sì, nuotava ogni giorno.
\item\relax Sì, ma non le piaceva.
\item\relax Non lo dice.
\item\relax No, nessuno le aveva insegnato.
\end{options}
\item\label{ex:story-lettera-3}\relax \textit{Se avessi la tua età, passerei tutte le giornate in spiaggia.} What kind of hypothesis is this?\par\nopagebreak\begin{options}
\item\relax a polite request
\item\relax an unreal present hypothesis (she is not young)
\item\relax a real future condition
\item\relax a counterfactual past
\end{options}
\item\label{ex:story-lettera-4}\relax \textit{Prima che tu \blank{}, ti manderò una mappa.}\par\nopagebreak\begin{options}
\item\relax partiresti
\item\relax parta
\item\relax parti
\item\relax partirai
\end{options}
\item\label{ex:story-lettera-5}\relax Make it negative (tu): \textit{Dimentica la crema!} → \blank{} la crema!
\item\label{ex:story-lettera-6}\relax \textit{Se farà bel tempo, andremo in montagna.} Which other version is also correct?\par\nopagebreak\begin{options}
\item\relax Se farebbe bel tempo, andremo in montagna.
\item\relax Se facesse bel tempo, andremo in montagna.
\item\relax Se fa bel tempo, andiamo in montagna.
\item\relax Se farà bel tempo, andassimo in montagna.
\end{options}
\end{exercises}

\chapter{Rhymes}\label{rhymes}
Original rhymes written for this course, each with a literal and a natural translation.\par

\section{La pioggia sul tetto}\label{rhyme:la-pioggia-sul-tetto}
\storymeta{The rain on the roof}{A1}{A cheerful rain rhyme: the cat sleeps, the wind sings, grandmother cooks soup. Form: Four rhyming couplets (AABB).}
\begin{verse}
Piove, piove sul tetto,\\
il gatto dorme nel letto.\\
Cade l'acqua piano piano\\
sulla strada e sul prato lontano.

Canta il vento alla finestra,\\
la nonna prepara la minestra.\\
Piove, piove, che allegria:\\
domani la nuvola va via.
\end{verse}
\subsection*{Line by line}
\begin{sentences}
\item\relax \textit{Piove, piove sul tetto,}\newline Literal: (It) rains, (it) rains on-the roof,\newline Natural: It's raining, raining on the roof,\newline {\small\color{muted}piove: impersonal weather verb — no subject. sul = su + il.}\newline {\footnotesize Verbs: \textit{Piove} = indicativo presente of piovere, lui/lei; main clause. (\S\ref{les:present-are}); \textit{piove} = indicativo presente of piovere, lui/lei; main clause. (\S\ref{les:present-are})}
\item\relax \textit{il gatto dorme nel letto.}\newline Literal: the cat sleeps in-the bed.\newline Natural: the cat is asleep in bed.\newline {\small\color{muted}tetto / letto: rhyme on -etto.}\newline {\footnotesize Verbs: \textit{dorme} = indicativo presente of dormire, lui/lei; main clause. (\S\ref{les:present-are})}
\item\relax \textit{Cade l'acqua piano piano}\newline Literal: Falls the water softly softly\newline Natural: The water falls gently, gently\newline {\small\color{muted}Verb before subject: common in poetry and with verbs like cadere. Doubling an adverb (piano piano) intensifies it.}\newline {\footnotesize Verbs: \textit{Cade} = indicativo presente of cadere, lui/lei; main clause. (\S\ref{les:present-are})}
\item\relax \textit{sulla strada e sul prato lontano.}\newline Literal: on-the street and on-the meadow far.\newline Natural: on the street and the distant meadow.\newline {\small\color{muted}lontano agrees with prato (masculine singular).}
\item\relax \textit{Canta il vento alla finestra,}\newline Literal: Sings the wind at-the window,\newline Natural: The wind sings at the window,\newline {\small\color{muted}alla = a + la.}\newline {\footnotesize Verbs: \textit{Canta} = indicativo presente of cantare, lui/lei; main clause. (\S\ref{les:present-are})}
\item\relax \textit{la nonna prepara la minestra.}\newline Literal: the grandmother prepares the soup.\newline Natural: Grandma is making soup.\newline {\small\color{muted}finestra / minestra: rhyme on -estra.}\newline {\footnotesize Verbs: \textit{prepara} = indicativo presente of preparare, lui/lei; main clause. (\S\ref{les:present-are})}
\item\relax \textit{Piove, piove, che allegria:}\newline Literal: (It) rains, (it) rains, what cheerfulness:\newline Natural: It's raining, raining — what fun!\newline {\small\color{muted}che + noun = what (a) …!}\newline {\footnotesize Verbs: \textit{Piove} = indicativo presente of piovere, lui/lei; main clause. (\S\ref{les:present-are}); \textit{piove} = indicativo presente of piovere, lui/lei; main clause. (\S\ref{les:present-are})}
\item\relax \textit{domani la nuvola va via.}\newline Literal: tomorrow the cloud goes away.\newline Natural: tomorrow the cloud will go away.\newline {\small\color{muted}Present tense for the near future. andare via = to go away.}\newline {\footnotesize Verbs: \textit{va} = indicativo presente of andare, lui/lei; main clause. (\S\ref{les:present-are})}
\end{sentences}
\subsection*{Grammar in this rhyme}
\begin{itemize}
\item\relax \textbf{Impersonal weather verbs.} piove = it rains / it is raining. Italian has no 'it' subject here. (\S\ref{les:word-order-negation-questions})
\item\relax \textbf{Subject after the verb.} Cade l'acqua, Canta il vento: the subject follows the verb, which gives a poetic rhythm. (\S\ref{les:word-order-negation-questions})
\item\relax \textbf{Articulated prepositions.} sul tetto, nel letto, sulla strada, alla finestra. (\S\ref{les:articulated-prepositions})
\end{itemize}
\relatedwords{\textit{piovere} to rain \textperiodcentered{} \textit{il tetto} roof \textperiodcentered{} \textit{il gatto} cat \textperiodcentered{} \textit{il letto} bed \textperiodcentered{} \textit{cadere} to fall \textperiodcentered{} \textit{l'acqua} water \textperiodcentered{} \textit{piano} slowly; softly; quietly \textperiodcentered{} \textit{la strada} street; road \textperiodcentered{} \textit{il prato} meadow; lawn \textperiodcentered{} \textit{lontano} far; distant \textperiodcentered{} \textit{il vento} wind \textperiodcentered{} \textit{la finestra} window \textperiodcentered{} \textit{la nonna} grandmother \textperiodcentered{} \textit{la minestra} soup \textperiodcentered{} \textit{l'allegria} cheerfulness; fun \textperiodcentered{} \textit{la nuvola} cloud \textperiodcentered{} \textit{domani} tomorrow}
\subsection*{Exercises}
{\color{muted}\small Answers: page~\pageref{ans:rhyme:la-pioggia-sul-tetto}.}

\begin{exercises}
\item\label{ex:rhyme-pioggia-1}\relax What does \textit{piano piano} mean in 'Cade l'acqua piano piano'?\par\nopagebreak\begin{options}
\item\relax upstairs
\item\relax gently, slowly
\item\relax very loudly
\item\relax on the floor
\end{options}
\item\label{ex:rhyme-pioggia-2}\relax Who is singing at the window?\par\nopagebreak\begin{options}
\item\relax the rain
\item\relax the wind
\item\relax the cat
\item\relax the grandmother
\end{options}
\item\label{ex:rhyme-pioggia-3}\relax Infinitive of \textit{cade}?\par\answerline
\item\label{ex:rhyme-pioggia-4}\relax 'domani la nuvola va via' — when does the cloud go away?\par\nopagebreak\begin{options}
\item\relax yesterday
\item\relax now
\item\relax never
\item\relax tomorrow
\end{options}
\item\label{ex:rhyme-pioggia-5}\relax Complete the rhyme: il gatto dorme nel \blank{}.
\end{exercises}

\section{Ninna nanna della luna}\label{rhyme:ninna-nanna-della-luna}
\storymeta{Lullaby of the moon}{A1}{A mother sings her little bear to sleep under the moon and stars. Form: Lullaby in rhyming couplets.}
\begin{verse}
Dormi, dormi, piccolo orso,\\
la luna guarda il bosco.\\
Brilla in cielo una stella,\\
la notte è calma e bella.

Il gufo canta piano piano,\\
il giorno è ancora lontano.\\
Chiudi gli occhi, vicino al cuore:\\
la mamma canta con amore.
\end{verse}
\subsection*{Line by line}
\begin{sentences}
\item\relax \textit{Dormi, dormi, piccolo orso,}\newline Literal: Sleep, sleep, little bear,\newline Natural: Sleep now, sleep, little bear,\newline {\small\color{muted}dormi here is an imperative (an instruction to tu). For -ere and -ire verbs the tu imperative is identical to the present: dormi = you sleep / sleep!}\newline {\footnotesize Verbs: \textit{Dormi} = indicativo presente of dormire, tu; main clause. (\S\ref{les:present-are}); \textit{dormi} = indicativo presente of dormire, tu; main clause. (\S\ref{les:present-are})}
\item\relax \textit{la luna guarda il bosco.}\newline Literal: the moon looks-at the wood.\newline Natural: the moon is watching over the wood.\newline {\small\color{muted}guardare takes a direct object: no preposition for 'at'.}\newline {\footnotesize Verbs: \textit{guarda} = indicativo presente of guardare, lui/lei; main clause. (\S\ref{les:present-are})}
\item\relax \textit{Brilla in cielo una stella,}\newline Literal: Shines in sky a star,\newline Natural: A star is shining in the sky,\newline {\small\color{muted}in cielo: no article in this fixed phrase.}\newline {\footnotesize Verbs: \textit{Brilla} = indicativo presente of brillare, lui/lei; main clause. (\S\ref{les:present-are})}
\item\relax \textit{la notte è calma e bella.}\newline Literal: the night is calm and beautiful.\newline Natural: the night is calm and lovely.\newline {\small\color{muted}stella / bella: rhyme on -ella. Both adjectives agree with notte (feminine).}\newline {\footnotesize Verbs: \textit{è} = indicativo presente of essere, lui/lei; main clause. (\S\ref{les:present-are})}
\item\relax \textit{Il gufo canta piano piano,}\newline Literal: The owl sings softly softly,\newline Natural: The owl sings very softly,\newline {\footnotesize Verbs: \textit{canta} = indicativo presente of cantare, lui/lei; main clause. (\S\ref{les:present-are})}
\item\relax \textit{il giorno è ancora lontano.}\newline Literal: the day is still far.\newline Natural: morning is still far away.\newline {\small\color{muted}ancora = still. piano / lontano rhyme.}\newline {\footnotesize Verbs: \textit{è} = indicativo presente of essere, lui/lei; main clause. (\S\ref{les:present-are})}
\item\relax \textit{Chiudi gli occhi, vicino al cuore:}\newline Literal: Close the eyes, near to-the heart:\newline Natural: Close your eyes, close to my heart:\newline {\small\color{muted}chiudi: imperative of chiudere (same form as the present tu). gli occhi: Italian uses the article, not the possessive, for body parts.}\newline {\footnotesize Verbs: \textit{Chiudi} = indicativo presente of chiudere, tu; main clause. (\S\ref{les:present-are})}
\item\relax \textit{la mamma canta con amore.}\newline Literal: the mum sings with love.\newline Natural: Mummy sings to you with love.\newline {\small\color{muted}cuore / amore: the most traditional rhyme in Italian song.}\newline {\footnotesize Verbs: \textit{canta} = indicativo presente of cantare, lui/lei; main clause. (\S\ref{les:present-are})}
\end{sentences}
\subsection*{Grammar in this rhyme}
\begin{itemize}
\item\relax \textbf{The tu imperative.} dormi! chiudi! For -ere and -ire verbs the tu imperative is identical to the present tense; for -are verbs it ends in -a (canta!). The full imperative, including the formal and negative forms, is in the imperativo lesson. (\S\ref{les:imperativo})
\item\relax \textbf{Body parts with the article.} Chiudi gli occhi = close your eyes. (\S\ref{les:reflexive-verbs})
\end{itemize}
\relatedwords{\textit{dormire} to sleep \textperiodcentered{} \textit{piccolo} small; little \textperiodcentered{} \textit{l'orso} bear \textperiodcentered{} \textit{la luna} moon \textperiodcentered{} \textit{guardare} to look at; to watch \textperiodcentered{} \textit{il bosco} wood; forest \textperiodcentered{} \textit{brillare} to shine; to sparkle \textperiodcentered{} \textit{il cielo} sky \textperiodcentered{} \textit{la stella} star \textperiodcentered{} \textit{la notte} night \textperiodcentered{} \textit{calmo} calm \textperiodcentered{} \textit{bello} beautiful; lovely; handsome \textperiodcentered{} \textit{il gufo} owl \textperiodcentered{} \textit{cantare} to sing \textperiodcentered{} \textit{il giorno} day \textperiodcentered{} \textit{ancora} still; again; yet \textperiodcentered{} \textit{lontano} far; distant \textperiodcentered{} \textit{chiudere} to close \textperiodcentered{} \textit{l'occhio} eye \textperiodcentered{} \textit{il cuore} heart \textperiodcentered{} \textit{la mamma} mum; mom \textperiodcentered{} \textit{l'amore} love}
\subsection*{Exercises}
{\color{muted}\small Answers: page~\pageref{ans:rhyme:ninna-nanna-della-luna}.}

\begin{exercises}
\item\label{ex:rhyme-ninna-1}\relax In \textit{Dormi, dormi, piccolo orso}, what is \textit{dormi}?\par\nopagebreak\begin{options}
\item\relax an adjective
\item\relax a past tense
\item\relax an instruction: 'sleep!'
\item\relax a noun meaning 'dream'
\end{options}
\item\label{ex:rhyme-ninna-2}\relax What does \textit{il giorno è ancora lontano} tell the little bear?\par\nopagebreak\begin{options}
\item\relax Tomorrow is a holiday.
\item\relax Morning is still a long way off.
\item\relax The day is over.
\item\relax The day is very short.
\end{options}
\item\label{ex:rhyme-ninna-3}\relax In \textit{Brilla in cielo una stella}, what is the subject?\par\nopagebreak\begin{options}
\item\relax in cielo
\item\relax brilla
\item\relax cielo
\item\relax una stella
\end{options}
\item\label{ex:rhyme-ninna-4}\relax Which word rhymes with \textit{cuore} in the last couplet?\par\answerline
\item\label{ex:rhyme-ninna-5}\relax Translate naturally: \textit{la notte è calma e bella}\par\answerline
\end{exercises}

\section{Conta delle stelle}\label{rhyme:conta-delle-stelle}
\storymeta{Counting the stars}{A0}{Count from one to ten with stars, little owls and fish. Form: Counting rhyme (filastrocca).}
\begin{verse}
Una stella, due stelle,\\
tre stelle tanto belle.\\
Quattro, cinque, sei e sette:\\
sul tetto ci sono le civette.

Otto, nove, dieci:\\
nel mare dormono i pesci.\\
Buonanotte, stelle belle,\\
a domani, care sorelle!
\end{verse}
\subsection*{Line by line}
\begin{sentences}
\item\relax \textit{Una stella, due stelle,}\newline Literal: One star, two stars,\newline Natural: One star, two stars,\newline {\small\color{muted}uno → una before a feminine noun; stella → stelle.}
\item\relax \textit{tre stelle tanto belle.}\newline Literal: three stars so beautiful.\newline Natural: three stars, so very pretty.\newline {\small\color{muted}belle agrees with stelle (feminine plural).}
\item\relax \textit{Quattro, cinque, sei e sette:}\newline Literal: Four, five, six and seven:\newline Natural: Four, five, six and seven:\newline {\footnotesize Verbs: \textit{sei}: verb, indicativo presente, 2nd person singular (tu) of essere}
\item\relax \textit{sul tetto ci sono le civette.}\newline Literal: on-the roof there are the little-owls.\newline Natural: there are little owls on the roof.\newline {\small\color{muted}ci sono + plural. sette / civette rhyme.}\newline {\footnotesize Verbs: \textit{sono} = indicativo presente of essere, io; main clause. (\S\ref{les:present-are})}
\item\relax \textit{Otto, nove, dieci:}\newline Literal: Eight, nine, ten:\newline Natural: Eight, nine, ten:
\item\relax \textit{nel mare dormono i pesci.}\newline Literal: in-the sea sleep the fish.\newline Natural: in the sea the fish are sleeping.\newline {\small\color{muted}Subject after the verb. dieci / pesci rhyme by assonance.}\newline {\footnotesize Verbs: \textit{dormono} = indicativo presente of dormire, loro; main clause. (\S\ref{les:present-are})}
\item\relax \textit{Buonanotte, stelle belle,}\newline Literal: Good-night, stars beautiful,\newline Natural: Good night, lovely stars,
\item\relax \textit{a domani, care sorelle!}\newline Literal: until tomorrow, dear sisters!\newline Natural: see you tomorrow, dear sisters!\newline {\small\color{muted}a domani = see you tomorrow. care: feminine plural of caro.}
\end{sentences}
\subsection*{Grammar in this rhyme}
\begin{itemize}
\item\relax \textbf{Plural of -a nouns.} stella → stelle; civetta → civette; sorella → sorelle. (\S\ref{les:noun-gender-number})
\item\relax \textbf{ci sono.} sul tetto ci sono le civette = there are little owls on the roof. (\S\ref{les:indefinite-articles})
\end{itemize}
\relatedwords{\textit{uno} one \textperiodcentered{} \textit{due} two \textperiodcentered{} \textit{tre} three \textperiodcentered{} \textit{quattro} four \textperiodcentered{} \textit{cinque} five \textperiodcentered{} \textit{sei} six \textperiodcentered{} \textit{sette} seven \textperiodcentered{} \textit{otto} eight \textperiodcentered{} \textit{nove} nine \textperiodcentered{} \textit{dieci} ten \textperiodcentered{} \textit{la stella} star \textperiodcentered{} \textit{bello} beautiful; lovely; handsome \textperiodcentered{} \textit{il tetto} roof \textperiodcentered{} \textit{la civetta} little owl \textperiodcentered{} \textit{il mare} sea \textperiodcentered{} \textit{il pesce} fish \textperiodcentered{} \textit{buonanotte} good night \textperiodcentered{} \textit{caro} dear; expensive \textperiodcentered{} \textit{la sorella} sister}
\subsection*{Exercises}
{\color{muted}\small Answers: page~\pageref{ans:rhyme:conta-delle-stelle}.}

\begin{exercises}
\item\label{ex:rhyme-conta-1}\relax Match the numbers.\par\smallskip\begin{tabularx}{\linewidth}{@{}L L@{}}
1.\ tre & (a)\ 9 \\
2.\ sette & (b)\ 4 \\
3.\ nove & (c)\ 7 \\
4.\ quattro & (d)\ 10 \\
5.\ dieci & (e)\ 3 \\
\end{tabularx}
\item\label{ex:rhyme-conta-2}\relax Where are the little owls?\par\nopagebreak\begin{options}
\item\relax in the sky
\item\relax in the wood
\item\relax in the sea
\item\relax on the roof
\end{options}
\item\label{ex:rhyme-conta-3}\relax Plural of \textit{la civetta}: le \blank{}
\item\label{ex:rhyme-conta-4}\relax What does \textit{a domani} mean?\par\nopagebreak\begin{options}
\item\relax goodbye forever
\item\relax today
\item\relax to the morning
\item\relax see you tomorrow
\end{options}
\item\label{ex:rhyme-conta-5}\relax Why \textit{care sorelle} and not \textit{cari sorelle}?\par\nopagebreak\begin{options}
\item\relax sorelle is feminine plural
\item\relax sorelle is singular
\item\relax care is invariable
\item\relax Both are correct
\end{options}
\end{exercises}

\section{Il valzer degli animali}\label{rhyme:il-valzer-degli-animali}
\storymeta{The animals' waltz}{A1}{Animals dance on the roof, in bed and in the meadow; the owl challenges everyone. Form: Dance rhyme in couplets.}
\begin{verse}
Balla il gatto sul tetto,\\
balla il topo sul letto.\\
La rana salta, la lumaca è lenta,\\
la farfalla vola ed è contenta.

Il gufo canta: «Uno, due, tre!\\
Chi balla più veloce di me?»\\
Tutti insieme, piano piano,\\
balliamo il valzer, mano nella mano.
\end{verse}
\subsection*{Line by line}
\begin{sentences}
\item\relax \textit{Balla il gatto sul tetto,}\newline Literal: Dances the cat on-the roof,\newline Natural: The cat is dancing on the roof,\newline {\small\color{muted}Verb–subject order for rhythm.}\newline {\footnotesize Verbs: \textit{Balla} = indicativo presente of ballare, lui/lei; main clause. (\S\ref{les:present-are})}
\item\relax \textit{balla il topo sul letto.}\newline Literal: dances the mouse on-the bed.\newline Natural: the mouse is dancing on the bed.\newline {\footnotesize Verbs: \textit{balla} = indicativo presente of ballare, lui/lei; main clause. (\S\ref{les:present-are})}
\item\relax \textit{La rana salta, la lumaca è lenta,}\newline Literal: The frog jumps, the snail is slow,\newline Natural: The frog jumps, the snail is slow,\newline {\footnotesize Verbs: \textit{salta} = indicativo presente of saltare, lui/lei; main clause. (\S\ref{les:present-are}); \textit{è} = indicativo presente of essere, lui/lei; main clause. (\S\ref{les:present-are})}
\item\relax \textit{la farfalla vola ed è contenta.}\newline Literal: the butterfly flies and is happy.\newline Natural: the butterfly flies and is happy.\newline {\small\color{muted}ed: e before a vowel (here è). lenta / contenta rhyme.}\newline {\footnotesize Verbs: \textit{vola} = indicativo presente of volare, lui/lei; main clause. (\S\ref{les:present-are}); \textit{è} = indicativo presente of essere, lui/lei; main clause. (\S\ref{les:present-are})}
\item\relax \textit{Il gufo canta: «Uno, due, tre!}\newline Literal: The owl sings: “One, two, three!\newline Natural: The owl sings: “One, two, three!\newline {\footnotesize Verbs: \textit{canta} = indicativo presente of cantare, lui/lei; main clause. (\S\ref{les:present-are})}
\item\relax \textit{Chi balla più veloce di me?»}\newline Literal: Who dances more fast than me?”\newline Natural: Who can dance faster than me?”\newline {\small\color{muted}più … di = more … than; di + stressed pronoun me. tre / me rhyme.}\newline {\footnotesize Verbs: \textit{balla} = indicativo presente of ballare, lui/lei; main clause. (\S\ref{les:present-are})}
\item\relax \textit{Tutti insieme, piano piano,}\newline Literal: All together, slowly slowly,\newline Natural: All together, nice and slow,
\item\relax \textit{balliamo il valzer, mano nella mano.}\newline Literal: we-dance the waltz, hand in-the hand.\newline Natural: let's dance the waltz, hand in hand.\newline {\small\color{muted}balliamo: noi form, here a suggestion (let's dance). mano nella mano = hand in hand.}\newline {\footnotesize Verbs: \textit{balliamo} = indicativo presente of ballare, noi; main clause. (\S\ref{les:present-are})}
\end{sentences}
\subsection*{Grammar in this rhyme}
\begin{itemize}
\item\relax \textbf{Comparison.} più veloce di me = faster than me. (\S\ref{les:bello-buono-comparison})
\item\relax \textbf{noi for suggestions.} balliamo = we dance / let's dance. (\S\ref{les:present-are})
\end{itemize}
\relatedwords{\textit{ballare} to dance \textperiodcentered{} \textit{il gatto} cat \textperiodcentered{} \textit{il tetto} roof \textperiodcentered{} \textit{il topo} mouse \textperiodcentered{} \textit{il letto} bed \textperiodcentered{} \textit{la rana} frog \textperiodcentered{} \textit{saltare} to jump \textperiodcentered{} \textit{la lumaca} snail \textperiodcentered{} \textit{lento} slow \textperiodcentered{} \textit{la farfalla} butterfly \textperiodcentered{} \textit{volare} to fly \textperiodcentered{} \textit{contento} glad; pleased \textperiodcentered{} \textit{il gufo} owl \textperiodcentered{} \textit{cantare} to sing \textperiodcentered{} \textit{veloce} fast; quick \textperiodcentered{} \textit{insieme} together \textperiodcentered{} \textit{il valzer} waltz \textperiodcentered{} \textit{la mano} hand}
\subsection*{Exercises}
{\color{muted}\small Answers: page~\pageref{ans:rhyme:il-valzer-degli-animali}.}

\begin{exercises}
\item\label{ex:rhyme-valzer-1}\relax What does the owl ask?\par\nopagebreak\begin{options}
\item\relax Where is the cat?
\item\relax Who wants to sleep?
\item\relax Who dances faster than me?
\item\relax Who sings better than me?
\end{options}
\item\label{ex:rhyme-valzer-2}\relax Infinitive of \textit{vola}?\par\answerline
\item\label{ex:rhyme-valzer-3}\relax \textit{mano nella mano} means…\par\nopagebreak\begin{options}
\item\relax handmade
\item\relax one hand only
\item\relax in the hand of the waltz
\item\relax hand in hand
\end{options}
\item\label{ex:rhyme-valzer-4}\relax Why does the text say \textit{ed è contenta} instead of \textit{e è contenta}?\par\nopagebreak\begin{options}
\item\relax It is a spelling mistake
\item\relax ed means 'but'
\item\relax e becomes ed before a vowel, especially e
\item\relax ed is the feminine of e
\end{options}
\item\label{ex:rhyme-valzer-5}\relax ballare, noi: \blank{}
\end{exercises}

\part{Practice}\label{part:practice}

\chapter{Diagnostic test}\label{diag}
Twenty short questions, one or two per stage of the course, from pronunciation to the congiuntivo. The result suggests where to start; it does not lock any lesson.\par
\begin{notebox}{tipcol}\notelabel{tipcol}{How to read the result}{}Each question tests one module; the answer key names it. Mark your answers, then find the \emph{first} module (in course order) in which you got a question wrong: start there, and skim the modules before it. If every answer is right, review the later modules and concentrate on the Reader.\end{notebox}
{\color{muted}\small Answers: page~\pageref{ans:diag}.}

\begin{exercises}
\item\label{ex:diag-1}\relax Which word has a hard k sound?\par\nopagebreak\begin{options}
\item\relax cielo
\item\relax chiave
\item\relax cena
\item\relax ciao
\end{options}
\item\label{ex:diag-2}\relax Write 'is': Il cane \blank{} grande.
\item\label{ex:diag-3}\relax \blank{} zaino\par\nopagebreak\begin{options}
\item\relax lo
\item\relax l'
\item\relax il
\item\relax la
\end{options}
\item\label{ex:diag-4}\relax Plural of la stella: le \blank{}
\item\label{ex:diag-5}\relax le case (bianco) → le case \blank{}
\item\label{ex:diag-6}\relax parlare, noi: \blank{}
\item\label{ex:diag-7}\relax andare, io: \blank{}
\item\label{ex:diag-8}\relax 'I never cry.'\par\nopagebreak\begin{options}
\item\relax Piango mai.
\item\relax Non piango mai.
\item\relax Non mai piango.
\item\relax Mai non piango.
\end{options}
\item\label{ex:diag-9}\relax in + il bosco → \blank{} bosco
\item\label{ex:diag-10}\relax Noi \blank{} svegliamo presto.\par\nopagebreak\begin{options}
\item\relax si
\item\relax mi
\item\relax vi
\item\relax ci
\end{options}
\item\label{ex:diag-11}\relax Ieri Sara \blank{} al mare.\par\nopagebreak\begin{options}
\item\relax è andata
\item\relax ha andato
\item\relax è andato
\item\relax va
\end{options}
\item\label{ex:diag-12}\relax Mentre \blank{}, è arrivato il nonno.\par\nopagebreak\begin{options}
\item\relax dormirò
\item\relax dormivo
\item\relax ho dormito
\item\relax dormo
\end{options}
\item\label{ex:diag-13}\relax (tu, negative) \blank{} la porta! (Don’t open the door!)\par\nopagebreak\begin{options}
\item\relax Non aprite
\item\relax Non aprire
\item\relax Non apra
\item\relax Non apri
\end{options}
\item\label{ex:diag-14}\relax Le chiavi? Le ho \blank{}.\par\nopagebreak\begin{options}
\item\relax persa
\item\relax perse
\item\relax persi
\item\relax perso
\end{options}
\item\label{ex:diag-15}\relax Quando sono arrivato alla stazione, il treno \blank{} già \blank{}.\par\nopagebreak\begin{options}
\item\relax sarà … partito
\item\relax era … partito
\item\relax è … partito
\item\relax partiva … —
\end{options}
\item\label{ex:diag-16}\relax Non risponde: \blank{} ancora a scuola. (He must still be at school.)\par\nopagebreak\begin{options}
\item\relax sarà
\item\relax sia
\item\relax fu
\item\relax era
\end{options}
\item\label{ex:diag-17}\relax Polite request: \blank{} un caffè, per favore.\par\nopagebreak\begin{options}
\item\relax Vorrò
\item\relax Vorrei
\item\relax Volli
\item\relax Voglio
\end{options}
\item\label{ex:diag-18}\relax In a story: \textit{Il re prese la spada.} What is the infinitive of \textit{prese}?\par\nopagebreak\begin{options}
\item\relax prestare
\item\relax prendere
\item\relax pregare
\item\relax presentare
\end{options}
\item\label{ex:diag-19}\relax Penso che Marco \blank{} ragione.\par\nopagebreak\begin{options}
\item\relax abbia
\item\relax ha
\item\relax avrebbe
\item\relax avesse
\end{options}
\item\label{ex:diag-20}\relax Se \blank{} tempo, verrei con te.\par\nopagebreak\begin{options}
\item\relax avessi
\item\relax abbia
\item\relax avevo
\item\relax avrei
\end{options}
\end{exercises}

\chapter{Mixed practice}\label{mixed}
These sets mix several lessons, so that you must first recognise which rule applies.\par

\section{Mixed practice: articles}\label{set:mixed-articles}
{\small Lessons: The definite article (\S\ref{les:definite-articles}); The indefinite article; c'è and ci sono (\S\ref{les:indefinite-articles}); Articulated prepositions (\S\ref{les:articulated-prepositions}).}\par
{\color{muted}\small Answers: page~\pageref{ans:set:mixed-articles}.}

\begin{exercises}
\item\label{ex:mix-art-1}\relax \blank{} specchio\par\nopagebreak\begin{options}
\item\relax il
\item\relax lo
\item\relax la
\item\relax l'
\end{options}
\item\label{ex:mix-art-2}\relax \blank{} estate (summer, feminine)\par\nopagebreak\begin{options}
\item\relax la
\item\relax lo
\item\relax il
\item\relax l'
\end{options}
\item\label{ex:mix-art-3}\relax \blank{} amica\par\nopagebreak\begin{options}
\item\relax una
\item\relax un
\item\relax un'
\item\relax uno
\end{options}
\item\label{ex:mix-art-4}\relax \blank{} occhi\par\nopagebreak\begin{options}
\item\relax l'
\item\relax i
\item\relax le
\item\relax gli
\end{options}
\item\label{ex:mix-art-5}\relax a + le finestre → \blank{} finestre
\item\label{ex:mix-art-6}\relax su + gli alberi → \blank{} alberi
\end{exercises}

\section{Mixed practice: the present tense}\label{set:mixed-present}
{\small Lessons: Present tense of -are verbs (\S\ref{les:present-are}); Present tense of -ere and -ire verbs (\S\ref{les:present-ere-ire}); essere and avere (\S\ref{les:essere-avere}); andare, fare, stare, dire (\S\ref{les:andare-fare-stare-dire}); venire, uscire, bere, dare and piacere (\S\ref{les:other-irregulars}).}\par
{\color{muted}\small Answers: page~\pageref{ans:set:mixed-present}.}

\begin{exercises}
\item\label{ex:mix-pres-1}\relax giocare, tu: \blank{}
\item\label{ex:mix-pres-2}\relax capire, loro: \blank{}
\item\label{ex:mix-pres-3}\relax fare, noi: \blank{}
\item\label{ex:mix-pres-4}\relax venire, tu: \blank{}
\item\label{ex:mix-pres-5}\relax Infinitive of vanno?\par\answerline
\item\label{ex:mix-pres-6}\relax Infinitive of esce?\par\answerline
\end{exercises}

\section{Mixed practice: past tenses}\label{set:mixed-past}
{\small Lessons: Passato prossimo with avere (\S\ref{les:passato-prossimo-avere}); Passato prossimo with essere; agreement (\S\ref{les:passato-prossimo-essere}); The imperfetto (\S\ref{les:imperfetto}); Imperfetto or passato prossimo? (\S\ref{les:imperfetto-passato}).}\par
{\color{muted}\small Answers: page~\pageref{ans:set:mixed-past}.}

\begin{exercises}
\item\label{ex:mix-past-1}\relax leggere, io, passato prossimo: \blank{}
\item\label{ex:mix-past-2}\relax partire, loro (masculine), passato prossimo: \blank{}
\item\label{ex:mix-past-3}\relax essere, noi, imperfetto: \blank{}
\item\label{ex:mix-past-4}\relax dire, lui/lei, imperfetto: \blank{}
\item\label{ex:mix-past-5}\relax Da bambino \blank{} (giocare) sempre in giardino.\par\nopagebreak\begin{options}
\item\relax giocavo
\item\relax giocato
\item\relax gioco
\item\relax ho giocato
\end{options}
\item\label{ex:mix-past-6}\relax Mi \blank{} (svegliarsi, female speaker) alle sette.\par\nopagebreak\begin{options}
\item\relax sono svegliato
\item\relax sono svegliata
\item\relax svegliavo
\item\relax ho svegliata
\end{options}
\end{exercises}

\section{Mixed practice: futuro and condizionale}\label{set:mixed-future-conditional}
{\small Lessons: Futuro semplice (\S\ref{les:futuro-semplice}); Futuro anteriore (\S\ref{les:futuro-anteriore}); Condizionale presente (\S\ref{les:condizionale-presente}); Condizionale passato and the future in the past (\S\ref{les:condizionale-passato}).}\par
{\color{muted}\small Answers: page~\pageref{ans:set:mixed-future-conditional}.}

\begin{exercises}
\item\label{ex:mix-fc-1}\relax Give the form.\par\smallskip{\begin{tabular}{@{}l l@{}}
venire, io, futuro & \blankcell \\
venire, io, condizionale & \blankcell \\
essere, noi, futuro & \blankcell \\
essere, noi, condizionale & \blankcell \\
fare, loro, futuro & \blankcell \\
fare, loro, condizionale & \blankcell \\
\end{tabular}}
\item\label{ex:mix-fc-2}\relax Marco ha detto che \blank{} domani. (future in the past)\par\nopagebreak\begin{options}
\item\relax venisse
\item\relax sarebbe venuto
\item\relax verrà
\item\relax verrebbe
\end{options}
\item\label{ex:mix-fc-3}\relax \textit{Secondo i giornali, il ministro si sarebbe dimesso.} What does the condizionale express?\par\nopagebreak\begin{options}
\item\relax a polite request
\item\relax a condition
\item\relax information that has not been confirmed
\item\relax a future event
\end{options}
\item\label{ex:mix-fc-4}\relax Quando \blank{} (finire, tu), usciremo.\par\nopagebreak\begin{options}
\item\relax avrai finito
\item\relax finivi
\item\relax finiresti
\item\relax avresti finito
\end{options}
\item\label{ex:mix-fc-5}\relax Translate: I would have called you, but I didn’t have your number.\par\answerline
\end{exercises}

\section{Mixed practice: indicativo, congiuntivo or condizionale?}\label{set:mixed-moods}
{\small Lessons: Congiuntivo presente: forms and main uses (\S\ref{les:congiuntivo-presente}); Congiuntivo imperfetto and trapassato (\S\ref{les:congiuntivo-imperfetto-trapassato}); Hypothetical sentences (periodo ipotetico) (\S\ref{les:periodo-ipotetico}); Sequence of tenses (concordanza dei tempi) (\S\ref{les:concordanza-dei-tempi}).}\par
{\color{muted}\small Answers: page~\pageref{ans:set:mixed-moods}.}

\begin{exercises}
\item\label{ex:mix-mood-1}\relax So che Luca \blank{} a Roma.\par\nopagebreak\begin{options}
\item\relax abitasse
\item\relax abiti
\item\relax abita
\item\relax abiterebbe
\end{options}
\item\label{ex:mix-mood-2}\relax Credo che Luca \blank{} a Roma.\par\nopagebreak\begin{options}
\item\relax abitava
\item\relax abita
\item\relax abiterebbe
\item\relax abiti
\end{options}
\item\label{ex:mix-mood-3}\relax Credevo che Luca \blank{} a Roma.\par\nopagebreak\begin{options}
\item\relax abita
\item\relax abiti
\item\relax abitasse
\item\relax abiterà
\end{options}
\item\label{ex:mix-mood-4}\relax Complete the hypotheses.\par\smallskip\begingroup\setlength{\parindent}{0pt}Se \clozeblank{1} (potere, io), ti aiuterei. Se avessi studiato, \clozeblank{2} (passare, tu) l’esame.\par\endgroup
\item\label{ex:mix-mood-5}\relax Cerco una casa che \blank{} un giardino.\par\nopagebreak\begin{options}
\item\relax avrà
\item\relax ha
\item\relax abbia
\item\relax avrebbe
\end{options}
\item\label{ex:mix-mood-6}\relax Benché \blank{} tardi, siamo usciti.\par\nopagebreak\begin{options}
\item\relax fosse
\item\relax è
\item\relax sarebbe
\item\relax era
\end{options}
\end{exercises}

\part{Grammar reference}\label{part:reference}

\chapter{The verb system at a glance}\label{ref:system}
Italian has four finite moods (indicativo, congiuntivo, condizionale, imperativo) and three non-finite forms (infinito, gerundio, participio). Every compound tense is the corresponding simple tense of \textit{avere} or \textit{essere} followed by the past participle. The level column gives an approximate teaching order, not a statement about frequency.\par

{\footnotesize\setlength{\LTpre}{2pt}
\tablecaption{The Italian verb system}
\begin{xltabular}{\linewidth}{@{}>{\raggedright\arraybackslash}p{1.6cm} >{\raggedright\arraybackslash}p{2.1cm} l l L l l@{}}
\toprule
\textbf{Mood} & \textbf{Tense} & \textbf{\textit{parlare}} (io) & \textbf{\textit{andare}} (io) & \textbf{Main use} & \textbf{Level} & \textbf{Lesson} \\
\midrule
\endhead
indicativo & \textbf{presente} & \textit{parlo} & \textit{vado} & Facts, habits and actions in progress; also the near future (domani parto) and, in narration, the "historic present". & A0 & \S\ref{les:present-are} \\
indicativo & \textbf{passato prossimo} {\color{muted}(compound)} & \textit{ho parlato} & \textit{sono andato} & Completed past events, especially with a link to the present. In the North it is the ordinary spoken past; see also passato remoto. & A1 & \S\ref{les:passato-prossimo-avere} \\
indicativo & \textbf{imperfetto} & \textit{parlavo} & \textit{andavo} & Background, description, habits and actions in progress in the past; also polite requests (volevo chiedere) and, colloquially, counterfactuals. & A2 & \S\ref{les:imperfetto} \\
indicativo & \textbf{trapassato prossimo} {\color{muted}(compound)} & \textit{avevo parlato} & \textit{ero andato} & An action completed before another past moment: quando sono arrivato, il treno era già partito. & A2 & \S\ref{les:trapassato-prossimo} \\
indicativo & \textbf{passato remoto} & \textit{parlai} & \textit{andai} & Completed past events viewed as cut off from the present: standard in written narration, history and fairy tales; in speech its use varies strongly by region. & B1 & \S\ref{les:passato-remoto} \\
indicativo & \textbf{trapassato remoto} {\color{muted}(compound)} & \textit{ebbi parlato} & \textit{fui andato} & Literary: immediate anteriority in a temporal clause (dopo che, appena, quando) whose main verb is in the passato remoto. & B2 & \S\ref{les:trapassato-remoto} \\
indicativo & \textbf{futuro semplice} & \textit{parlerò} & \textit{andrò} & Future events; also supposition about the present (sarà stanco = he is probably tired). & A2 & \S\ref{les:futuro-semplice} \\
indicativo & \textbf{futuro anteriore} {\color{muted}(compound)} & \textit{avrò parlato} & \textit{sarò andato} & An action completed before a future moment; also supposition about the past (avrà perso il treno). & B1 & \S\ref{les:futuro-anteriore} \\
congiuntivo & \textbf{congiuntivo presente} & \textit{parli} & \textit{vada} & In subordinate clauses after verbs of opinion, wish, emotion and doubt, impersonal expressions and certain conjunctions, when the main verb is present or future. & B1 & \S\ref{les:congiuntivo-presente} \\
congiuntivo & \textbf{congiuntivo passato} {\color{muted}(compound)} & \textit{abbia parlato} & \textit{sia andato} & As the congiuntivo presente, for an action prior to the main verb: penso che sia partito. & B1 & \S\ref{les:congiuntivo-passato} \\
congiuntivo & \textbf{congiuntivo imperfetto} & \textit{parlassi} & \textit{andassi} & Subjunctive contexts with a past or conditional main verb; the "se" clause of possible hypotheses; wishes (magari fosse vero). & B1 & \S\ref{les:congiuntivo-imperfetto-trapassato} \\
congiuntivo & \textbf{congiuntivo trapassato} {\color{muted}(compound)} & \textit{avessi parlato} & \textit{fossi andato} & Anteriority with a past main verb; the "se" clause of counterfactual hypotheses (se avessi saputo…). & B2 & \S\ref{les:congiuntivo-imperfetto-trapassato} \\
condizionale & \textbf{condizionale presente} & \textit{parlerei} & \textit{andrei} & Politeness, wishes, advice, hypothetical consequences, and reported unconfirmed information. & A2 & \S\ref{les:condizionale-presente} \\
condizionale & \textbf{condizionale passato} {\color{muted}(compound)} & \textit{avrei parlato} & \textit{sarei andato} & Unrealised past possibilities; and the "future in the past": disse che sarebbe venuto. & B1 & \S\ref{les:condizionale-passato} \\
imperativo & \textbf{imperativo} & \textit{parla (tu)} & \textit{va' / vai (tu)} & Commands, requests and instructions; formal Lei/Loro forms are taken from the congiuntivo presente; negative tu = non + infinitive. & A2 & \S\ref{les:imperativo} \\
\bottomrule
\end{xltabular}}

\section{Sequence of tenses (concordanza dei tempi)}\label{ref:concordanza}
In a subordinate clause that requires the congiuntivo, the tense depends on the main verb and on whether the subordinate action is simultaneous, earlier or later. Full treatment: \S\ref{les:concordanza-dei-tempi}.\par

{\small\setlength{\LTpre}{2pt}\setlength{\LTpost}{6pt}
\tablecaption{Choosing the subordinate tense}
\begin{xltabular}{\linewidth}{LllLL}
\toprule
\textbf{Main verb} & \textbf{Simultaneous} & \textbf{Earlier} & \textbf{Later} & \textbf{Example} \\
\midrule
present / future / imperative & congiuntivo presente & congiuntivo passato & congiuntivo presente or futuro & \textit{Penso che parta / sia partito / partirà.} \\
past (passato prossimo, imperfetto, passato remoto, trapassato) & congiuntivo imperfetto & congiuntivo trapassato & condizionale passato & \textit{Pensavo che partisse / fosse partito / sarebbe partito.} \\
condizionale (vorrei, preferirei …) & congiuntivo imperfetto & congiuntivo trapassato & congiuntivo imperfetto & \textit{Vorrei che tu venissi. Avrei voluto che tu fossi venuto.} \\
\bottomrule
\end{xltabular}}

\section{Hypothetical sentences (periodo ipotetico)}\label{ref:ipotetico}
Full treatment: \S\ref{les:periodo-ipotetico}.\par

{\small\setlength{\LTpre}{2pt}\setlength{\LTpost}{6pt}
\tablecaption{Types of hypothetical sentence}
\begin{xltabular}{\linewidth}{lLlL}
\toprule
\textbf{Type} & \textbf{se-clause (protasi)} & \textbf{Main clause (apodosi)} & \textbf{Example} \\
\midrule
realtà & se + indicativo & indicativo / imperativo & \textit{Se piove, resto a casa.} \\
possibilità & se + congiuntivo imperfetto & condizionale presente & \textit{Se avessi tempo, leggerei di più.} \\
irrealtà (passato) & se + congiuntivo trapassato & condizionale passato & \textit{Se fossi partito prima, sarei arrivato in tempo.} \\
irrealtà (mista) & se + congiuntivo trapassato & condizionale presente & \textit{Se avessi studiato, ora sarei più tranquillo.} \\
colloquiale (informal) & se + imperfetto & imperfetto & \textit{Se lo sapevo, venivo. (speech; not used in formal writing)} \\
\bottomrule
\end{xltabular}}
\begin{notebox}{tendcol}\notelabel{tendcol}{Register}{}The conditional is not used in the \textit{se}-clause in standard Italian (\textit{se avrei} is non-standard). The double imperfetto is common in conversation and attested in literature, but belongs to an informal register.\end{notebox}

\section{Articulated prepositions}\label{ref:articulated}

\tablecaption{Preposition + definite article}
{\small\begin{tabular}{@{}l l l l l l l l@{}}
\toprule
 & \textbf{\textit{il}} & \textbf{\textit{lo}} & \textbf{\textit{l'}} & \textbf{\textit{la}} & \textbf{\textit{i}} & \textbf{\textit{gli}} & \textbf{\textit{le}} \\
\midrule
\textbf{\textit{di}} & \textit{del} & \textit{dello} & \textit{dell'} & \textit{della} & \textit{dei} & \textit{degli} & \textit{delle} \\
\textbf{\textit{a}} & \textit{al} & \textit{allo} & \textit{all'} & \textit{alla} & \textit{ai} & \textit{agli} & \textit{alle} \\
\textbf{\textit{da}} & \textit{dal} & \textit{dallo} & \textit{dall'} & \textit{dalla} & \textit{dai} & \textit{dagli} & \textit{dalle} \\
\textbf{\textit{in}} & \textit{nel} & \textit{nello} & \textit{nell'} & \textit{nella} & \textit{nei} & \textit{negli} & \textit{nelle} \\
\textbf{\textit{su}} & \textit{sul} & \textit{sullo} & \textit{sull'} & \textit{sulla} & \textit{sui} & \textit{sugli} & \textit{sulle} \\
\bottomrule
\end{tabular}}\par\medskip

\chapter{Model conjugations}\label{ref:conjugations}
Each verb below is given in every tense and mood, as generated by the same grammar engine that drives the website. Forms printed in \irr{colour} differ from the regular pattern of the verb's family. Compound tenses of verbs with \textit{essere} are shown for a masculine subject; for a feminine subject change the participle ending (\textit{sono andata, siamo andate}). Where two forms are separated by a slash, both are standard.\par

\section{Auxiliaries}

\subsection{essere --- to be}\label{conj:essere}\index{essere@essere}
{\small Auxiliary: \textit{essere}. Completely irregular. The participle stato is shared with stare: sono stato means both 'I have been' and 'I stayed'.}\par

\moodheading{indicativo}

\par\noindent\begin{minipage}[t]{0.32\linewidth}\raggedright
\paradigmbox{indicativo presente}{%
\begin{tabular}{@{}l l@{}}
\textit{io} & \irr{sono} \\
\textit{tu} & \irr{sei} \\
\textit{lui/lei} & \irr{è} \\
\textit{noi} & \irr{siamo} \\
\textit{voi} & \irr{siete} \\
\textit{loro} & \irr{sono} \\
\end{tabular}}{\irr{Highlighted} forms differ from the regular pattern of the verb's family.}
\end{minipage}\hspace{0.02\linewidth}\begin{minipage}[t]{0.32\linewidth}\raggedright
\paradigmbox{indicativo passato prossimo (with essere, masculine subject)}{%
\begin{tabular}{@{}l l@{}}
\textit{io} & sono stato \\
\textit{tu} & sei stato \\
\textit{lui/lei} & è stato \\
\textit{noi} & siamo stati \\
\textit{voi} & siete stati \\
\textit{loro} & sono stati \\
\end{tabular}}{}
\end{minipage}\hspace{0.02\linewidth}\begin{minipage}[t]{0.32\linewidth}\raggedright
\paradigmbox{indicativo imperfetto}{%
\begin{tabular}{@{}l l@{}}
\textit{io} & \irr{ero} \\
\textit{tu} & \irr{eri} \\
\textit{lui/lei} & \irr{era} \\
\textit{noi} & \irr{eravamo} \\
\textit{voi} & \irr{eravate} \\
\textit{loro} & \irr{erano} \\
\end{tabular}}{\irr{Highlighted} forms differ from the regular pattern of the verb's family.}
\end{minipage}\par\medskip\noindent\begin{minipage}[t]{0.32\linewidth}\raggedright
\paradigmbox{indicativo trapassato prossimo (with essere, masculine subject)}{%
\begin{tabular}{@{}l l@{}}
\textit{io} & ero stato \\
\textit{tu} & eri stato \\
\textit{lui/lei} & era stato \\
\textit{noi} & eravamo stati \\
\textit{voi} & eravate stati \\
\textit{loro} & erano stati \\
\end{tabular}}{}
\end{minipage}\hspace{0.02\linewidth}\begin{minipage}[t]{0.32\linewidth}\raggedright
\paradigmbox{indicativo passato remoto}{%
\begin{tabular}{@{}l l@{}}
\textit{io} & \irr{fui} \\
\textit{tu} & \irr{fosti} \\
\textit{lui/lei} & \irr{fu} \\
\textit{noi} & \irr{fummo} \\
\textit{voi} & \irr{foste} \\
\textit{loro} & \irr{furono} \\
\end{tabular}}{\irr{Highlighted} forms differ from the regular pattern of the verb's family.}
\end{minipage}\hspace{0.02\linewidth}\begin{minipage}[t]{0.32\linewidth}\raggedright
\paradigmbox{indicativo trapassato remoto (with essere, masculine subject)}{%
\begin{tabular}{@{}l l@{}}
\textit{io} & fui stato \\
\textit{tu} & fosti stato \\
\textit{lui/lei} & fu stato \\
\textit{noi} & fummo stati \\
\textit{voi} & foste stati \\
\textit{loro} & furono stati \\
\end{tabular}}{}
\end{minipage}\par\medskip\noindent\begin{minipage}[t]{0.32\linewidth}\raggedright
\paradigmbox{indicativo futuro semplice}{%
\begin{tabular}{@{}l l@{}}
\textit{io} & \irr{sarò} \\
\textit{tu} & \irr{sarai} \\
\textit{lui/lei} & \irr{sarà} \\
\textit{noi} & \irr{saremo} \\
\textit{voi} & \irr{sarete} \\
\textit{loro} & \irr{saranno} \\
\end{tabular}}{\irr{Highlighted} forms differ from the regular pattern of the verb's family.}
\end{minipage}\hspace{0.02\linewidth}\begin{minipage}[t]{0.32\linewidth}\raggedright
\paradigmbox{indicativo futuro anteriore (with essere, masculine subject)}{%
\begin{tabular}{@{}l l@{}}
\textit{io} & sarò stato \\
\textit{tu} & sarai stato \\
\textit{lui/lei} & sarà stato \\
\textit{noi} & saremo stati \\
\textit{voi} & sarete stati \\
\textit{loro} & saranno stati \\
\end{tabular}}{}
\end{minipage}\par\medskip

\moodheading{congiuntivo}

\par\noindent\begin{minipage}[t]{0.32\linewidth}\raggedright
\paradigmbox{congiuntivo presente}{%
\begin{tabular}{@{}l l@{}}
\textit{io} & \irr{sia} \\
\textit{tu} & \irr{sia} \\
\textit{lui/lei} & \irr{sia} \\
\textit{noi} & \irr{siamo} \\
\textit{voi} & \irr{siate} \\
\textit{loro} & \irr{siano} \\
\end{tabular}}{\irr{Highlighted} forms differ from the regular pattern of the verb's family.}
\end{minipage}\hspace{0.02\linewidth}\begin{minipage}[t]{0.32\linewidth}\raggedright
\paradigmbox{congiuntivo passato (with essere, masculine subject)}{%
\begin{tabular}{@{}l l@{}}
\textit{io} & sia stato \\
\textit{tu} & sia stato \\
\textit{lui/lei} & sia stato \\
\textit{noi} & siamo stati \\
\textit{voi} & siate stati \\
\textit{loro} & siano stati \\
\end{tabular}}{}
\end{minipage}\hspace{0.02\linewidth}\begin{minipage}[t]{0.32\linewidth}\raggedright
\paradigmbox{congiuntivo imperfetto}{%
\begin{tabular}{@{}l l@{}}
\textit{io} & \irr{fossi} \\
\textit{tu} & \irr{fossi} \\
\textit{lui/lei} & \irr{fosse} \\
\textit{noi} & \irr{fossimo} \\
\textit{voi} & \irr{foste} \\
\textit{loro} & \irr{fossero} \\
\end{tabular}}{\irr{Highlighted} forms differ from the regular pattern of the verb's family.}
\end{minipage}\par\medskip\noindent\begin{minipage}[t]{0.32\linewidth}\raggedright
\paradigmbox{congiuntivo trapassato (with essere, masculine subject)}{%
\begin{tabular}{@{}l l@{}}
\textit{io} & fossi stato \\
\textit{tu} & fossi stato \\
\textit{lui/lei} & fosse stato \\
\textit{noi} & fossimo stati \\
\textit{voi} & foste stati \\
\textit{loro} & fossero stati \\
\end{tabular}}{}
\end{minipage}\par\medskip

\moodheading{condizionale}

\par\noindent\begin{minipage}[t]{0.32\linewidth}\raggedright
\paradigmbox{condizionale presente}{%
\begin{tabular}{@{}l l@{}}
\textit{io} & \irr{sarei} \\
\textit{tu} & \irr{saresti} \\
\textit{lui/lei} & \irr{sarebbe} \\
\textit{noi} & \irr{saremmo} \\
\textit{voi} & \irr{sareste} \\
\textit{loro} & \irr{sarebbero} \\
\end{tabular}}{\irr{Highlighted} forms differ from the regular pattern of the verb's family.}
\end{minipage}\hspace{0.02\linewidth}\begin{minipage}[t]{0.32\linewidth}\raggedright
\paradigmbox{condizionale passato (with essere, masculine subject)}{%
\begin{tabular}{@{}l l@{}}
\textit{io} & sarei stato \\
\textit{tu} & saresti stato \\
\textit{lui/lei} & sarebbe stato \\
\textit{noi} & saremmo stati \\
\textit{voi} & sareste stati \\
\textit{loro} & sarebbero stati \\
\end{tabular}}{}
\end{minipage}\par\medskip

\moodheading{imperativo}

\par\noindent\begin{minipage}[t]{0.32\linewidth}\raggedright
\paradigmbox{imperativo}{%
\begin{tabular}{@{}l l@{}}
\textit{tu} & \irr{sii} \\
\textit{Lei} & \irr{sia} \\
\textit{noi} & \irr{siamo} \\
\textit{voi} & \irr{siate} \\
\textit{Loro} & \irr{siano} \\
\end{tabular}}{\irr{Highlighted} forms differ from the regular pattern of the verb's family. Negative \textit{tu}: \textit{non essere}. \textit{Lei} and \textit{Loro} are the congiuntivo forms.}
\end{minipage}\par\medskip

\moodheading{forme non finite}
{\small\begin{tabular}{@{}l l@{}}
infinito presente & \textit{essere} \\
infinito passato & \textit{essere stato} \\
gerundio presente & \textit{essendo} \\
gerundio passato & \textit{essendo stato} \\
participio passato & \textit{stato} \\
\end{tabular}}\par
\clearpage

\subsection{avere --- to have}\label{conj:avere}\index{avere@avere}
{\small Auxiliary: \textit{avere}. The h in ho, hai, ha, hanno is silent; it distinguishes ho (I have) from o (or), ha (has) from a (to), hanno (they have) from anno (year). Noi and voi are regular.}\par

\moodheading{indicativo}

\par\noindent\begin{minipage}[t]{0.32\linewidth}\raggedright
\paradigmbox{indicativo presente}{%
\begin{tabular}{@{}l l@{}}
\textit{io} & \irr{ho} \\
\textit{tu} & \irr{hai} \\
\textit{lui/lei} & \irr{ha} \\
\textit{noi} & \irr{abbiamo} \\
\textit{voi} & avete \\
\textit{loro} & \irr{hanno} \\
\end{tabular}}{\irr{Highlighted} forms differ from the regular pattern of the verb's family.}
\end{minipage}\hspace{0.02\linewidth}\begin{minipage}[t]{0.32\linewidth}\raggedright
\paradigmbox{indicativo passato prossimo (with avere)}{%
\begin{tabular}{@{}l l@{}}
\textit{io} & ho avuto \\
\textit{tu} & hai avuto \\
\textit{lui/lei} & ha avuto \\
\textit{noi} & abbiamo avuto \\
\textit{voi} & avete avuto \\
\textit{loro} & hanno avuto \\
\end{tabular}}{}
\end{minipage}\hspace{0.02\linewidth}\begin{minipage}[t]{0.32\linewidth}\raggedright
\paradigmbox{indicativo imperfetto}{%
\begin{tabular}{@{}l l@{}}
\textit{io} & avevo \\
\textit{tu} & avevi \\
\textit{lui/lei} & aveva \\
\textit{noi} & avevamo \\
\textit{voi} & avevate \\
\textit{loro} & avevano \\
\end{tabular}}{}
\end{minipage}\par\medskip\noindent\begin{minipage}[t]{0.32\linewidth}\raggedright
\paradigmbox{indicativo trapassato prossimo (with avere)}{%
\begin{tabular}{@{}l l@{}}
\textit{io} & avevo avuto \\
\textit{tu} & avevi avuto \\
\textit{lui/lei} & aveva avuto \\
\textit{noi} & avevamo avuto \\
\textit{voi} & avevate avuto \\
\textit{loro} & avevano avuto \\
\end{tabular}}{}
\end{minipage}\hspace{0.02\linewidth}\begin{minipage}[t]{0.32\linewidth}\raggedright
\paradigmbox{indicativo passato remoto}{%
\begin{tabular}{@{}l l@{}}
\textit{io} & \irr{ebbi} \\
\textit{tu} & avesti \\
\textit{lui/lei} & \irr{ebbe} \\
\textit{noi} & avemmo \\
\textit{voi} & aveste \\
\textit{loro} & \irr{ebbero} \\
\end{tabular}}{\irr{Highlighted} forms differ from the regular pattern of the verb's family.}
\end{minipage}\hspace{0.02\linewidth}\begin{minipage}[t]{0.32\linewidth}\raggedright
\paradigmbox{indicativo trapassato remoto (with avere)}{%
\begin{tabular}{@{}l l@{}}
\textit{io} & ebbi avuto \\
\textit{tu} & avesti avuto \\
\textit{lui/lei} & ebbe avuto \\
\textit{noi} & avemmo avuto \\
\textit{voi} & aveste avuto \\
\textit{loro} & ebbero avuto \\
\end{tabular}}{}
\end{minipage}\par\medskip\noindent\begin{minipage}[t]{0.32\linewidth}\raggedright
\paradigmbox{indicativo futuro semplice}{%
\begin{tabular}{@{}l l@{}}
\textit{io} & \irr{avrò} \\
\textit{tu} & \irr{avrai} \\
\textit{lui/lei} & \irr{avrà} \\
\textit{noi} & \irr{avremo} \\
\textit{voi} & \irr{avrete} \\
\textit{loro} & \irr{avranno} \\
\end{tabular}}{\irr{Highlighted} forms differ from the regular pattern of the verb's family.}
\end{minipage}\hspace{0.02\linewidth}\begin{minipage}[t]{0.32\linewidth}\raggedright
\paradigmbox{indicativo futuro anteriore (with avere)}{%
\begin{tabular}{@{}l l@{}}
\textit{io} & avrò avuto \\
\textit{tu} & avrai avuto \\
\textit{lui/lei} & avrà avuto \\
\textit{noi} & avremo avuto \\
\textit{voi} & avrete avuto \\
\textit{loro} & avranno avuto \\
\end{tabular}}{}
\end{minipage}\par\medskip

\moodheading{congiuntivo}

\par\noindent\begin{minipage}[t]{0.32\linewidth}\raggedright
\paradigmbox{congiuntivo presente}{%
\begin{tabular}{@{}l l@{}}
\textit{io} & \irr{abbia} \\
\textit{tu} & \irr{abbia} \\
\textit{lui/lei} & \irr{abbia} \\
\textit{noi} & \irr{abbiamo} \\
\textit{voi} & \irr{abbiate} \\
\textit{loro} & \irr{abbiano} \\
\end{tabular}}{\irr{Highlighted} forms differ from the regular pattern of the verb's family.}
\end{minipage}\hspace{0.02\linewidth}\begin{minipage}[t]{0.32\linewidth}\raggedright
\paradigmbox{congiuntivo passato (with avere)}{%
\begin{tabular}{@{}l l@{}}
\textit{io} & abbia avuto \\
\textit{tu} & abbia avuto \\
\textit{lui/lei} & abbia avuto \\
\textit{noi} & abbiamo avuto \\
\textit{voi} & abbiate avuto \\
\textit{loro} & abbiano avuto \\
\end{tabular}}{}
\end{minipage}\hspace{0.02\linewidth}\begin{minipage}[t]{0.32\linewidth}\raggedright
\paradigmbox{congiuntivo imperfetto}{%
\begin{tabular}{@{}l l@{}}
\textit{io} & avessi \\
\textit{tu} & avessi \\
\textit{lui/lei} & avesse \\
\textit{noi} & avessimo \\
\textit{voi} & aveste \\
\textit{loro} & avessero \\
\end{tabular}}{}
\end{minipage}\par\medskip\noindent\begin{minipage}[t]{0.32\linewidth}\raggedright
\paradigmbox{congiuntivo trapassato (with avere)}{%
\begin{tabular}{@{}l l@{}}
\textit{io} & avessi avuto \\
\textit{tu} & avessi avuto \\
\textit{lui/lei} & avesse avuto \\
\textit{noi} & avessimo avuto \\
\textit{voi} & aveste avuto \\
\textit{loro} & avessero avuto \\
\end{tabular}}{}
\end{minipage}\par\medskip

\moodheading{condizionale}

\par\noindent\begin{minipage}[t]{0.32\linewidth}\raggedright
\paradigmbox{condizionale presente}{%
\begin{tabular}{@{}l l@{}}
\textit{io} & \irr{avrei} \\
\textit{tu} & \irr{avresti} \\
\textit{lui/lei} & \irr{avrebbe} \\
\textit{noi} & \irr{avremmo} \\
\textit{voi} & \irr{avreste} \\
\textit{loro} & \irr{avrebbero} \\
\end{tabular}}{\irr{Highlighted} forms differ from the regular pattern of the verb's family.}
\end{minipage}\hspace{0.02\linewidth}\begin{minipage}[t]{0.32\linewidth}\raggedright
\paradigmbox{condizionale passato (with avere)}{%
\begin{tabular}{@{}l l@{}}
\textit{io} & avrei avuto \\
\textit{tu} & avresti avuto \\
\textit{lui/lei} & avrebbe avuto \\
\textit{noi} & avremmo avuto \\
\textit{voi} & avreste avuto \\
\textit{loro} & avrebbero avuto \\
\end{tabular}}{}
\end{minipage}\par\medskip

\moodheading{imperativo}

\par\noindent\begin{minipage}[t]{0.32\linewidth}\raggedright
\paradigmbox{imperativo}{%
\begin{tabular}{@{}l l@{}}
\textit{tu} & \irr{abbi} \\
\textit{Lei} & \irr{abbia} \\
\textit{noi} & \irr{abbiamo} \\
\textit{voi} & \irr{abbiate} \\
\textit{Loro} & \irr{abbiano} \\
\end{tabular}}{\irr{Highlighted} forms differ from the regular pattern of the verb's family. Negative \textit{tu}: \textit{non avere}. \textit{Lei} and \textit{Loro} are the congiuntivo forms.}
\end{minipage}\par\medskip

\moodheading{forme non finite}
{\small\begin{tabular}{@{}l l@{}}
infinito presente & \textit{avere} \\
infinito passato & \textit{avere avuto} \\
gerundio presente & \textit{avendo} \\
gerundio passato & \textit{avendo avuto} \\
participio passato & \textit{avuto} \\
\end{tabular}}\par
\clearpage

\section{Regular models}

\subsection{parlare --- to speak; to talk}\label{conj:parlare}\index{parlare@parlare}
{\small Auxiliary: \textit{avere}. }\par

\moodheading{indicativo}

\par\noindent\begin{minipage}[t]{0.32\linewidth}\raggedright
\paradigmbox{indicativo presente}{%
\begin{tabular}{@{}l l@{}}
\textit{io} & parlo \\
\textit{tu} & parli \\
\textit{lui/lei} & parla \\
\textit{noi} & parliamo \\
\textit{voi} & parlate \\
\textit{loro} & parlano \\
\end{tabular}}{}
\end{minipage}\hspace{0.02\linewidth}\begin{minipage}[t]{0.32\linewidth}\raggedright
\paradigmbox{indicativo passato prossimo (with avere)}{%
\begin{tabular}{@{}l l@{}}
\textit{io} & ho parlato \\
\textit{tu} & hai parlato \\
\textit{lui/lei} & ha parlato \\
\textit{noi} & abbiamo parlato \\
\textit{voi} & avete parlato \\
\textit{loro} & hanno parlato \\
\end{tabular}}{}
\end{minipage}\hspace{0.02\linewidth}\begin{minipage}[t]{0.32\linewidth}\raggedright
\paradigmbox{indicativo imperfetto}{%
\begin{tabular}{@{}l l@{}}
\textit{io} & parlavo \\
\textit{tu} & parlavi \\
\textit{lui/lei} & parlava \\
\textit{noi} & parlavamo \\
\textit{voi} & parlavate \\
\textit{loro} & parlavano \\
\end{tabular}}{}
\end{minipage}\par\medskip\noindent\begin{minipage}[t]{0.32\linewidth}\raggedright
\paradigmbox{indicativo trapassato prossimo (with avere)}{%
\begin{tabular}{@{}l l@{}}
\textit{io} & avevo parlato \\
\textit{tu} & avevi parlato \\
\textit{lui/lei} & aveva parlato \\
\textit{noi} & avevamo parlato \\
\textit{voi} & avevate parlato \\
\textit{loro} & avevano parlato \\
\end{tabular}}{}
\end{minipage}\hspace{0.02\linewidth}\begin{minipage}[t]{0.32\linewidth}\raggedright
\paradigmbox{indicativo passato remoto}{%
\begin{tabular}{@{}l l@{}}
\textit{io} & parlai \\
\textit{tu} & parlasti \\
\textit{lui/lei} & parlò \\
\textit{noi} & parlammo \\
\textit{voi} & parlaste \\
\textit{loro} & parlarono \\
\end{tabular}}{}
\end{minipage}\hspace{0.02\linewidth}\begin{minipage}[t]{0.32\linewidth}\raggedright
\paradigmbox{indicativo trapassato remoto (with avere)}{%
\begin{tabular}{@{}l l@{}}
\textit{io} & ebbi parlato \\
\textit{tu} & avesti parlato \\
\textit{lui/lei} & ebbe parlato \\
\textit{noi} & avemmo parlato \\
\textit{voi} & aveste parlato \\
\textit{loro} & ebbero parlato \\
\end{tabular}}{}
\end{minipage}\par\medskip\noindent\begin{minipage}[t]{0.32\linewidth}\raggedright
\paradigmbox{indicativo futuro semplice}{%
\begin{tabular}{@{}l l@{}}
\textit{io} & parlerò \\
\textit{tu} & parlerai \\
\textit{lui/lei} & parlerà \\
\textit{noi} & parleremo \\
\textit{voi} & parlerete \\
\textit{loro} & parleranno \\
\end{tabular}}{}
\end{minipage}\hspace{0.02\linewidth}\begin{minipage}[t]{0.32\linewidth}\raggedright
\paradigmbox{indicativo futuro anteriore (with avere)}{%
\begin{tabular}{@{}l l@{}}
\textit{io} & avrò parlato \\
\textit{tu} & avrai parlato \\
\textit{lui/lei} & avrà parlato \\
\textit{noi} & avremo parlato \\
\textit{voi} & avrete parlato \\
\textit{loro} & avranno parlato \\
\end{tabular}}{}
\end{minipage}\par\medskip

\moodheading{congiuntivo}

\par\noindent\begin{minipage}[t]{0.32\linewidth}\raggedright
\paradigmbox{congiuntivo presente}{%
\begin{tabular}{@{}l l@{}}
\textit{io} & parli \\
\textit{tu} & parli \\
\textit{lui/lei} & parli \\
\textit{noi} & parliamo \\
\textit{voi} & parliate \\
\textit{loro} & parlino \\
\end{tabular}}{}
\end{minipage}\hspace{0.02\linewidth}\begin{minipage}[t]{0.32\linewidth}\raggedright
\paradigmbox{congiuntivo passato (with avere)}{%
\begin{tabular}{@{}l l@{}}
\textit{io} & abbia parlato \\
\textit{tu} & abbia parlato \\
\textit{lui/lei} & abbia parlato \\
\textit{noi} & abbiamo parlato \\
\textit{voi} & abbiate parlato \\
\textit{loro} & abbiano parlato \\
\end{tabular}}{}
\end{minipage}\hspace{0.02\linewidth}\begin{minipage}[t]{0.32\linewidth}\raggedright
\paradigmbox{congiuntivo imperfetto}{%
\begin{tabular}{@{}l l@{}}
\textit{io} & parlassi \\
\textit{tu} & parlassi \\
\textit{lui/lei} & parlasse \\
\textit{noi} & parlassimo \\
\textit{voi} & parlaste \\
\textit{loro} & parlassero \\
\end{tabular}}{}
\end{minipage}\par\medskip\noindent\begin{minipage}[t]{0.32\linewidth}\raggedright
\paradigmbox{congiuntivo trapassato (with avere)}{%
\begin{tabular}{@{}l l@{}}
\textit{io} & avessi parlato \\
\textit{tu} & avessi parlato \\
\textit{lui/lei} & avesse parlato \\
\textit{noi} & avessimo parlato \\
\textit{voi} & aveste parlato \\
\textit{loro} & avessero parlato \\
\end{tabular}}{}
\end{minipage}\par\medskip

\moodheading{condizionale}

\par\noindent\begin{minipage}[t]{0.32\linewidth}\raggedright
\paradigmbox{condizionale presente}{%
\begin{tabular}{@{}l l@{}}
\textit{io} & parlerei \\
\textit{tu} & parleresti \\
\textit{lui/lei} & parlerebbe \\
\textit{noi} & parleremmo \\
\textit{voi} & parlereste \\
\textit{loro} & parlerebbero \\
\end{tabular}}{}
\end{minipage}\hspace{0.02\linewidth}\begin{minipage}[t]{0.32\linewidth}\raggedright
\paradigmbox{condizionale passato (with avere)}{%
\begin{tabular}{@{}l l@{}}
\textit{io} & avrei parlato \\
\textit{tu} & avresti parlato \\
\textit{lui/lei} & avrebbe parlato \\
\textit{noi} & avremmo parlato \\
\textit{voi} & avreste parlato \\
\textit{loro} & avrebbero parlato \\
\end{tabular}}{}
\end{minipage}\par\medskip

\moodheading{imperativo}

\par\noindent\begin{minipage}[t]{0.32\linewidth}\raggedright
\paradigmbox{imperativo}{%
\begin{tabular}{@{}l l@{}}
\textit{tu} & parla \\
\textit{Lei} & parli \\
\textit{noi} & parliamo \\
\textit{voi} & parlate \\
\textit{Loro} & parlino \\
\end{tabular}}{Negative \textit{tu}: \textit{non parlare}. \textit{Lei} and \textit{Loro} are the congiuntivo forms.}
\end{minipage}\par\medskip

\moodheading{forme non finite}
{\small\begin{tabular}{@{}l l@{}}
infinito presente & \textit{parlare} \\
infinito passato & \textit{avere parlato} \\
gerundio presente & \textit{parlando} \\
gerundio passato & \textit{avendo parlato} \\
participio passato & \textit{parlato} \\
\end{tabular}}\par
\clearpage

\subsection{credere --- to believe}\label{conj:credere}\index{credere@credere}
{\small Auxiliary: \textit{avere}. }\par

\moodheading{indicativo}

\par\noindent\begin{minipage}[t]{0.32\linewidth}\raggedright
\paradigmbox{indicativo presente}{%
\begin{tabular}{@{}l l@{}}
\textit{io} & credo \\
\textit{tu} & credi \\
\textit{lui/lei} & crede \\
\textit{noi} & crediamo \\
\textit{voi} & credete \\
\textit{loro} & credono \\
\end{tabular}}{}
\end{minipage}\hspace{0.02\linewidth}\begin{minipage}[t]{0.32\linewidth}\raggedright
\paradigmbox{indicativo passato prossimo (with avere)}{%
\begin{tabular}{@{}l l@{}}
\textit{io} & ho creduto \\
\textit{tu} & hai creduto \\
\textit{lui/lei} & ha creduto \\
\textit{noi} & abbiamo creduto \\
\textit{voi} & avete creduto \\
\textit{loro} & hanno creduto \\
\end{tabular}}{}
\end{minipage}\hspace{0.02\linewidth}\begin{minipage}[t]{0.32\linewidth}\raggedright
\paradigmbox{indicativo imperfetto}{%
\begin{tabular}{@{}l l@{}}
\textit{io} & credevo \\
\textit{tu} & credevi \\
\textit{lui/lei} & credeva \\
\textit{noi} & credevamo \\
\textit{voi} & credevate \\
\textit{loro} & credevano \\
\end{tabular}}{}
\end{minipage}\par\medskip\noindent\begin{minipage}[t]{0.32\linewidth}\raggedright
\paradigmbox{indicativo trapassato prossimo (with avere)}{%
\begin{tabular}{@{}l l@{}}
\textit{io} & avevo creduto \\
\textit{tu} & avevi creduto \\
\textit{lui/lei} & aveva creduto \\
\textit{noi} & avevamo creduto \\
\textit{voi} & avevate creduto \\
\textit{loro} & avevano creduto \\
\end{tabular}}{}
\end{minipage}\hspace{0.02\linewidth}\begin{minipage}[t]{0.32\linewidth}\raggedright
\paradigmbox{indicativo passato remoto}{%
\begin{tabular}{@{}l l@{}}
\textit{io} & credei / credetti \\
\textit{tu} & credesti \\
\textit{lui/lei} & credé / credette \\
\textit{noi} & credemmo \\
\textit{voi} & credeste \\
\textit{loro} & crederono / credettero \\
\end{tabular}}{Where two forms are given, both are standard.}
\end{minipage}\hspace{0.02\linewidth}\begin{minipage}[t]{0.32\linewidth}\raggedright
\paradigmbox{indicativo trapassato remoto (with avere)}{%
\begin{tabular}{@{}l l@{}}
\textit{io} & ebbi creduto \\
\textit{tu} & avesti creduto \\
\textit{lui/lei} & ebbe creduto \\
\textit{noi} & avemmo creduto \\
\textit{voi} & aveste creduto \\
\textit{loro} & ebbero creduto \\
\end{tabular}}{}
\end{minipage}\par\medskip\noindent\begin{minipage}[t]{0.32\linewidth}\raggedright
\paradigmbox{indicativo futuro semplice}{%
\begin{tabular}{@{}l l@{}}
\textit{io} & crederò \\
\textit{tu} & crederai \\
\textit{lui/lei} & crederà \\
\textit{noi} & crederemo \\
\textit{voi} & crederete \\
\textit{loro} & crederanno \\
\end{tabular}}{}
\end{minipage}\hspace{0.02\linewidth}\begin{minipage}[t]{0.32\linewidth}\raggedright
\paradigmbox{indicativo futuro anteriore (with avere)}{%
\begin{tabular}{@{}l l@{}}
\textit{io} & avrò creduto \\
\textit{tu} & avrai creduto \\
\textit{lui/lei} & avrà creduto \\
\textit{noi} & avremo creduto \\
\textit{voi} & avrete creduto \\
\textit{loro} & avranno creduto \\
\end{tabular}}{}
\end{minipage}\par\medskip

\moodheading{congiuntivo}

\par\noindent\begin{minipage}[t]{0.32\linewidth}\raggedright
\paradigmbox{congiuntivo presente}{%
\begin{tabular}{@{}l l@{}}
\textit{io} & creda \\
\textit{tu} & creda \\
\textit{lui/lei} & creda \\
\textit{noi} & crediamo \\
\textit{voi} & crediate \\
\textit{loro} & credano \\
\end{tabular}}{}
\end{minipage}\hspace{0.02\linewidth}\begin{minipage}[t]{0.32\linewidth}\raggedright
\paradigmbox{congiuntivo passato (with avere)}{%
\begin{tabular}{@{}l l@{}}
\textit{io} & abbia creduto \\
\textit{tu} & abbia creduto \\
\textit{lui/lei} & abbia creduto \\
\textit{noi} & abbiamo creduto \\
\textit{voi} & abbiate creduto \\
\textit{loro} & abbiano creduto \\
\end{tabular}}{}
\end{minipage}\hspace{0.02\linewidth}\begin{minipage}[t]{0.32\linewidth}\raggedright
\paradigmbox{congiuntivo imperfetto}{%
\begin{tabular}{@{}l l@{}}
\textit{io} & credessi \\
\textit{tu} & credessi \\
\textit{lui/lei} & credesse \\
\textit{noi} & credessimo \\
\textit{voi} & credeste \\
\textit{loro} & credessero \\
\end{tabular}}{}
\end{minipage}\par\medskip\noindent\begin{minipage}[t]{0.32\linewidth}\raggedright
\paradigmbox{congiuntivo trapassato (with avere)}{%
\begin{tabular}{@{}l l@{}}
\textit{io} & avessi creduto \\
\textit{tu} & avessi creduto \\
\textit{lui/lei} & avesse creduto \\
\textit{noi} & avessimo creduto \\
\textit{voi} & aveste creduto \\
\textit{loro} & avessero creduto \\
\end{tabular}}{}
\end{minipage}\par\medskip

\moodheading{condizionale}

\par\noindent\begin{minipage}[t]{0.32\linewidth}\raggedright
\paradigmbox{condizionale presente}{%
\begin{tabular}{@{}l l@{}}
\textit{io} & crederei \\
\textit{tu} & crederesti \\
\textit{lui/lei} & crederebbe \\
\textit{noi} & crederemmo \\
\textit{voi} & credereste \\
\textit{loro} & crederebbero \\
\end{tabular}}{}
\end{minipage}\hspace{0.02\linewidth}\begin{minipage}[t]{0.32\linewidth}\raggedright
\paradigmbox{condizionale passato (with avere)}{%
\begin{tabular}{@{}l l@{}}
\textit{io} & avrei creduto \\
\textit{tu} & avresti creduto \\
\textit{lui/lei} & avrebbe creduto \\
\textit{noi} & avremmo creduto \\
\textit{voi} & avreste creduto \\
\textit{loro} & avrebbero creduto \\
\end{tabular}}{}
\end{minipage}\par\medskip

\moodheading{imperativo}

\par\noindent\begin{minipage}[t]{0.32\linewidth}\raggedright
\paradigmbox{imperativo}{%
\begin{tabular}{@{}l l@{}}
\textit{tu} & credi \\
\textit{Lei} & creda \\
\textit{noi} & crediamo \\
\textit{voi} & credete \\
\textit{Loro} & credano \\
\end{tabular}}{Negative \textit{tu}: \textit{non credere}. \textit{Lei} and \textit{Loro} are the congiuntivo forms.}
\end{minipage}\par\medskip

\moodheading{forme non finite}
{\small\begin{tabular}{@{}l l@{}}
infinito presente & \textit{credere} \\
infinito passato & \textit{avere creduto} \\
gerundio presente & \textit{credendo} \\
gerundio passato & \textit{avendo creduto} \\
participio passato & \textit{creduto} \\
\end{tabular}}\par
\clearpage

\subsection{dormire --- to sleep}\label{conj:dormire}\index{dormire@dormire}
{\small Auxiliary: \textit{avere}. }\par

\moodheading{indicativo}

\par\noindent\begin{minipage}[t]{0.32\linewidth}\raggedright
\paradigmbox{indicativo presente}{%
\begin{tabular}{@{}l l@{}}
\textit{io} & dormo \\
\textit{tu} & dormi \\
\textit{lui/lei} & dorme \\
\textit{noi} & dormiamo \\
\textit{voi} & dormite \\
\textit{loro} & dormono \\
\end{tabular}}{}
\end{minipage}\hspace{0.02\linewidth}\begin{minipage}[t]{0.32\linewidth}\raggedright
\paradigmbox{indicativo passato prossimo (with avere)}{%
\begin{tabular}{@{}l l@{}}
\textit{io} & ho dormito \\
\textit{tu} & hai dormito \\
\textit{lui/lei} & ha dormito \\
\textit{noi} & abbiamo dormito \\
\textit{voi} & avete dormito \\
\textit{loro} & hanno dormito \\
\end{tabular}}{}
\end{minipage}\hspace{0.02\linewidth}\begin{minipage}[t]{0.32\linewidth}\raggedright
\paradigmbox{indicativo imperfetto}{%
\begin{tabular}{@{}l l@{}}
\textit{io} & dormivo \\
\textit{tu} & dormivi \\
\textit{lui/lei} & dormiva \\
\textit{noi} & dormivamo \\
\textit{voi} & dormivate \\
\textit{loro} & dormivano \\
\end{tabular}}{}
\end{minipage}\par\medskip\noindent\begin{minipage}[t]{0.32\linewidth}\raggedright
\paradigmbox{indicativo trapassato prossimo (with avere)}{%
\begin{tabular}{@{}l l@{}}
\textit{io} & avevo dormito \\
\textit{tu} & avevi dormito \\
\textit{lui/lei} & aveva dormito \\
\textit{noi} & avevamo dormito \\
\textit{voi} & avevate dormito \\
\textit{loro} & avevano dormito \\
\end{tabular}}{}
\end{minipage}\hspace{0.02\linewidth}\begin{minipage}[t]{0.32\linewidth}\raggedright
\paradigmbox{indicativo passato remoto}{%
\begin{tabular}{@{}l l@{}}
\textit{io} & dormii \\
\textit{tu} & dormisti \\
\textit{lui/lei} & dormì \\
\textit{noi} & dormimmo \\
\textit{voi} & dormiste \\
\textit{loro} & dormirono \\
\end{tabular}}{}
\end{minipage}\hspace{0.02\linewidth}\begin{minipage}[t]{0.32\linewidth}\raggedright
\paradigmbox{indicativo trapassato remoto (with avere)}{%
\begin{tabular}{@{}l l@{}}
\textit{io} & ebbi dormito \\
\textit{tu} & avesti dormito \\
\textit{lui/lei} & ebbe dormito \\
\textit{noi} & avemmo dormito \\
\textit{voi} & aveste dormito \\
\textit{loro} & ebbero dormito \\
\end{tabular}}{}
\end{minipage}\par\medskip\noindent\begin{minipage}[t]{0.32\linewidth}\raggedright
\paradigmbox{indicativo futuro semplice}{%
\begin{tabular}{@{}l l@{}}
\textit{io} & dormirò \\
\textit{tu} & dormirai \\
\textit{lui/lei} & dormirà \\
\textit{noi} & dormiremo \\
\textit{voi} & dormirete \\
\textit{loro} & dormiranno \\
\end{tabular}}{}
\end{minipage}\hspace{0.02\linewidth}\begin{minipage}[t]{0.32\linewidth}\raggedright
\paradigmbox{indicativo futuro anteriore (with avere)}{%
\begin{tabular}{@{}l l@{}}
\textit{io} & avrò dormito \\
\textit{tu} & avrai dormito \\
\textit{lui/lei} & avrà dormito \\
\textit{noi} & avremo dormito \\
\textit{voi} & avrete dormito \\
\textit{loro} & avranno dormito \\
\end{tabular}}{}
\end{minipage}\par\medskip

\moodheading{congiuntivo}

\par\noindent\begin{minipage}[t]{0.32\linewidth}\raggedright
\paradigmbox{congiuntivo presente}{%
\begin{tabular}{@{}l l@{}}
\textit{io} & dorma \\
\textit{tu} & dorma \\
\textit{lui/lei} & dorma \\
\textit{noi} & dormiamo \\
\textit{voi} & dormiate \\
\textit{loro} & dormano \\
\end{tabular}}{}
\end{minipage}\hspace{0.02\linewidth}\begin{minipage}[t]{0.32\linewidth}\raggedright
\paradigmbox{congiuntivo passato (with avere)}{%
\begin{tabular}{@{}l l@{}}
\textit{io} & abbia dormito \\
\textit{tu} & abbia dormito \\
\textit{lui/lei} & abbia dormito \\
\textit{noi} & abbiamo dormito \\
\textit{voi} & abbiate dormito \\
\textit{loro} & abbiano dormito \\
\end{tabular}}{}
\end{minipage}\hspace{0.02\linewidth}\begin{minipage}[t]{0.32\linewidth}\raggedright
\paradigmbox{congiuntivo imperfetto}{%
\begin{tabular}{@{}l l@{}}
\textit{io} & dormissi \\
\textit{tu} & dormissi \\
\textit{lui/lei} & dormisse \\
\textit{noi} & dormissimo \\
\textit{voi} & dormiste \\
\textit{loro} & dormissero \\
\end{tabular}}{}
\end{minipage}\par\medskip\noindent\begin{minipage}[t]{0.32\linewidth}\raggedright
\paradigmbox{congiuntivo trapassato (with avere)}{%
\begin{tabular}{@{}l l@{}}
\textit{io} & avessi dormito \\
\textit{tu} & avessi dormito \\
\textit{lui/lei} & avesse dormito \\
\textit{noi} & avessimo dormito \\
\textit{voi} & aveste dormito \\
\textit{loro} & avessero dormito \\
\end{tabular}}{}
\end{minipage}\par\medskip

\moodheading{condizionale}

\par\noindent\begin{minipage}[t]{0.32\linewidth}\raggedright
\paradigmbox{condizionale presente}{%
\begin{tabular}{@{}l l@{}}
\textit{io} & dormirei \\
\textit{tu} & dormiresti \\
\textit{lui/lei} & dormirebbe \\
\textit{noi} & dormiremmo \\
\textit{voi} & dormireste \\
\textit{loro} & dormirebbero \\
\end{tabular}}{}
\end{minipage}\hspace{0.02\linewidth}\begin{minipage}[t]{0.32\linewidth}\raggedright
\paradigmbox{condizionale passato (with avere)}{%
\begin{tabular}{@{}l l@{}}
\textit{io} & avrei dormito \\
\textit{tu} & avresti dormito \\
\textit{lui/lei} & avrebbe dormito \\
\textit{noi} & avremmo dormito \\
\textit{voi} & avreste dormito \\
\textit{loro} & avrebbero dormito \\
\end{tabular}}{}
\end{minipage}\par\medskip

\moodheading{imperativo}

\par\noindent\begin{minipage}[t]{0.32\linewidth}\raggedright
\paradigmbox{imperativo}{%
\begin{tabular}{@{}l l@{}}
\textit{tu} & dormi \\
\textit{Lei} & dorma \\
\textit{noi} & dormiamo \\
\textit{voi} & dormite \\
\textit{Loro} & dormano \\
\end{tabular}}{Negative \textit{tu}: \textit{non dormire}. \textit{Lei} and \textit{Loro} are the congiuntivo forms.}
\end{minipage}\par\medskip

\moodheading{forme non finite}
{\small\begin{tabular}{@{}l l@{}}
infinito presente & \textit{dormire} \\
infinito passato & \textit{avere dormito} \\
gerundio presente & \textit{dormendo} \\
gerundio passato & \textit{avendo dormito} \\
participio passato & \textit{dormito} \\
\end{tabular}}\par
\clearpage

\subsection{finire --- to finish; to end}\label{conj:finire}\index{finire@finire}
{\small Auxiliary: \textit{avere}. }\par

\moodheading{indicativo}

\par\noindent\begin{minipage}[t]{0.32\linewidth}\raggedright
\paradigmbox{indicativo presente}{%
\begin{tabular}{@{}l l@{}}
\textit{io} & finisco \\
\textit{tu} & finisci \\
\textit{lui/lei} & finisce \\
\textit{noi} & finiamo \\
\textit{voi} & finite \\
\textit{loro} & finiscono \\
\end{tabular}}{}
\end{minipage}\hspace{0.02\linewidth}\begin{minipage}[t]{0.32\linewidth}\raggedright
\paradigmbox{indicativo passato prossimo (with avere)}{%
\begin{tabular}{@{}l l@{}}
\textit{io} & ho finito \\
\textit{tu} & hai finito \\
\textit{lui/lei} & ha finito \\
\textit{noi} & abbiamo finito \\
\textit{voi} & avete finito \\
\textit{loro} & hanno finito \\
\end{tabular}}{}
\end{minipage}\hspace{0.02\linewidth}\begin{minipage}[t]{0.32\linewidth}\raggedright
\paradigmbox{indicativo imperfetto}{%
\begin{tabular}{@{}l l@{}}
\textit{io} & finivo \\
\textit{tu} & finivi \\
\textit{lui/lei} & finiva \\
\textit{noi} & finivamo \\
\textit{voi} & finivate \\
\textit{loro} & finivano \\
\end{tabular}}{}
\end{minipage}\par\medskip\noindent\begin{minipage}[t]{0.32\linewidth}\raggedright
\paradigmbox{indicativo trapassato prossimo (with avere)}{%
\begin{tabular}{@{}l l@{}}
\textit{io} & avevo finito \\
\textit{tu} & avevi finito \\
\textit{lui/lei} & aveva finito \\
\textit{noi} & avevamo finito \\
\textit{voi} & avevate finito \\
\textit{loro} & avevano finito \\
\end{tabular}}{}
\end{minipage}\hspace{0.02\linewidth}\begin{minipage}[t]{0.32\linewidth}\raggedright
\paradigmbox{indicativo passato remoto}{%
\begin{tabular}{@{}l l@{}}
\textit{io} & finii \\
\textit{tu} & finisti \\
\textit{lui/lei} & finì \\
\textit{noi} & finimmo \\
\textit{voi} & finiste \\
\textit{loro} & finirono \\
\end{tabular}}{}
\end{minipage}\hspace{0.02\linewidth}\begin{minipage}[t]{0.32\linewidth}\raggedright
\paradigmbox{indicativo trapassato remoto (with avere)}{%
\begin{tabular}{@{}l l@{}}
\textit{io} & ebbi finito \\
\textit{tu} & avesti finito \\
\textit{lui/lei} & ebbe finito \\
\textit{noi} & avemmo finito \\
\textit{voi} & aveste finito \\
\textit{loro} & ebbero finito \\
\end{tabular}}{}
\end{minipage}\par\medskip\noindent\begin{minipage}[t]{0.32\linewidth}\raggedright
\paradigmbox{indicativo futuro semplice}{%
\begin{tabular}{@{}l l@{}}
\textit{io} & finirò \\
\textit{tu} & finirai \\
\textit{lui/lei} & finirà \\
\textit{noi} & finiremo \\
\textit{voi} & finirete \\
\textit{loro} & finiranno \\
\end{tabular}}{}
\end{minipage}\hspace{0.02\linewidth}\begin{minipage}[t]{0.32\linewidth}\raggedright
\paradigmbox{indicativo futuro anteriore (with avere)}{%
\begin{tabular}{@{}l l@{}}
\textit{io} & avrò finito \\
\textit{tu} & avrai finito \\
\textit{lui/lei} & avrà finito \\
\textit{noi} & avremo finito \\
\textit{voi} & avrete finito \\
\textit{loro} & avranno finito \\
\end{tabular}}{}
\end{minipage}\par\medskip

\moodheading{congiuntivo}

\par\noindent\begin{minipage}[t]{0.32\linewidth}\raggedright
\paradigmbox{congiuntivo presente}{%
\begin{tabular}{@{}l l@{}}
\textit{io} & finisca \\
\textit{tu} & finisca \\
\textit{lui/lei} & finisca \\
\textit{noi} & finiamo \\
\textit{voi} & finiate \\
\textit{loro} & finiscano \\
\end{tabular}}{}
\end{minipage}\hspace{0.02\linewidth}\begin{minipage}[t]{0.32\linewidth}\raggedright
\paradigmbox{congiuntivo passato (with avere)}{%
\begin{tabular}{@{}l l@{}}
\textit{io} & abbia finito \\
\textit{tu} & abbia finito \\
\textit{lui/lei} & abbia finito \\
\textit{noi} & abbiamo finito \\
\textit{voi} & abbiate finito \\
\textit{loro} & abbiano finito \\
\end{tabular}}{}
\end{minipage}\hspace{0.02\linewidth}\begin{minipage}[t]{0.32\linewidth}\raggedright
\paradigmbox{congiuntivo imperfetto}{%
\begin{tabular}{@{}l l@{}}
\textit{io} & finissi \\
\textit{tu} & finissi \\
\textit{lui/lei} & finisse \\
\textit{noi} & finissimo \\
\textit{voi} & finiste \\
\textit{loro} & finissero \\
\end{tabular}}{}
\end{minipage}\par\medskip\noindent\begin{minipage}[t]{0.32\linewidth}\raggedright
\paradigmbox{congiuntivo trapassato (with avere)}{%
\begin{tabular}{@{}l l@{}}
\textit{io} & avessi finito \\
\textit{tu} & avessi finito \\
\textit{lui/lei} & avesse finito \\
\textit{noi} & avessimo finito \\
\textit{voi} & aveste finito \\
\textit{loro} & avessero finito \\
\end{tabular}}{}
\end{minipage}\par\medskip

\moodheading{condizionale}

\par\noindent\begin{minipage}[t]{0.32\linewidth}\raggedright
\paradigmbox{condizionale presente}{%
\begin{tabular}{@{}l l@{}}
\textit{io} & finirei \\
\textit{tu} & finiresti \\
\textit{lui/lei} & finirebbe \\
\textit{noi} & finiremmo \\
\textit{voi} & finireste \\
\textit{loro} & finirebbero \\
\end{tabular}}{}
\end{minipage}\hspace{0.02\linewidth}\begin{minipage}[t]{0.32\linewidth}\raggedright
\paradigmbox{condizionale passato (with avere)}{%
\begin{tabular}{@{}l l@{}}
\textit{io} & avrei finito \\
\textit{tu} & avresti finito \\
\textit{lui/lei} & avrebbe finito \\
\textit{noi} & avremmo finito \\
\textit{voi} & avreste finito \\
\textit{loro} & avrebbero finito \\
\end{tabular}}{}
\end{minipage}\par\medskip

\moodheading{imperativo}

\par\noindent\begin{minipage}[t]{0.32\linewidth}\raggedright
\paradigmbox{imperativo}{%
\begin{tabular}{@{}l l@{}}
\textit{tu} & finisci \\
\textit{Lei} & finisca \\
\textit{noi} & finiamo \\
\textit{voi} & finite \\
\textit{Loro} & finiscano \\
\end{tabular}}{Negative \textit{tu}: \textit{non finire}. \textit{Lei} and \textit{Loro} are the congiuntivo forms.}
\end{minipage}\par\medskip

\moodheading{forme non finite}
{\small\begin{tabular}{@{}l l@{}}
infinito presente & \textit{finire} \\
infinito passato & \textit{avere finito} \\
gerundio presente & \textit{finendo} \\
gerundio passato & \textit{avendo finito} \\
participio passato & \textit{finito} \\
\end{tabular}}\par
\clearpage

\subsection{lavarsi --- to wash (oneself)}\label{conj:lavarsi}\index{lavarsi@lavarsi}
{\small Auxiliary: \textit{essere}. }\par

\moodheading{indicativo}

\par\noindent\begin{minipage}[t]{0.32\linewidth}\raggedright
\paradigmbox{indicativo presente}{%
\begin{tabular}{@{}l l@{}}
\textit{io} & mi lavo \\
\textit{tu} & ti lavi \\
\textit{lui/lei} & si lava \\
\textit{noi} & ci laviamo \\
\textit{voi} & vi lavate \\
\textit{loro} & si lavano \\
\end{tabular}}{}
\end{minipage}\hspace{0.02\linewidth}\begin{minipage}[t]{0.32\linewidth}\raggedright
\paradigmbox{indicativo passato prossimo (with essere, masculine subject)}{%
\begin{tabular}{@{}l l@{}}
\textit{io} & mi sono lavato \\
\textit{tu} & ti sei lavato \\
\textit{lui/lei} & si è lavato \\
\textit{noi} & ci siamo lavati \\
\textit{voi} & vi siete lavati \\
\textit{loro} & si sono lavati \\
\end{tabular}}{}
\end{minipage}\hspace{0.02\linewidth}\begin{minipage}[t]{0.32\linewidth}\raggedright
\paradigmbox{indicativo imperfetto}{%
\begin{tabular}{@{}l l@{}}
\textit{io} & mi lavavo \\
\textit{tu} & ti lavavi \\
\textit{lui/lei} & si lavava \\
\textit{noi} & ci lavavamo \\
\textit{voi} & vi lavavate \\
\textit{loro} & si lavavano \\
\end{tabular}}{}
\end{minipage}\par\medskip\noindent\begin{minipage}[t]{0.32\linewidth}\raggedright
\paradigmbox{indicativo trapassato prossimo (with essere, masculine subject)}{%
\begin{tabular}{@{}l l@{}}
\textit{io} & mi ero lavato \\
\textit{tu} & ti eri lavato \\
\textit{lui/lei} & si era lavato \\
\textit{noi} & ci eravamo lavati \\
\textit{voi} & vi eravate lavati \\
\textit{loro} & si erano lavati \\
\end{tabular}}{}
\end{minipage}\hspace{0.02\linewidth}\begin{minipage}[t]{0.32\linewidth}\raggedright
\paradigmbox{indicativo passato remoto}{%
\begin{tabular}{@{}l l@{}}
\textit{io} & mi lavai \\
\textit{tu} & ti lavasti \\
\textit{lui/lei} & si lavò \\
\textit{noi} & ci lavammo \\
\textit{voi} & vi lavaste \\
\textit{loro} & si lavarono \\
\end{tabular}}{}
\end{minipage}\hspace{0.02\linewidth}\begin{minipage}[t]{0.32\linewidth}\raggedright
\paradigmbox{indicativo trapassato remoto (with essere, masculine subject)}{%
\begin{tabular}{@{}l l@{}}
\textit{io} & mi fui lavato \\
\textit{tu} & ti fosti lavato \\
\textit{lui/lei} & si fu lavato \\
\textit{noi} & ci fummo lavati \\
\textit{voi} & vi foste lavati \\
\textit{loro} & si furono lavati \\
\end{tabular}}{}
\end{minipage}\par\medskip\noindent\begin{minipage}[t]{0.32\linewidth}\raggedright
\paradigmbox{indicativo futuro semplice}{%
\begin{tabular}{@{}l l@{}}
\textit{io} & mi laverò \\
\textit{tu} & ti laverai \\
\textit{lui/lei} & si laverà \\
\textit{noi} & ci laveremo \\
\textit{voi} & vi laverete \\
\textit{loro} & si laveranno \\
\end{tabular}}{}
\end{minipage}\hspace{0.02\linewidth}\begin{minipage}[t]{0.32\linewidth}\raggedright
\paradigmbox{indicativo futuro anteriore (with essere, masculine subject)}{%
\begin{tabular}{@{}l l@{}}
\textit{io} & mi sarò lavato \\
\textit{tu} & ti sarai lavato \\
\textit{lui/lei} & si sarà lavato \\
\textit{noi} & ci saremo lavati \\
\textit{voi} & vi sarete lavati \\
\textit{loro} & si saranno lavati \\
\end{tabular}}{}
\end{minipage}\par\medskip

\moodheading{congiuntivo}

\par\noindent\begin{minipage}[t]{0.32\linewidth}\raggedright
\paradigmbox{congiuntivo presente}{%
\begin{tabular}{@{}l l@{}}
\textit{io} & mi lavi \\
\textit{tu} & ti lavi \\
\textit{lui/lei} & si lavi \\
\textit{noi} & ci laviamo \\
\textit{voi} & vi laviate \\
\textit{loro} & si lavino \\
\end{tabular}}{}
\end{minipage}\hspace{0.02\linewidth}\begin{minipage}[t]{0.32\linewidth}\raggedright
\paradigmbox{congiuntivo passato (with essere, masculine subject)}{%
\begin{tabular}{@{}l l@{}}
\textit{io} & mi sia lavato \\
\textit{tu} & ti sia lavato \\
\textit{lui/lei} & si sia lavato \\
\textit{noi} & ci siamo lavati \\
\textit{voi} & vi siate lavati \\
\textit{loro} & si siano lavati \\
\end{tabular}}{}
\end{minipage}\hspace{0.02\linewidth}\begin{minipage}[t]{0.32\linewidth}\raggedright
\paradigmbox{congiuntivo imperfetto}{%
\begin{tabular}{@{}l l@{}}
\textit{io} & mi lavassi \\
\textit{tu} & ti lavassi \\
\textit{lui/lei} & si lavasse \\
\textit{noi} & ci lavassimo \\
\textit{voi} & vi lavaste \\
\textit{loro} & si lavassero \\
\end{tabular}}{}
\end{minipage}\par\medskip\noindent\begin{minipage}[t]{0.32\linewidth}\raggedright
\paradigmbox{congiuntivo trapassato (with essere, masculine subject)}{%
\begin{tabular}{@{}l l@{}}
\textit{io} & mi fossi lavato \\
\textit{tu} & ti fossi lavato \\
\textit{lui/lei} & si fosse lavato \\
\textit{noi} & ci fossimo lavati \\
\textit{voi} & vi foste lavati \\
\textit{loro} & si fossero lavati \\
\end{tabular}}{}
\end{minipage}\par\medskip

\moodheading{condizionale}

\par\noindent\begin{minipage}[t]{0.32\linewidth}\raggedright
\paradigmbox{condizionale presente}{%
\begin{tabular}{@{}l l@{}}
\textit{io} & mi laverei \\
\textit{tu} & ti laveresti \\
\textit{lui/lei} & si laverebbe \\
\textit{noi} & ci laveremmo \\
\textit{voi} & vi lavereste \\
\textit{loro} & si laverebbero \\
\end{tabular}}{}
\end{minipage}\hspace{0.02\linewidth}\begin{minipage}[t]{0.32\linewidth}\raggedright
\paradigmbox{condizionale passato (with essere, masculine subject)}{%
\begin{tabular}{@{}l l@{}}
\textit{io} & mi sarei lavato \\
\textit{tu} & ti saresti lavato \\
\textit{lui/lei} & si sarebbe lavato \\
\textit{noi} & ci saremmo lavati \\
\textit{voi} & vi sareste lavati \\
\textit{loro} & si sarebbero lavati \\
\end{tabular}}{}
\end{minipage}\par\medskip

\moodheading{imperativo}

\par\noindent\begin{minipage}[t]{0.32\linewidth}\raggedright
\paradigmbox{imperativo}{%
\begin{tabular}{@{}l l@{}}
\textit{tu} & lavati \\
\textit{Lei} & si lavi \\
\textit{noi} & laviamoci \\
\textit{voi} & lavatevi \\
\textit{Loro} & si lavino \\
\end{tabular}}{Negative \textit{tu}: \textit{non lavarti / non ti lavare}. \textit{Lei} and \textit{Loro} are the congiuntivo forms.}
\end{minipage}\par\medskip

\moodheading{forme non finite}
{\small\begin{tabular}{@{}l l@{}}
infinito presente & \textit{lavarsi} \\
infinito passato & \textit{essersi lavato} \\
gerundio presente & \textit{lavandosi} \\
gerundio passato & \textit{essendosi lavato} \\
participio passato & \textit{lavato} \\
\end{tabular}}\par
\clearpage

\section{Irregular verbs}

\subsection{andare --- to go}\label{conj:andare}\index{andare@andare}
{\small Auxiliary: \textit{essere}. Two stems: vad-/va- in the stressed forms (io, tu, lui/lei, loro), regular and- in noi and voi.}\par

\moodheading{indicativo}

\par\noindent\begin{minipage}[t]{0.32\linewidth}\raggedright
\paradigmbox{indicativo presente}{%
\begin{tabular}{@{}l l@{}}
\textit{io} & \irr{vado} \\
\textit{tu} & \irr{vai} \\
\textit{lui/lei} & \irr{va} \\
\textit{noi} & andiamo \\
\textit{voi} & andate \\
\textit{loro} & \irr{vanno} \\
\end{tabular}}{\irr{Highlighted} forms differ from the regular pattern of the verb's family.}
\end{minipage}\hspace{0.02\linewidth}\begin{minipage}[t]{0.32\linewidth}\raggedright
\paradigmbox{indicativo passato prossimo (with essere, masculine subject)}{%
\begin{tabular}{@{}l l@{}}
\textit{io} & sono andato \\
\textit{tu} & sei andato \\
\textit{lui/lei} & è andato \\
\textit{noi} & siamo andati \\
\textit{voi} & siete andati \\
\textit{loro} & sono andati \\
\end{tabular}}{}
\end{minipage}\hspace{0.02\linewidth}\begin{minipage}[t]{0.32\linewidth}\raggedright
\paradigmbox{indicativo imperfetto}{%
\begin{tabular}{@{}l l@{}}
\textit{io} & andavo \\
\textit{tu} & andavi \\
\textit{lui/lei} & andava \\
\textit{noi} & andavamo \\
\textit{voi} & andavate \\
\textit{loro} & andavano \\
\end{tabular}}{}
\end{minipage}\par\medskip\noindent\begin{minipage}[t]{0.32\linewidth}\raggedright
\paradigmbox{indicativo trapassato prossimo (with essere, masculine subject)}{%
\begin{tabular}{@{}l l@{}}
\textit{io} & ero andato \\
\textit{tu} & eri andato \\
\textit{lui/lei} & era andato \\
\textit{noi} & eravamo andati \\
\textit{voi} & eravate andati \\
\textit{loro} & erano andati \\
\end{tabular}}{}
\end{minipage}\hspace{0.02\linewidth}\begin{minipage}[t]{0.32\linewidth}\raggedright
\paradigmbox{indicativo passato remoto}{%
\begin{tabular}{@{}l l@{}}
\textit{io} & andai \\
\textit{tu} & andasti \\
\textit{lui/lei} & andò \\
\textit{noi} & andammo \\
\textit{voi} & andaste \\
\textit{loro} & andarono \\
\end{tabular}}{}
\end{minipage}\hspace{0.02\linewidth}\begin{minipage}[t]{0.32\linewidth}\raggedright
\paradigmbox{indicativo trapassato remoto (with essere, masculine subject)}{%
\begin{tabular}{@{}l l@{}}
\textit{io} & fui andato \\
\textit{tu} & fosti andato \\
\textit{lui/lei} & fu andato \\
\textit{noi} & fummo andati \\
\textit{voi} & foste andati \\
\textit{loro} & furono andati \\
\end{tabular}}{}
\end{minipage}\par\medskip\noindent\begin{minipage}[t]{0.32\linewidth}\raggedright
\paradigmbox{indicativo futuro semplice}{%
\begin{tabular}{@{}l l@{}}
\textit{io} & \irr{andrò} \\
\textit{tu} & \irr{andrai} \\
\textit{lui/lei} & \irr{andrà} \\
\textit{noi} & \irr{andremo} \\
\textit{voi} & \irr{andrete} \\
\textit{loro} & \irr{andranno} \\
\end{tabular}}{\irr{Highlighted} forms differ from the regular pattern of the verb's family.}
\end{minipage}\hspace{0.02\linewidth}\begin{minipage}[t]{0.32\linewidth}\raggedright
\paradigmbox{indicativo futuro anteriore (with essere, masculine subject)}{%
\begin{tabular}{@{}l l@{}}
\textit{io} & sarò andato \\
\textit{tu} & sarai andato \\
\textit{lui/lei} & sarà andato \\
\textit{noi} & saremo andati \\
\textit{voi} & sarete andati \\
\textit{loro} & saranno andati \\
\end{tabular}}{}
\end{minipage}\par\medskip

\moodheading{congiuntivo}

\par\noindent\begin{minipage}[t]{0.32\linewidth}\raggedright
\paradigmbox{congiuntivo presente}{%
\begin{tabular}{@{}l l@{}}
\textit{io} & \irr{vada} \\
\textit{tu} & \irr{vada} \\
\textit{lui/lei} & \irr{vada} \\
\textit{noi} & andiamo \\
\textit{voi} & andiate \\
\textit{loro} & \irr{vadano} \\
\end{tabular}}{\irr{Highlighted} forms differ from the regular pattern of the verb's family.}
\end{minipage}\hspace{0.02\linewidth}\begin{minipage}[t]{0.32\linewidth}\raggedright
\paradigmbox{congiuntivo passato (with essere, masculine subject)}{%
\begin{tabular}{@{}l l@{}}
\textit{io} & sia andato \\
\textit{tu} & sia andato \\
\textit{lui/lei} & sia andato \\
\textit{noi} & siamo andati \\
\textit{voi} & siate andati \\
\textit{loro} & siano andati \\
\end{tabular}}{}
\end{minipage}\hspace{0.02\linewidth}\begin{minipage}[t]{0.32\linewidth}\raggedright
\paradigmbox{congiuntivo imperfetto}{%
\begin{tabular}{@{}l l@{}}
\textit{io} & andassi \\
\textit{tu} & andassi \\
\textit{lui/lei} & andasse \\
\textit{noi} & andassimo \\
\textit{voi} & andaste \\
\textit{loro} & andassero \\
\end{tabular}}{}
\end{minipage}\par\medskip\noindent\begin{minipage}[t]{0.32\linewidth}\raggedright
\paradigmbox{congiuntivo trapassato (with essere, masculine subject)}{%
\begin{tabular}{@{}l l@{}}
\textit{io} & fossi andato \\
\textit{tu} & fossi andato \\
\textit{lui/lei} & fosse andato \\
\textit{noi} & fossimo andati \\
\textit{voi} & foste andati \\
\textit{loro} & fossero andati \\
\end{tabular}}{}
\end{minipage}\par\medskip

\moodheading{condizionale}

\par\noindent\begin{minipage}[t]{0.32\linewidth}\raggedright
\paradigmbox{condizionale presente}{%
\begin{tabular}{@{}l l@{}}
\textit{io} & \irr{andrei} \\
\textit{tu} & \irr{andresti} \\
\textit{lui/lei} & \irr{andrebbe} \\
\textit{noi} & \irr{andremmo} \\
\textit{voi} & \irr{andreste} \\
\textit{loro} & \irr{andrebbero} \\
\end{tabular}}{\irr{Highlighted} forms differ from the regular pattern of the verb's family.}
\end{minipage}\hspace{0.02\linewidth}\begin{minipage}[t]{0.32\linewidth}\raggedright
\paradigmbox{condizionale passato (with essere, masculine subject)}{%
\begin{tabular}{@{}l l@{}}
\textit{io} & sarei andato \\
\textit{tu} & saresti andato \\
\textit{lui/lei} & sarebbe andato \\
\textit{noi} & saremmo andati \\
\textit{voi} & sareste andati \\
\textit{loro} & sarebbero andati \\
\end{tabular}}{}
\end{minipage}\par\medskip

\moodheading{imperativo}

\par\noindent\begin{minipage}[t]{0.32\linewidth}\raggedright
\paradigmbox{imperativo}{%
\begin{tabular}{@{}l l@{}}
\textit{tu} & \irr{va' / vai} \\
\textit{Lei} & \irr{vada} \\
\textit{noi} & andiamo \\
\textit{voi} & andate \\
\textit{Loro} & \irr{vadano} \\
\end{tabular}}{\irr{Highlighted} forms differ from the regular pattern of the verb's family. Where two forms are given, both are standard. Negative \textit{tu}: \textit{non andare}. \textit{Lei} and \textit{Loro} are the congiuntivo forms.}
\end{minipage}\par\medskip

\moodheading{forme non finite}
{\small\begin{tabular}{@{}l l@{}}
infinito presente & \textit{andare} \\
infinito passato & \textit{essere andato} \\
gerundio presente & \textit{andando} \\
gerundio passato & \textit{essendo andato} \\
participio passato & \textit{andato} \\
\end{tabular}}\par
\clearpage

\subsection{fare --- to do; to make}\label{conj:fare}\index{fare@fare}
{\small Auxiliary: \textit{avere}. Historically from Latin facere: the stem fac- reappears in faccio, facciamo and the whole imperfetto (facevo).}\par

\moodheading{indicativo}

\par\noindent\begin{minipage}[t]{0.32\linewidth}\raggedright
\paradigmbox{indicativo presente}{%
\begin{tabular}{@{}l l@{}}
\textit{io} & \irr{faccio} \\
\textit{tu} & \irr{fai} \\
\textit{lui/lei} & fa \\
\textit{noi} & \irr{facciamo} \\
\textit{voi} & fate \\
\textit{loro} & \irr{fanno} \\
\end{tabular}}{\irr{Highlighted} forms differ from the regular pattern of the verb's family.}
\end{minipage}\hspace{0.02\linewidth}\begin{minipage}[t]{0.32\linewidth}\raggedright
\paradigmbox{indicativo passato prossimo (with avere)}{%
\begin{tabular}{@{}l l@{}}
\textit{io} & ho fatto \\
\textit{tu} & hai fatto \\
\textit{lui/lei} & ha fatto \\
\textit{noi} & abbiamo fatto \\
\textit{voi} & avete fatto \\
\textit{loro} & hanno fatto \\
\end{tabular}}{}
\end{minipage}\hspace{0.02\linewidth}\begin{minipage}[t]{0.32\linewidth}\raggedright
\paradigmbox{indicativo imperfetto}{%
\begin{tabular}{@{}l l@{}}
\textit{io} & \irr{facevo} \\
\textit{tu} & \irr{facevi} \\
\textit{lui/lei} & \irr{faceva} \\
\textit{noi} & \irr{facevamo} \\
\textit{voi} & \irr{facevate} \\
\textit{loro} & \irr{facevano} \\
\end{tabular}}{\irr{Highlighted} forms differ from the regular pattern of the verb's family.}
\end{minipage}\par\medskip\noindent\begin{minipage}[t]{0.32\linewidth}\raggedright
\paradigmbox{indicativo trapassato prossimo (with avere)}{%
\begin{tabular}{@{}l l@{}}
\textit{io} & avevo fatto \\
\textit{tu} & avevi fatto \\
\textit{lui/lei} & aveva fatto \\
\textit{noi} & avevamo fatto \\
\textit{voi} & avevate fatto \\
\textit{loro} & avevano fatto \\
\end{tabular}}{}
\end{minipage}\hspace{0.02\linewidth}\begin{minipage}[t]{0.32\linewidth}\raggedright
\paradigmbox{indicativo passato remoto}{%
\begin{tabular}{@{}l l@{}}
\textit{io} & \irr{feci} \\
\textit{tu} & \irr{facesti} \\
\textit{lui/lei} & \irr{fece} \\
\textit{noi} & \irr{facemmo} \\
\textit{voi} & \irr{faceste} \\
\textit{loro} & \irr{fecero} \\
\end{tabular}}{\irr{Highlighted} forms differ from the regular pattern of the verb's family.}
\end{minipage}\hspace{0.02\linewidth}\begin{minipage}[t]{0.32\linewidth}\raggedright
\paradigmbox{indicativo trapassato remoto (with avere)}{%
\begin{tabular}{@{}l l@{}}
\textit{io} & ebbi fatto \\
\textit{tu} & avesti fatto \\
\textit{lui/lei} & ebbe fatto \\
\textit{noi} & avemmo fatto \\
\textit{voi} & aveste fatto \\
\textit{loro} & ebbero fatto \\
\end{tabular}}{}
\end{minipage}\par\medskip\noindent\begin{minipage}[t]{0.32\linewidth}\raggedright
\paradigmbox{indicativo futuro semplice}{%
\begin{tabular}{@{}l l@{}}
\textit{io} & \irr{farò} \\
\textit{tu} & \irr{farai} \\
\textit{lui/lei} & \irr{farà} \\
\textit{noi} & \irr{faremo} \\
\textit{voi} & \irr{farete} \\
\textit{loro} & \irr{faranno} \\
\end{tabular}}{\irr{Highlighted} forms differ from the regular pattern of the verb's family.}
\end{minipage}\hspace{0.02\linewidth}\begin{minipage}[t]{0.32\linewidth}\raggedright
\paradigmbox{indicativo futuro anteriore (with avere)}{%
\begin{tabular}{@{}l l@{}}
\textit{io} & avrò fatto \\
\textit{tu} & avrai fatto \\
\textit{lui/lei} & avrà fatto \\
\textit{noi} & avremo fatto \\
\textit{voi} & avrete fatto \\
\textit{loro} & avranno fatto \\
\end{tabular}}{}
\end{minipage}\par\medskip

\moodheading{congiuntivo}

\par\noindent\begin{minipage}[t]{0.32\linewidth}\raggedright
\paradigmbox{congiuntivo presente}{%
\begin{tabular}{@{}l l@{}}
\textit{io} & \irr{faccia} \\
\textit{tu} & \irr{faccia} \\
\textit{lui/lei} & \irr{faccia} \\
\textit{noi} & \irr{facciamo} \\
\textit{voi} & \irr{facciate} \\
\textit{loro} & \irr{facciano} \\
\end{tabular}}{\irr{Highlighted} forms differ from the regular pattern of the verb's family.}
\end{minipage}\hspace{0.02\linewidth}\begin{minipage}[t]{0.32\linewidth}\raggedright
\paradigmbox{congiuntivo passato (with avere)}{%
\begin{tabular}{@{}l l@{}}
\textit{io} & abbia fatto \\
\textit{tu} & abbia fatto \\
\textit{lui/lei} & abbia fatto \\
\textit{noi} & abbiamo fatto \\
\textit{voi} & abbiate fatto \\
\textit{loro} & abbiano fatto \\
\end{tabular}}{}
\end{minipage}\hspace{0.02\linewidth}\begin{minipage}[t]{0.32\linewidth}\raggedright
\paradigmbox{congiuntivo imperfetto}{%
\begin{tabular}{@{}l l@{}}
\textit{io} & \irr{facessi} \\
\textit{tu} & \irr{facessi} \\
\textit{lui/lei} & \irr{facesse} \\
\textit{noi} & \irr{facessimo} \\
\textit{voi} & \irr{faceste} \\
\textit{loro} & \irr{facessero} \\
\end{tabular}}{\irr{Highlighted} forms differ from the regular pattern of the verb's family.}
\end{minipage}\par\medskip\noindent\begin{minipage}[t]{0.32\linewidth}\raggedright
\paradigmbox{congiuntivo trapassato (with avere)}{%
\begin{tabular}{@{}l l@{}}
\textit{io} & avessi fatto \\
\textit{tu} & avessi fatto \\
\textit{lui/lei} & avesse fatto \\
\textit{noi} & avessimo fatto \\
\textit{voi} & aveste fatto \\
\textit{loro} & avessero fatto \\
\end{tabular}}{}
\end{minipage}\par\medskip

\moodheading{condizionale}

\par\noindent\begin{minipage}[t]{0.32\linewidth}\raggedright
\paradigmbox{condizionale presente}{%
\begin{tabular}{@{}l l@{}}
\textit{io} & \irr{farei} \\
\textit{tu} & \irr{faresti} \\
\textit{lui/lei} & \irr{farebbe} \\
\textit{noi} & \irr{faremmo} \\
\textit{voi} & \irr{fareste} \\
\textit{loro} & \irr{farebbero} \\
\end{tabular}}{\irr{Highlighted} forms differ from the regular pattern of the verb's family.}
\end{minipage}\hspace{0.02\linewidth}\begin{minipage}[t]{0.32\linewidth}\raggedright
\paradigmbox{condizionale passato (with avere)}{%
\begin{tabular}{@{}l l@{}}
\textit{io} & avrei fatto \\
\textit{tu} & avresti fatto \\
\textit{lui/lei} & avrebbe fatto \\
\textit{noi} & avremmo fatto \\
\textit{voi} & avreste fatto \\
\textit{loro} & avrebbero fatto \\
\end{tabular}}{}
\end{minipage}\par\medskip

\moodheading{imperativo}

\par\noindent\begin{minipage}[t]{0.32\linewidth}\raggedright
\paradigmbox{imperativo}{%
\begin{tabular}{@{}l l@{}}
\textit{tu} & \irr{fa' / fai} \\
\textit{Lei} & \irr{faccia} \\
\textit{noi} & \irr{facciamo} \\
\textit{voi} & fate \\
\textit{Loro} & \irr{facciano} \\
\end{tabular}}{\irr{Highlighted} forms differ from the regular pattern of the verb's family. Where two forms are given, both are standard. Negative \textit{tu}: \textit{non fare}. \textit{Lei} and \textit{Loro} are the congiuntivo forms.}
\end{minipage}\par\medskip

\moodheading{forme non finite}
{\small\begin{tabular}{@{}l l@{}}
infinito presente & \textit{fare} \\
infinito passato & \textit{avere fatto} \\
gerundio presente & \textit{facendo} \\
gerundio passato & \textit{avendo fatto} \\
participio passato & \textit{fatto} \\
\end{tabular}}\par
\clearpage

\subsection{stare --- to stay; to be (of health)}\label{conj:stare}\index{stare@stare}
{\small Auxiliary: \textit{essere}. Short forms like those of dare and fare. Used for health (sto bene), staying, and in stare + gerund.}\par

\moodheading{indicativo}

\par\noindent\begin{minipage}[t]{0.32\linewidth}\raggedright
\paradigmbox{indicativo presente}{%
\begin{tabular}{@{}l l@{}}
\textit{io} & sto \\
\textit{tu} & \irr{stai} \\
\textit{lui/lei} & sta \\
\textit{noi} & stiamo \\
\textit{voi} & state \\
\textit{loro} & \irr{stanno} \\
\end{tabular}}{\irr{Highlighted} forms differ from the regular pattern of the verb's family.}
\end{minipage}\hspace{0.02\linewidth}\begin{minipage}[t]{0.32\linewidth}\raggedright
\paradigmbox{indicativo passato prossimo (with essere, masculine subject)}{%
\begin{tabular}{@{}l l@{}}
\textit{io} & sono stato \\
\textit{tu} & sei stato \\
\textit{lui/lei} & è stato \\
\textit{noi} & siamo stati \\
\textit{voi} & siete stati \\
\textit{loro} & sono stati \\
\end{tabular}}{}
\end{minipage}\hspace{0.02\linewidth}\begin{minipage}[t]{0.32\linewidth}\raggedright
\paradigmbox{indicativo imperfetto}{%
\begin{tabular}{@{}l l@{}}
\textit{io} & stavo \\
\textit{tu} & stavi \\
\textit{lui/lei} & stava \\
\textit{noi} & stavamo \\
\textit{voi} & stavate \\
\textit{loro} & stavano \\
\end{tabular}}{}
\end{minipage}\par\medskip\noindent\begin{minipage}[t]{0.32\linewidth}\raggedright
\paradigmbox{indicativo trapassato prossimo (with essere, masculine subject)}{%
\begin{tabular}{@{}l l@{}}
\textit{io} & ero stato \\
\textit{tu} & eri stato \\
\textit{lui/lei} & era stato \\
\textit{noi} & eravamo stati \\
\textit{voi} & eravate stati \\
\textit{loro} & erano stati \\
\end{tabular}}{}
\end{minipage}\hspace{0.02\linewidth}\begin{minipage}[t]{0.32\linewidth}\raggedright
\paradigmbox{indicativo passato remoto}{%
\begin{tabular}{@{}l l@{}}
\textit{io} & \irr{stetti} \\
\textit{tu} & \irr{stesti} \\
\textit{lui/lei} & \irr{stette} \\
\textit{noi} & \irr{stemmo} \\
\textit{voi} & \irr{steste} \\
\textit{loro} & \irr{stettero} \\
\end{tabular}}{\irr{Highlighted} forms differ from the regular pattern of the verb's family.}
\end{minipage}\hspace{0.02\linewidth}\begin{minipage}[t]{0.32\linewidth}\raggedright
\paradigmbox{indicativo trapassato remoto (with essere, masculine subject)}{%
\begin{tabular}{@{}l l@{}}
\textit{io} & fui stato \\
\textit{tu} & fosti stato \\
\textit{lui/lei} & fu stato \\
\textit{noi} & fummo stati \\
\textit{voi} & foste stati \\
\textit{loro} & furono stati \\
\end{tabular}}{}
\end{minipage}\par\medskip\noindent\begin{minipage}[t]{0.32\linewidth}\raggedright
\paradigmbox{indicativo futuro semplice}{%
\begin{tabular}{@{}l l@{}}
\textit{io} & \irr{starò} \\
\textit{tu} & \irr{starai} \\
\textit{lui/lei} & \irr{starà} \\
\textit{noi} & \irr{staremo} \\
\textit{voi} & \irr{starete} \\
\textit{loro} & \irr{staranno} \\
\end{tabular}}{\irr{Highlighted} forms differ from the regular pattern of the verb's family.}
\end{minipage}\hspace{0.02\linewidth}\begin{minipage}[t]{0.32\linewidth}\raggedright
\paradigmbox{indicativo futuro anteriore (with essere, masculine subject)}{%
\begin{tabular}{@{}l l@{}}
\textit{io} & sarò stato \\
\textit{tu} & sarai stato \\
\textit{lui/lei} & sarà stato \\
\textit{noi} & saremo stati \\
\textit{voi} & sarete stati \\
\textit{loro} & saranno stati \\
\end{tabular}}{}
\end{minipage}\par\medskip

\moodheading{congiuntivo}

\par\noindent\begin{minipage}[t]{0.32\linewidth}\raggedright
\paradigmbox{congiuntivo presente}{%
\begin{tabular}{@{}l l@{}}
\textit{io} & \irr{stia} \\
\textit{tu} & \irr{stia} \\
\textit{lui/lei} & \irr{stia} \\
\textit{noi} & stiamo \\
\textit{voi} & stiate \\
\textit{loro} & \irr{stiano} \\
\end{tabular}}{\irr{Highlighted} forms differ from the regular pattern of the verb's family.}
\end{minipage}\hspace{0.02\linewidth}\begin{minipage}[t]{0.32\linewidth}\raggedright
\paradigmbox{congiuntivo passato (with essere, masculine subject)}{%
\begin{tabular}{@{}l l@{}}
\textit{io} & sia stato \\
\textit{tu} & sia stato \\
\textit{lui/lei} & sia stato \\
\textit{noi} & siamo stati \\
\textit{voi} & siate stati \\
\textit{loro} & siano stati \\
\end{tabular}}{}
\end{minipage}\hspace{0.02\linewidth}\begin{minipage}[t]{0.32\linewidth}\raggedright
\paradigmbox{congiuntivo imperfetto}{%
\begin{tabular}{@{}l l@{}}
\textit{io} & \irr{stessi} \\
\textit{tu} & \irr{stessi} \\
\textit{lui/lei} & \irr{stesse} \\
\textit{noi} & \irr{stessimo} \\
\textit{voi} & \irr{steste} \\
\textit{loro} & \irr{stessero} \\
\end{tabular}}{\irr{Highlighted} forms differ from the regular pattern of the verb's family.}
\end{minipage}\par\medskip\noindent\begin{minipage}[t]{0.32\linewidth}\raggedright
\paradigmbox{congiuntivo trapassato (with essere, masculine subject)}{%
\begin{tabular}{@{}l l@{}}
\textit{io} & fossi stato \\
\textit{tu} & fossi stato \\
\textit{lui/lei} & fosse stato \\
\textit{noi} & fossimo stati \\
\textit{voi} & foste stati \\
\textit{loro} & fossero stati \\
\end{tabular}}{}
\end{minipage}\par\medskip

\moodheading{condizionale}

\par\noindent\begin{minipage}[t]{0.32\linewidth}\raggedright
\paradigmbox{condizionale presente}{%
\begin{tabular}{@{}l l@{}}
\textit{io} & \irr{starei} \\
\textit{tu} & \irr{staresti} \\
\textit{lui/lei} & \irr{starebbe} \\
\textit{noi} & \irr{staremmo} \\
\textit{voi} & \irr{stareste} \\
\textit{loro} & \irr{starebbero} \\
\end{tabular}}{\irr{Highlighted} forms differ from the regular pattern of the verb's family.}
\end{minipage}\hspace{0.02\linewidth}\begin{minipage}[t]{0.32\linewidth}\raggedright
\paradigmbox{condizionale passato (with essere, masculine subject)}{%
\begin{tabular}{@{}l l@{}}
\textit{io} & sarei stato \\
\textit{tu} & saresti stato \\
\textit{lui/lei} & sarebbe stato \\
\textit{noi} & saremmo stati \\
\textit{voi} & sareste stati \\
\textit{loro} & sarebbero stati \\
\end{tabular}}{}
\end{minipage}\par\medskip

\moodheading{imperativo}

\par\noindent\begin{minipage}[t]{0.32\linewidth}\raggedright
\paradigmbox{imperativo}{%
\begin{tabular}{@{}l l@{}}
\textit{tu} & \irr{sta' / stai} \\
\textit{Lei} & \irr{stia} \\
\textit{noi} & stiamo \\
\textit{voi} & state \\
\textit{Loro} & \irr{stiano} \\
\end{tabular}}{\irr{Highlighted} forms differ from the regular pattern of the verb's family. Where two forms are given, both are standard. Negative \textit{tu}: \textit{non stare}. \textit{Lei} and \textit{Loro} are the congiuntivo forms.}
\end{minipage}\par\medskip

\moodheading{forme non finite}
{\small\begin{tabular}{@{}l l@{}}
infinito presente & \textit{stare} \\
infinito passato & \textit{essere stato} \\
gerundio presente & \textit{stando} \\
gerundio passato & \textit{essendo stato} \\
participio passato & \textit{stato} \\
\end{tabular}}\par
\clearpage

\subsection{dare --- to give}\label{conj:dare}\index{dare@dare}
{\small Auxiliary: \textit{avere}. dà (gives) carries an accent to distinguish it from the preposition da (from).}\par

\moodheading{indicativo}

\par\noindent\begin{minipage}[t]{0.32\linewidth}\raggedright
\paradigmbox{indicativo presente}{%
\begin{tabular}{@{}l l@{}}
\textit{io} & do \\
\textit{tu} & \irr{dai} \\
\textit{lui/lei} & \irr{dà} \\
\textit{noi} & diamo \\
\textit{voi} & date \\
\textit{loro} & \irr{danno} \\
\end{tabular}}{\irr{Highlighted} forms differ from the regular pattern of the verb's family.}
\end{minipage}\hspace{0.02\linewidth}\begin{minipage}[t]{0.32\linewidth}\raggedright
\paradigmbox{indicativo passato prossimo (with avere)}{%
\begin{tabular}{@{}l l@{}}
\textit{io} & ho dato \\
\textit{tu} & hai dato \\
\textit{lui/lei} & ha dato \\
\textit{noi} & abbiamo dato \\
\textit{voi} & avete dato \\
\textit{loro} & hanno dato \\
\end{tabular}}{}
\end{minipage}\hspace{0.02\linewidth}\begin{minipage}[t]{0.32\linewidth}\raggedright
\paradigmbox{indicativo imperfetto}{%
\begin{tabular}{@{}l l@{}}
\textit{io} & davo \\
\textit{tu} & davi \\
\textit{lui/lei} & dava \\
\textit{noi} & davamo \\
\textit{voi} & davate \\
\textit{loro} & davano \\
\end{tabular}}{}
\end{minipage}\par\medskip\noindent\begin{minipage}[t]{0.32\linewidth}\raggedright
\paradigmbox{indicativo trapassato prossimo (with avere)}{%
\begin{tabular}{@{}l l@{}}
\textit{io} & avevo dato \\
\textit{tu} & avevi dato \\
\textit{lui/lei} & aveva dato \\
\textit{noi} & avevamo dato \\
\textit{voi} & avevate dato \\
\textit{loro} & avevano dato \\
\end{tabular}}{}
\end{minipage}\hspace{0.02\linewidth}\begin{minipage}[t]{0.32\linewidth}\raggedright
\paradigmbox{indicativo passato remoto}{%
\begin{tabular}{@{}l l@{}}
\textit{io} & \irr{diedi / detti} \\
\textit{tu} & \irr{desti} \\
\textit{lui/lei} & \irr{diede / dette} \\
\textit{noi} & \irr{demmo} \\
\textit{voi} & \irr{deste} \\
\textit{loro} & \irr{diedero / dettero} \\
\end{tabular}}{\irr{Highlighted} forms differ from the regular pattern of the verb's family. Where two forms are given, both are standard.}
\end{minipage}\hspace{0.02\linewidth}\begin{minipage}[t]{0.32\linewidth}\raggedright
\paradigmbox{indicativo trapassato remoto (with avere)}{%
\begin{tabular}{@{}l l@{}}
\textit{io} & ebbi dato \\
\textit{tu} & avesti dato \\
\textit{lui/lei} & ebbe dato \\
\textit{noi} & avemmo dato \\
\textit{voi} & aveste dato \\
\textit{loro} & ebbero dato \\
\end{tabular}}{}
\end{minipage}\par\medskip\noindent\begin{minipage}[t]{0.32\linewidth}\raggedright
\paradigmbox{indicativo futuro semplice}{%
\begin{tabular}{@{}l l@{}}
\textit{io} & \irr{darò} \\
\textit{tu} & \irr{darai} \\
\textit{lui/lei} & \irr{darà} \\
\textit{noi} & \irr{daremo} \\
\textit{voi} & \irr{darete} \\
\textit{loro} & \irr{daranno} \\
\end{tabular}}{\irr{Highlighted} forms differ from the regular pattern of the verb's family.}
\end{minipage}\hspace{0.02\linewidth}\begin{minipage}[t]{0.32\linewidth}\raggedright
\paradigmbox{indicativo futuro anteriore (with avere)}{%
\begin{tabular}{@{}l l@{}}
\textit{io} & avrò dato \\
\textit{tu} & avrai dato \\
\textit{lui/lei} & avrà dato \\
\textit{noi} & avremo dato \\
\textit{voi} & avrete dato \\
\textit{loro} & avranno dato \\
\end{tabular}}{}
\end{minipage}\par\medskip

\moodheading{congiuntivo}

\par\noindent\begin{minipage}[t]{0.32\linewidth}\raggedright
\paradigmbox{congiuntivo presente}{%
\begin{tabular}{@{}l l@{}}
\textit{io} & \irr{dia} \\
\textit{tu} & \irr{dia} \\
\textit{lui/lei} & \irr{dia} \\
\textit{noi} & diamo \\
\textit{voi} & diate \\
\textit{loro} & \irr{diano} \\
\end{tabular}}{\irr{Highlighted} forms differ from the regular pattern of the verb's family.}
\end{minipage}\hspace{0.02\linewidth}\begin{minipage}[t]{0.32\linewidth}\raggedright
\paradigmbox{congiuntivo passato (with avere)}{%
\begin{tabular}{@{}l l@{}}
\textit{io} & abbia dato \\
\textit{tu} & abbia dato \\
\textit{lui/lei} & abbia dato \\
\textit{noi} & abbiamo dato \\
\textit{voi} & abbiate dato \\
\textit{loro} & abbiano dato \\
\end{tabular}}{}
\end{minipage}\hspace{0.02\linewidth}\begin{minipage}[t]{0.32\linewidth}\raggedright
\paradigmbox{congiuntivo imperfetto}{%
\begin{tabular}{@{}l l@{}}
\textit{io} & \irr{dessi} \\
\textit{tu} & \irr{dessi} \\
\textit{lui/lei} & \irr{desse} \\
\textit{noi} & \irr{dessimo} \\
\textit{voi} & \irr{deste} \\
\textit{loro} & \irr{dessero} \\
\end{tabular}}{\irr{Highlighted} forms differ from the regular pattern of the verb's family.}
\end{minipage}\par\medskip\noindent\begin{minipage}[t]{0.32\linewidth}\raggedright
\paradigmbox{congiuntivo trapassato (with avere)}{%
\begin{tabular}{@{}l l@{}}
\textit{io} & avessi dato \\
\textit{tu} & avessi dato \\
\textit{lui/lei} & avesse dato \\
\textit{noi} & avessimo dato \\
\textit{voi} & aveste dato \\
\textit{loro} & avessero dato \\
\end{tabular}}{}
\end{minipage}\par\medskip

\moodheading{condizionale}

\par\noindent\begin{minipage}[t]{0.32\linewidth}\raggedright
\paradigmbox{condizionale presente}{%
\begin{tabular}{@{}l l@{}}
\textit{io} & \irr{darei} \\
\textit{tu} & \irr{daresti} \\
\textit{lui/lei} & \irr{darebbe} \\
\textit{noi} & \irr{daremmo} \\
\textit{voi} & \irr{dareste} \\
\textit{loro} & \irr{darebbero} \\
\end{tabular}}{\irr{Highlighted} forms differ from the regular pattern of the verb's family.}
\end{minipage}\hspace{0.02\linewidth}\begin{minipage}[t]{0.32\linewidth}\raggedright
\paradigmbox{condizionale passato (with avere)}{%
\begin{tabular}{@{}l l@{}}
\textit{io} & avrei dato \\
\textit{tu} & avresti dato \\
\textit{lui/lei} & avrebbe dato \\
\textit{noi} & avremmo dato \\
\textit{voi} & avreste dato \\
\textit{loro} & avrebbero dato \\
\end{tabular}}{}
\end{minipage}\par\medskip

\moodheading{imperativo}

\par\noindent\begin{minipage}[t]{0.32\linewidth}\raggedright
\paradigmbox{imperativo}{%
\begin{tabular}{@{}l l@{}}
\textit{tu} & \irr{da' / dai} \\
\textit{Lei} & \irr{dia} \\
\textit{noi} & diamo \\
\textit{voi} & date \\
\textit{Loro} & \irr{diano} \\
\end{tabular}}{\irr{Highlighted} forms differ from the regular pattern of the verb's family. Where two forms are given, both are standard. Negative \textit{tu}: \textit{non dare}. \textit{Lei} and \textit{Loro} are the congiuntivo forms.}
\end{minipage}\par\medskip

\moodheading{forme non finite}
{\small\begin{tabular}{@{}l l@{}}
infinito presente & \textit{dare} \\
infinito passato & \textit{avere dato} \\
gerundio presente & \textit{dando} \\
gerundio passato & \textit{avendo dato} \\
participio passato & \textit{dato} \\
\end{tabular}}\par
\clearpage

\subsection{dire --- to say; to tell}\label{conj:dire}\index{dire@dire}
{\small Auxiliary: \textit{avere}. From Latin dicere: the stem dic- appears everywhere except dite. The c is hard in dico, dicono and soft in dici, dice.}\par

\moodheading{indicativo}

\par\noindent\begin{minipage}[t]{0.32\linewidth}\raggedright
\paradigmbox{indicativo presente}{%
\begin{tabular}{@{}l l@{}}
\textit{io} & \irr{dico} \\
\textit{tu} & dici \\
\textit{lui/lei} & \irr{dice} \\
\textit{noi} & \irr{diciamo} \\
\textit{voi} & dite \\
\textit{loro} & \irr{dicono} \\
\end{tabular}}{\irr{Highlighted} forms differ from the regular pattern of the verb's family.}
\end{minipage}\hspace{0.02\linewidth}\begin{minipage}[t]{0.32\linewidth}\raggedright
\paradigmbox{indicativo passato prossimo (with avere)}{%
\begin{tabular}{@{}l l@{}}
\textit{io} & ho detto \\
\textit{tu} & hai detto \\
\textit{lui/lei} & ha detto \\
\textit{noi} & abbiamo detto \\
\textit{voi} & avete detto \\
\textit{loro} & hanno detto \\
\end{tabular}}{}
\end{minipage}\hspace{0.02\linewidth}\begin{minipage}[t]{0.32\linewidth}\raggedright
\paradigmbox{indicativo imperfetto}{%
\begin{tabular}{@{}l l@{}}
\textit{io} & \irr{dicevo} \\
\textit{tu} & \irr{dicevi} \\
\textit{lui/lei} & \irr{diceva} \\
\textit{noi} & \irr{dicevamo} \\
\textit{voi} & \irr{dicevate} \\
\textit{loro} & \irr{dicevano} \\
\end{tabular}}{\irr{Highlighted} forms differ from the regular pattern of the verb's family.}
\end{minipage}\par\medskip\noindent\begin{minipage}[t]{0.32\linewidth}\raggedright
\paradigmbox{indicativo trapassato prossimo (with avere)}{%
\begin{tabular}{@{}l l@{}}
\textit{io} & avevo detto \\
\textit{tu} & avevi detto \\
\textit{lui/lei} & aveva detto \\
\textit{noi} & avevamo detto \\
\textit{voi} & avevate detto \\
\textit{loro} & avevano detto \\
\end{tabular}}{}
\end{minipage}\hspace{0.02\linewidth}\begin{minipage}[t]{0.32\linewidth}\raggedright
\paradigmbox{indicativo passato remoto}{%
\begin{tabular}{@{}l l@{}}
\textit{io} & \irr{dissi} \\
\textit{tu} & \irr{dicesti} \\
\textit{lui/lei} & \irr{disse} \\
\textit{noi} & \irr{dicemmo} \\
\textit{voi} & \irr{diceste} \\
\textit{loro} & \irr{dissero} \\
\end{tabular}}{\irr{Highlighted} forms differ from the regular pattern of the verb's family.}
\end{minipage}\hspace{0.02\linewidth}\begin{minipage}[t]{0.32\linewidth}\raggedright
\paradigmbox{indicativo trapassato remoto (with avere)}{%
\begin{tabular}{@{}l l@{}}
\textit{io} & ebbi detto \\
\textit{tu} & avesti detto \\
\textit{lui/lei} & ebbe detto \\
\textit{noi} & avemmo detto \\
\textit{voi} & aveste detto \\
\textit{loro} & ebbero detto \\
\end{tabular}}{}
\end{minipage}\par\medskip\noindent\begin{minipage}[t]{0.32\linewidth}\raggedright
\paradigmbox{indicativo futuro semplice}{%
\begin{tabular}{@{}l l@{}}
\textit{io} & dirò \\
\textit{tu} & dirai \\
\textit{lui/lei} & dirà \\
\textit{noi} & diremo \\
\textit{voi} & direte \\
\textit{loro} & diranno \\
\end{tabular}}{}
\end{minipage}\hspace{0.02\linewidth}\begin{minipage}[t]{0.32\linewidth}\raggedright
\paradigmbox{indicativo futuro anteriore (with avere)}{%
\begin{tabular}{@{}l l@{}}
\textit{io} & avrò detto \\
\textit{tu} & avrai detto \\
\textit{lui/lei} & avrà detto \\
\textit{noi} & avremo detto \\
\textit{voi} & avrete detto \\
\textit{loro} & avranno detto \\
\end{tabular}}{}
\end{minipage}\par\medskip

\moodheading{congiuntivo}

\par\noindent\begin{minipage}[t]{0.32\linewidth}\raggedright
\paradigmbox{congiuntivo presente}{%
\begin{tabular}{@{}l l@{}}
\textit{io} & \irr{dica} \\
\textit{tu} & \irr{dica} \\
\textit{lui/lei} & \irr{dica} \\
\textit{noi} & \irr{diciamo} \\
\textit{voi} & \irr{diciate} \\
\textit{loro} & \irr{dicano} \\
\end{tabular}}{\irr{Highlighted} forms differ from the regular pattern of the verb's family.}
\end{minipage}\hspace{0.02\linewidth}\begin{minipage}[t]{0.32\linewidth}\raggedright
\paradigmbox{congiuntivo passato (with avere)}{%
\begin{tabular}{@{}l l@{}}
\textit{io} & abbia detto \\
\textit{tu} & abbia detto \\
\textit{lui/lei} & abbia detto \\
\textit{noi} & abbiamo detto \\
\textit{voi} & abbiate detto \\
\textit{loro} & abbiano detto \\
\end{tabular}}{}
\end{minipage}\hspace{0.02\linewidth}\begin{minipage}[t]{0.32\linewidth}\raggedright
\paradigmbox{congiuntivo imperfetto}{%
\begin{tabular}{@{}l l@{}}
\textit{io} & \irr{dicessi} \\
\textit{tu} & \irr{dicessi} \\
\textit{lui/lei} & \irr{dicesse} \\
\textit{noi} & \irr{dicessimo} \\
\textit{voi} & \irr{diceste} \\
\textit{loro} & \irr{dicessero} \\
\end{tabular}}{\irr{Highlighted} forms differ from the regular pattern of the verb's family.}
\end{minipage}\par\medskip\noindent\begin{minipage}[t]{0.32\linewidth}\raggedright
\paradigmbox{congiuntivo trapassato (with avere)}{%
\begin{tabular}{@{}l l@{}}
\textit{io} & avessi detto \\
\textit{tu} & avessi detto \\
\textit{lui/lei} & avesse detto \\
\textit{noi} & avessimo detto \\
\textit{voi} & aveste detto \\
\textit{loro} & avessero detto \\
\end{tabular}}{}
\end{minipage}\par\medskip

\moodheading{condizionale}

\par\noindent\begin{minipage}[t]{0.32\linewidth}\raggedright
\paradigmbox{condizionale presente}{%
\begin{tabular}{@{}l l@{}}
\textit{io} & direi \\
\textit{tu} & diresti \\
\textit{lui/lei} & direbbe \\
\textit{noi} & diremmo \\
\textit{voi} & direste \\
\textit{loro} & direbbero \\
\end{tabular}}{}
\end{minipage}\hspace{0.02\linewidth}\begin{minipage}[t]{0.32\linewidth}\raggedright
\paradigmbox{condizionale passato (with avere)}{%
\begin{tabular}{@{}l l@{}}
\textit{io} & avrei detto \\
\textit{tu} & avresti detto \\
\textit{lui/lei} & avrebbe detto \\
\textit{noi} & avremmo detto \\
\textit{voi} & avreste detto \\
\textit{loro} & avrebbero detto \\
\end{tabular}}{}
\end{minipage}\par\medskip

\moodheading{imperativo}

\par\noindent\begin{minipage}[t]{0.32\linewidth}\raggedright
\paradigmbox{imperativo}{%
\begin{tabular}{@{}l l@{}}
\textit{tu} & \irr{di'} \\
\textit{Lei} & \irr{dica} \\
\textit{noi} & \irr{diciamo} \\
\textit{voi} & dite \\
\textit{Loro} & \irr{dicano} \\
\end{tabular}}{\irr{Highlighted} forms differ from the regular pattern of the verb's family. Negative \textit{tu}: \textit{non dire}. \textit{Lei} and \textit{Loro} are the congiuntivo forms.}
\end{minipage}\par\medskip

\moodheading{forme non finite}
{\small\begin{tabular}{@{}l l@{}}
infinito presente & \textit{dire} \\
infinito passato & \textit{avere detto} \\
gerundio presente & \textit{dicendo} \\
gerundio passato & \textit{avendo detto} \\
participio passato & \textit{detto} \\
\end{tabular}}\par
\clearpage

\subsection{venire --- to come}\label{conj:venire}\index{venire@venire}
{\small Auxiliary: \textit{essere}. A g appears in io and loro (vengo, vengono); the stressed e becomes ie in tu and lui/lei (vieni, viene). The participle is venuto, not *venito.}\par

\moodheading{indicativo}

\par\noindent\begin{minipage}[t]{0.32\linewidth}\raggedright
\paradigmbox{indicativo presente}{%
\begin{tabular}{@{}l l@{}}
\textit{io} & \irr{vengo} \\
\textit{tu} & \irr{vieni} \\
\textit{lui/lei} & \irr{viene} \\
\textit{noi} & veniamo \\
\textit{voi} & venite \\
\textit{loro} & \irr{vengono} \\
\end{tabular}}{\irr{Highlighted} forms differ from the regular pattern of the verb's family.}
\end{minipage}\hspace{0.02\linewidth}\begin{minipage}[t]{0.32\linewidth}\raggedright
\paradigmbox{indicativo passato prossimo (with essere, masculine subject)}{%
\begin{tabular}{@{}l l@{}}
\textit{io} & sono venuto \\
\textit{tu} & sei venuto \\
\textit{lui/lei} & è venuto \\
\textit{noi} & siamo venuti \\
\textit{voi} & siete venuti \\
\textit{loro} & sono venuti \\
\end{tabular}}{}
\end{minipage}\hspace{0.02\linewidth}\begin{minipage}[t]{0.32\linewidth}\raggedright
\paradigmbox{indicativo imperfetto}{%
\begin{tabular}{@{}l l@{}}
\textit{io} & venivo \\
\textit{tu} & venivi \\
\textit{lui/lei} & veniva \\
\textit{noi} & venivamo \\
\textit{voi} & venivate \\
\textit{loro} & venivano \\
\end{tabular}}{}
\end{minipage}\par\medskip\noindent\begin{minipage}[t]{0.32\linewidth}\raggedright
\paradigmbox{indicativo trapassato prossimo (with essere, masculine subject)}{%
\begin{tabular}{@{}l l@{}}
\textit{io} & ero venuto \\
\textit{tu} & eri venuto \\
\textit{lui/lei} & era venuto \\
\textit{noi} & eravamo venuti \\
\textit{voi} & eravate venuti \\
\textit{loro} & erano venuti \\
\end{tabular}}{}
\end{minipage}\hspace{0.02\linewidth}\begin{minipage}[t]{0.32\linewidth}\raggedright
\paradigmbox{indicativo passato remoto}{%
\begin{tabular}{@{}l l@{}}
\textit{io} & \irr{venni} \\
\textit{tu} & venisti \\
\textit{lui/lei} & \irr{venne} \\
\textit{noi} & venimmo \\
\textit{voi} & veniste \\
\textit{loro} & \irr{vennero} \\
\end{tabular}}{\irr{Highlighted} forms differ from the regular pattern of the verb's family.}
\end{minipage}\hspace{0.02\linewidth}\begin{minipage}[t]{0.32\linewidth}\raggedright
\paradigmbox{indicativo trapassato remoto (with essere, masculine subject)}{%
\begin{tabular}{@{}l l@{}}
\textit{io} & fui venuto \\
\textit{tu} & fosti venuto \\
\textit{lui/lei} & fu venuto \\
\textit{noi} & fummo venuti \\
\textit{voi} & foste venuti \\
\textit{loro} & furono venuti \\
\end{tabular}}{}
\end{minipage}\par\medskip\noindent\begin{minipage}[t]{0.32\linewidth}\raggedright
\paradigmbox{indicativo futuro semplice}{%
\begin{tabular}{@{}l l@{}}
\textit{io} & \irr{verrò} \\
\textit{tu} & \irr{verrai} \\
\textit{lui/lei} & \irr{verrà} \\
\textit{noi} & \irr{verremo} \\
\textit{voi} & \irr{verrete} \\
\textit{loro} & \irr{verranno} \\
\end{tabular}}{\irr{Highlighted} forms differ from the regular pattern of the verb's family.}
\end{minipage}\hspace{0.02\linewidth}\begin{minipage}[t]{0.32\linewidth}\raggedright
\paradigmbox{indicativo futuro anteriore (with essere, masculine subject)}{%
\begin{tabular}{@{}l l@{}}
\textit{io} & sarò venuto \\
\textit{tu} & sarai venuto \\
\textit{lui/lei} & sarà venuto \\
\textit{noi} & saremo venuti \\
\textit{voi} & sarete venuti \\
\textit{loro} & saranno venuti \\
\end{tabular}}{}
\end{minipage}\par\medskip

\moodheading{congiuntivo}

\par\noindent\begin{minipage}[t]{0.32\linewidth}\raggedright
\paradigmbox{congiuntivo presente}{%
\begin{tabular}{@{}l l@{}}
\textit{io} & \irr{venga} \\
\textit{tu} & \irr{venga} \\
\textit{lui/lei} & \irr{venga} \\
\textit{noi} & veniamo \\
\textit{voi} & veniate \\
\textit{loro} & \irr{vengano} \\
\end{tabular}}{\irr{Highlighted} forms differ from the regular pattern of the verb's family.}
\end{minipage}\hspace{0.02\linewidth}\begin{minipage}[t]{0.32\linewidth}\raggedright
\paradigmbox{congiuntivo passato (with essere, masculine subject)}{%
\begin{tabular}{@{}l l@{}}
\textit{io} & sia venuto \\
\textit{tu} & sia venuto \\
\textit{lui/lei} & sia venuto \\
\textit{noi} & siamo venuti \\
\textit{voi} & siate venuti \\
\textit{loro} & siano venuti \\
\end{tabular}}{}
\end{minipage}\hspace{0.02\linewidth}\begin{minipage}[t]{0.32\linewidth}\raggedright
\paradigmbox{congiuntivo imperfetto}{%
\begin{tabular}{@{}l l@{}}
\textit{io} & venissi \\
\textit{tu} & venissi \\
\textit{lui/lei} & venisse \\
\textit{noi} & venissimo \\
\textit{voi} & veniste \\
\textit{loro} & venissero \\
\end{tabular}}{}
\end{minipage}\par\medskip\noindent\begin{minipage}[t]{0.32\linewidth}\raggedright
\paradigmbox{congiuntivo trapassato (with essere, masculine subject)}{%
\begin{tabular}{@{}l l@{}}
\textit{io} & fossi venuto \\
\textit{tu} & fossi venuto \\
\textit{lui/lei} & fosse venuto \\
\textit{noi} & fossimo venuti \\
\textit{voi} & foste venuti \\
\textit{loro} & fossero venuti \\
\end{tabular}}{}
\end{minipage}\par\medskip

\moodheading{condizionale}

\par\noindent\begin{minipage}[t]{0.32\linewidth}\raggedright
\paradigmbox{condizionale presente}{%
\begin{tabular}{@{}l l@{}}
\textit{io} & \irr{verrei} \\
\textit{tu} & \irr{verresti} \\
\textit{lui/lei} & \irr{verrebbe} \\
\textit{noi} & \irr{verremmo} \\
\textit{voi} & \irr{verreste} \\
\textit{loro} & \irr{verrebbero} \\
\end{tabular}}{\irr{Highlighted} forms differ from the regular pattern of the verb's family.}
\end{minipage}\hspace{0.02\linewidth}\begin{minipage}[t]{0.32\linewidth}\raggedright
\paradigmbox{condizionale passato (with essere, masculine subject)}{%
\begin{tabular}{@{}l l@{}}
\textit{io} & sarei venuto \\
\textit{tu} & saresti venuto \\
\textit{lui/lei} & sarebbe venuto \\
\textit{noi} & saremmo venuti \\
\textit{voi} & sareste venuti \\
\textit{loro} & sarebbero venuti \\
\end{tabular}}{}
\end{minipage}\par\medskip

\moodheading{imperativo}

\par\noindent\begin{minipage}[t]{0.32\linewidth}\raggedright
\paradigmbox{imperativo}{%
\begin{tabular}{@{}l l@{}}
\textit{tu} & \irr{vieni} \\
\textit{Lei} & \irr{venga} \\
\textit{noi} & veniamo \\
\textit{voi} & venite \\
\textit{Loro} & \irr{vengano} \\
\end{tabular}}{\irr{Highlighted} forms differ from the regular pattern of the verb's family. Negative \textit{tu}: \textit{non venire}. \textit{Lei} and \textit{Loro} are the congiuntivo forms.}
\end{minipage}\par\medskip

\moodheading{forme non finite}
{\small\begin{tabular}{@{}l l@{}}
infinito presente & \textit{venire} \\
infinito passato & \textit{essere venuto} \\
gerundio presente & \textit{venendo} \\
gerundio passato & \textit{essendo venuto} \\
participio passato & \textit{venuto} \\
\end{tabular}}\par
\clearpage

\subsection{tenere --- to hold; to keep}\label{conj:tenere}\index{tenere@tenere}
{\small Auxiliary: \textit{avere}. Exactly the same pattern as venire: tengo, tieni, tiene, …, tengono.}\par

\moodheading{indicativo}

\par\noindent\begin{minipage}[t]{0.32\linewidth}\raggedright
\paradigmbox{indicativo presente}{%
\begin{tabular}{@{}l l@{}}
\textit{io} & \irr{tengo} \\
\textit{tu} & \irr{tieni} \\
\textit{lui/lei} & \irr{tiene} \\
\textit{noi} & teniamo \\
\textit{voi} & tenete \\
\textit{loro} & \irr{tengono} \\
\end{tabular}}{\irr{Highlighted} forms differ from the regular pattern of the verb's family.}
\end{minipage}\hspace{0.02\linewidth}\begin{minipage}[t]{0.32\linewidth}\raggedright
\paradigmbox{indicativo passato prossimo (with avere)}{%
\begin{tabular}{@{}l l@{}}
\textit{io} & ho tenuto \\
\textit{tu} & hai tenuto \\
\textit{lui/lei} & ha tenuto \\
\textit{noi} & abbiamo tenuto \\
\textit{voi} & avete tenuto \\
\textit{loro} & hanno tenuto \\
\end{tabular}}{}
\end{minipage}\hspace{0.02\linewidth}\begin{minipage}[t]{0.32\linewidth}\raggedright
\paradigmbox{indicativo imperfetto}{%
\begin{tabular}{@{}l l@{}}
\textit{io} & tenevo \\
\textit{tu} & tenevi \\
\textit{lui/lei} & teneva \\
\textit{noi} & tenevamo \\
\textit{voi} & tenevate \\
\textit{loro} & tenevano \\
\end{tabular}}{}
\end{minipage}\par\medskip\noindent\begin{minipage}[t]{0.32\linewidth}\raggedright
\paradigmbox{indicativo trapassato prossimo (with avere)}{%
\begin{tabular}{@{}l l@{}}
\textit{io} & avevo tenuto \\
\textit{tu} & avevi tenuto \\
\textit{lui/lei} & aveva tenuto \\
\textit{noi} & avevamo tenuto \\
\textit{voi} & avevate tenuto \\
\textit{loro} & avevano tenuto \\
\end{tabular}}{}
\end{minipage}\hspace{0.02\linewidth}\begin{minipage}[t]{0.32\linewidth}\raggedright
\paradigmbox{indicativo passato remoto}{%
\begin{tabular}{@{}l l@{}}
\textit{io} & \irr{tenni} \\
\textit{tu} & tenesti \\
\textit{lui/lei} & \irr{tenne} \\
\textit{noi} & tenemmo \\
\textit{voi} & teneste \\
\textit{loro} & \irr{tennero} \\
\end{tabular}}{\irr{Highlighted} forms differ from the regular pattern of the verb's family.}
\end{minipage}\hspace{0.02\linewidth}\begin{minipage}[t]{0.32\linewidth}\raggedright
\paradigmbox{indicativo trapassato remoto (with avere)}{%
\begin{tabular}{@{}l l@{}}
\textit{io} & ebbi tenuto \\
\textit{tu} & avesti tenuto \\
\textit{lui/lei} & ebbe tenuto \\
\textit{noi} & avemmo tenuto \\
\textit{voi} & aveste tenuto \\
\textit{loro} & ebbero tenuto \\
\end{tabular}}{}
\end{minipage}\par\medskip\noindent\begin{minipage}[t]{0.32\linewidth}\raggedright
\paradigmbox{indicativo futuro semplice}{%
\begin{tabular}{@{}l l@{}}
\textit{io} & \irr{terrò} \\
\textit{tu} & \irr{terrai} \\
\textit{lui/lei} & \irr{terrà} \\
\textit{noi} & \irr{terremo} \\
\textit{voi} & \irr{terrete} \\
\textit{loro} & \irr{terranno} \\
\end{tabular}}{\irr{Highlighted} forms differ from the regular pattern of the verb's family.}
\end{minipage}\hspace{0.02\linewidth}\begin{minipage}[t]{0.32\linewidth}\raggedright
\paradigmbox{indicativo futuro anteriore (with avere)}{%
\begin{tabular}{@{}l l@{}}
\textit{io} & avrò tenuto \\
\textit{tu} & avrai tenuto \\
\textit{lui/lei} & avrà tenuto \\
\textit{noi} & avremo tenuto \\
\textit{voi} & avrete tenuto \\
\textit{loro} & avranno tenuto \\
\end{tabular}}{}
\end{minipage}\par\medskip

\moodheading{congiuntivo}

\par\noindent\begin{minipage}[t]{0.32\linewidth}\raggedright
\paradigmbox{congiuntivo presente}{%
\begin{tabular}{@{}l l@{}}
\textit{io} & \irr{tenga} \\
\textit{tu} & \irr{tenga} \\
\textit{lui/lei} & \irr{tenga} \\
\textit{noi} & teniamo \\
\textit{voi} & teniate \\
\textit{loro} & \irr{tengano} \\
\end{tabular}}{\irr{Highlighted} forms differ from the regular pattern of the verb's family.}
\end{minipage}\hspace{0.02\linewidth}\begin{minipage}[t]{0.32\linewidth}\raggedright
\paradigmbox{congiuntivo passato (with avere)}{%
\begin{tabular}{@{}l l@{}}
\textit{io} & abbia tenuto \\
\textit{tu} & abbia tenuto \\
\textit{lui/lei} & abbia tenuto \\
\textit{noi} & abbiamo tenuto \\
\textit{voi} & abbiate tenuto \\
\textit{loro} & abbiano tenuto \\
\end{tabular}}{}
\end{minipage}\hspace{0.02\linewidth}\begin{minipage}[t]{0.32\linewidth}\raggedright
\paradigmbox{congiuntivo imperfetto}{%
\begin{tabular}{@{}l l@{}}
\textit{io} & tenessi \\
\textit{tu} & tenessi \\
\textit{lui/lei} & tenesse \\
\textit{noi} & tenessimo \\
\textit{voi} & teneste \\
\textit{loro} & tenessero \\
\end{tabular}}{}
\end{minipage}\par\medskip\noindent\begin{minipage}[t]{0.32\linewidth}\raggedright
\paradigmbox{congiuntivo trapassato (with avere)}{%
\begin{tabular}{@{}l l@{}}
\textit{io} & avessi tenuto \\
\textit{tu} & avessi tenuto \\
\textit{lui/lei} & avesse tenuto \\
\textit{noi} & avessimo tenuto \\
\textit{voi} & aveste tenuto \\
\textit{loro} & avessero tenuto \\
\end{tabular}}{}
\end{minipage}\par\medskip

\moodheading{condizionale}

\par\noindent\begin{minipage}[t]{0.32\linewidth}\raggedright
\paradigmbox{condizionale presente}{%
\begin{tabular}{@{}l l@{}}
\textit{io} & \irr{terrei} \\
\textit{tu} & \irr{terresti} \\
\textit{lui/lei} & \irr{terrebbe} \\
\textit{noi} & \irr{terremmo} \\
\textit{voi} & \irr{terreste} \\
\textit{loro} & \irr{terrebbero} \\
\end{tabular}}{\irr{Highlighted} forms differ from the regular pattern of the verb's family.}
\end{minipage}\hspace{0.02\linewidth}\begin{minipage}[t]{0.32\linewidth}\raggedright
\paradigmbox{condizionale passato (with avere)}{%
\begin{tabular}{@{}l l@{}}
\textit{io} & avrei tenuto \\
\textit{tu} & avresti tenuto \\
\textit{lui/lei} & avrebbe tenuto \\
\textit{noi} & avremmo tenuto \\
\textit{voi} & avreste tenuto \\
\textit{loro} & avrebbero tenuto \\
\end{tabular}}{}
\end{minipage}\par\medskip

\moodheading{imperativo}

\par\noindent\begin{minipage}[t]{0.32\linewidth}\raggedright
\paradigmbox{imperativo}{%
\begin{tabular}{@{}l l@{}}
\textit{tu} & \irr{tieni} \\
\textit{Lei} & \irr{tenga} \\
\textit{noi} & teniamo \\
\textit{voi} & tenete \\
\textit{Loro} & \irr{tengano} \\
\end{tabular}}{\irr{Highlighted} forms differ from the regular pattern of the verb's family. Negative \textit{tu}: \textit{non tenere}. \textit{Lei} and \textit{Loro} are the congiuntivo forms.}
\end{minipage}\par\medskip

\moodheading{forme non finite}
{\small\begin{tabular}{@{}l l@{}}
infinito presente & \textit{tenere} \\
infinito passato & \textit{avere tenuto} \\
gerundio presente & \textit{tenendo} \\
gerundio passato & \textit{avendo tenuto} \\
participio passato & \textit{tenuto} \\
\end{tabular}}\par
\clearpage

\subsection{uscire --- to go out; to leave}\label{conj:uscire}\index{uscire@uscire}
{\small Auxiliary: \textit{essere}. The stressed forms begin with e- (esco, esci, esce, escono); noi and voi keep u- (usciamo, uscite).}\par

\moodheading{indicativo}

\par\noindent\begin{minipage}[t]{0.32\linewidth}\raggedright
\paradigmbox{indicativo presente}{%
\begin{tabular}{@{}l l@{}}
\textit{io} & \irr{esco} \\
\textit{tu} & \irr{esci} \\
\textit{lui/lei} & \irr{esce} \\
\textit{noi} & usciamo \\
\textit{voi} & uscite \\
\textit{loro} & \irr{escono} \\
\end{tabular}}{\irr{Highlighted} forms differ from the regular pattern of the verb's family.}
\end{minipage}\hspace{0.02\linewidth}\begin{minipage}[t]{0.32\linewidth}\raggedright
\paradigmbox{indicativo passato prossimo (with essere, masculine subject)}{%
\begin{tabular}{@{}l l@{}}
\textit{io} & sono uscito \\
\textit{tu} & sei uscito \\
\textit{lui/lei} & è uscito \\
\textit{noi} & siamo usciti \\
\textit{voi} & siete usciti \\
\textit{loro} & sono usciti \\
\end{tabular}}{}
\end{minipage}\hspace{0.02\linewidth}\begin{minipage}[t]{0.32\linewidth}\raggedright
\paradigmbox{indicativo imperfetto}{%
\begin{tabular}{@{}l l@{}}
\textit{io} & uscivo \\
\textit{tu} & uscivi \\
\textit{lui/lei} & usciva \\
\textit{noi} & uscivamo \\
\textit{voi} & uscivate \\
\textit{loro} & uscivano \\
\end{tabular}}{}
\end{minipage}\par\medskip\noindent\begin{minipage}[t]{0.32\linewidth}\raggedright
\paradigmbox{indicativo trapassato prossimo (with essere, masculine subject)}{%
\begin{tabular}{@{}l l@{}}
\textit{io} & ero uscito \\
\textit{tu} & eri uscito \\
\textit{lui/lei} & era uscito \\
\textit{noi} & eravamo usciti \\
\textit{voi} & eravate usciti \\
\textit{loro} & erano usciti \\
\end{tabular}}{}
\end{minipage}\hspace{0.02\linewidth}\begin{minipage}[t]{0.32\linewidth}\raggedright
\paradigmbox{indicativo passato remoto}{%
\begin{tabular}{@{}l l@{}}
\textit{io} & uscii \\
\textit{tu} & uscisti \\
\textit{lui/lei} & uscì \\
\textit{noi} & uscimmo \\
\textit{voi} & usciste \\
\textit{loro} & uscirono \\
\end{tabular}}{}
\end{minipage}\hspace{0.02\linewidth}\begin{minipage}[t]{0.32\linewidth}\raggedright
\paradigmbox{indicativo trapassato remoto (with essere, masculine subject)}{%
\begin{tabular}{@{}l l@{}}
\textit{io} & fui uscito \\
\textit{tu} & fosti uscito \\
\textit{lui/lei} & fu uscito \\
\textit{noi} & fummo usciti \\
\textit{voi} & foste usciti \\
\textit{loro} & furono usciti \\
\end{tabular}}{}
\end{minipage}\par\medskip\noindent\begin{minipage}[t]{0.32\linewidth}\raggedright
\paradigmbox{indicativo futuro semplice}{%
\begin{tabular}{@{}l l@{}}
\textit{io} & uscirò \\
\textit{tu} & uscirai \\
\textit{lui/lei} & uscirà \\
\textit{noi} & usciremo \\
\textit{voi} & uscirete \\
\textit{loro} & usciranno \\
\end{tabular}}{}
\end{minipage}\hspace{0.02\linewidth}\begin{minipage}[t]{0.32\linewidth}\raggedright
\paradigmbox{indicativo futuro anteriore (with essere, masculine subject)}{%
\begin{tabular}{@{}l l@{}}
\textit{io} & sarò uscito \\
\textit{tu} & sarai uscito \\
\textit{lui/lei} & sarà uscito \\
\textit{noi} & saremo usciti \\
\textit{voi} & sarete usciti \\
\textit{loro} & saranno usciti \\
\end{tabular}}{}
\end{minipage}\par\medskip

\moodheading{congiuntivo}

\par\noindent\begin{minipage}[t]{0.32\linewidth}\raggedright
\paradigmbox{congiuntivo presente}{%
\begin{tabular}{@{}l l@{}}
\textit{io} & \irr{esca} \\
\textit{tu} & \irr{esca} \\
\textit{lui/lei} & \irr{esca} \\
\textit{noi} & usciamo \\
\textit{voi} & usciate \\
\textit{loro} & \irr{escano} \\
\end{tabular}}{\irr{Highlighted} forms differ from the regular pattern of the verb's family.}
\end{minipage}\hspace{0.02\linewidth}\begin{minipage}[t]{0.32\linewidth}\raggedright
\paradigmbox{congiuntivo passato (with essere, masculine subject)}{%
\begin{tabular}{@{}l l@{}}
\textit{io} & sia uscito \\
\textit{tu} & sia uscito \\
\textit{lui/lei} & sia uscito \\
\textit{noi} & siamo usciti \\
\textit{voi} & siate usciti \\
\textit{loro} & siano usciti \\
\end{tabular}}{}
\end{minipage}\hspace{0.02\linewidth}\begin{minipage}[t]{0.32\linewidth}\raggedright
\paradigmbox{congiuntivo imperfetto}{%
\begin{tabular}{@{}l l@{}}
\textit{io} & uscissi \\
\textit{tu} & uscissi \\
\textit{lui/lei} & uscisse \\
\textit{noi} & uscissimo \\
\textit{voi} & usciste \\
\textit{loro} & uscissero \\
\end{tabular}}{}
\end{minipage}\par\medskip\noindent\begin{minipage}[t]{0.32\linewidth}\raggedright
\paradigmbox{congiuntivo trapassato (with essere, masculine subject)}{%
\begin{tabular}{@{}l l@{}}
\textit{io} & fossi uscito \\
\textit{tu} & fossi uscito \\
\textit{lui/lei} & fosse uscito \\
\textit{noi} & fossimo usciti \\
\textit{voi} & foste usciti \\
\textit{loro} & fossero usciti \\
\end{tabular}}{}
\end{minipage}\par\medskip

\moodheading{condizionale}

\par\noindent\begin{minipage}[t]{0.32\linewidth}\raggedright
\paradigmbox{condizionale presente}{%
\begin{tabular}{@{}l l@{}}
\textit{io} & uscirei \\
\textit{tu} & usciresti \\
\textit{lui/lei} & uscirebbe \\
\textit{noi} & usciremmo \\
\textit{voi} & uscireste \\
\textit{loro} & uscirebbero \\
\end{tabular}}{}
\end{minipage}\hspace{0.02\linewidth}\begin{minipage}[t]{0.32\linewidth}\raggedright
\paradigmbox{condizionale passato (with essere, masculine subject)}{%
\begin{tabular}{@{}l l@{}}
\textit{io} & sarei uscito \\
\textit{tu} & saresti uscito \\
\textit{lui/lei} & sarebbe uscito \\
\textit{noi} & saremmo usciti \\
\textit{voi} & sareste usciti \\
\textit{loro} & sarebbero usciti \\
\end{tabular}}{}
\end{minipage}\par\medskip

\moodheading{imperativo}

\par\noindent\begin{minipage}[t]{0.32\linewidth}\raggedright
\paradigmbox{imperativo}{%
\begin{tabular}{@{}l l@{}}
\textit{tu} & \irr{esci} \\
\textit{Lei} & \irr{esca} \\
\textit{noi} & usciamo \\
\textit{voi} & uscite \\
\textit{Loro} & \irr{escano} \\
\end{tabular}}{\irr{Highlighted} forms differ from the regular pattern of the verb's family. Negative \textit{tu}: \textit{non uscire}. \textit{Lei} and \textit{Loro} are the congiuntivo forms.}
\end{minipage}\par\medskip

\moodheading{forme non finite}
{\small\begin{tabular}{@{}l l@{}}
infinito presente & \textit{uscire} \\
infinito passato & \textit{essere uscito} \\
gerundio presente & \textit{uscendo} \\
gerundio passato & \textit{essendo uscito} \\
participio passato & \textit{uscito} \\
\end{tabular}}\par
\clearpage

\subsection{bere --- to drink}\label{conj:bere}\index{bere@bere}
{\small Auxiliary: \textit{avere}. Short infinitive of an older bevere: once you know the stem bev-, the forms are regular -ere forms.}\par

\moodheading{indicativo}

\par\noindent\begin{minipage}[t]{0.32\linewidth}\raggedright
\paradigmbox{indicativo presente}{%
\begin{tabular}{@{}l l@{}}
\textit{io} & \irr{bevo} \\
\textit{tu} & \irr{bevi} \\
\textit{lui/lei} & \irr{beve} \\
\textit{noi} & \irr{beviamo} \\
\textit{voi} & \irr{bevete} \\
\textit{loro} & \irr{bevono} \\
\end{tabular}}{\irr{Highlighted} forms differ from the regular pattern of the verb's family.}
\end{minipage}\hspace{0.02\linewidth}\begin{minipage}[t]{0.32\linewidth}\raggedright
\paradigmbox{indicativo passato prossimo (with avere)}{%
\begin{tabular}{@{}l l@{}}
\textit{io} & ho bevuto \\
\textit{tu} & hai bevuto \\
\textit{lui/lei} & ha bevuto \\
\textit{noi} & abbiamo bevuto \\
\textit{voi} & avete bevuto \\
\textit{loro} & hanno bevuto \\
\end{tabular}}{}
\end{minipage}\hspace{0.02\linewidth}\begin{minipage}[t]{0.32\linewidth}\raggedright
\paradigmbox{indicativo imperfetto}{%
\begin{tabular}{@{}l l@{}}
\textit{io} & \irr{bevevo} \\
\textit{tu} & \irr{bevevi} \\
\textit{lui/lei} & \irr{beveva} \\
\textit{noi} & \irr{bevevamo} \\
\textit{voi} & \irr{bevevate} \\
\textit{loro} & \irr{bevevano} \\
\end{tabular}}{\irr{Highlighted} forms differ from the regular pattern of the verb's family.}
\end{minipage}\par\medskip\noindent\begin{minipage}[t]{0.32\linewidth}\raggedright
\paradigmbox{indicativo trapassato prossimo (with avere)}{%
\begin{tabular}{@{}l l@{}}
\textit{io} & avevo bevuto \\
\textit{tu} & avevi bevuto \\
\textit{lui/lei} & aveva bevuto \\
\textit{noi} & avevamo bevuto \\
\textit{voi} & avevate bevuto \\
\textit{loro} & avevano bevuto \\
\end{tabular}}{}
\end{minipage}\hspace{0.02\linewidth}\begin{minipage}[t]{0.32\linewidth}\raggedright
\paradigmbox{indicativo passato remoto}{%
\begin{tabular}{@{}l l@{}}
\textit{io} & \irr{bevvi} \\
\textit{tu} & \irr{bevesti} \\
\textit{lui/lei} & \irr{bevve} \\
\textit{noi} & \irr{bevemmo} \\
\textit{voi} & \irr{beveste} \\
\textit{loro} & \irr{bevvero} \\
\end{tabular}}{\irr{Highlighted} forms differ from the regular pattern of the verb's family.}
\end{minipage}\hspace{0.02\linewidth}\begin{minipage}[t]{0.32\linewidth}\raggedright
\paradigmbox{indicativo trapassato remoto (with avere)}{%
\begin{tabular}{@{}l l@{}}
\textit{io} & ebbi bevuto \\
\textit{tu} & avesti bevuto \\
\textit{lui/lei} & ebbe bevuto \\
\textit{noi} & avemmo bevuto \\
\textit{voi} & aveste bevuto \\
\textit{loro} & ebbero bevuto \\
\end{tabular}}{}
\end{minipage}\par\medskip\noindent\begin{minipage}[t]{0.32\linewidth}\raggedright
\paradigmbox{indicativo futuro semplice}{%
\begin{tabular}{@{}l l@{}}
\textit{io} & \irr{berrò} \\
\textit{tu} & \irr{berrai} \\
\textit{lui/lei} & \irr{berrà} \\
\textit{noi} & \irr{berremo} \\
\textit{voi} & \irr{berrete} \\
\textit{loro} & \irr{berranno} \\
\end{tabular}}{\irr{Highlighted} forms differ from the regular pattern of the verb's family.}
\end{minipage}\hspace{0.02\linewidth}\begin{minipage}[t]{0.32\linewidth}\raggedright
\paradigmbox{indicativo futuro anteriore (with avere)}{%
\begin{tabular}{@{}l l@{}}
\textit{io} & avrò bevuto \\
\textit{tu} & avrai bevuto \\
\textit{lui/lei} & avrà bevuto \\
\textit{noi} & avremo bevuto \\
\textit{voi} & avrete bevuto \\
\textit{loro} & avranno bevuto \\
\end{tabular}}{}
\end{minipage}\par\medskip

\moodheading{congiuntivo}

\par\noindent\begin{minipage}[t]{0.32\linewidth}\raggedright
\paradigmbox{congiuntivo presente}{%
\begin{tabular}{@{}l l@{}}
\textit{io} & \irr{beva} \\
\textit{tu} & \irr{beva} \\
\textit{lui/lei} & \irr{beva} \\
\textit{noi} & \irr{beviamo} \\
\textit{voi} & \irr{beviate} \\
\textit{loro} & \irr{bevano} \\
\end{tabular}}{\irr{Highlighted} forms differ from the regular pattern of the verb's family.}
\end{minipage}\hspace{0.02\linewidth}\begin{minipage}[t]{0.32\linewidth}\raggedright
\paradigmbox{congiuntivo passato (with avere)}{%
\begin{tabular}{@{}l l@{}}
\textit{io} & abbia bevuto \\
\textit{tu} & abbia bevuto \\
\textit{lui/lei} & abbia bevuto \\
\textit{noi} & abbiamo bevuto \\
\textit{voi} & abbiate bevuto \\
\textit{loro} & abbiano bevuto \\
\end{tabular}}{}
\end{minipage}\hspace{0.02\linewidth}\begin{minipage}[t]{0.32\linewidth}\raggedright
\paradigmbox{congiuntivo imperfetto}{%
\begin{tabular}{@{}l l@{}}
\textit{io} & \irr{bevessi} \\
\textit{tu} & \irr{bevessi} \\
\textit{lui/lei} & \irr{bevesse} \\
\textit{noi} & \irr{bevessimo} \\
\textit{voi} & \irr{beveste} \\
\textit{loro} & \irr{bevessero} \\
\end{tabular}}{\irr{Highlighted} forms differ from the regular pattern of the verb's family.}
\end{minipage}\par\medskip\noindent\begin{minipage}[t]{0.32\linewidth}\raggedright
\paradigmbox{congiuntivo trapassato (with avere)}{%
\begin{tabular}{@{}l l@{}}
\textit{io} & avessi bevuto \\
\textit{tu} & avessi bevuto \\
\textit{lui/lei} & avesse bevuto \\
\textit{noi} & avessimo bevuto \\
\textit{voi} & aveste bevuto \\
\textit{loro} & avessero bevuto \\
\end{tabular}}{}
\end{minipage}\par\medskip

\moodheading{condizionale}

\par\noindent\begin{minipage}[t]{0.32\linewidth}\raggedright
\paradigmbox{condizionale presente}{%
\begin{tabular}{@{}l l@{}}
\textit{io} & \irr{berrei} \\
\textit{tu} & \irr{berresti} \\
\textit{lui/lei} & \irr{berrebbe} \\
\textit{noi} & \irr{berremmo} \\
\textit{voi} & \irr{berreste} \\
\textit{loro} & \irr{berrebbero} \\
\end{tabular}}{\irr{Highlighted} forms differ from the regular pattern of the verb's family.}
\end{minipage}\hspace{0.02\linewidth}\begin{minipage}[t]{0.32\linewidth}\raggedright
\paradigmbox{condizionale passato (with avere)}{%
\begin{tabular}{@{}l l@{}}
\textit{io} & avrei bevuto \\
\textit{tu} & avresti bevuto \\
\textit{lui/lei} & avrebbe bevuto \\
\textit{noi} & avremmo bevuto \\
\textit{voi} & avreste bevuto \\
\textit{loro} & avrebbero bevuto \\
\end{tabular}}{}
\end{minipage}\par\medskip

\moodheading{imperativo}

\par\noindent\begin{minipage}[t]{0.32\linewidth}\raggedright
\paradigmbox{imperativo}{%
\begin{tabular}{@{}l l@{}}
\textit{tu} & \irr{bevi} \\
\textit{Lei} & \irr{beva} \\
\textit{noi} & \irr{beviamo} \\
\textit{voi} & \irr{bevete} \\
\textit{Loro} & \irr{bevano} \\
\end{tabular}}{\irr{Highlighted} forms differ from the regular pattern of the verb's family. Negative \textit{tu}: \textit{non bere}. \textit{Lei} and \textit{Loro} are the congiuntivo forms.}
\end{minipage}\par\medskip

\moodheading{forme non finite}
{\small\begin{tabular}{@{}l l@{}}
infinito presente & \textit{bere} \\
infinito passato & \textit{avere bevuto} \\
gerundio presente & \textit{bevendo} \\
gerundio passato & \textit{avendo bevuto} \\
participio passato & \textit{bevuto} \\
\end{tabular}}\par
\clearpage

\subsection{potere --- to be able to; can; may}\label{conj:potere}\index{potere@potere}
{\small Auxiliary: \textit{avere}. Modal verb: usually followed by an infinitive (posso venire). In compound tenses a modal often takes the auxiliary of the following verb (sono potuto venire).}\par

\moodheading{indicativo}

\par\noindent\begin{minipage}[t]{0.32\linewidth}\raggedright
\paradigmbox{indicativo presente}{%
\begin{tabular}{@{}l l@{}}
\textit{io} & \irr{posso} \\
\textit{tu} & \irr{puoi} \\
\textit{lui/lei} & \irr{può} \\
\textit{noi} & \irr{possiamo} \\
\textit{voi} & potete \\
\textit{loro} & \irr{possono} \\
\end{tabular}}{\irr{Highlighted} forms differ from the regular pattern of the verb's family.}
\end{minipage}\hspace{0.02\linewidth}\begin{minipage}[t]{0.32\linewidth}\raggedright
\paradigmbox{indicativo passato prossimo (with avere)}{%
\begin{tabular}{@{}l l@{}}
\textit{io} & ho potuto \\
\textit{tu} & hai potuto \\
\textit{lui/lei} & ha potuto \\
\textit{noi} & abbiamo potuto \\
\textit{voi} & avete potuto \\
\textit{loro} & hanno potuto \\
\end{tabular}}{}
\end{minipage}\hspace{0.02\linewidth}\begin{minipage}[t]{0.32\linewidth}\raggedright
\paradigmbox{indicativo imperfetto}{%
\begin{tabular}{@{}l l@{}}
\textit{io} & potevo \\
\textit{tu} & potevi \\
\textit{lui/lei} & poteva \\
\textit{noi} & potevamo \\
\textit{voi} & potevate \\
\textit{loro} & potevano \\
\end{tabular}}{}
\end{minipage}\par\medskip\noindent\begin{minipage}[t]{0.32\linewidth}\raggedright
\paradigmbox{indicativo trapassato prossimo (with avere)}{%
\begin{tabular}{@{}l l@{}}
\textit{io} & avevo potuto \\
\textit{tu} & avevi potuto \\
\textit{lui/lei} & aveva potuto \\
\textit{noi} & avevamo potuto \\
\textit{voi} & avevate potuto \\
\textit{loro} & avevano potuto \\
\end{tabular}}{}
\end{minipage}\hspace{0.02\linewidth}\begin{minipage}[t]{0.32\linewidth}\raggedright
\paradigmbox{indicativo passato remoto}{%
\begin{tabular}{@{}l l@{}}
\textit{io} & potei / potetti \\
\textit{tu} & potesti \\
\textit{lui/lei} & poté / potette \\
\textit{noi} & potemmo \\
\textit{voi} & poteste \\
\textit{loro} & poterono / potettero \\
\end{tabular}}{Where two forms are given, both are standard.}
\end{minipage}\hspace{0.02\linewidth}\begin{minipage}[t]{0.32\linewidth}\raggedright
\paradigmbox{indicativo trapassato remoto (with avere)}{%
\begin{tabular}{@{}l l@{}}
\textit{io} & ebbi potuto \\
\textit{tu} & avesti potuto \\
\textit{lui/lei} & ebbe potuto \\
\textit{noi} & avemmo potuto \\
\textit{voi} & aveste potuto \\
\textit{loro} & ebbero potuto \\
\end{tabular}}{}
\end{minipage}\par\medskip\noindent\begin{minipage}[t]{0.32\linewidth}\raggedright
\paradigmbox{indicativo futuro semplice}{%
\begin{tabular}{@{}l l@{}}
\textit{io} & \irr{potrò} \\
\textit{tu} & \irr{potrai} \\
\textit{lui/lei} & \irr{potrà} \\
\textit{noi} & \irr{potremo} \\
\textit{voi} & \irr{potrete} \\
\textit{loro} & \irr{potranno} \\
\end{tabular}}{\irr{Highlighted} forms differ from the regular pattern of the verb's family.}
\end{minipage}\hspace{0.02\linewidth}\begin{minipage}[t]{0.32\linewidth}\raggedright
\paradigmbox{indicativo futuro anteriore (with avere)}{%
\begin{tabular}{@{}l l@{}}
\textit{io} & avrò potuto \\
\textit{tu} & avrai potuto \\
\textit{lui/lei} & avrà potuto \\
\textit{noi} & avremo potuto \\
\textit{voi} & avrete potuto \\
\textit{loro} & avranno potuto \\
\end{tabular}}{}
\end{minipage}\par\medskip

\moodheading{congiuntivo}

\par\noindent\begin{minipage}[t]{0.32\linewidth}\raggedright
\paradigmbox{congiuntivo presente}{%
\begin{tabular}{@{}l l@{}}
\textit{io} & \irr{possa} \\
\textit{tu} & \irr{possa} \\
\textit{lui/lei} & \irr{possa} \\
\textit{noi} & \irr{possiamo} \\
\textit{voi} & \irr{possiate} \\
\textit{loro} & \irr{possano} \\
\end{tabular}}{\irr{Highlighted} forms differ from the regular pattern of the verb's family.}
\end{minipage}\hspace{0.02\linewidth}\begin{minipage}[t]{0.32\linewidth}\raggedright
\paradigmbox{congiuntivo passato (with avere)}{%
\begin{tabular}{@{}l l@{}}
\textit{io} & abbia potuto \\
\textit{tu} & abbia potuto \\
\textit{lui/lei} & abbia potuto \\
\textit{noi} & abbiamo potuto \\
\textit{voi} & abbiate potuto \\
\textit{loro} & abbiano potuto \\
\end{tabular}}{}
\end{minipage}\hspace{0.02\linewidth}\begin{minipage}[t]{0.32\linewidth}\raggedright
\paradigmbox{congiuntivo imperfetto}{%
\begin{tabular}{@{}l l@{}}
\textit{io} & potessi \\
\textit{tu} & potessi \\
\textit{lui/lei} & potesse \\
\textit{noi} & potessimo \\
\textit{voi} & poteste \\
\textit{loro} & potessero \\
\end{tabular}}{}
\end{minipage}\par\medskip\noindent\begin{minipage}[t]{0.32\linewidth}\raggedright
\paradigmbox{congiuntivo trapassato (with avere)}{%
\begin{tabular}{@{}l l@{}}
\textit{io} & avessi potuto \\
\textit{tu} & avessi potuto \\
\textit{lui/lei} & avesse potuto \\
\textit{noi} & avessimo potuto \\
\textit{voi} & aveste potuto \\
\textit{loro} & avessero potuto \\
\end{tabular}}{}
\end{minipage}\par\medskip

\moodheading{condizionale}

\par\noindent\begin{minipage}[t]{0.32\linewidth}\raggedright
\paradigmbox{condizionale presente}{%
\begin{tabular}{@{}l l@{}}
\textit{io} & \irr{potrei} \\
\textit{tu} & \irr{potresti} \\
\textit{lui/lei} & \irr{potrebbe} \\
\textit{noi} & \irr{potremmo} \\
\textit{voi} & \irr{potreste} \\
\textit{loro} & \irr{potrebbero} \\
\end{tabular}}{\irr{Highlighted} forms differ from the regular pattern of the verb's family.}
\end{minipage}\hspace{0.02\linewidth}\begin{minipage}[t]{0.32\linewidth}\raggedright
\paradigmbox{condizionale passato (with avere)}{%
\begin{tabular}{@{}l l@{}}
\textit{io} & avrei potuto \\
\textit{tu} & avresti potuto \\
\textit{lui/lei} & avrebbe potuto \\
\textit{noi} & avremmo potuto \\
\textit{voi} & avreste potuto \\
\textit{loro} & avrebbero potuto \\
\end{tabular}}{}
\end{minipage}\par\medskip

\moodheading{imperativo}

\par\noindent\begin{minipage}[t]{0.32\linewidth}\raggedright
\paradigmbox{imperativo}{%
\begin{tabular}{@{}l l@{}}
\textit{tu} & --- \\
\textit{Lei} & --- \\
\textit{noi} & --- \\
\textit{voi} & --- \\
\textit{Loro} & --- \\
\end{tabular}}{This verb has no imperative in normal use.}
\end{minipage}\par\medskip

\moodheading{forme non finite}
{\small\begin{tabular}{@{}l l@{}}
infinito presente & \textit{potere} \\
infinito passato & \textit{avere potuto} \\
gerundio presente & \textit{potendo} \\
gerundio passato & \textit{avendo potuto} \\
participio passato & \textit{potuto} \\
\end{tabular}}\par
\clearpage

\subsection{volere --- to want}\label{conj:volere}\index{volere@volere}
{\small Auxiliary: \textit{avere}. Modal verb. Can also take a direct object: voglio un gelato. Polite request: vorrei (conditional).}\par

\moodheading{indicativo}

\par\noindent\begin{minipage}[t]{0.32\linewidth}\raggedright
\paradigmbox{indicativo presente}{%
\begin{tabular}{@{}l l@{}}
\textit{io} & \irr{voglio} \\
\textit{tu} & \irr{vuoi} \\
\textit{lui/lei} & \irr{vuole} \\
\textit{noi} & \irr{vogliamo} \\
\textit{voi} & volete \\
\textit{loro} & \irr{vogliono} \\
\end{tabular}}{\irr{Highlighted} forms differ from the regular pattern of the verb's family.}
\end{minipage}\hspace{0.02\linewidth}\begin{minipage}[t]{0.32\linewidth}\raggedright
\paradigmbox{indicativo passato prossimo (with avere)}{%
\begin{tabular}{@{}l l@{}}
\textit{io} & ho voluto \\
\textit{tu} & hai voluto \\
\textit{lui/lei} & ha voluto \\
\textit{noi} & abbiamo voluto \\
\textit{voi} & avete voluto \\
\textit{loro} & hanno voluto \\
\end{tabular}}{}
\end{minipage}\hspace{0.02\linewidth}\begin{minipage}[t]{0.32\linewidth}\raggedright
\paradigmbox{indicativo imperfetto}{%
\begin{tabular}{@{}l l@{}}
\textit{io} & volevo \\
\textit{tu} & volevi \\
\textit{lui/lei} & voleva \\
\textit{noi} & volevamo \\
\textit{voi} & volevate \\
\textit{loro} & volevano \\
\end{tabular}}{}
\end{minipage}\par\medskip\noindent\begin{minipage}[t]{0.32\linewidth}\raggedright
\paradigmbox{indicativo trapassato prossimo (with avere)}{%
\begin{tabular}{@{}l l@{}}
\textit{io} & avevo voluto \\
\textit{tu} & avevi voluto \\
\textit{lui/lei} & aveva voluto \\
\textit{noi} & avevamo voluto \\
\textit{voi} & avevate voluto \\
\textit{loro} & avevano voluto \\
\end{tabular}}{}
\end{minipage}\hspace{0.02\linewidth}\begin{minipage}[t]{0.32\linewidth}\raggedright
\paradigmbox{indicativo passato remoto}{%
\begin{tabular}{@{}l l@{}}
\textit{io} & \irr{volli} \\
\textit{tu} & volesti \\
\textit{lui/lei} & \irr{volle} \\
\textit{noi} & volemmo \\
\textit{voi} & voleste \\
\textit{loro} & \irr{vollero} \\
\end{tabular}}{\irr{Highlighted} forms differ from the regular pattern of the verb's family.}
\end{minipage}\hspace{0.02\linewidth}\begin{minipage}[t]{0.32\linewidth}\raggedright
\paradigmbox{indicativo trapassato remoto (with avere)}{%
\begin{tabular}{@{}l l@{}}
\textit{io} & ebbi voluto \\
\textit{tu} & avesti voluto \\
\textit{lui/lei} & ebbe voluto \\
\textit{noi} & avemmo voluto \\
\textit{voi} & aveste voluto \\
\textit{loro} & ebbero voluto \\
\end{tabular}}{}
\end{minipage}\par\medskip\noindent\begin{minipage}[t]{0.32\linewidth}\raggedright
\paradigmbox{indicativo futuro semplice}{%
\begin{tabular}{@{}l l@{}}
\textit{io} & \irr{vorrò} \\
\textit{tu} & \irr{vorrai} \\
\textit{lui/lei} & \irr{vorrà} \\
\textit{noi} & \irr{vorremo} \\
\textit{voi} & \irr{vorrete} \\
\textit{loro} & \irr{vorranno} \\
\end{tabular}}{\irr{Highlighted} forms differ from the regular pattern of the verb's family.}
\end{minipage}\hspace{0.02\linewidth}\begin{minipage}[t]{0.32\linewidth}\raggedright
\paradigmbox{indicativo futuro anteriore (with avere)}{%
\begin{tabular}{@{}l l@{}}
\textit{io} & avrò voluto \\
\textit{tu} & avrai voluto \\
\textit{lui/lei} & avrà voluto \\
\textit{noi} & avremo voluto \\
\textit{voi} & avrete voluto \\
\textit{loro} & avranno voluto \\
\end{tabular}}{}
\end{minipage}\par\medskip

\moodheading{congiuntivo}

\par\noindent\begin{minipage}[t]{0.32\linewidth}\raggedright
\paradigmbox{congiuntivo presente}{%
\begin{tabular}{@{}l l@{}}
\textit{io} & \irr{voglia} \\
\textit{tu} & \irr{voglia} \\
\textit{lui/lei} & \irr{voglia} \\
\textit{noi} & \irr{vogliamo} \\
\textit{voi} & \irr{vogliate} \\
\textit{loro} & \irr{vogliano} \\
\end{tabular}}{\irr{Highlighted} forms differ from the regular pattern of the verb's family.}
\end{minipage}\hspace{0.02\linewidth}\begin{minipage}[t]{0.32\linewidth}\raggedright
\paradigmbox{congiuntivo passato (with avere)}{%
\begin{tabular}{@{}l l@{}}
\textit{io} & abbia voluto \\
\textit{tu} & abbia voluto \\
\textit{lui/lei} & abbia voluto \\
\textit{noi} & abbiamo voluto \\
\textit{voi} & abbiate voluto \\
\textit{loro} & abbiano voluto \\
\end{tabular}}{}
\end{minipage}\hspace{0.02\linewidth}\begin{minipage}[t]{0.32\linewidth}\raggedright
\paradigmbox{congiuntivo imperfetto}{%
\begin{tabular}{@{}l l@{}}
\textit{io} & volessi \\
\textit{tu} & volessi \\
\textit{lui/lei} & volesse \\
\textit{noi} & volessimo \\
\textit{voi} & voleste \\
\textit{loro} & volessero \\
\end{tabular}}{}
\end{minipage}\par\medskip\noindent\begin{minipage}[t]{0.32\linewidth}\raggedright
\paradigmbox{congiuntivo trapassato (with avere)}{%
\begin{tabular}{@{}l l@{}}
\textit{io} & avessi voluto \\
\textit{tu} & avessi voluto \\
\textit{lui/lei} & avesse voluto \\
\textit{noi} & avessimo voluto \\
\textit{voi} & aveste voluto \\
\textit{loro} & avessero voluto \\
\end{tabular}}{}
\end{minipage}\par\medskip

\moodheading{condizionale}

\par\noindent\begin{minipage}[t]{0.32\linewidth}\raggedright
\paradigmbox{condizionale presente}{%
\begin{tabular}{@{}l l@{}}
\textit{io} & \irr{vorrei} \\
\textit{tu} & \irr{vorresti} \\
\textit{lui/lei} & \irr{vorrebbe} \\
\textit{noi} & \irr{vorremmo} \\
\textit{voi} & \irr{vorreste} \\
\textit{loro} & \irr{vorrebbero} \\
\end{tabular}}{\irr{Highlighted} forms differ from the regular pattern of the verb's family.}
\end{minipage}\hspace{0.02\linewidth}\begin{minipage}[t]{0.32\linewidth}\raggedright
\paradigmbox{condizionale passato (with avere)}{%
\begin{tabular}{@{}l l@{}}
\textit{io} & avrei voluto \\
\textit{tu} & avresti voluto \\
\textit{lui/lei} & avrebbe voluto \\
\textit{noi} & avremmo voluto \\
\textit{voi} & avreste voluto \\
\textit{loro} & avrebbero voluto \\
\end{tabular}}{}
\end{minipage}\par\medskip

\moodheading{imperativo}

\par\noindent\begin{minipage}[t]{0.32\linewidth}\raggedright
\paradigmbox{imperativo}{%
\begin{tabular}{@{}l l@{}}
\textit{tu} & \irr{vogli} \\
\textit{Lei} & \irr{voglia} \\
\textit{noi} & \irr{vogliamo} \\
\textit{voi} & \irr{vogliate} \\
\textit{Loro} & \irr{vogliano} \\
\end{tabular}}{\irr{Highlighted} forms differ from the regular pattern of the verb's family. Negative \textit{tu}: \textit{non volere}. \textit{Lei} and \textit{Loro} are the congiuntivo forms.}
\end{minipage}\par\medskip

\moodheading{forme non finite}
{\small\begin{tabular}{@{}l l@{}}
infinito presente & \textit{volere} \\
infinito passato & \textit{avere voluto} \\
gerundio presente & \textit{volendo} \\
gerundio passato & \textit{avendo voluto} \\
participio passato & \textit{voluto} \\
\end{tabular}}\par
\clearpage

\subsection{dovere --- to have to; must}\label{conj:dovere}\index{dovere@dovere}
{\small Auxiliary: \textit{avere}. Modal verb: dev- in the stressed forms, dobb- in noi, dov- in voi. Alternative forms debbo, debbono exist but are rare today.}\par

\moodheading{indicativo}

\par\noindent\begin{minipage}[t]{0.32\linewidth}\raggedright
\paradigmbox{indicativo presente}{%
\begin{tabular}{@{}l l@{}}
\textit{io} & \irr{devo} \\
\textit{tu} & \irr{devi} \\
\textit{lui/lei} & \irr{deve} \\
\textit{noi} & \irr{dobbiamo} \\
\textit{voi} & dovete \\
\textit{loro} & \irr{devono} \\
\end{tabular}}{\irr{Highlighted} forms differ from the regular pattern of the verb's family.}
\end{minipage}\hspace{0.02\linewidth}\begin{minipage}[t]{0.32\linewidth}\raggedright
\paradigmbox{indicativo passato prossimo (with avere)}{%
\begin{tabular}{@{}l l@{}}
\textit{io} & ho dovuto \\
\textit{tu} & hai dovuto \\
\textit{lui/lei} & ha dovuto \\
\textit{noi} & abbiamo dovuto \\
\textit{voi} & avete dovuto \\
\textit{loro} & hanno dovuto \\
\end{tabular}}{}
\end{minipage}\hspace{0.02\linewidth}\begin{minipage}[t]{0.32\linewidth}\raggedright
\paradigmbox{indicativo imperfetto}{%
\begin{tabular}{@{}l l@{}}
\textit{io} & dovevo \\
\textit{tu} & dovevi \\
\textit{lui/lei} & doveva \\
\textit{noi} & dovevamo \\
\textit{voi} & dovevate \\
\textit{loro} & dovevano \\
\end{tabular}}{}
\end{minipage}\par\medskip\noindent\begin{minipage}[t]{0.32\linewidth}\raggedright
\paradigmbox{indicativo trapassato prossimo (with avere)}{%
\begin{tabular}{@{}l l@{}}
\textit{io} & avevo dovuto \\
\textit{tu} & avevi dovuto \\
\textit{lui/lei} & aveva dovuto \\
\textit{noi} & avevamo dovuto \\
\textit{voi} & avevate dovuto \\
\textit{loro} & avevano dovuto \\
\end{tabular}}{}
\end{minipage}\hspace{0.02\linewidth}\begin{minipage}[t]{0.32\linewidth}\raggedright
\paradigmbox{indicativo passato remoto}{%
\begin{tabular}{@{}l l@{}}
\textit{io} & dovei / dovetti \\
\textit{tu} & dovesti \\
\textit{lui/lei} & dové / dovette \\
\textit{noi} & dovemmo \\
\textit{voi} & doveste \\
\textit{loro} & doverono / dovettero \\
\end{tabular}}{Where two forms are given, both are standard.}
\end{minipage}\hspace{0.02\linewidth}\begin{minipage}[t]{0.32\linewidth}\raggedright
\paradigmbox{indicativo trapassato remoto (with avere)}{%
\begin{tabular}{@{}l l@{}}
\textit{io} & ebbi dovuto \\
\textit{tu} & avesti dovuto \\
\textit{lui/lei} & ebbe dovuto \\
\textit{noi} & avemmo dovuto \\
\textit{voi} & aveste dovuto \\
\textit{loro} & ebbero dovuto \\
\end{tabular}}{}
\end{minipage}\par\medskip\noindent\begin{minipage}[t]{0.32\linewidth}\raggedright
\paradigmbox{indicativo futuro semplice}{%
\begin{tabular}{@{}l l@{}}
\textit{io} & \irr{dovrò} \\
\textit{tu} & \irr{dovrai} \\
\textit{lui/lei} & \irr{dovrà} \\
\textit{noi} & \irr{dovremo} \\
\textit{voi} & \irr{dovrete} \\
\textit{loro} & \irr{dovranno} \\
\end{tabular}}{\irr{Highlighted} forms differ from the regular pattern of the verb's family.}
\end{minipage}\hspace{0.02\linewidth}\begin{minipage}[t]{0.32\linewidth}\raggedright
\paradigmbox{indicativo futuro anteriore (with avere)}{%
\begin{tabular}{@{}l l@{}}
\textit{io} & avrò dovuto \\
\textit{tu} & avrai dovuto \\
\textit{lui/lei} & avrà dovuto \\
\textit{noi} & avremo dovuto \\
\textit{voi} & avrete dovuto \\
\textit{loro} & avranno dovuto \\
\end{tabular}}{}
\end{minipage}\par\medskip

\moodheading{congiuntivo}

\par\noindent\begin{minipage}[t]{0.32\linewidth}\raggedright
\paradigmbox{congiuntivo presente}{%
\begin{tabular}{@{}l l@{}}
\textit{io} & \irr{debba / deva} \\
\textit{tu} & \irr{debba / deva} \\
\textit{lui/lei} & \irr{debba / deva} \\
\textit{noi} & \irr{dobbiamo} \\
\textit{voi} & \irr{dobbiate} \\
\textit{loro} & \irr{debbano / devano} \\
\end{tabular}}{\irr{Highlighted} forms differ from the regular pattern of the verb's family. Where two forms are given, both are standard.}
\end{minipage}\hspace{0.02\linewidth}\begin{minipage}[t]{0.32\linewidth}\raggedright
\paradigmbox{congiuntivo passato (with avere)}{%
\begin{tabular}{@{}l l@{}}
\textit{io} & abbia dovuto \\
\textit{tu} & abbia dovuto \\
\textit{lui/lei} & abbia dovuto \\
\textit{noi} & abbiamo dovuto \\
\textit{voi} & abbiate dovuto \\
\textit{loro} & abbiano dovuto \\
\end{tabular}}{}
\end{minipage}\hspace{0.02\linewidth}\begin{minipage}[t]{0.32\linewidth}\raggedright
\paradigmbox{congiuntivo imperfetto}{%
\begin{tabular}{@{}l l@{}}
\textit{io} & dovessi \\
\textit{tu} & dovessi \\
\textit{lui/lei} & dovesse \\
\textit{noi} & dovessimo \\
\textit{voi} & doveste \\
\textit{loro} & dovessero \\
\end{tabular}}{}
\end{minipage}\par\medskip\noindent\begin{minipage}[t]{0.32\linewidth}\raggedright
\paradigmbox{congiuntivo trapassato (with avere)}{%
\begin{tabular}{@{}l l@{}}
\textit{io} & avessi dovuto \\
\textit{tu} & avessi dovuto \\
\textit{lui/lei} & avesse dovuto \\
\textit{noi} & avessimo dovuto \\
\textit{voi} & aveste dovuto \\
\textit{loro} & avessero dovuto \\
\end{tabular}}{}
\end{minipage}\par\medskip

\moodheading{condizionale}

\par\noindent\begin{minipage}[t]{0.32\linewidth}\raggedright
\paradigmbox{condizionale presente}{%
\begin{tabular}{@{}l l@{}}
\textit{io} & \irr{dovrei} \\
\textit{tu} & \irr{dovresti} \\
\textit{lui/lei} & \irr{dovrebbe} \\
\textit{noi} & \irr{dovremmo} \\
\textit{voi} & \irr{dovreste} \\
\textit{loro} & \irr{dovrebbero} \\
\end{tabular}}{\irr{Highlighted} forms differ from the regular pattern of the verb's family.}
\end{minipage}\hspace{0.02\linewidth}\begin{minipage}[t]{0.32\linewidth}\raggedright
\paradigmbox{condizionale passato (with avere)}{%
\begin{tabular}{@{}l l@{}}
\textit{io} & avrei dovuto \\
\textit{tu} & avresti dovuto \\
\textit{lui/lei} & avrebbe dovuto \\
\textit{noi} & avremmo dovuto \\
\textit{voi} & avreste dovuto \\
\textit{loro} & avrebbero dovuto \\
\end{tabular}}{}
\end{minipage}\par\medskip

\moodheading{imperativo}

\par\noindent\begin{minipage}[t]{0.32\linewidth}\raggedright
\paradigmbox{imperativo}{%
\begin{tabular}{@{}l l@{}}
\textit{tu} & --- \\
\textit{Lei} & --- \\
\textit{noi} & --- \\
\textit{voi} & --- \\
\textit{Loro} & --- \\
\end{tabular}}{This verb has no imperative in normal use.}
\end{minipage}\par\medskip

\moodheading{forme non finite}
{\small\begin{tabular}{@{}l l@{}}
infinito presente & \textit{dovere} \\
infinito passato & \textit{avere dovuto} \\
gerundio presente & \textit{dovendo} \\
gerundio passato & \textit{avendo dovuto} \\
participio passato & \textit{dovuto} \\
\end{tabular}}\par
\clearpage

\subsection{sapere --- to know (a fact); to know how to}\label{conj:sapere}\index{sapere@sapere}
{\small Auxiliary: \textit{avere}. Means 'to know' (a fact) and, with an infinitive, 'to know how to': so nuotare = I can (know how to) swim.}\par

\moodheading{indicativo}

\par\noindent\begin{minipage}[t]{0.32\linewidth}\raggedright
\paradigmbox{indicativo presente}{%
\begin{tabular}{@{}l l@{}}
\textit{io} & \irr{so} \\
\textit{tu} & \irr{sai} \\
\textit{lui/lei} & \irr{sa} \\
\textit{noi} & \irr{sappiamo} \\
\textit{voi} & sapete \\
\textit{loro} & \irr{sanno} \\
\end{tabular}}{\irr{Highlighted} forms differ from the regular pattern of the verb's family.}
\end{minipage}\hspace{0.02\linewidth}\begin{minipage}[t]{0.32\linewidth}\raggedright
\paradigmbox{indicativo passato prossimo (with avere)}{%
\begin{tabular}{@{}l l@{}}
\textit{io} & ho saputo \\
\textit{tu} & hai saputo \\
\textit{lui/lei} & ha saputo \\
\textit{noi} & abbiamo saputo \\
\textit{voi} & avete saputo \\
\textit{loro} & hanno saputo \\
\end{tabular}}{}
\end{minipage}\hspace{0.02\linewidth}\begin{minipage}[t]{0.32\linewidth}\raggedright
\paradigmbox{indicativo imperfetto}{%
\begin{tabular}{@{}l l@{}}
\textit{io} & sapevo \\
\textit{tu} & sapevi \\
\textit{lui/lei} & sapeva \\
\textit{noi} & sapevamo \\
\textit{voi} & sapevate \\
\textit{loro} & sapevano \\
\end{tabular}}{}
\end{minipage}\par\medskip\noindent\begin{minipage}[t]{0.32\linewidth}\raggedright
\paradigmbox{indicativo trapassato prossimo (with avere)}{%
\begin{tabular}{@{}l l@{}}
\textit{io} & avevo saputo \\
\textit{tu} & avevi saputo \\
\textit{lui/lei} & aveva saputo \\
\textit{noi} & avevamo saputo \\
\textit{voi} & avevate saputo \\
\textit{loro} & avevano saputo \\
\end{tabular}}{}
\end{minipage}\hspace{0.02\linewidth}\begin{minipage}[t]{0.32\linewidth}\raggedright
\paradigmbox{indicativo passato remoto}{%
\begin{tabular}{@{}l l@{}}
\textit{io} & \irr{seppi} \\
\textit{tu} & sapesti \\
\textit{lui/lei} & \irr{seppe} \\
\textit{noi} & sapemmo \\
\textit{voi} & sapeste \\
\textit{loro} & \irr{seppero} \\
\end{tabular}}{\irr{Highlighted} forms differ from the regular pattern of the verb's family.}
\end{minipage}\hspace{0.02\linewidth}\begin{minipage}[t]{0.32\linewidth}\raggedright
\paradigmbox{indicativo trapassato remoto (with avere)}{%
\begin{tabular}{@{}l l@{}}
\textit{io} & ebbi saputo \\
\textit{tu} & avesti saputo \\
\textit{lui/lei} & ebbe saputo \\
\textit{noi} & avemmo saputo \\
\textit{voi} & aveste saputo \\
\textit{loro} & ebbero saputo \\
\end{tabular}}{}
\end{minipage}\par\medskip\noindent\begin{minipage}[t]{0.32\linewidth}\raggedright
\paradigmbox{indicativo futuro semplice}{%
\begin{tabular}{@{}l l@{}}
\textit{io} & \irr{saprò} \\
\textit{tu} & \irr{saprai} \\
\textit{lui/lei} & \irr{saprà} \\
\textit{noi} & \irr{sapremo} \\
\textit{voi} & \irr{saprete} \\
\textit{loro} & \irr{sapranno} \\
\end{tabular}}{\irr{Highlighted} forms differ from the regular pattern of the verb's family.}
\end{minipage}\hspace{0.02\linewidth}\begin{minipage}[t]{0.32\linewidth}\raggedright
\paradigmbox{indicativo futuro anteriore (with avere)}{%
\begin{tabular}{@{}l l@{}}
\textit{io} & avrò saputo \\
\textit{tu} & avrai saputo \\
\textit{lui/lei} & avrà saputo \\
\textit{noi} & avremo saputo \\
\textit{voi} & avrete saputo \\
\textit{loro} & avranno saputo \\
\end{tabular}}{}
\end{minipage}\par\medskip

\moodheading{congiuntivo}

\par\noindent\begin{minipage}[t]{0.32\linewidth}\raggedright
\paradigmbox{congiuntivo presente}{%
\begin{tabular}{@{}l l@{}}
\textit{io} & \irr{sappia} \\
\textit{tu} & \irr{sappia} \\
\textit{lui/lei} & \irr{sappia} \\
\textit{noi} & \irr{sappiamo} \\
\textit{voi} & \irr{sappiate} \\
\textit{loro} & \irr{sappiano} \\
\end{tabular}}{\irr{Highlighted} forms differ from the regular pattern of the verb's family.}
\end{minipage}\hspace{0.02\linewidth}\begin{minipage}[t]{0.32\linewidth}\raggedright
\paradigmbox{congiuntivo passato (with avere)}{%
\begin{tabular}{@{}l l@{}}
\textit{io} & abbia saputo \\
\textit{tu} & abbia saputo \\
\textit{lui/lei} & abbia saputo \\
\textit{noi} & abbiamo saputo \\
\textit{voi} & abbiate saputo \\
\textit{loro} & abbiano saputo \\
\end{tabular}}{}
\end{minipage}\hspace{0.02\linewidth}\begin{minipage}[t]{0.32\linewidth}\raggedright
\paradigmbox{congiuntivo imperfetto}{%
\begin{tabular}{@{}l l@{}}
\textit{io} & sapessi \\
\textit{tu} & sapessi \\
\textit{lui/lei} & sapesse \\
\textit{noi} & sapessimo \\
\textit{voi} & sapeste \\
\textit{loro} & sapessero \\
\end{tabular}}{}
\end{minipage}\par\medskip\noindent\begin{minipage}[t]{0.32\linewidth}\raggedright
\paradigmbox{congiuntivo trapassato (with avere)}{%
\begin{tabular}{@{}l l@{}}
\textit{io} & avessi saputo \\
\textit{tu} & avessi saputo \\
\textit{lui/lei} & avesse saputo \\
\textit{noi} & avessimo saputo \\
\textit{voi} & aveste saputo \\
\textit{loro} & avessero saputo \\
\end{tabular}}{}
\end{minipage}\par\medskip

\moodheading{condizionale}

\par\noindent\begin{minipage}[t]{0.32\linewidth}\raggedright
\paradigmbox{condizionale presente}{%
\begin{tabular}{@{}l l@{}}
\textit{io} & \irr{saprei} \\
\textit{tu} & \irr{sapresti} \\
\textit{lui/lei} & \irr{saprebbe} \\
\textit{noi} & \irr{sapremmo} \\
\textit{voi} & \irr{sapreste} \\
\textit{loro} & \irr{saprebbero} \\
\end{tabular}}{\irr{Highlighted} forms differ from the regular pattern of the verb's family.}
\end{minipage}\hspace{0.02\linewidth}\begin{minipage}[t]{0.32\linewidth}\raggedright
\paradigmbox{condizionale passato (with avere)}{%
\begin{tabular}{@{}l l@{}}
\textit{io} & avrei saputo \\
\textit{tu} & avresti saputo \\
\textit{lui/lei} & avrebbe saputo \\
\textit{noi} & avremmo saputo \\
\textit{voi} & avreste saputo \\
\textit{loro} & avrebbero saputo \\
\end{tabular}}{}
\end{minipage}\par\medskip

\moodheading{imperativo}

\par\noindent\begin{minipage}[t]{0.32\linewidth}\raggedright
\paradigmbox{imperativo}{%
\begin{tabular}{@{}l l@{}}
\textit{tu} & \irr{sappi} \\
\textit{Lei} & \irr{sappia} \\
\textit{noi} & \irr{sappiamo} \\
\textit{voi} & \irr{sappiate} \\
\textit{Loro} & \irr{sappiano} \\
\end{tabular}}{\irr{Highlighted} forms differ from the regular pattern of the verb's family. Negative \textit{tu}: \textit{non sapere}. \textit{Lei} and \textit{Loro} are the congiuntivo forms.}
\end{minipage}\par\medskip

\moodheading{forme non finite}
{\small\begin{tabular}{@{}l l@{}}
infinito presente & \textit{sapere} \\
infinito passato & \textit{avere saputo} \\
gerundio presente & \textit{sapendo} \\
gerundio passato & \textit{avendo saputo} \\
participio passato & \textit{saputo} \\
\end{tabular}}\par
\clearpage

\subsection{rimanere --- to remain; to stay}\label{conj:rimanere}\index{rimanere@rimanere}
{\small Auxiliary: \textit{essere}. A g appears in io and loro (rimango, rimangono). Irregular participle rimasto.}\par

\moodheading{indicativo}

\par\noindent\begin{minipage}[t]{0.32\linewidth}\raggedright
\paradigmbox{indicativo presente}{%
\begin{tabular}{@{}l l@{}}
\textit{io} & \irr{rimango} \\
\textit{tu} & rimani \\
\textit{lui/lei} & rimane \\
\textit{noi} & rimaniamo \\
\textit{voi} & rimanete \\
\textit{loro} & \irr{rimangono} \\
\end{tabular}}{\irr{Highlighted} forms differ from the regular pattern of the verb's family.}
\end{minipage}\hspace{0.02\linewidth}\begin{minipage}[t]{0.32\linewidth}\raggedright
\paradigmbox{indicativo passato prossimo (with essere, masculine subject)}{%
\begin{tabular}{@{}l l@{}}
\textit{io} & sono rimasto \\
\textit{tu} & sei rimasto \\
\textit{lui/lei} & è rimasto \\
\textit{noi} & siamo rimasti \\
\textit{voi} & siete rimasti \\
\textit{loro} & sono rimasti \\
\end{tabular}}{}
\end{minipage}\hspace{0.02\linewidth}\begin{minipage}[t]{0.32\linewidth}\raggedright
\paradigmbox{indicativo imperfetto}{%
\begin{tabular}{@{}l l@{}}
\textit{io} & rimanevo \\
\textit{tu} & rimanevi \\
\textit{lui/lei} & rimaneva \\
\textit{noi} & rimanevamo \\
\textit{voi} & rimanevate \\
\textit{loro} & rimanevano \\
\end{tabular}}{}
\end{minipage}\par\medskip\noindent\begin{minipage}[t]{0.32\linewidth}\raggedright
\paradigmbox{indicativo trapassato prossimo (with essere, masculine subject)}{%
\begin{tabular}{@{}l l@{}}
\textit{io} & ero rimasto \\
\textit{tu} & eri rimasto \\
\textit{lui/lei} & era rimasto \\
\textit{noi} & eravamo rimasti \\
\textit{voi} & eravate rimasti \\
\textit{loro} & erano rimasti \\
\end{tabular}}{}
\end{minipage}\hspace{0.02\linewidth}\begin{minipage}[t]{0.32\linewidth}\raggedright
\paradigmbox{indicativo passato remoto}{%
\begin{tabular}{@{}l l@{}}
\textit{io} & \irr{rimasi} \\
\textit{tu} & rimanesti \\
\textit{lui/lei} & \irr{rimase} \\
\textit{noi} & rimanemmo \\
\textit{voi} & rimaneste \\
\textit{loro} & \irr{rimasero} \\
\end{tabular}}{\irr{Highlighted} forms differ from the regular pattern of the verb's family.}
\end{minipage}\hspace{0.02\linewidth}\begin{minipage}[t]{0.32\linewidth}\raggedright
\paradigmbox{indicativo trapassato remoto (with essere, masculine subject)}{%
\begin{tabular}{@{}l l@{}}
\textit{io} & fui rimasto \\
\textit{tu} & fosti rimasto \\
\textit{lui/lei} & fu rimasto \\
\textit{noi} & fummo rimasti \\
\textit{voi} & foste rimasti \\
\textit{loro} & furono rimasti \\
\end{tabular}}{}
\end{minipage}\par\medskip\noindent\begin{minipage}[t]{0.32\linewidth}\raggedright
\paradigmbox{indicativo futuro semplice}{%
\begin{tabular}{@{}l l@{}}
\textit{io} & \irr{rimarrò} \\
\textit{tu} & \irr{rimarrai} \\
\textit{lui/lei} & \irr{rimarrà} \\
\textit{noi} & \irr{rimarremo} \\
\textit{voi} & \irr{rimarrete} \\
\textit{loro} & \irr{rimarranno} \\
\end{tabular}}{\irr{Highlighted} forms differ from the regular pattern of the verb's family.}
\end{minipage}\hspace{0.02\linewidth}\begin{minipage}[t]{0.32\linewidth}\raggedright
\paradigmbox{indicativo futuro anteriore (with essere, masculine subject)}{%
\begin{tabular}{@{}l l@{}}
\textit{io} & sarò rimasto \\
\textit{tu} & sarai rimasto \\
\textit{lui/lei} & sarà rimasto \\
\textit{noi} & saremo rimasti \\
\textit{voi} & sarete rimasti \\
\textit{loro} & saranno rimasti \\
\end{tabular}}{}
\end{minipage}\par\medskip

\moodheading{congiuntivo}

\par\noindent\begin{minipage}[t]{0.32\linewidth}\raggedright
\paradigmbox{congiuntivo presente}{%
\begin{tabular}{@{}l l@{}}
\textit{io} & \irr{rimanga} \\
\textit{tu} & \irr{rimanga} \\
\textit{lui/lei} & \irr{rimanga} \\
\textit{noi} & rimaniamo \\
\textit{voi} & rimaniate \\
\textit{loro} & \irr{rimangano} \\
\end{tabular}}{\irr{Highlighted} forms differ from the regular pattern of the verb's family.}
\end{minipage}\hspace{0.02\linewidth}\begin{minipage}[t]{0.32\linewidth}\raggedright
\paradigmbox{congiuntivo passato (with essere, masculine subject)}{%
\begin{tabular}{@{}l l@{}}
\textit{io} & sia rimasto \\
\textit{tu} & sia rimasto \\
\textit{lui/lei} & sia rimasto \\
\textit{noi} & siamo rimasti \\
\textit{voi} & siate rimasti \\
\textit{loro} & siano rimasti \\
\end{tabular}}{}
\end{minipage}\hspace{0.02\linewidth}\begin{minipage}[t]{0.32\linewidth}\raggedright
\paradigmbox{congiuntivo imperfetto}{%
\begin{tabular}{@{}l l@{}}
\textit{io} & rimanessi \\
\textit{tu} & rimanessi \\
\textit{lui/lei} & rimanesse \\
\textit{noi} & rimanessimo \\
\textit{voi} & rimaneste \\
\textit{loro} & rimanessero \\
\end{tabular}}{}
\end{minipage}\par\medskip\noindent\begin{minipage}[t]{0.32\linewidth}\raggedright
\paradigmbox{congiuntivo trapassato (with essere, masculine subject)}{%
\begin{tabular}{@{}l l@{}}
\textit{io} & fossi rimasto \\
\textit{tu} & fossi rimasto \\
\textit{lui/lei} & fosse rimasto \\
\textit{noi} & fossimo rimasti \\
\textit{voi} & foste rimasti \\
\textit{loro} & fossero rimasti \\
\end{tabular}}{}
\end{minipage}\par\medskip

\moodheading{condizionale}

\par\noindent\begin{minipage}[t]{0.32\linewidth}\raggedright
\paradigmbox{condizionale presente}{%
\begin{tabular}{@{}l l@{}}
\textit{io} & \irr{rimarrei} \\
\textit{tu} & \irr{rimarresti} \\
\textit{lui/lei} & \irr{rimarrebbe} \\
\textit{noi} & \irr{rimarremmo} \\
\textit{voi} & \irr{rimarreste} \\
\textit{loro} & \irr{rimarrebbero} \\
\end{tabular}}{\irr{Highlighted} forms differ from the regular pattern of the verb's family.}
\end{minipage}\hspace{0.02\linewidth}\begin{minipage}[t]{0.32\linewidth}\raggedright
\paradigmbox{condizionale passato (with essere, masculine subject)}{%
\begin{tabular}{@{}l l@{}}
\textit{io} & sarei rimasto \\
\textit{tu} & saresti rimasto \\
\textit{lui/lei} & sarebbe rimasto \\
\textit{noi} & saremmo rimasti \\
\textit{voi} & sareste rimasti \\
\textit{loro} & sarebbero rimasti \\
\end{tabular}}{}
\end{minipage}\par\medskip

\moodheading{imperativo}

\par\noindent\begin{minipage}[t]{0.32\linewidth}\raggedright
\paradigmbox{imperativo}{%
\begin{tabular}{@{}l l@{}}
\textit{tu} & rimani \\
\textit{Lei} & \irr{rimanga} \\
\textit{noi} & rimaniamo \\
\textit{voi} & rimanete \\
\textit{Loro} & \irr{rimangano} \\
\end{tabular}}{\irr{Highlighted} forms differ from the regular pattern of the verb's family. Negative \textit{tu}: \textit{non rimanere}. \textit{Lei} and \textit{Loro} are the congiuntivo forms.}
\end{minipage}\par\medskip

\moodheading{forme non finite}
{\small\begin{tabular}{@{}l l@{}}
infinito presente & \textit{rimanere} \\
infinito passato & \textit{essere rimasto} \\
gerundio presente & \textit{rimanendo} \\
gerundio passato & \textit{essendo rimasto} \\
participio passato & \textit{rimasto} \\
\end{tabular}}\par
\clearpage

\subsection{salire --- to go up; to climb; to get on}\label{conj:salire}\index{salire@salire}
{\small Auxiliary: \textit{essere}. A g appears in io and loro (salgo, salgono). Auxiliary essere when it means 'to go up' (sono salito sull'albero); avere with a direct object (ho salito le scale).}\par

\moodheading{indicativo}

\par\noindent\begin{minipage}[t]{0.32\linewidth}\raggedright
\paradigmbox{indicativo presente}{%
\begin{tabular}{@{}l l@{}}
\textit{io} & \irr{salgo} \\
\textit{tu} & sali \\
\textit{lui/lei} & sale \\
\textit{noi} & saliamo \\
\textit{voi} & salite \\
\textit{loro} & \irr{salgono} \\
\end{tabular}}{\irr{Highlighted} forms differ from the regular pattern of the verb's family.}
\end{minipage}\hspace{0.02\linewidth}\begin{minipage}[t]{0.32\linewidth}\raggedright
\paradigmbox{indicativo passato prossimo (with essere, masculine subject)}{%
\begin{tabular}{@{}l l@{}}
\textit{io} & sono salito \\
\textit{tu} & sei salito \\
\textit{lui/lei} & è salito \\
\textit{noi} & siamo saliti \\
\textit{voi} & siete saliti \\
\textit{loro} & sono saliti \\
\end{tabular}}{}
\end{minipage}\hspace{0.02\linewidth}\begin{minipage}[t]{0.32\linewidth}\raggedright
\paradigmbox{indicativo imperfetto}{%
\begin{tabular}{@{}l l@{}}
\textit{io} & salivo \\
\textit{tu} & salivi \\
\textit{lui/lei} & saliva \\
\textit{noi} & salivamo \\
\textit{voi} & salivate \\
\textit{loro} & salivano \\
\end{tabular}}{}
\end{minipage}\par\medskip\noindent\begin{minipage}[t]{0.32\linewidth}\raggedright
\paradigmbox{indicativo trapassato prossimo (with essere, masculine subject)}{%
\begin{tabular}{@{}l l@{}}
\textit{io} & ero salito \\
\textit{tu} & eri salito \\
\textit{lui/lei} & era salito \\
\textit{noi} & eravamo saliti \\
\textit{voi} & eravate saliti \\
\textit{loro} & erano saliti \\
\end{tabular}}{}
\end{minipage}\hspace{0.02\linewidth}\begin{minipage}[t]{0.32\linewidth}\raggedright
\paradigmbox{indicativo passato remoto}{%
\begin{tabular}{@{}l l@{}}
\textit{io} & salii \\
\textit{tu} & salisti \\
\textit{lui/lei} & salì \\
\textit{noi} & salimmo \\
\textit{voi} & saliste \\
\textit{loro} & salirono \\
\end{tabular}}{}
\end{minipage}\hspace{0.02\linewidth}\begin{minipage}[t]{0.32\linewidth}\raggedright
\paradigmbox{indicativo trapassato remoto (with essere, masculine subject)}{%
\begin{tabular}{@{}l l@{}}
\textit{io} & fui salito \\
\textit{tu} & fosti salito \\
\textit{lui/lei} & fu salito \\
\textit{noi} & fummo saliti \\
\textit{voi} & foste saliti \\
\textit{loro} & furono saliti \\
\end{tabular}}{}
\end{minipage}\par\medskip\noindent\begin{minipage}[t]{0.32\linewidth}\raggedright
\paradigmbox{indicativo futuro semplice}{%
\begin{tabular}{@{}l l@{}}
\textit{io} & salirò \\
\textit{tu} & salirai \\
\textit{lui/lei} & salirà \\
\textit{noi} & saliremo \\
\textit{voi} & salirete \\
\textit{loro} & saliranno \\
\end{tabular}}{}
\end{minipage}\hspace{0.02\linewidth}\begin{minipage}[t]{0.32\linewidth}\raggedright
\paradigmbox{indicativo futuro anteriore (with essere, masculine subject)}{%
\begin{tabular}{@{}l l@{}}
\textit{io} & sarò salito \\
\textit{tu} & sarai salito \\
\textit{lui/lei} & sarà salito \\
\textit{noi} & saremo saliti \\
\textit{voi} & sarete saliti \\
\textit{loro} & saranno saliti \\
\end{tabular}}{}
\end{minipage}\par\medskip

\moodheading{congiuntivo}

\par\noindent\begin{minipage}[t]{0.32\linewidth}\raggedright
\paradigmbox{congiuntivo presente}{%
\begin{tabular}{@{}l l@{}}
\textit{io} & \irr{salga} \\
\textit{tu} & \irr{salga} \\
\textit{lui/lei} & \irr{salga} \\
\textit{noi} & saliamo \\
\textit{voi} & saliate \\
\textit{loro} & \irr{salgano} \\
\end{tabular}}{\irr{Highlighted} forms differ from the regular pattern of the verb's family.}
\end{minipage}\hspace{0.02\linewidth}\begin{minipage}[t]{0.32\linewidth}\raggedright
\paradigmbox{congiuntivo passato (with essere, masculine subject)}{%
\begin{tabular}{@{}l l@{}}
\textit{io} & sia salito \\
\textit{tu} & sia salito \\
\textit{lui/lei} & sia salito \\
\textit{noi} & siamo saliti \\
\textit{voi} & siate saliti \\
\textit{loro} & siano saliti \\
\end{tabular}}{}
\end{minipage}\hspace{0.02\linewidth}\begin{minipage}[t]{0.32\linewidth}\raggedright
\paradigmbox{congiuntivo imperfetto}{%
\begin{tabular}{@{}l l@{}}
\textit{io} & salissi \\
\textit{tu} & salissi \\
\textit{lui/lei} & salisse \\
\textit{noi} & salissimo \\
\textit{voi} & saliste \\
\textit{loro} & salissero \\
\end{tabular}}{}
\end{minipage}\par\medskip\noindent\begin{minipage}[t]{0.32\linewidth}\raggedright
\paradigmbox{congiuntivo trapassato (with essere, masculine subject)}{%
\begin{tabular}{@{}l l@{}}
\textit{io} & fossi salito \\
\textit{tu} & fossi salito \\
\textit{lui/lei} & fosse salito \\
\textit{noi} & fossimo saliti \\
\textit{voi} & foste saliti \\
\textit{loro} & fossero saliti \\
\end{tabular}}{}
\end{minipage}\par\medskip

\moodheading{condizionale}

\par\noindent\begin{minipage}[t]{0.32\linewidth}\raggedright
\paradigmbox{condizionale presente}{%
\begin{tabular}{@{}l l@{}}
\textit{io} & salirei \\
\textit{tu} & saliresti \\
\textit{lui/lei} & salirebbe \\
\textit{noi} & saliremmo \\
\textit{voi} & salireste \\
\textit{loro} & salirebbero \\
\end{tabular}}{}
\end{minipage}\hspace{0.02\linewidth}\begin{minipage}[t]{0.32\linewidth}\raggedright
\paradigmbox{condizionale passato (with essere, masculine subject)}{%
\begin{tabular}{@{}l l@{}}
\textit{io} & sarei salito \\
\textit{tu} & saresti salito \\
\textit{lui/lei} & sarebbe salito \\
\textit{noi} & saremmo saliti \\
\textit{voi} & sareste saliti \\
\textit{loro} & sarebbero saliti \\
\end{tabular}}{}
\end{minipage}\par\medskip

\moodheading{imperativo}

\par\noindent\begin{minipage}[t]{0.32\linewidth}\raggedright
\paradigmbox{imperativo}{%
\begin{tabular}{@{}l l@{}}
\textit{tu} & sali \\
\textit{Lei} & \irr{salga} \\
\textit{noi} & saliamo \\
\textit{voi} & salite \\
\textit{Loro} & \irr{salgano} \\
\end{tabular}}{\irr{Highlighted} forms differ from the regular pattern of the verb's family. Negative \textit{tu}: \textit{non salire}. \textit{Lei} and \textit{Loro} are the congiuntivo forms.}
\end{minipage}\par\medskip

\moodheading{forme non finite}
{\small\begin{tabular}{@{}l l@{}}
infinito presente & \textit{salire} \\
infinito passato & \textit{essere salito} \\
gerundio presente & \textit{salendo} \\
gerundio passato & \textit{essendo salito} \\
participio passato & \textit{salito} \\
\end{tabular}}\par
\clearpage

\subsection{piacere --- to please; to be liked}\label{conj:piacere}\index{piacere@piacere}
{\small Auxiliary: \textit{essere}. The thing liked is the grammatical subject: mi piace il gelato (ice cream pleases me), mi piacciono le mele (apples please me). The cc appears in io, noi and loro.}\par

\moodheading{indicativo}

\par\noindent\begin{minipage}[t]{0.32\linewidth}\raggedright
\paradigmbox{indicativo presente}{%
\begin{tabular}{@{}l l@{}}
\textit{io} & \irr{piaccio} \\
\textit{tu} & piaci \\
\textit{lui/lei} & piace \\
\textit{noi} & \irr{piacciamo} \\
\textit{voi} & piacete \\
\textit{loro} & \irr{piacciono} \\
\end{tabular}}{\irr{Highlighted} forms differ from the regular pattern of the verb's family.}
\end{minipage}\hspace{0.02\linewidth}\begin{minipage}[t]{0.32\linewidth}\raggedright
\paradigmbox{indicativo passato prossimo (with essere, masculine subject)}{%
\begin{tabular}{@{}l l@{}}
\textit{io} & sono piaciuto \\
\textit{tu} & sei piaciuto \\
\textit{lui/lei} & è piaciuto \\
\textit{noi} & siamo piaciuti \\
\textit{voi} & siete piaciuti \\
\textit{loro} & sono piaciuti \\
\end{tabular}}{}
\end{minipage}\hspace{0.02\linewidth}\begin{minipage}[t]{0.32\linewidth}\raggedright
\paradigmbox{indicativo imperfetto}{%
\begin{tabular}{@{}l l@{}}
\textit{io} & piacevo \\
\textit{tu} & piacevi \\
\textit{lui/lei} & piaceva \\
\textit{noi} & piacevamo \\
\textit{voi} & piacevate \\
\textit{loro} & piacevano \\
\end{tabular}}{}
\end{minipage}\par\medskip\noindent\begin{minipage}[t]{0.32\linewidth}\raggedright
\paradigmbox{indicativo trapassato prossimo (with essere, masculine subject)}{%
\begin{tabular}{@{}l l@{}}
\textit{io} & ero piaciuto \\
\textit{tu} & eri piaciuto \\
\textit{lui/lei} & era piaciuto \\
\textit{noi} & eravamo piaciuti \\
\textit{voi} & eravate piaciuti \\
\textit{loro} & erano piaciuti \\
\end{tabular}}{}
\end{minipage}\hspace{0.02\linewidth}\begin{minipage}[t]{0.32\linewidth}\raggedright
\paradigmbox{indicativo passato remoto}{%
\begin{tabular}{@{}l l@{}}
\textit{io} & \irr{piacqui} \\
\textit{tu} & piacesti \\
\textit{lui/lei} & \irr{piacque} \\
\textit{noi} & piacemmo \\
\textit{voi} & piaceste \\
\textit{loro} & \irr{piacquero} \\
\end{tabular}}{\irr{Highlighted} forms differ from the regular pattern of the verb's family.}
\end{minipage}\hspace{0.02\linewidth}\begin{minipage}[t]{0.32\linewidth}\raggedright
\paradigmbox{indicativo trapassato remoto (with essere, masculine subject)}{%
\begin{tabular}{@{}l l@{}}
\textit{io} & fui piaciuto \\
\textit{tu} & fosti piaciuto \\
\textit{lui/lei} & fu piaciuto \\
\textit{noi} & fummo piaciuti \\
\textit{voi} & foste piaciuti \\
\textit{loro} & furono piaciuti \\
\end{tabular}}{}
\end{minipage}\par\medskip\noindent\begin{minipage}[t]{0.32\linewidth}\raggedright
\paradigmbox{indicativo futuro semplice}{%
\begin{tabular}{@{}l l@{}}
\textit{io} & piacerò \\
\textit{tu} & piacerai \\
\textit{lui/lei} & piacerà \\
\textit{noi} & piaceremo \\
\textit{voi} & piacerete \\
\textit{loro} & piaceranno \\
\end{tabular}}{}
\end{minipage}\hspace{0.02\linewidth}\begin{minipage}[t]{0.32\linewidth}\raggedright
\paradigmbox{indicativo futuro anteriore (with essere, masculine subject)}{%
\begin{tabular}{@{}l l@{}}
\textit{io} & sarò piaciuto \\
\textit{tu} & sarai piaciuto \\
\textit{lui/lei} & sarà piaciuto \\
\textit{noi} & saremo piaciuti \\
\textit{voi} & sarete piaciuti \\
\textit{loro} & saranno piaciuti \\
\end{tabular}}{}
\end{minipage}\par\medskip

\moodheading{congiuntivo}

\par\noindent\begin{minipage}[t]{0.32\linewidth}\raggedright
\paradigmbox{congiuntivo presente}{%
\begin{tabular}{@{}l l@{}}
\textit{io} & \irr{piaccia} \\
\textit{tu} & \irr{piaccia} \\
\textit{lui/lei} & \irr{piaccia} \\
\textit{noi} & \irr{piacciamo} \\
\textit{voi} & \irr{piacciate} \\
\textit{loro} & \irr{piacciano} \\
\end{tabular}}{\irr{Highlighted} forms differ from the regular pattern of the verb's family.}
\end{minipage}\hspace{0.02\linewidth}\begin{minipage}[t]{0.32\linewidth}\raggedright
\paradigmbox{congiuntivo passato (with essere, masculine subject)}{%
\begin{tabular}{@{}l l@{}}
\textit{io} & sia piaciuto \\
\textit{tu} & sia piaciuto \\
\textit{lui/lei} & sia piaciuto \\
\textit{noi} & siamo piaciuti \\
\textit{voi} & siate piaciuti \\
\textit{loro} & siano piaciuti \\
\end{tabular}}{}
\end{minipage}\hspace{0.02\linewidth}\begin{minipage}[t]{0.32\linewidth}\raggedright
\paradigmbox{congiuntivo imperfetto}{%
\begin{tabular}{@{}l l@{}}
\textit{io} & piacessi \\
\textit{tu} & piacessi \\
\textit{lui/lei} & piacesse \\
\textit{noi} & piacessimo \\
\textit{voi} & piaceste \\
\textit{loro} & piacessero \\
\end{tabular}}{}
\end{minipage}\par\medskip\noindent\begin{minipage}[t]{0.32\linewidth}\raggedright
\paradigmbox{congiuntivo trapassato (with essere, masculine subject)}{%
\begin{tabular}{@{}l l@{}}
\textit{io} & fossi piaciuto \\
\textit{tu} & fossi piaciuto \\
\textit{lui/lei} & fosse piaciuto \\
\textit{noi} & fossimo piaciuti \\
\textit{voi} & foste piaciuti \\
\textit{loro} & fossero piaciuti \\
\end{tabular}}{}
\end{minipage}\par\medskip

\moodheading{condizionale}

\par\noindent\begin{minipage}[t]{0.32\linewidth}\raggedright
\paradigmbox{condizionale presente}{%
\begin{tabular}{@{}l l@{}}
\textit{io} & piacerei \\
\textit{tu} & piaceresti \\
\textit{lui/lei} & piacerebbe \\
\textit{noi} & piaceremmo \\
\textit{voi} & piacereste \\
\textit{loro} & piacerebbero \\
\end{tabular}}{}
\end{minipage}\hspace{0.02\linewidth}\begin{minipage}[t]{0.32\linewidth}\raggedright
\paradigmbox{condizionale passato (with essere, masculine subject)}{%
\begin{tabular}{@{}l l@{}}
\textit{io} & sarei piaciuto \\
\textit{tu} & saresti piaciuto \\
\textit{lui/lei} & sarebbe piaciuto \\
\textit{noi} & saremmo piaciuti \\
\textit{voi} & sareste piaciuti \\
\textit{loro} & sarebbero piaciuti \\
\end{tabular}}{}
\end{minipage}\par\medskip

\moodheading{imperativo}

\par\noindent\begin{minipage}[t]{0.32\linewidth}\raggedright
\paradigmbox{imperativo}{%
\begin{tabular}{@{}l l@{}}
\textit{tu} & piaci \\
\textit{Lei} & \irr{piaccia} \\
\textit{noi} & \irr{piacciamo} \\
\textit{voi} & piacete \\
\textit{Loro} & \irr{piacciano} \\
\end{tabular}}{\irr{Highlighted} forms differ from the regular pattern of the verb's family. Negative \textit{tu}: \textit{non piacere}. \textit{Lei} and \textit{Loro} are the congiuntivo forms.}
\end{minipage}\par\medskip

\moodheading{forme non finite}
{\small\begin{tabular}{@{}l l@{}}
infinito presente & \textit{piacere} \\
infinito passato & \textit{essere piaciuto} \\
gerundio presente & \textit{piacendo} \\
gerundio passato & \textit{essendo piaciuto} \\
participio passato & \textit{piaciuto} \\
\end{tabular}}\par
\clearpage

\subsection{sedersi --- to sit down}\label{conj:sedersi}\index{sedersi@sedersi}
{\small Auxiliary: \textit{essere}. The stressed e becomes ie (siedo, siedi, siede, siedono). Normally used reflexively: sedersi, mi siedo.}\par

\moodheading{indicativo}

\par\noindent\begin{minipage}[t]{0.32\linewidth}\raggedright
\paradigmbox{indicativo presente}{%
\begin{tabular}{@{}l l@{}}
\textit{io} & \irr{mi siedo} \\
\textit{tu} & \irr{ti siedi} \\
\textit{lui/lei} & \irr{si siede} \\
\textit{noi} & ci sediamo \\
\textit{voi} & vi sedete \\
\textit{loro} & \irr{si siedono} \\
\end{tabular}}{\irr{Highlighted} forms differ from the regular pattern of the verb's family.}
\end{minipage}\hspace{0.02\linewidth}\begin{minipage}[t]{0.32\linewidth}\raggedright
\paradigmbox{indicativo passato prossimo (with essere, masculine subject)}{%
\begin{tabular}{@{}l l@{}}
\textit{io} & mi sono seduto \\
\textit{tu} & ti sei seduto \\
\textit{lui/lei} & si è seduto \\
\textit{noi} & ci siamo seduti \\
\textit{voi} & vi siete seduti \\
\textit{loro} & si sono seduti \\
\end{tabular}}{}
\end{minipage}\hspace{0.02\linewidth}\begin{minipage}[t]{0.32\linewidth}\raggedright
\paradigmbox{indicativo imperfetto}{%
\begin{tabular}{@{}l l@{}}
\textit{io} & mi sedevo \\
\textit{tu} & ti sedevi \\
\textit{lui/lei} & si sedeva \\
\textit{noi} & ci sedevamo \\
\textit{voi} & vi sedevate \\
\textit{loro} & si sedevano \\
\end{tabular}}{}
\end{minipage}\par\medskip\noindent\begin{minipage}[t]{0.32\linewidth}\raggedright
\paradigmbox{indicativo trapassato prossimo (with essere, masculine subject)}{%
\begin{tabular}{@{}l l@{}}
\textit{io} & mi ero seduto \\
\textit{tu} & ti eri seduto \\
\textit{lui/lei} & si era seduto \\
\textit{noi} & ci eravamo seduti \\
\textit{voi} & vi eravate seduti \\
\textit{loro} & si erano seduti \\
\end{tabular}}{}
\end{minipage}\hspace{0.02\linewidth}\begin{minipage}[t]{0.32\linewidth}\raggedright
\paradigmbox{indicativo passato remoto}{%
\begin{tabular}{@{}l l@{}}
\textit{io} & mi sedei / mi sedetti \\
\textit{tu} & ti sedesti \\
\textit{lui/lei} & si sedé / si sedette \\
\textit{noi} & ci sedemmo \\
\textit{voi} & vi sedeste \\
\textit{loro} & si sederono / si sedettero \\
\end{tabular}}{Where two forms are given, both are standard.}
\end{minipage}\hspace{0.02\linewidth}\begin{minipage}[t]{0.32\linewidth}\raggedright
\paradigmbox{indicativo trapassato remoto (with essere, masculine subject)}{%
\begin{tabular}{@{}l l@{}}
\textit{io} & mi fui seduto \\
\textit{tu} & ti fosti seduto \\
\textit{lui/lei} & si fu seduto \\
\textit{noi} & ci fummo seduti \\
\textit{voi} & vi foste seduti \\
\textit{loro} & si furono seduti \\
\end{tabular}}{}
\end{minipage}\par\medskip\noindent\begin{minipage}[t]{0.32\linewidth}\raggedright
\paradigmbox{indicativo futuro semplice}{%
\begin{tabular}{@{}l l@{}}
\textit{io} & mi sederò / mi siederò \\
\textit{tu} & ti sederai / ti siederai \\
\textit{lui/lei} & si sederà / si siederà \\
\textit{noi} & ci sederemo / ci siederemo \\
\textit{voi} & vi sederete / vi siederete \\
\textit{loro} & si sederanno / si siederanno \\
\end{tabular}}{Where two forms are given, both are standard.}
\end{minipage}\hspace{0.02\linewidth}\begin{minipage}[t]{0.32\linewidth}\raggedright
\paradigmbox{indicativo futuro anteriore (with essere, masculine subject)}{%
\begin{tabular}{@{}l l@{}}
\textit{io} & mi sarò seduto \\
\textit{tu} & ti sarai seduto \\
\textit{lui/lei} & si sarà seduto \\
\textit{noi} & ci saremo seduti \\
\textit{voi} & vi sarete seduti \\
\textit{loro} & si saranno seduti \\
\end{tabular}}{}
\end{minipage}\par\medskip

\moodheading{congiuntivo}

\par\noindent\begin{minipage}[t]{0.32\linewidth}\raggedright
\paradigmbox{congiuntivo presente}{%
\begin{tabular}{@{}l l@{}}
\textit{io} & \irr{mi sieda} \\
\textit{tu} & \irr{ti sieda} \\
\textit{lui/lei} & \irr{si sieda} \\
\textit{noi} & ci sediamo \\
\textit{voi} & vi sediate \\
\textit{loro} & \irr{si siedano} \\
\end{tabular}}{\irr{Highlighted} forms differ from the regular pattern of the verb's family.}
\end{minipage}\hspace{0.02\linewidth}\begin{minipage}[t]{0.32\linewidth}\raggedright
\paradigmbox{congiuntivo passato (with essere, masculine subject)}{%
\begin{tabular}{@{}l l@{}}
\textit{io} & mi sia seduto \\
\textit{tu} & ti sia seduto \\
\textit{lui/lei} & si sia seduto \\
\textit{noi} & ci siamo seduti \\
\textit{voi} & vi siate seduti \\
\textit{loro} & si siano seduti \\
\end{tabular}}{}
\end{minipage}\hspace{0.02\linewidth}\begin{minipage}[t]{0.32\linewidth}\raggedright
\paradigmbox{congiuntivo imperfetto}{%
\begin{tabular}{@{}l l@{}}
\textit{io} & mi sedessi \\
\textit{tu} & ti sedessi \\
\textit{lui/lei} & si sedesse \\
\textit{noi} & ci sedessimo \\
\textit{voi} & vi sedeste \\
\textit{loro} & si sedessero \\
\end{tabular}}{}
\end{minipage}\par\medskip\noindent\begin{minipage}[t]{0.32\linewidth}\raggedright
\paradigmbox{congiuntivo trapassato (with essere, masculine subject)}{%
\begin{tabular}{@{}l l@{}}
\textit{io} & mi fossi seduto \\
\textit{tu} & ti fossi seduto \\
\textit{lui/lei} & si fosse seduto \\
\textit{noi} & ci fossimo seduti \\
\textit{voi} & vi foste seduti \\
\textit{loro} & si fossero seduti \\
\end{tabular}}{}
\end{minipage}\par\medskip

\moodheading{condizionale}

\par\noindent\begin{minipage}[t]{0.32\linewidth}\raggedright
\paradigmbox{condizionale presente}{%
\begin{tabular}{@{}l l@{}}
\textit{io} & mi sederei / mi siederei \\
\textit{tu} & ti sederesti / ti siederesti \\
\textit{lui/lei} & si sederebbe / si siederebbe \\
\textit{noi} & ci sederemmo / ci siederemmo \\
\textit{voi} & vi sedereste / vi siedereste \\
\textit{loro} & si sederebbero / si siederebbero \\
\end{tabular}}{Where two forms are given, both are standard.}
\end{minipage}\hspace{0.02\linewidth}\begin{minipage}[t]{0.32\linewidth}\raggedright
\paradigmbox{condizionale passato (with essere, masculine subject)}{%
\begin{tabular}{@{}l l@{}}
\textit{io} & mi sarei seduto \\
\textit{tu} & ti saresti seduto \\
\textit{lui/lei} & si sarebbe seduto \\
\textit{noi} & ci saremmo seduti \\
\textit{voi} & vi sareste seduti \\
\textit{loro} & si sarebbero seduti \\
\end{tabular}}{}
\end{minipage}\par\medskip

\moodheading{imperativo}

\par\noindent\begin{minipage}[t]{0.32\linewidth}\raggedright
\paradigmbox{imperativo}{%
\begin{tabular}{@{}l l@{}}
\textit{tu} & \irr{siediti} \\
\textit{Lei} & \irr{si sieda} \\
\textit{noi} & sediamoci \\
\textit{voi} & sedetevi \\
\textit{Loro} & \irr{si siedano} \\
\end{tabular}}{\irr{Highlighted} forms differ from the regular pattern of the verb's family. Negative \textit{tu}: \textit{non sederti / non ti sedere}. \textit{Lei} and \textit{Loro} are the congiuntivo forms.}
\end{minipage}\par\medskip

\moodheading{forme non finite}
{\small\begin{tabular}{@{}l l@{}}
infinito presente & \textit{sedersi} \\
infinito passato & \textit{essersi seduto} \\
gerundio presente & \textit{sedendosi} \\
gerundio passato & \textit{essendosi seduto} \\
participio passato & \textit{seduto} \\
\end{tabular}}\par
\clearpage

\subsection{scegliere --- to choose}\label{conj:scegliere}\index{scegliere@scegliere}
{\small Auxiliary: \textit{avere}. gli becomes lg before -o: scelgo, scelgono. Same pattern: raccogliere (raccolgo), togliere (tolgo).}\par

\moodheading{indicativo}

\par\noindent\begin{minipage}[t]{0.32\linewidth}\raggedright
\paradigmbox{indicativo presente}{%
\begin{tabular}{@{}l l@{}}
\textit{io} & \irr{scelgo} \\
\textit{tu} & scegli \\
\textit{lui/lei} & sceglie \\
\textit{noi} & scegliamo \\
\textit{voi} & scegliete \\
\textit{loro} & \irr{scelgono} \\
\end{tabular}}{\irr{Highlighted} forms differ from the regular pattern of the verb's family.}
\end{minipage}\hspace{0.02\linewidth}\begin{minipage}[t]{0.32\linewidth}\raggedright
\paradigmbox{indicativo passato prossimo (with avere)}{%
\begin{tabular}{@{}l l@{}}
\textit{io} & ho scelto \\
\textit{tu} & hai scelto \\
\textit{lui/lei} & ha scelto \\
\textit{noi} & abbiamo scelto \\
\textit{voi} & avete scelto \\
\textit{loro} & hanno scelto \\
\end{tabular}}{}
\end{minipage}\hspace{0.02\linewidth}\begin{minipage}[t]{0.32\linewidth}\raggedright
\paradigmbox{indicativo imperfetto}{%
\begin{tabular}{@{}l l@{}}
\textit{io} & sceglievo \\
\textit{tu} & sceglievi \\
\textit{lui/lei} & sceglieva \\
\textit{noi} & sceglievamo \\
\textit{voi} & sceglievate \\
\textit{loro} & sceglievano \\
\end{tabular}}{}
\end{minipage}\par\medskip\noindent\begin{minipage}[t]{0.32\linewidth}\raggedright
\paradigmbox{indicativo trapassato prossimo (with avere)}{%
\begin{tabular}{@{}l l@{}}
\textit{io} & avevo scelto \\
\textit{tu} & avevi scelto \\
\textit{lui/lei} & aveva scelto \\
\textit{noi} & avevamo scelto \\
\textit{voi} & avevate scelto \\
\textit{loro} & avevano scelto \\
\end{tabular}}{}
\end{minipage}\hspace{0.02\linewidth}\begin{minipage}[t]{0.32\linewidth}\raggedright
\paradigmbox{indicativo passato remoto}{%
\begin{tabular}{@{}l l@{}}
\textit{io} & \irr{scelsi} \\
\textit{tu} & scegliesti \\
\textit{lui/lei} & \irr{scelse} \\
\textit{noi} & scegliemmo \\
\textit{voi} & sceglieste \\
\textit{loro} & \irr{scelsero} \\
\end{tabular}}{\irr{Highlighted} forms differ from the regular pattern of the verb's family.}
\end{minipage}\hspace{0.02\linewidth}\begin{minipage}[t]{0.32\linewidth}\raggedright
\paradigmbox{indicativo trapassato remoto (with avere)}{%
\begin{tabular}{@{}l l@{}}
\textit{io} & ebbi scelto \\
\textit{tu} & avesti scelto \\
\textit{lui/lei} & ebbe scelto \\
\textit{noi} & avemmo scelto \\
\textit{voi} & aveste scelto \\
\textit{loro} & ebbero scelto \\
\end{tabular}}{}
\end{minipage}\par\medskip\noindent\begin{minipage}[t]{0.32\linewidth}\raggedright
\paradigmbox{indicativo futuro semplice}{%
\begin{tabular}{@{}l l@{}}
\textit{io} & sceglierò \\
\textit{tu} & sceglierai \\
\textit{lui/lei} & sceglierà \\
\textit{noi} & sceglieremo \\
\textit{voi} & sceglierete \\
\textit{loro} & sceglieranno \\
\end{tabular}}{}
\end{minipage}\hspace{0.02\linewidth}\begin{minipage}[t]{0.32\linewidth}\raggedright
\paradigmbox{indicativo futuro anteriore (with avere)}{%
\begin{tabular}{@{}l l@{}}
\textit{io} & avrò scelto \\
\textit{tu} & avrai scelto \\
\textit{lui/lei} & avrà scelto \\
\textit{noi} & avremo scelto \\
\textit{voi} & avrete scelto \\
\textit{loro} & avranno scelto \\
\end{tabular}}{}
\end{minipage}\par\medskip

\moodheading{congiuntivo}

\par\noindent\begin{minipage}[t]{0.32\linewidth}\raggedright
\paradigmbox{congiuntivo presente}{%
\begin{tabular}{@{}l l@{}}
\textit{io} & \irr{scelga} \\
\textit{tu} & \irr{scelga} \\
\textit{lui/lei} & \irr{scelga} \\
\textit{noi} & scegliamo \\
\textit{voi} & scegliate \\
\textit{loro} & \irr{scelgano} \\
\end{tabular}}{\irr{Highlighted} forms differ from the regular pattern of the verb's family.}
\end{minipage}\hspace{0.02\linewidth}\begin{minipage}[t]{0.32\linewidth}\raggedright
\paradigmbox{congiuntivo passato (with avere)}{%
\begin{tabular}{@{}l l@{}}
\textit{io} & abbia scelto \\
\textit{tu} & abbia scelto \\
\textit{lui/lei} & abbia scelto \\
\textit{noi} & abbiamo scelto \\
\textit{voi} & abbiate scelto \\
\textit{loro} & abbiano scelto \\
\end{tabular}}{}
\end{minipage}\hspace{0.02\linewidth}\begin{minipage}[t]{0.32\linewidth}\raggedright
\paradigmbox{congiuntivo imperfetto}{%
\begin{tabular}{@{}l l@{}}
\textit{io} & scegliessi \\
\textit{tu} & scegliessi \\
\textit{lui/lei} & scegliesse \\
\textit{noi} & scegliessimo \\
\textit{voi} & sceglieste \\
\textit{loro} & scegliessero \\
\end{tabular}}{}
\end{minipage}\par\medskip\noindent\begin{minipage}[t]{0.32\linewidth}\raggedright
\paradigmbox{congiuntivo trapassato (with avere)}{%
\begin{tabular}{@{}l l@{}}
\textit{io} & avessi scelto \\
\textit{tu} & avessi scelto \\
\textit{lui/lei} & avesse scelto \\
\textit{noi} & avessimo scelto \\
\textit{voi} & aveste scelto \\
\textit{loro} & avessero scelto \\
\end{tabular}}{}
\end{minipage}\par\medskip

\moodheading{condizionale}

\par\noindent\begin{minipage}[t]{0.32\linewidth}\raggedright
\paradigmbox{condizionale presente}{%
\begin{tabular}{@{}l l@{}}
\textit{io} & sceglierei \\
\textit{tu} & sceglieresti \\
\textit{lui/lei} & sceglierebbe \\
\textit{noi} & sceglieremmo \\
\textit{voi} & scegliereste \\
\textit{loro} & sceglierebbero \\
\end{tabular}}{}
\end{minipage}\hspace{0.02\linewidth}\begin{minipage}[t]{0.32\linewidth}\raggedright
\paradigmbox{condizionale passato (with avere)}{%
\begin{tabular}{@{}l l@{}}
\textit{io} & avrei scelto \\
\textit{tu} & avresti scelto \\
\textit{lui/lei} & avrebbe scelto \\
\textit{noi} & avremmo scelto \\
\textit{voi} & avreste scelto \\
\textit{loro} & avrebbero scelto \\
\end{tabular}}{}
\end{minipage}\par\medskip

\moodheading{imperativo}

\par\noindent\begin{minipage}[t]{0.32\linewidth}\raggedright
\paradigmbox{imperativo}{%
\begin{tabular}{@{}l l@{}}
\textit{tu} & scegli \\
\textit{Lei} & \irr{scelga} \\
\textit{noi} & scegliamo \\
\textit{voi} & scegliete \\
\textit{Loro} & \irr{scelgano} \\
\end{tabular}}{\irr{Highlighted} forms differ from the regular pattern of the verb's family. Negative \textit{tu}: \textit{non scegliere}. \textit{Lei} and \textit{Loro} are the congiuntivo forms.}
\end{minipage}\par\medskip

\moodheading{forme non finite}
{\small\begin{tabular}{@{}l l@{}}
infinito presente & \textit{scegliere} \\
infinito passato & \textit{avere scelto} \\
gerundio presente & \textit{scegliendo} \\
gerundio passato & \textit{avendo scelto} \\
participio passato & \textit{scelto} \\
\end{tabular}}\par
\clearpage

\subsection{raccogliere --- to pick (up); to gather}\label{conj:raccogliere}\index{raccogliere@raccogliere}
{\small Auxiliary: \textit{avere}. Same pattern as scegliere: raccolgo, raccolgono; participle raccolto.}\par

\moodheading{indicativo}

\par\noindent\begin{minipage}[t]{0.32\linewidth}\raggedright
\paradigmbox{indicativo presente}{%
\begin{tabular}{@{}l l@{}}
\textit{io} & \irr{raccolgo} \\
\textit{tu} & raccogli \\
\textit{lui/lei} & raccoglie \\
\textit{noi} & raccogliamo \\
\textit{voi} & raccogliete \\
\textit{loro} & \irr{raccolgono} \\
\end{tabular}}{\irr{Highlighted} forms differ from the regular pattern of the verb's family.}
\end{minipage}\hspace{0.02\linewidth}\begin{minipage}[t]{0.32\linewidth}\raggedright
\paradigmbox{indicativo passato prossimo (with avere)}{%
\begin{tabular}{@{}l l@{}}
\textit{io} & ho raccolto \\
\textit{tu} & hai raccolto \\
\textit{lui/lei} & ha raccolto \\
\textit{noi} & abbiamo raccolto \\
\textit{voi} & avete raccolto \\
\textit{loro} & hanno raccolto \\
\end{tabular}}{}
\end{minipage}\hspace{0.02\linewidth}\begin{minipage}[t]{0.32\linewidth}\raggedright
\paradigmbox{indicativo imperfetto}{%
\begin{tabular}{@{}l l@{}}
\textit{io} & raccoglievo \\
\textit{tu} & raccoglievi \\
\textit{lui/lei} & raccoglieva \\
\textit{noi} & raccoglievamo \\
\textit{voi} & raccoglievate \\
\textit{loro} & raccoglievano \\
\end{tabular}}{}
\end{minipage}\par\medskip\noindent\begin{minipage}[t]{0.32\linewidth}\raggedright
\paradigmbox{indicativo trapassato prossimo (with avere)}{%
\begin{tabular}{@{}l l@{}}
\textit{io} & avevo raccolto \\
\textit{tu} & avevi raccolto \\
\textit{lui/lei} & aveva raccolto \\
\textit{noi} & avevamo raccolto \\
\textit{voi} & avevate raccolto \\
\textit{loro} & avevano raccolto \\
\end{tabular}}{}
\end{minipage}\hspace{0.02\linewidth}\begin{minipage}[t]{0.32\linewidth}\raggedright
\paradigmbox{indicativo passato remoto}{%
\begin{tabular}{@{}l l@{}}
\textit{io} & \irr{raccolsi} \\
\textit{tu} & raccogliesti \\
\textit{lui/lei} & \irr{raccolse} \\
\textit{noi} & raccogliemmo \\
\textit{voi} & raccoglieste \\
\textit{loro} & \irr{raccolsero} \\
\end{tabular}}{\irr{Highlighted} forms differ from the regular pattern of the verb's family.}
\end{minipage}\hspace{0.02\linewidth}\begin{minipage}[t]{0.32\linewidth}\raggedright
\paradigmbox{indicativo trapassato remoto (with avere)}{%
\begin{tabular}{@{}l l@{}}
\textit{io} & ebbi raccolto \\
\textit{tu} & avesti raccolto \\
\textit{lui/lei} & ebbe raccolto \\
\textit{noi} & avemmo raccolto \\
\textit{voi} & aveste raccolto \\
\textit{loro} & ebbero raccolto \\
\end{tabular}}{}
\end{minipage}\par\medskip\noindent\begin{minipage}[t]{0.32\linewidth}\raggedright
\paradigmbox{indicativo futuro semplice}{%
\begin{tabular}{@{}l l@{}}
\textit{io} & raccoglierò \\
\textit{tu} & raccoglierai \\
\textit{lui/lei} & raccoglierà \\
\textit{noi} & raccoglieremo \\
\textit{voi} & raccoglierete \\
\textit{loro} & raccoglieranno \\
\end{tabular}}{}
\end{minipage}\hspace{0.02\linewidth}\begin{minipage}[t]{0.32\linewidth}\raggedright
\paradigmbox{indicativo futuro anteriore (with avere)}{%
\begin{tabular}{@{}l l@{}}
\textit{io} & avrò raccolto \\
\textit{tu} & avrai raccolto \\
\textit{lui/lei} & avrà raccolto \\
\textit{noi} & avremo raccolto \\
\textit{voi} & avrete raccolto \\
\textit{loro} & avranno raccolto \\
\end{tabular}}{}
\end{minipage}\par\medskip

\moodheading{congiuntivo}

\par\noindent\begin{minipage}[t]{0.32\linewidth}\raggedright
\paradigmbox{congiuntivo presente}{%
\begin{tabular}{@{}l l@{}}
\textit{io} & \irr{raccolga} \\
\textit{tu} & \irr{raccolga} \\
\textit{lui/lei} & \irr{raccolga} \\
\textit{noi} & raccogliamo \\
\textit{voi} & raccogliate \\
\textit{loro} & \irr{raccolgano} \\
\end{tabular}}{\irr{Highlighted} forms differ from the regular pattern of the verb's family.}
\end{minipage}\hspace{0.02\linewidth}\begin{minipage}[t]{0.32\linewidth}\raggedright
\paradigmbox{congiuntivo passato (with avere)}{%
\begin{tabular}{@{}l l@{}}
\textit{io} & abbia raccolto \\
\textit{tu} & abbia raccolto \\
\textit{lui/lei} & abbia raccolto \\
\textit{noi} & abbiamo raccolto \\
\textit{voi} & abbiate raccolto \\
\textit{loro} & abbiano raccolto \\
\end{tabular}}{}
\end{minipage}\hspace{0.02\linewidth}\begin{minipage}[t]{0.32\linewidth}\raggedright
\paradigmbox{congiuntivo imperfetto}{%
\begin{tabular}{@{}l l@{}}
\textit{io} & raccogliessi \\
\textit{tu} & raccogliessi \\
\textit{lui/lei} & raccogliesse \\
\textit{noi} & raccogliessimo \\
\textit{voi} & raccoglieste \\
\textit{loro} & raccogliessero \\
\end{tabular}}{}
\end{minipage}\par\medskip\noindent\begin{minipage}[t]{0.32\linewidth}\raggedright
\paradigmbox{congiuntivo trapassato (with avere)}{%
\begin{tabular}{@{}l l@{}}
\textit{io} & avessi raccolto \\
\textit{tu} & avessi raccolto \\
\textit{lui/lei} & avesse raccolto \\
\textit{noi} & avessimo raccolto \\
\textit{voi} & aveste raccolto \\
\textit{loro} & avessero raccolto \\
\end{tabular}}{}
\end{minipage}\par\medskip

\moodheading{condizionale}

\par\noindent\begin{minipage}[t]{0.32\linewidth}\raggedright
\paradigmbox{condizionale presente}{%
\begin{tabular}{@{}l l@{}}
\textit{io} & raccoglierei \\
\textit{tu} & raccoglieresti \\
\textit{lui/lei} & raccoglierebbe \\
\textit{noi} & raccoglieremmo \\
\textit{voi} & raccogliereste \\
\textit{loro} & raccoglierebbero \\
\end{tabular}}{}
\end{minipage}\hspace{0.02\linewidth}\begin{minipage}[t]{0.32\linewidth}\raggedright
\paradigmbox{condizionale passato (with avere)}{%
\begin{tabular}{@{}l l@{}}
\textit{io} & avrei raccolto \\
\textit{tu} & avresti raccolto \\
\textit{lui/lei} & avrebbe raccolto \\
\textit{noi} & avremmo raccolto \\
\textit{voi} & avreste raccolto \\
\textit{loro} & avrebbero raccolto \\
\end{tabular}}{}
\end{minipage}\par\medskip

\moodheading{imperativo}

\par\noindent\begin{minipage}[t]{0.32\linewidth}\raggedright
\paradigmbox{imperativo}{%
\begin{tabular}{@{}l l@{}}
\textit{tu} & raccogli \\
\textit{Lei} & \irr{raccolga} \\
\textit{noi} & raccogliamo \\
\textit{voi} & raccogliete \\
\textit{Loro} & \irr{raccolgano} \\
\end{tabular}}{\irr{Highlighted} forms differ from the regular pattern of the verb's family. Negative \textit{tu}: \textit{non raccogliere}. \textit{Lei} and \textit{Loro} are the congiuntivo forms.}
\end{minipage}\par\medskip

\moodheading{forme non finite}
{\small\begin{tabular}{@{}l l@{}}
infinito presente & \textit{raccogliere} \\
infinito passato & \textit{avere raccolto} \\
gerundio presente & \textit{raccogliendo} \\
gerundio passato & \textit{avendo raccolto} \\
participio passato & \textit{raccolto} \\
\end{tabular}}\par
\clearpage

\subsection{riuscire --- to manage; to succeed}\label{conj:riuscire}\index{riuscire@riuscire}
{\small Auxiliary: \textit{essere}. Conjugated like uscire: riesco, riesci, riesce, riusciamo, riuscite, riescono. riuscire a + infinitive = to manage to.}\par

\moodheading{indicativo}

\par\noindent\begin{minipage}[t]{0.32\linewidth}\raggedright
\paradigmbox{indicativo presente}{%
\begin{tabular}{@{}l l@{}}
\textit{io} & \irr{riesco} \\
\textit{tu} & \irr{riesci} \\
\textit{lui/lei} & \irr{riesce} \\
\textit{noi} & riusciamo \\
\textit{voi} & riuscite \\
\textit{loro} & \irr{riescono} \\
\end{tabular}}{\irr{Highlighted} forms differ from the regular pattern of the verb's family.}
\end{minipage}\hspace{0.02\linewidth}\begin{minipage}[t]{0.32\linewidth}\raggedright
\paradigmbox{indicativo passato prossimo (with essere, masculine subject)}{%
\begin{tabular}{@{}l l@{}}
\textit{io} & sono riuscito \\
\textit{tu} & sei riuscito \\
\textit{lui/lei} & è riuscito \\
\textit{noi} & siamo riusciti \\
\textit{voi} & siete riusciti \\
\textit{loro} & sono riusciti \\
\end{tabular}}{}
\end{minipage}\hspace{0.02\linewidth}\begin{minipage}[t]{0.32\linewidth}\raggedright
\paradigmbox{indicativo imperfetto}{%
\begin{tabular}{@{}l l@{}}
\textit{io} & riuscivo \\
\textit{tu} & riuscivi \\
\textit{lui/lei} & riusciva \\
\textit{noi} & riuscivamo \\
\textit{voi} & riuscivate \\
\textit{loro} & riuscivano \\
\end{tabular}}{}
\end{minipage}\par\medskip\noindent\begin{minipage}[t]{0.32\linewidth}\raggedright
\paradigmbox{indicativo trapassato prossimo (with essere, masculine subject)}{%
\begin{tabular}{@{}l l@{}}
\textit{io} & ero riuscito \\
\textit{tu} & eri riuscito \\
\textit{lui/lei} & era riuscito \\
\textit{noi} & eravamo riusciti \\
\textit{voi} & eravate riusciti \\
\textit{loro} & erano riusciti \\
\end{tabular}}{}
\end{minipage}\hspace{0.02\linewidth}\begin{minipage}[t]{0.32\linewidth}\raggedright
\paradigmbox{indicativo passato remoto}{%
\begin{tabular}{@{}l l@{}}
\textit{io} & riuscii \\
\textit{tu} & riuscisti \\
\textit{lui/lei} & riuscì \\
\textit{noi} & riuscimmo \\
\textit{voi} & riusciste \\
\textit{loro} & riuscirono \\
\end{tabular}}{}
\end{minipage}\hspace{0.02\linewidth}\begin{minipage}[t]{0.32\linewidth}\raggedright
\paradigmbox{indicativo trapassato remoto (with essere, masculine subject)}{%
\begin{tabular}{@{}l l@{}}
\textit{io} & fui riuscito \\
\textit{tu} & fosti riuscito \\
\textit{lui/lei} & fu riuscito \\
\textit{noi} & fummo riusciti \\
\textit{voi} & foste riusciti \\
\textit{loro} & furono riusciti \\
\end{tabular}}{}
\end{minipage}\par\medskip\noindent\begin{minipage}[t]{0.32\linewidth}\raggedright
\paradigmbox{indicativo futuro semplice}{%
\begin{tabular}{@{}l l@{}}
\textit{io} & riuscirò \\
\textit{tu} & riuscirai \\
\textit{lui/lei} & riuscirà \\
\textit{noi} & riusciremo \\
\textit{voi} & riuscirete \\
\textit{loro} & riusciranno \\
\end{tabular}}{}
\end{minipage}\hspace{0.02\linewidth}\begin{minipage}[t]{0.32\linewidth}\raggedright
\paradigmbox{indicativo futuro anteriore (with essere, masculine subject)}{%
\begin{tabular}{@{}l l@{}}
\textit{io} & sarò riuscito \\
\textit{tu} & sarai riuscito \\
\textit{lui/lei} & sarà riuscito \\
\textit{noi} & saremo riusciti \\
\textit{voi} & sarete riusciti \\
\textit{loro} & saranno riusciti \\
\end{tabular}}{}
\end{minipage}\par\medskip

\moodheading{congiuntivo}

\par\noindent\begin{minipage}[t]{0.32\linewidth}\raggedright
\paradigmbox{congiuntivo presente}{%
\begin{tabular}{@{}l l@{}}
\textit{io} & \irr{riesca} \\
\textit{tu} & \irr{riesca} \\
\textit{lui/lei} & \irr{riesca} \\
\textit{noi} & riusciamo \\
\textit{voi} & riusciate \\
\textit{loro} & \irr{riescano} \\
\end{tabular}}{\irr{Highlighted} forms differ from the regular pattern of the verb's family.}
\end{minipage}\hspace{0.02\linewidth}\begin{minipage}[t]{0.32\linewidth}\raggedright
\paradigmbox{congiuntivo passato (with essere, masculine subject)}{%
\begin{tabular}{@{}l l@{}}
\textit{io} & sia riuscito \\
\textit{tu} & sia riuscito \\
\textit{lui/lei} & sia riuscito \\
\textit{noi} & siamo riusciti \\
\textit{voi} & siate riusciti \\
\textit{loro} & siano riusciti \\
\end{tabular}}{}
\end{minipage}\hspace{0.02\linewidth}\begin{minipage}[t]{0.32\linewidth}\raggedright
\paradigmbox{congiuntivo imperfetto}{%
\begin{tabular}{@{}l l@{}}
\textit{io} & riuscissi \\
\textit{tu} & riuscissi \\
\textit{lui/lei} & riuscisse \\
\textit{noi} & riuscissimo \\
\textit{voi} & riusciste \\
\textit{loro} & riuscissero \\
\end{tabular}}{}
\end{minipage}\par\medskip\noindent\begin{minipage}[t]{0.32\linewidth}\raggedright
\paradigmbox{congiuntivo trapassato (with essere, masculine subject)}{%
\begin{tabular}{@{}l l@{}}
\textit{io} & fossi riuscito \\
\textit{tu} & fossi riuscito \\
\textit{lui/lei} & fosse riuscito \\
\textit{noi} & fossimo riusciti \\
\textit{voi} & foste riusciti \\
\textit{loro} & fossero riusciti \\
\end{tabular}}{}
\end{minipage}\par\medskip

\moodheading{condizionale}

\par\noindent\begin{minipage}[t]{0.32\linewidth}\raggedright
\paradigmbox{condizionale presente}{%
\begin{tabular}{@{}l l@{}}
\textit{io} & riuscirei \\
\textit{tu} & riusciresti \\
\textit{lui/lei} & riuscirebbe \\
\textit{noi} & riusciremmo \\
\textit{voi} & riuscireste \\
\textit{loro} & riuscirebbero \\
\end{tabular}}{}
\end{minipage}\hspace{0.02\linewidth}\begin{minipage}[t]{0.32\linewidth}\raggedright
\paradigmbox{condizionale passato (with essere, masculine subject)}{%
\begin{tabular}{@{}l l@{}}
\textit{io} & sarei riuscito \\
\textit{tu} & saresti riuscito \\
\textit{lui/lei} & sarebbe riuscito \\
\textit{noi} & saremmo riusciti \\
\textit{voi} & sareste riusciti \\
\textit{loro} & sarebbero riusciti \\
\end{tabular}}{}
\end{minipage}\par\medskip

\moodheading{imperativo}

\par\noindent\begin{minipage}[t]{0.32\linewidth}\raggedright
\paradigmbox{imperativo}{%
\begin{tabular}{@{}l l@{}}
\textit{tu} & \irr{riesci} \\
\textit{Lei} & \irr{riesca} \\
\textit{noi} & riusciamo \\
\textit{voi} & riuscite \\
\textit{Loro} & \irr{riescano} \\
\end{tabular}}{\irr{Highlighted} forms differ from the regular pattern of the verb's family. Negative \textit{tu}: \textit{non riuscire}. \textit{Lei} and \textit{Loro} are the congiuntivo forms.}
\end{minipage}\par\medskip

\moodheading{forme non finite}
{\small\begin{tabular}{@{}l l@{}}
infinito presente & \textit{riuscire} \\
infinito passato & \textit{essere riuscito} \\
gerundio presente & \textit{riuscendo} \\
gerundio passato & \textit{essendo riuscito} \\
participio passato & \textit{riuscito} \\
\end{tabular}}\par
\clearpage

\subsection{morire --- to die}\label{conj:morire}\index{morire@morire}
{\small Auxiliary: \textit{essere}. The stressed o becomes uo (muoio, muori, muore, muoiono). Futuro morirò, also morrò (less common).}\par

\moodheading{indicativo}

\par\noindent\begin{minipage}[t]{0.32\linewidth}\raggedright
\paradigmbox{indicativo presente}{%
\begin{tabular}{@{}l l@{}}
\textit{io} & \irr{muoio} \\
\textit{tu} & \irr{muori} \\
\textit{lui/lei} & \irr{muore} \\
\textit{noi} & moriamo \\
\textit{voi} & morite \\
\textit{loro} & \irr{muoiono} \\
\end{tabular}}{\irr{Highlighted} forms differ from the regular pattern of the verb's family.}
\end{minipage}\hspace{0.02\linewidth}\begin{minipage}[t]{0.32\linewidth}\raggedright
\paradigmbox{indicativo passato prossimo (with essere, masculine subject)}{%
\begin{tabular}{@{}l l@{}}
\textit{io} & sono morto \\
\textit{tu} & sei morto \\
\textit{lui/lei} & è morto \\
\textit{noi} & siamo morti \\
\textit{voi} & siete morti \\
\textit{loro} & sono morti \\
\end{tabular}}{}
\end{minipage}\hspace{0.02\linewidth}\begin{minipage}[t]{0.32\linewidth}\raggedright
\paradigmbox{indicativo imperfetto}{%
\begin{tabular}{@{}l l@{}}
\textit{io} & morivo \\
\textit{tu} & morivi \\
\textit{lui/lei} & moriva \\
\textit{noi} & morivamo \\
\textit{voi} & morivate \\
\textit{loro} & morivano \\
\end{tabular}}{}
\end{minipage}\par\medskip\noindent\begin{minipage}[t]{0.32\linewidth}\raggedright
\paradigmbox{indicativo trapassato prossimo (with essere, masculine subject)}{%
\begin{tabular}{@{}l l@{}}
\textit{io} & ero morto \\
\textit{tu} & eri morto \\
\textit{lui/lei} & era morto \\
\textit{noi} & eravamo morti \\
\textit{voi} & eravate morti \\
\textit{loro} & erano morti \\
\end{tabular}}{}
\end{minipage}\hspace{0.02\linewidth}\begin{minipage}[t]{0.32\linewidth}\raggedright
\paradigmbox{indicativo passato remoto}{%
\begin{tabular}{@{}l l@{}}
\textit{io} & morii \\
\textit{tu} & moristi \\
\textit{lui/lei} & morì \\
\textit{noi} & morimmo \\
\textit{voi} & moriste \\
\textit{loro} & morirono \\
\end{tabular}}{}
\end{minipage}\hspace{0.02\linewidth}\begin{minipage}[t]{0.32\linewidth}\raggedright
\paradigmbox{indicativo trapassato remoto (with essere, masculine subject)}{%
\begin{tabular}{@{}l l@{}}
\textit{io} & fui morto \\
\textit{tu} & fosti morto \\
\textit{lui/lei} & fu morto \\
\textit{noi} & fummo morti \\
\textit{voi} & foste morti \\
\textit{loro} & furono morti \\
\end{tabular}}{}
\end{minipage}\par\medskip\noindent\begin{minipage}[t]{0.32\linewidth}\raggedright
\paradigmbox{indicativo futuro semplice}{%
\begin{tabular}{@{}l l@{}}
\textit{io} & morirò / morrò \\
\textit{tu} & morirai / morrai \\
\textit{lui/lei} & morirà / morrà \\
\textit{noi} & moriremo / morremo \\
\textit{voi} & morirete / morrete \\
\textit{loro} & moriranno / morranno \\
\end{tabular}}{Where two forms are given, both are standard.}
\end{minipage}\hspace{0.02\linewidth}\begin{minipage}[t]{0.32\linewidth}\raggedright
\paradigmbox{indicativo futuro anteriore (with essere, masculine subject)}{%
\begin{tabular}{@{}l l@{}}
\textit{io} & sarò morto \\
\textit{tu} & sarai morto \\
\textit{lui/lei} & sarà morto \\
\textit{noi} & saremo morti \\
\textit{voi} & sarete morti \\
\textit{loro} & saranno morti \\
\end{tabular}}{}
\end{minipage}\par\medskip

\moodheading{congiuntivo}

\par\noindent\begin{minipage}[t]{0.32\linewidth}\raggedright
\paradigmbox{congiuntivo presente}{%
\begin{tabular}{@{}l l@{}}
\textit{io} & \irr{muoia} \\
\textit{tu} & \irr{muoia} \\
\textit{lui/lei} & \irr{muoia} \\
\textit{noi} & moriamo \\
\textit{voi} & moriate \\
\textit{loro} & \irr{muoiano} \\
\end{tabular}}{\irr{Highlighted} forms differ from the regular pattern of the verb's family.}
\end{minipage}\hspace{0.02\linewidth}\begin{minipage}[t]{0.32\linewidth}\raggedright
\paradigmbox{congiuntivo passato (with essere, masculine subject)}{%
\begin{tabular}{@{}l l@{}}
\textit{io} & sia morto \\
\textit{tu} & sia morto \\
\textit{lui/lei} & sia morto \\
\textit{noi} & siamo morti \\
\textit{voi} & siate morti \\
\textit{loro} & siano morti \\
\end{tabular}}{}
\end{minipage}\hspace{0.02\linewidth}\begin{minipage}[t]{0.32\linewidth}\raggedright
\paradigmbox{congiuntivo imperfetto}{%
\begin{tabular}{@{}l l@{}}
\textit{io} & morissi \\
\textit{tu} & morissi \\
\textit{lui/lei} & morisse \\
\textit{noi} & morissimo \\
\textit{voi} & moriste \\
\textit{loro} & morissero \\
\end{tabular}}{}
\end{minipage}\par\medskip\noindent\begin{minipage}[t]{0.32\linewidth}\raggedright
\paradigmbox{congiuntivo trapassato (with essere, masculine subject)}{%
\begin{tabular}{@{}l l@{}}
\textit{io} & fossi morto \\
\textit{tu} & fossi morto \\
\textit{lui/lei} & fosse morto \\
\textit{noi} & fossimo morti \\
\textit{voi} & foste morti \\
\textit{loro} & fossero morti \\
\end{tabular}}{}
\end{minipage}\par\medskip

\moodheading{condizionale}

\par\noindent\begin{minipage}[t]{0.32\linewidth}\raggedright
\paradigmbox{condizionale presente}{%
\begin{tabular}{@{}l l@{}}
\textit{io} & morirei / morrei \\
\textit{tu} & moriresti / morresti \\
\textit{lui/lei} & morirebbe / morrebbe \\
\textit{noi} & moriremmo / morremmo \\
\textit{voi} & morireste / morreste \\
\textit{loro} & morirebbero / morrebbero \\
\end{tabular}}{Where two forms are given, both are standard.}
\end{minipage}\hspace{0.02\linewidth}\begin{minipage}[t]{0.32\linewidth}\raggedright
\paradigmbox{condizionale passato (with essere, masculine subject)}{%
\begin{tabular}{@{}l l@{}}
\textit{io} & sarei morto \\
\textit{tu} & saresti morto \\
\textit{lui/lei} & sarebbe morto \\
\textit{noi} & saremmo morti \\
\textit{voi} & sareste morti \\
\textit{loro} & sarebbero morti \\
\end{tabular}}{}
\end{minipage}\par\medskip

\moodheading{imperativo}

\par\noindent\begin{minipage}[t]{0.32\linewidth}\raggedright
\paradigmbox{imperativo}{%
\begin{tabular}{@{}l l@{}}
\textit{tu} & \irr{muori} \\
\textit{Lei} & \irr{muoia} \\
\textit{noi} & moriamo \\
\textit{voi} & morite \\
\textit{Loro} & \irr{muoiano} \\
\end{tabular}}{\irr{Highlighted} forms differ from the regular pattern of the verb's family. Negative \textit{tu}: \textit{non morire}. \textit{Lei} and \textit{Loro} are the congiuntivo forms.}
\end{minipage}\par\medskip

\moodheading{forme non finite}
{\small\begin{tabular}{@{}l l@{}}
infinito presente & \textit{morire} \\
infinito passato & \textit{essere morto} \\
gerundio presente & \textit{morendo} \\
gerundio passato & \textit{essendo morto} \\
participio passato & \textit{morto} \\
\end{tabular}}\par
\clearpage

\subsection{spegnere --- to switch off; to put out}\label{conj:spegnere}\index{spegnere@spegnere}
{\small Auxiliary: \textit{avere}. gn becomes ng before -o: spengo, spengono (like tengo, vengo). Participle spento; passato remoto spensi.}\par

\moodheading{indicativo}

\par\noindent\begin{minipage}[t]{0.32\linewidth}\raggedright
\paradigmbox{indicativo presente}{%
\begin{tabular}{@{}l l@{}}
\textit{io} & \irr{spengo} \\
\textit{tu} & spegni \\
\textit{lui/lei} & spegne \\
\textit{noi} & spegniamo \\
\textit{voi} & spegnete \\
\textit{loro} & \irr{spengono} \\
\end{tabular}}{\irr{Highlighted} forms differ from the regular pattern of the verb's family.}
\end{minipage}\hspace{0.02\linewidth}\begin{minipage}[t]{0.32\linewidth}\raggedright
\paradigmbox{indicativo passato prossimo (with avere)}{%
\begin{tabular}{@{}l l@{}}
\textit{io} & ho spento \\
\textit{tu} & hai spento \\
\textit{lui/lei} & ha spento \\
\textit{noi} & abbiamo spento \\
\textit{voi} & avete spento \\
\textit{loro} & hanno spento \\
\end{tabular}}{}
\end{minipage}\hspace{0.02\linewidth}\begin{minipage}[t]{0.32\linewidth}\raggedright
\paradigmbox{indicativo imperfetto}{%
\begin{tabular}{@{}l l@{}}
\textit{io} & spegnevo \\
\textit{tu} & spegnevi \\
\textit{lui/lei} & spegneva \\
\textit{noi} & spegnevamo \\
\textit{voi} & spegnevate \\
\textit{loro} & spegnevano \\
\end{tabular}}{}
\end{minipage}\par\medskip\noindent\begin{minipage}[t]{0.32\linewidth}\raggedright
\paradigmbox{indicativo trapassato prossimo (with avere)}{%
\begin{tabular}{@{}l l@{}}
\textit{io} & avevo spento \\
\textit{tu} & avevi spento \\
\textit{lui/lei} & aveva spento \\
\textit{noi} & avevamo spento \\
\textit{voi} & avevate spento \\
\textit{loro} & avevano spento \\
\end{tabular}}{}
\end{minipage}\hspace{0.02\linewidth}\begin{minipage}[t]{0.32\linewidth}\raggedright
\paradigmbox{indicativo passato remoto}{%
\begin{tabular}{@{}l l@{}}
\textit{io} & \irr{spensi} \\
\textit{tu} & spegnesti \\
\textit{lui/lei} & \irr{spense} \\
\textit{noi} & spegnemmo \\
\textit{voi} & spegneste \\
\textit{loro} & \irr{spensero} \\
\end{tabular}}{\irr{Highlighted} forms differ from the regular pattern of the verb's family.}
\end{minipage}\hspace{0.02\linewidth}\begin{minipage}[t]{0.32\linewidth}\raggedright
\paradigmbox{indicativo trapassato remoto (with avere)}{%
\begin{tabular}{@{}l l@{}}
\textit{io} & ebbi spento \\
\textit{tu} & avesti spento \\
\textit{lui/lei} & ebbe spento \\
\textit{noi} & avemmo spento \\
\textit{voi} & aveste spento \\
\textit{loro} & ebbero spento \\
\end{tabular}}{}
\end{minipage}\par\medskip\noindent\begin{minipage}[t]{0.32\linewidth}\raggedright
\paradigmbox{indicativo futuro semplice}{%
\begin{tabular}{@{}l l@{}}
\textit{io} & spegnerò \\
\textit{tu} & spegnerai \\
\textit{lui/lei} & spegnerà \\
\textit{noi} & spegneremo \\
\textit{voi} & spegnerete \\
\textit{loro} & spegneranno \\
\end{tabular}}{}
\end{minipage}\hspace{0.02\linewidth}\begin{minipage}[t]{0.32\linewidth}\raggedright
\paradigmbox{indicativo futuro anteriore (with avere)}{%
\begin{tabular}{@{}l l@{}}
\textit{io} & avrò spento \\
\textit{tu} & avrai spento \\
\textit{lui/lei} & avrà spento \\
\textit{noi} & avremo spento \\
\textit{voi} & avrete spento \\
\textit{loro} & avranno spento \\
\end{tabular}}{}
\end{minipage}\par\medskip

\moodheading{congiuntivo}

\par\noindent\begin{minipage}[t]{0.32\linewidth}\raggedright
\paradigmbox{congiuntivo presente}{%
\begin{tabular}{@{}l l@{}}
\textit{io} & \irr{spenga} \\
\textit{tu} & \irr{spenga} \\
\textit{lui/lei} & \irr{spenga} \\
\textit{noi} & spegniamo \\
\textit{voi} & spegniate \\
\textit{loro} & \irr{spengano} \\
\end{tabular}}{\irr{Highlighted} forms differ from the regular pattern of the verb's family.}
\end{minipage}\hspace{0.02\linewidth}\begin{minipage}[t]{0.32\linewidth}\raggedright
\paradigmbox{congiuntivo passato (with avere)}{%
\begin{tabular}{@{}l l@{}}
\textit{io} & abbia spento \\
\textit{tu} & abbia spento \\
\textit{lui/lei} & abbia spento \\
\textit{noi} & abbiamo spento \\
\textit{voi} & abbiate spento \\
\textit{loro} & abbiano spento \\
\end{tabular}}{}
\end{minipage}\hspace{0.02\linewidth}\begin{minipage}[t]{0.32\linewidth}\raggedright
\paradigmbox{congiuntivo imperfetto}{%
\begin{tabular}{@{}l l@{}}
\textit{io} & spegnessi \\
\textit{tu} & spegnessi \\
\textit{lui/lei} & spegnesse \\
\textit{noi} & spegnessimo \\
\textit{voi} & spegneste \\
\textit{loro} & spegnessero \\
\end{tabular}}{}
\end{minipage}\par\medskip\noindent\begin{minipage}[t]{0.32\linewidth}\raggedright
\paradigmbox{congiuntivo trapassato (with avere)}{%
\begin{tabular}{@{}l l@{}}
\textit{io} & avessi spento \\
\textit{tu} & avessi spento \\
\textit{lui/lei} & avesse spento \\
\textit{noi} & avessimo spento \\
\textit{voi} & aveste spento \\
\textit{loro} & avessero spento \\
\end{tabular}}{}
\end{minipage}\par\medskip

\moodheading{condizionale}

\par\noindent\begin{minipage}[t]{0.32\linewidth}\raggedright
\paradigmbox{condizionale presente}{%
\begin{tabular}{@{}l l@{}}
\textit{io} & spegnerei \\
\textit{tu} & spegneresti \\
\textit{lui/lei} & spegnerebbe \\
\textit{noi} & spegneremmo \\
\textit{voi} & spegnereste \\
\textit{loro} & spegnerebbero \\
\end{tabular}}{}
\end{minipage}\hspace{0.02\linewidth}\begin{minipage}[t]{0.32\linewidth}\raggedright
\paradigmbox{condizionale passato (with avere)}{%
\begin{tabular}{@{}l l@{}}
\textit{io} & avrei spento \\
\textit{tu} & avresti spento \\
\textit{lui/lei} & avrebbe spento \\
\textit{noi} & avremmo spento \\
\textit{voi} & avreste spento \\
\textit{loro} & avrebbero spento \\
\end{tabular}}{}
\end{minipage}\par\medskip

\moodheading{imperativo}

\par\noindent\begin{minipage}[t]{0.32\linewidth}\raggedright
\paradigmbox{imperativo}{%
\begin{tabular}{@{}l l@{}}
\textit{tu} & spegni \\
\textit{Lei} & \irr{spenga} \\
\textit{noi} & spegniamo \\
\textit{voi} & spegnete \\
\textit{Loro} & \irr{spengano} \\
\end{tabular}}{\irr{Highlighted} forms differ from the regular pattern of the verb's family. Negative \textit{tu}: \textit{non spegnere}. \textit{Lei} and \textit{Loro} are the congiuntivo forms.}
\end{minipage}\par\medskip

\moodheading{forme non finite}
{\small\begin{tabular}{@{}l l@{}}
infinito presente & \textit{spegnere} \\
infinito passato & \textit{avere spento} \\
gerundio presente & \textit{spegnendo} \\
gerundio passato & \textit{avendo spento} \\
participio passato & \textit{spento} \\
\end{tabular}}\par
\clearpage

\section{Impersonal verbs}

\subsection{piovere --- to rain}\label{conj:piovere}\index{piovere@piovere}
{\small Auxiliary: \textit{essere}. Weather verb used without a subject: piove = it is raining. Both è piovuto and ha piovuto are correct. Regular in form; listed here because only the third person is normally used.}\par

\moodheading{indicativo}

\par\noindent\begin{minipage}[t]{0.32\linewidth}\raggedright
\paradigmbox{indicativo presente}{%
\begin{tabular}{@{}l l@{}}
\textit{io} & --- \\
\textit{tu} & --- \\
\textit{lui/lei} & piove \\
\textit{noi} & --- \\
\textit{voi} & --- \\
\textit{loro} & piovono \\
\end{tabular}}{}
\end{minipage}\hspace{0.02\linewidth}\begin{minipage}[t]{0.32\linewidth}\raggedright
\paradigmbox{indicativo passato prossimo (with essere, masculine subject)}{%
\begin{tabular}{@{}l l@{}}
\textit{io} & --- \\
\textit{tu} & --- \\
\textit{lui/lei} & è piovuto \\
\textit{noi} & --- \\
\textit{voi} & --- \\
\textit{loro} & sono piovuti \\
\end{tabular}}{}
\end{minipage}\hspace{0.02\linewidth}\begin{minipage}[t]{0.32\linewidth}\raggedright
\paradigmbox{indicativo imperfetto}{%
\begin{tabular}{@{}l l@{}}
\textit{io} & --- \\
\textit{tu} & --- \\
\textit{lui/lei} & pioveva \\
\textit{noi} & --- \\
\textit{voi} & --- \\
\textit{loro} & piovevano \\
\end{tabular}}{}
\end{minipage}\par\medskip\noindent\begin{minipage}[t]{0.32\linewidth}\raggedright
\paradigmbox{indicativo trapassato prossimo (with essere, masculine subject)}{%
\begin{tabular}{@{}l l@{}}
\textit{io} & --- \\
\textit{tu} & --- \\
\textit{lui/lei} & era piovuto \\
\textit{noi} & --- \\
\textit{voi} & --- \\
\textit{loro} & erano piovuti \\
\end{tabular}}{}
\end{minipage}\hspace{0.02\linewidth}\begin{minipage}[t]{0.32\linewidth}\raggedright
\paradigmbox{indicativo passato remoto}{%
\begin{tabular}{@{}l l@{}}
\textit{io} & --- \\
\textit{tu} & --- \\
\textit{lui/lei} & \irr{piovve} \\
\textit{noi} & --- \\
\textit{voi} & --- \\
\textit{loro} & \irr{piovvero} \\
\end{tabular}}{\irr{Highlighted} forms differ from the regular pattern of the verb's family.}
\end{minipage}\hspace{0.02\linewidth}\begin{minipage}[t]{0.32\linewidth}\raggedright
\paradigmbox{indicativo trapassato remoto (with essere, masculine subject)}{%
\begin{tabular}{@{}l l@{}}
\textit{io} & --- \\
\textit{tu} & --- \\
\textit{lui/lei} & fu piovuto \\
\textit{noi} & --- \\
\textit{voi} & --- \\
\textit{loro} & furono piovuti \\
\end{tabular}}{}
\end{minipage}\par\medskip\noindent\begin{minipage}[t]{0.32\linewidth}\raggedright
\paradigmbox{indicativo futuro semplice}{%
\begin{tabular}{@{}l l@{}}
\textit{io} & --- \\
\textit{tu} & --- \\
\textit{lui/lei} & pioverà \\
\textit{noi} & --- \\
\textit{voi} & --- \\
\textit{loro} & pioveranno \\
\end{tabular}}{}
\end{minipage}\hspace{0.02\linewidth}\begin{minipage}[t]{0.32\linewidth}\raggedright
\paradigmbox{indicativo futuro anteriore (with essere, masculine subject)}{%
\begin{tabular}{@{}l l@{}}
\textit{io} & --- \\
\textit{tu} & --- \\
\textit{lui/lei} & sarà piovuto \\
\textit{noi} & --- \\
\textit{voi} & --- \\
\textit{loro} & saranno piovuti \\
\end{tabular}}{}
\end{minipage}\par\medskip

\moodheading{congiuntivo}

\par\noindent\begin{minipage}[t]{0.32\linewidth}\raggedright
\paradigmbox{congiuntivo presente}{%
\begin{tabular}{@{}l l@{}}
\textit{io} & --- \\
\textit{tu} & --- \\
\textit{lui/lei} & piova \\
\textit{noi} & --- \\
\textit{voi} & --- \\
\textit{loro} & piovano \\
\end{tabular}}{}
\end{minipage}\hspace{0.02\linewidth}\begin{minipage}[t]{0.32\linewidth}\raggedright
\paradigmbox{congiuntivo passato (with essere, masculine subject)}{%
\begin{tabular}{@{}l l@{}}
\textit{io} & --- \\
\textit{tu} & --- \\
\textit{lui/lei} & sia piovuto \\
\textit{noi} & --- \\
\textit{voi} & --- \\
\textit{loro} & siano piovuti \\
\end{tabular}}{}
\end{minipage}\hspace{0.02\linewidth}\begin{minipage}[t]{0.32\linewidth}\raggedright
\paradigmbox{congiuntivo imperfetto}{%
\begin{tabular}{@{}l l@{}}
\textit{io} & --- \\
\textit{tu} & --- \\
\textit{lui/lei} & piovesse \\
\textit{noi} & --- \\
\textit{voi} & --- \\
\textit{loro} & piovessero \\
\end{tabular}}{}
\end{minipage}\par\medskip\noindent\begin{minipage}[t]{0.32\linewidth}\raggedright
\paradigmbox{congiuntivo trapassato (with essere, masculine subject)}{%
\begin{tabular}{@{}l l@{}}
\textit{io} & --- \\
\textit{tu} & --- \\
\textit{lui/lei} & fosse piovuto \\
\textit{noi} & --- \\
\textit{voi} & --- \\
\textit{loro} & fossero piovuti \\
\end{tabular}}{}
\end{minipage}\par\medskip

\moodheading{condizionale}

\par\noindent\begin{minipage}[t]{0.32\linewidth}\raggedright
\paradigmbox{condizionale presente}{%
\begin{tabular}{@{}l l@{}}
\textit{io} & --- \\
\textit{tu} & --- \\
\textit{lui/lei} & pioverebbe \\
\textit{noi} & --- \\
\textit{voi} & --- \\
\textit{loro} & pioverebbero \\
\end{tabular}}{}
\end{minipage}\hspace{0.02\linewidth}\begin{minipage}[t]{0.32\linewidth}\raggedright
\paradigmbox{condizionale passato (with essere, masculine subject)}{%
\begin{tabular}{@{}l l@{}}
\textit{io} & --- \\
\textit{tu} & --- \\
\textit{lui/lei} & sarebbe piovuto \\
\textit{noi} & --- \\
\textit{voi} & --- \\
\textit{loro} & sarebbero piovuti \\
\end{tabular}}{}
\end{minipage}\par\medskip

\moodheading{imperativo}

\par\noindent\begin{minipage}[t]{0.32\linewidth}\raggedright
\paradigmbox{imperativo}{%
\begin{tabular}{@{}l l@{}}
\textit{tu} & --- \\
\textit{Lei} & --- \\
\textit{noi} & --- \\
\textit{voi} & --- \\
\textit{Loro} & --- \\
\end{tabular}}{This verb has no imperative in normal use.}
\end{minipage}\par\medskip

\moodheading{forme non finite}
{\small\begin{tabular}{@{}l l@{}}
infinito presente & \textit{piovere} \\
infinito passato & \textit{essere piovuto} \\
gerundio presente & \textit{piovendo} \\
gerundio passato & \textit{essendo piovuto} \\
participio passato & \textit{piovuto} \\
\end{tabular}}\par
\clearpage

\subsection{nevicare --- to snow}\label{conj:nevicare}\index{nevicare@nevicare}
{\small Auxiliary: \textit{essere}. Weather verb: nevica = it is snowing. Both è nevicato and ha nevicato are correct.}\par

\moodheading{indicativo}

\par\noindent\begin{minipage}[t]{0.32\linewidth}\raggedright
\paradigmbox{indicativo presente}{%
\begin{tabular}{@{}l l@{}}
\textit{io} & --- \\
\textit{tu} & --- \\
\textit{lui/lei} & nevica \\
\textit{noi} & --- \\
\textit{voi} & --- \\
\textit{loro} & --- \\
\end{tabular}}{}
\end{minipage}\hspace{0.02\linewidth}\begin{minipage}[t]{0.32\linewidth}\raggedright
\paradigmbox{indicativo passato prossimo (with essere, masculine subject)}{%
\begin{tabular}{@{}l l@{}}
\textit{io} & --- \\
\textit{tu} & --- \\
\textit{lui/lei} & è nevicato \\
\textit{noi} & --- \\
\textit{voi} & --- \\
\textit{loro} & --- \\
\end{tabular}}{}
\end{minipage}\hspace{0.02\linewidth}\begin{minipage}[t]{0.32\linewidth}\raggedright
\paradigmbox{indicativo imperfetto}{%
\begin{tabular}{@{}l l@{}}
\textit{io} & --- \\
\textit{tu} & --- \\
\textit{lui/lei} & nevicava \\
\textit{noi} & --- \\
\textit{voi} & --- \\
\textit{loro} & --- \\
\end{tabular}}{}
\end{minipage}\par\medskip\noindent\begin{minipage}[t]{0.32\linewidth}\raggedright
\paradigmbox{indicativo trapassato prossimo (with essere, masculine subject)}{%
\begin{tabular}{@{}l l@{}}
\textit{io} & --- \\
\textit{tu} & --- \\
\textit{lui/lei} & era nevicato \\
\textit{noi} & --- \\
\textit{voi} & --- \\
\textit{loro} & --- \\
\end{tabular}}{}
\end{minipage}\hspace{0.02\linewidth}\begin{minipage}[t]{0.32\linewidth}\raggedright
\paradigmbox{indicativo passato remoto}{%
\begin{tabular}{@{}l l@{}}
\textit{io} & --- \\
\textit{tu} & --- \\
\textit{lui/lei} & nevicò \\
\textit{noi} & --- \\
\textit{voi} & --- \\
\textit{loro} & --- \\
\end{tabular}}{}
\end{minipage}\hspace{0.02\linewidth}\begin{minipage}[t]{0.32\linewidth}\raggedright
\paradigmbox{indicativo trapassato remoto (with essere, masculine subject)}{%
\begin{tabular}{@{}l l@{}}
\textit{io} & --- \\
\textit{tu} & --- \\
\textit{lui/lei} & fu nevicato \\
\textit{noi} & --- \\
\textit{voi} & --- \\
\textit{loro} & --- \\
\end{tabular}}{}
\end{minipage}\par\medskip\noindent\begin{minipage}[t]{0.32\linewidth}\raggedright
\paradigmbox{indicativo futuro semplice}{%
\begin{tabular}{@{}l l@{}}
\textit{io} & --- \\
\textit{tu} & --- \\
\textit{lui/lei} & nevicherà \\
\textit{noi} & --- \\
\textit{voi} & --- \\
\textit{loro} & --- \\
\end{tabular}}{}
\end{minipage}\hspace{0.02\linewidth}\begin{minipage}[t]{0.32\linewidth}\raggedright
\paradigmbox{indicativo futuro anteriore (with essere, masculine subject)}{%
\begin{tabular}{@{}l l@{}}
\textit{io} & --- \\
\textit{tu} & --- \\
\textit{lui/lei} & sarà nevicato \\
\textit{noi} & --- \\
\textit{voi} & --- \\
\textit{loro} & --- \\
\end{tabular}}{}
\end{minipage}\par\medskip

\moodheading{congiuntivo}

\par\noindent\begin{minipage}[t]{0.32\linewidth}\raggedright
\paradigmbox{congiuntivo presente}{%
\begin{tabular}{@{}l l@{}}
\textit{io} & --- \\
\textit{tu} & --- \\
\textit{lui/lei} & nevichi \\
\textit{noi} & --- \\
\textit{voi} & --- \\
\textit{loro} & --- \\
\end{tabular}}{}
\end{minipage}\hspace{0.02\linewidth}\begin{minipage}[t]{0.32\linewidth}\raggedright
\paradigmbox{congiuntivo passato (with essere, masculine subject)}{%
\begin{tabular}{@{}l l@{}}
\textit{io} & --- \\
\textit{tu} & --- \\
\textit{lui/lei} & sia nevicato \\
\textit{noi} & --- \\
\textit{voi} & --- \\
\textit{loro} & --- \\
\end{tabular}}{}
\end{minipage}\hspace{0.02\linewidth}\begin{minipage}[t]{0.32\linewidth}\raggedright
\paradigmbox{congiuntivo imperfetto}{%
\begin{tabular}{@{}l l@{}}
\textit{io} & --- \\
\textit{tu} & --- \\
\textit{lui/lei} & nevicasse \\
\textit{noi} & --- \\
\textit{voi} & --- \\
\textit{loro} & --- \\
\end{tabular}}{}
\end{minipage}\par\medskip\noindent\begin{minipage}[t]{0.32\linewidth}\raggedright
\paradigmbox{congiuntivo trapassato (with essere, masculine subject)}{%
\begin{tabular}{@{}l l@{}}
\textit{io} & --- \\
\textit{tu} & --- \\
\textit{lui/lei} & fosse nevicato \\
\textit{noi} & --- \\
\textit{voi} & --- \\
\textit{loro} & --- \\
\end{tabular}}{}
\end{minipage}\par\medskip

\moodheading{condizionale}

\par\noindent\begin{minipage}[t]{0.32\linewidth}\raggedright
\paradigmbox{condizionale presente}{%
\begin{tabular}{@{}l l@{}}
\textit{io} & --- \\
\textit{tu} & --- \\
\textit{lui/lei} & nevicherebbe \\
\textit{noi} & --- \\
\textit{voi} & --- \\
\textit{loro} & --- \\
\end{tabular}}{}
\end{minipage}\hspace{0.02\linewidth}\begin{minipage}[t]{0.32\linewidth}\raggedright
\paradigmbox{condizionale passato (with essere, masculine subject)}{%
\begin{tabular}{@{}l l@{}}
\textit{io} & --- \\
\textit{tu} & --- \\
\textit{lui/lei} & sarebbe nevicato \\
\textit{noi} & --- \\
\textit{voi} & --- \\
\textit{loro} & --- \\
\end{tabular}}{}
\end{minipage}\par\medskip

\moodheading{imperativo}

\par\noindent\begin{minipage}[t]{0.32\linewidth}\raggedright
\paradigmbox{imperativo}{%
\begin{tabular}{@{}l l@{}}
\textit{tu} & --- \\
\textit{Lei} & --- \\
\textit{noi} & --- \\
\textit{voi} & --- \\
\textit{Loro} & --- \\
\end{tabular}}{This verb has no imperative in normal use.}
\end{minipage}\par\medskip

\moodheading{forme non finite}
{\small\begin{tabular}{@{}l l@{}}
infinito presente & \textit{nevicare} \\
infinito passato & \textit{essere nevicato} \\
gerundio presente & \textit{nevicando} \\
gerundio passato & \textit{essendo nevicato} \\
participio passato & \textit{nevicato} \\
\end{tabular}}\par
\clearpage

\subsection{bisognare --- to be necessary (impersonal)}\label{conj:bisognare}\index{bisognare@bisognare}
{\small Auxiliary: \textit{essere}. Impersonal: bisogna + infinitive (bisogna partire) or bisogna che + congiuntivo (bisogna che tu parta). Compound tenses are rare.}\par

\moodheading{indicativo}

\par\noindent\begin{minipage}[t]{0.32\linewidth}\raggedright
\paradigmbox{indicativo presente}{%
\begin{tabular}{@{}l l@{}}
\textit{io} & --- \\
\textit{tu} & --- \\
\textit{lui/lei} & bisogna \\
\textit{noi} & --- \\
\textit{voi} & --- \\
\textit{loro} & --- \\
\end{tabular}}{}
\end{minipage}\hspace{0.02\linewidth}\begin{minipage}[t]{0.32\linewidth}\raggedright
\paradigmbox{indicativo passato prossimo (with essere, masculine subject)}{%
\begin{tabular}{@{}l l@{}}
\textit{io} & --- \\
\textit{tu} & --- \\
\textit{lui/lei} & è bisognato \\
\textit{noi} & --- \\
\textit{voi} & --- \\
\textit{loro} & --- \\
\end{tabular}}{}
\end{minipage}\hspace{0.02\linewidth}\begin{minipage}[t]{0.32\linewidth}\raggedright
\paradigmbox{indicativo imperfetto}{%
\begin{tabular}{@{}l l@{}}
\textit{io} & --- \\
\textit{tu} & --- \\
\textit{lui/lei} & bisognava \\
\textit{noi} & --- \\
\textit{voi} & --- \\
\textit{loro} & --- \\
\end{tabular}}{}
\end{minipage}\par\medskip\noindent\begin{minipage}[t]{0.32\linewidth}\raggedright
\paradigmbox{indicativo trapassato prossimo (with essere, masculine subject)}{%
\begin{tabular}{@{}l l@{}}
\textit{io} & --- \\
\textit{tu} & --- \\
\textit{lui/lei} & era bisognato \\
\textit{noi} & --- \\
\textit{voi} & --- \\
\textit{loro} & --- \\
\end{tabular}}{}
\end{minipage}\hspace{0.02\linewidth}\begin{minipage}[t]{0.32\linewidth}\raggedright
\paradigmbox{indicativo passato remoto}{%
\begin{tabular}{@{}l l@{}}
\textit{io} & --- \\
\textit{tu} & --- \\
\textit{lui/lei} & bisognò \\
\textit{noi} & --- \\
\textit{voi} & --- \\
\textit{loro} & --- \\
\end{tabular}}{}
\end{minipage}\hspace{0.02\linewidth}\begin{minipage}[t]{0.32\linewidth}\raggedright
\paradigmbox{indicativo trapassato remoto (with essere, masculine subject)}{%
\begin{tabular}{@{}l l@{}}
\textit{io} & --- \\
\textit{tu} & --- \\
\textit{lui/lei} & fu bisognato \\
\textit{noi} & --- \\
\textit{voi} & --- \\
\textit{loro} & --- \\
\end{tabular}}{}
\end{minipage}\par\medskip\noindent\begin{minipage}[t]{0.32\linewidth}\raggedright
\paradigmbox{indicativo futuro semplice}{%
\begin{tabular}{@{}l l@{}}
\textit{io} & --- \\
\textit{tu} & --- \\
\textit{lui/lei} & bisognerà \\
\textit{noi} & --- \\
\textit{voi} & --- \\
\textit{loro} & --- \\
\end{tabular}}{}
\end{minipage}\hspace{0.02\linewidth}\begin{minipage}[t]{0.32\linewidth}\raggedright
\paradigmbox{indicativo futuro anteriore (with essere, masculine subject)}{%
\begin{tabular}{@{}l l@{}}
\textit{io} & --- \\
\textit{tu} & --- \\
\textit{lui/lei} & sarà bisognato \\
\textit{noi} & --- \\
\textit{voi} & --- \\
\textit{loro} & --- \\
\end{tabular}}{}
\end{minipage}\par\medskip

\moodheading{congiuntivo}

\par\noindent\begin{minipage}[t]{0.32\linewidth}\raggedright
\paradigmbox{congiuntivo presente}{%
\begin{tabular}{@{}l l@{}}
\textit{io} & --- \\
\textit{tu} & --- \\
\textit{lui/lei} & bisogni \\
\textit{noi} & --- \\
\textit{voi} & --- \\
\textit{loro} & --- \\
\end{tabular}}{}
\end{minipage}\hspace{0.02\linewidth}\begin{minipage}[t]{0.32\linewidth}\raggedright
\paradigmbox{congiuntivo passato (with essere, masculine subject)}{%
\begin{tabular}{@{}l l@{}}
\textit{io} & --- \\
\textit{tu} & --- \\
\textit{lui/lei} & sia bisognato \\
\textit{noi} & --- \\
\textit{voi} & --- \\
\textit{loro} & --- \\
\end{tabular}}{}
\end{minipage}\hspace{0.02\linewidth}\begin{minipage}[t]{0.32\linewidth}\raggedright
\paradigmbox{congiuntivo imperfetto}{%
\begin{tabular}{@{}l l@{}}
\textit{io} & --- \\
\textit{tu} & --- \\
\textit{lui/lei} & bisognasse \\
\textit{noi} & --- \\
\textit{voi} & --- \\
\textit{loro} & --- \\
\end{tabular}}{}
\end{minipage}\par\medskip\noindent\begin{minipage}[t]{0.32\linewidth}\raggedright
\paradigmbox{congiuntivo trapassato (with essere, masculine subject)}{%
\begin{tabular}{@{}l l@{}}
\textit{io} & --- \\
\textit{tu} & --- \\
\textit{lui/lei} & fosse bisognato \\
\textit{noi} & --- \\
\textit{voi} & --- \\
\textit{loro} & --- \\
\end{tabular}}{}
\end{minipage}\par\medskip

\moodheading{condizionale}

\par\noindent\begin{minipage}[t]{0.32\linewidth}\raggedright
\paradigmbox{condizionale presente}{%
\begin{tabular}{@{}l l@{}}
\textit{io} & --- \\
\textit{tu} & --- \\
\textit{lui/lei} & bisognerebbe \\
\textit{noi} & --- \\
\textit{voi} & --- \\
\textit{loro} & --- \\
\end{tabular}}{}
\end{minipage}\hspace{0.02\linewidth}\begin{minipage}[t]{0.32\linewidth}\raggedright
\paradigmbox{condizionale passato (with essere, masculine subject)}{%
\begin{tabular}{@{}l l@{}}
\textit{io} & --- \\
\textit{tu} & --- \\
\textit{lui/lei} & sarebbe bisognato \\
\textit{noi} & --- \\
\textit{voi} & --- \\
\textit{loro} & --- \\
\end{tabular}}{}
\end{minipage}\par\medskip

\moodheading{imperativo}

\par\noindent\begin{minipage}[t]{0.32\linewidth}\raggedright
\paradigmbox{imperativo}{%
\begin{tabular}{@{}l l@{}}
\textit{tu} & --- \\
\textit{Lei} & --- \\
\textit{noi} & --- \\
\textit{voi} & --- \\
\textit{Loro} & --- \\
\end{tabular}}{This verb has no imperative in normal use.}
\end{minipage}\par\medskip

\moodheading{forme non finite}
{\small\begin{tabular}{@{}l l@{}}
infinito presente & \textit{bisognare} \\
infinito passato & \textit{essere bisognato} \\
gerundio presente & \textit{bisognando} \\
gerundio passato & \textit{essendo bisognato} \\
participio passato & \textit{bisognato} \\
\end{tabular}}\par
\clearpage

\chapter{Principal parts of every verb in the course}\label{ref:principal}
For each verb: auxiliary, the \textit{io} and \textit{noi} forms of the presente, the \textit{io} and \textit{lui/lei} forms of the passato remoto, the \textit{io} form of the futuro and of the congiuntivo presente, and the past participle. Every other form can be built from these (see Chapter~\ref{ref:system}). Irregular forms are printed in \irr{colour}.\par

{\footnotesize\setlength{\tabcolsep}{3pt}
\begin{xltabular}{\linewidth}{@{}>{\raggedright\arraybackslash}p{2.7cm} l L L l >{\raggedright\arraybackslash}p{1.9cm} l@{}}
\toprule
\textbf{Infinitive} & \textbf{Auxiliary} & \textbf{Presente} & \textbf{Passato remoto} & \textbf{Futuro} & \textbf{Congiuntivo presente} & \textbf{Participle} \\
\midrule
\endhead
\textbf{abitare}\index{abitare@abitare}\newline {\color{muted}to live (reside)} & avere & \textit{abito, abitiamo} & \textit{abitai, abitò} & \textit{abiterò} & \textit{abiti} & \textit{abitato} \\
\textbf{accendere}\index{accendere@accendere}\newline {\color{muted}to switch on; to light} & avere & \textit{accendo, accendiamo} & \textit{\irr{accesi}, \irr{accese}} & \textit{accenderò} & \textit{accenda} & \textit{acceso} \\
\textbf{addormentarsi}\index{addormentarsi@addormentarsi}\newline {\color{muted}to fall asleep} & essere & \textit{mi addormento, ci addormentiamo} & \textit{mi addormentai, si addormentò} & \textit{mi addormenterò} & \textit{mi addormenti} & \textit{addormentato} \\
\textbf{aiutare}\index{aiutare@aiutare}\newline {\color{muted}to help} & avere & \textit{aiuto, aiutiamo} & \textit{aiutai, aiutò} & \textit{aiuterò} & \textit{aiuti} & \textit{aiutato} \\
\textbf{alzarsi}\index{alzarsi@alzarsi}\newline {\color{muted}to get up} & essere & \textit{mi alzo, ci alziamo} & \textit{mi alzai, si alzò} & \textit{mi alzerò} & \textit{mi alzi} & \textit{alzato} \\
\textbf{amare}\index{amare@amare}\newline {\color{muted}to love} & avere & \textit{amo, amiamo} & \textit{amai, amò} & \textit{amerò} & \textit{ami} & \textit{amato} \\
\textbf{andare}\index{andare@andare}\newline {\color{muted}to go} & essere & \textit{\irr{vado}, andiamo} & \textit{andai, andò} & \textit{\irr{andrò}} & \textit{\irr{vada}} & \textit{andato} \\
\textbf{annunciare}\index{annunciare@annunciare}\newline {\color{muted}to announce} & avere & \textit{annuncio, annunciamo} & \textit{annunciai, annunciò} & \textit{annuncerò} & \textit{annunci} & \textit{annunciato} \\
\textbf{aprire}\index{aprire@aprire}\newline {\color{muted}to open} & avere & \textit{apro, apriamo} & \textit{aprii, aprì} & \textit{aprirò} & \textit{apra} & \textit{aperto} \\
\textbf{arrivare}\index{arrivare@arrivare}\newline {\color{muted}to arrive} & essere & \textit{arrivo, arriviamo} & \textit{arrivai, arrivò} & \textit{arriverò} & \textit{arrivi} & \textit{arrivato} \\
\textbf{ascoltare}\index{ascoltare@ascoltare}\newline {\color{muted}to listen to} & avere & \textit{ascolto, ascoltiamo} & \textit{ascoltai, ascoltò} & \textit{ascolterò} & \textit{ascolti} & \textit{ascoltato} \\
\textbf{aspettare}\index{aspettare@aspettare}\newline {\color{muted}to wait (for)} & avere & \textit{aspetto, aspettiamo} & \textit{aspettai, aspettò} & \textit{aspetterò} & \textit{aspetti} & \textit{aspettato} \\
\textbf{avere}\index{avere@avere}\newline {\color{muted}to have} & avere & \textit{\irr{ho}, \irr{abbiamo}} & \textit{\irr{ebbi}, \irr{ebbe}} & \textit{\irr{avrò}} & \textit{\irr{abbia}} & \textit{avuto} \\
\textbf{ballare}\index{ballare@ballare}\newline {\color{muted}to dance} & avere & \textit{ballo, balliamo} & \textit{ballai, ballò} & \textit{ballerò} & \textit{balli} & \textit{ballato} \\
\textbf{bere}\index{bere@bere}\newline {\color{muted}to drink} & avere & \textit{\irr{bevo}, \irr{beviamo}} & \textit{\irr{bevvi}, \irr{bevve}} & \textit{\irr{berrò}} & \textit{\irr{beva}} & \textit{bevuto} \\
\textbf{bisognare}\index{bisognare@bisognare}\newline {\color{muted}to be necessary (impersonal)} & essere & \textit{bisogna} & \textit{bisognò} & \textit{bisognerà} & \textit{bisogni} & \textit{bisognato} \\
\textbf{brillare}\index{brillare@brillare}\newline {\color{muted}to shine; to sparkle} & avere & \textit{brillo, brilliamo} & \textit{brillai, brillò} & \textit{brillerò} & \textit{brilli} & \textit{brillato} \\
\textbf{bussare}\index{bussare@bussare}\newline {\color{muted}to knock} & avere & \textit{busso, bussiamo} & \textit{bussai, bussò} & \textit{busserò} & \textit{bussi} & \textit{bussato} \\
\textbf{cadere}\index{cadere@cadere}\newline {\color{muted}to fall} & essere & \textit{cado, cadiamo} & \textit{\irr{caddi}, \irr{cadde}} & \textit{\irr{cadrò}} & \textit{cada} & \textit{caduto} \\
\textbf{cambiare}\index{cambiare@cambiare}\newline {\color{muted}to change} & avere & \textit{cambio, cambiamo} & \textit{cambiai, cambiò} & \textit{cambierò} & \textit{cambi} & \textit{cambiato} \\
\textbf{camminare}\index{camminare@camminare}\newline {\color{muted}to walk} & avere & \textit{cammino, camminiamo} & \textit{camminai, camminò} & \textit{camminerò} & \textit{cammini} & \textit{camminato} \\
\textbf{cantare}\index{cantare@cantare}\newline {\color{muted}to sing} & avere & \textit{canto, cantiamo} & \textit{cantai, cantò} & \textit{canterò} & \textit{canti} & \textit{cantato} \\
\textbf{capire}\index{capire@capire}\newline {\color{muted}to understand} & avere & \textit{capisco, capiamo} & \textit{capii, capì} & \textit{capirò} & \textit{capisca} & \textit{capito} \\
\textbf{cercare}\index{cercare@cercare}\newline {\color{muted}to look for; to try} & avere & \textit{cerco, cerchiamo} & \textit{cercai, cercò} & \textit{cercherò} & \textit{cerchi} & \textit{cercato} \\
\textbf{chiamare}\index{chiamare@chiamare}\newline {\color{muted}to call} & avere & \textit{chiamo, chiamiamo} & \textit{chiamai, chiamò} & \textit{chiamerò} & \textit{chiami} & \textit{chiamato} \\
\textbf{chiamarsi}\index{chiamarsi@chiamarsi}\newline {\color{muted}to be called} & essere & \textit{mi chiamo, ci chiamiamo} & \textit{mi chiamai, si chiamò} & \textit{mi chiamerò} & \textit{mi chiami} & \textit{chiamato} \\
\textbf{chiedere}\index{chiedere@chiedere}\newline {\color{muted}to ask} & avere & \textit{chiedo, chiediamo} & \textit{\irr{chiesi}, \irr{chiese}} & \textit{chiederò} & \textit{chieda} & \textit{chiesto} \\
\textbf{chiudere}\index{chiudere@chiudere}\newline {\color{muted}to close} & avere & \textit{chiudo, chiudiamo} & \textit{\irr{chiusi}, \irr{chiuse}} & \textit{chiuderò} & \textit{chiuda} & \textit{chiuso} \\
\textbf{cominciare}\index{cominciare@cominciare}\newline {\color{muted}to begin; to start} & avere & \textit{comincio, cominciamo} & \textit{cominciai, cominciò} & \textit{comincerò} & \textit{cominci} & \textit{cominciato} \\
\textbf{comprare}\index{comprare@comprare}\newline {\color{muted}to buy} & avere & \textit{compro, compriamo} & \textit{comprai, comprò} & \textit{comprerò} & \textit{compri} & \textit{comprato} \\
\textbf{conoscere}\index{conoscere@conoscere}\newline {\color{muted}to know (a person or place)} & avere & \textit{conosco, conosciamo} & \textit{\irr{conobbi}, \irr{conobbe}} & \textit{conoscerò} & \textit{conosca} & \textit{conosciuto} \\
\textbf{continuare}\index{continuare@continuare}\newline {\color{muted}to continue} & avere & \textit{continuo, continuiamo} & \textit{continuai, continuò} & \textit{continuerò} & \textit{continui} & \textit{continuato} \\
\textbf{correre}\index{correre@correre}\newline {\color{muted}to run} & avere & \textit{corro, corriamo} & \textit{\irr{corsi}, \irr{corse}} & \textit{correrò} & \textit{corra} & \textit{corso} \\
\textbf{costare}\index{costare@costare}\newline {\color{muted}to cost} & essere & \textit{costo, costiamo} & \textit{costai, costò} & \textit{costerò} & \textit{costi} & \textit{costato} \\
\textbf{costruire}\index{costruire@costruire}\newline {\color{muted}to build} & avere & \textit{costruisco, costruiamo} & \textit{costruii, costruì} & \textit{costruirò} & \textit{costruisca} & \textit{costruito} \\
\textbf{credere}\index{credere@credere}\newline {\color{muted}to believe} & avere & \textit{credo, crediamo} & \textit{credei / credetti, credé / credette} & \textit{crederò} & \textit{creda} & \textit{creduto} \\
\textbf{crescere}\index{crescere@crescere}\newline {\color{muted}to grow} & essere & \textit{cresco, cresciamo} & \textit{\irr{crebbi}, \irr{crebbe}} & \textit{crescerò} & \textit{cresca} & \textit{cresciuto} \\
\textbf{cucinare}\index{cucinare@cucinare}\newline {\color{muted}to cook} & avere & \textit{cucino, cuciniamo} & \textit{cucinai, cucinò} & \textit{cucinerò} & \textit{cucini} & \textit{cucinato} \\
\textbf{dare}\index{dare@dare}\newline {\color{muted}to give} & avere & \textit{do, diamo} & \textit{\irr{diedi / detti}, \irr{diede / dette}} & \textit{\irr{darò}} & \textit{\irr{dia}} & \textit{dato} \\
\textbf{decidere}\index{decidere@decidere}\newline {\color{muted}to decide} & avere & \textit{decido, decidiamo} & \textit{\irr{decisi}, \irr{decise}} & \textit{deciderò} & \textit{decida} & \textit{deciso} \\
\textbf{dimenticare}\index{dimenticare@dimenticare}\newline {\color{muted}to forget} & avere & \textit{dimentico, dimentichiamo} & \textit{dimenticai, dimenticò} & \textit{dimenticherò} & \textit{dimentichi} & \textit{dimenticato} \\
\textbf{dire}\index{dire@dire}\newline {\color{muted}to say; to tell} & avere & \textit{\irr{dico}, \irr{diciamo}} & \textit{\irr{dissi}, \irr{disse}} & \textit{dirò} & \textit{\irr{dica}} & \textit{detto} \\
\textbf{divertirsi}\index{divertirsi@divertirsi}\newline {\color{muted}to have fun; to enjoy oneself} & essere & \textit{mi diverto, ci divertiamo} & \textit{mi divertii, si divertì} & \textit{mi divertirò} & \textit{mi diverta} & \textit{divertito} \\
\textbf{dormire}\index{dormire@dormire}\newline {\color{muted}to sleep} & avere & \textit{dormo, dormiamo} & \textit{dormii, dormì} & \textit{dormirò} & \textit{dorma} & \textit{dormito} \\
\textbf{dovere}\index{dovere@dovere}\newline {\color{muted}to have to; must} & avere & \textit{\irr{devo}, \irr{dobbiamo}} & \textit{dovei / dovetti, dové / dovette} & \textit{\irr{dovrò}} & \textit{\irr{debba / deva}} & \textit{dovuto} \\
\textbf{dubitare}\index{dubitare@dubitare}\newline {\color{muted}to doubt} & avere & \textit{dubito, dubitiamo} & \textit{dubitai, dubitò} & \textit{dubiterò} & \textit{dubiti} & \textit{dubitato} \\
\textbf{entrare}\index{entrare@entrare}\newline {\color{muted}to enter; to go in} & essere & \textit{entro, entriamo} & \textit{entrai, entrò} & \textit{entrerò} & \textit{entri} & \textit{entrato} \\
\textbf{essere}\index{essere@essere}\newline {\color{muted}to be} & essere & \textit{\irr{sono}, \irr{siamo}} & \textit{\irr{fui}, \irr{fu}} & \textit{\irr{sarò}} & \textit{\irr{sia}} & \textit{stato} \\
\textbf{fare}\index{fare@fare}\newline {\color{muted}to do; to make} & avere & \textit{\irr{faccio}, \irr{facciamo}} & \textit{\irr{feci}, \irr{fece}} & \textit{\irr{farò}} & \textit{\irr{faccia}} & \textit{fatto} \\
\textbf{fermarsi}\index{fermarsi@fermarsi}\newline {\color{muted}to stop (oneself)} & essere & \textit{mi fermo, ci fermiamo} & \textit{mi fermai, si fermò} & \textit{mi fermerò} & \textit{mi fermi} & \textit{fermato} \\
\textbf{finire}\index{finire@finire}\newline {\color{muted}to finish; to end} & avere & \textit{finisco, finiamo} & \textit{finii, finì} & \textit{finirò} & \textit{finisca} & \textit{finito} \\
\textbf{giocare}\index{giocare@giocare}\newline {\color{muted}to play} & avere & \textit{gioco, giochiamo} & \textit{giocai, giocò} & \textit{giocherò} & \textit{giochi} & \textit{giocato} \\
\textbf{gridare}\index{gridare@gridare}\newline {\color{muted}to shout} & avere & \textit{grido, gridiamo} & \textit{gridai, gridò} & \textit{griderò} & \textit{gridi} & \textit{gridato} \\
\textbf{guardare}\index{guardare@guardare}\newline {\color{muted}to look at; to watch} & avere & \textit{guardo, guardiamo} & \textit{guardai, guardò} & \textit{guarderò} & \textit{guardi} & \textit{guardato} \\
\textbf{guidare}\index{guidare@guidare}\newline {\color{muted}to drive; to guide} & avere & \textit{guido, guidiamo} & \textit{guidai, guidò} & \textit{guiderò} & \textit{guidi} & \textit{guidato} \\
\textbf{imparare}\index{imparare@imparare}\newline {\color{muted}to learn} & avere & \textit{imparo, impariamo} & \textit{imparai, imparò} & \textit{imparerò} & \textit{impari} & \textit{imparato} \\
\textbf{incontrare}\index{incontrare@incontrare}\newline {\color{muted}to meet} & avere & \textit{incontro, incontriamo} & \textit{incontrai, incontrò} & \textit{incontrerò} & \textit{incontri} & \textit{incontrato} \\
\textbf{insegnare}\index{insegnare@insegnare}\newline {\color{muted}to teach} & avere & \textit{insegno, insegniamo} & \textit{insegnai, insegnò} & \textit{insegnerò} & \textit{insegni} & \textit{insegnato} \\
\textbf{invitare}\index{invitare@invitare}\newline {\color{muted}to invite} & avere & \textit{invito, invitiamo} & \textit{invitai, invitò} & \textit{inviterò} & \textit{inviti} & \textit{invitato} \\
\textbf{lasciare}\index{lasciare@lasciare}\newline {\color{muted}to leave (something); to let} & avere & \textit{lascio, lasciamo} & \textit{lasciai, lasciò} & \textit{lascerò} & \textit{lasci} & \textit{lasciato} \\
\textbf{lavare}\index{lavare@lavare}\newline {\color{muted}to wash} & avere & \textit{lavo, laviamo} & \textit{lavai, lavò} & \textit{laverò} & \textit{lavi} & \textit{lavato} \\
\textbf{lavarsi}\index{lavarsi@lavarsi}\newline {\color{muted}to wash (oneself)} & essere & \textit{mi lavo, ci laviamo} & \textit{mi lavai, si lavò} & \textit{mi laverò} & \textit{mi lavi} & \textit{lavato} \\
\textbf{lavorare}\index{lavorare@lavorare}\newline {\color{muted}to work} & avere & \textit{lavoro, lavoriamo} & \textit{lavorai, lavorò} & \textit{lavorerò} & \textit{lavori} & \textit{lavorato} \\
\textbf{leggere}\index{leggere@leggere}\newline {\color{muted}to read} & avere & \textit{leggo, leggiamo} & \textit{\irr{lessi}, \irr{lesse}} & \textit{leggerò} & \textit{legga} & \textit{letto} \\
\textbf{mancare}\index{mancare@mancare}\newline {\color{muted}to be missing; to miss} & essere & \textit{manco, manchiamo} & \textit{mancai, mancò} & \textit{mancherò} & \textit{manchi} & \textit{mancato} \\
\textbf{mandare}\index{mandare@mandare}\newline {\color{muted}to send} & avere & \textit{mando, mandiamo} & \textit{mandai, mandò} & \textit{manderò} & \textit{mandi} & \textit{mandato} \\
\textbf{mangiare}\index{mangiare@mangiare}\newline {\color{muted}to eat} & avere & \textit{mangio, mangiamo} & \textit{mangiai, mangiò} & \textit{mangerò} & \textit{mangi} & \textit{mangiato} \\
\textbf{mettere}\index{mettere@mettere}\newline {\color{muted}to put} & avere & \textit{metto, mettiamo} & \textit{\irr{misi}, \irr{mise}} & \textit{metterò} & \textit{metta} & \textit{messo} \\
\textbf{morire}\index{morire@morire}\newline {\color{muted}to die} & essere & \textit{\irr{muoio}, moriamo} & \textit{morii, morì} & \textit{morirò / morrò} & \textit{\irr{muoia}} & \textit{morto} \\
\textbf{nascere}\index{nascere@nascere}\newline {\color{muted}to be born} & essere & \textit{nasco, nasciamo} & \textit{\irr{nacqui}, \irr{nacque}} & \textit{nascerò} & \textit{nasca} & \textit{nato} \\
\textbf{nascondere}\index{nascondere@nascondere}\newline {\color{muted}to hide} & avere & \textit{nascondo, nascondiamo} & \textit{\irr{nascosi}, \irr{nascose}} & \textit{nasconderò} & \textit{nasconda} & \textit{nascosto} \\
\textbf{nevicare}\index{nevicare@nevicare}\newline {\color{muted}to snow} & essere & \textit{nevica} & \textit{nevicò} & \textit{nevicherà} & \textit{nevichi} & \textit{nevicato} \\
\textbf{nuotare}\index{nuotare@nuotare}\newline {\color{muted}to swim} & avere & \textit{nuoto, nuotiamo} & \textit{nuotai, nuotò} & \textit{nuoterò} & \textit{nuoti} & \textit{nuotato} \\
\textbf{offrire}\index{offrire@offrire}\newline {\color{muted}to offer} & avere & \textit{offro, offriamo} & \textit{offrii, offrì} & \textit{offrirò} & \textit{offra} & \textit{offerto} \\
\textbf{pagare}\index{pagare@pagare}\newline {\color{muted}to pay (for)} & avere & \textit{pago, paghiamo} & \textit{pagai, pagò} & \textit{pagherò} & \textit{paghi} & \textit{pagato} \\
\textbf{parlare}\index{parlare@parlare}\newline {\color{muted}to speak; to talk} & avere & \textit{parlo, parliamo} & \textit{parlai, parlò} & \textit{parlerò} & \textit{parli} & \textit{parlato} \\
\textbf{partire}\index{partire@partire}\newline {\color{muted}to leave; to set off} & essere & \textit{parto, partiamo} & \textit{partii, partì} & \textit{partirò} & \textit{parta} & \textit{partito} \\
\textbf{passare}\index{passare@passare}\newline {\color{muted}to pass; to spend (time); to drop by} & essere & \textit{passo, passiamo} & \textit{passai, passò} & \textit{passerò} & \textit{passi} & \textit{passato} \\
\textbf{pensare}\index{pensare@pensare}\newline {\color{muted}to think} & avere & \textit{penso, pensiamo} & \textit{pensai, pensò} & \textit{penserò} & \textit{pensi} & \textit{pensato} \\
\textbf{perdere}\index{perdere@perdere}\newline {\color{muted}to lose; to miss} & avere & \textit{perdo, perdiamo} & \textit{\irr{persi}, \irr{perse}} & \textit{perderò} & \textit{perda} & \textit{perso} \\
\textbf{permettere}\index{permettere@permettere}\newline {\color{muted}to allow} & avere & \textit{permetto, permettiamo} & \textit{\irr{permisi}, \irr{permise}} & \textit{permetterò} & \textit{permetta} & \textit{permesso} \\
\textbf{pescare}\index{pescare@pescare}\newline {\color{muted}to fish; to catch (fish)} & avere & \textit{pesco, peschiamo} & \textit{pescai, pescò} & \textit{pescherò} & \textit{peschi} & \textit{pescato} \\
\textbf{piacere}\index{piacere@piacere}\newline {\color{muted}to please; to be liked} & essere & \textit{\irr{piaccio}, \irr{piacciamo}} & \textit{\irr{piacqui}, \irr{piacque}} & \textit{piacerò} & \textit{\irr{piaccia}} & \textit{piaciuto} \\
\textbf{piangere}\index{piangere@piangere}\newline {\color{muted}to cry} & avere & \textit{piango, piangiamo} & \textit{\irr{piansi}, \irr{pianse}} & \textit{piangerò} & \textit{pianga} & \textit{pianto} \\
\textbf{piovere}\index{piovere@piovere}\newline {\color{muted}to rain} & essere & \textit{piove} & \textit{\irr{piovve}} & \textit{pioverà} & \textit{piova} & \textit{piovuto} \\
\textbf{portare}\index{portare@portare}\newline {\color{muted}to bring; to carry; to wear} & avere & \textit{porto, portiamo} & \textit{portai, portò} & \textit{porterò} & \textit{porti} & \textit{portato} \\
\textbf{potere}\index{potere@potere}\newline {\color{muted}to be able to; can; may} & avere & \textit{\irr{posso}, \irr{possiamo}} & \textit{potei / potetti, poté / potette} & \textit{\irr{potrò}} & \textit{\irr{possa}} & \textit{potuto} \\
\textbf{preferire}\index{preferire@preferire}\newline {\color{muted}to prefer} & avere & \textit{preferisco, preferiamo} & \textit{preferii, preferì} & \textit{preferirò} & \textit{preferisca} & \textit{preferito} \\
\textbf{prendere}\index{prendere@prendere}\newline {\color{muted}to take; to catch; to have (food)} & avere & \textit{prendo, prendiamo} & \textit{\irr{presi}, \irr{prese}} & \textit{prenderò} & \textit{prenda} & \textit{preso} \\
\textbf{preparare}\index{preparare@preparare}\newline {\color{muted}to prepare; to make} & avere & \textit{preparo, prepariamo} & \textit{preparai, preparò} & \textit{preparerò} & \textit{prepari} & \textit{preparato} \\
\textbf{promettere}\index{promettere@promettere}\newline {\color{muted}to promise} & avere & \textit{prometto, promettiamo} & \textit{\irr{promisi}, \irr{promise}} & \textit{prometterò} & \textit{prometta} & \textit{promesso} \\
\textbf{pulire}\index{pulire@pulire}\newline {\color{muted}to clean} & avere & \textit{pulisco, puliamo} & \textit{pulii, pulì} & \textit{pulirò} & \textit{pulisca} & \textit{pulito} \\
\textbf{raccogliere}\index{raccogliere@raccogliere}\newline {\color{muted}to pick (up); to gather} & avere & \textit{\irr{raccolgo}, raccogliamo} & \textit{\irr{raccolsi}, \irr{raccolse}} & \textit{raccoglierò} & \textit{\irr{raccolga}} & \textit{raccolto} \\
\textbf{raccontare}\index{raccontare@raccontare}\newline {\color{muted}to tell (a story)} & avere & \textit{racconto, raccontiamo} & \textit{raccontai, raccontò} & \textit{racconterò} & \textit{racconti} & \textit{raccontato} \\
\textbf{restare}\index{restare@restare}\newline {\color{muted}to stay; to remain} & essere & \textit{resto, restiamo} & \textit{restai, restò} & \textit{resterò} & \textit{resti} & \textit{restato} \\
\textbf{ricevere}\index{ricevere@ricevere}\newline {\color{muted}to receive} & avere & \textit{ricevo, riceviamo} & \textit{ricevei / ricevetti, ricevé / ricevette} & \textit{riceverò} & \textit{riceva} & \textit{ricevuto} \\
\textbf{ricordare}\index{ricordare@ricordare}\newline {\color{muted}to remember; to remind} & avere & \textit{ricordo, ricordiamo} & \textit{ricordai, ricordò} & \textit{ricorderò} & \textit{ricordi} & \textit{ricordato} \\
\textbf{ridere}\index{ridere@ridere}\newline {\color{muted}to laugh} & avere & \textit{rido, ridiamo} & \textit{\irr{risi}, \irr{rise}} & \textit{riderò} & \textit{rida} & \textit{riso} \\
\textbf{rimanere}\index{rimanere@rimanere}\newline {\color{muted}to remain; to stay} & essere & \textit{\irr{rimango}, rimaniamo} & \textit{\irr{rimasi}, \irr{rimase}} & \textit{\irr{rimarrò}} & \textit{\irr{rimanga}} & \textit{rimasto} \\
\textbf{rispondere}\index{rispondere@rispondere}\newline {\color{muted}to answer; to reply} & avere & \textit{rispondo, rispondiamo} & \textit{\irr{risposi}, \irr{rispose}} & \textit{risponderò} & \textit{risponda} & \textit{risposto} \\
\textbf{riuscire}\index{riuscire@riuscire}\newline {\color{muted}to manage; to succeed} & essere & \textit{\irr{riesco}, riusciamo} & \textit{riuscii, riuscì} & \textit{riuscirò} & \textit{\irr{riesca}} & \textit{riuscito} \\
\textbf{rompere}\index{rompere@rompere}\newline {\color{muted}to break} & avere & \textit{rompo, rompiamo} & \textit{\irr{ruppi}, \irr{ruppe}} & \textit{romperò} & \textit{rompa} & \textit{rotto} \\
\textbf{salire}\index{salire@salire}\newline {\color{muted}to go up; to climb; to get on} & essere & \textit{\irr{salgo}, saliamo} & \textit{salii, salì} & \textit{salirò} & \textit{\irr{salga}} & \textit{salito} \\
\textbf{saltare}\index{saltare@saltare}\newline {\color{muted}to jump} & avere & \textit{salto, saltiamo} & \textit{saltai, saltò} & \textit{salterò} & \textit{salti} & \textit{saltato} \\
\textbf{salutare}\index{salutare@salutare}\newline {\color{muted}to greet; to say goodbye} & avere & \textit{saluto, salutiamo} & \textit{salutai, salutò} & \textit{saluterò} & \textit{saluti} & \textit{salutato} \\
\textbf{sapere}\index{sapere@sapere}\newline {\color{muted}to know (a fact); to know how to} & avere & \textit{\irr{so}, \irr{sappiamo}} & \textit{\irr{seppi}, \irr{seppe}} & \textit{\irr{saprò}} & \textit{\irr{sappia}} & \textit{saputo} \\
\textbf{scappare}\index{scappare@scappare}\newline {\color{muted}to run away; to escape} & essere & \textit{scappo, scappiamo} & \textit{scappai, scappò} & \textit{scapperò} & \textit{scappi} & \textit{scappato} \\
\textbf{scegliere}\index{scegliere@scegliere}\newline {\color{muted}to choose} & avere & \textit{\irr{scelgo}, scegliamo} & \textit{\irr{scelsi}, \irr{scelse}} & \textit{sceglierò} & \textit{\irr{scelga}} & \textit{scelto} \\
\textbf{scendere}\index{scendere@scendere}\newline {\color{muted}to go down; to get off} & essere & \textit{scendo, scendiamo} & \textit{\irr{scesi}, \irr{scese}} & \textit{scenderò} & \textit{scenda} & \textit{sceso} \\
\textbf{scoprire}\index{scoprire@scoprire}\newline {\color{muted}to discover; to uncover} & avere & \textit{scopro, scopriamo} & \textit{scoprii, scoprì} & \textit{scoprirò} & \textit{scopra} & \textit{scoperto} \\
\textbf{scrivere}\index{scrivere@scrivere}\newline {\color{muted}to write} & avere & \textit{scrivo, scriviamo} & \textit{\irr{scrissi}, \irr{scrisse}} & \textit{scriverò} & \textit{scriva} & \textit{scritto} \\
\textbf{sedersi}\index{sedersi@sedersi}\newline {\color{muted}to sit down} & essere & \textit{\irr{mi siedo}, ci sediamo} & \textit{mi sedei / mi sedetti, si sedé / si sedette} & \textit{mi sederò / mi siederò} & \textit{\irr{mi sieda}} & \textit{seduto} \\
\textbf{seguire}\index{seguire@seguire}\newline {\color{muted}to follow} & avere & \textit{seguo, seguiamo} & \textit{seguii, seguì} & \textit{seguirò} & \textit{segua} & \textit{seguito} \\
\textbf{sembrare}\index{sembrare@sembrare}\newline {\color{muted}to seem} & essere & \textit{sembro, sembriamo} & \textit{sembrai, sembrò} & \textit{sembrerò} & \textit{sembri} & \textit{sembrato} \\
\textbf{sentire}\index{sentire@sentire}\newline {\color{muted}to hear; to feel} & avere & \textit{sento, sentiamo} & \textit{sentii, sentì} & \textit{sentirò} & \textit{senta} & \textit{sentito} \\
\textbf{servire}\index{servire@servire}\newline {\color{muted}to serve; to be needed} & avere & \textit{servo, serviamo} & \textit{servii, servì} & \textit{servirò} & \textit{serva} & \textit{servito} \\
\textbf{smettere}\index{smettere@smettere}\newline {\color{muted}to stop (doing)} & avere & \textit{smetto, smettiamo} & \textit{\irr{smisi}, \irr{smise}} & \textit{smetterò} & \textit{smetta} & \textit{smesso} \\
\textbf{soffiare}\index{soffiare@soffiare}\newline {\color{muted}to blow} & avere & \textit{soffio, soffiamo} & \textit{soffiai, soffiò} & \textit{soffierò} & \textit{soffi} & \textit{soffiato} \\
\textbf{soffrire}\index{soffrire@soffrire}\newline {\color{muted}to suffer} & avere & \textit{soffro, soffriamo} & \textit{soffrii, soffrì} & \textit{soffrirò} & \textit{soffra} & \textit{sofferto} \\
\textbf{sognare}\index{sognare@sognare}\newline {\color{muted}to dream} & avere & \textit{sogno, sogniamo} & \textit{sognai, sognò} & \textit{sognerò} & \textit{sogni} & \textit{sognato} \\
\textbf{sorridere}\index{sorridere@sorridere}\newline {\color{muted}to smile} & avere & \textit{sorrido, sorridiamo} & \textit{\irr{sorrisi}, \irr{sorrise}} & \textit{sorriderò} & \textit{sorrida} & \textit{sorriso} \\
\textbf{spegnere}\index{spegnere@spegnere}\newline {\color{muted}to switch off; to put out} & avere & \textit{\irr{spengo}, spegniamo} & \textit{\irr{spensi}, \irr{spense}} & \textit{spegnerò} & \textit{\irr{spenga}} & \textit{spento} \\
\textbf{sperare}\index{sperare@sperare}\newline {\color{muted}to hope} & avere & \textit{spero, speriamo} & \textit{sperai, sperò} & \textit{spererò} & \textit{speri} & \textit{sperato} \\
\textbf{spiegare}\index{spiegare@spiegare}\newline {\color{muted}to explain} & avere & \textit{spiego, spieghiamo} & \textit{spiegai, spiegò} & \textit{spiegherò} & \textit{spieghi} & \textit{spiegato} \\
\textbf{stare}\index{stare@stare}\newline {\color{muted}to stay; to be (of health)} & essere & \textit{sto, stiamo} & \textit{\irr{stetti}, \irr{stette}} & \textit{\irr{starò}} & \textit{\irr{stia}} & \textit{stato} \\
\textbf{studiare}\index{studiare@studiare}\newline {\color{muted}to study} & avere & \textit{studio, studiamo} & \textit{studiai, studiò} & \textit{studierò} & \textit{studi} & \textit{studiato} \\
\textbf{succedere}\index{succedere@succedere}\newline {\color{muted}to happen} & essere & \textit{succedo, succediamo} & \textit{\irr{successi}, \irr{successe}} & \textit{succederò} & \textit{succeda} & \textit{successo} \\
\textbf{svegliarsi}\index{svegliarsi@svegliarsi}\newline {\color{muted}to wake up} & essere & \textit{mi sveglio, ci svegliamo} & \textit{mi svegliai, si svegliò} & \textit{mi sveglierò} & \textit{mi svegli} & \textit{svegliato} \\
\textbf{telefonare}\index{telefonare@telefonare}\newline {\color{muted}to phone} & avere & \textit{telefono, telefoniamo} & \textit{telefonai, telefonò} & \textit{telefonerò} & \textit{telefoni} & \textit{telefonato} \\
\textbf{temere}\index{temere@temere}\newline {\color{muted}to fear; to be afraid} & avere & \textit{temo, temiamo} & \textit{temei / temetti, temé / temette} & \textit{temerò} & \textit{tema} & \textit{temuto} \\
\textbf{tenere}\index{tenere@tenere}\newline {\color{muted}to hold; to keep} & avere & \textit{\irr{tengo}, teniamo} & \textit{\irr{tenni}, \irr{tenne}} & \textit{\irr{terrò}} & \textit{\irr{tenga}} & \textit{tenuto} \\
\textbf{tornare}\index{tornare@tornare}\newline {\color{muted}to return; to come back} & essere & \textit{torno, torniamo} & \textit{tornai, tornò} & \textit{tornerò} & \textit{torni} & \textit{tornato} \\
\textbf{trovare}\index{trovare@trovare}\newline {\color{muted}to find} & avere & \textit{trovo, troviamo} & \textit{trovai, trovò} & \textit{troverò} & \textit{trovi} & \textit{trovato} \\
\textbf{uscire}\index{uscire@uscire}\newline {\color{muted}to go out; to leave} & essere & \textit{\irr{esco}, usciamo} & \textit{uscii, uscì} & \textit{uscirò} & \textit{\irr{esca}} & \textit{uscito} \\
\textbf{vedere}\index{vedere@vedere}\newline {\color{muted}to see} & avere & \textit{vedo, vediamo} & \textit{\irr{vidi}, \irr{vide}} & \textit{\irr{vedrò}} & \textit{veda} & \textit{visto} \\
\textbf{vendere}\index{vendere@vendere}\newline {\color{muted}to sell} & avere & \textit{vendo, vendiamo} & \textit{vendei / vendetti, vendé / vendette} & \textit{venderò} & \textit{venda} & \textit{venduto} \\
\textbf{venire}\index{venire@venire}\newline {\color{muted}to come} & essere & \textit{\irr{vengo}, veniamo} & \textit{\irr{venni}, \irr{venne}} & \textit{\irr{verrò}} & \textit{\irr{venga}} & \textit{venuto} \\
\textbf{vestirsi}\index{vestirsi@vestirsi}\newline {\color{muted}to get dressed} & essere & \textit{mi vesto, ci vestiamo} & \textit{mi vestii, si vestì} & \textit{mi vestirò} & \textit{mi vesta} & \textit{vestito} \\
\textbf{viaggiare}\index{viaggiare@viaggiare}\newline {\color{muted}to travel} & avere & \textit{viaggio, viaggiamo} & \textit{viaggiai, viaggiò} & \textit{viaggerò} & \textit{viaggi} & \textit{viaggiato} \\
\textbf{vincere}\index{vincere@vincere}\newline {\color{muted}to win} & avere & \textit{vinco, vinciamo} & \textit{\irr{vinsi}, \irr{vinse}} & \textit{vincerò} & \textit{vinca} & \textit{vinto} \\
\textbf{visitare}\index{visitare@visitare}\newline {\color{muted}to visit (a place)} & avere & \textit{visito, visitiamo} & \textit{visitai, visitò} & \textit{visiterò} & \textit{visiti} & \textit{visitato} \\
\textbf{vivere}\index{vivere@vivere}\newline {\color{muted}to live} & avere & \textit{vivo, viviamo} & \textit{\irr{vissi}, \irr{visse}} & \textit{\irr{vivrò}} & \textit{viva} & \textit{vissuto} \\
\textbf{volare}\index{volare@volare}\newline {\color{muted}to fly} & avere & \textit{volo, voliamo} & \textit{volai, volò} & \textit{volerò} & \textit{voli} & \textit{volato} \\
\textbf{volere}\index{volere@volere}\newline {\color{muted}to want} & avere & \textit{\irr{voglio}, \irr{vogliamo}} & \textit{\irr{volli}, \irr{volle}} & \textit{\irr{vorrò}} & \textit{\irr{voglia}} & \textit{voluto} \\
\bottomrule
\end{xltabular}}
{\footnotesize\color{muted}Impersonal verbs (\textit{piovere, nevicare, bisognare}) are given in the third person singular only.}\par

\appendix
\addtocontents{toc}{\protect\addvspace{6pt}}

\chapter{Answer key}\label{answers}
Answers are given with the explanation shown on the website. Where the website accepts several answers, all of them are listed. Answers are compared as whole words: accents, articles, endings and double consonants count; capitalisation and final punctuation do not.\par

\section*{Chapter \ref{mod:m01}: Pronunciation and orthography}

\answerheading{\S\ref{les:alphabet-vowels} The alphabet and the vowels}\label{ans:alphabet-vowels}
\begin{answers}
\item\relax \textbf{(c) It is silent}\newline h is never pronounced in Italian. ho sounds exactly like o (or); the h only distinguishes the two in writing. {\color{muted}(See \S\ref{les:alphabet-vowels:alphabet}.)}
\item\relax \textbf{1--a, 2--d, 3--e, 4--b, 5--c {\color{muted}(h = acca; z = zeta; q = cu; r = erre; g = gi)}}\newline Consonant names follow two patterns: vowel + consonant + e (effe, elle, emme, enne, erre, esse) or consonant + i (bi, ci, di, gi, pi, ti). h (acca), q (cu) and z (zeta) are special. {\color{muted}(See \S\ref{les:alphabet-vowels:alphabet}.)}
\item\relax \textbf{(c) k}\newline j, k, w, x and y are not part of the native Italian alphabet; they appear in words such as il kiwi or il weekend. {\color{muted}(See \S\ref{les:alphabet-vowels:alphabet}.)}
\item\relax \textbf{(b) Number (singular or plural)}\newline gatto = cat, gatti = cats. Final vowels carry gender and number, which is why they are always pronounced clearly. {\color{muted}(See \S\ref{les:alphabet-vowels:vowels}.)}
\item\relax \textbf{(c) 3}\newline Every Italian vowel is pronounced: i-de-a. {\color{muted}(See \S\ref{les:alphabet-vowels:spelling}.)}
\item\relax \textbf{\textit{the little girl} {\color{muted}(also accepted: \textit{the girl}; \textit{the child}; \textit{the baby girl}; \textit{the young girl})}}\newline bambina ends in -a: a girl. bambino (-o) is a boy. {\color{muted}(See \S\ref{les:alphabet-vowels:vowels}.)}
\end{answers}

\answerheading{\S\ref{les:c-and-g} Hard and soft c and g; sc, gn and gli}\label{ans:c-and-g}
\begin{answers}
\item\relax \textbf{(a) chiave}\newline ch before i or e is hard (k). In cena and cielo the c is soft before e/i, and in ciao the i signals a soft c before a. {\color{muted}(See \S\ref{les:c-and-g:c-g}.)}
\item\relax \textbf{(d) giorno}\newline gi before o = soft g. gatto and gufo have g before a/u (hard); laghi uses gh to keep g hard before i. {\color{muted}(See \S\ref{les:c-and-g:c-g}.)}
\item\relax \textbf{(c) i laghi}\newline Plural -i would make g soft (lagi = laji). The h keeps the hard sound: laghi.\newline {\color{muted}Common mistake: \textit{i lagi} --- lagi would be pronounced 'laji'. Insert h to keep g hard before i: laghi.} {\color{muted}(See \S\ref{les:c-and-g:c-g}.)}
\item\relax \textbf{1--a, 2--d, 3--b, 4--c {\color{muted}(scuola = sk; pesce = sh; montagna = ny; famiglia = ly (palatal l))}}\newline sc + u = sk; sc + e = sh; gn = ny; gli = palatal l. {\color{muted}(See \S\ref{les:c-and-g:sc}.)}
\item\relax \textbf{(1) \textit{ch}; (2) \textit{gi}; (3) \textit{sc}; (4) \textit{sch}}\newline chiave (ch = k before i), giallo (gi = j before a), scuola (sc = sk before u), scherzo (sch = sk before e). {\color{muted}(See \S\ref{les:c-and-g:c-g}.)}
\item\relax \textbf{(c) As a w: kwando}\newline After q, u is always pronounced as a w sound. {\color{muted}(See \S\ref{les:c-and-g:gn-gli}.)}
\end{answers}

\answerheading{\S\ref{les:stress-doubles-accents} Syllables, stress, double consonants and accents}\label{ans:stress-doubles-accents}
\begin{answers}
\item\relax \textbf{\textit{è}}\newline è (with an accent) is the verb 'is'. e without an accent means 'and'.\newline {\color{muted}Common mistake: \textit{e} --- Without the accent, e means 'and'. 'Is' is è.}\newline {\color{muted}Common mistake: \textit{e'} --- An apostrophe is not an accent: write è.} {\color{muted}(See \S\ref{les:stress-doubles-accents:accents}.)}
\item\relax \textbf{(b) grandfather}\newline nonno (double n) = grandfather; nono (single n) = ninth. {\color{muted}(See \S\ref{les:stress-doubles-accents:doubles}.)}
\item\relax \textbf{(c) un'amica}\newline una elides to un' before a vowel, with no space: un'amica. Without the apostrophe (un amica) it would be the masculine form. {\color{muted}(See \S\ref{les:stress-doubles-accents:elision}.)}
\item\relax \textbf{(a) un amico}\newline Masculine un never takes an apostrophe: un amico, un albero.\newline {\color{muted}Common mistake: \textit{un'amico} --- Only the feminine una is elided (un'amica). Masculine un never takes an apostrophe.} {\color{muted}(See \S\ref{les:stress-doubles-accents:elision}.)}
\item\relax \textbf{(b) albero}\newline AL-be-ro. bam-BI-no and fi-NE-stra follow the usual penultimate stress; cit-TÀ has written final stress. {\color{muted}(See \S\ref{les:stress-doubles-accents:syllables-stress}.)}
\item\relax \textbf{(1) \textit{l'}; (2) \textit{c'è}; (3) \textit{dov'è}}\newline la → l' before a vowel; ci + è → c'è; dove + è → dov'è. {\color{muted}(See \S\ref{les:stress-doubles-accents:elision}.)}
\end{answers}

\section*{Chapter \ref{mod:m02}: Nouns and articles}

\answerheading{\S\ref{les:noun-gender-number} Noun gender and number}\label{ans:noun-gender-number}
\begin{answers}
\item\relax \textbf{\textit{libri}}\newline Masculine -o → -i. {\color{muted}(See \S\ref{les:noun-gender-number:plural}.)}
\item\relax \textbf{\textit{stelle}}\newline Feminine -a → -e. {\color{muted}(See \S\ref{les:noun-gender-number:plural}.)}
\item\relax \textbf{\textit{notti}}\newline Nouns in -e make their plural in -i, whatever their gender.\newline {\color{muted}Common mistake: \textit{notte} --- notte ends in -e, so its plural ends in -i (even though it is feminine): le notti.} {\color{muted}(See \S\ref{les:noun-gender-number:plural}.)}
\item\relax \textbf{\textit{città}}\newline Nouns with a stressed final vowel do not change: le città.\newline {\color{muted}Common mistake: \textit{citte} --- città ends in a stressed (accented) vowel, so it is invariable: le città.} {\color{muted}(See \S\ref{les:noun-gender-number:plural}.)}
\item\relax \textbf{(b) feminine}\newline la mano is a common exception: feminine despite -o. Plural: le mani. {\color{muted}(See \S\ref{les:noun-gender-number:gender}.)}
\item\relax \textbf{(a) Nothing certain: -e nouns can be masculine or feminine}\newline il fiore is masculine, la notte is feminine. Learn every -e noun with its article. {\color{muted}(See \S\ref{les:noun-gender-number:gender}.)}
\item\relax \textbf{1--b, 2--d, 3--c, 4--a {\color{muted}(il gatto = i gatti; la mela = le mele; il cane = i cani; il re = i re)}}\newline -o → -i, -a → -e, -e → -i; one-syllable re is invariable. {\color{muted}(See \S\ref{les:noun-gender-number:plural}.)}
\end{answers}

\answerheading{\S\ref{les:definite-articles} The definite article}\label{ans:definite-articles}
\begin{answers}
\item\relax \textbf{(c) lo}\newline Masculine nouns beginning with z take lo: lo zaino.\newline {\color{muted}Common mistake: \textit{il} --- il is used before most consonants, but z (like s + consonant, gn, ps) requires lo.} {\color{muted}(See \S\ref{les:definite-articles:forms}.)}
\item\relax \textbf{(b) l'}\newline Before a vowel, lo elides to l': l'albero. {\color{muted}(See \S\ref{les:definite-articles:forms}.)}
\item\relax \textbf{(a) le}\newline The feminine plural article is always le, even before a vowel.\newline {\color{muted}Common mistake: \textit{l'} --- l' is singular only. The feminine plural is always le: le api.} {\color{muted}(See \S\ref{les:definite-articles:forms}.)}
\item\relax \textbf{il cane: \textit{i cani}; lo specchio: \textit{gli specchi}; l'orso: \textit{gli orsi}; la foglia: \textit{le foglie}; l'isola: \textit{le isole}}\newline il → i; lo → gli; masculine l' → gli; la → le; feminine l' → le. {\color{muted}(See \S\ref{les:definite-articles:forms}.)}
\item\relax \textbf{(b) il}\newline The article is chosen by the word that immediately follows it: vecchio begins with v, so il vecchio zio, even though lo zio. (With a possessive and a singular family noun the article is normally dropped: mio zio — see the possessives lesson.) {\color{muted}(See \S\ref{les:definite-articles:forms}.)}
\item\relax \textbf{\textit{Le api fanno il miele.} {\color{muted}(also accepted: \textit{Le api fanno il miele})}}\newline General statements use the article in Italian: le api, il miele. {\color{muted}(See \S\ref{les:definite-articles:use}.)}
\end{answers}

\answerheading{\S\ref{les:indefinite-articles} The indefinite article; c'è and ci sono}\label{ans:indefinite-articles}
\begin{answers}
\item\relax \textbf{(b) uno}\newline z requires uno (compare lo zaino). {\color{muted}(See \S\ref{les:indefinite-articles:forms}.)}
\item\relax \textbf{(a) un'}\newline Feminine before a vowel: un'isola. {\color{muted}(See \S\ref{les:indefinite-articles:forms}.)}
\item\relax \textbf{(a) un}\newline Masculine before a vowel: un orso — never un'.\newline {\color{muted}Common mistake: \textit{un'} --- un' is only for feminine nouns. orso is masculine: un orso.} {\color{muted}(See \S\ref{les:indefinite-articles:forms}.)}
\item\relax \textbf{(1) \textit{c'è}; (2) \textit{ci sono}}\newline c'è + singular, ci sono + plural. {\color{muted}(See \S\ref{les:indefinite-articles:ce}.)}
\item\relax \textbf{\textit{There is a bee on the flower.} {\color{muted}(also accepted: \textit{There's a bee on the flower})}}\newline c'è = there is; un'ape = a bee; sul fiore = on the flower. {\color{muted}(See \S\ref{les:indefinite-articles:ce}.)}
\item\relax \textbf{\textit{Ci sono tre gatti.}}\newline ci sono + plural noun. {\color{muted}(See \S\ref{les:indefinite-articles:ce}.)}
\end{answers}

\answerheading{\S\ref{les:irregular-nouns} Spelling changes and irregular plurals}\label{ans:irregular-nouns}
\begin{answers}
\item\relax \textbf{\textit{barche}}\newline -ca → -che, always.\newline {\color{muted}Common mistake: \textit{barce} --- In barce the c before e would be soft, as in English 'church'. Keep the hard k sound with h: barche.} {\color{muted}(See \S\ref{les:irregular-nouns:ca-ga}.)}
\item\relax \textbf{\textit{amici}}\newline amico is an exception to the -co → -chi tendency.\newline {\color{muted}Common mistake: \textit{amichi} --- amico is a well-known exception: gli amici.} {\color{muted}(See \S\ref{les:irregular-nouns:co-go}.)}
\item\relax \textbf{\textit{giochi}}\newline GIO-co is stressed on the second-to-last syllable: giochi. {\color{muted}(See \S\ref{les:irregular-nouns:co-go}.)}
\item\relax \textbf{l'uomo: \textit{gli uomini}; l'uovo: \textit{le uova}; la mano: \textit{le mani}; lo zio: \textit{gli zii}; la spiaggia: \textit{le spiagge}}\newline uomini (irregular), le uova (feminine plural), le mani, zii (stressed ì), spiagge (-gia after a consonant). {\color{muted}(See \S\ref{les:irregular-nouns:irregular}.)}
\item\relax \textbf{(b) -ca always becomes -che, but -co becomes -chi or -ci depending largely on stress.}\newline The -ca rule has no exceptions; the -co pattern is a tendency with exceptions (amico, medico…). {\color{muted}(See \S\ref{les:irregular-nouns:co-go}.)}
\item\relax \textbf{\textit{Il gatto ha gli occhi verdi.}}\newline occhio → occhi (unstressed -io → -i); gli before a vowel. {\color{muted}(See \S\ref{les:irregular-nouns:cia-gia-io}.)}
\end{answers}

\section*{Chapter \ref{mod:m03}: Adjectives}

\answerheading{\S\ref{les:adjective-agreement} Adjectives: agreement and position}\label{ans:adjective-agreement}
\begin{answers}
\item\relax \textbf{\textit{rossa}}\newline Feminine singular of an -o adjective: rossa. {\color{muted}(See \S\ref{les:adjective-agreement:classes}.)}
\item\relax \textbf{\textit{fredde}}\newline notte is feminine; plural notti; feminine plural adjective: fredde.\newline {\color{muted}Common mistake: \textit{freddi} --- notti ends in -i but is feminine (la notte), so the adjective is feminine plural: fredde.} {\color{muted}(See \S\ref{les:adjective-agreement:classes}.)}
\item\relax \textbf{\textit{felici}}\newline -e adjectives have one plural form for both genders: felici. {\color{muted}(See \S\ref{les:adjective-agreement:classes}.)}
\item\relax \textbf{(d) le farfalle rosa}\newline rosa (the colour) is invariable. {\color{muted}(See \S\ref{les:adjective-agreement:classes}.)}
\item\relax \textbf{masculine singular: \textit{bianco}; feminine singular: \textit{bianca}; masculine plural: \textit{bianchi}; feminine plural: \textit{bianche}}\newline -co adjectives keep the hard sound: bianchi, bianche. {\color{muted}(See \S\ref{les:adjective-agreement:classes}.)}
\item\relax \textbf{(b) a great man}\newline Before the noun grande tends to mean 'great'; after it, 'big': un uomo grande. {\color{muted}(See \S\ref{les:adjective-agreement:position}.)}
\item\relax \textbf{\textit{Il gatto nero dorme}}\newline Colour adjectives follow the noun. {\color{muted}(See \S\ref{les:adjective-agreement:position}.)}
\end{answers}

\answerheading{\S\ref{les:demonstratives-possessives} Demonstratives and possessives}\label{ans:demonstratives-possessives}
\begin{answers}
\item\relax \textbf{(a) quegli}\newline alberi takes gli, so quello becomes quegli. {\color{muted}(See \S\ref{les:demonstratives-possessives:questo-quello}.)}
\item\relax \textbf{(d) quel}\newline il gatto → quel gatto. {\color{muted}(See \S\ref{les:demonstratives-possessives:questo-quello}.)}
\item\relax \textbf{(d) il suo}\newline libro is masculine, so il suo — Sara's gender does not matter.\newline {\color{muted}Common mistake: \textit{la sua} --- The possessive agrees with libro (masculine), not with Sara: il suo libro.} {\color{muted}(See \S\ref{les:demonstratives-possessives:possessives}.)}
\item\relax \textbf{(d) mio padre}\newline Singular family noun without adjective: no article. {\color{muted}(See \S\ref{les:demonstratives-possessives:possessives}.)}
\item\relax \textbf{(d) la loro casa}\newline loro is invariable and always takes the article. {\color{muted}(See \S\ref{les:demonstratives-possessives:possessives}.)}
\item\relax \textbf{\textit{I miei nonni abitano al mare.} {\color{muted}(also accepted: \textit{I miei nonni vivono al mare.})}}\newline Plural family noun: the article is used (i miei nonni). al mare = by the sea / at the seaside. {\color{muted}(See \S\ref{les:demonstratives-possessives:possessives}.)}
\end{answers}

\answerheading{\S\ref{les:bello-buono-comparison} bello and buono; comparing things}\label{ans:bello-buono-comparison}
\begin{answers}
\item\relax \textbf{(d) bell'}\newline l'albero → bell'albero. {\color{muted}(See \S\ref{les:bello-buono-comparison:bello-buono}.)}
\item\relax \textbf{(d) bel}\newline il giorno → bel giorno. {\color{muted}(See \S\ref{les:bello-buono-comparison:bello-buono}.)}
\item\relax \textbf{(b) belli}\newline After essere bello is a regular adjective: belli. bei/begli are only for before a noun. {\color{muted}(See \S\ref{les:bello-buono-comparison:bello-buono}.)}
\item\relax \textbf{(1) \textit{più}; (2) \textit{della}}\newline più … di; di + la = della. {\color{muted}(See \S\ref{les:bello-buono-comparison:comparison}.)}
\item\relax \textbf{\textit{Luca è alto come suo padre.} {\color{muted}(also accepted: \textit{Luca è così alto come suo padre.}; \textit{Luca è alto quanto suo padre.}; \textit{Luca è tanto alto quanto suo padre.})}}\newline (così) … come or (tanto) … quanto; suo padre without article. {\color{muted}(See \S\ref{les:bello-buono-comparison:comparison}.)}
\item\relax \textbf{(d) migliore}\newline migliore is the adjective (una torta migliore); meglio is the adverb (sto meglio). {\color{muted}(See \S\ref{les:bello-buono-comparison:comparison}.)}
\end{answers}

\section*{Chapter \ref{mod:m04}: Regular verbs in the present}

\answerheading{\S\ref{les:verb-system-overview} Map of the verb system}\label{ans:verb-system-overview}
\begin{answers}
\item\relax \textbf{1--b, 2--c, 3--a, 4--d {\color{muted}(indicativo = as a fact; congiuntivo = as something thought, wished or doubted; condizionale = as dependent on a condition; imperativo = as a command)}}\newline These are the central uses; each lesson adds the details and exceptions. {\color{muted}(See \S\ref{les:verb-system-overview:moods}.)}
\item\relax \textbf{(b) avere or essere + past participle}\newline The auxiliary carries person, tense and mood; the participle carries the meaning. stare + gerundio is a periphrasis, not a tense. {\color{muted}(See \S\ref{les:verb-system-overview:simple-compound}.)}
\item\relax \textbf{1--a, 2--c, 3--e, 4--b, 5--d {\color{muted}(ho / sono = passato prossimo; avevo / ero = trapassato prossimo; avrò / sarò = futuro anteriore; avrei / sarei = condizionale passato; abbia / sia = congiuntivo passato)}}\newline The compound tense takes its name from the simple tense of the auxiliary. {\color{muted}(See \S\ref{les:verb-system-overview:simple-compound}.)}
\item\relax \textbf{(b) 1st person plural, imperfetto, indicativo}\newline -ava- marks the imperfetto of -are verbs; -mo marks noi. {\color{muted}(See \S\ref{les:verb-system-overview:what-a-form-encodes}.)}
\item\relax \textbf{(c) parlando}\newline parlando is the gerundio. parlano (indicativo), parlino (congiuntivo) and parlerò (futuro) all show person and number. {\color{muted}(See \S\ref{les:verb-system-overview:moods}.)}
\item\relax \textbf{(c) tu (indicativo) or io/tu/lui/lei (congiuntivo presente)}\newline In the congiuntivo presente the three singular persons share one form; context (e.g. \textit{penso che …}) decides. {\color{muted}(See \S\ref{les:verb-system-overview:what-a-form-encodes}.)}
\end{answers}

\answerheading{\S\ref{les:infinitives-pronouns} Infinitives, stems and subject pronouns}\label{ans:infinitives-pronouns}
\begin{answers}
\item\relax \textbf{(b) cant-}\newline Remove -are: cant-. {\color{muted}(See \S\ref{les:infinitives-pronouns:families}.)}
\item\relax \textbf{1--b, 2--d, 3--a, 4--b, 5--d {\color{muted}(leggere = -ere; aprire = -ire; guardare = -are; credere = -ere; sentire = -ire)}}\newline The family is the last three letters of the infinitive. {\color{muted}(See \S\ref{les:infinitives-pronouns:families}.)}
\item\relax \textbf{1--c, 2--d, 3--b, 4--a, 5--e {\color{muted}(io = I; noi = we; voi = you (plural); loro = they; Lei = you (formal))}}\newline Lei with a capital L in writing is the formal 'you'. {\color{muted}(See \S\ref{les:infinitives-pronouns:pronouns}.)}
\item\relax \textbf{(b) The ending -o already shows the subject is io.}\newline Subject pronouns are optional because the verb ending identifies the person. {\color{muted}(See \S\ref{les:infinitives-pronouns:omission}.)}
\item\relax \textbf{(d) Io parlo italiano, lui parla inglese.}\newline Contrast between two subjects requires the pronouns. In the others they are optional. {\color{muted}(See \S\ref{les:infinitives-pronouns:omission}.)}
\item\relax \textbf{\textit{dormire}}\newline dorm- + -ire. {\color{muted}(See \S\ref{les:infinitives-pronouns:families}.)}
\end{answers}

\answerheading{\S\ref{les:present-are} Present tense of -are verbs}\label{ans:present-are}
\begin{answers}
\item\relax \textbf{io: \textit{canto}; tu: \textit{canti}; lui/lei: \textit{canta}; noi: \textit{cantiamo}; voi: \textit{cantate}; loro: \textit{cantano}}\newline cant- + -o, -i, -a, -iamo, -ate, -ano. {\color{muted}(See \S\ref{les:present-are:endings}.)}
\item\relax \textbf{\textit{cerchi}}\newline Add h to keep the hard c: cerchi.\newline {\color{muted}Common mistake: \textit{cerci} --- In cerci the c before i would be soft, as in English 'church'. Add h to keep the hard k sound: cerchi.} {\color{muted}(See \S\ref{les:present-are:spelling}.)}
\item\relax \textbf{\textit{mangiamo}}\newline The stem mangi- already ends in i: mangi + amo = mangiamo.\newline {\color{muted}Common mistake: \textit{mangiiamo} --- Do not double the i: mangiamo.} {\color{muted}(See \S\ref{les:present-are:spelling}.)}
\item\relax \textbf{\textit{giochiamo}}\newline giocare adds h before i: giochiamo. {\color{muted}(See \S\ref{les:present-are:spelling}.)}
\item\relax \textbf{(1) \textit{giocano}; (2) \textit{guarda}; (3) \textit{ascoltiamo}}\newline loro → -ano; lui/lei → -a; noi → -iamo. {\color{muted}(See \S\ref{les:present-are:endings}.)}
\item\relax \textbf{\textit{Che cosa cerchi?} {\color{muted}(also accepted: \textit{Cosa cerchi?}; \textit{Che cerchi?})}}\newline Italian has no separate continuous form here: cerchi = you look for / you are looking for. {\color{muted}(See \S\ref{les:present-are:meaning}.)}
\item\relax \textbf{\textit{studiare}}\newline studi- + -amo; infinitive studiare. {\color{muted}(See \S\ref{les:present-are:spelling}.)}
\end{answers}

\answerheading{\S\ref{les:present-ere-ire} Present tense of -ere and -ire verbs}\label{ans:present-ere-ire}
\begin{answers}
\item\relax \textbf{io: \textit{prendo}; tu: \textit{prendi}; lui/lei: \textit{prende}; noi: \textit{prendiamo}; voi: \textit{prendete}; loro: \textit{prendono}}\newline prend- + -o, -i, -e, -iamo, -ete, -ono. {\color{muted}(See \S\ref{les:present-ere-ire:ere}.)}
\item\relax \textbf{io: \textit{finisco}; tu: \textit{finisci}; lui/lei: \textit{finisce}; noi: \textit{finiamo}; voi: \textit{finite}; loro: \textit{finiscono}}\newline -isc- in all forms except noi and voi. {\color{muted}(See \S\ref{les:present-ere-ire:ire}.)}
\item\relax \textbf{\textit{dormite}}\newline voi ending follows the infinitive vowel: -ite.\newline {\color{muted}Common mistake: \textit{dormete} --- -ete is the -ere ending. dormire is an -ire verb: dormite.} {\color{muted}(See \S\ref{les:present-ere-ire:comparison}.)}
\item\relax \textbf{(d) costruiscono}\newline costruire is an -isc- verb: costruiscono. {\color{muted}(See \S\ref{les:present-ere-ire:ire}.)}
\item\relax \textbf{\textit{capire}}\newline cap- + -isc- + -e → capire. {\color{muted}(See \S\ref{les:present-ere-ire:ire}.)}
\item\relax \textbf{(1) \textit{leggo}; (2) \textit{scrivi}; (3) \textit{partono}}\newline leggo, scrivi, partono. {\color{muted}(See \S\ref{les:present-ere-ire:comparison}.)}
\item\relax \textbf{\textit{I prefer tea.} {\color{muted}(also accepted: \textit{I prefer the tea})}}\newline preferire is an -isc- verb: preferisco. {\color{muted}(See \S\ref{les:present-ere-ire:ire}.)}
\end{answers}

\section*{Chapter \ref{mod:m05}: Essential irregular verbs}

\answerheading{\S\ref{les:essere-avere} essere and avere}\label{ans:essere-avere}
\begin{answers}
\item\relax \textbf{io: \textit{sono}; tu: \textit{sei}; lui/lei: \textit{è}; noi: \textit{siamo}; voi: \textit{siete}; loro: \textit{sono}}\newline Fully irregular; note the accent on è. {\color{muted}(See \S\ref{les:essere-avere:essere}.)}
\item\relax \textbf{io: \textit{ho}; tu: \textit{hai}; lui/lei: \textit{ha}; noi: \textit{abbiamo}; voi: \textit{avete}; loro: \textit{hanno}}\newline Remember the silent h in ho, hai, ha, hanno. {\color{muted}(See \S\ref{les:essere-avere:avere}.)}
\item\relax \textbf{(b) ha}\newline Age uses avere: ha dieci anni.\newline {\color{muted}Common mistake: \textit{è} --- In Italian you 'have' years: Luca ha dieci anni.}\newline {\color{muted}Common mistake: \textit{a} --- a without h is the preposition 'to'. The verb is ha.} {\color{muted}(See \S\ref{les:essere-avere:avere}.)}
\item\relax \textbf{(b) ha}\newline avere paura = to be afraid. {\color{muted}(See \S\ref{les:essere-avere:avere}.)}
\item\relax \textbf{(1) \textit{abbiamo}; (2) \textit{siete}; (3) \textit{hanno}}\newline avere fame; essere stanchi; avere una casa. {\color{muted}(See \S\ref{les:essere-avere:comparison}.)}
\item\relax \textbf{\textit{Hai fame?}}\newline avere fame. {\color{muted}(See \S\ref{les:essere-avere:avere}.)}
\item\relax \textbf{\textit{We are from Milan.} {\color{muted}(also accepted: \textit{We're from Milan}; \textit{We come from Milan})}}\newline essere di + city = to be from. {\color{muted}(See \S\ref{les:essere-avere:essere}.)}
\end{answers}

\answerheading{\S\ref{les:andare-fare-stare-dire} andare, fare, stare, dire}\label{ans:andare-fare-stare-dire}
\begin{answers}
\item\relax \textbf{io: \textit{vado}; tu: \textit{vai}; lui/lei: \textit{va}; noi: \textit{andiamo}; voi: \textit{andate}; loro: \textit{vanno}}\newline vad-/va- in stressed forms, and- in noi/voi. {\color{muted}(See \S\ref{les:andare-fare-stare-dire:andare}.)}
\item\relax \textbf{io: \textit{faccio}; tu: \textit{fai}; lui/lei: \textit{fa}; noi: \textit{facciamo}; voi: \textit{fate}; loro: \textit{fanno}}\newline faccio and facciamo with double c. {\color{muted}(See \S\ref{les:andare-fare-stare-dire:fare-stare}.)}
\item\relax \textbf{(b) a}\newline Cities take a: vado a Roma. {\color{muted}(See \S\ref{les:andare-fare-stare-dire:andare}.)}
\item\relax \textbf{(b) in}\newline Fixed phrase: andare in montagna (but al mare). {\color{muted}(See \S\ref{les:andare-fare-stare-dire:andare}.)}
\item\relax \textbf{(1) \textit{stai}; (2) \textit{Sto} {\color{muted}(also accepted: \textit{sto})}; (3) \textit{fa}}\newline stai, sto; fa caldo. {\color{muted}(See \S\ref{les:andare-fare-stare-dire:fare-stare}.)}
\item\relax \textbf{\textit{dicono}}\newline dic- stem: dicono.\newline {\color{muted}Common mistake: \textit{dirono} --- dire uses the stem dic-: dicono.} {\color{muted}(See \S\ref{les:andare-fare-stare-dire:dire}.)}
\item\relax \textbf{\textit{Facciamo colazione insieme.}}\newline fare colazione. {\color{muted}(See \S\ref{les:andare-fare-stare-dire:fare-stare}.)}
\end{answers}

\answerheading{\S\ref{les:other-irregulars} venire, uscire, bere, dare and piacere}\label{ans:other-irregulars}
\begin{answers}
\item\relax \textbf{io: \textit{vengo}; tu: \textit{vieni}; lui/lei: \textit{viene}; noi: \textit{veniamo}; voi: \textit{venite}; loro: \textit{vengono}}\newline g in io and loro; ie in tu and lui/lei. {\color{muted}(See \S\ref{les:other-irregulars:pattern}.)}
\item\relax \textbf{\textit{esco}}\newline Stressed forms begin with e-: esco, esci, esce, escono.\newline {\color{muted}Common mistake: \textit{usco} --- The stressed forms of uscire begin with e-: esco.} {\color{muted}(See \S\ref{les:other-irregulars:uscire-bere-dare}.)}
\item\relax \textbf{\textit{beviamo}}\newline bere uses the stem bev-: beviamo. {\color{muted}(See \S\ref{les:other-irregulars:uscire-bere-dare}.)}
\item\relax \textbf{(d) piacciono}\newline le fragole is plural, so piacciono.\newline {\color{muted}Common mistake: \textit{piace} --- The thing liked is the subject. le fragole is plural: mi piacciono.} {\color{muted}(See \S\ref{les:other-irregulars:piacere}.)}
\item\relax \textbf{(d) piace}\newline An infinitive counts as singular: ti piace cantare. {\color{muted}(See \S\ref{les:other-irregulars:piacere}.)}
\item\relax \textbf{\textit{tenere}}\newline tengo, tieni, tiene, teniamo, tenete, tengono → tenere. {\color{muted}(See \S\ref{les:other-irregulars:pattern}.)}
\item\relax \textbf{\textit{Non mi piace la pioggia.}}\newline non goes before mi piace. {\color{muted}(See \S\ref{les:other-irregulars:piacere}.)}
\end{answers}

\section*{Chapter \ref{mod:m06}: Building sentences}

\answerheading{\S\ref{les:word-order-negation-questions} Word order, negation and yes/no questions}\label{ans:word-order-negation-questions}
\begin{answers}
\item\relax \textbf{\textit{Il gatto non dorme}}\newline non directly before the verb. {\color{muted}(See \S\ref{les:word-order-negation-questions:negation}.)}
\item\relax \textbf{\textit{Non mi sveglio presto}}\newline non comes before the reflexive pronoun mi. {\color{muted}(See \S\ref{les:word-order-negation-questions:negation}.)}
\item\relax \textbf{\textit{Non piango mai.}}\newline non … mai: both parts are needed.\newline {\color{muted}Common mistake: \textit{Piango mai.} --- Italian needs non before the verb with mai: Non piango mai.} {\color{muted}(See \S\ref{les:word-order-negation-questions:negation}.)}
\item\relax \textbf{\textit{Parli italiano?} {\color{muted}(also accepted: \textit{Tu parli italiano?}; \textit{Parli l'italiano?})}}\newline Same word order as the statement, with rising intonation. {\color{muted}(See \S\ref{les:word-order-negation-questions:questions}.)}
\item\relax \textbf{(c) Non vedo niente.}\newline non + verb + niente. {\color{muted}(See \S\ref{les:word-order-negation-questions:negation}.)}
\item\relax \textbf{(c) the subject}\newline The subject can follow the verb, especially with verbs of arriving and appearing. arriva agrees with il treno. {\color{muted}(See \S\ref{les:word-order-negation-questions:order}.)}
\end{answers}

\answerheading{\S\ref{les:question-words-conjunctions} Question words and conjunctions}\label{ans:question-words-conjunctions}
\begin{answers}
\item\relax \textbf{1--a, 2--d, 3--c, 4--e, 5--b {\color{muted}(Dove abiti? = A Roma.; Quanti anni hai? = Sette.; Chi è? = È la nonna.; Quando parti? = Domani.; Perché piangi? = Perché sono stanco.)}}\newline dove → place; quanti → number; chi → person; quando → time; perché → reason. {\color{muted}(See \S\ref{les:question-words-conjunctions:question-words}.)}
\item\relax \textbf{(d) Quante}\newline stelle is feminine plural: quante. {\color{muted}(See \S\ref{les:question-words-conjunctions:question-words}.)}
\item\relax \textbf{(b) Qual è}\newline qual è is written without an apostrophe (quale is truncated, not elided). {\color{muted}(See \S\ref{les:question-words-conjunctions:question-words}.)}
\item\relax \textbf{(1) \textit{perché}; (2) \textit{Se} {\color{muted}(also accepted: \textit{se})}; (3) \textit{ma}; (4) \textit{mentre}}\newline reason → perché; condition → se; contrast → ma; simultaneity → mentre. {\color{muted}(See \S\ref{les:question-words-conjunctions:conjunctions}.)}
\item\relax \textbf{\textit{Dove abita la volpe?} {\color{muted}(also accepted: \textit{Dove vive la volpe?})}}\newline The subject follows the verb in the question. {\color{muted}(See \S\ref{les:question-words-conjunctions:question-words}.)}
\item\relax \textbf{\textit{What is your name?} {\color{muted}(also accepted: \textit{What's your name?}; \textit{What are you called?})}}\newline Literally 'How do you call yourself?' (chiamarsi). {\color{muted}(See \S\ref{les:question-words-conjunctions:question-words}.)}
\end{answers}

\section*{Chapter \ref{mod:m07}: Prepositions}

\answerheading{\S\ref{les:simple-prepositions} The simple prepositions}\label{ans:simple-prepositions}
\begin{answers}
\item\relax \textbf{(a) a}\newline Cities take a. {\color{muted}(See \S\ref{les:simple-prepositions:overview}.)}
\item\relax \textbf{(b) in}\newline Countries take in. {\color{muted}(See \S\ref{les:simple-prepositions:overview}.)}
\item\relax \textbf{(1) \textit{di}; (2) \textit{da}; (3) \textit{con}; (4) \textit{per}}\newline di = possession; da = origin; con = with; per = for. {\color{muted}(See \S\ref{les:simple-prepositions:overview}.)}
\item\relax \textbf{(a) fra}\newline tra/fra + time = in … from now. per un'ora would mean 'for an hour'.\newline {\color{muted}Common mistake: \textit{per} --- per un'ora = for (the duration of) an hour. 'In an hour from now' is fra (or tra) un'ora.} {\color{muted}(See \S\ref{les:simple-prepositions:time}.)}
\item\relax \textbf{\textit{Il cane di Sara è nero.}}\newline Italian has no 's: il cane di Sara. {\color{muted}(See \S\ref{les:simple-prepositions:overview}.)}
\item\relax \textbf{\textit{I have lived here for two years.} {\color{muted}(also accepted: \textit{I've lived here for two years}; \textit{I have been living here for two years}; \textit{I've been living here for two years})}}\newline da + time with the present tense = an action that started in the past and continues. {\color{muted}(See \S\ref{les:simple-prepositions:time}.)}
\end{answers}

\answerheading{\S\ref{les:articulated-prepositions} Articulated prepositions}\label{ans:articulated-prepositions}
\begin{answers}
\item\relax \textbf{\textit{nel}}\newline in becomes ne-: ne + l = nel.\newline {\color{muted}Common mistake: \textit{in il} --- in + il always contracts: nel.} {\color{muted}(See \S\ref{les:articulated-prepositions:table}.)}
\item\relax \textbf{\textit{sulla}}\newline su + la = sulla (l doubled). {\color{muted}(See \S\ref{les:articulated-prepositions:table}.)}
\item\relax \textbf{di + il: \textit{del}; di + lo: \textit{dello}; di + l': \textit{dell'}; di + la: \textit{della}; di + i: \textit{dei}; di + gli: \textit{degli}; di + le: \textit{delle}}\newline di → de- + article; l doubled before articles beginning with l. {\color{muted}(See \S\ref{les:articulated-prepositions:table}.)}
\item\relax \textbf{(c) sugli}\newline alberi takes gli: su + gli = sugli. {\color{muted}(See \S\ref{les:articulated-prepositions:table}.)}
\item\relax \textbf{(d) a}\newline City names take no article, so no contraction: a Roma. {\color{muted}(See \S\ref{les:articulated-prepositions:use}.)}
\item\relax \textbf{(1) \textit{della}; (2) \textit{al}; (3) \textit{dai}}\newline della, al, dai. {\color{muted}(See \S\ref{les:articulated-prepositions:use}.)}
\item\relax \textbf{\textit{La volpe torna nella tana.} {\color{muted}(also accepted: \textit{La volpe torna nella sua tana.})}}\newline in + la = nella. Italian often omits the possessive when ownership is obvious. {\color{muted}(See \S\ref{les:articulated-prepositions:use}.)}
\end{answers}

\section*{Chapter \ref{mod:m08}: Modal verbs}

\answerheading{\S\ref{les:modal-verbs} Modal verbs: potere, volere, dovere, sapere}\label{ans:modal-verbs}
\begin{answers}
\item\relax \textbf{io: \textit{devo}; tu: \textit{devi}; lui/lei: \textit{deve}; noi: \textit{dobbiamo}; voi: \textit{dovete}; loro: \textit{devono}}\newline dev- (stressed), dobb- (noi), dov- (voi). {\color{muted}(See \S\ref{les:modal-verbs:forms}.)}
\item\relax \textbf{\textit{può}}\newline può carries an accent.\newline {\color{muted}Common mistake: \textit{puo} --- può needs the accent on the final o.} {\color{muted}(See \S\ref{les:modal-verbs:forms}.)}
\item\relax \textbf{\textit{vogliamo}}\newline vogli- + -amo. {\color{muted}(See \S\ref{les:modal-verbs:forms}.)}
\item\relax \textbf{(a) So}\newline A learned skill: sapere + infinitive.\newline {\color{muted}Common mistake: \textit{Posso} --- posso = I am able / allowed (circumstances). For a learned skill use so.}\newline {\color{muted}Common mistake: \textit{Conosco} --- conoscere means to know a person or place, not a skill.} {\color{muted}(See \S\ref{les:modal-verbs:potere-sapere}.)}
\item\relax \textbf{\textit{Devo alzarmi presto}}\newline The reflexive pronoun can attach to the infinitive: alzarmi. (Also correct: Mi devo alzare presto.) {\color{muted}(See \S\ref{les:modal-verbs:construction}.)}
\item\relax \textbf{\textit{Non posso venire oggi.} {\color{muted}(also accepted: \textit{Oggi non posso venire.})}}\newline non before the modal; infinitive venire. {\color{muted}(See \S\ref{les:modal-verbs:construction}.)}
\item\relax \textbf{\textit{Vorrei un bicchiere d'acqua, per favore.} {\color{muted}(also accepted: \textit{Vorrei un bicchiere di acqua, per favore.})}}\newline vorrei is the polite form; d'acqua = di acqua. {\color{muted}(See \S\ref{les:modal-verbs:construction}.)}
\end{answers}

\section*{Chapter \ref{mod:m09}: Reflexive verbs}

\answerheading{\S\ref{les:reflexive-verbs} Reflexive verbs and daily routine}\label{ans:reflexive-verbs}
\begin{answers}
\item\relax \textbf{io: \textit{mi alzo}; tu: \textit{ti alzi}; lui/lei: \textit{si alza}; noi: \textit{ci alziamo}; voi: \textit{vi alzate}; loro: \textit{si alzano}}\newline Pronoun + regular -are forms. {\color{muted}(See \S\ref{les:reflexive-verbs:form}.)}
\item\relax \textbf{(b) ci}\newline noi → ci.\newline {\color{muted}Common mistake: \textit{si} --- si is for lui/lei and loro. With noi use ci.} {\color{muted}(See \S\ref{les:reflexive-verbs:form}.)}
\item\relax \textbf{\textit{mi siedo}}\newline sedere has the stressed stem sied-: mi siedo. {\color{muted}(See \S\ref{les:reflexive-verbs:form}.)}
\item\relax \textbf{\textit{Non mi alzo presto}}\newline non + pronoun + verb. {\color{muted}(See \S\ref{les:reflexive-verbs:form}.)}
\item\relax \textbf{\textit{Mi lavo le mani.}}\newline Body parts: reflexive pronoun + article, no possessive.\newline {\color{muted}Common mistake: \textit{Lavo le mie mani.} --- Grammatically possible but unidiomatic. Italian says Mi lavo le mani.} {\color{muted}(See \S\ref{les:reflexive-verbs:meaning}.)}
\item\relax \textbf{\textit{addormentarsi}}\newline addormentare + si. {\color{muted}(See \S\ref{les:reflexive-verbs:form}.)}
\item\relax \textbf{(1) \textit{mi sveglio}; (2) \textit{mi vesto}; (3) \textit{si addormentano}}\newline mi sveglio, mi vesto, si addormentano. {\color{muted}(See \S\ref{les:reflexive-verbs:routine}.)}
\end{answers}

\section*{Chapter \ref{mod:m15}: Imperativo}

\answerheading{\S\ref{les:imperativo} The imperative}\label{ans:imperativo}
\begin{answers}
\item\relax \textbf{parlare: \textit{parla}; scrivere: \textit{scrivi}; finire: \textit{finisci}; essere: \textit{sii}; avere: \textit{abbi}; dire: \textit{di'}}\newline -are → -a; -ere/-ire → -i (with -isc- where the present has it); irregular sii, abbi, di'. {\color{muted}(See \S\ref{les:imperativo:forms}.)}
\item\relax \textbf{(c) Non andare!}\newline The negative tu imperative is non + infinitive. Non vada! is the formal Lei form. {\color{muted}(See \S\ref{les:imperativo:negative}.)}
\item\relax \textbf{\textit{venga}}\newline Formal Lei uses the congiuntivo presente: venga.\newline {\color{muted}Common mistake: \textit{vieni} --- vieni is the informal tu form. Formal Lei: venga.} {\color{muted}(See \S\ref{les:imperativo:forms}.)}
\item\relax \textbf{\textit{fallo} {\color{muted}(also accepted: \textit{fallo!})}}\newline After fa' the l of lo doubles: fallo.\newline {\color{muted}Common mistake: \textit{falo} --- After the short form fa' the consonant doubles: fallo.} {\color{muted}(See \S\ref{les:imperativo:pronouns}.)}
\item\relax \textbf{(b) Sedetevi!}\newline voi imperative + attached reflexive pronoun vi: sedetevi. {\color{muted}(See \S\ref{les:imperativo:pronouns}.)}
\item\relax \textbf{\textit{Non dirlo!} {\color{muted}(also accepted: \textit{Non lo dire!})}}\newline Both positions of the pronoun are standard with the negative tu imperative. {\color{muted}(See \S\ref{les:imperativo:pronouns}.)}
\item\relax \textbf{tu: \textit{alzati}; Lei: \textit{si alzi}; noi: \textit{alziamoci}; voi: \textit{alzatevi}}\newline Pronoun attached to tu/noi/voi; before the verb with Lei. {\color{muted}(See \S\ref{les:imperativo:pronouns}.)}
\end{answers}

\section*{Chapter \ref{mod:m10}: Passato prossimo and auxiliaries}

\answerheading{\S\ref{les:passato-prossimo-avere} Passato prossimo with avere}\label{ans:passato-prossimo-avere}
\begin{answers}
\item\relax \textbf{giocare: \textit{giocato}; credere: \textit{creduto}; sentire: \textit{sentito}; fare: \textit{fatto}; vedere: \textit{visto} {\color{muted}(also accepted: \textit{veduto})}; scrivere: \textit{scritto}}\newline Regular: -ato, -uto, -ito. fatto, visto (also veduto, less common), scritto are irregular. {\color{muted}(See \S\ref{les:passato-prossimo-avere:irregular}.)}
\item\relax \textbf{\textit{abbiamo mangiato}}\newline abbiamo + mangiato. {\color{muted}(See \S\ref{les:passato-prossimo-avere:formation}.)}
\item\relax \textbf{\textit{ha preso}}\newline Irregular participle preso.\newline {\color{muted}Common mistake: \textit{ha prenduto} --- prendere has an irregular participle: preso.} {\color{muted}(See \S\ref{les:passato-prossimo-avere:irregular}.)}
\item\relax \textbf{\textit{dire}}\newline dire → detto. {\color{muted}(See \S\ref{les:passato-prossimo-avere:irregular}.)}
\item\relax \textbf{\textit{Non ho mai visto un lupo}}\newline mai goes between the auxiliary and the participle. {\color{muted}(See \S\ref{les:passato-prossimo-avere:use}.)}
\item\relax \textbf{\textit{Ieri il nonno ha letto una storia.} {\color{muted}(also accepted: \textit{Il nonno ieri ha letto una storia.}; \textit{Il nonno ha letto una storia ieri.})}}\newline leggere → letto. {\color{muted}(See \S\ref{les:passato-prossimo-avere:irregular}.)}
\item\relax \textbf{(1) \textit{ho fatto}; (2) \textit{hanno bevuto}; (3) \textit{hai visto} {\color{muted}(also accepted: \textit{hai veduto})}}\newline ho fatto, hanno bevuto, hai visto. {\color{muted}(See \S\ref{les:passato-prossimo-avere:formation}.)}
\end{answers}

\answerheading{\S\ref{les:passato-prossimo-essere} Passato prossimo with essere; agreement}\label{ans:passato-prossimo-essere}
\begin{answers}
\item\relax \textbf{(b) è andata}\newline andare takes essere; the participle agrees with Sara (feminine singular).\newline {\color{muted}Common mistake: \textit{ha andato} --- andare is a verb of movement that takes essere: è andata.}\newline {\color{muted}Common mistake: \textit{è andato} --- With essere the participle agrees with the subject: Sara → andata.} {\color{muted}(See \S\ref{les:passato-prossimo-essere:agreement}.)}
\item\relax \textbf{1--a, 2--b, 3--a, 4--b, 5--a {\color{muted}(arrivare = essere; mangiare = avere; partire = essere; camminare = avere; alzarsi = essere)}}\newline Movement to/from a place and reflexives → essere; camminare (moving without a destination) → avere. {\color{muted}(See \S\ref{les:passato-prossimo-essere:which}.)}
\item\relax \textbf{\textit{sono tornate}}\newline essere + feminine plural participle. {\color{muted}(See \S\ref{les:passato-prossimo-essere:agreement}.)}
\item\relax \textbf{\textit{mi sono svegliato}}\newline Reflexive pronoun + essere + agreeing participle.\newline {\color{muted}Common mistake: \textit{mi ho svegliato} --- Reflexive verbs take essere: mi sono svegliato.} {\color{muted}(See \S\ref{les:passato-prossimo-essere:reflexive}.)}
\item\relax \textbf{(1) \textit{è caduta}; (2) \textit{ci siamo divertiti} {\color{muted}(also accepted: \textit{ci siamo divertite})}; (3) \textit{è salito}}\newline foglia (f) → caduta; noi → divertiti (or divertite for all women); gatto → salito. {\color{muted}(See \S\ref{les:passato-prossimo-essere:agreement}.)}
\item\relax \textbf{\textit{I bambini sono tornati a casa.}}\newline tornare → essere; bambini masculine plural → tornati. {\color{muted}(See \S\ref{les:passato-prossimo-essere:agreement}.)}
\end{answers}

\answerheading{\S\ref{les:ausiliari-accordo} Choosing the auxiliary and agreeing the participle}\label{ans:ausiliari-accordo}
\begin{answers}
\item\relax \textbf{1--a, 2--d, 3--a, 4--c {\color{muted}(Il concerto \blank{} cominciato. = è; Noi \blank{} cominciato il lavoro. = abbiamo; Il prezzo \blank{} cambiato. = è; Loro \blank{} cambiato idea. = hanno)}}\newline With an object (il lavoro, idea) the verb is transitive → avere; without, describing what happened to the subject → essere. {\color{muted}(See \S\ref{les:ausiliari-accordo:both}.)}
\item\relax \textbf{(1) \textit{è nata}; (2) \textit{è cresciuta}; (3) \textit{è piaciuto}}\newline nascere, crescere (intransitive) and piacere take essere; the participle agrees with the subject (Sara, il film). {\color{muted}(See \S\ref{les:ausiliari-accordo:avere-essere}.)}
\item\relax \textbf{\textit{Le ho comprate.}}\newline le before the verb → the participle agrees: comprate.\newline {\color{muted}Common mistake: \textit{Le ho comprato.} --- With a preceding le the participle must agree: comprate.} {\color{muted}(See \S\ref{les:ausiliari-accordo:agreement}.)}
\item\relax \textbf{(c) scritto}\newline Le here is an indirect object (to you); the direct object una lettera follows the verb. No agreement. {\color{muted}(See \S\ref{les:ausiliari-accordo:agreement}.)}
\item\relax \textbf{\textit{hanno visti} {\color{muted}(also accepted: \textit{hanno visto}; \textit{hanno veduti}; \textit{hanno veduto})}}\newline With ci agreement is optional: ci hanno visti and ci hanno visto are both correct. {\color{muted}(See \S\ref{les:ausiliari-accordo:agreement}.)}
\item\relax \textbf{(a) Non ci hanno voluti andare.}\newline When ci (linked to the essere-verb andare) is placed before the modal, the modal takes essere: non ci sono voluti andare. With ci attached to the infinitive, avere is normal: non hanno voluto andarci. {\color{muted}(See \S\ref{les:ausiliari-accordo:modals}.)}
\item\relax \textbf{\textit{Sono dovuta partire.} {\color{muted}(also accepted: \textit{Ho dovuto partire.}; \textit{Io sono dovuta partire.}; \textit{Io ho dovuto partire.}; \textit{Sono dovuta andare via.}; \textit{Ho dovuto andare via.})}}\newline Both auxiliaries are correct with modal + essere-verb; with essere the participle agrees (dovuta).\newline {\color{muted}Common mistake: \textit{Sono dovuto partire.} --- With essere the participle agrees with a female speaker: dovuta.} {\color{muted}(See \S\ref{les:ausiliari-accordo:modals}.)}
\item\relax \textbf{(1) \textit{fatte} {\color{muted}(also accepted: \textit{fatto})}; (2) \textit{conosciuto}}\newline With ne agreement is usual (fatte; fatto also found). With relative che the modern standard is invariable (conosciuto). {\color{muted}(See \S\ref{les:ausiliari-accordo:agreement}.)}
\end{answers}

\section*{Chapter \ref{mod:m11}: Imperfetto and trapassato prossimo}

\answerheading{\S\ref{les:imperfetto} The imperfetto}\label{ans:imperfetto}
\begin{answers}
\item\relax \textbf{io: \textit{cantavo}; tu: \textit{cantavi}; lui/lei: \textit{cantava}; noi: \textit{cantavamo}; voi: \textit{cantavate}; loro: \textit{cantavano}}\newline canta- + -vo, -vi, -va, -vamo, -vate, -vano. {\color{muted}(See \S\ref{les:imperfetto:formation}.)}
\item\relax \textbf{io: \textit{ero}; tu: \textit{eri}; lui/lei: \textit{era}; noi: \textit{eravamo}; voi: \textit{eravate}; loro: \textit{erano}}\newline Irregular stem er-. {\color{muted}(See \S\ref{les:imperfetto:formation}.)}
\item\relax \textbf{\textit{faceva}}\newline fare uses the stem face-: faceva.\newline {\color{muted}Common mistake: \textit{fava} --- fare uses its old stem face-: faceva.} {\color{muted}(See \S\ref{les:imperfetto:formation}.)}
\item\relax \textbf{\textit{dire}}\newline dice- + -vano → dire. {\color{muted}(See \S\ref{les:imperfetto:formation}.)}
\item\relax \textbf{(1) \textit{Era} {\color{muted}(also accepted: \textit{era})}; (2) \textit{era}; (3) \textit{faceva}; (4) \textit{aveva}}\newline Descriptions of time, place, weather and states use the imperfetto. {\color{muted}(See \S\ref{les:imperfetto:uses}.)}
\item\relax \textbf{\textit{Da bambino abitavo al mare.} {\color{muted}(also accepted: \textit{Da bambino vivevo al mare.}; \textit{Quando ero bambino abitavo al mare.}; \textit{Quando ero bambino vivevo al mare.})}}\newline da bambino = as a child; a past situation → imperfetto. {\color{muted}(See \S\ref{les:imperfetto:uses}.)}
\item\relax \textbf{\textit{Once upon a time there was a little bear.} {\color{muted}(also accepted: \textit{Once upon a time there was a small bear.})}}\newline C'era una volta = once upon a time there was. {\color{muted}(See \S\ref{les:imperfetto:uses}.)}
\end{answers}

\answerheading{\S\ref{les:imperfetto-passato} Imperfetto or passato prossimo?}\label{ans:imperfetto-passato}
\begin{answers}
\item\relax \textbf{(a) dormivo}\newline Action in progress, interrupted → imperfetto. {\color{muted}(See \S\ref{les:imperfetto-passato:mentre-quando}.)}
\item\relax \textbf{(b) è entrato}\newline The interrupting event → passato prossimo, with essere (entrare).\newline {\color{muted}Common mistake: \textit{ha entrato} --- Right tense, wrong auxiliary: entrare takes essere — è entrato.} {\color{muted}(See \S\ref{les:imperfetto-passato:mentre-quando}.)}
\item\relax \textbf{(d) andavamo}\newline A repeated habit → imperfetto. {\color{muted}(See \S\ref{les:imperfetto-passato:contrast}.)}
\item\relax \textbf{(c) siamo andati}\newline One specific occasion → passato prossimo. {\color{muted}(See \S\ref{les:imperfetto-passato:contrast}.)}
\item\relax \textbf{(1) \textit{Era} {\color{muted}(also accepted: \textit{era})}; (2) \textit{pioveva}; (3) \textit{ha sentito}}\newline Scene-setting → imperfetto; a sudden event → passato prossimo. {\color{muted}(See \S\ref{les:imperfetto-passato:contrast}.)}
\item\relax \textbf{\textit{Mentre la mamma cucinava, io leggevo.} {\color{muted}(also accepted: \textit{Mentre la mamma cucinava, leggevo.}; \textit{Io leggevo mentre la mamma cucinava.}; \textit{Leggevo mentre la mamma cucinava.})}}\newline Two simultaneous background actions → both imperfetto. {\color{muted}(See \S\ref{les:imperfetto-passato:mentre-quando}.)}
\end{answers}

\answerheading{\S\ref{les:trapassato-prossimo} Trapassato prossimo}\label{ans:trapassato-prossimo}
\begin{answers}
\item\relax \textbf{io: \textit{ero uscito}; tu: \textit{eri uscito}; lui: \textit{era uscito}; noi: \textit{eravamo usciti}; voi: \textit{eravate usciti}; loro: \textit{erano usciti}}\newline uscire takes essere: imperfetto of essere + uscito/usciti. {\color{muted}(See \S\ref{les:trapassato-prossimo:formation}.)}
\item\relax \textbf{(1) \textit{avevano mangiato}; (2) \textit{era andata}}\newline Both events precede my arrival. andare takes essere; andata agrees with Sara. {\color{muted}(See \S\ref{les:trapassato-prossimo:uses}.)}
\item\relax \textbf{(c) Luca went out}\newline The trapassato (era uscito) marks the earlier event. {\color{muted}(See \S\ref{les:trapassato-prossimo:uses}.)}
\item\relax \textbf{\textit{aveva finito}}\newline Past reporting verb: passato prossimo → trapassato prossimo.\newline {\color{muted}Common mistake: \textit{ha finito} --- This is possible in casual speech, but the standard sequence after ha detto is the trapassato: aveva finito.} {\color{muted}(See \S\ref{les:trapassato-prossimo:uses}.)}
\item\relax \textbf{\textit{Non avevo mai visto un lupo.} {\color{muted}(also accepted: \textit{Io non avevo mai visto un lupo.}; \textit{Non avevo mai veduto un lupo.})}}\newline mai goes between the auxiliary and the participle, as in the passato prossimo. {\color{muted}(See \S\ref{les:trapassato-prossimo:formation}.)}
\item\relax \textbf{(b) ebbe scritto}\newline With a passato remoto main clause, dopo che + trapassato remoto is the literary choice; aveva scritto is also correct. {\color{muted}(See \S\ref{les:trapassato-prossimo:uses}.)}
\end{answers}

\section*{Chapter \ref{mod:m12}: Pronouns and complex sentences}

\answerheading{\S\ref{les:direct-object-pronouns} Direct object pronouns}\label{ans:direct-object-pronouns}
\begin{answers}
\item\relax \textbf{(d) Le}\newline le stelle: feminine plural → le. {\color{muted}(See \S\ref{les:direct-object-pronouns:forms}.)}
\item\relax \textbf{(c) Li}\newline i bambini: masculine plural → li.\newline {\color{muted}Common mistake: \textit{Gli} --- gli is an indirect object pronoun ('to him/to them'). The direct object 'them' (masculine) is li.} {\color{muted}(See \S\ref{les:direct-object-pronouns:forms}.)}
\item\relax \textbf{\textit{Non lo vedo}}\newline non + pronoun + verb. {\color{muted}(See \S\ref{les:direct-object-pronouns:position}.)}
\item\relax \textbf{\textit{l'ho vista}}\newline la elides before ho, and the participle agrees with la: vista.\newline {\color{muted}Common mistake: \textit{l'ho visto} --- The pronoun la (la volpe) is feminine, so the participle agrees: l'ho vista.} {\color{muted}(See \S\ref{les:direct-object-pronouns:agreement}.)}
\item\relax \textbf{\textit{vederlo}}\newline vedere loses -e and takes lo: vederlo. {\color{muted}(See \S\ref{les:direct-object-pronouns:position}.)}
\item\relax \textbf{(a) direct object pronoun (it, feminine)}\newline Before a verb, la is an object pronoun; before a noun (la mela) it is the article. {\color{muted}(See \S\ref{les:direct-object-pronouns:forms}.)}
\item\relax \textbf{(1) \textit{Ne}; (2) \textit{Lo}}\newline A quantity (due) → ne; the whole definite object (il libro) → lo. {\color{muted}(See \S\ref{les:direct-object-pronouns:ne}.)}
\item\relax \textbf{\textit{Ne ho mangiate} {\color{muted}(also accepted: \textit{Ne ho mangiato})}}\newline With ne the participle normally agrees with the noun (mele → mangiate); the invariable form is also found. {\color{muted}(See \S\ref{les:direct-object-pronouns:ne}.)}
\end{answers}

\answerheading{\S\ref{les:indirect-object-pronouns} Indirect object pronouns}\label{ans:indirect-object-pronouns}
\begin{answers}
\item\relax \textbf{(b) Le}\newline a Sara = to her → le.\newline {\color{muted}Common mistake: \textit{La} --- la is the direct object (her). 'To her' is le.} {\color{muted}(See \S\ref{les:indirect-object-pronouns:forms}.)}
\item\relax \textbf{(d) Gli}\newline a Marco = to him → gli. {\color{muted}(See \S\ref{les:indirect-object-pronouns:forms}.)}
\item\relax \textbf{(d) Lo}\newline aspettare takes a direct object (no preposition): lo aspetto (often l'aspetto).\newline {\color{muted}Common mistake: \textit{Gli} --- aspettare takes a direct object in Italian (aspettare il treno), so use lo.} {\color{muted}(See \S\ref{les:indirect-object-pronouns:verbs}.)}
\item\relax \textbf{(1) \textit{le}; (2) \textit{gli}; (3) \textit{ci}}\newline to her → le; to him → gli; to us → ci. {\color{muted}(See \S\ref{les:indirect-object-pronouns:piacere}.)}
\item\relax \textbf{\textit{Il nonno ci racconta una storia.}}\newline ci = to us, before the verb. {\color{muted}(See \S\ref{les:indirect-object-pronouns:verbs}.)}
\item\relax \textbf{\textit{They like cats.} {\color{muted}(also accepted: \textit{They like the cats})}}\newline Literally 'Cats please them'. {\color{muted}(See \S\ref{les:indirect-object-pronouns:piacere}.)}
\item\relax \textbf{\textit{Glielo do.}}\newline a Marco → gli; il libro → lo; gli + lo = glielo.\newline {\color{muted}Common mistake: \textit{Gli lo do.} --- gli + lo fuse into one word: glielo.} {\color{muted}(See \S\ref{les:indirect-object-pronouns:combined}.)}
\item\relax \textbf{(d) te la}\newline ti + la → te la (la penna is feminine). {\color{muted}(See \S\ref{les:indirect-object-pronouns:combined}.)}
\end{answers}

\answerheading{\S\ref{les:relative-clauses-adverbs} Relative clauses, adverbs and expressions in stories}\label{ans:relative-clauses-adverbs}
\begin{answers}
\item\relax \textbf{(d) che}\newline Direct object without preposition → che. {\color{muted}(See \S\ref{les:relative-clauses-adverbs:relatives}.)}
\item\relax \textbf{(c) cui}\newline After a preposition (con) → cui.\newline {\color{muted}Common mistake: \textit{che} --- After a preposition Italian uses cui: con cui.} {\color{muted}(See \S\ref{les:relative-clauses-adverbs:relatives}.)}
\item\relax \textbf{\textit{velocemente}}\newline veloce + -mente. {\color{muted}(See \S\ref{les:relative-clauses-adverbs:adverbs}.)}
\item\relax \textbf{\textit{gentilmente}}\newline -le adjectives drop the e: gentilmente.\newline {\color{muted}Common mistake: \textit{gentilemente} --- Adjectives in -le and -re drop the final e: gentilmente.} {\color{muted}(See \S\ref{les:relative-clauses-adverbs:adverbs}.)}
\item\relax \textbf{(d) di}\newline cercare di + infinitive = to try to. {\color{muted}(See \S\ref{les:relative-clauses-adverbs:infinitives-idioms}.)}
\item\relax \textbf{(d) a}\newline cominciare a + infinitive. {\color{muted}(See \S\ref{les:relative-clauses-adverbs:infinitives-idioms}.)}
\item\relax \textbf{\textit{Il ragazzo che canta è mio fratello.}}\newline che cannot be omitted. {\color{muted}(See \S\ref{les:relative-clauses-adverbs:relatives}.)}
\end{answers}

\section*{Chapter \ref{mod:m13}: Futuro}

\answerheading{\S\ref{les:futuro-semplice} Futuro semplice}\label{ans:futuro-semplice}
\begin{answers}
\item\relax \textbf{io: \textit{cercherò}; tu: \textit{cercherai}; lui/lei: \textit{cercherà}; noi: \textit{cercheremo}; voi: \textit{cercherete}; loro: \textit{cercheranno}}\newline cerc- + h + -erò …: the h keeps the c hard before e. {\color{muted}(See \S\ref{les:futuro-semplice:formation}.)}
\item\relax \textbf{\textit{mangerò}}\newline -giare verbs drop the i: mangerò.\newline {\color{muted}Common mistake: \textit{mangierò} --- The i of mangi- is only needed before a, o, u. Before e it is dropped: mangerò.}\newline {\color{muted}Common mistake: \textit{mangero} --- The io form carries a written accent: mangerò.} {\color{muted}(See \S\ref{les:futuro-semplice:formation}.)}
\item\relax \textbf{essere: \textit{sarò}; venire: \textit{verrò}; andare: \textit{andrò}; volere: \textit{vorrò}; fare: \textit{farò}; vedere: \textit{vedrò}}\newline Irregular stems: sar-, verr-, andr-, vorr-, far-, vedr-. {\color{muted}(See \S\ref{les:futuro-semplice:irregular}.)}
\item\relax \textbf{(d) a supposition about the present}\newline The futuro of supposition: 'he must be / is probably at the gym (now)'. {\color{muted}(See \S\ref{les:futuro-semplice:uses}.)}
\item\relax \textbf{(c) A quest'ora sarà in ufficio.}\newline A quest'ora (at this time, now) + futuro = a guess about the present: 'He must be at the office by now'. The others refer to future time. {\color{muted}(See \S\ref{les:futuro-semplice:uses}.)}
\item\relax \textbf{(1) \textit{andremo}; (2) \textit{rimarrete}; (3) \textit{saranno}}\newline andr-emo, rimarr-ete, sar-anno. {\color{muted}(See \S\ref{les:futuro-semplice:irregular}.)}
\item\relax \textbf{\textit{Saranno le otto.} {\color{muted}(also accepted: \textit{Saranno circa le otto.}; \textit{Saranno più o meno le otto.})}}\newline Le otto is plural, so the verb is saranno. {\color{muted}(See \S\ref{les:futuro-semplice:uses}.)}
\end{answers}

\answerheading{\S\ref{les:futuro-anteriore} Futuro anteriore}\label{ans:futuro-anteriore}
\begin{answers}
\item\relax \textbf{\textit{avrai finito}}\newline Transitive finire (finire qualcosa) takes avere. {\color{muted}(See \S\ref{les:futuro-anteriore:formation}.)}
\item\relax \textbf{\textit{saremo arrivati}}\newline arrivare takes essere; the participle agrees with noi.\newline {\color{muted}Common mistake: \textit{avremo arrivato} --- arrivare takes essere: saremo arrivati.} {\color{muted}(See \S\ref{les:futuro-anteriore:formation}.)}
\item\relax \textbf{(d) a guess about the past}\newline Futuro anteriore of supposition: 'there must have been a fault'. {\color{muted}(See \S\ref{les:futuro-anteriore:uses}.)}
\item\relax \textbf{(b) avrete finito}\newline An action completed before a future action: futuro anteriore. (Informally also: quando avete finito, andiamo al parco.) {\color{muted}(See \S\ref{les:futuro-anteriore:uses}.)}
\item\relax \textbf{\textit{Avrà perso il treno.} {\color{muted}(also accepted: \textit{Avrà perduto il treno.})}}\newline perdere → perso (also perduto). {\color{muted}(See \S\ref{les:futuro-anteriore:uses}.)}
\item\relax \textbf{(1) \textit{avrò scritto}; (2) \textit{sarà andata}}\newline avrò scritto (anteriority); sarà andata (supposition, essere + agreement with Sara). {\color{muted}(See \S\ref{les:futuro-anteriore:uses}.)}
\end{answers}

\section*{Chapter \ref{mod:m14}: Condizionale}

\answerheading{\S\ref{les:condizionale-presente} Condizionale presente}\label{ans:condizionale-presente}
\begin{answers}
\item\relax \textbf{io: \textit{potrei}; tu: \textit{potresti}; lui/lei: \textit{potrebbe}; noi: \textit{potremmo}; voi: \textit{potreste}; loro: \textit{potrebbero}}\newline Future stem potr- + conditional endings. {\color{muted}(See \S\ref{les:condizionale-presente:formation}.)}
\item\relax \textbf{(b) andremmo}\newline Conditional: -emmo (double m). andremo is the future, andassimo the congiuntivo imperfetto.\newline {\color{muted}Common mistake: \textit{andremo} --- andremo (one m) is the futuro: we will go. The conditional has -emmo.} {\color{muted}(See \S\ref{les:condizionale-presente:formation}.)}
\item\relax \textbf{(b) The president is reportedly ill (not confirmed).}\newline The journalistic conditional marks information the writer does not vouch for. {\color{muted}(See \S\ref{les:condizionale-presente:uses}.)}
\item\relax \textbf{(c) potessi}\newline Standard Italian uses the congiuntivo imperfetto in the se-clause; the conditional goes in the main clause (verrei).\newline {\color{muted}Common mistake: \textit{potrei} --- The conditional is not used after se in standard Italian: se potessi.} {\color{muted}(See \S\ref{les:condizionale-presente:uses}.)}
\item\relax \textbf{\textit{Dovresti dormire di più.}}\newline dovere in the conditional = ought to / should. {\color{muted}(See \S\ref{les:condizionale-presente:uses}.)}
\item\relax \textbf{(1) \textit{piacerebbe}; (2) \textit{Vorrei} {\color{muted}(also accepted: \textit{vorrei})}; (3) \textit{verreste}}\newline piacer-ebbe; vorr-ei; verr-este. {\color{muted}(See \S\ref{les:condizionale-presente:formation}.)}
\end{answers}

\answerheading{\S\ref{les:condizionale-passato} Condizionale passato and the future in the past}\label{ans:condizionale-passato}
\begin{answers}
\item\relax \textbf{\textit{sarei andato}}\newline andare takes essere; the participle agrees with the subject.\newline {\color{muted}Common mistake: \textit{avrei andato} --- andare takes essere: sarei andato.} {\color{muted}(See \S\ref{les:condizionale-passato:formation}.)}
\item\relax \textbf{(b) sarebbe venuta}\newline A future seen from a past point of view: condizionale passato. 'verrà' is possible only if the speaker stresses that the event is still in the future now; verrebbe is not standard here.\newline {\color{muted}Common mistake: \textit{verrebbe} --- Italian uses the condizionale passato for the future in the past: sarebbe venuta.} {\color{muted}(See \S\ref{les:condizionale-passato:future-in-past}.)}
\item\relax \textbf{\textit{Avresti dovuto dirmelo.} {\color{muted}(also accepted: \textit{Me lo avresti dovuto dire.}; \textit{Me l'avresti dovuto dire.}; \textit{Avresti dovuto dirlo a me.})}}\newline condizionale passato of dovere + infinitive with the combined pronoun me lo attached (dirmelo). {\color{muted}(See \S\ref{les:condizionale-passato:unrealised}.)}
\item\relax \textbf{(1) \textit{avrebbe scritto}; (2) \textit{sarebbe piovuto} {\color{muted}(also accepted: \textit{avrebbe piovuto})}}\newline Future in the past = condizionale passato. Weather verbs accept both auxiliaries. {\color{muted}(See \S\ref{les:condizionale-passato:future-in-past}.)}
\item\relax \textbf{(c) No}\newline The condizionale passato presents a possibility that was not realised. {\color{muted}(See \S\ref{les:condizionale-passato:unrealised}.)}
\item\relax \textbf{1--d, 2--b, 3--c, 4--a {\color{muted}(Avrei voluto un tè. = polite past wish; Avresti dovuto studiare. = reproach; Disse che sarebbe tornato. = future in the past; Il ladro sarebbe fuggito. = unconfirmed report)}}\newline The condizionale passato covers unrealised wishes and obligations, the future in the past and reported information. {\color{muted}(See \S\ref{les:condizionale-passato:unrealised}.)}
\end{answers}

\section*{Chapter \ref{mod:m16}: Infinitive, gerund and participle}

\answerheading{\S\ref{les:infinito-gerundio} Infinitive and gerund}\label{ans:infinito-gerundio}
\begin{answers}
\item\relax \textbf{(d) Dopo aver finito}\newline dopo requires the past infinitive. {\color{muted}(See \S\ref{les:infinito-gerundio:infinitive}.)}
\item\relax \textbf{(a) Mi piace nuotare.}\newline The thing liked is an action: infinitive, with piace (singular). {\color{muted}(See \S\ref{les:infinito-gerundio:gerund}.)}
\item\relax \textbf{cantare: \textit{cantando}; leggere: \textit{leggendo}; partire: \textit{partendo}; fare: \textit{facendo}; dire: \textit{dicendo}; bere: \textit{bevendo}}\newline -ando / -endo; fare, dire, bere use their extended stems (fac-, dic-, bev-). {\color{muted}(See \S\ref{les:infinito-gerundio:gerund}.)}
\item\relax \textbf{\textit{Stavamo cenando.} {\color{muted}(also accepted: \textit{Noi stavamo cenando.})}}\newline stare in the imperfetto + gerundio. (cenare = to have dinner; also: Stavamo mangiando la cena.) {\color{muted}(See \S\ref{les:infinito-gerundio:gerund}.)}
\item\relax \textbf{(c) the speaker}\newline The gerund's subject is normally the main verb's subject (io). To say Sara was: Ho visto Sara che usciva dal negozio. {\color{muted}(See \S\ref{les:infinito-gerundio:gerund}.)}
\item\relax \textbf{(1) \textit{essere arrivata} {\color{muted}(also accepted: \textit{esser arrivata})}; (2) \textit{aver letto} {\color{muted}(also accepted: \textit{avere letto})}}\newline arrivare → essere + agreement (arrivata); leggere → avere. Truncated aver/esser are equally correct. {\color{muted}(See \S\ref{les:infinito-gerundio:infinitive}.)}
\end{answers}

\answerheading{\S\ref{les:participio} The participle}\label{ans:participio}
\begin{answers}
\item\relax \textbf{(c) Scritte}\newline Transitive participle in a participial clause agrees with the object it refers to: le lettere → scritte. {\color{muted}(See \S\ref{les:participio:past}.)}
\item\relax \textbf{(b) Arrivata}\newline Intransitive verb: the participle agrees with the subject (Sara). {\color{muted}(See \S\ref{les:participio:past}.)}
\item\relax \textbf{(d) a past participle used as an adjective}\newline aprire → aperto/a: here it describes the noun. {\color{muted}(See \S\ref{les:participio:past}.)}
\item\relax \textbf{1--d, 2--c, 3--a, 4--b {\color{muted}(cantante = singer; insegnante = teacher; bollente = boiling hot; corrente = running, current)}}\newline Present participles now mostly live as nouns and adjectives. {\color{muted}(See \S\ref{les:participio:present}.)}
\item\relax \textbf{\textit{When the lesson was over, the children went out.} {\color{muted}(also accepted: \textit{After the lesson ended, the children went out.}; \textit{When the lesson ended, the children went out.}; \textit{Once the lesson was over, the children went out.}; \textit{The lesson over, the children went out.}; \textit{After the lesson was over, the children went out.})}}\newline A participial clause expressing an earlier action; uscirono is passato remoto. {\color{muted}(See \S\ref{les:participio:past}.)}
\end{answers}

\section*{Chapter \ref{mod:m17}: Passato remoto}

\answerheading{\S\ref{les:passato-remoto} Passato remoto: forms, meaning and regional use}\label{ans:passato-remoto}
\begin{answers}
\item\relax \textbf{io: \textit{cantai}; tu: \textit{cantasti}; lui/lei: \textit{cantò}; noi: \textit{cantammo}; voi: \textit{cantaste}; loro: \textit{cantarono}}\newline -ai, -asti, -ò, -ammo, -aste, -arono. {\color{muted}(See \S\ref{les:passato-remoto:regular}.)}
\item\relax \textbf{\textit{vendetti} {\color{muted}(also accepted: \textit{vendei})}}\newline Both vendei and vendetti are standard. {\color{muted}(See \S\ref{les:passato-remoto:regular}.)}
\item\relax \textbf{ebbe: \textit{avere}; fece: \textit{fare}; disse: \textit{dire}; vide: \textit{vedere}; nacque: \textit{nascere}; fu: \textit{essere}}\newline Strong forms hide the infinitive: learn them in sets (ebbi/ebbe/ebbero). {\color{muted}(See \S\ref{les:passato-remoto:irregular}.)}
\item\relax \textbf{io: \textit{presi}; tu: \textit{prendesti}; lui/lei: \textit{prese}; noi: \textit{prendemmo}; voi: \textit{prendeste}; loro: \textit{presero}}\newline 1-3-3: strong stem pres- in io, lui/lei, loro; regular prend- elsewhere. {\color{muted}(See \S\ref{les:passato-remoto:irregular}.)}
\item\relax \textbf{(b) ho visto}\newline stamattina is part of today, linked to the present: passato prossimo. In southern regional speech vidi is also heard.\newline {\color{muted}Common mistake: \textit{vidi} --- In standard usage a time period that includes the present (stamattina, oggi) takes the passato prossimo. vidi here is typical of some southern regional speech.} {\color{muted}(See \S\ref{les:passato-remoto:meaning}.)}
\item\relax \textbf{(c) It is standard in written narration; in speech its use varies strongly by region.}\newline Rare in northern conversation, traditional in Tuscany, very common in southern speech, and the norm in written narrative everywhere. {\color{muted}(See \S\ref{les:passato-remoto:meaning}.)}
\item\relax \textbf{(1) \textit{entrò}; (2) \textit{vide}; (3) \textit{ebbe}; (4) \textit{scappò}}\newline entrò, vide (strong), ebbe (strong), scappò. {\color{muted}(See \S\ref{les:passato-remoto:irregular}.)}
\end{answers}

\answerheading{\S\ref{les:trapassato-remoto} Trapassato remoto}\label{ans:trapassato-remoto}
\begin{answers}
\item\relax \textbf{\textit{ebbe mangiato}}\newline Passato remoto of avere + participle. {\color{muted}(See \S\ref{les:trapassato-remoto:formation}.)}
\item\relax \textbf{\textit{furono arrivati}}\newline arrivare takes essere: furono arrivati. {\color{muted}(See \S\ref{les:trapassato-remoto:formation}.)}
\item\relax \textbf{(b) ebbe letto}\newline Temporal clause with dopo che and main verb in the passato remoto (pianse): trapassato remoto. (aveva letto is also found in modern prose, but the traditional norm is ebbe letto.) {\color{muted}(See \S\ref{les:trapassato-remoto:use}.)}
\item\relax \textbf{(a) Appena ebbe finito, se ne andò.}\newline Only in a temporal clause before a main verb in the passato remoto (andò). {\color{muted}(See \S\ref{les:trapassato-remoto:use}.)}
\item\relax \textbf{\textit{When they had finished, they went home.} {\color{muted}(also accepted: \textit{When they had finished, they returned home.}; \textit{When they had finished, they went back home.}; \textit{After they had finished, they went home.})}}\newline ebbero finito: trapassato remoto; tornarono: passato remoto. {\color{muted}(See \S\ref{les:trapassato-remoto:use}.)}
\end{answers}

\section*{Chapter \ref{mod:m18}: Congiuntivo}

\answerheading{\S\ref{les:congiuntivo-presente} Congiuntivo presente: forms and main uses}\label{ans:congiuntivo-presente}
\begin{answers}
\item\relax \textbf{parlare: \textit{parli}; prendere: \textit{prenda}; finire: \textit{finisca}; fare: \textit{faccia}; andare: \textit{vada}; essere: \textit{sia}}\newline -are → -i; -ere/-ire → -a (finisca keeps -isc-); irregulars from the io form (faccio → faccia, vado → vada); essere → sia. {\color{muted}(See \S\ref{les:congiuntivo-presente:formation}.)}
\item\relax \textbf{(b) abbia}\newline Verb of opinion in the present: congiuntivo presente. (ha is frequent in informal speech but not the standard written form.)\newline {\color{muted}Common mistake: \textit{ha} --- Understandable and common in casual speech, but after credo che standard Italian uses the congiuntivo: abbia.} {\color{muted}(See \S\ref{les:congiuntivo-presente:uses}.)}
\item\relax \textbf{(b) Spero di vederti presto.}\newline Same subject (io spero, io vedo): di + infinitive. {\color{muted}(See \S\ref{les:congiuntivo-presente:uses}.)}
\item\relax \textbf{1--c, 2--e, 3--b, 4--d, 5--a {\color{muted}(So che… = indicative; Penso che… = congiuntivo; Benché… = congiuntivo (conjunction); È certo che… = indicative (certainty); Bisogna che… = congiuntivo (impersonal))}}\newline Certainty → indicative; opinion, necessity, concession → congiuntivo. {\color{muted}(See \S\ref{les:congiuntivo-presente:uses}.)}
\item\relax \textbf{(1) \textit{veniate}; (2) \textit{dorma}; (3) \textit{sappiano}}\newline veniate (from veniamo), dorma, sappiano (irregular). {\color{muted}(See \S\ref{les:congiuntivo-presente:formation}.)}
\item\relax \textbf{\textit{Voglio che tu venga con me.} {\color{muted}(also accepted: \textit{Voglio che venga con me.})}}\newline Different subjects: voglio che + congiuntivo; tu is usually expressed because venga is the same for io, tu and lui.\newline {\color{muted}Common mistake: \textit{Voglio te venire con me.} --- Italian does not use 'want + person + infinitive'. Use voglio che tu venga.} {\color{muted}(See \S\ref{les:congiuntivo-presente:uses}.)}
\item\relax \textbf{(a) affinché}\newline affinché (so that) requires the congiuntivo. perché can also mean 'so that' with the congiuntivo, but as 'because' it takes the indicative. {\color{muted}(See \S\ref{les:congiuntivo-presente:uses}.)}
\end{answers}

\answerheading{\S\ref{les:congiuntivo-passato} Congiuntivo passato}\label{ans:congiuntivo-passato}
\begin{answers}
\item\relax \textbf{\textit{abbia fatto}}\newline abbia + fatto. {\color{muted}(See \S\ref{les:congiuntivo-passato:formation}.)}
\item\relax \textbf{(d) sia andato}\newline Present main verb + earlier action: congiuntivo passato. {\color{muted}(See \S\ref{les:congiuntivo-passato:use}.)}
\item\relax \textbf{(b) sia}\newline Simultaneous action: congiuntivo presente. {\color{muted}(See \S\ref{les:congiuntivo-passato:use}.)}
\item\relax \textbf{(1) \textit{siate arrivati} {\color{muted}(also accepted: \textit{siate arrivate})}; (2) \textit{abbia telefonato}}\newline arrivare → essere (agreement); telefonare → avere. {\color{muted}(See \S\ref{les:congiuntivo-passato:formation}.)}
\item\relax \textbf{\textit{È il libro più bello che io abbia mai letto.} {\color{muted}(also accepted: \textit{È il libro più bello che abbia mai letto.})}}\newline Superlative + che + congiuntivo passato. {\color{muted}(See \S\ref{les:congiuntivo-passato:use}.)}
\end{answers}

\answerheading{\S\ref{les:congiuntivo-imperfetto-trapassato} Congiuntivo imperfetto and trapassato}\label{ans:congiuntivo-imperfetto-trapassato}
\begin{answers}
\item\relax \textbf{parlare: \textit{parlassi}; avere: \textit{avessi}; essere: \textit{fossi}; fare: \textit{facessi}; dire: \textit{dicessi}; stare: \textit{stessi}}\newline From the imperfetto stem (parla-, ave-, face-, dice-); fossi and stessi are irregular. {\color{muted}(See \S\ref{les:congiuntivo-imperfetto-trapassato:formation}.)}
\item\relax \textbf{(c) lavorasse}\newline Past main verb (credevo) + simultaneous action: congiuntivo imperfetto.\newline {\color{muted}Common mistake: \textit{lavori} --- After a past main verb the congiuntivo presente is not used: lavorasse.} {\color{muted}(See \S\ref{les:congiuntivo-imperfetto-trapassato:use}.)}
\item\relax \textbf{(b) foste partiti}\newline Past main verb + earlier action: congiuntivo trapassato (foste partiti). partiste would be passato remoto or congiuntivo imperfetto (simultaneous). {\color{muted}(See \S\ref{les:congiuntivo-imperfetto-trapassato:use}.)}
\item\relax \textbf{(1) \textit{fosse}; (2) \textit{venissi}; (3) \textit{sapesse}}\newline Wishes, conditional main verb and come se all take the congiuntivo imperfetto. {\color{muted}(See \S\ref{les:congiuntivo-imperfetto-trapassato:use}.)}
\item\relax \textbf{\textit{Pensavo che fossi già partito.} {\color{muted}(also accepted: \textit{Pensavo che tu fossi già partito.}; \textit{Credevo che fossi già partito.}; \textit{Credevo che tu fossi già partito.})}}\newline Past main verb + earlier action: congiuntivo trapassato. {\color{muted}(See \S\ref{les:congiuntivo-imperfetto-trapassato:use}.)}
\item\relax \textbf{(d) passato remoto or congiuntivo imperfetto (voi)}\newline The voi forms coincide; context decides. {\color{muted}(See \S\ref{les:congiuntivo-imperfetto-trapassato:formation}.)}
\end{answers}

\section*{Chapter \ref{mod:m19}: Complex sentences}

\answerheading{\S\ref{les:periodo-ipotetico} Hypothetical sentences (periodo ipotetico)}\label{ans:periodo-ipotetico}
\begin{answers}
\item\relax \textbf{(b) avessi}\newline Possible type: se + congiuntivo imperfetto, condizionale presente.\newline {\color{muted}Common mistake: \textit{avrei} --- The conditional is not used in the se-clause in standard Italian: se avessi.} {\color{muted}(See \S\ref{les:periodo-ipotetico:types}.)}
\item\relax \textbf{(a) avrei superato}\newline Unreal past: se + congiuntivo trapassato, condizionale passato. (superavo would be the colloquial pattern, but only with avevo studiato.) {\color{muted}(See \S\ref{les:periodo-ipotetico:types}.)}
\item\relax \textbf{(1) \textit{fossi}; (2) \textit{partirei}; (3) \textit{piove} {\color{muted}(also accepted: \textit{pioverà})}}\newline fossi / partirei (possible-unreal present); piove or pioverà (real, future reference). {\color{muted}(See \S\ref{les:periodo-ipotetico:types}.)}
\item\relax \textbf{1--c, 2--d, 3--b, 4--e, 5--a {\color{muted}(Se hai fame, mangia. = realtà; Se potessi, verrei. = possibilità; Se fossi venuto, ti saresti divertito. = irrealtà (past); Se l'avessi saputo, ora sarei lì. = mixed; Se lo sapevo, venivo. = colloquial)}}\newline Look at the mood and tense in each clause. {\color{muted}(See \S\ref{les:periodo-ipotetico:types}.)}
\item\relax \textbf{\textit{Se l'avessi saputo, sarei venuto.} {\color{muted}(also accepted: \textit{Se avessi saputo, sarei venuto.}; \textit{Se lo avessi saputo, sarei venuto.})}}\newline Congiuntivo trapassato + condizionale passato. 'Se lo sapevo, venivo' is the colloquial equivalent.\newline {\color{muted}Common mistake: \textit{Se lo sapevo, venivo.} --- Correct in informal speech, but the task asks for standard written Italian: Se l'avessi saputo, sarei venuto.} {\color{muted}(See \S\ref{les:periodo-ipotetico:types}.)}
\item\relax \textbf{(b) Mi chiedo se Marco verrebbe.}\newline In an indirect question (se = whether) the conditional is fine. In hypothetical se-clauses it is not standard. {\color{muted}(See \S\ref{les:periodo-ipotetico:register}.)}
\end{answers}

\answerheading{\S\ref{les:concordanza-dei-tempi} Sequence of tenses (concordanza dei tempi)}\label{ans:concordanza-dei-tempi}
\begin{answers}
\item\relax \textbf{earlier: \textit{fosse partito}; simultaneous: \textit{partisse}; later: \textit{sarebbe partito}}\newline Past main verb: congiuntivo trapassato / imperfetto / condizionale passato. {\color{muted}(See \S\ref{les:concordanza-dei-tempi:subjunctive}.)}
\item\relax \textbf{earlier: \textit{sia partito}; simultaneous: \textit{parta}; later: \textit{partirà} {\color{muted}(also accepted: \textit{parta})}}\newline Present main verb: congiuntivo passato / presente / futuro (or congiuntivo presente). {\color{muted}(See \S\ref{les:concordanza-dei-tempi:subjunctive}.)}
\item\relax \textbf{(c) restaste}\newline Conditional main verb of wish → congiuntivo imperfetto.\newline {\color{muted}Common mistake: \textit{restiate} --- After vorrei the congiuntivo imperfetto is required: restaste.} {\color{muted}(See \S\ref{les:concordanza-dei-tempi:subjunctive}.)}
\item\relax \textbf{(c) avrebbe vinto}\newline Posteriority in the past: condizionale passato. {\color{muted}(See \S\ref{les:concordanza-dei-tempi:subjunctive}.)}
\item\relax \textbf{(1) \textit{era}; (2) \textit{aveva dormito}}\newline Indicative context: imperfetto (simultaneous), trapassato prossimo (earlier). {\color{muted}(See \S\ref{les:concordanza-dei-tempi:indicative}.)}
\item\relax \textbf{(a) Yes: the fact is still true now, so the present can be kept.}\newline Both è sposato and era sposato are correct; the present stresses that it is still true. {\color{muted}(See \S\ref{les:concordanza-dei-tempi:indicative}.)}
\end{answers}

\answerheading{\S\ref{les:discorso-indiretto} Reported speech (discorso indiretto)}\label{ans:discorso-indiretto}
\begin{answers}
\item\relax \textbf{\textit{aveva}}\newline presente → imperfetto after a past verb of saying. {\color{muted}(See \S\ref{les:discorso-indiretto:tenses}.)}
\item\relax \textbf{\textit{sarebbe partito}}\newline futuro → condizionale passato.\newline {\color{muted}Common mistake: \textit{partirebbe} --- The future in the past is the condizionale passato: sarebbe partito.} {\color{muted}(See \S\ref{les:discorso-indiretto:tenses}.)}
\item\relax \textbf{\textit{aveva visto}}\newline passato prossimo → trapassato prossimo. {\color{muted}(See \S\ref{les:discorso-indiretto:tenses}.)}
\item\relax \textbf{1--b, 2--e, 3--d, 4--a, 5--c {\color{muted}(oggi = quel giorno; domani = il giorno dopo; ieri = il giorno prima; qui = lì; questo = quello)}}\newline Deictic words shift to the reported situation. {\color{muted}(See \S\ref{les:discorso-indiretto:deixis}.)}
\item\relax \textbf{\textit{He asked me if I was hungry.} {\color{muted}(also accepted: \textit{She asked me if I was hungry.}; \textit{He asked me whether I was hungry.}; \textit{She asked me whether I was hungry.})}}\newline Yes/no question reported with se; avere fame = to be hungry. {\color{muted}(See \S\ref{les:discorso-indiretto:questions}.)}
\item\relax \textbf{(a) di chiudere}\newline Reported command: di + infinitive (or che chiudessi). {\color{muted}(See \S\ref{les:discorso-indiretto:questions}.)}
\end{answers}

\section*{Chapter \ref{mod:m20}: Passive and impersonal constructions}

\answerheading{\S\ref{les:forma-passiva} The passive voice}\label{ans:forma-passiva}
\begin{answers}
\item\relax \textbf{\textit{è raccontata} {\color{muted}(also accepted: \textit{viene raccontata})}}\newline Present passive: essere (or venire) + participle agreeing with la storia. {\color{muted}(See \S\ref{les:forma-passiva:essere}.)}
\item\relax \textbf{\textit{sono state mangiate}}\newline Passato prossimo passive: sono state + mangiate (agreement with le mele).\newline {\color{muted}Common mistake: \textit{sono venute mangiate} --- venire is used only in simple tenses: sono state mangiate.} {\color{muted}(See \S\ref{les:forma-passiva:essere}.)}
\item\relax \textbf{(c) si vendono}\newline si passivante: the verb agrees with libri (plural).\newline {\color{muted}Common mistake: \textit{si vende} --- With the si passivante the verb agrees with the noun: libri → si vendono.} {\color{muted}(See \S\ref{les:forma-passiva:si-passivante}.)}
\item\relax \textbf{(d) The form must be signed.}\newline andare + participle = obligation. {\color{muted}(See \S\ref{les:forma-passiva:venire-andare}.)}
\item\relax \textbf{(c) arrivare}\newline Only transitive verbs have a passive; arrivare is intransitive. {\color{muted}(See \S\ref{les:forma-passiva:essere}.)}
\item\relax \textbf{\textit{Il castello fu costruito nel 1300.}}\newline fu (passato remoto of essere) + costruito. {\color{muted}(See \S\ref{les:forma-passiva:essere}.)}
\end{answers}

\answerheading{\S\ref{les:si-impersonale} Impersonal si}\label{ans:si-impersonale}
\begin{answers}
\item\relax \textbf{(d) si è}\newline Impersonal si + singular verb; the adjective giovani is plural. {\color{muted}(See \S\ref{les:si-impersonale:basics}.)}
\item\relax \textbf{(d) ci si alza}\newline Reflexive verb + impersonal si → ci si.\newline {\color{muted}Common mistake: \textit{si si alza} --- Two si are not possible: the first becomes ci (ci si alza).} {\color{muted}(See \S\ref{les:si-impersonale:basics}.)}
\item\relax \textbf{\textit{stanchi}}\newline Predicate adjectives after impersonal si are plural.\newline {\color{muted}Common mistake: \textit{stanco} --- Although the verb is singular, the adjective is plural: si è stanchi.} {\color{muted}(See \S\ref{les:si-impersonale:basics}.)}
\item\relax \textbf{(c) si è arrivati}\newline Impersonal si takes essere; with an essere-verb the participle is plural. {\color{muted}(See \S\ref{les:si-impersonale:compound}.)}
\item\relax \textbf{1--a, 2--b, 3--a, 4--b {\color{muted}(Si mangia bene qui. = impersonal; Si mangiano molte mele. = passivante; Si vive bene. = impersonal; Si vende la casa. = passivante)}}\newline With a direct-object noun the si is passivante and the verb agrees with it. {\color{muted}(See \S\ref{les:si-impersonale:compound}.)}
\item\relax \textbf{\textit{Come si dice «cat» in italiano?} {\color{muted}(also accepted: \textit{Come si dice cat in italiano?}; \textit{Come si dice 'cat' in italiano?})}}\newline Impersonal si + dire. {\color{muted}(See \S\ref{les:si-impersonale:basics}.)}
\end{answers}

\section*{Stories}

\answerheading{\S\ref{story:il-gatto-e-la-luna} Il gatto e la luna}\label{ans:story:il-gatto-e-la-luna}
\begin{answers}
\item\relax \textbf{(c) nero}\newline Tobia è un gatto nero e piccolo.
\item\relax \textbf{(c) sul letto}\newline Di giorno Tobia dorme sul letto.
\item\relax \textbf{(c) Non risponde.}\newline La luna non risponde.
\item\relax \textbf{(a) milk}\newline Tobia beve il suo latte.
\item\relax \textbf{\textit{The moon is big and white.} {\color{muted}(also accepted: \textit{The moon is large and white})}}\newline grande = big, bianca = white (feminine).
\end{answers}

\answerheading{\S\ref{story:la-famiglia-di-marco} La famiglia di Marco}\label{ans:story:la-famiglia-di-marco}
\begin{answers}
\item\relax \textbf{(b) sette}\newline Mi chiamo Marco e ho sette anni.
\item\relax \textbf{(a) È maestra.}\newline Mia madre si chiama Anna ed è maestra.
\item\relax \textbf{(d) il cane}\newline Abbiamo un cane, Bruno.
\item\relax \textbf{(c) al mare}\newline Il nonno e la nonna abitano al mare.
\item\relax \textbf{\textit{Mia sorella ama i gatti.}}\newline mia sorella without article; i gatti with article (general category).
\end{answers}

\answerheading{\S\ref{story:la-colazione} La colazione}\label{ans:story:la-colazione}
\begin{answers}
\item\relax \textbf{(c) la mamma}\newline La mamma prepara la colazione.
\item\relax \textbf{(d) pane e marmellata}\newline Vorrei pane e marmellata… Sara mangia il pane.
\item\relax \textbf{(a) un po' di latte}\newline La mamma dà al gatto un po' di latte.
\item\relax \textbf{(a) to school}\newline Poi Sara prende lo zaino e va a scuola.
\item\relax \textbf{\textit{bere}}\newline bere: bevo, bevi, beve…
\end{answers}

\answerheading{\S\ref{story:il-cane-al-mare} Il cane al mare}\label{ans:story:il-cane-al-mare}
\begin{answers}
\item\relax \textbf{(a) Fa caldo.}\newline Oggi fa caldo.
\item\relax \textbf{(a) un castello di sabbia}\newline Luca fa un castello di sabbia.
\item\relax \textbf{(b) sotto un sasso}\newline Il granchio scappa sotto un sasso.
\item\relax \textbf{(c) la conchiglia}\newline Luca trova una conchiglia e la porta a casa: la = the shell.
\item\relax \textbf{(a) tired, wet and happy}\newline La sera Pippo è stanco e bagnato, ma felice.
\end{answers}

\answerheading{\S\ref{story:nel-bosco} Nel bosco}\label{ans:story:nel-bosco}
\begin{answers}
\item\relax \textbf{(d) con il nonno}\newline Giulia e suo fratello Luca vanno nel bosco con il nonno.
\item\relax \textbf{(d) tre}\newline Su un ramo c'è un nido con tre piccoli uccelli.
\item\relax \textbf{(d) nella sua tana}\newline La volpe … torna nella sua tana.
\item\relax \textbf{(a) so as not to make noise}\newline per non fare rumore.
\item\relax \textbf{\textit{salire}}\newline salire: salgo, sali, sale, saliamo, salite, salgono.
\item\relax \textbf{\textit{di}}\newline di for possession/belonging.
\end{answers}

\answerheading{\S\ref{story:la-tartaruga-lenta} La tartaruga lenta}\label{ans:story:la-tartaruga-lenta}
\begin{answers}
\item\relax \textbf{(b) la lepre}\newline Ogni giorno la lepre ride della tartaruga.
\item\relax \textbf{(a) al grande albero}\newline «La gara finisce al grande albero» dice il gufo.
\item\relax \textbf{(b) Perché si addormenta.}\newline La lepre si siede sotto un albero e si addormenta.
\item\relax \textbf{\textit{sedersi}}\newline sedersi: mi siedo, ti siedi, si siede…
\item\relax \textbf{\textit{The tortoise never stops.}}\newline non … mai = never; fermarsi = to stop.
\end{answers}

\answerheading{\S\ref{story:una-giornata-di-pioggia} Una giornata di pioggia}\label{ans:story:una-giornata-di-pioggia}
\begin{answers}
\item\relax \textbf{(a) Piove.}\newline Oggi piove e il cielo è grigio.
\item\relax \textbf{(d) gli stivali}\newline «Mettiamo gli stivali e prendiamo l'ombrello!»
\item\relax \textbf{(b) bagnati ma contenti}\newline …i bambini sono bagnati ma contenti.
\item\relax \textbf{\textit{uscire}}\newline uscire: esco … escono.
\item\relax \textbf{\textit{Non possiamo giocare in giardino.}}\newline non + possiamo + infinitive.
\end{answers}

\answerheading{\S\ref{story:lorso-che-non-poteva-dormire} L'orso che non poteva dormire}\label{ans:story:lorso-che-non-poteva-dormire}
\begin{answers}
\item\relax \textbf{(a) Aveva la pancia vuota.}\newline «Ho la pancia vuota.»
\item\relax \textbf{(d) un gufo}\newline Sopra un ramo c'era un gufo.
\item\relax \textbf{(a) Ha cantato una canzone.}\newline Il gufo … ha cominciato a cantare piano piano.
\item\relax \textbf{(d) It describes the background.}\newline Descriptions of time and weather use the imperfetto.
\item\relax \textbf{\textit{ridere}}\newline ridere → riso (irregular participle).
\item\relax \textbf{(a) until spring}\newline Ha dormito fino alla primavera.
\end{answers}

\answerheading{\S\ref{story:il-segreto-del-faro} Il segreto del faro}\label{ans:story:il-segreto-del-faro}
\begin{answers}
\item\relax \textbf{(a) Accendeva la luce del faro.}\newline «Ogni sera Tommaso saliva sul faro e accendeva la luce.»
\item\relax \textbf{(c) una piccola luce in mezzo alle onde}\newline «Elena … vide una piccola luce in mezzo alle onde.»
\item\relax \textbf{(b) Perché suo padre non era tornato una notte in cui il faro era spento.}\newline Tommaso tells Elena about his father in the last paragraph.
\item\relax \textbf{(c) trapassato prossimo: it happened before disse}\newline After a past reporting verb, an earlier event is expressed with the trapassato prossimo.
\item\relax \textbf{\textit{sorridere}}\newline sorridere → sorrisi, sorridesti, sorrise … (like ridere).
\item\relax \textbf{(a) an unreal past consequence (he did not die)}\newline Without the light he would have died: a consequence that did not happen.
\end{answers}

\answerheading{\S\ref{story:una-lettera-dalla-nonna} Una lettera dalla nonna}\label{ans:story:una-lettera-dalla-nonna}
\begin{answers}
\item\relax \textbf{(c) Insegnarle a nuotare.}\newline «ti porterò al mare e ti insegnerò a nuotare.»
\item\relax \textbf{(d) No, nessuno le aveva insegnato.}\newline «Quando ero giovane, non sapevo nuotare. Se qualcuno mi avesse insegnato …»
\item\relax \textbf{(b) an unreal present hypothesis (she is not young)}\newline se + congiuntivo imperfetto, condizionale presente: a hypothesis that is not true now.
\item\relax \textbf{(b) parta}\newline prima che requires the congiuntivo: parta.
\item\relax \textbf{\textit{Non dimenticare}}\newline The negative imperativo of tu is non + infinitive.
\item\relax \textbf{(c) Se fa bel tempo, andiamo in montagna.}\newline For a real condition, se + presente / presente is common and correct. The conditional after se (se farebbe) is non-standard.
\end{answers}

\section*{Rhymes}

\answerheading{\S\ref{rhyme:la-pioggia-sul-tetto} La pioggia sul tetto}\label{ans:rhyme:la-pioggia-sul-tetto}
\begin{answers}
\item\relax \textbf{(b) gently, slowly}\newline piano as an adverb = softly, slowly; doubling it intensifies.
\item\relax \textbf{(b) the wind}\newline Canta il vento alla finestra: il vento is the subject, placed after the verb.
\item\relax \textbf{\textit{cadere}}\newline cadere: cado, cadi, cade…
\item\relax \textbf{(d) tomorrow}\newline The present tense with domani expresses the near future.
\item\relax \textbf{\textit{letto}}\newline tetto / letto.
\end{answers}

\answerheading{\S\ref{rhyme:ninna-nanna-della-luna} Ninna nanna della luna}\label{ans:rhyme:ninna-nanna-della-luna}
\begin{answers}
\item\relax \textbf{(c) an instruction: 'sleep!'}\newline dormi is the tu imperative here (identical in form to the present tu).
\item\relax \textbf{(b) Morning is still a long way off.}\newline ancora = still; lontano = far: it is still the middle of the night.
\item\relax \textbf{(d) una stella}\newline The subject follows the verb: una stella brilla.
\item\relax \textbf{\textit{amore}}\newline cuore / amore.
\item\relax \textbf{\textit{the night is calm and beautiful} {\color{muted}(also accepted: \textit{the night is calm and lovely}; \textit{the night is quiet and beautiful}; \textit{the night is quiet and lovely})}}\newline calma and bella agree with notte.
\end{answers}

\answerheading{\S\ref{rhyme:conta-delle-stelle} Conta delle stelle}\label{ans:rhyme:conta-delle-stelle}
\begin{answers}
\item\relax \textbf{1--e, 2--c, 3--a, 4--b, 5--d {\color{muted}(tre = 3; sette = 7; nove = 9; quattro = 4; dieci = 10)}}\newline uno, due, tre, quattro, cinque, sei, sette, otto, nove, dieci.
\item\relax \textbf{(d) on the roof}\newline sul tetto ci sono le civette.
\item\relax \textbf{\textit{civette}}\newline -a → -e.
\item\relax \textbf{(d) see you tomorrow}\newline a domani is a common way to say goodbye until the next day.
\item\relax \textbf{(a) sorelle is feminine plural}\newline The adjective agrees: caro → care (feminine plural).
\end{answers}

\answerheading{\S\ref{rhyme:il-valzer-degli-animali} Il valzer degli animali}\label{ans:rhyme:il-valzer-degli-animali}
\begin{answers}
\item\relax \textbf{(c) Who dances faster than me?}\newline Chi balla più veloce di me?
\item\relax \textbf{\textit{volare}}\newline volare: volo, voli, vola…
\item\relax \textbf{(d) hand in hand}\newline nella = in + la.
\item\relax \textbf{(c) e becomes ed before a vowel, especially e}\newline The euphonic d avoids two identical vowels in a row.
\item\relax \textbf{\textit{balliamo}}\newline ball- + -iamo.
\end{answers}

\section*{Diagnostic test}\label{ans:diag}
\begin{answers}
\item\relax \textbf{(b) chiave}\newline ch before i/e is hard. {\color{muted}(Chapter~\ref{mod:m01}: Pronunciation and orthography)}
\item\relax \textbf{\textit{è}}\newline è (is) has an accent; e means 'and'. {\color{muted}(Chapter~\ref{mod:m01}: Pronunciation and orthography)}
\item\relax \textbf{(a) lo}\newline lo before z. {\color{muted}(Chapter~\ref{mod:m02}: Nouns and articles)}
\item\relax \textbf{\textit{stelle}}\newline -a → -e. {\color{muted}(Chapter~\ref{mod:m02}: Nouns and articles)}
\item\relax \textbf{\textit{bianche}}\newline Feminine plural of bianco: bianche. {\color{muted}(Chapter~\ref{mod:m03}: Adjectives)}
\item\relax \textbf{\textit{parliamo}}\newline parl- + -iamo. {\color{muted}(Chapter~\ref{mod:m04}: Regular verbs in the present)}
\item\relax \textbf{\textit{vado}}\newline Irregular: vado. {\color{muted}(Chapter~\ref{mod:m05}: Essential irregular verbs)}
\item\relax \textbf{(b) Non piango mai.}\newline non … mai. {\color{muted}(Chapter~\ref{mod:m06}: Building sentences)}
\item\relax \textbf{\textit{nel}}\newline in + il = nel. {\color{muted}(Chapter~\ref{mod:m07}: Prepositions)}
\item\relax \textbf{(d) ci}\newline noi → ci. {\color{muted}(Chapter~\ref{mod:m09}: Reflexive verbs)}
\item\relax \textbf{(a) è andata}\newline andare → essere; agreement with Sara. {\color{muted}(Chapter~\ref{mod:m10}: Passato prossimo and auxiliaries)}
\item\relax \textbf{(b) dormivo}\newline Background action in progress → imperfetto. {\color{muted}(Chapter~\ref{mod:m11}: Imperfetto and trapassato prossimo)}
\item\relax \textbf{(b) Non aprire}\newline Negative imperative for tu: non + infinitive. {\color{muted}(Chapter~\ref{mod:m15}: Imperativo)}
\item\relax \textbf{(b) perse}\newline A preceding le makes the participle agree: perse (also perdute). {\color{muted}(Chapter~\ref{mod:m10}: Passato prossimo and auxiliaries)}
\item\relax \textbf{(b) era … partito}\newline An action before another past action: trapassato prossimo. {\color{muted}(Chapter~\ref{mod:m11}: Imperfetto and trapassato prossimo)}
\item\relax \textbf{(a) sarà}\newline The futuro expresses a supposition about the present. {\color{muted}(Chapter~\ref{mod:m13}: Futuro)}
\item\relax \textbf{(b) Vorrei}\newline Condizionale presente of volere: the standard polite form (volevo is also used). {\color{muted}(Chapter~\ref{mod:m14}: Condizionale)}
\item\relax \textbf{(b) prendere}\newline prese is the passato remoto (3rd person singular) of prendere. {\color{muted}(Chapter~\ref{mod:m17}: Passato remoto)}
\item\relax \textbf{(a) abbia}\newline pensare che + congiuntivo presente (abbia). In informal speech the indicative ha is also heard; abbia is the standard. {\color{muted}(Chapter~\ref{mod:m18}: Congiuntivo)}
\item\relax \textbf{(a) avessi}\newline Hypothesis about the present: se + congiuntivo imperfetto, condizionale in the main clause. {\color{muted}(Chapter~\ref{mod:m19}: Complex sentences)}
\end{answers}

\section*{Mixed practice}

\answerheading{\S\ref{set:mixed-articles} Mixed practice: articles}\label{ans:set:mixed-articles}
\begin{answers}
\item\relax \textbf{(b) lo}\newline s + consonant → lo.
\item\relax \textbf{(d) l'}\newline Feminine before a vowel → l'.
\item\relax \textbf{(c) un'}\newline Feminine before a vowel → un'.
\item\relax \textbf{(d) gli}\newline Masculine plural before a vowel → gli.
\item\relax \textbf{\textit{alle}}\newline a + le = alle.
\item\relax \textbf{\textit{sugli}}\newline su + gli = sugli.
\end{answers}

\answerheading{\S\ref{set:mixed-present} Mixed practice: the present tense}\label{ans:set:mixed-present}
\begin{answers}
\item\relax \textbf{\textit{giochi}}\newline h keeps g hard: giochi.
\item\relax \textbf{\textit{capiscono}}\newline -isc- verb.
\item\relax \textbf{\textit{facciamo}}\newline fare: facciamo.
\item\relax \textbf{\textit{vieni}}\newline venire: vieni.
\item\relax \textbf{\textit{andare}}\newline andare: vanno.
\item\relax \textbf{\textit{uscire}}\newline uscire: esce.
\end{answers}

\answerheading{\S\ref{set:mixed-past} Mixed practice: past tenses}\label{ans:set:mixed-past}
\begin{answers}
\item\relax \textbf{\textit{ho letto}}\newline leggere → letto.
\item\relax \textbf{\textit{sono partiti}}\newline partire → essere; agreement.
\item\relax \textbf{\textit{eravamo}}\newline Irregular imperfetto.
\item\relax \textbf{\textit{diceva}}\newline dice- + va.
\item\relax \textbf{(a) giocavo}\newline Habit in the past → imperfetto.
\item\relax \textbf{(b) sono svegliata}\newline Reflexive → essere; agreement with a female speaker.
\end{answers}

\answerheading{\S\ref{set:mixed-future-conditional} Mixed practice: futuro and condizionale}\label{ans:set:mixed-future-conditional}
\begin{answers}
\item\relax \textbf{venire, io, futuro: \textit{verrò}; venire, io, condizionale: \textit{verrei}; essere, noi, futuro: \textit{saremo}; essere, noi, condizionale: \textit{saremmo}; fare, loro, futuro: \textit{faranno}; fare, loro, condizionale: \textit{farebbero}}\newline Futuro and condizionale share the stem: verr-, sar-, far-. Note the double m in saremmo (condizionale) vs saremo (futuro).
\item\relax \textbf{(b) sarebbe venuto}\newline After a past verb of saying, a later event is expressed with the condizionale passato. In everyday speech verrà is also used when the event is still in the future; verrebbe is not standard here.
\item\relax \textbf{(c) information that has not been confirmed}\newline The condizionale of reported information (typical of journalism): "the minister is said to have resigned".
\item\relax \textbf{(a) avrai finito}\newline Anteriority in the future: futuro anteriore. Informally also: quando hai finito, usciamo.
\item\relax \textbf{\textit{Ti avrei chiamato, ma non avevo il tuo numero.} {\color{muted}(also accepted: \textit{Ti avrei chiamata, ma non avevo il tuo numero.}; \textit{Ti avrei telefonato, ma non avevo il tuo numero.}; \textit{Io ti avrei chiamato, ma non avevo il tuo numero.}; \textit{Ti avrei chiamato ma non avevo il tuo numero.}; \textit{Ti avrei chiamata ma non avevo il tuo numero.})}}\newline Unrealised past action: condizionale passato. chiamato/chiamata: agreement with ti is optional.
\end{answers}

\answerheading{\S\ref{set:mixed-moods} Mixed practice: indicativo, congiuntivo or condizionale?}\label{ans:set:mixed-moods}
\begin{answers}
\item\relax \textbf{(c) abita}\newline sapere expresses certainty: indicativo.
\item\relax \textbf{(d) abiti}\newline credere che expresses opinion: congiuntivo in the standard. The indicative abita is common in informal speech.
\item\relax \textbf{(c) abitasse}\newline Past main verb + simultaneous action: congiuntivo imperfetto.
\item\relax \textbf{(1) \textit{potessi}; (2) \textit{avresti passato}}\newline Present unreal hypothesis: se + congiuntivo imperfetto (potessi). Past counterfactual: condizionale passato in the main clause; passare un esame is transitive, so it takes avere (avresti passato).
\item\relax \textbf{(c) abbia}\newline A relative clause describing something sought but not yet identified usually takes the congiuntivo. With the indicativo (che ha) the speaker refers to a specific house that exists.
\item\relax \textbf{(a) fosse}\newline benché always takes the congiuntivo; with a past main verb, the congiuntivo imperfetto.
\end{answers}

\chapter{Vocabulary by topic}\label{glossary}
All 503 words of the course, grouped by topic and then alphabetically. Nouns are given with their article and plural, verbs with their auxiliary and participle (full forms: Chapter~\ref{ref:principal}), adjectives with their four forms. The level is the approximate stage at which the word is introduced. The index at the end of the book lists every Italian word alphabetically.\par

\section{Actions}
\begin{glosslist}
\glossentry{\textbf{accendere}\index{accendere@accendere}}{verb, -ere, auxiliary \textit{avere}, participle \textit{acceso}}{to switch on; to light}{Accendo la lampada.}{I switch on the lamp.}
\glossentry{\textbf{aiutare}\index{aiutare@aiutare}}{verb, -are, auxiliary \textit{avere}, participle \textit{aiutato}}{to help}{Aiuto la mamma in cucina.}{I help Mum in the kitchen.}\glossnote{Constructions: aiutare qualcuno a + infinitive.}
\glossentry{\textbf{aprire}\index{aprire@aprire}}{verb, -ire, auxiliary \textit{avere}, participle \textit{aperto}}{to open}{Apro la finestra.}{I open the window.}
\glossentry{\textbf{aspettare}\index{aspettare@aspettare}}{verb, -are, auxiliary \textit{avere}, participle \textit{aspettato}}{to wait (for)}{Aspetto il treno.}{I'm waiting for the train.}\glossnote{Constructions: Takes a direct object: aspettare qualcuno (no preposition).}
\glossentry{\textbf{bisognare}\index{bisognare@bisognare}}{verb, -are, auxiliary \textit{essere}, participle \textit{bisognato}}{to be necessary (impersonal)}{Bisogna partire presto.}{We need to leave early.}\glossnote{Constructions: bisogna + infinitive; bisogna che + congiuntivo. Impersonal: bisogna + infinitive (bisogna partire) or bisogna che + congiuntivo (bisogna che tu parta). Compound tenses are rare.}
\glossentry{\textbf{bussare}\index{bussare@bussare}}{verb, -are, auxiliary \textit{avere}, participle \textit{bussato}}{to knock}{Qualcuno bussa alla porta.}{Someone is knocking at the door.}\glossnote{Constructions: bussare a + door/window.}
\glossentry{\textbf{cambiare}\index{cambiare@cambiare}}{verb, -are, auxiliary \textit{avere}, participle \textit{cambiato}}{to change}{Il tempo è cambiato.}{The weather has changed.}\glossnote{Constructions: Auxiliary essere when intransitive (il tempo è cambiato), avere when transitive (ho cambiato idea).}
\glossentry{\textbf{cercare}\index{cercare@cercare}}{verb, -are, auxiliary \textit{avere}, participle \textit{cercato}}{to look for; to try}{Il cane cerca la palla.}{The dog is looking for the ball.}\glossnote{Constructions: cercare di + infinitive = to try to.}
\glossentry{\textbf{chiudere}\index{chiudere@chiudere}}{verb, -ere, auxiliary \textit{avere}, participle \textit{chiuso}}{to close}{Chiudo la porta.}{I close the door.}
\glossentry{\textbf{cominciare}\index{cominciare@cominciare}}{verb, -are, auxiliary \textit{avere}, participle \textit{cominciato}}{to begin; to start}{La scuola comincia domani.}{School starts tomorrow.}\glossnote{Constructions: cominciare a + infinitive.}
\glossentry{\textbf{continuare}\index{continuare@continuare}}{verb, -are, auxiliary \textit{avere}, participle \textit{continuato}}{to continue}{Continua a piovere.}{It keeps raining.}\glossnote{Constructions: continuare a + infinitive.}
\glossentry{\textbf{costruire}\index{costruire@costruire}}{verb, -ire (-isc-), auxiliary \textit{avere}, participle \textit{costruito}}{to build}{L'uccello costruisce il nido.}{The bird builds its nest.}
\glossentry{\textbf{decidere}\index{decidere@decidere}}{verb, -ere, auxiliary \textit{avere}, participle \textit{deciso}}{to decide}{Abbiamo deciso di partire.}{We have decided to leave.}\glossnote{Constructions: decidere di + infinitive.}
\glossentry{\textbf{finire}\index{finire@finire}}{verb, -ire (-isc-), auxiliary \textit{avere}, participle \textit{finito}}{to finish; to end}{La lezione finisce alle due.}{The lesson finishes at two.}\glossnote{Constructions: Auxiliary avere with an object (ho finito il libro), essere without (la lezione è finita).}
\glossentry{\textbf{giocare}\index{giocare@giocare}}{verb, -are, auxiliary \textit{avere}, participle \textit{giocato}}{to play}{I bambini giocano in giardino.}{The children play in the garden.}\glossnote{Constructions: giocare a + game: giocare a palla.}
\glossentry{\textbf{guardare}\index{guardare@guardare}}{verb, -are, auxiliary \textit{avere}, participle \textit{guardato}}{to look at; to watch}{Guardo le stelle.}{I look at the stars.}
\glossentry{\textbf{insegnare}\index{insegnare@insegnare}}{verb, -are, auxiliary \textit{avere}, participle \textit{insegnato}}{to teach}{La maestra insegna l'italiano.}{The teacher teaches Italian.}\glossnote{Constructions: insegnare qualcosa a qualcuno; insegnare a qualcuno a + infinitive: ti insegno a nuotare.}
\glossentry{\textbf{lasciare}\index{lasciare@lasciare}}{verb, -are, auxiliary \textit{avere}, participle \textit{lasciato}}{to leave (something); to let}{Ho lasciato l'ombrello a casa.}{I left the umbrella at home.}\glossnote{Constructions: lasciare + infinitive = to let (lascia parlare).}
\glossentry{\textbf{mettere}\index{mettere@mettere}}{verb, -ere, auxiliary \textit{avere}, participle \textit{messo}}{to put}{Metto il libro sul tavolo.}{I put the book on the table.}\glossnote{Constructions: mettersi = to put on (clothes).}
\glossentry{\textbf{nascondere}\index{nascondere@nascondere}}{verb, -ere, auxiliary \textit{avere}, participle \textit{nascosto}}{to hide}{Lo scoiattolo nasconde le noci.}{The squirrel hides the nuts.}
\glossentry{\textbf{offrire}\index{offrire@offrire}}{verb, -ire, auxiliary \textit{avere}, participle \textit{offerto}}{to offer}{Ti offro un caffè.}{I'll buy you a coffee.}
\glossentry{\textbf{perdere}\index{perdere@perdere}}{verb, -ere, auxiliary \textit{avere}, participle \textit{perso}}{to lose; to miss}{Perdo sempre la chiave.}{I always lose the key.}
\glossentry{\textbf{pescare}\index{pescare@pescare}}{verb, -are, auxiliary \textit{avere}, participle \textit{pescato}}{to fish; to catch (fish)}{Il nonno pesca nel lago.}{Grandpa fishes in the lake.}
\glossentry{\textbf{portare}\index{portare@portare}}{verb, -are, auxiliary \textit{avere}, participle \textit{portato}}{to bring; to carry; to wear}{Porto un regalo.}{I'm bringing a present.}
\glossentry{\textbf{prendere}\index{prendere@prendere}}{verb, -ere, auxiliary \textit{avere}, participle \textit{preso}}{to take; to catch; to have (food)}{Prendo l'ombrello.}{I take the umbrella.}
\glossentry{\textbf{riuscire}\index{riuscire@riuscire}}{verb, -ire, auxiliary \textit{essere}, participle \textit{riuscito}; irregular: \textit{io riesco, tu riesci, lui/lei riesce, loro riescono}}{to manage; to succeed}{Non riesco a dormire.}{I can't manage to sleep.}\glossnote{Constructions: riuscire a + infinitive. Conjugated like uscire: riesco, riesci, riesce, riusciamo, riuscite, riescono. riuscire a + infinitive = to manage to.}
\glossentry{\textbf{rompere}\index{rompere@rompere}}{verb, -ere, auxiliary \textit{avere}, participle \textit{rotto}}{to break}{Il bambino ha rotto il bicchiere.}{The child broke the glass.}
\glossentry{\textbf{scegliere}\index{scegliere@scegliere}}{verb, -ere, auxiliary \textit{avere}, participle \textit{scelto}; irregular: \textit{io scelgo, loro scelgono}}{to choose}{Scegli tu il libro.}{You choose the book.}\glossnote{gli becomes lg before -o: scelgo, scelgono. Same pattern: raccogliere (raccolgo), togliere (tolgo).}
\glossentry{\textbf{scoprire}\index{scoprire@scoprire}}{verb, -ire, auxiliary \textit{avere}, participle \textit{scoperto}}{to discover; to uncover}{Hanno scoperto un'isola.}{They discovered an island.}
\glossentry{\textbf{sentire}\index{sentire@sentire}}{verb, -ire, auxiliary \textit{avere}, participle \textit{sentito}}{to hear; to feel}{Sento un rumore.}{I hear a noise.}
\glossentry{\textbf{servire}\index{servire@servire}}{verb, -ire, auxiliary \textit{avere}, participle \textit{servito}}{to serve; to be needed}{Mi serve una penna.}{I need a pen.}\glossnote{Constructions: servire (intransitive, essere) = to be useful/needed: mi è servito.}
\glossentry{\textbf{smettere}\index{smettere@smettere}}{verb, -ere, auxiliary \textit{avere}, participle \textit{smesso}}{to stop (doing)}{Ha smesso di piovere.}{It has stopped raining.}\glossnote{Constructions: smettere di + infinitive; participle smesso.}
\glossentry{\textbf{spegnere}\index{spegnere@spegnere}}{verb, -ere, auxiliary \textit{avere}, participle \textit{spento}; irregular: \textit{io spengo, loro spengono}}{to switch off; to put out}{Spegni la luce, per favore.}{Switch off the light, please.}\glossnote{Constructions: Present: spengo, spegni, spegne, spegniamo, spegnete, spengono. gn becomes ng before -o: spengo, spengono (like tengo, vengo). Participle spento; passato remoto spensi.}
\glossentry{\textbf{succedere}\index{succedere@succedere}}{verb, -ere, auxiliary \textit{essere}, participle \textit{successo}}{to happen}{Che cosa è successo?}{What happened?}\glossnote{Constructions: Mostly third person: succede, è successo.}
\glossentry{\textbf{tenere}\index{tenere@tenere}}{verb, -ere, auxiliary \textit{avere}, participle \textit{tenuto}; irregular: \textit{io tengo, tu tieni, lui/lei tiene, loro tengono}}{to hold; to keep}{Tengo la chiave in mano.}{I am holding the key in my hand.}\glossnote{Exactly the same pattern as venire: tengo, tieni, tiene, …, tengono.}
\glossentry{\textbf{trovare}\index{trovare@trovare}}{verb, -are, auxiliary \textit{avere}, participle \textit{trovato}}{to find}{Trovo una conchiglia.}{I find a shell.}
\glossentry{\textbf{vedere}\index{vedere@vedere}}{verb, -ere, auxiliary \textit{avere}, participle \textit{visto}}{to see}{Vedo il mare.}{I see the sea.}
\glossentry{\textbf{vincere}\index{vincere@vincere}}{verb, -ere, auxiliary \textit{avere}, participle \textit{vinto}}{to win}{La tartaruga vince la gara.}{The tortoise wins the race.}
\end{glosslist}

\section{Adverbs}
\begin{glosslist}
\glossentry{\textbf{almeno}\index{almeno@almeno}}{adverb}{at least}{Dormi almeno otto ore.}{Sleep at least eight hours.}
\glossentry{\textbf{bene}\index{bene@bene}}{adverb}{well}{Sto bene, grazie.}{I'm fine, thanks.}
\glossentry{\textbf{finalmente}\index{finalmente@finalmente}}{adverb}{at last; finally}{Finalmente arriva la primavera.}{At last spring arrives.}
\glossentry{\textbf{insieme}\index{insieme@insieme}}{adverb}{together}{Giochiamo insieme.}{Let's play together.}
\glossentry{\textbf{là}\index{la@là}}{adverb}{there}{Il cane è là.}{The dog is there.}\glossnote{Written with an accent to distinguish it from the article/pronoun la. lì has the same meaning.}
\glossentry{\textbf{lentamente}\index{lentamente@lentamente}}{adverb}{slowly}{La lumaca cammina lentamente.}{The snail moves slowly.}\glossnote{Formed from the feminine of the adjective + -mente: lenta → lentamente.}
\glossentry{\textbf{lì}\index{li@lì}}{adverb}{there}{Il libro è lì, sul tavolo.}{The book is there, on the table.}\glossnote{lì and là both mean 'there'; là is often a little more distant or vague. Accented lì/là are adverbs; li without an accent is a pronoun (them).}
\glossentry{\textbf{male}\index{male@male}}{adverb}{badly; ill}{Il cane sta male.}{The dog is unwell.}
\glossentry{\textbf{molto}\index{molto@molto}}{adverb}{very; a lot; much}{Il gatto è molto bello.}{The cat is very beautiful.}\glossnote{As an adverb it is invariable (molto bella). As an adjective it agrees: molti gatti, molte stelle.}
\glossentry{\textbf{non}\index{non@non}}{adverb}{not}{Non ho fame.}{I'm not hungry.}\glossnote{Always placed before the conjugated verb (and before any object pronouns): non lo so.}
\glossentry{\textbf{piano}\index{piano@piano}}{adverb}{slowly; softly; quietly}{Parla piano, il bambino dorme.}{Speak quietly, the child is asleep.}\glossnote{Also a noun: il piano = floor (storey), piano(forte).}
\glossentry{\textbf{poco}\index{poco@poco}}{adverb}{little; not much; not very}{Mangio poco.}{I don't eat much.}\glossnote{As an adjective it agrees: pochi amici, poche parole.}
\glossentry{\textbf{qui}\index{qui@qui}}{adverb}{here}{Il gatto è qui.}{The cat is here.}\glossnote{qua has the same meaning.}
\glossentry{\textbf{solo}\index{solo@solo}}{adverb}{only}{Ho solo un fratello.}{I only have one brother.}\glossnote{Also an adjective (alone): sono solo, è sola. As an adverb solo = soltanto.}
\end{glosslist}

\section{Animals}
\begin{glosslist}
\glossentry{\textbf{l'anatra}\index{anatra@anatra}}{feminine; plural \textit{le anatre}}{duck}{L'anatra nuota nel lago.}{The duck swims in the lake.}
\glossentry{\textbf{l'animale}\index{animale@animale}}{masculine; plural \textit{gli animali}}{animal}{Il bosco è la casa di molti animali.}{The wood is home to many animals.}
\glossentry{\textbf{l'ape}\index{ape@ape}}{feminine; plural \textit{le api}}{bee}{L'ape fa il miele.}{The bee makes honey.}
\glossentry{\textbf{l'asino}\index{asino@asino}}{masculine; plural \textit{gli asini}}{donkey}{L'asino porta il cesto.}{The donkey carries the basket.}
\glossentry{\textbf{il cane}\index{cane@cane}}{masculine; plural \textit{i cani}}{dog}{Il cane corre nel prato.}{The dog runs in the meadow.}
\glossentry{\textbf{il cavallo}\index{cavallo@cavallo}}{masculine; plural \textit{i cavalli}}{horse}{Il cavallo mangia l'erba.}{The horse eats the grass.}
\glossentry{\textbf{la civetta}\index{civetta@civetta}}{feminine; plural \textit{le civette}}{little owl}{La civetta canta di notte.}{The little owl calls at night.}
\glossentry{\textbf{il coniglio}\index{coniglio@coniglio}}{masculine; plural \textit{i conigli}}{rabbit}{Il coniglio mangia una carota.}{The rabbit eats a carrot.}
\glossentry{\textbf{il delfino}\index{delfino@delfino}}{masculine; plural \textit{i delfini}}{dolphin}{Il delfino salta nel mare.}{The dolphin jumps in the sea.}
\glossentry{\textbf{la farfalla}\index{farfalla@farfalla}}{feminine; plural \textit{le farfalle}}{butterfly}{La farfalla vola sui fiori.}{The butterfly flies over the flowers.}
\glossentry{\textbf{la formica}\index{formica@formica}}{feminine; plural \textit{le formiche}}{ant}{Le formiche lavorano insieme.}{The ants work together.}
\glossentry{\textbf{la gallina}\index{gallina@gallina}}{feminine; plural \textit{le galline}}{hen}{La gallina fa un uovo.}{The hen lays an egg.}
\glossentry{\textbf{il gatto}\index{gatto@gatto}}{masculine; plural \textit{i gatti}}{cat}{Il gatto dorme sul letto.}{The cat sleeps on the bed.}
\glossentry{\textbf{il granchio}\index{granchio@granchio}}{masculine; plural \textit{i granchi}}{crab}{Il granchio cammina sulla sabbia.}{The crab walks on the sand.}
\glossentry{\textbf{il gufo}\index{gufo@gufo}}{masculine; plural \textit{i gufi}}{owl}{Il gufo guarda la luna.}{The owl looks at the moon.}
\glossentry{\textbf{il leone}\index{leone@leone}}{masculine; plural \textit{i leoni}}{lion}{Il leone è forte.}{The lion is strong.}
\glossentry{\textbf{la lepre}\index{lepre@lepre}}{feminine; plural \textit{le lepri}}{hare}{La lepre corre veloce.}{The hare runs fast.}
\glossentry{\textbf{la lumaca}\index{lumaca@lumaca}}{feminine; plural \textit{le lumache}}{snail}{La lumaca è lenta.}{The snail is slow.}
\glossentry{\textbf{il lupo}\index{lupo@lupo}}{masculine; plural \textit{i lupi}}{wolf}{Il lupo cammina di notte.}{The wolf walks at night.}
\glossentry{\textbf{il maiale}\index{maiale@maiale}}{masculine; plural \textit{i maiali}}{pig}{Il maiale mangia molto.}{The pig eats a lot.}
\glossentry{\textbf{la mucca}\index{mucca@mucca}}{feminine; plural \textit{le mucche}}{cow}{La mucca dà il latte.}{The cow gives milk.}
\glossentry{\textbf{l'orso}\index{orso@orso}}{masculine; plural \textit{gli orsi}}{bear}{L'orso dorme in inverno.}{The bear sleeps in winter.}
\glossentry{\textbf{la pecora}\index{pecora@pecora}}{feminine; plural \textit{le pecore}}{sheep}{Le pecore sono nel prato.}{The sheep are in the meadow.}
\glossentry{\textbf{il pesce}\index{pesce@pesce}}{masculine; plural \textit{i pesci}}{fish}{Il pesce nuota nel fiume.}{The fish swims in the river.}
\glossentry{\textbf{la rana}\index{rana@rana}}{feminine; plural \textit{le rane}}{frog}{La rana salta nel lago.}{The frog jumps into the lake.}
\glossentry{\textbf{il riccio}\index{riccio@riccio}}{masculine; plural \textit{i ricci}}{hedgehog}{Il riccio dorme sotto le foglie.}{The hedgehog sleeps under the leaves.}
\glossentry{\textbf{lo scoiattolo}\index{scoiattolo@scoiattolo}}{masculine; plural \textit{gli scoiattoli}}{squirrel}{Lo scoiattolo sale sull'albero.}{The squirrel climbs the tree.}
\glossentry{\textbf{la tartaruga}\index{tartaruga@tartaruga}}{feminine; plural \textit{le tartarughe}}{tortoise; turtle}{La tartaruga è molto lenta.}{The tortoise is very slow.}
\glossentry{\textbf{il topo}\index{topo@topo}}{masculine; plural \textit{i topi}}{mouse}{Il topo ha paura del gatto.}{The mouse is afraid of the cat.}
\glossentry{\textbf{l'uccello}\index{uccello@uccello}}{masculine; plural \textit{gli uccelli}}{bird}{Un uccello canta sull'albero.}{A bird sings in the tree.}
\glossentry{\textbf{la volpe}\index{volpe@volpe}}{feminine; plural \textit{le volpi}}{fox}{La volpe vive nel bosco.}{The fox lives in the wood.}
\end{glosslist}

\section{Body}
\begin{glosslist}
\glossentry{\textbf{il braccio}\index{braccio@braccio}}{masculine; plural \textit{le braccia} (feminine)}{arm}{La mamma ha il bambino in braccio.}{Mum holds the child in her arms.}
\glossentry{\textbf{la mano}\index{mano@mano}}{feminine; plural \textit{le mani}}{hand}{Lavo le mani prima di mangiare.}{I wash my hands before eating.}
\glossentry{\textbf{l'occhio}\index{occhio@occhio}}{masculine; plural \textit{gli occhi}}{eye}{Il gatto ha gli occhi verdi.}{The cat has green eyes.}
\glossentry{\textbf{la pancia}\index{pancia@pancia}}{feminine; plural \textit{le pance}}{belly; tummy}{L'orso ha la pancia vuota.}{The bear has an empty belly.}
\glossentry{\textbf{il piede}\index{piede@piede}}{masculine; plural \textit{i piedi}}{foot}{Vado a scuola a piedi.}{I go to school on foot.}
\glossentry{\textbf{la testa}\index{testa@testa}}{feminine; plural \textit{le teste}}{head}{Il gufo gira la testa.}{The owl turns its head.}
\end{glosslist}

\section{Colours}
\begin{glosslist}
\glossentry{\textbf{azzurro}\index{azzurro@azzurro}}{adjective: \textit{azzurro, azzurra, azzurri, azzurre}}{light blue; sky blue}{Il cielo è azzurro.}{The sky is blue.}
\glossentry{\textbf{bianco}\index{bianco@bianco}}{adjective: \textit{bianco, bianca, bianchi, bianche}}{white}{La neve è bianca.}{Snow is white.}
\glossentry{\textbf{blu}\index{blu@blu}}{adjective: \textit{blu, blu, blu, blu}}{blue (dark)}{Il mare è blu.}{The sea is blue.}\glossnote{Invariable: il mare blu, le barche blu.}
\glossentry{\textbf{giallo}\index{giallo@giallo}}{adjective: \textit{giallo, gialla, gialli, gialle}}{yellow}{Il sole è giallo.}{The sun is yellow.}
\glossentry{\textbf{grigio}\index{grigio@grigio}}{adjective: \textit{grigio, grigia, grigi, grigie}}{grey}{Il cielo è grigio.}{The sky is grey.}\glossnote{Feminine plural grigie (also written grige).}
\glossentry{\textbf{nero}\index{nero@nero}}{adjective: \textit{nero, nera, neri, nere}}{black}{Il gatto nero dorme.}{The black cat is sleeping.}
\glossentry{\textbf{rosa}\index{rosa@rosa}}{adjective: \textit{rosa, rosa, rosa, rosa}}{pink}{La bambina ha un vestito rosa.}{The little girl has a pink dress.}\glossnote{Invariable: le farfalle rosa.}
\glossentry{\textbf{rosso}\index{rosso@rosso}}{adjective: \textit{rosso, rossa, rossi, rosse}}{red}{La fragola è rossa.}{The strawberry is red.}
\glossentry{\textbf{verde}\index{verde@verde}}{adjective: \textit{verde, verde, verdi, verdi}}{green}{L'erba è verde.}{The grass is green.}
\end{glosslist}

\section{Communication}
\begin{glosslist}
\glossentry{\textbf{annunciare}\index{annunciare@annunciare}}{verb, -are, auxiliary \textit{avere}, participle \textit{annunciato}}{to announce}{Hanno annunciato la notizia.}{They announced the news.}
\glossentry{\textbf{capire}\index{capire@capire}}{verb, -ire (-isc-), auxiliary \textit{avere}, participle \textit{capito}}{to understand}{Capisco l'italiano.}{I understand Italian.}
\glossentry{\textbf{chiamare}\index{chiamare@chiamare}}{verb, -are, auxiliary \textit{avere}, participle \textit{chiamato}}{to call}{La mamma chiama i bambini.}{Mum calls the children.}
\glossentry{\textbf{chiamarsi}\index{chiamarsi@chiamarsi}}{verb, -are, auxiliary \textit{essere}, participle \textit{chiamato}}{to be called}{Mi chiamo Sara.}{My name is Sara.}\glossnote{Constructions: reflexive.}
\glossentry{\textbf{chiedere}\index{chiedere@chiedere}}{verb, -ere, auxiliary \textit{avere}, participle \textit{chiesto}}{to ask}{Il bambino chiede un gelato.}{The child asks for an ice cream.}\glossnote{Constructions: chiedere qualcosa a qualcuno.}
\glossentry{\textbf{conoscere}\index{conoscere@conoscere}}{verb, -ere, auxiliary \textit{avere}, participle \textit{conosciuto}}{to know (a person or place)}{Conosco Sara.}{I know Sara.}
\glossentry{\textbf{credere}\index{credere@credere}}{verb, -ere, auxiliary \textit{avere}, participle \textit{creduto}}{to believe}{Credo di sì.}{I think so.}\glossnote{Constructions: credere a/in = to believe (in).}
\glossentry{\textbf{dimenticare}\index{dimenticare@dimenticare}}{verb, -are, auxiliary \textit{avere}, participle \textit{dimenticato}}{to forget}{Ho dimenticato la chiave.}{I forgot the key.}\glossnote{Constructions: dimenticare qualcosa / dimenticarsi di qualcosa.}
\glossentry{\textbf{dire}\index{dire@dire}}{verb, -ire, auxiliary \textit{avere}, participle \textit{detto}; irregular: \textit{io dico, tu dici, lui/lei dice, noi diciamo, loro dicono}}{to say; to tell}{Il bambino dice la verità.}{The child tells the truth.}\glossnote{Constructions: dire qualcosa a qualcuno = to say something to someone. From Latin dicere: the stem dic- appears everywhere except dite. The c is hard in dico, dicono and soft in dici, dice.}
\glossentry{\textbf{dubitare}\index{dubitare@dubitare}}{verb, -are, auxiliary \textit{avere}, participle \textit{dubitato}}{to doubt}{Dubito che sia vero.}{I doubt that it is true.}\glossnote{Constructions: dubitare che + congiuntivo; dubitare di + noun.}
\glossentry{\textbf{gridare}\index{gridare@gridare}}{verb, -are, auxiliary \textit{avere}, participle \textit{gridato}}{to shout}{Non gridare, per favore!}{Don't shout, please!}
\glossentry{\textbf{l'italiano}\index{italiano@italiano}}{masculine; no plural}{Italian (the language)}{Studio l'italiano.}{I'm studying Italian.}\glossnote{Names of languages are masculine and written with a small letter. After parlare the article is usually dropped: parlo italiano.}
\glossentry{\textbf{la lettera}\index{lettera@lettera}}{feminine; plural \textit{le lettere}}{letter}{Scrivo una lettera alla nonna.}{I am writing a letter to Grandma.}
\glossentry{\textbf{mandare}\index{mandare@mandare}}{verb, -are, auxiliary \textit{avere}, participle \textit{mandato}}{to send}{Ti mando una lettera.}{I'm sending you a letter.}\glossnote{Constructions: mandare qualcosa a qualcuno.}
\glossentry{\textbf{parlare}\index{parlare@parlare}}{verb, -are, auxiliary \textit{avere}, participle \textit{parlato}}{to speak; to talk}{Parlo italiano.}{I speak Italian.}\glossnote{Constructions: parlare con = to talk to; parlare di = to talk about.}
\glossentry{\textbf{pensare}\index{pensare@pensare}}{verb, -are, auxiliary \textit{avere}, participle \textit{pensato}}{to think}{Penso a te.}{I'm thinking of you.}\glossnote{Constructions: pensare a = to think about; pensare di + infinitive = to plan to.}
\glossentry{\textbf{permettere}\index{permettere@permettere}}{verb, -ere, auxiliary \textit{avere}, participle \textit{permesso}}{to allow}{Il papà mi permette di uscire.}{Dad allows me to go out.}\glossnote{Constructions: permettere a qualcuno di + infinitive.}
\glossentry{\textbf{promettere}\index{promettere@promettere}}{verb, -ere, auxiliary \textit{avere}, participle \textit{promesso}}{to promise}{Ti prometto che tornerò.}{I promise you I'll come back.}\glossnote{Constructions: promettere di + infinitive.}
\glossentry{\textbf{raccontare}\index{raccontare@raccontare}}{verb, -are, auxiliary \textit{avere}, participle \textit{raccontato}}{to tell (a story)}{Il nonno racconta una storia.}{Grandfather tells a story.}\glossnote{Constructions: raccontare qualcosa a qualcuno.}
\glossentry{\textbf{ricevere}\index{ricevere@ricevere}}{verb, -ere, auxiliary \textit{avere}, participle \textit{ricevuto}}{to receive}{Ho ricevuto la tua lettera.}{I received your letter.}
\glossentry{\textbf{ricordare}\index{ricordare@ricordare}}{verb, -are, auxiliary \textit{avere}, participle \textit{ricordato}}{to remember; to remind}{Ricordo quella sera.}{I remember that evening.}\glossnote{Constructions: ricordare qualcosa / ricordarsi di qualcosa.}
\glossentry{\textbf{rispondere}\index{rispondere@rispondere}}{verb, -ere, auxiliary \textit{avere}, participle \textit{risposto}}{to answer; to reply}{Rispondo alla maestra.}{I answer the teacher.}\glossnote{Constructions: rispondere a qualcuno.}
\glossentry{\textbf{il segreto}\index{segreto@segreto}}{masculine; plural \textit{i segreti}}{secret}{Ti racconto un segreto.}{I'll tell you a secret.}
\glossentry{\textbf{spiegare}\index{spiegare@spiegare}}{verb, -are, auxiliary \textit{avere}, participle \textit{spiegato}}{to explain}{La maestra spiega la regola.}{The teacher explains the rule.}\glossnote{Constructions: spiegare qualcosa a qualcuno.}
\glossentry{\textbf{telefonare}\index{telefonare@telefonare}}{verb, -are, auxiliary \textit{avere}, participle \textit{telefonato}}{to phone}{Telefono alla nonna.}{I phone Grandma.}\glossnote{Constructions: telefonare a qualcuno (indirect object: le telefono).}
\end{glosslist}

\section{Conjunctions}
\begin{glosslist}
\glossentry{\textbf{allora}\index{allora@allora}}{adverb}{so; then; at that time}{Allora, andiamo?}{So, shall we go?}
\glossentry{\textbf{anche}\index{anche@anche}}{adverb}{also; too; even}{Anche io voglio un gelato.}{I want an ice cream too.}
\glossentry{\textbf{e}\index{e@e}}{conjunction}{and}{Il gatto e il cane.}{The cat and the dog.}\glossnote{Before a vowel it may become ed: ed ecco.}
\glossentry{\textbf{ma}\index{ma@ma}}{conjunction}{but}{È piccolo ma forte.}{He is small but strong.}
\glossentry{\textbf{mentre}\index{mentre@mentre}}{conjunction}{while}{Canto mentre lavo i piatti.}{I sing while I wash the dishes.}
\glossentry{\textbf{o}\index{o@o}}{conjunction}{or}{Vuoi il tè o il caffè?}{Do you want tea or coffee?}
\glossentry{\textbf{perché}\index{perche@perché}}{conjunction}{because; why}{Perché piangi? Perché sono stanco.}{Why are you crying? Because I'm tired.}\glossnote{Written with an acute accent (é). Means both 'why' (in questions) and 'because'.}
\glossentry{\textbf{poi}\index{poi@poi}}{adverb}{then; afterwards}{Mangio e poi dormo.}{I eat and then I sleep.}
\glossentry{\textbf{quando}\index{quando@quando}}{conjunction}{when}{Quando piove, resto a casa.}{When it rains, I stay at home.}\glossnote{Also a question word: Quando parti?}
\glossentry{\textbf{se}\index{se@se}}{conjunction}{if; whether}{Se piove, prendo l'ombrello.}{If it rains, I'll take the umbrella.}
\end{glosslist}

\section{Daily life}
\begin{glosslist}
\glossentry{\textbf{abitare}\index{abitare@abitare}}{verb, -are, auxiliary \textit{avere}, participle \textit{abitato}}{to live (reside)}{Abito in una piccola casa.}{I live in a small house.}\glossnote{Constructions: abitare a + city; abitare in + street, region, country.}
\glossentry{\textbf{addormentarsi}\index{addormentarsi@addormentarsi}}{verb, -are, auxiliary \textit{essere}, participle \textit{addormentato}}{to fall asleep}{L'orso si addormenta.}{The bear falls asleep.}\glossnote{Constructions: reflexive.}
\glossentry{\textbf{alzarsi}\index{alzarsi@alzarsi}}{verb, -are, auxiliary \textit{essere}, participle \textit{alzato}}{to get up}{Mi alzo presto.}{I get up early.}\glossnote{Constructions: reflexive.}
\glossentry{\textbf{comprare}\index{comprare@comprare}}{verb, -are, auxiliary \textit{avere}, participle \textit{comprato}}{to buy}{Compro il pane.}{I buy bread.}
\glossentry{\textbf{costare}\index{costare@costare}}{verb, -are, auxiliary \textit{essere}, participle \textit{costato}}{to cost}{Quanto costa?}{How much does it cost?}
\glossentry{\textbf{dormire}\index{dormire@dormire}}{verb, -ire, auxiliary \textit{avere}, participle \textit{dormito}}{to sleep}{Il gatto dorme.}{The cat sleeps.}
\glossentry{\textbf{imparare}\index{imparare@imparare}}{verb, -are, auxiliary \textit{avere}, participle \textit{imparato}}{to learn}{Imparo una canzone.}{I am learning a song.}\glossnote{Constructions: imparare a + infinitive.}
\glossentry{\textbf{lavare}\index{lavare@lavare}}{verb, -are, auxiliary \textit{avere}, participle \textit{lavato}}{to wash}{Lavo i piatti.}{I wash the dishes.}
\glossentry{\textbf{lavarsi}\index{lavarsi@lavarsi}}{verb, -are, auxiliary \textit{essere}, participle \textit{lavato}}{to wash (oneself)}{Il gatto si lava.}{The cat washes itself.}\glossnote{Constructions: reflexive; lavarsi le mani = to wash one's hands.}
\glossentry{\textbf{lavorare}\index{lavorare@lavorare}}{verb, -are, auxiliary \textit{avere}, participle \textit{lavorato}}{to work}{Il papà lavora in città.}{Dad works in town.}
\glossentry{\textbf{leggere}\index{leggere@leggere}}{verb, -ere, auxiliary \textit{avere}, participle \textit{letto}}{to read}{Leggo un libro.}{I read a book.}
\glossentry{\textbf{morire}\index{morire@morire}}{verb, -ire, auxiliary \textit{essere}, participle \textit{morto}; irregular: \textit{io muoio, tu muori, lui/lei muore, loro muoiono}}{to die}{L'albero è morto in inverno.}{The tree died in winter.}\glossnote{The stressed o becomes uo (muoio, muori, muore, muoiono). Futuro morirò, also morrò (less common).}
\glossentry{\textbf{nascere}\index{nascere@nascere}}{verb, -ere, auxiliary \textit{essere}, participle \textit{nato}}{to be born}{Sono nato a Roma.}{I was born in Rome.}
\glossentry{\textbf{pagare}\index{pagare@pagare}}{verb, -are, auxiliary \textit{avere}, participle \textit{pagato}}{to pay (for)}{Pago il caffè.}{I pay for the coffee.}\glossnote{Constructions: Takes a direct object: pagare il conto (no preposition for 'for').}
\glossentry{\textbf{pulire}\index{pulire@pulire}}{verb, -ire (-isc-), auxiliary \textit{avere}, participle \textit{pulito}}{to clean}{Pulisco la cucina.}{I clean the kitchen.}
\glossentry{\textbf{scrivere}\index{scrivere@scrivere}}{verb, -ere, auxiliary \textit{avere}, participle \textit{scritto}}{to write}{Scrivo una lettera.}{I write a letter.}
\glossentry{\textbf{sedersi}\index{sedersi@sedersi}}{verb, -ere, auxiliary \textit{essere}, participle \textit{seduto}; irregular: \textit{io siedo, tu siedi, lui/lei siede, loro siedono}}{to sit down}{Mi siedo sulla sedia.}{I sit down on the chair.}\glossnote{Constructions: reflexive; irregular stem sied-. The stressed e becomes ie (siedo, siedi, siede, siedono). Normally used reflexively: sedersi, mi siedo.}
\glossentry{\textbf{studiare}\index{studiare@studiare}}{verb, -are, auxiliary \textit{avere}, participle \textit{studiato}}{to study}{Studio l'italiano.}{I'm studying Italian.}
\glossentry{\textbf{svegliarsi}\index{svegliarsi@svegliarsi}}{verb, -are, auxiliary \textit{essere}, participle \textit{svegliato}}{to wake up}{Mi sveglio alle sette.}{I wake up at seven.}\glossnote{Constructions: reflexive.}
\glossentry{\textbf{vendere}\index{vendere@vendere}}{verb, -ere, auxiliary \textit{avere}, participle \textit{venduto}}{to sell}{Il pescatore vende il pesce.}{The fisherman sells the fish.}
\glossentry{\textbf{vestirsi}\index{vestirsi@vestirsi}}{verb, -ire, auxiliary \textit{essere}, participle \textit{vestito}}{to get dressed}{Mi vesto in fretta.}{I get dressed quickly.}\glossnote{Constructions: reflexive.}
\glossentry{\textbf{vivere}\index{vivere@vivere}}{verb, -ere, auxiliary \textit{avere}, participle \textit{vissuto}}{to live}{La volpe vive nel bosco.}{The fox lives in the wood.}\glossnote{Constructions: Compound tenses: both sono vissuto and ho vissuto are used.}
\end{glosslist}

\section{Descriptions}
\begin{glosslist}
\glossentry{\textbf{alto}\index{alto@alto}}{adjective: \textit{alto, alta, alti, alte}}{tall; high}{L'albero è alto.}{The tree is tall.}
\glossentry{\textbf{basso}\index{basso@basso}}{adjective: \textit{basso, bassa, bassi, basse}}{short; low}{Il tavolo è basso.}{The table is low.}
\glossentry{\textbf{bello}\index{bello@bello}}{adjective: \textit{bello, bella, belli, belle}}{beautiful; lovely; handsome}{Che bel fiore!}{What a beautiful flower!}\glossnote{Before a noun it changes like the definite article: bel, bello, bell', bella; bei, begli, belle.}
\glossentry{\textbf{brutto}\index{brutto@brutto}}{adjective: \textit{brutto, brutta, brutti, brutte}}{ugly; bad (weather)}{Oggi fa brutto tempo.}{The weather is bad today.}
\glossentry{\textbf{buio}\index{buio@buio}}{adjective: \textit{buio, buia, bui, buie}}{dark}{Il bosco è buio.}{The wood is dark.}\glossnote{Also a noun: il buio = the dark.}
\glossentry{\textbf{buono}\index{buono@buono}}{adjective: \textit{buono, buona, buoni, buone}}{good; tasty; kind}{La torta è buona.}{The cake is good.}\glossnote{Before a singular noun: buon (buon giorno), buono (buono zucchero), buona, buon' (buon'amica).}
\glossentry{\textbf{caldo}\index{caldo@caldo}}{adjective: \textit{caldo, calda, caldi, calde}}{hot; warm}{Il latte è caldo.}{The milk is hot.}\glossnote{Also a noun: fa caldo = it is hot (weather).}
\glossentry{\textbf{calmo}\index{calmo@calmo}}{adjective: \textit{calmo, calma, calmi, calme}}{calm}{Il mare è calmo.}{The sea is calm.}
\glossentry{\textbf{caro}\index{caro@caro}}{adjective: \textit{caro, cara, cari, care}}{dear; expensive}{Cara nonna, ti voglio bene.}{Dear Grandma, I love you.}\glossnote{Before a noun: dear (caro amico); after a noun or with essere often: expensive.}
\glossentry{\textbf{cattivo}\index{cattivo@cattivo}}{adjective: \textit{cattivo, cattiva, cattivi, cattive}}{bad; naughty; nasty}{Il lupo è cattivo.}{The wolf is bad.}
\glossentry{\textbf{corto}\index{corto@corto}}{adjective: \textit{corto, corta, corti, corte}}{short (length)}{La storia è corta.}{The story is short.}
\glossentry{\textbf{la cosa}\index{cosa@cosa}}{feminine; plural \textit{le cose}}{thing}{Che cosa mangi?}{What are you eating?}
\glossentry{\textbf{dolce}\index{dolce@dolce}}{adjective: \textit{dolce, dolce, dolci, dolci}}{sweet; gentle}{Il miele è dolce.}{Honey is sweet.}\glossnote{Also a noun: il dolce = dessert.}
\glossentry{\textbf{forte}\index{forte@forte}}{adjective: \textit{forte, forte, forti, forti}}{strong; loud}{Il vento è forte.}{The wind is strong.}
\glossentry{\textbf{freddo}\index{freddo@freddo}}{adjective: \textit{freddo, fredda, freddi, fredde}}{cold}{L'acqua è fredda.}{The water is cold.}\glossnote{Also a noun: fa freddo = it is cold (weather).}
\glossentry{\textbf{la gara}\index{gara@gara}}{feminine; plural \textit{le gare}}{race; competition}{La lepre e la tartaruga fanno una gara.}{The hare and the tortoise have a race.}
\glossentry{\textbf{gentile}\index{gentile@gentile}}{adjective: \textit{gentile, gentile, gentili, gentili}}{kind; polite}{La maestra è gentile.}{The teacher is kind.}
\glossentry{\textbf{giovane}\index{giovane@giovane}}{adjective: \textit{giovane, giovane, giovani, giovani}}{young}{La regina è giovane.}{The queen is young.}
\glossentry{\textbf{grande}\index{grande@grande}}{adjective: \textit{grande, grande, grandi, grandi}}{big; large; great}{La casa è grande.}{The house is big.}\glossnote{Before a noun it can mean 'great': un grande re.}
\glossentry{\textbf{l'idea}\index{idea@idea}}{feminine; plural \textit{le idee}}{idea}{Ho un'idea!}{I have an idea!}
\glossentry{\textbf{importante}\index{importante@importante}}{adjective: \textit{importante, importante, importanti, importanti}}{important}{È importante dormire bene.}{It is important to sleep well.}\glossnote{È importante che + congiuntivo: è importante che tu venga.}
\glossentry{\textbf{italiano}\index{italiano@italiano}}{adjective: \textit{italiano, italiana, italiani, italiane}}{Italian}{La pizza è italiana.}{Pizza is Italian.}\glossnote{Adjectives of nationality are written with a small letter and normally follow the noun: una canzone italiana.}
\glossentry{\textbf{il lavoro}\index{lavoro@lavoro}}{masculine; plural \textit{i lavori}}{work; job}{Il papà va al lavoro.}{Dad goes to work.}
\glossentry{\textbf{lento}\index{lento@lento}}{adjective: \textit{lento, lenta, lenti, lente}}{slow}{La tartaruga è lenta.}{The tortoise is slow.}
\glossentry{\textbf{lontano}\index{lontano@lontano}}{adjective: \textit{lontano, lontana, lontani, lontane}}{far; distant}{La luna è lontana.}{The moon is far away.}\glossnote{Also an adverb: abito lontano = I live far away; lontano da = far from.}
\glossentry{\textbf{lungo}\index{lungo@lungo}}{adjective: \textit{lungo, lunga, lunghi, lunghe}}{long}{Il fiume è lungo.}{The river is long.}
\glossentry{\textbf{nuovo}\index{nuovo@nuovo}}{adjective: \textit{nuovo, nuova, nuovi, nuove}}{new}{Ho uno zaino nuovo.}{I have a new backpack.}
\glossentry{\textbf{pesante}\index{pesante@pesante}}{adjective: \textit{pesante, pesante, pesanti, pesanti}}{heavy}{Lo zaino è pesante.}{The backpack is heavy.}
\glossentry{\textbf{piccolo}\index{piccolo@piccolo}}{adjective: \textit{piccolo, piccola, piccoli, piccole}}{small; little}{Il gatto è piccolo.}{The cat is small.}
\glossentry{\textbf{pieno}\index{pieno@pieno}}{adjective: \textit{pieno, piena, pieni, piene}}{full}{Il cesto è pieno di mele.}{The basket is full of apples.}
\glossentry{\textbf{primo}\index{primo@primo}}{adjective: \textit{primo, prima, primi, prime}}{first}{È il primo giorno di scuola.}{It is the first day of school.}\glossnote{Ordinal number; agrees like an -o adjective.}
\glossentry{\textbf{il problema}\index{problema@problema}}{masculine; plural \textit{i problemi}}{problem}{Non c'è nessun problema.}{There is no problem.}
\glossentry{\textbf{il rumore}\index{rumore@rumore}}{masculine; plural \textit{i rumori}}{noise}{Che cos'è questo rumore?}{What is this noise?}
\glossentry{\textbf{sembrare}\index{sembrare@sembrare}}{verb, -are, auxiliary \textit{essere}, participle \textit{sembrato}}{to seem}{Sembra stanco.}{He seems tired.}\glossnote{Constructions: sembrare + adjective; mi sembra che + congiuntivo.}
\glossentry{\textbf{il silenzio}\index{silenzio@silenzio}}{masculine; plural \textit{i silenzi}}{silence}{Nel bosco c'è silenzio.}{There is silence in the wood.}
\glossentry{\textbf{simpatico}\index{simpatico@simpatico}}{adjective: \textit{simpatico, simpatica, simpatici, simpatiche}}{nice; likeable}{Luca è simpatico.}{Luca is nice.}\glossnote{Stress on the third-to-last syllable (sim-PA-ti-co), so the masculine plural is simpatici, not *simpatichi.}
\glossentry{\textbf{tutto}\index{tutto@tutto}}{adjective: \textit{tutto, tutta, tutti, tutte}}{all; whole}{Tutti i bambini cantano.}{All the children sing.}\glossnote{Usually followed by the definite article: tutto il giorno, tutte le stelle.}
\glossentry{\textbf{ultimo}\index{ultimo@ultimo}}{adjective: \textit{ultimo, ultima, ultimi, ultime}}{last}{È l'ultima pagina.}{It is the last page.}
\glossentry{\textbf{vecchio}\index{vecchio@vecchio}}{adjective: \textit{vecchio, vecchia, vecchi, vecchie}}{old}{Il nonno ha un cane vecchio.}{Grandfather has an old dog.}
\glossentry{\textbf{veloce}\index{veloce@veloce}}{adjective: \textit{veloce, veloce, veloci, veloci}}{fast; quick}{La lepre è veloce.}{The hare is fast.}
\glossentry{\textbf{vicino}\index{vicino@vicino}}{adjective: \textit{vicino, vicina, vicini, vicine}}{near; close}{La scuola è vicina.}{The school is close.}\glossnote{As an adjective it agrees (la casa è vicina); vicino a … can also be an adverbial phrase: abito vicino al mare.}
\glossentry{\textbf{vuoto}\index{vuoto@vuoto}}{adjective: \textit{vuoto, vuota, vuoti, vuote}}{empty}{Il piatto è vuoto.}{The plate is empty.}
\end{glosslist}

\section{Essential}
\begin{glosslist}
\glossentry{\textbf{avere}\index{avere@avere}}{verb, -ere, auxiliary \textit{avere}, participle \textit{avuto}; irregular: \textit{io ho, tu hai, lui/lei ha, noi abbiamo, loro hanno}}{to have}{Ho un cane.}{I have a dog.}\glossnote{Constructions: avere fame/sete/sonno/paura = to be hungry/thirsty/sleepy/afraid; avere … anni = to be … years old. The h in ho, hai, ha, hanno is silent; it distinguishes ho (I have) from o (or), ha (has) from a (to), hanno (they have) from anno (year). Noi and voi are regular.}
\glossentry{\textbf{dare}\index{dare@dare}}{verb, -are, auxiliary \textit{avere}, participle \textit{dato}; irregular: \textit{tu dai, lui/lei dà, loro danno}}{to give}{La mucca dà il latte.}{The cow gives milk.}\glossnote{Constructions: dare qualcosa a qualcuno. dà (gives) carries an accent to distinguish it from the preposition da (from).}
\glossentry{\textbf{essere}\index{essere@essere}}{verb, -ere, auxiliary \textit{essere}, participle \textit{stato}; irregular: \textit{io sono, tu sei, lui/lei è, noi siamo, voi siete, loro sono}}{to be}{Il gatto è nero.}{The cat is black.}\glossnote{Constructions: essere + adjective; c'è / ci sono = there is / there are. Completely irregular. The participle stato is shared with stare: sono stato means both 'I have been' and 'I stayed'.}
\glossentry{\textbf{fare}\index{fare@fare}}{verb, -are, auxiliary \textit{avere}, participle \textit{fatto}; irregular: \textit{io faccio, tu fai, noi facciamo, loro fanno}}{to do; to make}{La nonna fa una torta.}{Grandmother makes a cake.}\glossnote{Constructions: fare colazione = to have breakfast; fa caldo/freddo = it is hot/cold. Historically from Latin facere: the stem fac- reappears in faccio, facciamo and the whole imperfetto (facevo).}
\glossentry{\textbf{stare}\index{stare@stare}}{verb, -are, auxiliary \textit{essere}, participle \textit{stato}; irregular: \textit{tu stai, loro stanno}}{to stay; to be (of health)}{Come stai? Sto bene.}{How are you? I'm well.}\glossnote{Constructions: stare bene/male = to be well/unwell; stare a casa = to stay at home. Short forms like those of dare and fare. Used for health (sto bene), staying, and in stare + gerund.}
\end{glosslist}

\section{Expressions}
\begin{glosslist}
\glossentry{\textbf{buonanotte}\index{buonanotte@buonanotte}}{interjection}{good night}{Buonanotte, piccolo orso.}{Good night, little bear.}
\glossentry{\textbf{buongiorno}\index{buongiorno@buongiorno}}{interjection}{good morning; good day}{Buongiorno, signora!}{Good morning, madam!}
\glossentry{\textbf{c'era una volta}\index{c'era una volta@c'era una volta}}{expression}{once upon a time (there was)}{C'era una volta un piccolo orso.}{Once upon a time there was a little bear.}\glossnote{Traditional opening of fairy tales: c' (= ci) + era (imperfetto of essere).}
\glossentry{\textbf{ciao}\index{ciao@ciao}}{interjection}{hello; hi; bye}{Ciao, Marco!}{Hi, Marco!}\glossnote{Informal greeting for both arriving and leaving.}
\glossentry{\textbf{ecco}\index{ecco@ecco}}{adverb}{here is; here are; there is}{Ecco il gatto!}{Here's the cat!}\glossnote{Presents something: ecco la mamma = here's Mum.}
\glossentry{\textbf{il favore}\index{favore@favore}}{masculine; plural \textit{i favori}}{favour; favor}{Mi fai un favore?}{Will you do me a favour?}\glossnote{per favore = please (literally: as a favour).}
\glossentry{\textbf{grazie}\index{grazie@grazie}}{interjection}{thank you}{Grazie per il regalo!}{Thank you for the present!}
\glossentry{\textbf{no}\index{no@no}}{adverb}{no}{No, grazie.}{No, thank you.}
\glossentry{\textbf{per favore}\index{per favore@per favore}}{expression}{please}{Un bicchiere d'acqua, per favore.}{A glass of water, please.}
\glossentry{\textbf{prego}\index{prego@prego}}{interjection}{you're welcome; please (go ahead)}{Grazie! — Prego.}{Thank you! — You're welcome.}
\glossentry{\textbf{scusa}\index{scusa@scusa}}{interjection}{sorry; excuse me (informal)}{Scusa, dov'è la scuola?}{Excuse me, where is the school?}\glossnote{Formal: scusi.}
\glossentry{\textbf{sì}\index{si@sì}}{adverb}{yes}{Sì, ho fame.}{Yes, I'm hungry.}\glossnote{Written with an accent to distinguish it from the pronoun si.}
\end{glosslist}

\section{Family}
\begin{glosslist}
\glossentry{\textbf{il cugino}\index{cugino@cugino}}{masculine; plural \textit{i cugini}}{cousin (male)}{Mio cugino gioca a calcio.}{My cousin plays football.}
\glossentry{\textbf{la famiglia}\index{famiglia@famiglia}}{feminine; plural \textit{le famiglie}}{family}{La mia famiglia è grande.}{My family is big.}
\glossentry{\textbf{la figlia}\index{figlia@figlia}}{feminine; plural \textit{le figlie}}{daughter}{La figlia somiglia alla madre.}{The daughter looks like her mother.}
\glossentry{\textbf{il figlio}\index{figlio@figlio}}{masculine; plural \textit{i figli}}{son}{Il figlio aiuta il padre.}{The son helps his father.}
\glossentry{\textbf{il fratello}\index{fratello@fratello}}{masculine; plural \textit{i fratelli}}{brother}{Mio fratello ha dieci anni.}{My brother is ten years old.}
\glossentry{\textbf{la madre}\index{madre@madre}}{feminine; plural \textit{le madri}}{mother}{Mia madre canta in cucina.}{My mother sings in the kitchen.}
\glossentry{\textbf{la mamma}\index{mamma@mamma}}{feminine; plural \textit{le mamme}}{mum; mom}{La mamma prepara la cena.}{Mum makes dinner.}
\glossentry{\textbf{il nipote}\index{nipote@nipote}}{masculine; plural \textit{i nipoti}}{grandchild; nephew/niece}{La nonna gioca con i nipoti.}{Grandma plays with her grandchildren.}\glossnote{Common gender: il nipote (male), la nipote (female); plural i nipoti / le nipoti. It means both grandchild and nephew/niece; context decides.}
\glossentry{\textbf{la nonna}\index{nonna@nonna}}{feminine; plural \textit{le nonne}}{grandmother}{La nonna fa una torta.}{Grandmother makes a cake.}
\glossentry{\textbf{il nonno}\index{nonno@nonno}}{masculine; plural \textit{i nonni}}{grandfather}{Il nonno racconta una storia.}{Grandfather tells a story.}
\glossentry{\textbf{il padre}\index{padre@padre}}{masculine; plural \textit{i padri}}{father}{Mio padre legge un libro.}{My father reads a book.}
\glossentry{\textbf{il papà}\index{papa@papà}}{masculine; plural \textit{i papà}}{dad}{Il papà gioca con i bambini.}{Dad plays with the children.}
\glossentry{\textbf{la sorella}\index{sorella@sorella}}{feminine; plural \textit{le sorelle}}{sister}{Mia sorella ama i gatti.}{My sister loves cats.}
\glossentry{\textbf{la zia}\index{zia@zia}}{feminine; plural \textit{le zie}}{aunt}{La zia porta un regalo.}{My aunt brings a present.}
\glossentry{\textbf{lo zio}\index{zio@zio}}{masculine; plural \textit{gli zii}}{uncle}{Lo zio abita al mare.}{My uncle lives by the sea.}
\end{glosslist}

\section{Feelings}
\begin{glosslist}
\glossentry{\textbf{l'abbraccio}\index{abbraccio@abbraccio}}{masculine; plural \textit{gli abbracci}}{hug}{Un abbraccio forte!}{A big hug!}
\glossentry{\textbf{l'allegria}\index{allegria@allegria}}{feminine; plural \textit{le allegrie}}{cheerfulness; fun}{Che allegria!}{What fun!}
\glossentry{\textbf{allegro}\index{allegro@allegro}}{adjective: \textit{allegro, allegra, allegri, allegre}}{cheerful; merry}{La canzone è allegra.}{The song is cheerful.}
\glossentry{\textbf{amare}\index{amare@amare}}{verb, -are, auxiliary \textit{avere}, participle \textit{amato}}{to love}{La mamma ama i suoi bambini.}{Mum loves her children.}
\glossentry{\textbf{l'amore}\index{amore@amore}}{masculine; plural \textit{gli amori}}{love}{L'amore della mamma è grande.}{A mother's love is great.}
\glossentry{\textbf{contento}\index{contento@contento}}{adjective: \textit{contento, contenta, contenti, contente}}{glad; pleased}{Sono contento di vederti.}{I am glad to see you.}
\glossentry{\textbf{il cuore}\index{cuore@cuore}}{masculine; plural \textit{i cuori}}{heart}{Il mio cuore è felice.}{My heart is happy.}
\glossentry{\textbf{divertirsi}\index{divertirsi@divertirsi}}{verb, -ire, auxiliary \textit{essere}, participle \textit{divertito}}{to have fun; to enjoy oneself}{Ci divertiamo al mare.}{We have fun at the seaside.}\glossnote{Constructions: reflexive.}
\glossentry{\textbf{felice}\index{felice@felice}}{adjective: \textit{felice, felice, felici, felici}}{happy}{La bambina è felice.}{The little girl is happy.}
\glossentry{\textbf{la felicità}\index{felicita@felicità}}{feminine; plural \textit{le felicità}}{happiness}{La felicità è semplice.}{Happiness is simple.}
\glossentry{\textbf{la gioia}\index{gioia@gioia}}{feminine; plural \textit{le gioie}}{joy}{Che gioia!}{What joy!}
\glossentry{\textbf{mancare}\index{mancare@mancare}}{verb, -are, auxiliary \textit{essere}, participle \textit{mancato}}{to be missing; to miss}{Mi manchi.}{I miss you.}\glossnote{Constructions: The person missed is the subject: mi manchi = I miss you (literally: you are lacking to me).}
\glossentry{\textbf{la paura}\index{paura@paura}}{feminine; plural \textit{le paure}}{fear}{Il topo ha paura.}{The mouse is afraid.}
\glossentry{\textbf{piacere}\index{piacere@piacere}}{verb, -ere, auxiliary \textit{essere}, participle \textit{piaciuto}; irregular: \textit{io piaccio, noi piacciamo, loro piacciono}}{to please; to be liked}{Mi piace il gelato.}{I like ice cream.}\glossnote{Constructions: mi piace + singular / mi piacciono + plural. The thing liked is the grammatical subject: mi piace il gelato (ice cream pleases me), mi piacciono le mele (apples please me). The cc appears in io, noi and loro.}
\glossentry{\textbf{piangere}\index{piangere@piangere}}{verb, -ere, auxiliary \textit{avere}, participle \textit{pianto}}{to cry}{Il bambino piange.}{The child is crying.}
\glossentry{\textbf{preferire}\index{preferire@preferire}}{verb, -ire (-isc-), auxiliary \textit{avere}, participle \textit{preferito}}{to prefer}{Preferisco il tè.}{I prefer tea.}
\glossentry{\textbf{ridere}\index{ridere@ridere}}{verb, -ere, auxiliary \textit{avere}, participle \textit{riso}}{to laugh}{I bambini ridono.}{The children laugh.}
\glossentry{\textbf{soffrire}\index{soffrire@soffrire}}{verb, -ire, auxiliary \textit{avere}, participle \textit{sofferto}}{to suffer}{Il cane soffre il caldo.}{The dog suffers in the heat.}
\glossentry{\textbf{sognare}\index{sognare@sognare}}{verb, -are, auxiliary \textit{avere}, participle \textit{sognato}}{to dream}{Il bambino sogna il mare.}{The child dreams of the sea.}
\glossentry{\textbf{il sogno}\index{sogno@sogno}}{masculine; plural \textit{i sogni}}{dream}{Ho fatto un bel sogno.}{I had a lovely dream.}
\glossentry{\textbf{il sonno}\index{sonno@sonno}}{masculine; no plural}{sleep; sleepiness}{L'orso ha sonno.}{The bear is sleepy.}
\glossentry{\textbf{sorridere}\index{sorridere@sorridere}}{verb, -ere, auxiliary \textit{avere}, participle \textit{sorriso}}{to smile}{La nonna sorride.}{Grandmother smiles.}\glossnote{Constructions: sorridere a qualcuno.}
\glossentry{\textbf{sperare}\index{sperare@sperare}}{verb, -are, auxiliary \textit{avere}, participle \textit{sperato}}{to hope}{Spero che tu stia bene.}{I hope you are well.}\glossnote{Constructions: sperare di + infinitive (same subject); sperare che + congiuntivo (also futuro: spero che verrai).}
\glossentry{\textbf{stanco}\index{stanco@stanco}}{adjective: \textit{stanco, stanca, stanchi, stanche}}{tired}{L'orso è stanco.}{The bear is tired.}
\glossentry{\textbf{temere}\index{temere@temere}}{verb, -ere, auxiliary \textit{avere}, participle \textit{temuto}}{to fear; to be afraid}{Temo che piova.}{I'm afraid it will rain.}\glossnote{Constructions: temere che + congiuntivo; temere di + infinitive.}
\glossentry{\textbf{triste}\index{triste@triste}}{adjective: \textit{triste, triste, tristi, tristi}}{sad}{Il cane è triste.}{The dog is sad.}
\glossentry{\textbf{la tristezza}\index{tristezza@tristezza}}{feminine; plural \textit{le tristezze}}{sadness}{C'è tristezza nella sua voce.}{There is sadness in his voice.}
\end{glosslist}

\section{Food}
\begin{glosslist}
\glossentry{\textbf{l'acqua}\index{acqua@acqua}}{feminine; plural \textit{le acque}}{water}{L'acqua è fredda.}{The water is cold.}
\glossentry{\textbf{l'arancia}\index{arancia@arancia}}{feminine; plural \textit{le arance}}{orange (fruit)}{L'arancia è dolce.}{The orange is sweet.}
\glossentry{\textbf{bere}\index{bere@bere}}{verb, -ere, auxiliary \textit{avere}, participle \textit{bevuto}; irregular: \textit{io bevo, tu bevi, lui/lei beve, noi beviamo, voi bevete, loro bevono}}{to drink}{Bevo un bicchiere di latte.}{I drink a glass of milk.}\glossnote{Short infinitive of an older bevere: once you know the stem bev-, the forms are regular -ere forms.}
\glossentry{\textbf{il biscotto}\index{biscotto@biscotto}}{masculine; plural \textit{i biscotti}}{biscuit; cookie}{Il bambino mangia un biscotto.}{The child eats a biscuit.}
\glossentry{\textbf{il caffè}\index{caffe@caffè}}{masculine; plural \textit{i caffè}}{coffee}{Il papà beve il caffè.}{Dad drinks coffee.}
\glossentry{\textbf{la carne}\index{carne@carne}}{feminine; plural \textit{le carni}}{meat}{Il lupo mangia la carne.}{The wolf eats meat.}
\glossentry{\textbf{la carota}\index{carota@carota}}{feminine; plural \textit{le carote}}{carrot}{Il coniglio mangia le carote.}{The rabbit eats carrots.}
\glossentry{\textbf{la cena}\index{cena@cena}}{feminine; plural \textit{le cene}}{dinner}{La cena è pronta.}{Dinner is ready.}
\glossentry{\textbf{la colazione}\index{colazione@colazione}}{feminine; plural \textit{le colazioni}}{breakfast}{Facciamo colazione insieme.}{We have breakfast together.}
\glossentry{\textbf{cucinare}\index{cucinare@cucinare}}{verb, -are, auxiliary \textit{avere}, participle \textit{cucinato}}{to cook}{La nonna cucina la minestra.}{Grandmother cooks the soup.}
\glossentry{\textbf{il formaggio}\index{formaggio@formaggio}}{masculine; plural \textit{i formaggi}}{cheese}{Il topo mangia il formaggio.}{The mouse eats the cheese.}
\glossentry{\textbf{la fragola}\index{fragola@fragola}}{feminine; plural \textit{le fragole}}{strawberry}{Le fragole sono rosse.}{Strawberries are red.}
\glossentry{\textbf{la frutta}\index{frutta@frutta}}{feminine; no plural}{fruit (collective)}{Mangio la frutta ogni giorno.}{I eat fruit every day.}
\glossentry{\textbf{il gelato}\index{gelato@gelato}}{masculine; plural \textit{i gelati}}{ice cream}{Il gelato è freddo.}{Ice cream is cold.}
\glossentry{\textbf{l'insalata}\index{insalata@insalata}}{feminine; plural \textit{le insalate}}{salad}{L'insalata è verde.}{The salad is green.}
\glossentry{\textbf{il latte}\index{latte@latte}}{masculine; no plural}{milk}{Il gatto beve il latte.}{The cat drinks the milk.}
\glossentry{\textbf{mangiare}\index{mangiare@mangiare}}{verb, -are, auxiliary \textit{avere}, participle \textit{mangiato}}{to eat}{Mangiamo una pizza.}{We eat a pizza.}
\glossentry{\textbf{la marmellata}\index{marmellata@marmellata}}{feminine; plural \textit{le marmellate}}{jam}{Il pane con la marmellata è buono.}{Bread with jam is good.}
\glossentry{\textbf{la mela}\index{mela@mela}}{feminine; plural \textit{le mele}}{apple}{La mela è rossa.}{The apple is red.}
\glossentry{\textbf{il miele}\index{miele@miele}}{masculine; no plural}{honey}{L'orso ama il miele.}{The bear loves honey.}
\glossentry{\textbf{la minestra}\index{minestra@minestra}}{feminine; plural \textit{le minestre}}{soup}{La nonna prepara la minestra.}{Grandmother makes soup.}
\glossentry{\textbf{il pane}\index{pane@pane}}{masculine; plural \textit{i pani}}{bread}{Il pane è caldo.}{The bread is warm.}
\glossentry{\textbf{la pasta}\index{pasta@pasta}}{feminine; plural \textit{le paste}}{pasta; pastry}{Mangiamo la pasta a pranzo.}{We eat pasta at lunch.}
\glossentry{\textbf{la pera}\index{pera@pera}}{feminine; plural \textit{le pere}}{pear}{La pera è dolce.}{The pear is sweet.}
\glossentry{\textbf{la pizza}\index{pizza@pizza}}{feminine; plural \textit{le pizze}}{pizza}{La pizza è calda.}{The pizza is hot.}
\glossentry{\textbf{il pomodoro}\index{pomodoro@pomodoro}}{masculine; plural \textit{i pomodori}}{tomato}{Il pomodoro è rosso.}{The tomato is red.}
\glossentry{\textbf{il pranzo}\index{pranzo@pranzo}}{masculine; plural \textit{i pranzi}}{lunch}{Il pranzo è alle due.}{Lunch is at two.}
\glossentry{\textbf{preparare}\index{preparare@preparare}}{verb, -are, auxiliary \textit{avere}, participle \textit{preparato}}{to prepare; to make}{La mamma prepara la colazione.}{Mum makes breakfast.}
\glossentry{\textbf{il sale}\index{sale@sale}}{masculine; no plural}{salt}{La nonna mette il sale nella minestra.}{Grandmother puts salt in the soup.}
\glossentry{\textbf{il succo}\index{succo@succo}}{masculine; plural \textit{i succhi}}{juice}{Bevo un succo di frutta.}{I drink a fruit juice.}
\glossentry{\textbf{il tè}\index{te@tè}}{masculine; plural \textit{i tè}}{tea}{La nonna beve il tè.}{Grandmother drinks tea.}
\glossentry{\textbf{la torta}\index{torta@torta}}{feminine; plural \textit{le torte}}{cake}{La torta è per la festa.}{The cake is for the party.}
\glossentry{\textbf{l'uovo}\index{uovo@uovo}}{masculine; plural \textit{le uova} (feminine)}{egg}{Mangio un uovo a colazione.}{I eat an egg at breakfast.}
\glossentry{\textbf{l'uva}\index{uva@uva}}{feminine; no plural}{grapes}{L'uva è buona.}{The grapes are good.}
\glossentry{\textbf{lo zucchero}\index{zucchero@zucchero}}{masculine; no plural}{sugar}{Lo zucchero è dolce.}{Sugar is sweet.}
\end{glosslist}

\section{Home}
\begin{glosslist}
\glossentry{\textbf{la barca}\index{barca@barca}}{feminine; plural \textit{le barche}}{boat}{La barca è sul mare.}{The boat is on the sea.}
\glossentry{\textbf{il bicchiere}\index{bicchiere@bicchiere}}{masculine; plural \textit{i bicchieri}}{(drinking) glass}{Il bicchiere è pieno d'acqua.}{The glass is full of water.}
\glossentry{\textbf{la borsa}\index{borsa@borsa}}{feminine; plural \textit{le borse}}{bag}{La borsa è sulla sedia.}{The bag is on the chair.}
\glossentry{\textbf{la camera}\index{camera@camera}}{feminine; plural \textit{le camere}}{room; bedroom}{La mia camera è blu.}{My bedroom is blue.}
\glossentry{\textbf{il cappello}\index{cappello@cappello}}{masculine; plural \textit{i cappelli}}{hat}{Metti il cappello: c'è il sole.}{Put your hat on: it's sunny.}
\glossentry{\textbf{la casa}\index{casa@casa}}{feminine; plural \textit{le case}}{house; home}{La casa è piccola.}{The house is small.}
\glossentry{\textbf{il castello}\index{castello@castello}}{masculine; plural \textit{i castelli}}{castle}{I bambini fanno un castello di sabbia.}{The children make a sandcastle.}
\glossentry{\textbf{il cesto}\index{cesto@cesto}}{masculine; plural \textit{i cesti}}{basket}{Nel cesto ci sono le mele.}{There are apples in the basket.}
\glossentry{\textbf{la chiave}\index{chiave@chiave}}{feminine; plural \textit{le chiavi}}{key}{La chiave è nella borsa.}{The key is in the bag.}
\glossentry{\textbf{la città}\index{citta@città}}{feminine; plural \textit{le città}}{city; town}{La città è grande.}{The city is big.}
\glossentry{\textbf{la coperta}\index{coperta@coperta}}{feminine; plural \textit{le coperte}}{blanket}{L'orso ha una coperta calda.}{The bear has a warm blanket.}
\glossentry{\textbf{la crema}\index{crema@crema}}{feminine; plural \textit{le creme}}{cream}{Metti la crema per il sole.}{Put on the sun cream.}
\glossentry{\textbf{la cucina}\index{cucina@cucina}}{feminine; plural \textit{le cucine}}{kitchen}{La mamma è in cucina.}{Mum is in the kitchen.}
\glossentry{\textbf{la festa}\index{festa@festa}}{feminine; plural \textit{le feste}}{party; holiday}{Domani c'è una festa.}{Tomorrow there is a party.}
\glossentry{\textbf{la finestra}\index{finestra@finestra}}{feminine; plural \textit{le finestre}}{window}{Il gatto guarda dalla finestra.}{The cat looks out of the window.}
\glossentry{\textbf{il giardino}\index{giardino@giardino}}{masculine; plural \textit{i giardini}}{garden}{I bambini giocano in giardino.}{The children play in the garden.}
\glossentry{\textbf{il gioco}\index{gioco@gioco}}{masculine; plural \textit{i giochi}}{game; toy}{Il gioco è divertente.}{The game is fun.}
\glossentry{\textbf{la lampada}\index{lampada@lampada}}{feminine; plural \textit{le lampade}}{lamp}{La lampada è accesa.}{The lamp is on.}
\glossentry{\textbf{la lanterna}\index{lanterna@lanterna}}{feminine; plural \textit{le lanterne}}{lantern}{Il nonno prende la lanterna.}{Grandpa takes the lantern.}
\glossentry{\textbf{il letto}\index{letto@letto}}{masculine; plural \textit{i letti}}{bed}{Il bambino va a letto.}{The child goes to bed.}
\glossentry{\textbf{il libro}\index{libro@libro}}{masculine; plural \textit{i libri}}{book}{Leggo un libro.}{I am reading a book.}
\glossentry{\textbf{la luce}\index{luce@luce}}{feminine; plural \textit{le luci}}{light}{Accendi la luce, per favore.}{Switch on the light, please.}
\glossentry{\textbf{la macchina}\index{macchina@macchina}}{feminine; plural \textit{le macchine}}{car}{La macchina è rossa.}{The car is red.}
\glossentry{\textbf{l'ombrello}\index{ombrello@ombrello}}{masculine; plural \textit{gli ombrelli}}{umbrella}{Prendo l'ombrello perché piove.}{I take the umbrella because it is raining.}
\glossentry{\textbf{il paese}\index{paese@paese}}{masculine; plural \textit{i paesi}}{village; country}{Il nonno vive in un piccolo paese.}{Grandfather lives in a small village.}
\glossentry{\textbf{la palla}\index{palla@palla}}{feminine; plural \textit{le palle}}{ball}{Il cane prende la palla.}{The dog catches the ball.}
\glossentry{\textbf{il piatto}\index{piatto@piatto}}{masculine; plural \textit{i piatti}}{plate; dish}{Il piatto è sul tavolo.}{The plate is on the table.}
\glossentry{\textbf{la porta}\index{porta@porta}}{feminine; plural \textit{le porte}}{door}{La porta è aperta.}{The door is open.}
\glossentry{\textbf{il regalo}\index{regalo@regalo}}{masculine; plural \textit{i regali}}{present; gift}{Il regalo è per te.}{The present is for you.}
\glossentry{\textbf{la scuola}\index{scuola@scuola}}{feminine; plural \textit{le scuole}}{school}{Vado a scuola a piedi.}{I go to school on foot.}
\glossentry{\textbf{la sedia}\index{sedia@sedia}}{feminine; plural \textit{le sedie}}{chair}{La sedia è vicino alla finestra.}{The chair is near the window.}
\glossentry{\textbf{lo specchio}\index{specchio@specchio}}{masculine; plural \textit{gli specchi}}{mirror}{La regina guarda nello specchio.}{The queen looks in the mirror.}
\glossentry{\textbf{lo stivale}\index{stivale@stivale}}{masculine; plural \textit{gli stivali}}{boot}{Sara mette gli stivali gialli.}{Sara puts on her yellow boots.}
\glossentry{\textbf{la strada}\index{strada@strada}}{feminine; plural \textit{le strade}}{street; road}{La strada è bagnata.}{The street is wet.}
\glossentry{\textbf{il tavolo}\index{tavolo@tavolo}}{masculine; plural \textit{i tavoli}}{table}{Il libro è sul tavolo.}{The book is on the table.}
\glossentry{\textbf{la tazza}\index{tazza@tazza}}{feminine; plural \textit{le tazze}}{cup; mug}{La tazza di latte è calda.}{The cup of milk is hot.}
\glossentry{\textbf{il tetto}\index{tetto@tetto}}{masculine; plural \textit{i tetti}}{roof}{La pioggia cade sul tetto.}{The rain falls on the roof.}
\glossentry{\textbf{il treno}\index{treno@treno}}{masculine; plural \textit{i treni}}{train}{Il treno arriva alle tre.}{The train arrives at three.}
\glossentry{\textbf{lo zaino}\index{zaino@zaino}}{masculine; plural \textit{gli zaini}}{backpack; rucksack}{Lo zaino è pesante.}{The backpack is heavy.}
\end{glosslist}

\section{Modal}
\begin{glosslist}
\glossentry{\textbf{dovere}\index{dovere@dovere}}{verb, -ere, auxiliary \textit{avere}, participle \textit{dovuto}; irregular: \textit{io devo, tu devi, lui/lei deve, noi dobbiamo, loro devono}}{to have to; must}{Devo andare a scuola.}{I have to go to school.}\glossnote{Constructions: modal: dovere + infinitive. Modal verb: dev- in the stressed forms, dobb- in noi, dov- in voi. Alternative forms debbo, debbono exist but are rare today.}
\glossentry{\textbf{potere}\index{potere@potere}}{verb, -ere, auxiliary \textit{avere}, participle \textit{potuto}; irregular: \textit{io posso, tu puoi, lui/lei può, noi possiamo, loro possono}}{to be able to; can; may}{Posso entrare?}{May I come in?}\glossnote{Constructions: modal: potere + infinitive. Modal verb: usually followed by an infinitive (posso venire). In compound tenses a modal often takes the auxiliary of the following verb (sono potuto venire).}
\glossentry{\textbf{sapere}\index{sapere@sapere}}{verb, -ere, auxiliary \textit{avere}, participle \textit{saputo}; irregular: \textit{io so, tu sai, lui/lei sa, noi sappiamo, loro sanno}}{to know (a fact); to know how to}{So nuotare.}{I can swim.}\glossnote{Constructions: sapere + infinitive = to know how to. Means 'to know' (a fact) and, with an infinitive, 'to know how to': so nuotare = I can (know how to) swim.}
\glossentry{\textbf{volere}\index{volere@volere}}{verb, -ere, auxiliary \textit{avere}, participle \textit{voluto}; irregular: \textit{io voglio, tu vuoi, lui/lei vuole, noi vogliamo, loro vogliono}}{to want}{Voglio un gelato.}{I want an ice cream.}\glossnote{Constructions: volere + noun / + infinitive; polite: vorrei. Modal verb. Can also take a direct object: voglio un gelato. Polite request: vorrei (conditional).}
\end{glosslist}

\section{Movement}
\begin{glosslist}
\glossentry{\textbf{andare}\index{andare@andare}}{verb, -are, auxiliary \textit{essere}, participle \textit{andato}; irregular: \textit{io vado, tu vai, lui/lei va, loro vanno}}{to go}{Vado a scuola.}{I go to school.}\glossnote{Constructions: andare a + city / infinitive; andare in + country, region or many places (in giardino, in cucina). Two stems: vad-/va- in the stressed forms (io, tu, lui/lei, loro), regular and- in noi and voi.}
\glossentry{\textbf{arrivare}\index{arrivare@arrivare}}{verb, -are, auxiliary \textit{essere}, participle \textit{arrivato}}{to arrive}{Il treno arriva alle tre.}{The train arrives at three.}
\glossentry{\textbf{cadere}\index{cadere@cadere}}{verb, -ere, auxiliary \textit{essere}, participle \textit{caduto}}{to fall}{La foglia cade.}{The leaf falls.}
\glossentry{\textbf{camminare}\index{camminare@camminare}}{verb, -are, auxiliary \textit{avere}, participle \textit{camminato}}{to walk}{Camminiamo nel bosco.}{We walk in the wood.}
\glossentry{\textbf{correre}\index{correre@correre}}{verb, -ere, auxiliary \textit{avere}, participle \textit{corso}}{to run}{Il cane corre.}{The dog runs.}\glossnote{Constructions: Auxiliary avere (ho corso), but essere with a destination (sono corso a casa).}
\glossentry{\textbf{la direzione}\index{direzione@direzione}}{feminine; plural \textit{le direzioni}}{direction}{Il vento ha cambiato direzione.}{The wind has changed direction.}
\glossentry{\textbf{entrare}\index{entrare@entrare}}{verb, -are, auxiliary \textit{essere}, participle \textit{entrato}}{to enter; to go in}{Il gatto entra in casa.}{The cat comes into the house.}\glossnote{Constructions: entrare in.}
\glossentry{\textbf{fermarsi}\index{fermarsi@fermarsi}}{verb, -are, auxiliary \textit{essere}, participle \textit{fermato}}{to stop (oneself)}{Il treno si ferma qui.}{The train stops here.}\glossnote{Constructions: fermarsi a + infinitive = to stop to do something; fermare (non-reflexive) = to stop something or someone.}
\glossentry{\textbf{guidare}\index{guidare@guidare}}{verb, -are, auxiliary \textit{avere}, participle \textit{guidato}}{to drive; to guide}{Il papà guida la macchina.}{Dad drives the car.}
\glossentry{\textbf{la mappa}\index{mappa@mappa}}{feminine; plural \textit{le mappe}}{map}{Guardiamo la mappa della città.}{Let's look at the map of the city.}
\glossentry{\textbf{nuotare}\index{nuotare@nuotare}}{verb, -are, auxiliary \textit{avere}, participle \textit{nuotato}}{to swim}{Il pesce nuota.}{The fish swims.}
\glossentry{\textbf{partire}\index{partire@partire}}{verb, -ire, auxiliary \textit{essere}, participle \textit{partito}}{to leave; to set off}{Il treno parte alle otto.}{The train leaves at eight.}
\glossentry{\textbf{passare}\index{passare@passare}}{verb, -are, auxiliary \textit{essere}, participle \textit{passato}}{to pass; to spend (time); to drop by}{Passo da te stasera.}{I'll drop by your place this evening.}\glossnote{Constructions: Auxiliary avere when transitive: ho passato l'estate al mare.}
\glossentry{\textbf{restare}\index{restare@restare}}{verb, -are, auxiliary \textit{essere}, participle \textit{restato}}{to stay; to remain}{Resto a casa.}{I'm staying at home.}
\glossentry{\textbf{rimanere}\index{rimanere@rimanere}}{verb, -ere, auxiliary \textit{essere}, participle \textit{rimasto}; irregular: \textit{io rimango, loro rimangono}}{to remain; to stay}{Rimango a casa oggi.}{I am staying at home today.}\glossnote{A g appears in io and loro (rimango, rimangono). Irregular participle rimasto.}
\glossentry{\textbf{salire}\index{salire@salire}}{verb, -ire, auxiliary \textit{essere}, participle \textit{salito}; irregular: \textit{io salgo, loro salgono}}{to go up; to climb; to get on}{Lo scoiattolo sale sull'albero.}{The squirrel climbs up the tree.}\glossnote{Constructions: salire su = to climb onto / get on. A g appears in io and loro (salgo, salgono). Auxiliary essere when it means 'to go up' (sono salito sull'albero); avere with a direct object (ho salito le scale).}
\glossentry{\textbf{saltare}\index{saltare@saltare}}{verb, -are, auxiliary \textit{avere}, participle \textit{saltato}}{to jump}{La rana salta.}{The frog jumps.}
\glossentry{\textbf{scappare}\index{scappare@scappare}}{verb, -are, auxiliary \textit{essere}, participle \textit{scappato}}{to run away; to escape}{Il topo scappa.}{The mouse runs away.}
\glossentry{\textbf{scendere}\index{scendere@scendere}}{verb, -ere, auxiliary \textit{essere}, participle \textit{sceso}}{to go down; to get off}{Scendo dal treno.}{I get off the train.}\glossnote{Constructions: Auxiliary avere with a direct object: ho sceso le scale.}
\glossentry{\textbf{seguire}\index{seguire@seguire}}{verb, -ire, auxiliary \textit{avere}, participle \textit{seguito}}{to follow}{Il cane segue il bambino.}{The dog follows the child.}
\glossentry{\textbf{tornare}\index{tornare@tornare}}{verb, -are, auxiliary \textit{essere}, participle \textit{tornato}}{to return; to come back}{Torno a casa la sera.}{I come back home in the evening.}
\glossentry{\textbf{uscire}\index{uscire@uscire}}{verb, -ire, auxiliary \textit{essere}, participle \textit{uscito}; irregular: \textit{io esco, tu esci, lui/lei esce, loro escono}}{to go out; to leave}{Esco con gli amici.}{I'm going out with my friends.}\glossnote{Constructions: uscire di casa = to leave the house. The stressed forms begin with e- (esco, esci, esce, escono); noi and voi keep u- (usciamo, uscite).}
\glossentry{\textbf{venire}\index{venire@venire}}{verb, -ire, auxiliary \textit{essere}, participle \textit{venuto}; irregular: \textit{io vengo, tu vieni, lui/lei viene, loro vengono}}{to come}{Vieni al mare con noi?}{Are you coming to the seaside with us?}\glossnote{Constructions: venire da = to come from. A g appears in io and loro (vengo, vengono); the stressed e becomes ie in tu and lui/lei (vieni, viene). The participle is venuto, not *venito.}
\glossentry{\textbf{viaggiare}\index{viaggiare@viaggiare}}{verb, -are, auxiliary \textit{avere}, participle \textit{viaggiato}}{to travel}{Viaggiamo in treno.}{We travel by train.}
\glossentry{\textbf{visitare}\index{visitare@visitare}}{verb, -are, auxiliary \textit{avere}, participle \textit{visitato}}{to visit (a place)}{Visitiamo il museo.}{We visit the museum.}\glossnote{Constructions: Visiting a person: andare a trovare qualcuno.}
\glossentry{\textbf{volare}\index{volare@volare}}{verb, -are, auxiliary \textit{avere}, participle \textit{volato}}{to fly}{L'uccello vola alto.}{The bird flies high.}
\end{glosslist}

\section{Nature}
\begin{glosslist}
\glossentry{\textbf{l'albero}\index{albero@albero}}{masculine; plural \textit{gli alberi}}{tree}{L'albero è alto.}{The tree is tall.}
\glossentry{\textbf{il bosco}\index{bosco@bosco}}{masculine; plural \textit{i boschi}}{wood; forest}{Nel bosco ci sono molti alberi.}{There are many trees in the wood.}
\glossentry{\textbf{brillare}\index{brillare@brillare}}{verb, -are, auxiliary \textit{avere}, participle \textit{brillato}}{to shine; to sparkle}{Le stelle brillano.}{The stars are shining.}
\glossentry{\textbf{il cielo}\index{cielo@cielo}}{masculine; plural \textit{i cieli}}{sky}{Il cielo è azzurro.}{The sky is blue.}
\glossentry{\textbf{la cima}\index{cima@cima}}{feminine; plural \textit{le cime}}{top; summit}{Siamo arrivati in cima alla montagna.}{We reached the top of the mountain.}
\glossentry{\textbf{la collina}\index{collina@collina}}{feminine; plural \textit{le colline}}{hill}{La casa è sulla collina.}{The house is on the hill.}
\glossentry{\textbf{la conchiglia}\index{conchiglia@conchiglia}}{feminine; plural \textit{le conchiglie}}{shell; seashell}{Trovo una conchiglia sulla spiaggia.}{I find a shell on the beach.}
\glossentry{\textbf{crescere}\index{crescere@crescere}}{verb, -ere, auxiliary \textit{essere}, participle \textit{cresciuto}}{to grow}{L'albero cresce lentamente.}{The tree grows slowly.}\glossnote{Constructions: Auxiliary avere when transitive (= to raise): ha cresciuto tre figli.}
\glossentry{\textbf{l'erba}\index{erba@erba}}{feminine; plural \textit{le erbe}}{grass; herb}{L'erba è verde.}{The grass is green.}
\glossentry{\textbf{il faro}\index{faro@faro}}{masculine; plural \textit{i fari}}{lighthouse}{Il faro è sulla scogliera.}{The lighthouse is on the cliff.}
\glossentry{\textbf{il fiore}\index{fiore@fiore}}{masculine; plural \textit{i fiori}}{flower}{Il fiore è giallo.}{The flower is yellow.}
\glossentry{\textbf{il fiume}\index{fiume@fiume}}{masculine; plural \textit{i fiumi}}{river}{Il fiume è lungo.}{The river is long.}
\glossentry{\textbf{la foglia}\index{foglia@foglia}}{feminine; plural \textit{le foglie}}{leaf}{La foglia cade dall'albero.}{The leaf falls from the tree.}
\glossentry{\textbf{la foresta}\index{foresta@foresta}}{feminine; plural \textit{le foreste}}{forest}{La foresta è grande e buia.}{The forest is big and dark.}
\glossentry{\textbf{la grotta}\index{grotta@grotta}}{feminine; plural \textit{le grotte}}{cave}{L'orso dorme in una grotta.}{The bear sleeps in a cave.}
\glossentry{\textbf{l'isola}\index{isola@isola}}{feminine; plural \textit{le isole}}{island}{L'isola è piccola.}{The island is small.}
\glossentry{\textbf{il lago}\index{lago@lago}}{masculine; plural \textit{i laghi}}{lake}{Il lago è blu.}{The lake is blue.}
\glossentry{\textbf{la luna}\index{luna@luna}}{feminine; plural \textit{le lune}}{moon}{La luna è bianca.}{The moon is white.}
\glossentry{\textbf{il mare}\index{mare@mare}}{masculine; plural \textit{i mari}}{sea}{Il mare è calmo.}{The sea is calm.}
\glossentry{\textbf{la montagna}\index{montagna@montagna}}{feminine; plural \textit{le montagne}}{mountain}{La montagna è alta.}{The mountain is high.}
\glossentry{\textbf{la nave}\index{nave@nave}}{feminine; plural \textit{le navi}}{ship}{La nave entra nel porto.}{The ship enters the harbour.}
\glossentry{\textbf{il nido}\index{nido@nido}}{masculine; plural \textit{i nidi}}{nest}{L'uccello fa il nido.}{The bird builds its nest.}
\glossentry{\textbf{l'onda}\index{onda@onda}}{feminine; plural \textit{le onde}}{wave}{Le onde sono alte oggi.}{The waves are high today.}
\glossentry{\textbf{il prato}\index{prato@prato}}{masculine; plural \textit{i prati}}{meadow; lawn}{Le pecore sono nel prato.}{The sheep are in the meadow.}
\glossentry{\textbf{raccogliere}\index{raccogliere@raccogliere}}{verb, -ere, auxiliary \textit{avere}, participle \textit{raccolto}; irregular: \textit{io raccolgo, loro raccolgono}}{to pick (up); to gather}{Raccogliamo le fragole.}{We pick strawberries.}\glossnote{Same pattern as scegliere: raccolgo, raccolgono; participle raccolto.}
\glossentry{\textbf{il ramo}\index{ramo@ramo}}{masculine; plural \textit{i rami}}{branch}{Il gufo è sul ramo.}{The owl is on the branch.}
\glossentry{\textbf{la sabbia}\index{sabbia@sabbia}}{feminine; plural \textit{le sabbie}}{sand}{La sabbia è calda.}{The sand is hot.}
\glossentry{\textbf{il sasso}\index{sasso@sasso}}{masculine; plural \textit{i sassi}}{stone; pebble}{Il bambino lancia un sasso nel lago.}{The child throws a stone into the lake.}
\glossentry{\textbf{lo scoglio}\index{scoglio@scoglio}}{masculine; plural \textit{gli scogli}}{rock (in the sea)}{I bambini giocano sugli scogli.}{The children play on the rocks.}
\glossentry{\textbf{il sentiero}\index{sentiero@sentiero}}{masculine; plural \textit{i sentieri}}{path; trail}{Il sentiero entra nel bosco.}{The path goes into the wood.}
\glossentry{\textbf{il sole}\index{sole@sole}}{masculine; plural \textit{i soli}}{sun}{Il sole è caldo.}{The sun is hot.}
\glossentry{\textbf{la spiaggia}\index{spiaggia@spiaggia}}{feminine; plural \textit{le spiagge}}{beach}{Andiamo in spiaggia.}{Let's go to the beach.}
\glossentry{\textbf{la stella}\index{stella@stella}}{feminine; plural \textit{le stelle}}{star}{Ci sono molte stelle in cielo.}{There are many stars in the sky.}
\glossentry{\textbf{la tana}\index{tana@tana}}{feminine; plural \textit{le tane}}{den; burrow}{La volpe torna nella tana.}{The fox goes back into its den.}
\glossentry{\textbf{la terra}\index{terra@terra}}{feminine; plural \textit{le terre}}{earth; ground; land}{La mela cade a terra.}{The apple falls to the ground.}
\end{glosslist}

\section{Numbers}
\begin{glosslist}
\glossentry{\textbf{cinque}\index{cinque@cinque}}{numeral}{five}{Cinque stelle in cielo.}{Five stars in the sky.}
\glossentry{\textbf{dieci}\index{dieci@dieci}}{numeral}{ten}{Ho dieci biscotti.}{I have ten biscuits.}
\glossentry{\textbf{due}\index{due@due}}{numeral}{two}{Ho due gatti.}{I have two cats.}
\glossentry{\textbf{nove}\index{nove@nove}}{numeral}{nine}{Nove uccelli sul ramo.}{Nine birds on the branch.}
\glossentry{\textbf{otto}\index{otto@otto}}{numeral}{eight}{Il treno parte alle otto.}{The train leaves at eight.}
\glossentry{\textbf{quattro}\index{quattro@quattro}}{numeral}{four}{Il tavolo ha quattro gambe.}{The table has four legs.}
\glossentry{\textbf{sei}\index{sei@sei}}{numeral}{six}{Luca ha sei anni.}{Luca is six years old.}\glossnote{Also the verb form 'you are' (tu sei).}
\glossentry{\textbf{sette}\index{sette@sette}}{numeral}{seven}{Una settimana ha sette giorni.}{A week has seven days.}
\glossentry{\textbf{tre}\index{tre@tre}}{numeral}{three}{Tre orsi nel bosco.}{Three bears in the wood.}\glossnote{Ventitré, trentatré … carry an accent; tre alone does not.}
\glossentry{\textbf{uno}\index{uno@uno}}{numeral}{one}{Ho un fratello.}{I have one brother.}\glossnote{Before a noun it behaves like the indefinite article: un, uno, una, un'.}
\end{glosslist}

\section{People}
\begin{glosslist}
\glossentry{\textbf{l'amica}\index{amica@amica}}{feminine; plural \textit{le amiche}}{friend (female)}{Sara è la mia amica.}{Sara is my friend.}
\glossentry{\textbf{l'amico}\index{amico@amico}}{masculine; plural \textit{gli amici}}{friend (male)}{Marco è il mio amico.}{Marco is my friend.}
\glossentry{\textbf{la bambina}\index{bambina@bambina}}{feminine; plural \textit{le bambine}}{little girl}{La bambina mangia una mela.}{The little girl eats an apple.}
\glossentry{\textbf{il bambino}\index{bambino@bambino}}{masculine; plural \textit{i bambini}}{child; little boy}{Il bambino gioca con la palla.}{The little boy plays with the ball.}
\glossentry{\textbf{la donna}\index{donna@donna}}{feminine; plural \textit{le donne}}{woman}{La donna compra il pane.}{The woman buys bread.}
\glossentry{\textbf{la gente}\index{gente@gente}}{feminine; no plural}{people (collective, singular)}{La gente cammina sulla spiaggia.}{People are walking on the beach.}
\glossentry{\textbf{incontrare}\index{incontrare@incontrare}}{verb, -are, auxiliary \textit{avere}, participle \textit{incontrato}}{to meet}{Ho incontrato Sara al mercato.}{I met Sara at the market.}
\glossentry{\textbf{invitare}\index{invitare@invitare}}{verb, -are, auxiliary \textit{avere}, participle \textit{invitato}}{to invite}{Ti invito alla festa.}{I'm inviting you to the party.}
\glossentry{\textbf{la maestra}\index{maestra@maestra}}{feminine; plural \textit{le maestre}}{(primary-school) teacher (female)}{La maestra legge una storia.}{The teacher reads a story.}
\glossentry{\textbf{il nome}\index{nome@nome}}{masculine; plural \textit{i nomi}}{name}{Il mio nome è Sara.}{My name is Sara.}
\glossentry{\textbf{la persona}\index{persona@persona}}{feminine; plural \textit{le persone}}{person}{Ci sono tre persone in barca.}{There are three people in the boat.}
\glossentry{\textbf{il pescatore}\index{pescatore@pescatore}}{masculine; plural \textit{i pescatori}}{fisherman}{Il pescatore torna con la barca.}{The fisherman comes back with the boat.}
\glossentry{\textbf{la principessa}\index{principessa@principessa}}{feminine; plural \textit{le principesse}}{princess}{La principessa canta alla finestra.}{The princess sings at the window.}
\glossentry{\textbf{la ragazza}\index{ragazza@ragazza}}{feminine; plural \textit{le ragazze}}{girl; young woman}{La ragazza canta una canzone.}{The girl sings a song.}
\glossentry{\textbf{il ragazzo}\index{ragazzo@ragazzo}}{masculine; plural \textit{i ragazzi}}{boy; young man}{Il ragazzo legge un libro.}{The boy reads a book.}
\glossentry{\textbf{il re}\index{re@re}}{masculine; plural \textit{i re}}{king}{Il re vive in un castello.}{The king lives in a castle.}
\glossentry{\textbf{la regina}\index{regina@regina}}{feminine; plural \textit{le regine}}{queen}{La regina ha una corona.}{The queen has a crown.}
\glossentry{\textbf{salutare}\index{salutare@salutare}}{verb, -are, auxiliary \textit{avere}, participle \textit{salutato}}{to greet; to say goodbye}{Saluto i vicini.}{I greet the neighbours.}
\glossentry{\textbf{la signora}\index{signora@signora}}{feminine; plural \textit{le signore}}{lady; Mrs}{La signora ha un cane piccolo.}{The lady has a small dog.}
\glossentry{\textbf{l'uomo}\index{uomo@uomo}}{masculine; plural \textit{gli uomini}}{man}{L'uomo apre la porta.}{The man opens the door.}
\end{glosslist}

\section{Prepositions}
\begin{glosslist}
\glossentry{\textbf{a}\index{a@a}}{preposition}{to; at; in (cities)}{Vado a Roma.}{I'm going to Rome.}\glossnote{Before a vowel it may become ad: ad Anna. Combines with the article: al, allo, all', alla, ai, agli, alle.}
\glossentry{\textbf{con}\index{con@con}}{preposition}{with}{Gioco con il cane.}{I play with the dog.}\glossnote{Normally not contracted in modern writing: con il, con la.}
\glossentry{\textbf{da}\index{da@da}}{preposition}{from; at (someone's place); since}{Vengo da Milano.}{I come from Milan.}\glossnote{da + person = at someone's place: vado dal nonno. Never elides. Combines: dal, dallo, dall', dalla, dai, dagli, dalle.}
\glossentry{\textbf{di}\index{di@di}}{preposition}{of; from (origin with essere)}{Il libro di Marco.}{Marco's book.}\glossnote{Possession: il cane di Sara = Sara's dog. Elides before a vowel: d'estate. Combines: del, dello, dell', della, dei, degli, delle.}
\glossentry{\textbf{fra}\index{fra@fra}}{preposition}{between; among; in (time from now)}{Arrivo fra un'ora.}{I'll arrive in an hour.}\glossnote{Same meaning as tra.}
\glossentry{\textbf{in}\index{in@in}}{preposition}{in; into; to (countries, rooms)}{Il gatto è in cucina.}{The cat is in the kitchen.}\glossnote{Combines with the article: nel, nello, nell', nella, nei, negli, nelle.}
\glossentry{\textbf{per}\index{per@per}}{preposition}{for; through; in order to}{Questo regalo è per te.}{This present is for you.}\glossnote{per + infinitive = in order to: studio per imparare.}
\glossentry{\textbf{sopra}\index{sopra@sopra}}{preposition}{over; above; on top of}{La lampada è sopra il tavolo.}{The lamp is above the table.}
\glossentry{\textbf{sotto}\index{sotto@sotto}}{preposition}{under; below}{Il cane dorme sotto il tavolo.}{The dog sleeps under the table.}
\glossentry{\textbf{su}\index{su@su}}{preposition}{on; onto; about}{Il libro è sul tavolo.}{The book is on the table.}\glossnote{Combines with the article: sul, sullo, sull', sulla, sui, sugli, sulle.}
\glossentry{\textbf{tra}\index{tra@tra}}{preposition}{between; among; in (time from now)}{Il gatto dorme tra i fiori.}{The cat sleeps among the flowers.}\glossnote{tra and fra mean the same; choose the one that sounds better: tra fratelli, fra tre giorni.}
\glossentry{\textbf{verso}\index{verso@verso}}{preposition}{towards; around (time)}{Il cane corre verso il mare.}{The dog runs towards the sea.}\glossnote{verso + time = around: verso le otto. Before a personal pronoun: verso di me.}
\glossentry{\textbf{vicino}\index{vicino@vicino}}{adverb}{near; nearby}{La scuola è vicino a casa.}{The school is near home.}\glossnote{vicino a + noun = near (to)}
\end{glosslist}

\section{Pronouns}
\begin{glosslist}
\glossentry{\textbf{io}\index{io@io}}{pronoun}{I}{Io sono Marco.}{I am Marco.}\glossnote{Usually omitted, because the verb ending already shows the person; used for emphasis or contrast.}
\glossentry{\textbf{lei}\index{lei@lei}}{pronoun}{she; her; you (formal)}{Lei è mia sorella.}{She is my sister.}\glossnote{Lei (capitalised in writing) is also the formal 'you', used with third-person verb forms.}
\glossentry{\textbf{loro}\index{loro@loro}}{pronoun}{they; them; their}{Loro giocano in giardino.}{They play in the garden.}\glossnote{Also a possessive, preceded by the article and invariable: il loro cane, la loro casa.}
\glossentry{\textbf{lui}\index{lui@lui}}{pronoun}{he; him}{Lui è mio fratello.}{He is my brother.}
\glossentry{\textbf{noi}\index{noi@noi}}{pronoun}{we; us}{Noi andiamo al mare.}{We are going to the seaside.}
\glossentry{\textbf{tu}\index{tu@tu}}{pronoun}{you (singular, informal)}{Tu sei mio amico.}{You are my friend.}
\glossentry{\textbf{voi}\index{voi@voi}}{pronoun}{you (plural)}{Voi siete gentili.}{You are kind.}
\end{glosslist}

\section{Question-words}
\begin{glosslist}
\glossentry{\textbf{che}\index{che@che}}{pronoun}{that; which; who; what}{Il libro che leggo è bello.}{The book (that) I'm reading is good.}\glossnote{Relative pronoun (the book that…), conjunction (penso che…) and question word (che cosa?). Never omitted in Italian, unlike English 'that'.}
\glossentry{\textbf{chi}\index{chi@chi}}{pronoun}{who; whom}{Chi è?}{Who is it?}
\glossentry{\textbf{come}\index{come@come}}{adverb}{how; like; as}{Come ti chiami?}{What is your name?}\glossnote{Before è it elides: com'è?}
\glossentry{\textbf{cosa}\index{cosa@cosa}}{pronoun}{what (question)}{Cosa fai?}{What are you doing?}\glossnote{Che cosa?, cosa? and che? all mean 'what?'.}
\glossentry{\textbf{dove}\index{dove@dove}}{adverb}{where}{Dove abiti?}{Where do you live?}\glossnote{Before è it elides: dov'è?}
\glossentry{\textbf{quale}\index{quale@quale}}{adjective: \textit{quale, quale, quali, quali}}{which; what}{Quale libro preferisci?}{Which book do you prefer?}\glossnote{Plural quali. Before è: qual è (no apostrophe).}
\glossentry{\textbf{quanto}\index{quanto@quanto}}{adjective: \textit{quanto, quanta, quanti, quante}}{how much; how many}{Quanti anni hai?}{How old are you?}\glossnote{Agrees with the noun: quanto, quanta, quanti, quante.}
\end{glosslist}

\section{Songs}
\begin{glosslist}
\glossentry{\textbf{ascoltare}\index{ascoltare@ascoltare}}{verb, -are, auxiliary \textit{avere}, participle \textit{ascoltato}}{to listen to}{Ascoltiamo una canzone.}{We listen to a song.}\glossnote{Constructions: Takes a direct object: ascoltare la musica (no preposition).}
\glossentry{\textbf{ballare}\index{ballare@ballare}}{verb, -are, auxiliary \textit{avere}, participle \textit{ballato}}{to dance}{Balliamo alla festa.}{We dance at the party.}
\glossentry{\textbf{cantare}\index{cantare@cantare}}{verb, -are, auxiliary \textit{avere}, participle \textit{cantato}}{to sing}{La ragazza canta una canzone.}{The girl sings a song.}
\glossentry{\textbf{la canzone}\index{canzone@canzone}}{feminine; plural \textit{le canzoni}}{song}{La nonna canta una canzone.}{Grandmother sings a song.}
\glossentry{\textbf{la musica}\index{musica@musica}}{feminine; plural \textit{le musiche}}{music}{Ascolto la musica.}{I listen to music.}
\glossentry{\textbf{la parola}\index{parola@parola}}{feminine; plural \textit{le parole}}{word}{Questa parola è nuova.}{This word is new.}
\glossentry{\textbf{la storia}\index{storia@storia}}{feminine; plural \textit{le storie}}{story; history}{Il nonno racconta una storia.}{Grandfather tells a story.}
\glossentry{\textbf{il valzer}\index{valzer@valzer}}{masculine; plural \textit{i valzer}}{waltz}{Balliamo un valzer.}{Let's dance a waltz.}
\glossentry{\textbf{la voce}\index{voce@voce}}{feminine; plural \textit{le voci}}{voice}{La ragazza ha una bella voce.}{The girl has a beautiful voice.}
\end{glosslist}

\section{Time}
\begin{glosslist}
\glossentry{\textbf{adesso}\index{adesso@adesso}}{adverb}{now}{Adesso mangiamo.}{Now we're eating.}\glossnote{ora has the same meaning.}
\glossentry{\textbf{ancora}\index{ancora@ancora}}{adverb}{still; again; yet}{Dormi ancora?}{Are you still asleep?}\glossnote{non … ancora = not yet}
\glossentry{\textbf{l'anno}\index{anno@anno}}{masculine; plural \textit{gli anni}}{year}{Luca ha sette anni.}{Luca is seven years old.}
\glossentry{\textbf{l'aprile}\index{aprile@aprile}}{masculine; no plural}{April}{Ad aprile piove spesso.}{It often rains in April.}
\glossentry{\textbf{l'autunno}\index{autunno@autunno}}{masculine; plural \textit{gli autunni}}{autumn}{In autunno le foglie cadono.}{In autumn the leaves fall.}
\glossentry{\textbf{domani}\index{domani@domani}}{adverb}{tomorrow}{Domani andiamo al mare.}{Tomorrow we are going to the seaside.}
\glossentry{\textbf{l'estate}\index{estate@estate}}{feminine; plural \textit{le estati}}{summer}{In estate andiamo al mare.}{In summer we go to the seaside.}
\glossentry{\textbf{l'età}\index{eta@età}}{feminine; plural \textit{le età}}{age}{Abbiamo la stessa età.}{We are the same age.}
\glossentry{\textbf{già}\index{gia@già}}{adverb}{already}{Ho già mangiato.}{I have already eaten.}
\glossentry{\textbf{la giornata}\index{giornata@giornata}}{feminine; plural \textit{le giornate}}{day (its duration, the daytime)}{Oggi è una bella giornata.}{Today is a lovely day.}
\glossentry{\textbf{il giorno}\index{giorno@giorno}}{masculine; plural \textit{i giorni}}{day}{Oggi è un bel giorno.}{Today is a lovely day.}
\glossentry{\textbf{il giugno}\index{giugno@giugno}}{masculine; no plural}{June}{A giugno finisce la scuola.}{School ends in June.}
\glossentry{\textbf{ieri}\index{ieri@ieri}}{adverb}{yesterday}{Ieri ho visto un gufo.}{Yesterday I saw an owl.}
\glossentry{\textbf{l'inverno}\index{inverno@inverno}}{masculine; plural \textit{gli inverni}}{winter}{In inverno fa freddo.}{In winter it is cold.}
\glossentry{\textbf{mai}\index{mai@mai}}{adverb}{never; ever}{Non piango mai.}{I never cry.}\glossnote{With a verb it needs non: non … mai.}
\glossentry{\textbf{la mattina}\index{mattina@mattina}}{feminine; plural \textit{le mattine}}{morning}{La mattina bevo il latte.}{In the morning I drink milk.}
\glossentry{\textbf{il mese}\index{mese@mese}}{masculine; plural \textit{i mesi}}{month}{Dicembre è un mese freddo.}{December is a cold month.}
\glossentry{\textbf{il mezzo}\index{mezzo@mezzo}}{masculine; plural \textit{i mezzi}}{middle; means}{C'è un albero in mezzo al prato.}{There is a tree in the middle of the meadow.}
\glossentry{\textbf{la notte}\index{notte@notte}}{feminine; plural \textit{le notti}}{night}{La notte è buia.}{The night is dark.}
\glossentry{\textbf{oggi}\index{oggi@oggi}}{adverb}{today}{Oggi piove.}{Today it is raining.}
\glossentry{\textbf{l'ora}\index{ora@ora}}{feminine; plural \textit{le ore}}{hour}{Il viaggio dura un'ora.}{The journey lasts an hour.}
\glossentry{\textbf{presto}\index{presto@presto}}{adverb}{early; soon; quickly}{Mi alzo presto.}{I get up early.}
\glossentry{\textbf{la primavera}\index{primavera@primavera}}{feminine; plural \textit{le primavere}}{spring}{In primavera ci sono molti fiori.}{In spring there are many flowers.}
\glossentry{\textbf{sempre}\index{sempre@sempre}}{adverb}{always}{Il gatto dorme sempre.}{The cat always sleeps.}
\glossentry{\textbf{la sera}\index{sera@sera}}{feminine; plural \textit{le sere}}{evening}{La sera leggo un libro.}{In the evening I read a book.}
\glossentry{\textbf{la settimana}\index{settimana@settimana}}{feminine; plural \textit{le settimane}}{week}{Una settimana ha sette giorni.}{A week has seven days.}
\glossentry{\textbf{subito}\index{subito@subito}}{adverb}{immediately; at once}{Vengo subito.}{I'm coming right away.}
\glossentry{\textbf{tardi}\index{tardi@tardi}}{adverb}{late}{È tardi!}{It's late!}
\glossentry{\textbf{la vita}\index{vita@vita}}{feminine; plural \textit{le vite}}{life}{La vita al mare è bella.}{Life by the sea is lovely.}
\glossentry{\textbf{la volta}\index{volta@volta}}{feminine; plural \textit{le volte}}{time (occasion)}{C'era una volta un orso.}{Once upon a time there was a bear.}
\end{glosslist}

\section{Weather}
\begin{glosslist}
\glossentry{\textbf{bagnato}\index{bagnato@bagnato}}{adjective: \textit{bagnato, bagnata, bagnati, bagnate}}{wet}{Il cane è bagnato.}{The dog is wet.}
\glossentry{\textbf{la neve}\index{neve@neve}}{feminine; plural \textit{le nevi}}{snow}{La neve è bianca.}{Snow is white.}
\glossentry{\textbf{nevicare}\index{nevicare@nevicare}}{verb, -are, auxiliary \textit{essere}, participle \textit{nevicato}}{to snow}{In montagna nevica.}{It is snowing in the mountains.}\glossnote{Constructions: impersonal: used without a subject. Weather verb: nevica = it is snowing. Both è nevicato and ha nevicato are correct.}
\glossentry{\textbf{la nuvola}\index{nuvola@nuvola}}{feminine; plural \textit{le nuvole}}{cloud}{Una nuvola nasconde il sole.}{A cloud hides the sun.}
\glossentry{\textbf{la pioggia}\index{pioggia@pioggia}}{feminine; plural \textit{le piogge}}{rain}{La pioggia cade piano.}{The rain falls gently.}
\glossentry{\textbf{piovere}\index{piovere@piovere}}{verb, -ere, auxiliary \textit{essere}, participle \textit{piovuto}}{to rain}{Oggi piove.}{It is raining today.}\glossnote{Constructions: impersonal: used without a subject. Weather verb used without a subject: piove = it is raining. Both è piovuto and ha piovuto are correct. Regular in form; listed here because only the third person is normally used.}
\glossentry{\textbf{la pozzanghera}\index{pozzanghera@pozzanghera}}{feminine; plural \textit{le pozzanghere}}{puddle}{Il bambino salta nella pozzanghera.}{The child jumps in the puddle.}
\glossentry{\textbf{soffiare}\index{soffiare@soffiare}}{verb, -are, auxiliary \textit{avere}, participle \textit{soffiato}}{to blow}{Il vento soffia forte.}{The wind is blowing hard.}
\glossentry{\textbf{la tempesta}\index{tempesta@tempesta}}{feminine; plural \textit{le tempeste}}{storm}{Stanotte c'è stata una tempesta.}{There was a storm last night.}
\glossentry{\textbf{il tempo}\index{tempo@tempo}}{masculine; plural \textit{i tempi}}{weather; time}{Che tempo fa oggi?}{What is the weather like today?}
\glossentry{\textbf{il temporale}\index{temporale@temporale}}{masculine; plural \textit{i temporali}}{thunderstorm}{Stanotte c'è un temporale.}{There is a thunderstorm tonight.}
\glossentry{\textbf{il vento}\index{vento@vento}}{masculine; plural \textit{i venti}}{wind}{Il vento muove le foglie.}{The wind moves the leaves.}
\end{glosslist}

\chapter{Sources and verification}\label{sources}
\begin{notebox}{exccol}\notelabel{exccol}{Not reviewed by a native speaker}{}No native speaker, Italian teacher or linguist has reviewed this course. Statements marked 'reference-checked' were compared with the named references by the author of the content; everything else is the author's own work based on standard grammar and should be treated as unreviewed. Corrections are welcome.\end{notebox}
\section{Verification register}
Status values: \emph{checked against references} --- the specific claims listed below were compared with the named reference; \emph{forms tested automatically} --- the forms are generated by code and compared with hand-written expected forms in the website's test suite; \emph{author-written} --- written from standard grammar but not compared line by line with a cited source.\par

{\small\begin{xltabular}{\linewidth}{@{}p{3.6cm} p{2.6cm} L@{}}
\toprule
\textbf{Topic} & \textbf{Status} & \textbf{Lessons and open points} \\
\midrule
\endhead
Pronunciation and spelling & Author-written, not checked against a reference & \S\ref{les:alphabet-vowels}, \ref{les:c-and-g}, \ref{les:stress-doubles-accents}. Seven-vowel system and regional variation stated; individual open/closed vowel claims should be checked against DiPI. \\
Nouns, articles, adjectives, prepositions & Forms tested automatically & \S\ref{les:noun-gender-number}, \ref{les:definite-articles}, \ref{les:indefinite-articles}, \ref{les:irregular-nouns}, \ref{les:adjective-agreement}, \ref{les:demonstratives-possessives}, \ref{les:bello-buono-comparison}, \ref{les:simple-prepositions}, \ref{les:articulated-prepositions}. Article selection, plurals and adjective agreement are generated by code and covered by unit tests; the prose explanations are author-written. \\
Present tense and irregular verbs & Forms tested automatically & \S\ref{les:verb-system-overview}, \ref{les:infinitives-pronouns}, \ref{les:present-are}, \ref{les:present-ere-ire}, \ref{les:essere-avere}, \ref{les:andare-fare-stare-dire}, \ref{les:other-irregulars}, \ref{les:modal-verbs}, \ref{les:reflexive-verbs}. Every vocabulary verb is conjugated by the engine in every tense; irregular paradigms are stored explicitly and spot-checked in tests. \\
Sentence building: word order, negation, questions, conjunctions & Author-written, not checked against a reference & \S\ref{les:word-order-negation-questions}, \ref{les:question-words-conjunctions}. Word-order statements describe neutral order; marked orders (dislocation, focus) are only touched on. \\
Passato prossimo, auxiliaries and participle agreement & Checked against references & \S\ref{les:passato-prossimo-avere}, \ref{les:passato-prossimo-essere}, \ref{les:ausiliari-accordo}. Agreement with ne ('normally agrees, invariable also found') is not covered by the cited Treccani article and is still unverified. Weather verbs with either auxiliary and ho vissuto / sono vissuto are author-written. References: Participio; La scelta degli ausiliari (consulenza linguistica); Verbi servili (entry on the auxiliary of modal verbs); Passato prossimo. \\
Imperfetto and trapassato prossimo & Author-written, not checked against a reference & \S\ref{les:imperfetto}, \ref{les:imperfetto-passato}, \ref{les:trapassato-prossimo}. Colloquial counterfactual imperfetto checked against treccani-periodo-ipotetico; the rest is author-written. \\
Pronouns, ne, combined pronouns, relative clauses & Author-written, not checked against a reference & \S\ref{les:direct-object-pronouns}, \ref{les:indirect-object-pronouns}, \ref{les:relative-clauses-adverbs}. Combined-pronoun table is standard; usage notes on il quale should be reviewed. \\
Imperativo & Forms tested automatically & \S\ref{les:imperativo}. Forms generated and tested (tested cases include va'/vai, di', essere, avere, reflexive and negative forms; formal forms come from the congiuntivo). \\
Futuro and condizionale & Forms tested automatically & \S\ref{les:futuro-semplice}, \ref{les:futuro-anteriore}, \ref{les:condizionale-presente}, \ref{les:condizionale-passato}. Forms tested. Uses (epistemic future, reportative conditional, future in the past) are author-written. 'sperare che + futuro' is mentioned as an alternative and has not been checked against a reference. \\
Infinito, gerundio, participio & Forms tested automatically & \S\ref{les:infinito-gerundio}, \ref{les:participio}. Forms tested; participio assoluto description checked against treccani-participio. \\
Passato remoto and trapassato remoto & Checked against references & \S\ref{les:passato-remoto}, \ref{les:trapassato-remoto}. Regional statements are kept general on purpose. Forms are engine-generated; strong-verb list and -ei/-etti alternatives are tested for the verbs in the vocabulary. Less common alternatives (e.g. for verbs not in the vocabulary) are not covered. References: Passato remoto; Passato prossimo. \\
Congiuntivo & Checked against references & \S\ref{les:congiuntivo-presente}, \ref{les:congiuntivo-passato}, \ref{les:congiuntivo-imperfetto-trapassato}. Register statements and relative-clause mood choice checked. Alternatives such as debba/deva and the full lists of verbs and conjunctions that require the congiuntivo are author-written. References: Congiuntivo. \\
Periodo ipotetico, concordanza dei tempi, discorso indiretto & Checked against references & \S\ref{les:periodo-ipotetico}, \ref{les:concordanza-dei-tempi}, \ref{les:discorso-indiretto}. Hypothetical types checked. Sequence-of-tenses tables and indirect-speech transformations are author-written and should be reviewed. References: Periodo ipotetico; Congiuntivo. \\
Passive and impersonal si & Author-written, not checked against a reference & \S\ref{les:forma-passiva}, \ref{les:si-impersonale}. venire/andare passives, si passivante agreement and plural agreement after impersonal si (si è stanchi) should be reviewed. \\
Stories and rhymes & Author-written, not checked against a reference & All texts are original and every word is checked automatically against the dictionary. Naturalness of the Italian has not been reviewed by a native speaker. \\
Sentence analyser & Forms tested automatically & Heuristic: it lists every dictionary analysis of a form and chooses one by context; it refuses sentences with unknown words. Its choices are tested on a fixed set of sentences only, and it can be wrong on others. \\
\bottomrule
\end{xltabular}}
\section{References and what was checked against them}
\begin{itemize}
\item\relax \textbf{Passato remoto} --- Eva Wiberg, \textit{Enciclopedia dell'Italiano, Treccani (2011)}. \url{https://www.treccani.it/enciclopedia/passato-remoto_(Enciclopedia-dell'Italiano)/}\newline -ei/-etti alternatives (vendei o vendetti); narrative (foreground) function; regional Italian largely mirrors dialect usage; strong preference for the passato remoto in Sicilian, Calabrian and Salentine dialect areas.
\item\relax \textbf{Passato prossimo} --- Eva Wiberg, \textit{Enciclopedia dell'Italiano, Treccani (2011)}. \url{https://www.treccani.it/enciclopedia/passato-prossimo_(Enciclopedia-dell'Italiano)/}\newline Aoristic passato prossimo spreading from the North, except in Tuscany; in the South the passato remoto is used where the standard would have the passato prossimo.
\item\relax \textbf{Participio} --- Verner Egerland, \textit{Enciclopedia dell'Italiano, Treccani (2011)}. \url{https://www.treccani.it/enciclopedia/participio_(Enciclopedia-dell'Italiano)/}\newline Agreement with a preceding direct-object clitic, especially 3rd person; less regular with 1st/2nd person (ci hanno visti / visto); agreement with relative che no longer standard (literary, antiquated, dialectal); participio assoluto. The article does not discuss ne.
\item\relax \textbf{Periodo ipotetico} --- Michele Prandi, \textit{Enciclopedia dell'Italiano, Treccani (2011)}. \url{https://www.treccani.it/enciclopedia/periodo-ipotetico_(Enciclopedia-dell'Italiano)/}\newline Three standard types (reality, possibility, unreality); substandard double imperfetto, attested in Dante, Machiavelli, Manzoni; conditional in the protasis non-standard ('doppio condizionale').
\item\relax \textbf{Congiuntivo} --- Michele Prandi, \textit{Enciclopedia dell'Italiano, Treccani (2010)}. \url{https://www.treccani.it/enciclopedia/congiuntivo_(Enciclopedia-dell'Italiano)/}\newline Indicativo for congiuntivo in completive clauses = lower register, same meaning; in restrictive relatives the mood choice changes the meaning.
\item\relax \textbf{La scelta degli ausiliari (consulenza linguistica)} --- Accademia della Crusca, \textit{Accademia della Crusca, Consulenza linguistica}. \url{https://accademiadellacrusca.it/it/consulenza}\newline Modal + essere-verb: both sono dovuto andare and ho dovuto andare acceptable; avere with essere or a passive infinitive; clitic before the modal → essere (non ci sono voluti andare). The link points to the consultation index; search it for the title.
\item\relax \textbf{Verbi servili (entry on the auxiliary of modal verbs)} --- DICO — Dubbi sull'italiano consulenza online, Università di Messina, \textit{online usage service}.\newline Modal takes the auxiliary of the infinitive; reports Serianni's view that avere is required when the infinitive is essere. Serianni's Grammatica italiana was not consulted directly.
\end{itemize}
\section{Further references}
\begin{itemize}
\item\relax \textbf{Treccani — Vocabolario online}: \url{https://www.treccani.it/vocabolario/}\newline Authoritative Italian dictionary (Italian only).
\item\relax \textbf{Wiktionary — Italian}: \url{https://en.wiktionary.org/wiki/Category:Italian_language}\newline Free dictionary with conjugation tables and pronunciation (CC BY-SA).
\item\relax \textbf{DiPI — Dizionario di pronuncia italiana}: \url{https://www.dipionline.it/}\newline Pronunciation and stress of Italian words.
\end{itemize}
\section{Provenance and licences}
\begin{itemize}
\item\relax \textbf{Italian Learning Lab original educational content} (content): All lessons, explanations, example sentences, exercises, vocabulary entries, stories and rhymes bundled with the site. Licence: CC BY 4.0. Author: Italian Learning Lab contributors. Written for this project. Grammar descriptions follow standard Italian reference grammar; no text was copied from third-party textbooks, dictionaries or songs.
\item\relax \textbf{Italian Learning Lab source code} (code): HTML, CSS and JavaScript in this repository. Licence: MIT. Author: Italian Learning Lab contributors. 
\item\relax \textbf{Lucide icons} (dependency (bundled at build time)): Interface icons. Licence: ISC (icons derived from Feather: MIT). Author: Lucide Icons and Contributors; Cole Bemis (Feather). 
\item\relax \textbf{Vite} (development tool (not shipped)): Development server and production build. Licence: MIT. Author: Evan You and Vite contributors. 
\item\relax \textbf{Vitest and jsdom} (development tools (not shipped)): Automated tests. Licence: MIT. Author: Vitest contributors; jsdom contributors. 
\item\relax \textbf{Browser speech synthesis (Web Speech API)} (optional browser feature): Read-aloud buttons. Uses voices installed in your browser or operating system; nothing is downloaded by this site. Some browsers may route synthesis through their vendor's online service — check your browser's documentation if this matters to you. Licence: Provided by your browser.  
\end{itemize}
No commercial song lyrics are bundled. The Songs workspace lets you paste lyrics that you have the right to use; they are saved only in your own browser and never uploaded. Audio files you open are played locally and are not stored or transmitted.\par

\backmatter
\printindex
\end{document}
